% 本文件由 LaTeX 排版 Agent 根据有效 translation.json 自动生成。
% author=Augustine of Hippo; work_id=1603; title=新约讲道集

\chapter{\Emph{新约讲道集}}

% structure: p0001 source=h1 target=omitted reason=page-title-already-rendered-by-parent

\Emph{讲道1}\newline{}\Emph{讲道2}\newline{}\Emph{讲道3}\newline{}\Emph{讲道4}\newline{}\Emph{讲道5}\newline{}\Emph{讲道6}\newline{}\Emph{讲道7}\newline{}\Emph{讲道8}\newline{}\Emph{讲道9}\newline{}\Emph{讲道10}\newline{}\Emph{讲道11}\newline{}\Emph{讲道12}\newline{}\Emph{讲道13}\newline{}\Emph{讲道14}\newline{}\Emph{讲道15}\newline{}\Emph{讲道16}\newline{}\Emph{讲道17}\newline{}\Emph{讲道18}\newline{}\Emph{讲道19}\newline{}\Emph{讲道20}\newline{}\Emph{讲道21}\newline{}\Emph{讲道22}\newline{}\Emph{讲道23}\newline{}\Emph{讲道24}\newline{}\Emph{讲道25}\newline{}\Emph{讲道26}\newline{}\Emph{讲道27}\newline{}\Emph{讲道28}\newline{}\Emph{讲道29}\newline{}\Emph{讲道30}\newline{}\Emph{讲道31}\newline{}\Emph{讲道32}\newline{}\Emph{讲道33}\newline{}\Emph{讲道34}\newline{}\Emph{讲道35}\newline{}\Emph{讲道36}\newline{}\Emph{讲道37}\newline{}\Emph{讲道38}\newline{}\Emph{讲道39}\newline{}\Emph{讲道40}\newline{}\Emph{讲道41}\newline{}\Emph{讲道42}\newline{}\Emph{讲道43}\newline{}\Emph{讲道44}\newline{}\Emph{讲道45}\newline{}\Emph{讲道46}\newline{}\Emph{讲道47}\newline{}\Emph{讲道48}\newline{}\Emph{讲道49}\newline{}\Emph{讲道50}\newline{}\Emph{讲道51}\newline{}\Emph{讲道52}\newline{}\Emph{讲道53}\newline{}\Emph{讲道54}\newline{}\Emph{讲道55}\newline{}\Emph{讲道56}\newline{}\Emph{讲道57}\newline{}\Emph{讲道58}\newline{}\Emph{讲道59}\newline{}\Emph{讲道60}\newline{}\Emph{讲道61}\newline{}\Emph{讲道62}\newline{}\Emph{讲道63}\newline{}\Emph{讲道64}\newline{}\Emph{讲道65}\newline{}\Emph{讲道66}\newline{}\Emph{讲道67}\newline{}\Emph{讲道68}\newline{}\Emph{讲道69}\newline{}\Emph{讲道70}\newline{}\Emph{讲道71}\newline{}\Emph{讲道72}\newline{}\Emph{讲道73}\newline{}\Emph{讲道74}\newline{}\Emph{讲道75}\newline{}\Emph{讲道76}\newline{}\Emph{讲道77}\newline{}\Emph{讲道78}\newline{}\Emph{讲道79}\newline{}\Emph{讲道80}\newline{}\Emph{讲道81}\newline{}\Emph{讲道82}\newline{}\Emph{讲道83}\newline{}\Emph{讲道84}\newline{}\Emph{讲道85}\newline{}\Emph{讲道86}\newline{}\Emph{讲道87}\newline{}\Emph{讲道88}\newline{}\Emph{讲道89}\newline{}\Emph{讲道90}\newline{}\Emph{讲道91}\newline{}\Emph{讲道92}\newline{}\Emph{讲道93}\newline{}\Emph{讲道94}\newline{}\Emph{讲道95}\newline{}\Emph{讲道96}\newline{}\Emph{讲道97}\par

\SourceRecord{Translated by R.G. MacMullen.；Nicene and Post-Nicene Fathers, First Series；Vol. 6.；Edited by Philip Schaff.；Buffalo, NY: Christian Literature Publishing Co.,；1888.；<http://www.newadvent.org/fathers/1603.htm>.}

% structure: p0001 source=h1 target=section reason=page-title

\section{新约讲道集 第一讲}

[LI. 本笃会版]\par

\Emph{论圣史玛窦与路加所记主世代的一致性}\par

1. 愿那激起你们期望的祂，亲爱的弟兄们，满足你们的期望：因为虽然我确信我所说的并非我的，而是天主的，但我更理所当然地这样说：宗徒在谦卑中说：\BibleQuote{我们在瓦器中存有这宝贝，为彰显那卓越的能力是属于天主，并非出于我们。} 因此，我毫不怀疑你们记得我的承诺；在祂内我作出了这承诺，通过祂我现在实现了它，因为当我作出承诺时，我曾向主祈求，现在当我实现它时，我从祂领受。现在你们会记得，亲爱的，那是在主圣诞节的晨祷中，我推迟了我提出要解决的问题，因为有许多人与我们一同去参加那一天的惯常庆典，而天主圣言对他们通常是一种负担；但现在我想象，来到这里的没有别人，只有那些渴望聆听的人，所以我并不是在对聋耳和蔑视话语的心灵说话，而是你们这热切的期待就是为我祈祷。还有一个考虑；因为公共表演的日子使许多人从这里散去，为了他们的得救，我恳请你们分担我极大的焦虑，并请你们以全部的热切之心，向天主祈求那些尚未专注于真理的景观、而是完全投入于肉身的景观的人；因为我知道并确信，你们中有些人今天藐视了这些表演，挣脱了他们根深蒂固的习惯的束缚；因为人会变得更好或更坏。通过这些日常的例子，我们时而欢喜，时而悲伤；我们为悔改者欢喜，为败坏者悲伤；因此主并没有说开始的人将得救，而是说：\BibleQuote{惟有坚持到底的，才可得救。}\par

2. 现在有什么比我们的主耶稣基督、天主子同时也是人子（因为祂也屈尊成为人子）赐予我们更奇妙、更伟大的事呢？祂不仅将那些愚蠢表演的观众，甚至将一些表演者本人聚集到祂的羊圈中；因为祂不仅与那些热爱人与野兽搏斗的人作战直到得救，也与那些搏斗者本人作战，因为祂自己也成了一个景观。请听如何。祂亲自告诉我们，并在祂成为景观之前预言了它，以预言的话语预先宣布了将要发生的事，就像已经发生了一样，在圣咏中说：\BibleQuote{他们刺穿了我的手和我的脚，他们数说了我全部的骨头。} 看！祂是如何成为一个景观，以至于祂的骨头被数说！而这个景观祂更清楚地表达：\BibleQuote{他们观察并注视着我。} 祂成了一个景观和嘲弄的对象，被那些在景观中并不对祂施恩、反而对祂发怒的人所观看，正如起初祂使祂的殉道者们成为景观一样；正如宗徒所说：\BibleQuote{我们成了一台戏，给世界、天使和世人观看。} 现在有两种人观看这样的景观：一种是属血肉的人，另一种是属神的人。属血肉的人观看，认为那些被扔给野兽、或斩首、或投入烈火的殉道者是可怜的人，他们憎恶和厌恶他们；但另一些人观看，如同圣天使一样，不关注他们身体的撕裂，而是钦佩他们信德的完好纯洁。在破碎的身体中，一个完整的内心向心中的眼睛展示出何等壮观的景象！当这些事在教会中被宣读时，你们用这些心中的眼睛快乐地观看它们，因为如果你们什么也没看见，你们就什么也听不见；所以你们看见，你们今天并没有忽视表演，而是做出了选择。愿天主与你们同在，并赐给你们恩宠，以温柔的劝说，向你们的朋友报告你们的表演，你们今天看到他们跑去竞技场而不愿来教会，这使你们痛苦；以便他们也开始藐视那些因爱它们而使他们自己变得可鄙的事物，并与你们一起爱天主，凡爱祂的人永远不会羞愧，因为他们爱的是那不能被战胜的一位：让他们像你们一样爱基督，祂正是借着那看似被战胜的同一件事，战胜了全世界。因为祂已经战胜了全世界，正如我们所见，我的弟兄们；祂征服了所有的权势，祂征服了君王，不是凭借军队的骄傲，而是借着十字架的耻辱；不是靠剑的愤怒，而是靠悬在木架上，在身体内受苦，在圣神内工作。祂的身体被高举在十字架上，由此祂将灵魂征服给十字架；现在，在他们王冠上的宝石中，有什么比君王前额上的基督十字架更珍贵的呢？爱祂，你们永远不会羞愧。然而，从竞技场回来的人，有多少是被征服的，因为他们所疯狂支持的人被征服了！如果他们获胜，他们更会被征服。因为他们会沦为虚妄的快乐和堕落欲望的狂喜的奴隶，而他们正是因跑去看这些表演而被征服的。因为有多少人，我的弟兄们，你们认为今天曾在犹豫是来这里还是去那里？那些在犹豫中，把心思转向基督、跑到教会来的人，已经战胜了，不是任何一个人，而是魔鬼本身，那追逐全世界灵魂的魔鬼。但那些在犹豫中，选择跑去竞技场的人，肯定被那另一些人战胜的魔鬼所战胜——那些人在祂内得胜，祂说：\BibleQuote{放心，我已经战胜了世界。} 因为这位统帅允许自己受试探，只是为了教导祂的士兵作战。\par

3. 为成就此事，我们的主耶稣基督由女人所生，成为人子。但有人会说：\Quote{倘若祂不由童贞玛利亚所生，难道祂就不算是完整的人吗？}\Quote{祂愿意成为人，这很好；祂本可以如此，却不由女人所生，因为祂最初造的第一个人也非从女人所出。}现在，且看我如何回应。你说：祂为何选择由女人所生？我答：祂为何要避免由女人所生？即便我无法证明祂选择由女人所生，你也该指出祂必须避免的理由。我已在别处说过，倘若祂避开了女人的子宫，似乎暗示祂可能从她那里沾染污秽；然而，祂的本性越是不可玷污，就越不必担心女人的子宫会玷污祂。但祂由女人所生，旨在向我们启示某种高深的奥义。的确，弟兄们，我们也承认，若主愿意成为人而不由女人所生，这对祂至高无上的尊威而言轻而易举。因为正如祂能无须男人而由女人所生，祂也能不由女人所生。祂这样向我们表明，是为了使人类中任何性别都不至绝望于自己的得救，因为人的性别有男有女。因此，如果祂作为人——这确实是祂应当成为的——却不由女人所生，女人可能会因想起第一个罪而绝望，因为第一个男人正是被女人所欺骗，她们会认为自己决无希望在基督内。所以祂以男人身份而来，特别拣选了那一性别；又由女人所生，以安慰女性，仿佛祂在对她们说：\Quote{为使你们知道，天主的任何受造物都不是坏的，而是无节制的快乐败坏了它。当初我造人时，造了男女。我不谴责我所造的受造物。看，我成了人，且由女人所生；因此我所谴责的不是我所造的受造物，而是我所未造的罪恶。}那么，让每一性别立刻看到自己的尊荣，承认自己的不义，并双双期望得救。毒害男人的毒药是女人递给他的；同样，让女人藉着生育基督，呈上拯救人类复原的救恩。为此，女人最先向宗徒们宣布了天主的复活。乐园中的女人向丈夫宣布了死亡，而教会中的女人向男人宣布了救恩；宗徒们应向万民宣布基督的复活，而女人先向宗徒们宣布了此事。因此，不要有人以基督由女人所生来责备祂；救主不会因这一性别而受玷污，造物主本来就定意要尊荣这一性别。\par

4. 但他们说：\Quote{我们怎能相信基督是由女人所生？}我要回答：凭着那已向普世宣讲、且仍在宣讲的福音。然而这些人，自己眼瞎，又企图使他人眼瞎，看不见所当见的，试图动摇所当信的，竟想在一件如今已遍及全地被相信的事上强加疑问。因为他们回应说：\Quote{不要以为你能用全世界的权威压倒我们——让我们查看圣经本身，不要用单纯人数的论点来反对我们，因为被迷惑的群众偏袒你。}对此，我首先回答：\Quote{被迷惑的群众偏袒我吗？}这群众曾是一小撮。这群众从何处增长，以致其壮大早已被预言？因为这如今被看见增长了的，正是那先前被预见到的。我不必说它是一小撮；起初只有亚巴郎。请想想，弟兄们：那时普世只有亚巴郎一人；在全世界、全人类、万民中，只有亚巴郎一人。他对他说过：\BibleQuote{在祢的后裔中，万民都要蒙受祝福。}而他独自一人凭己信德所信的，如今在众多后裔中显明给许多人。那时它未被看见，却被相信；如今它被看见了，却被争议；那当初对一人说、并由那人相信的，如今在许多人身上实现时，却受到少数人的争辩。那使门徒成为捕人的渔夫的那一位，将各种权威都收进了祂的网中。若相信多数，还有什么比教会更广泛地散布于普世？若相信富人，请看祂网住了多少富人；若相信穷人，请看成千上万的穷人；若相信贵族，几乎所有贵族都在教会内；若相信君王，请看他们全都臣服于基督；若相信更善辩、更有智慧、更有学问的人，请看有多少演说家、科学家和世间哲学家被那些渔夫捕获，从深渊里被拉上来得救；请想一想那一位，祂为医治人类灵魂那极大的恶——骄傲，以自身谦卑为榜样，降世为人，\BibleQuote{拣选了世上懦弱的，以羞辱强健的；拣选了世上愚拙的，以羞辱智慧的}（并非真正智慧，而是自以为智慧），\BibleQuote{又拣选了世上卑贱的和那无有的，为废除那有的。}\par

5. \Quote{任你说什么，} 他们说，\Quote{我们发现，在读到的基督诞生之处，福音彼此矛盾，而两件矛盾的事不可能都是真的；} 因为，有人说，\Quote{一旦我证明了这种矛盾，我就可以正当地拒绝相信它，或者，至少，你们这些接受相信它的人，请表明一致。} 而我要问，你要证明什么矛盾？\Quote{一个明显的矛盾，} 他说，\Quote{无人能否认。} 弟兄们，你们听到这些时有多么稳妥，因为你们是信徒！亲爱的，请留意，看宗徒给予何等有益的劝告，他说：\BibleQuote{你们既然接受了主基督耶稣，就当在他内行事，在他内生根建造，在信德上坚固；} 因为我们必须以这种简单而确定的信德在他内坚定不移，好使他亲自向信者开启隐藏在他内的事物；因为正如同一宗徒所说：\BibleQuote{在他内蕴藏着一切智慧和知识的宝藏；} 他隐藏它们并非为了拒绝，而是为了激发对隐藏事物的渴望。这就是隐藏的好处。那么，尊重在你内尚未理解的事，并且你看到的帷幕越多，就越要尊重：因为一个人的尊位越高，他的宫殿悬挂的帷幕就越多。帷幕使隐藏的东西受到尊敬，对于那些尊敬它的人，帷幕会被揭开；但那些嘲笑帷幕的人，甚至被赶离接近它们。因为我们 \BibleQuote{转向基督时，帷幕就被除去。}\par

6. 于是他们提出他们的吹毛求疵，说：\Quote{你承认玛窦是位福音作者。} 我们回答：是的，确实如此，带着虔诚的宣认和虔诚的心，毫无疑惑，我们明确回答：玛窦是位福音作者。\Quote{你相信他吗？} 他们说。谁不会回答：我相信？您那虔诚的低语传达了多么清晰的赞同！所以，弟兄们，你们完全相信它；你们没有理由为此羞愧。我是对你们说话，我曾经受骗，当我在童年早期，我选择将批评的微妙带到神圣的圣经，而不是一个寻求真理者的虔诚态度：由于我放荡的生活，我关闭了主的大门：当我本该敲门让它打开时，我却继续使它更紧密地关闭，因为我骄傲地寻找只有谦卑者才能发现的东西。现在你们是多么有福，怀着何等确定的信心学习，在何等安全中，你们仍是信德之巢中的幼雏，接受属神的食物；而我，可悲的我，自以为能飞，离开了巢，在飞之前就跌落：但仁慈的主扶起了我，免得我被路人踩死，并让我再次回到巢中；因为那些曾经困扰我的事情，现在我平静安稳地奉主的名向你们提出并解释。\par

7. 正如我开始所说，他们就这样吹毛求疵。\Quote{玛窦，}他们说，\Quote{是位福音作者，你相信他吗？} 一旦我们承认他是福音作者，我们就必然相信他。那么，请留意基督的族谱，玛窦所记载的。\BibleQuote{亚巴郎之子，达味之子耶稣基督的族谱。} 如何是达味之子，亚巴郎之子？除非通过世代传承，否则无法显示；因为当主由童贞玛利亚诞生时，亚巴郎和达味都不在世，你能说同一个人既是达味之子，又是亚巴郎之子吗？让我们对玛窦说：证明你的话，因为我等待基督的世代传承。\BibleQuote{亚巴郎生依撒格；依撒格生雅各伯；雅各伯生犹大和他的兄弟们；犹大由塔玛尔生培勒兹和则辣黑；培勒兹生赫兹龙；赫兹龙生兰；兰生阿米纳达布；阿米纳达布生纳赫雄；纳赫雄生撒尔孟；撒尔孟由辣哈布生波阿次；波阿次由卢德生敖贝得；敖贝得生叶瑟；叶瑟生达味王。} 现在看看从这点起，族谱如何从达味下传到基督，他被称为亚巴郎之子和达味之子。\BibleQuote{达味由乌黎雅的妻子生撒罗满；撒罗满生勒哈贝罕；勒哈贝罕生阿彼雅；阿彼雅生阿撒；阿撒生约沙法特；约沙法特生约兰；约兰生乌齐雅；乌齐雅生约堂；约堂生阿哈次；阿哈次生希则克雅；希则克雅生默纳舍；默纳舍生阿孟；阿孟生约史雅；约史雅生耶苛尼雅和他的兄弟们，约在巴比伦充军的时候；充军巴比伦以后，耶苛尼雅生沙耳提耳；沙耳提耳生则鲁巴贝耳；则鲁巴贝耳生阿彼乌得；阿彼乌得生厄里雅金；厄里雅金生阿佐尔；阿佐尔生匝多克；匝多克生阿歆；阿歆生厄里乌得；厄里乌得生厄肋阿匝尔；厄肋阿匝尔生玛堂；玛堂生雅各伯；雅各伯生若瑟，玛利亚的丈夫，由她生耶稣，称为基督。} 这样，通过父亲和祖先的次序和传承，基督被发现是达味之子，亚巴郎之子。\par

8. 现在，在如此忠实地叙述之后，他们提出的第一个刁难是，同一位玛窦接着说：\BibleQuote{从亚巴郎到达味共十四代；从达味直到流徙巴比伦共十四代；从流徙巴比伦直到基督共十四代。}然后，为了告诉我们基督如何由童贞玛利亚诞生，他继续说：\BibleQuote{耶稣基督的诞生是这样的：}因为藉着家谱的线索，他已表明了为何基督被称为达味之子、亚巴郎之子。但现在需要表明祂如何诞生并在人间显现：于是紧接着那段叙述，藉此我们相信我们的主耶稣基督不但生于永远的天主，与生祂者在一切时间、一切受造之前同永，万物是藉祂而造的；而且现在也由圣神、从童贞玛利亚诞生，我们承认这与前者同等；因为你们记得并知道（因为我是在对天主教徒、我的弟兄们说话），这就是我们的信德，这就是我们所宣认并承认的；为了这信德，千万殉道者在全世界被杀。\par

9. 他们接下来喜欢嘲笑这一点，他们想摧毁福音书的权威，以便似乎表明我们毫无理由地相信了所说的话：\BibleQuote{祂的母亲玛利亚已许配给若瑟，在他们同居以前，她因圣神有孕的事已显示出来。她的丈夫若瑟，因是义人，不愿公开羞辱她，有意暗暗地休退她。}因为他知道她不是由他怀孕，就认为她必须被视为一个奸妇。\BibleQuote{因是义人，}正如圣经所说，\BibleQuote{不愿公开羞辱她，}（即泄露此事，因为许多抄本如此），\BibleQuote{有意暗暗地休退她。}这位丈夫确实困扰，但作为义人，他并不严厉处理；因为如此大的正义被归给这人，以至于他既不愿保留一个奸妇，也不能狠心惩罚和暴露她。\BibleQuote{有意暗暗地休退她，}因为他不仅不愿惩罚，甚至不愿揭露她；请注意他真正的正义：他不想饶恕她，是因为他渴望保留她；因为许多人出于肉欲的爱饶恕他们的奸妇，宁愿保留她们，即使她们是奸妇，以便因肉欲的渴望享受她们。但这义人并不想保留她，所以他没有任何肉欲的爱；然而他也不愿惩罚她；所以出于怜悯，他饶恕了她。这人是多么真正正义啊！他既不愿保留一个奸妇，以免显得出于不洁的情感而饶恕她，也不愿惩罚或揭露她。他确实被拣选作他妻子贞洁的见证是正确的；因此，那个因人性软弱而困惑的人，得到了神圣权威的保证。\par

10. 因为福音作者接着说：\BibleQuote{当他思虑这事时，看，上主的天使在梦中显现给他，说：若瑟，达味之子，不要怕娶你的妻子玛利亚，因为那在她内受生的，是出于圣神。她要生一个儿子，你要给祂起名叫耶稣。}为什么叫耶稣？\BibleQuote{因为祂要把自己的民族从他们的罪恶中拯救出来。}众所周知，在希伯来语中，\OriginalTerm{耶稣}{Jesus}在拉丁语中被解释为\OriginalTerm{救世主}{Saviour}，我们从这名字的解释中看出来；因为好像有人问：\Quote{为什么叫耶稣？}他立即补充说明这词的理由：\BibleQuote{因为祂要把自己的民族从他们的罪恶中拯救出来。}因此，我们虔诚地相信，最坚定地持守：基督是由圣神、从童贞玛利亚诞生的。\par

11. 那么我们的对手说什么？\Quote{如果，}有人说，\Quote{我发现了一个谎言，你当然不会相信全部；而我发现了这样一个谎言。}让我们看看：我将列出世代；因为他们通过诽谤性的刁难邀请并促使我们这样做。是的，如果我们虔诚生活，如果我们相信基督，如果我们不想在时机未到之前飞出巢穴，他们只会促使我们——了解奥迹。注意，圣洁的弟兄们，异端者的用处；也就是说，在天主计划中的用处，祂甚至善用恶人；而就他们自身而言，他们自己谋划的果实归于他们，而不是天主从他们中带出的善。如同在犹达斯的情况中，他做了多么大的坏事！藉着主的苦难，万民得救；但为使主受难，犹达斯出卖了祂。天主既藉着祂圣子的苦难拯救万民，又因犹达斯自己的邪恶惩罚他。因为隐藏在圣经中的奥迹，没有谁满足于信德的单纯会好奇地筛分，因此无人会筛分，无人会发现它们，除非有刁难者迫使我们。因为当异端者刁难时，小孩子们被扰乱；当他们被扰乱时，他们寻求，他们的寻求就像在母亲乳房上叩头，以便为这些小孩子们流出足够的奶。于是他们寻求，因为他们受困扰；而那些知道并学到了这些事的人，因为他们已经研究过，并且天主在他们敲门时开启了，他们转而向受困扰的人开启。因此，异端者对于真理的发现是有用的，而他们刁难是为了引诱人进入错误。因为真理若不具有说谎的对手，就不会被更仔细地寻求；\BibleQuote{因为在你们中间也不免有异端，}好像我们要问原因，他立即补充说：\BibleQuote{好叫那些被核准的人，在你们中间显出来。}\par

12. 那么，他们说什么呢？看，玛窦列举了族谱，说：\Quote{从亚巴郎到达味共十四代，从达味到流徙巴比伦共十四代，从流徙巴比伦到基督共十四代。}\newline{} 如今三次十四是四十二；但他们一数，发现是四十一代，立刻便提出他们的吹毛求疵和侮辱性的嘲弄，说：\Quote{福音中说有三组十四代，但数在一起时，发现不是四十二，而是四十一，这是什么意思？}\newline{} 毫无疑问这里有一个伟大的奥秘：我们很高兴，我们感谢主，借着吹毛求疵者的机会，我们发现了某种东西，它在被发现时给我们带来更多的喜悦，其程度与它作为寻求对象时的晦涩成比例；因为正如我之前所说的，我们正在向你们的思想展示一个景观。那么，从亚巴郎到达味是十四代；之后，从撒罗满开始计数，因为达味生了撒罗满；我说，从撒罗满开始计数，直到耶苛尼雅，在他的有生之年发生了流徙巴比伦；这样就有了另外十四代，通过将撒罗满算在第二组的开头，也将耶苛尼雅算在内，该计数以他结束以凑足十四之数；第三组则从同一位耶苛尼雅开始。\par

13. 圣洁的弟兄们，请留意这个既神秘又令人愉快的情况；因为我向你们承认我内心的感受，我相信当我把它讲出来，你们品尝到它时，你们也会对它做出同样的评价。请留意。在第三组中，从这位耶苛尼雅直到主耶稣基督，共发现十四代；因为这位耶苛尼雅被算了两次，作为前一组的最后一个和下一组的第一个。\newline{} 但有人会说：\Quote{为什么耶苛尼雅被算了两次？}在以色列百姓中，古时所发生的一切，没有一样不是未来事物的神秘预象；事实上，耶苛尼雅被算两次并非没有充分的理由，因为如果两块田地之间有一个边界，无论是石头还是任何分界墙，那么处于一边的人会量到同一堵墙，而处于另一边的人则从那同一堵墙重新开始测量。\newline{} 但为什么在分组的第一连接点没有这样做呢？当我们从亚巴郎数到达味共十四代，然后开始数另外十四代时，不是从达味重新开始，而是从撒罗满开始，就必须给出一个包含重要奥秘的理由。那么，请留意。流徙巴比伦发生在耶苛尼雅被立为王以接替他已故的父亲的时候。王国从他那里被夺走，另一个人被立来代替他；但流徙到外邦人那里确实发生在耶苛尼雅有生之年，没有提到耶苛尼雅的过错致使他失去王位；倒是那些继任者的罪过被指出来。所以后来就有了被掳和流徙巴比伦；恶人并不是独自前往，圣人们也和他们同去：因为在那次被掳中，有先知厄则克耳和达尼尔，以及被投入火焰中而闻名的那三个孩子。他们都按照耶肋米亚先知的预言去了。\par

14. 那么你们要记住，耶苛尼雅被废弃，并非因他有何过错，他停止了统治，并在巴比伦充军发生时，转向了外邦人。现在观察这件事所预先昭示的、关于在吾主耶稣基督内将来之事的预像。因为犹太人不想我们的主耶稣基督统治他们，却在他身上找不到任何过错。他被弃绝，不仅在他自身，也在他的仆人们身上，所以他们预像性地转向外邦人，如同进入巴比伦。因为耶肋米亚也曾预言，主命令他们进入巴比伦；而其他劝诫百姓不要进入巴比伦的先知，他斥之为假先知。阅读\WorkTitle{圣经}的人，愿你们像我一样记住此事；不读的人，则请信任我们。于是耶肋米亚代表天主威胁那些不愿进入巴比伦的人，而对那些应当前往的人，他在那里许诺安息，并在培植葡萄园、栽种园囿以及丰收果实方面有一种幸福。那么以色列百姓，如今不再是预像而是真实地，是如何转向巴比伦的呢？宗徒们来自哪里？他们不也是犹太民族的吗？保禄本人来自哪里？因为他说：\BibleQuote{我亦是以色列人，亚巴郎的后裔，本雅明支派。} 于是许多犹太人信了主；宗徒们从他们中被选；超过五百位弟兄从他们中而来，他们蒙恩在主复活后看见主；一百二十位在楼上的人从他们中而来，当圣神降临时。但宗徒在\BibleBook{宗徒}{大事录}中，当犹太人拒绝真理之言时说了什么？\BibleQuote{我们本是奉派到你们这里来，但既然你们拒绝了天主的圣言，看，我们转向外邦人。} 因此，真正的转向巴比伦，原是耶肋米亚时代预像的，发生在主降生成人的精神恩宠时期。但是耶肋米亚对这些巴比伦人，对那些转向他们的人说了什么？\BibleQuote{因为他们的平安将成为你们的平安。} 当以色列通过基督和宗徒们也转向巴比伦，即福音传到外邦人时，宗徒如同借着古耶肋米亚的口说了什么？\BibleQuote{所以我劝勉，第一要为一切人恳求、祈祷、转求和谢恩，为君王和一切有权位的人，使我们能过安宁平和的生活，在一切虔敬和端庄中。} 因为他们那时还不是基督徒君王，然而他为他们祈祷。于是以色列在巴比伦祈祷已被垂听；教会的祈祷已被垂听，君王们成了基督徒，你们现在看到那当时在预像中所说的应验了：\BibleQuote{在他们的平安中你们的平安，} 因为他们已接受了基督的平安，并停止了迫害基督徒，以致如今在平安的安稳宁静中，教会得以建立，万民被栽种在天主的园囿中，并且所有民族能在信德、望德和爱德中结出果实，这一切是在基督内。\par

15. 古时通过耶苛尼雅发生了充军至巴比伦，他不被允许在犹太民族中统治，作为基督的预像，基督是犹太人不想让他统治他们的。以色列转向了外邦人，即福音的宣讲者转向了外邦人民。那么，耶苛尼雅被计算两次，有什么稀奇呢？因为如果他是基督从犹太人转向外邦人的预像，只需考虑基督在犹太人和外邦人之间是什么。难道他不是那角石吗？在一块角石上，你看到一堵墙的终点和另一堵墙的起点；量到那块石头，你量完一堵墙，又从它开始量另一堵墙；因此，连接两堵墙的角石被计算两次。那么耶苛尼雅，作为主的预像，可以说是角石的预像；正如耶苛尼雅不被允许统治犹太人，反而他们去了巴比伦，同样，基督，\BibleQuote{匠人弃而不用的石头，反而成了屋角的基石，} 以便福音能传到外邦人。那么不必迟疑将角石的头计算两次，你立即就得到所写的数目：如此在三种划分中各有十四代，然而所有世代不是四十二，而是四十一；因为正如当石头的排列沿直线运行时，它们都只被计算一次，但当从直线偏转成一个角时，那个偏转开始处的石头必须被计算两次，因为它既属于那条在它那里结束的线，又属于那条从它开始的线；同样，当世代的次序在犹太民族中延续时，它在十四代的常规划分中没有形成角；但当线转向，使百姓能转向巴比伦时，在耶苛尼雅那里形成了一个角，因此有必要将他计算两次，作为那可敬的角石的预像。\par

16. 他们还有另一种异议。\Quote{基督的世代，}他们说，\Quote{是通过若瑟计算的，而不是通过玛利亚。} 请稍听，神圣的弟兄们。\Quote{不应该这样，}他们说，\Quote{通过若瑟。} 为什么不应该呢？若瑟不是玛利亚的丈夫吗？他们说：\Quote{不是。} 谁这样说的？因为\WorkTitle{圣经}借天使的权威说他是她的丈夫：\BibleQuote{不要怕娶你的妻子玛利亚，因为那在她内受孕的是出于圣神。} 此外，他被命令给那孩子命名，虽不是他精血所生：\BibleQuote{她要生一个儿子，你要给他起名叫耶稣。} \WorkTitle{圣经}意在表明，他不是由若瑟的精血所生，当他在疑虑她怀孕时被告知：\BibleQuote{他是出于圣神；} 然而，他为人父的权柄并未被剥夺，因为他被命令给孩子命名；此外，童贞玛利亚自己也深知，她怀孕基督并非由于他，却仍称他为基督的父亲。\par

17. 请思考这事发生在何时。主耶稣就其人性而言，当时十二岁（因为就其天主性而言，祂在万有之前，且不受时间限制），祂留在圣殿中，与长老们辩论，他们惊奇于祂的教导；祂的父母从耶路撒冷返回，在同行的人中寻找祂，即那些与他们一起旅行的人，当他们找不到祂时，便忧心忡忡地回到耶路撒冷，在圣殿里找到祂与长老们辩论，那时祂正如我所说，十二岁。但这有什么奇怪呢？圣言从不沉默，尽管并非总被听见。于是祂在圣殿中被找到，祂的母亲对祂说：\BibleQuote{祢为什么这样对待我们？祢的父亲和我一直忧伤地寻找祢}；祂回答说：\BibleQuote{难道你们不知道我必须在我父的事务上服务吗？}祂这样说，是因为天主子在天主的圣殿里，因为那圣殿不是若瑟的，而是天主的。看，有人说：\Quote{祂不承认自己是若瑟的儿子。}弟兄们，请稍忍耐，因为时间紧迫，以便有足够时间讲完我要说的话。当玛利亚说过：\BibleQuote{祢的父亲和我一直忧伤地寻找祢}，祂回答说：\BibleQuote{难道你们不知道我必须在我父的事务上服务吗？}因为祂不愿意做他们的儿子到如此程度，以至于不被理解也是天主子。因为祂是天主子——永远是，也是创造那对祂说话者的那一位；但祂在时间内是人子；由童贞女所生，未经丈夫的作为，却又是双亲的儿子。我们从哪里证明这点呢？我们已经用玛利亚的话证明了：\BibleQuote{祢的父亲和我一直忧伤地寻找祢。}\par

18. 首先，为训导妇女，我们的姊妹，不可忽略童贞玛利亚如此圣洁的谦逊，弟兄们。她生了基督——天使曾来到她面前，说：\BibleQuote{看，你将在你的腹中怀孕，生一个儿子，并要给他起名叫耶稣。他将成为伟大的人物，并被称为至高者的儿子。}她被认为配得上生产至高者的儿子，然而她却极其谦卑；她甚至不把自己置于丈夫之前，即在称呼的顺序上，以至于不说\Quote{我和你父亲}，而是说\BibleQuote{祢的父亲和我}。她不看重自己腹中的崇高尊荣，而是看重婚姻的秩序，因为谦卑的基督不会教导祂的母亲骄傲。\BibleQuote{祢的父亲和我一直忧伤地寻找祢。}祢的父亲和我，她说，\BibleQuote{因为丈夫是女人的头。}那么其他妇女更应当多么谦卑！因为玛利亚自己也被称为女人，并非因失去童贞，而是因她本国的特殊表达方式；因为关于主耶稣，宗徒也曾说：\BibleQuote{由女人所生}，但这并不中断我们信经的次序和连贯，其中我们宣认\Quote{祂由圣神和童贞玛利亚所生}。因为她以童贞怀了祂，以童贞生了祂，并保持童贞；但他们把所有女性都称为\Quote{\OriginalTerm{女人}{woman}}，这是希伯来语的一个特点。听一个非常明显的例子：天主造的第一个女人，是从男人的肋骨中取出，被称为女人，在她\Quote{\OriginalTerm{认识}{knew}}丈夫之前，我们知道那是在他们离开乐园之后，因为圣经说：\BibleQuote{祂造了她成为一个女人。}\par

19. 主耶稣基督的回答，\BibleQuote{我必须在我父的事务上服务}，并不是以这样的方式宣告天主是祂的父，以至于否认若瑟也是祂的父亲；我们从哪里证明这点呢？由圣经，它如此说：\BibleQuote{祂对他们说：难道你们不知道我必须在我父的事务上服务吗？但他们不明白祂对他们所说的话：祂同他们下去，来到\Place{纳匝肋}，并服从他们。}它并没有说\Quote{祂服从祂的母亲}，或\Quote{服从她}，而是\BibleQuote{祂服从他们}。祂服从谁呢？难道不是祂的父母吗？祂服从祂的双亲，是由于相同的自谦，即祂是人子的自谦。不久前，妇女接受了训令。现在让儿女接受他们的训令——服从父母，并服从他们。世界服从基督，而基督服从祂的父母。\par

20. 那么，弟兄们，你们看，祂并没有说：\Quote{我必须从事我父的事务}，以此让我们理解为祂是在说：\Quote{你们不是我的父母。}从时间来说，他们是祂的父母；从永恒来说，天主是祂的父。他们是人子的父母——祂是\OriginalTerm{圣言}{Word}的父，也是\OriginalTerm{智慧}{Wisdom}与\OriginalTerm{能力}{Power}的父，祂藉此创造了万有。但如果万物是由那智慧创造的，这智慧\Quote{刚强地从地极达到地极，从容地治理万有}，那么万物也是由天主圣子所造，而祂自己后来作为人子要服从于这圣子；宗徒说祂是达味之子，\Quote{祂按肉身是从达味的后裔生的}。然而主自己向犹太人提出一个问题，宗徒正是用这些话解决了这个问题：因为当他说\Quote{祂按肉身是从达味的后裔生的}时，他补充说\Quote{按肉身}，以便让人明白：祂按天主性不是达味之子，而是天主圣子是达味的主；因为宗徒在另一处论及犹太民族的特权时说：\Quote{他们也是先祖，基督按肉身也是从他们来的，祂是万有之上的天主，永受赞美。}所以，按肉身，祂是达味之子；但作为\Quote{万有之上的天主，永受赞美}，祂是达味的主。于是主对犹太人说：\Quote{你们说基督是谁的儿子？}他们回答说：\Quote{达味之子。}因为他们知道这一点，并容易从先知的宣讲中学到；的确，祂是达味的后裔，但\Quote{按肉身}，经由与若瑟订婚的童贞玛利亚。当他们回答说基督是达味之子时，耶稣对他们说：\Quote{那么，达味怎么在圣神内称祂为主，说：\InnerQuote{主对我主说：你坐在我右边，等我把你的仇敌放在你脚下。}如果达味在圣神内称祂为主，祂怎么又是达味的儿子呢？}犹太人不能回答祂。我们在福音中读到这些。祂没有否认祂是达味之子，以便他们不能理解祂也是达味的主。因为他们只承认基督在时间中成为什么，却不理解祂从永恒中是什么。因此，为了教导他们祂的天主性，祂提出了一个关于祂人性的问题，仿佛在说：\Quote{你们知道基督是达味之子，回答我，祂怎么也是达味的主？}为了不让他们说\Quote{祂不是达味的主}，祂引用了达味自己的见证。那么达味说了什么呢？他确实说了真话。因为你在圣咏中看到天主对达味说：\Quote{我会将你身子的果子安置在你的宝座上。}所以这里祂是达味之子。但达味之子怎么是达味的主呢？\Quote{主对我主说：你坐在我右边。}你能奇怪达味之子是他的主吗？当你看到玛利亚是她儿子的母亲时，难道她的儿子不是她的主吗？那么，祂作为天主是达味的主，作为万有的主；作为人子，祂是达味之子。同时是主和子。达味的主，\Quote{祂虽具有天主的形体，却没有将与自己同等视为应当紧抓不放的}；达味之子，因为\Quote{祂空虚了自己，取了奴仆的形体。}\par

21. 那么，若瑟并没有因为不认识主的母亲而不再是祂的父亲，仿佛情欲而非夫妻之情构成了婚姻的纽带。请留意，圣洁的弟兄们；基督的宗徒后来要在教会中说：\Quote{那有妻子的，要像没有妻子一样。}我们知道许多弟兄藉着恩宠结出果实，为基督的名，他们通过彼此同意实行完全的节制，却仍然经受着真正夫妻之情的约束。是的，前者越受抑制，后者就越被加强和巩固。那么，这样生活的人难道不是已婚的人吗？他们不向对方要求任何肉体的满足，也不求任何身体欲望的满足。然而，妻子顺服丈夫，因为这是适宜的，而且她顺服得越多，她的贞洁就越大；丈夫也真心爱他的妻子，正如经上所记：\Quote{在尊贵和圣化中}，作为恩宠的共同继承人；正如宗徒所说：\Quote{基督爱了教会。}如果这是一种结合，一种婚姻；如果它仍然是婚姻，尽管他们之间没有发生那种即使在未婚者之间也可能发生、但不合法的事；（但愿所有的人都能够这样生活，但许多人没有这个能力！）那么，至少不要分离那些有能力这样做的人，也不要否认那人是丈夫或那人是妻子，因为没有肉体的交往，只有他们之间的心灵结合。\par

22. 所以，诸位弟兄，请理解圣经中关于我们那些古老祖先的含意：他们在婚姻中的唯一目的是通过妻子生育子女。因为即使那些按照他们时代和民族的习俗拥有多位妻子的人，也以如此贞洁的方式与她们生活，以至于若非为了我提到的那个原因，绝不接近她们的床榻，从而确实以尊敬对待她们。但是，若有人超越这规则为达成婚姻目的而设定的限度，便违反了娶妻时所立的契约本身。那份契约被宣读，在所有见证人面前宣读；其中明确写有他们结婚是\Quote{为了生育子女}；这被称为婚姻契约。如果不是为此目的，妻子被给予和娶回，那么哪个父亲能毫无羞愧地将自己的女儿交给任何男人的情欲呢？但现在，为了父母不羞愧，为了他们能体面地嫁女，而非蒙羞，契约得以宣读。其中读到了什么？——条款：\Quote{为了生育子女}。听到这，父母的眉头舒展、心境平复。让我们再思考娶妻的丈夫的感受。丈夫本人也会因任何其他目的而接受她感到羞愧，正如父亲会因其他任何目的而嫁女感到羞愧一样。尽管如此，如果他们不能自制（正如我在其他场合所说），就让他们索取应得的，但不要从其他人那里索取，而只能从应尽义务的一方那里索取。让男女双方都在自己身上寻求对软弱的缓解。不要丈夫去找别的女人，也不要女人去找别的男人，因为通奸因此得名，仿佛是\Quote{去往另一个}。如果他们超越了婚姻契约的界限，至少不要超越夫妻忠贞的界限。已婚者互相索取超出\Quote{生育子女}这一目的所需的，岂不是罪吗？这无疑是罪，尽管是小罪。宗徒说：\Quote{但这是我说的宽免的话，}他在论述此事时这样说道。\Quote{你们彼此不可亏负，除非两相情愿，暂时分房，为专务祈祷；以后再合，免得撒殚因你们不能自制而诱惑你们。}这是什么意思？就是你们不要给自己加添超乎能力的负担，不要因互相的节欲而陷入通奸。\Quote{免得撒殚因你们不能自制而诱惑你们。}为了不让人以为他在命令，而只是允许（因为向德行的力量发布诫命是一回事，向软弱做出宽免是另一回事），他随即补充说：\Quote{但这是我说的宽免的话，不是命令。因为我愿意众人都像我一样。}似乎他要说：我不是命令你们这样做，而是宽恕你们若这样做。\par

23. 所以，诸位弟兄，请留意。那些只为生育子女而娶妻的著名人物——正如我们通过许多证据，通过圣经清晰明确的见证而读到的圣祖们那样；无论他们是谁，仅为这一目的而娶妻，如果能让他们无需与妻子同房就能有子女，他们岂不会以难以言表的喜悦拥抱如此巨大的恩赐吗？他们岂不会以极大的欢欣接受它吗？因为有两种肉体的活动维系人类，贤哲圣洁的人出于义务而参与其中，但愚昧者却因情欲而一头扎进去；这两者截然不同。那么，维系人类的这两件事是什么？第一件关乎我们自身，与摄取食物有关（当然，若无肉体的些许满足就无法进食），即吃和喝；若不做，你就会死。因此，人类种族凭着这一维生之道——吃和喝——按照其本性的法则而生存。但由此，人类仅能维持自身；因为通过吃喝，他们并不提供后代的延续，而是通过娶妻来实现。因为人类种族是这样保存的：首先，凭生命之道；但由于无论他们如何努力，当然不能永远活着，于是有第二种预备：让新生者代替逝者。因为人类种族，如经上所记载，像树上的叶子，比如橄榄树、月桂树或此类树木，它们从不缺叶，但叶子却不总是一样的。正如经上所说：\Quote{它长出一些，又落下一些，}因为新发的嫩叶替代了掉落的旧叶，树不断落叶，却又始终披着叶子。同样，人类种族并不感觉天天逝去者的损失，因为新生者的供应；整个人类种族按其自身的法则得以维持，正如树上总是有叶子，大地也充满了人。而如果只有死亡，没有新生，大地将失去所有居民，如同某些树完全失去叶子。\par

24. 既然人类如此生存，以致那两种支撑——现已充分论述——对人类是必需的，那么明智、有悟性且有信德的人便出于本分降至于它们，并非因贪欲而堕落其中。然而，有多少人贪婪地冲向饮食，将其视为生活的全部，仿佛它们就是活着的理由？因为人本当为活而吃，他们却以为活着就是为了吃。每个明智的人都会谴责这些人，而圣经尤其谴责所有饕餮之徒、酗酒者、贪食者，\Quote{他们的神是自己的肚腹}。唯有肉体的贪欲，而非补充精力的需要，将他们引向餐桌。这些人于是扑向他们的饮食。然而，那些出于维持生命的本分而降至饮食的人，并非为吃而活，而是为活而吃。因此，若向这些明智而节制的人提出——他们可以无需饮食而生活——他们会何等欢欣地领受这一恩赐！如此，他们甚至不再被迫降至他们从未习惯堕落其中的事物，而是能永远在主内被举扬，且无需因修复身体损耗的需要而放下对祂恒定的专注。你们认为，圣\Person{厄里亚}是如何接受那壶水和那块饼，以维持他四十天的饱足？无疑，带着极大的喜乐，因为他吃喝是为活，而非服务于贪欲。但请尝试将此带给某人——假若你能——那人如同厩中的牲畜，将其全部福乐与幸福置于餐桌上。他会憎恨你的恩赐，将其推开，视之为惩罚。同样，在婚姻的另一义务中，好色之徒寻求妻子只是为了满足其情欲，因此最终连妻子也难满足。唉！但愿他们若不能或不愿治疗其情欲，至少不使其超越夫妻义务所规定的界限——我指的是甚至那为软弱而许可的界限。然而，若你对这样的人说：\Quote{你为何结婚？}他或许因羞愧而回答：\Quote{为了子女。}但若有任何一个他毫不怀疑其可信的人对他说：\Quote{天主能够，而且确实会赐给你子女，且无需你与妻子有任何交合；}他必定被迫承认，他寻求妻子并非为了子女。那么，让他承认自己的软弱，并如此接受他假借本分之名而接受的事物。\par

25. 古时的圣人们，那些天主的人，正是如此寻求并渴望子女。他们与妻子的交合与结合——唯一目的——就是生育子女。正因如此，他们被允许拥有多位妻子。因为若在这些欲望上的过度能取悦于天主，那么当时一个女子有多个丈夫就会像一个男子有多个妻子一样被允许。为何所有贞洁的女子只拥有一个丈夫，而一个男子却可拥有多个妻子？除非是因为一个男子有多个妻子有助于家族的繁衍，而一个女子即使有再多丈夫也不会生育更多子女。因此，弟兄们，如果我们的祖先与妻子的结合与交合，除了生育子女别无他目的，那么若他们能在无此交合的情况下获得子女，这对他们将是极大的喜乐，因为为了拥有子女，他们只是出于本分降至于交合，而非因贪欲而冲入其中。那么，\Person{若瑟}不也是父亲吗？因为他得了儿子却毫无肉体的贪欲。愿基督徒的贞洁绝不怀有连犹太人的贞洁都不曾怀有的想法！因此，要爱你们的妻子，但要贞洁地爱。在与她们的交往中，保持为生育子女所必需的界限。因为你们不能以其他方式获得子女，所以要带着遗憾降至于此。因为这一需要是对我们出自的\Person{亚当}的惩罚。我们不要以我们的惩罚为骄傲。那是他的惩罚，因为他因罪成了可朽的，被判只生育可朽的后代。天主并未撤回这一惩罚，好使人记念他从何种状态被召离，以及他被召往何种状态，并寻求那没有朽坏的交合。\par

26. 在那些子民中，由于在基督来临之前必须有丰盛的繁衍，藉着那子民的增多——在其中预像了为教会训导所需的一切——娶妻是一项责任，藉着妻子，那将在其中预表教会的子民可以增多。然而，当万民之王亲自降生时，贞洁的荣誉便始于主的母亲，她享有生育圣子而不失其童贞纯洁的特恩。因此，既然那是一个真正的婚姻，且是毫无玷污的婚姻，那么丈夫为何不该贞洁地接受妻子贞洁所生之子呢？因为正如她是贞洁的妻子，他也是贞洁的丈夫；正如她是贞洁的母亲，他也是贞洁的父亲。因此，任何人若说他因不是以通常方式生养圣子而不应被称为父亲，则他所注目的是生育子女时情欲的满足，而非情感的自然倾向。别人在肉身中所渴望实现的，他以更卓越的方式在精神中实现了。因为那些收养子女的人，以更大的贞洁在心中生育了他们，而在肉身中却无法生育。弟兄们，请思考收养的法则：一个人如何成为另一个非自己所生之子，以至于收养者的选择比生养者的本性在他身上拥有更大的权利？因此，\Person{若瑟}不仅应被称为父亲，而且是以最卓越的方式为父亲。因为男人也能从不是自己妻子的女人生育子女，他们被称为私生子，而合法婚姻的子女则高于他们。就出生方式而言，他们都是同样出生的，那么后者为何高于前者？岂不是因为妻子的爱——所生子女出自妻子——更为纯洁？在这件事上，两性的结合不予考虑，因为这在两个女人中是相同的。妻子在哪里具有优越性？岂不就在于她的忠贞、她的婚姻之爱、她的更真实更纯洁的情感？如果男人能无需交合便从妻子得子，他岂不应当因为所爱之人的更大贞洁而更加喜悦吗？\par

27. 由此也可看到，一个人如何可能不仅有两个儿子，而且有两个父亲。因为藉着收养的提及，你们或许想到事情可能如此。因为常有人说：\Quote{一个人可以有两个儿子，但不能有两个父亲。}但真理是，我们发现他也可以有两个父亲，如果一人由肉身生了他，另一人以爱收养了他。那么，如果一个人可以有两个父亲，\Person{若瑟}也可以有两个父亲；可能由一人生育，由另一人收养。如果真是这样，那么那些坚持\Person{玛窦}遵循了一套谱系、\Person{路加}遵循了另一套谱系的人的吹毛求疵又有什么意义呢？事实上我们发现确实如此，因为\Person{玛窦}给若瑟的父亲是\Person{雅各伯}，\Person{路加}则是\Person{赫里}。诚然，这似乎可能是指同一个人——若瑟是他的儿子——有两个名字。但由于祖辈以及他们所列举的所有其他祖先都不同，而且代际数目本身也一个多一个少，因此若瑟显然由此表明有两个父亲。既然这个问题所挑剔的难点已被解决，因为清楚的理由已表明，生育孩子的人可能是一个父亲，而收养他的人可以是另一个：假设有两个父亲，那么祖辈、曾祖辈以及向上所列举的其他祖先不同——因为他们来自不同的父亲——也就不足为奇了。\par

28. 不要让义子的法律在你们看来与圣经陌生，好像它只被承认于人间法律中，而不能与神圣书籍的权威相合。因为这是一件自古以来就确立的事，在教会书籍中也常听到：不仅自然的生育方式，而且意志的自由选择也能产生子女。因为妇女们，若自己没有孩子，常收养丈夫与婢女所生的孩子，甚至迫使丈夫以这种方式给她们孩子，如撒辣、辣黑耳和肋阿。这样行时，丈夫并没有犯奸淫，因为他们在这与夫妻本分有关的事上听从了妻子，正如宗徒所说：\BibleQuote{妻子对自己的身体没有主权，而是丈夫；同样，丈夫对自己的身体也没有主权，而是妻子。}梅瑟生自希伯来母亲，被弃后，被法郎的女儿收养。那时法律的形式虽与如今不同，但意志的选择被视为法律准则，正如宗徒在另一处也说：\BibleQuote{没有法律的外邦人，顺着本性去行法律上的事。}如果妇女被允许将没有生育的孩子认作自己的子女，为何男人不能将非自己身体所生，而是出于义子之爱的孩子认作自己的呢？因为我们读到圣祖雅各伯虽子女众多，却将他的孙子若瑟的儿子们认作自己的儿子，用这些话：\BibleQuote{这两个归我，他们要同他们的兄弟一样分得产业；至于你以后生的，都算是你的。}但也许有人会说，\Quote{义子}这个词并未出现在圣经中。仿佛名称有多重要，毕竟事物本身就在那里——即妇女拥有非自己所生的孩子，或男人拥有非自己所生的孩子。若有人不承认若瑟是义子，只要他承认若瑟是某人非所生的儿子，我也不反对。然而，宗徒保禄确实不断使用\Quote{义子}这个词，以表达一个伟大的奥迹。因为虽然圣经为证，我们的主耶稣基督是天主的独生子，却说他恩赐的弟兄和同继承者，是藉由一种通过神圣恩宠的义子关系而成。\BibleQuote{时期一满，天主就派遣了自己的儿子，生于女人，生于法律之下，为把在法律之下的人赎出来，使我们获得义子的名分。}在另一处：\BibleQuote{我们也在心中叹息，等待着义子名分的实现，就是我们身体的救赎。}再者，当他论及犹太人时：\BibleQuote{我曾祈求过，宁愿自己与基督隔绝，为我的弟兄，我的骨肉之亲而受诅咒。他们是以色列人，基于义子的名分、光荣、盟约、法律的颁布，都是他们的；圣祖也是他们的，而且按照血统，基督也是从他们来的，他是万有之上的天主，世世代代应受赞美。}他显示，\Quote{义子}这个词，或至少它所表示的事物，在犹太人中自古使用，正如他同时提到的盟约和法律的颁布一样。\par

29. 此外，还有另一种犹太人的特有方式，人可成为另一个非按肉身所生之人的儿子。因为亲属常娶无子女而死的近亲之妻，为死者立后。如此，所生之子既是他亲生父亲之子，也是他继承序列之人的儿子。这一切之所以说，是怕有人以为一人不可能有两个父亲，而诬蔑地认为叙述主族谱的两位福音作者中有一位在撒谎；尤其是我们看到，他们的用词本身就警告我们不要如此。因为玛窦（他被认为是提及若瑟所生之父的那位）如此列举世代：\Quote{这个生那个}，直至最后他说：\BibleQuote{雅各伯生若瑟。}但是路加——因为经由收养或因死者的继承从妻子所生的人，不能恰当地被称为\Quote{被生}——并未说\Quote{赫里生若瑟}或\Quote{若瑟为赫里所生}，而是说\BibleQuote{他是赫里的儿子}，无论他是被收养的，还是作为死者的近亲后代而生的。\par

30. 现在已经说得足够多了，以表明为何族谱是通过\Person{若瑟}而非通过\Person{玛利亚}计算这个问题，不应困扰我们；因为她作为一个母亲，没有肉欲；同样，他作为一个父亲，也没有任何肉身交合。\newline{}那么，就让族谱通过他上升和下降。我们不要因为他没有这种肉欲，就排除他作为父亲的身份。让他更大的纯洁反而证实他的父亲关系，免得圣\Person{玛利亚}本人责备我们。\newline{}因为她不会把自己的名字放在丈夫之前；却说：\BibleQuote{你的父亲和我，忧苦地寻找你。}\newline{}那么，不要让这些悖逆的埋怨者，做出\Person{若瑟}那贞洁的配偶没有做过的事。我们就通过\Person{若瑟}来计算吧，因为正如他因贞洁而成为丈夫，同样他也因贞洁而成为父亲。\newline{}并且，我们按照自然秩序和天主的律法，将男人置于女人之前。\newline{}因为如果我们抛弃他而只留下她，他会理直气壮地说：\Quote{你们为何排除我？为何族谱不通过我上升和下降？}\newline{}我们是否要对他说：\Quote{因为你没有通过你肉身的运作生祂吗？}他必将回答：\Quote{难道童贞女是通过肉身的运作生下祂的吗？圣神所成就的，祂是为双方成就的。}\newline{}福音说：\BibleQuote{他是一个义人。}\newline{}那么，丈夫是义的，妻子也是义的。圣神安息在他们双方的义德中，赐给了双方一个儿子。\newline{}在自然适合生育的性别中，祂成就了那也属于丈夫的生产。\newline{}因此，天使吩咐他们双方给孩子起名，由此双方父母的权威得以确立。\newline{}因为当\Person{匝加利亚}还哑着时，母亲给她新生的儿子起了名。当在场的人\BibleQuote{向他的父亲打手势，问他愿意叫他什么名字，他拿了一块写字板，写下了}她已说出的名字。\newline{}同样，天使对\Person{玛利亚}说：\BibleQuote{看，你将怀孕生子，要给祂起名叫\Person{耶稣}。}\newline{}天使也对\Person{若瑟}说：\BibleQuote{\Person{若瑟}，\Person{达味}之子，不要怕娶你的妻子\Person{玛利亚}，因为那在她内受孕的是出于圣神。她将生一个儿子，你要给祂起名叫\Person{耶稣}，因为祂要把自己的百姓从罪恶中拯救出来。}\newline{}经上又说：\BibleQuote{她给他生了一个儿子。}借此他被确立为父亲，确实不是按肉身，而是按爱情。\newline{}那么，让我们承认他是父亲，就如他确实是那样。因为福音作者们最谨慎、最明智地通过他来计算，无论是\Person{玛窦}从\Person{亚巴郎}下降到\Person{基督}，还是\Person{路加}从\Person{基督}通过\Person{亚巴郎}上升到天主。一个按下降顺序，另一个按上升顺序；但都是通过\Person{若瑟}。\newline{}为什么呢？\newline{}因为他是父亲。\newline{}如何是父亲？因为他越贞洁地是父亲，就越无可否认地是父亲。\newline{}诚然，他曾被以另一种方式认为是我们的主\Person{耶稣基督}的父亲，即如同其他父母按照肉身的生育，而非通过完全属神的爱情的丰盈。\newline{}因为\Person{路加}说：\BibleQuote{祂被认为是\Person{耶稣}的父亲。}为何被认为？因为人的思想和猜想指向通常发生在人身上的情况。\newline{}那么，主并非出自\Person{若瑟}的血统，尽管祂被以为如此；然而，童贞\Person{玛利亚}之子，他也是天主之子，是为\Person{若瑟}而生的，是他虔诚与爱情的果实。\par

31. 然而，为什么圣\Person{玛窦}按下降顺序计算，而圣\Person{路加}按上升顺序计算？我请你们留心听主在这件事上帮助我说的话；现在你们心中安逸，解除了所有这些吹毛求疵的困惑。\newline{}\Person{玛窦}通过他的族谱下降，以表示我们的主\Person{耶稣基督}下降来承担我们的罪，为使万民因\Person{亚巴郎}的后裔而受祝福。\newline{}因此，他不从\Person{亚当}开始，因为全人类都出自他。也不从\Person{诺厄}开始，因为洪水之后，全人类又从他的家族而出。\newline{}\Person{耶稣基督}这个人，作为从\Person{亚当}所出（所有人都从他而出），不能应验预言；同样，作为从\Person{诺厄}所出（所有人也都从他而出），也不能；而只能作为从\Person{亚巴郎}所出，他在当时被拣选，为使万民因他的后裔而受祝福，那时大地已充满万民。\newline{}但\Person{路加}按上升顺序计算，他没有从我们的主诞生叙述的开头开始列举族谱，而是从记述祂受\Person{若翰}的圣洗的地方开始。\newline{}如今，正如在主降生成人时，人类的罪被加在祂身上由祂承担，同样在祂受圣洗祝圣时，这些罪也被加在祂身上予以清除。\newline{}因此，圣\Person{玛窦}代表祂下降来承担我们的罪，按下降顺序列举族谱；而另一位则代表罪的清除（当然不是祂自己的罪，而是我们的罪），按上升顺序列举。\newline{}此外，圣\Person{玛窦}通过\Person{撒罗满}下降，\Person{达味}因\Person{撒罗满}的母亲而犯了罪；圣\Person{路加}通过\Person{纳堂}（同一位\Person{达味}的另一个儿子）上升，\Person{达味}借着他从罪中得以净化。\newline{}因为我们读到，\Person{纳堂}被派去责备他，好使他能通过悔改而痊愈。\newline{}两位福音作者在\Person{达味}那里相遇；一个下降，一个上升；从\Person{达味}到\Person{亚巴郎}，或从\Person{亚巴郎}到\Person{达味}，每一代都没有差别。\newline{}因此，\Person{基督}——既是\Person{达味}之子，也是\Person{亚巴郎}之子——上升到天主那里。因为当我们借着圣洗圣事更新，从罪的废除中，我们必须被带回天主面前。\par

32. 在玛窦所列举的世代中，占主导数字的是四十。因为圣经的习惯是不计整十数以外的余数。例如，以色列子民出埃及的年份说是四百年，实际上是四百三十年。同样，这里超过四十的那一代并不影响这个数字的主导地位。这个数字象征着我们在世上劳苦的生命，就是我们远离主的时候，在此期间，真理宣道的暂时性安排是必要的。因为数字十代表真福的圆满，乘以四（因季节的四分和世界的四分）就得四十。所以梅瑟、厄里亚以及中保本身、我们的主耶稣基督都禁食四十天，因为在此生中，需要戒绝肉体的诱惑。百姓也在旷野中流浪了四十年。洪水泛滥了四十天。主复活后四十天与门徒相处，使他们确信祂复活身体的真实性，以此表明在此生中——\Quote{我们远离主}（如已说，四十神秘地象征着此）——我们需要纪念主的身体，这在教会中实行，直到祂来临。因此，我们的主降生到此生，并且\BibleQuote{圣言成了血肉，为我们的罪恶被交付，又为我们的成义而复活}\BibleRef{若 1:14}\BibleRef{罗 4:25}，玛窦遵循了数字四十，所以超过的那一代并不妨碍其主导地位——正如那三十年并不妨碍四百的完美数字——或者还有这层深意：主本身，因祂的加入而构成四十一，祂降生到此生承担我们的罪恶，却因祂独特而卓越的位格，祂是人同时也是天主，被视为此生的例外。因为只有关于祂才说过，而且从未能或将要能对任何圣善之人（无论其在智慧和正义上如何完美）说过的话：\BibleQuote{圣言成了血肉}\BibleRef{若 1:14}。\par

33. 但路加从主的圣洗开始向上追溯世代，得出数字七十七，他从我们的主耶稣基督本身开始，经过若瑟，直到亚当，并上达天主。这是因为这个数字象征在圣洗中所涤除的一切罪恶。并非主自己在圣洗中有什么需要赦免的，而是祂以谦卑向我们显明了圣洗的益处。虽然那只是若翰的洗礼，但在其中以有形的方式显现了圣父、圣子及圣神的天主圣三，并以此祝圣了基督自己的圣洗，基督徒正是以此圣洗受圣洗。圣父藉从天上来的声音，圣子藉中保本身，圣神藉鸽子显现。\par

34. 至于为何数字七十七包含在圣洗中得以赦免的一切罪恶，有一个可能的理由：数字十代表一切义德和真福的圆满，当以七所象征的受造物依附于造物主的天主圣三时便是如此；因此法律的十诫也以十条诫命而祝圣。数字十一代表对\Quote{十}的\Quote{逾越}，而罪被认为是逾越，即人因追求\Quote{更多}而超出正义的规范。因此宗徒称贪婪为\BibleQuote{万恶的根源}\BibleRef{弟前 6:10}。对于那背离天主而行的灵魂，在同一主的位格中说过：\Quote{你曾希望，你若离开我，会拥有更多。}因为犯罪者在逾越中，即在罪中，只顾及自己——他希图藉某种私有的善来自我满足（为此那些\BibleQuote{寻求自己的事，而非耶稣基督的事}\BibleRef{斐 2:21}的人受到责备；而爱德受到赞扬，\BibleQuote{不求己益}\BibleRef{格前 13:5}）；因此，这代表逾越的数字十一，不乘以十，而乘以七，得七十七。因为逾越不指向造物主的天主圣三，而指向受造物，即人自己，这个受造物由数字七代表：三是因为灵魂，其中有造物主的天主圣三的某种肖像（因为人在灵魂中按天主的肖像被造）；四是因为身体。构成身体的四元素是众所周知的。即使有人不知，也容易记起，这个世界的身体（我们的身体在其中活动）可以说有四个主要部分，圣经也经常提及：东、西、南、北。因为罪要么由心灵（只在意志中）犯下，要么也由身体的行为而犯下，从而可见；因此，亚毛斯先知不断引入天主威胁说：\BibleQuote{为了三和四的过犯，我决不收回成命}\BibleRef{亚 1:3}，即我不掩饰我的愤怒。三是因为灵魂的本性，四是因为身体的本性；人由这两者组成。\par

所以，七乘十一，即如前所述，只涉及罪人自身的不义之过犯，构成七十七这个数目，其中表明圣洗中所赦免的一切罪都包含在内。因此，\Person{路加}经由七十七代上达于天主，显示人因所有罪的废除而与天主和好。因此，当\Person{伯多禄}问主该多少次宽恕兄弟时，主亲自说：\BibleQuote{我不对你说七次，而是到七十个七次。}凡能从这些天主奥秘的深处与宝库中，由比我更勤奋、更配得之人所抽取出来的，你们都要领受。然而，我已按我薄弱的才能，照主所助佑、所赐予我的能力，并尽我所能，也考虑到我时间不多，讲说了这些。如果你们中有人能领悟更多，就请向那位敲叩，我从祂那里也领受了我所能领受和讲说的。但最重要的是，要记住这一点：不要因你们尚未理解的圣经而困扰，也不要因你们所理解的而自高自大；对于你们不理解的，要谦逊等待；对于你们所理解的，要以爱德持守。\par

\SourceRecord{Translated by R.G. MacMullen.；Nicene and Post-Nicene Fathers, First Series；Vol. 6.；Edited by Philip Schaff.；Buffalo, NY: Christian Literature Publishing Co.,；1888.；<http://www.newadvent.org/fathers/160301.htm>.}

% structure: p0001 source=h1 target=section reason=page-title

\section{新约讲道集 第二篇}

\Emph{[LII. Ben.]}\par

\Emph{论玛窦福音3:13的话}，\BibleRef{玛 3:13}, \BibleQuote{那时，耶稣从加里肋亚来到约但河若翰那里，为受他的洗。} \Emph{论天主圣三}。\par

1. 福音的读经在我面前摆了一个话题，要给你们讲论，亲爱的，好像是因了主的命令，实在是因了祂的命令。因为我的心等待着祂的命令来讲话，为使我因此明白，祂愿意我讲论那祂也愿意读给你们听的内容。那么，愿你们的热情和虔诚侧耳倾听，并在主我们的天主面前，亲自帮助我的劳作。因为我们看见，仿佛在一个神圣的场景中展现在我们面前的，我们的天主圣三的启示，在约但河传给了我们。因为当耶稣来了，并由若翰受洗，主由祂的仆人（祂这样做是为谦逊的榜样；因为祂表明，在这种谦逊中，义德得以完成，正如若翰对祂说：\BibleQuote{我需要受祢的洗，而祢却到我这里来？} 祂回答说：\BibleQuote{你暂且容许吧，因为我们应当这样，以履行全部义德。} ），当祂受洗时，天就开了，圣神以鸽子的形状降在祂身上；随后，从高处传来声音说：\BibleQuote{这是我的爱子，我所喜悦的。} 所以我们在这里就有了天主圣三在某种程度上的区分。圣父在声音中，圣子在人身上，圣神在鸽子中。只需要提一下这一点，因为它是显而易见的。因为天主圣三的启示在这里清楚无误地传递给我们，没有留下任何怀疑或犹豫的余地。因为主基督亲自以仆人的形状来到若翰那里，毫无疑问是圣子：因为不能说那是圣父或圣神。经上说：\BibleQuote{耶稣来了}，即天主子。谁对鸽子有怀疑呢？或谁说：\Quote{什么是鸽子？} 而福音本身最清楚地证明：\BibleQuote{圣神以鸽子的形状降在祂身上。} 同样，关于那声音，毫无疑问那是圣父的，当祂说：\BibleQuote{祢是我的儿子。} 这样我们就区分了天主圣三。\par

2. 如果我们考虑场所，我自信地说（虽是恐惧地说），天主圣三在某种程度是可分离的。当耶稣来到约但河时，祂从一个地方到另一个地方；鸽子从天降下到地上，从一个地方到另一个地方；而圣父的声音既不是从地上，也不是从水上，而是从天上传来的；这三个似乎是在场所、职分和工作上分离的。但有人会对我说：\Quote{请你证明天主圣三是不可分离的吧。记住，你讲话的是一位天主教徒，并且是对天主教徒讲话。} 因为我们的信德是这样教导的，即那真实、纯正的天主教信德，不是由个人判断的意见收集的，而是由圣经的证据，不受异端鲁莽的波动，而是奠基于宗徒的真理：这是我们所知道的，这是我们所相信的。虽然我们尚未用眼睛看见，也尚未用心看见，只要我们尚在因信德而被净化，但藉着这信德，我们最坚定、最有力地持守：圣父、圣子、圣神是一个不可分离的天主圣三，一个天主，不是三个神。然而，是一个天主，如同圣子不是圣父，圣父不是圣子，圣神既不是圣父也不是圣子，而是圣父和圣子的圣神。这不可言喻的天主性，常存于自身，使万物更新，创造，再造，派遣，召回，审判，释放，我说，这个天主圣三既是不可言喻的，又是不可分离的。\par

3. 那么我要做什么呢？看：圣子单独地来到人身上；圣神单独地从天以鸽子的形状降下；圣父的声音单独地从天传出：\BibleQuote{这是我的儿子。} 那么这不可分离的天主圣三在哪里呢？天主藉着我的言语使你们专注。请为我祈祷，展开你们心灵的内褶，愿祂赐予你们，使你们敞开的内心得以充满。与我一同分担我的劳苦。因为你们知道我承担了什么；不仅是什么，而且是什么人承担了它，我想讲什么，以及我在哪里和我的处境是什么，甚至在这\BibleQuote{可朽坏的身体压迫着灵魂，属地的居所压伏着那多虑的心神。} 因此，当我把我的心思从事物的多样性中抽象出来，并把它集中于独一天主、不可分离的天主圣三时，为使我看见一些可以论说它的东西，你们想我在这个\BibleQuote{压迫灵魂的身体}中，是否能够说：\BibleQuote{上主，我向祢举起了我的灵魂。} 愿祂助我，愿祂与我一同举起它。因为我在祂面前是过于软弱的，而祂在我面前是过于强大的。\par

4. 这是热心弟兄们常常提出的一个问题，也常出现在热爱天主圣言者的谈话中；因为这样常向天主敲门，人们说：\Quote{圣父所做的是否是圣子不做的？或者圣子所做的是否是圣父不做的？} 让我们先谈论圣父和圣子。当祂——我们向祂说：\Quote{祢做我的助佑，不要离弃我}——赐予我们这个尝试的成功时，我们就会明白圣神也绝不与圣父和圣子的作为分离。那么，关于圣父和圣子，弟兄们，请听。圣父是否做任何事而不藉着圣子？我们回答：不。你们怀疑吗？因为祂不藉着祂\BibleQuote{万物是藉着祂而造的}，能做什么呢？圣经说：\BibleQuote{万物是藉着祂而造的。} 为了充分教导迟钝、固执和好辩的人，又加了一句：\BibleQuote{凡受造的，没有一样不是藉着祂而造的。}\par

5. 那麼，弟兄們，\BibleQuote{萬物是藉着祂造成的。} 如此我們明白，由聖子造成的整個受造界，是聖父藉着祂的聖言——天主，藉着祂的能力和智慧造成的。那麼，我們能說，萬物在被造時，\BibleQuote{是藉着祂造成的}，但如今聖父不藉着祂行事嗎？決不！但願這樣的思想遠離信徒的心；但願它從虔誠者的心意中被驅逐，從敬神者的理解中被逐出！祂藉着祂造成，卻不藉着祂管理，這是不可能的。決不！願存在之物若沒有祂的管理，因着祂才得以存在。但讓我們用同一聖經的見證來證明，不僅萬物是藉着祂被創造和造成，正如我們從福音引用的，\BibleQuote{萬物是藉着祂造成的，沒有祂，什麼也沒有造成}，而且被造成之物也由祂治理和安排。那麼，你承認基督是天主的德能和智慧；也要承認論及智慧所說的：\BibleQuote{她大有能力地從一端達到另一端，溫和地安排萬物。} 那麼，讓我們不要懷疑：萬物由祂治理，萬物由祂造成。所以，聖父沒有聖子什麼也不做，聖子沒有聖父什麼也不做。\par

6. 但是，有一個難題向我們迎面而來，我們要藉着主的名和祂的旨意來解決它。如果聖父沒有聖子什麼也不做，聖子沒有聖父什麼也不做，豈不是說聖父也由童貞瑪利亞誕生，聖父在般雀比拉多治下受苦，聖父復活並升了天？決不！我們不這樣說，因為我們不信這話。\BibleQuote{因為我信了，所以就說：我們也信，因此也說。} 信經中說什麼？聖子由貞女誕生，而非聖父。信經中說什麼？聖子在般雀比拉多治下受苦並死了，而非聖父。我們豈忘記了，有些人因誤解而稱為\OriginalTerm{聖父受苦論者}{Patripassians}，他們說聖父自己由女子誕生，聖父自己受苦，聖父就是聖子，兩者只是兩個名稱，而不是兩個實體？這些人大公教會已將他們隔離於聖人的共融之外，免得他們欺騙人，只教他們在教會之外爭論。\par

7. 那麼，讓我們回想問題的難點。有人會對我說：\Quote{你說過聖父沒有聖子什麼也不做，聖子沒有聖父什麼也不做，你又從聖經中引出見證，說聖父沒有聖子什麼也不做，因為\BibleQuote{萬物是藉着祂造成的}；又說被造之物沒有聖子也不被治理，因為祂是聖父的智慧，\BibleQuote{大有能力地從一端達到另一端，溫和地安排萬物。}現在你卻對我說，彷彿自相矛盾，說聖子由貞女誕生，而非聖父；聖子受苦，而非聖父；聖子復活，而非聖父。那麼看，這裡我看到聖子做了聖父不做的事。所以你要麼承認聖子做了聖父不做的事，要麼承認聖父也誕生、受苦、死亡並復活了。從兩者中選一個吧。} 不：我都不選，我既不說這個也不說那個。我既不說聖子做了聖父不做的事，因為我若這樣說就是說謊；也不說聖父誕生、受苦、死亡並復活，因為我若這樣說同樣是說謊。\Quote{那麼，他說道，你將如何從這些困境中解脫出來呢？}\par

8. 問題的提出令你們滿意。願天主賜予幫助，使問題的解決也令你們滿意。看，我求祂，使祂釋放我和你。因為我們在基督的名下站在同一信仰裡；在同一屋簷下生活，在同一主下；在同一身體中，我們是同一頭下的肢體；並藉着同一聖神而生活。那麼，願主使我這說話的，和你們這聽話的，都從這個最令人困惑的問題的困境中解脫出來，我這樣說：確實是聖子，而非聖父，由童貞瑪利亞誕生；但聖子的這個誕生本身，而非聖父的誕生，是聖父和聖子共同的工作。聖父確實沒有受苦，但聖子受苦了，然而聖子的受苦是聖父和聖子的工作。聖父沒有復活，但聖子復活了，然而聖子的復活是聖父和聖子的工作。我們似乎已經從這個問題中解脫了，但也許只是憑我自己的話；讓我們看看是否也憑天主的話。那麼，我的職責是用聖經的見證證明：聖子的誕生、苦難和復活，在某種意義上是聖父和聖子的工作，以至於雖然只是聖子的誕生、苦難和復活，但這三件只屬於聖子的事，既不是單由聖父做的，也不是單由聖子做的，而是由聖父和聖子共同做的。讓我們逐點證明，你們作為裁判來聽；案情已經展開；現在讓證人出場。願你們的判斷對我說，如同對訴訟中的辯護人常說的那樣：\Quote{證實你所應許的。} 我一定會，藉着主的幫助，引用天上法律的書卷。你們已專心聽了我提出問題，現在請更加專心聽我證明。\par

9. 我先要教导你们关于基督的诞生，它既是圣父也是圣子的作为，尽管圣父和圣子所完成的作为只属于圣子。我要引用\Person{保禄}，一位精通神圣律法的人。我要引用的就是这位\Person{保禄}，他规定的是和平的法律，而非诉讼的法律；因为今天的律师也有一位\Person{保禄}，他制定的是法庭的法律，而非基督徒的法律。那么，让圣宗徒向我们显示，圣子的诞生如何是圣父的作为。他说：\BibleQuote{但时期一满，天主就派遣了自己的儿子，生于女人，生于法律之下，为把在法律之下的人赎出来。} 你们这样听了，因它明白而清晰，就懂了。看，圣父使圣子由童贞女诞生。因为\BibleQuote{时期一满，天主就派遣了自己的儿子}；圣父派遣了祂的基督。祂如何派遣了祂？\BibleQuote{生于女人，生于法律之下。} 圣父于是使祂由女人而生，且处于法律之下。\par

10. 也许你们会感到困惑，我说\Quote{由童贞女}，而\Person{保禄}说\Quote{由女人}？不要让这困惑你们；我们不要在此停留，因为我并非对没有受过教导的人说话。圣经两者都说，既说\Quote{由童贞女}，也说\Quote{由女人}。哪里说\Quote{由童贞女}？\BibleQuote{看，童贞女将怀孕生子。} 而\Quote{由女人}，正如你们刚才所听见的；这里没有矛盾。因为希伯来语的特性，将\Quote{女人}这个称呼不适用于那些失去童贞身份的人，而是泛指女性。你们在\BibleBook{创世}{纪}中有一个明确的段落，当\Person{厄娃}本人最初被造时，\BibleQuote{他造了她成为一个女人。} 圣经在另一处也说，天主命令那些\Quote{\Emph{女人}}（即\Quote{未曾与男子同寝的}）被分别出来。那么，这一点现在应该已牢固确立，不应再耽搁我们，好使我们得以借主的助佑，来解释那些值得耽搁我们的。\par

11. 我们已经证明了圣子的诞生是圣父的作为；现在让我们证明它也是圣子的作为。那么，圣子由童贞玛利亚的诞生是什么呢？当然就是祂在童贞女的胎中取了奴仆的形体。圣子的诞生岂非别样，正是取了奴仆的形体在童贞女的胎中？现在请听，这如何也是圣子的作为：\BibleQuote{祂虽具有\OriginalTerm{天主的形体}{form of God}，却没有以自己与天主同等为应当把持的，反而\OriginalTerm{空虚自己}{emptied Himself}，\OriginalTerm{取了奴仆的形体}{taking the form of a servant}。} \BibleQuote{时期一满，天主就派遣了自己的儿子，生于女人，} 祂是\BibleQuote{按血统由达味后裔所生。} 在这一点上，我们看见圣子的诞生是圣父的作为；但圣子自己\BibleQuote{\OriginalTerm{空虚自己}{emptied Himself}，\OriginalTerm{取了奴仆的形体}{taking the form of a servant}}，我们看见圣子的诞生也是圣子自己的作为。那么这一点已经证明；所以让我们由此点继续前进，你们要留心听取接下来顺序的内容。\par

12. 让我们证明圣子的受难也是圣父和圣子的作为。我们看见圣子的受难是圣父的作为，因为经上记载：\BibleQuote{祂没有怜惜自己的儿子，反而为我们众人把祂交出了}；并且圣子的受难也是祂自己的作为：\BibleQuote{祂爱了我，为我舍弃了自己。} 圣父交出了圣子，圣子也交出了自己。这一受难是为同一事而成就，却由两者共同完成。因此，正如基督的诞生，同样，祂的受难既不是没有圣父的圣子的作为，也不是没有圣子的圣父的作为。圣父交出了圣子，圣子也交出了自己。犹达斯在其中做了什么，不过是他的罪而已？那么，让我们也由这一点前进，来到复活。\par

13. 让我们看圣子确实复活了，而不是圣父，但圣父和圣子共同完成了圣子的复活。圣子的复活是圣父的作为；因为经上记载：\BibleQuote{因此，天主高高举扬了祂，赐给了祂一个名字，超越一切名字。} 圣父于是藉着举扬祂，并从死者中唤醒祂，使圣子复活了。那么，圣子是否也使自己复活了呢？当然，祂这样做了。因为祂论及圣殿，作为自己身体的预像，说：\BibleQuote{你们拆毁这座圣殿，三天之内，我要把它重建起来。} 最后，正如舍弃生命与受难有关，同样，再取回生命与复活有关。那么，我们来看圣子是否真的舍弃了自己的生命，而圣父将生命归还给祂，而不是祂自己归还。因为圣父归还了生命是明显的。因为圣咏这样说：\BibleQuote{求你使我复起，我将报复他们。} 但你们为何等待我来证明圣子也使自己生命复原呢？让祂自己说：\BibleQuote{我有权舍弃我的生命，也有权再取回它；没有人能夺去我的生命，而是我甘愿舍弃它，并且再取回它。}\par

14. 我已履行了所许诺的；我认为，我已用最有力的证据和见证确立了主张。那么，请坚守你们所听到的。我将简要复述一遍，并把它作为我认为最有用的东西，托付给你们珍藏于心。圣父并非由童贞玛利亚所生；然而，圣子由童贞玛利亚的诞生，却是圣父和圣子共同的工作。圣父并未在十字架上受难；然而，圣子的苦难，却是圣父和圣子共同的工作。圣父并未从死者中复活；然而，圣子的复活，却是圣父和圣子共同的工作。因此，你们看到了位格的区分，以及行动的不可分离性。所以，我们不要说圣父没有圣子做了什么事，或圣子没有圣父做了什么事。但也许你们对耶稣所行的奇迹有疑问，恐怕祂做了些圣父没有做的事！那么，那句\BibleQuote{住在我内的圣父，祂自己作这些事？}的话在哪里呢？我刚才所说的一切都是明白的；只需稍加提及；无需费很大力气去理解，只需留心，以便唤起你们的记忆。\par

15. 我想再说些什么，并在此真诚地恳求你们更专心地注意，并向上主保持虔诚。因为只有物体才被容纳或包裹在适合物体本性的地方。天主性超越所有这类地方：不要有人像在空间中那样寻找它。它无所不在，无形且不可分离；不是一部分大，另一部分小；而是整个无所不在，无处可分。谁能看见？谁能领会？让我们约束自己：让我们记住我们是谁，以及我们谈论的是谁。让这一切，或任何属于天主本性的事物，以虔诚的信德拥抱，以神圣的敬意接纳，并尽我们所能，尽我们可能，以不可言喻的方式理解。让言语静默，让舌头缄默，让心灵被唤醒，让心灵被提升到那里。因为它的本性并非能够上升到人心里；而是人心本身必须上升到它那里。让我们默想受造物（\BibleQuote{因为自从创造世界以来，祂那看不见的美善，都可凭理智，在祂所造的万物上，得以认识}），希望在天主所造的、我们与之有些熟悉交往的事物中，我们或许能找到某种相似，用以证明有些三样东西可以分开地呈现，而它们的作用却是不可分离的。\par

16. 来吧，弟兄们，请全神贯注。但首先请考虑我所许诺的是什么：如果在受造物中我能找到任何相似，因为造物主远远高于我们。也许我们中有些人，其心灵被真理的光芒如同火花般击中，能说出那些话，\BibleQuote{我曾在神魂超拔中说。}---你在神魂超拔中说了什么？---\BibleQuote{我由你面前被抛弃。}因为在我看来，说这话的人似乎曾将灵魂高举到天主面前，并被带到自己之外，当时人们天天对他说，\BibleQuote{你的天主在哪里？}---他通过一种属灵的接触，到达了那不变的光明，却因自己视力的软弱而无法承受它，于是又跌回到他自己那病弱的状态，并将自己与那光明比较，感到他心灵的眼睛还未能适应天主智慧的光。因为他是在神魂超拔中做了这事，从他肉身的知觉中匆匆被带走，并被提到天主那里，当他仿佛从天主那里被召回到人那里时，他说，\BibleQuote{我曾在神魂超拔中说。}因为我在神魂超拔中见到了我不知道的什么，我无法长久忍受；恢复到我必死的状态，以及来自那压迫灵魂的肉身的多种必死思想时，我说了什么呢？\BibleQuote{我由你面前被抛弃。}你远在高处，而我远在低处。那么，弟兄们，我们关于天主能说什么呢？因为如果你能领会你所要说的，那便不是天主；如果你能领会它，那你就领会了别的什么而不是天主。如果你自以为领会了祂，那你就欺骗了自己。所以，这不是天主，如果你领会了它；但如果它是天主，你便没有领会它。那么，你怎能谈论你所不能领会的东西呢？\par

17. 让我们看看，是否我们能在受造物中找到某种三个事物，它们分别展现，但其运作不可分离。可我们要往哪里去？到天上去，争辩太阳、月亮和星宿？到地上去，争辩灌木、树木和充满大地的动物？或是争辩天和地本身，它们包含了天地间的一切？人啊，你还要在受造物中漫游多久？回归你自己，看，思考，审视你自己。你在受造物中寻找某种三个事物，它们分别展现，但其运作不可分离；如果你在受造物中寻找这个，就先在你自己内寻找。因为你并非别物，也是一个受造物。你寻找的是一个肖像。难道你要在牲畜中寻找它？因为你曾谈论的是天主，当你寻找这肖像时。你谈论的是不可言喻的威严的天主圣三，但因你在默观天主性时失败，并因适当的谦卑承认了你的软弱，你便下降到人性；那么，在那里继续你的探究。你会在牲畜、太阳或星宿中寻找吗？这些中哪一个曾按天主的肖像和模样被造？你可以在你自己内寻找更熟悉、比这一切更卓越的东西。因为天主按他自己的肖像和模样造了人。那么，在你自己内寻找，看这肖像是否保留了一些天主圣三的痕迹。这肖像是什么？它是一个与其原型大不相同的肖像；然而虽不相同，它仍是一个肖像和模样，但不是像圣子那样是肖像，圣子与圣父是同一的。因为肖像在儿子中是一种方式，在镜子中是另一种方式。它们之间有巨大的差别。你在你儿子内的肖像就是你自身，因为儿子按本性就是你之所是。与你同一实体，但位格不同。那么，人不是像独生子那样的肖像，而是按某种肖像和模样造的。让他自己在内寻找，是否可以找到，某种三个事物，它们分别展现，但其运作不可分离。我将寻找，你们也与我一起寻找。我不会在你们内寻找，但你们在你们内寻找，我在我内寻找。我们一起寻找，并一起讨论我们共同的本性和实体。\par

18. 看，人啊，并思考我所说的是否真实。你有身体和血肉吗？我有，你说。因为我如何在这我现在占据的地方，又如何从一个地方移动到另一个地方？我如何听见说话者的话，难道不是靠我身体的耳朵吗？我如何看见说话者的嘴，难道不是靠我身体的眼睛吗？那么显然你有身体，无需为如此明显的事烦恼。再考虑另一点，考虑是什么通过这身体行动。因为你通过耳朵听见，但听见的不是耳朵。有别的内在的东西通过耳朵听见。你通过眼睛看——检查这眼睛。怎么！你认出了房子，却不注意居住者？眼睛自己看吗？不是有另一个通过眼睛看吗？我不说死人的眼睛，显然居住者已离开身体，它不看，甚至任何心不在焉的人的眼睛也不看面前对象的形状。那么，请看看你的内在之人。因为在那里更应寻找某种三个事物的肖像，它们分别展现，但其运作不可分离。那么在你心中有什么？或许如果我寻找，我会在那里找到许多东西，但有些东西非常接近，更容易理解。那么在你灵魂中有什么？回忆起它，默想它。因为我不要求在我将要说的方面得到信任；如果你在自己内找不到，就不要承认。那么向内看；但首先让我们看看我是否忽略了什么：人不是仅圣子或仅圣父的肖像，而是圣父和圣子的肖像，因而当然也是圣神的肖像。\BibleBook{}{创世纪}中的话是说：\BibleQuote{讓我們按我們自己的肖像和模樣造人。}因此，圣父不离开圣子行动，圣子也不离开圣父。\BibleQuote{讓我們按我們自己的肖像和模樣造人。讓我們造，}不是\Quote{我要造，}或\Quote{你造，}或\Quote{他造，}而是\BibleQuote{按我們的肖像。}\par

19. 我要问，我要说——请记住，这只是一种遥远的相似。所以没有人可以说：看，他将什么与天主相比！我已经预先提醒过你们，并且事先既警戒了你们，也警戒了我自己。两者确实相去甚远，如同最低与最高、可变的与不变的、受造的与造物主、人性与神性。看！我事先就告知你们这点，以免有人因我所要说的对象之间如此巨大的差异而反对我。那么，在我求取你们的耳朵时，恐怕你们中有人正磨牙以待，请记住我仅仅旨在表明：有三样东西，它们各自被呈现，但其运作不可分割。这些东西与全能的\OriginalTerm{天主圣三}{Trinity}何其相似或何其不相似，我目前并不关心；但在受造物中，在最卑微、可变的事物中，我们确实发现三样东西，它们各自被呈现，但其运作不可分割。哦，肉性的想象！顽固、不信的良心！为什么关于那不可言喻的\OriginalTerm{威严}{Majesty}，你们怀疑在你们自身中所能发现的事？因为我要问你，人啊，你有记忆吗？如果没有，你怎么能记住我所说的话？但也许你已经忘记了我刚才所说的话。然而，正是这些话——\Quote{我说}——这两个音节，除非通过记忆，否则你无法保留。因为你若在第二个音节响起时忘记了第一个，又怎会知道它们是两个呢？但我为何在此停留？为何如此急切？为何如此强迫人信服？因为你有记忆，这是明显的。那么我在寻找别的。你有理解吗？\Quote{我有}，你会说。因为如果你没有记忆，你就不能保留我所说的话；如果你没有理解，你就不能理解你所保留的。那么你也有这个，和另一个一样。你将你的理解回溯到你内里所保留的，从而看见它，并通过看见而被塑造为据说认识的状态。但我在寻找第三样东西。你有记忆，用以保留所说的话；你有理解，用以理解所保留的；但关于这两样，我再问你：你不是用你的意愿去保留和理解的吗？无疑，是用我的意愿，你会说。那么你就有意愿。\par

这就是我承诺要带到你们耳中、心中的三样东西。这三样东西在你们内，你们可以数算，但不能分开。那么这三样——记忆、理解、意愿——我说，这三样，思想它们如何各自被呈现，但其运作不可分割。\par

20. 主要作我眼前的帮助，我看祂就在帮助我；藉你们理解我所说的，我看祂就在帮助我。因为我从你们的声音感知到你们理解了我，我确实相信祂将继续帮助我们，使你们能完全领悟。我承诺向你们展示三样东西，它们各自被呈现，但其运作不可分割。看，我不知道你们心中所想，而你们通过说 \Quote{记忆} 向我表明。这个字、这个声音、这个表达从你们的心来到我的耳中。因为在此之前，你们对记忆有沉默的概念，但没有表达出来。它在你们内，但尚未到我这里。但为了使那在你们内的能传递给我，你们说出了那个字，即 \Quote{记忆}。我听见了，我听见了这个字中的三个音节：\Quote{记忆}。它是一个名词，一个三个音节的词，它响了，来到我的耳中，并在我的心中印下某个观念。声音过去了，但传达观念的那个词和观念本身仍然存在。但我问，当你们说出 \Quote{记忆} 这一词时，你们当然看到它只指向记忆。因为另外两样东西有它们自己的专名。一个叫做 \Quote{理解}，另一个叫做 \Quote{意愿}，而不是 \Quote{记忆}，唯独那个叫做 \Quote{记忆}。然而，你们藉着什么运作来表达它，以产生这三个音节？这个只指向记忆的词，记忆在你们内运作，使你们记住所说的；理解运作，使你们知道所记住的；意愿运作，使你们表达所知道的。感谢主我们的天主！祂帮助了我们，既你们也我。因为我告诉你们实话，亲爱的，我以极大的恐惧承担了对此主题的考察和解释。因为我恐怕可能使心胸开阔的人快乐，而给较迟钝者带来痛苦的疲倦。但现在我既从你们聆听的专注，也从你们理解我的敏捷中看到，你们不仅领会了我所说的，而且预见了我的话。感谢主！\par

21. 请你们看，从此以后我安全地讲论你们已经领悟的事；我不是在灌输未知的教训，而是藉着复述传递你们已领受的。在这三样中，只提到并表达了其中一样；在这三样中，只有\Quote{记忆}是其中一样的名称，然而这三样共同产生了这单独一样的名称。单独\Quote{记忆}一词，若不藉着意志、理解和记忆的运作，就不能被表达。单独\Quote{理解}一词，若不藉着记忆、意志和理解的运作，就不能被表达；单独\Quote{意志}一词，若不藉着记忆、理解和意志的运作，就不能被表达。因此，我所许诺的，我想已经解释清楚了：我分别说出它们，却不分离地理解它们。这三样共同产生了其中每一个，然而这三样所产生的那一个，并不指向三者，而是指向其中之一。这三样共同产生了\Quote{记忆}一词，但这词只指向记忆。这三样共同产生了\Quote{理解}一词，但这词只指向理解。这三样共同产生了\Quote{意志}一词，但这词只指向意志。同样，天主圣三共同参与了基督身体的形成，但这身体只属于基督。天主圣三共同参与了从天上降下的鸽子的形成，但这鸽子只属于圣神。天主圣三形成了从天上发出的声音，但这声音只属于圣父。\par

22. 那么，谁也不可对我说，谁也不可用不公的刁难来压迫我的软弱，说：\Quote{你指出在我们心灵或灵魂中的这三样，哪一样对应圣父，即如同圣父的肖像？哪一样对应圣子的肖像？哪一样对应圣神的肖像？}我说不出——我无法解释这一点。让我们保留一些给默想和静默。进入你自身；远离一切喧嚣。审视你的内在；看看你是否在那里有一个良心的甜美退隐处，那里没有噪音、没有争论、没有争斗或辩论；那里没有分裂和固执争论的念头。要温顺地听道，好使你领悟。也许你不久就会说：\Quote{祢将使我听到欢喜和快乐的声音，我的骨骸将要欢跃}；骨骸，即那些谦卑的，而不是那些高傲的。\par

23. 因此，我指出有三样事物被分别显示，但其运作不可分离，这已足够。如果你在自己内发现了这一点；如果你在人身上发现了这一点；如果你在行走于大地上、带着那脆弱\Quote{身体，这身体重压着灵魂}的受造物身上发现了这一点；那么请相信：圣父、圣子和圣神可以通过某些可见的标记、某些取自受造物的形式被人分别显示，而它们的运作仍然不可分离。这就够了。我不说\Quote{记忆}是圣父——\Quote{理解}是圣子——\Quote{意志}是圣神；我不这样说；让人随自己的意思去理解吧。我不敢这样说。让我们为那些能够承受更大真理的人保留它们；但我自己软弱，只能向软弱者传达符合我们能力的事。我不说这些事物在任何方式下等同于天主圣三，或按照类比——即某种精确的比较规则——来加以等同。我不这样说。但我说什么呢？请看。我在你内发现了三样事物，它们被分别显示，而其运作不可分离；在这三样中，每一个单独的名称都是由三者共同产生的；然而这名称并不属于三者，而是属于其中之一。因此，请相信你无法看见的天主圣三，如果你在自己内已经听见、看见并保存了它。因为你自己内的事物你可以认识；但造你的那位内的事物，无论是什么，你怎能认识呢？即使你将来能够认识，现在你还不能。即使你将来能够认识，你怎能像天主认识自己那样认识天主呢？因此，亲爱的诸位，这对你们就足够了：我已经说了我所能说的；我按你们的要求履行了我的许诺。至于其余必须添加的，好使你们的理解得以进步，当从主那里寻求。\par

\SourceRecord{Translated by R.G. MacMullen.；Nicene and Post-Nicene Fathers, First Series；Vol. 6.；Edited by Philip Schaff.；Buffalo, NY: Christian Literature Publishing Co.,；1888.；<http://www.newadvent.org/fathers/160302.htm>.}

% structure: p0001 source=h1 target=section reason=page-title

\section{新约讲道三}

\Emph{[LIII. Ben.]}\par

\Emph{论福音之言，\BibleRef{玛 5:3-8}，\BibleQuote{神贫的人是有福的，}等等，但特别是那句，\BibleQuote{心里洁净的人是有福的，因为他们要看见天主。}}\par

因一位圣洁贞女的纪念日再次来临，她为基督作了见证，也配得基督的见证；她公开被杀，无形中得冠；这提醒我，亲爱的诸位，要向你们谈论主刚刚在福音中发出的劝言，他向我们保证，有福的生命有许多源头，没有一个人不渴望有福的生命。肯定找不到一个人不希望有福。但是，唉！但愿人渴望奖赏，也如同他们不愿回避通往奖赏的工作！谁不会欣然奔跑，若有人告诉他：\BibleQuote{你将是有福的}？那么，当有人说\BibleQuote{你们若这样做，便是有福的}时，让他也乐意而留心倾听。如果奖赏被爱慕，就不要回避竞赛；藉着奖赏的展示，让心灵被点燃，热切地执行工作。我们所渴望、所愿望、所寻求的，将在将来；但我们为了将来而奉命要做的事，必须在现在。那么，现在就开始忆起神圣的言语，以及福音的诫命和奖赏。\BibleQuote{神贫的人是有福的，因为天国是他们的。}天国将来是你们的；现在要神贫。你愿意将来天国是你们的吗？好好看看你现在是谁的。要神贫。你也许问我：\Quote{什么是神贫？}骄傲的人不是神贫的；所以谦卑的人就是神贫的。天国是高傲的；但\BibleQuote{自谦自卑者必被高举。}\par

注意以下的话：\Quote{有福的，}祂说，\BibleQuote{是温良的人，因为他们要承受土地。}你希望现在就拥有土地；要小心免得被土地所拥有。如果你温良，你将拥有它；如果你不温和，你将被它拥有。当你听到所提议的奖赏时，不要为了拥有土地而展开贪婪的衣襟，藉此你暂时拥有土地，甚至不惜以任何方式排除你的近人；不要让这样的幻想欺骗你。当你依附于创造天地的主时，你才能真正拥有土地。因为温良就是不抵抗你的天主，这样，在你做好事时，祂会喜悦你，而不是你喜悦自己；在你公正地受苦时，祂不会不喜悦你，而是你喜悦自己。因为当你对自己不满意时，你将使祂喜悦，这并非小事；而如果你对自己满意，你将使祂不喜悦。\par

注意第三端教训：\BibleQuote{哀恸的人是有福的，因为他们要受安慰。}工作在于哀恸，奖赏在于安慰；因为那些按肉体哀恸的人，有什么安慰呢？可怜的安慰，倒不如说是可怕的对象。在那里，哀恸者被那些使他害怕再次哀恸的事物所安慰。例如，儿子的死亡给父亲带来悲伤，儿子的出生带来喜乐。一个他送去了埋葬，另一个他带到了世上；在前者中有悲伤的缘由，在后者中有恐惧；所以两者都没有安慰。所以真正的安慰将是那样，在其中将赐予那不可失去的东西，以致那些现在为流放而哀恸的人，将为以后的安慰而喜乐。\par

让我们来到第四项工作及其奖赏：\BibleQuote{饥渴慕义的人是有福的，因为他们要得饱饫。}你希望被充满吗？用什么？如果肉体渴望饱足，消化之后你会再次饥饿。所以祂说：\BibleQuote{凡喝这水的，还要再渴。}如果用于伤口的药膏治愈了它，就没有更多疼痛；但用于饥饿的补救，即食物，只是暂时缓解。因为饱足过去后，饥饿又回来。这种饱足的补救天天应用，但软弱的伤口却没有痊愈。因此，让我们\BibleQuote{饥渴慕义，好使我们得饱饫}，我们将被我们现在所饥渴的义所充满。因为我们饥渴的东西会使我们充满。那么，让我们的内心人饥渴，因为它有其适当的食粮和饮料。祂说：\BibleQuote{我是从天上降下来的食粮。}这也是饥饿者的食粮；也渴望干渴者的饮料，\BibleQuote{因为在你那里有生命的泉源。}\par

注意接下来：\BibleQuote{怜悯人的人是有福的，因为他们要受怜悯。}这样做，这样对待你，你对待别人，愿天主也这样对待你。因为你同时处于丰裕和匮乏之中——在暂时之物上丰裕，在永恒之物上匮乏。你所听到的那个人是一个乞丐，而你自己是天主的乞丐。人向你请求，而你也向天主请求。正如你对待你的请求者，天主也会对待祂的请求者。你同时是满的和空的；用你的丰满充满空虚的，以便你的空虚能被天主的丰满所充满。\par

6. 请注意接下来所说的：\BibleQuote{心地纯洁的人是有福的，因为他们要看见天主。} 这就是我们爱德的终结；一个使我们得以完善而非消灭的终结。因为食物有终结，衣服也有终结；食物的终结在于被吃尽，衣服的终结在于织成。两者都有终结，但一个是消灭的终结，另一个是完善的终结。我们现在所做的一切，我们现在所善的一切，我们现在所追求的一切，或在值得称赞的热心中所渴望的一切，或在无可指摘的愿望中所追求的一切，当我们达到看见天主的境界时，就不需要再有什么了。因为天主与他同在的人，还需要寻求什么？或者天主不能使之满足的人，还有什么能使他满足？我们渴望看见天主，我们追求，我们因渴望看见祂而炽热。谁不如此呢？但请注意所说的：\BibleQuote{心地纯洁的人是有福的，因为他们要看见天主。} 那么请为你自己准备那能使你看见祂的。因为（按肉体来说）你怎能用微弱的眼睛渴望太阳的升起？让眼睛健康，那光就是喜乐；若不健康，它只会是折磨。因为不被允许以不洁的心看见那只能被纯洁的心看见的。你将被排斥，被赶回，无法看见它。因为\BibleQuote{心地纯洁的人是有福的，因为他们要看见天主。} 祂已经多少次列举了有福的人，以及他们有福的原因，他们的工作和回报，他们的功绩和奖赏！但从未说过\Quote{他们要看见天主。} \BibleQuote{神贫的人是有福的，因为天国是他们的。} \BibleQuote{温良的人是有福的，因为他们要承受土地。} \BibleQuote{哀恸的人是有福的，因为他们要受安慰。} \BibleQuote{饥渴慕义的人是有福的，因为他们要得饱足。} \BibleQuote{怜悯人的人是有福的，因为他们要受怜悯。} 在这些当中，从未说过\Quote{他们要看见天主。} 当我们来到\BibleQuote{心地纯洁}时，那里应许了看见天主。而且不是没有原因的；因为在那里，在心中，有眼睛，通过它们看见天主。论及这些眼睛，宗徒保禄说：\BibleQuote{照亮你们心中的眼睛。} 如今这些眼睛被照亮，适合它们的软弱，借着信德；以后适合它们的力量，将借着看见被照亮。\BibleQuote{因为我们住在身内，是离家出外，离开主，因为我们现今是凭信德往来，并非凭目睹。} 当我们仍处于这信德的状态时，关于我们说了什么？\BibleQuote{我们现在是借着镜子观看，模糊不清，到那时，就要面对面地看见了。}\par

7. 此处不要想到身体的面容。因为如果你因渴望看见天主而燃烧，已经准备好你的身体面容去看祂，你也会在神身上寻找这样的面容。但是，如果你现在对天主的观念至少是足够精神的，不至于想象祂是有形体的（关于这主题我昨天讲了很长时间，如果还没有白费的话），如果我已经成功地在你们心中，如同在天主的圣殿里，摧毁了那人形的形象；如果宗徒所说的那些话——\BibleQuote{他们自称聪明，却成了愚妄，将不朽坏的天主的荣耀变为仿佛必朽坏的人的形像}——已深入你们的心思，占据你们内心；如果你们现在憎恶并厌恶这样的不敬，如果你们为创造主保持祂自己的圣殿洁净，如果你们愿意祂来与你们同住，\BibleQuote{你们应以善良的心怀念上主，以诚实的心寻找祂。} 请留意，你们向谁说，如果你们确实说了，并且真诚地说：\BibleQuote{我的心向我说：请寻求祂的仪容。} 让你们的心里也说，并加上：\BibleQuote{上主，我要寻求祢的仪容。} 因为这样你们将好好寻求，因为你们用心寻求。圣经提到了\Quote{天主的仪容、天主的手臂、天主的手、天主的脚、天主的宝座}和祂的脚凳；但不要在其中想到人的肢体。如果你们要成为真理的圣殿，就要拆毁谎言的偶像。天主的手是祂的能力。天主的仪容是对天主的认识。天主的脚是祂的亲临。天主的宝座，如果你们愿意，就是你们自己。但或许你们敢否认基督是天主！\Quote{不敢。}你们说。你们也承认这点：基督是天主的德能、天主的智慧吗？\Quote{我承认。}你们说。那么请听：\Quote{义人的灵魂是智慧的宝座。} \Quote{是的。}因为天主在哪里有祂的宝座，不就是祂居住的地方吗？祂住在哪里，不就是祂的圣殿里吗？\BibleQuote{天主的圣殿是神圣的，你们就是那座圣殿。} 所以要注意你们如何接待天主。\BibleQuote{天主是神，应当以心神以真理去朝拜祂。} 现在让约柜进入你们的心吧，如果你们愿意，让达贡倒塌。现在立刻专心听，学会渴慕天主；学会准备那使你们能看见天主的条件。祂说：\BibleQuote{心地纯洁的人是有福的，因为他们要看见天主。} 你们为什么准备身体的眼睛？如果能借它们看见，那被看见的东西必然被空间所容纳。但那位完全无处不在者并不被空间所容纳。请洁净那能看见祂的（心）。\par

8. 请听并明白，如果依靠祂的帮助我能够解释的话；愿祂帮助我们理解上述所有的行为与赏报，看相称的赏报如何分派给相应的职责。因为哪里提到一种赏报，不与它的工作相配并和谐呢？因为谦卑的人似乎像天国的异乡人，祂说：\BibleQuote{神贫的人是有福的，因为天国是他们的。}因为温良的人容易失去他们的土地，祂说：\BibleQuote{温良的人是有福的，因为他们要承受土地。}其余的立刻显而易见；它们自身就能明白，不需要人长篇大论地阐释；只需有人提及它们。\BibleQuote{哀恸的人是有福的。}现在哪个哀恸的人不渴望安慰呢？\BibleQuote{他们，}祂说，\BibleQuote{必要受安慰。}\BibleQuote{饥渴慕义的人是有福的。}哪个饥渴的人不寻求饱饫呢？\BibleQuote{他们，}祂说，\BibleQuote{必要得饱饫。}\BibleQuote{怜悯人的人是有福的。}哪个怜悯人的人不希望天主以祂自己的行为回报他，好使别人对他所做的，如同他对穷人所做的那样？\BibleQuote{有福的，}祂说，\BibleQuote{是怜悯人的人，因为他们要受怜悯。}每一项职责都有相应的赏报；且赏报中没有引入任何不符合诫命的东西！因为诫命是要你\BibleQuote{神贫}；赏报是你将拥有\BibleQuote{天国}。诫命是要你\BibleQuote{温良}；赏报是你将\BibleQuote{承受土地}。诫命是要你\BibleQuote{哀恸}；赏报是你将\BibleQuote{受安慰}。诫命是要你\BibleQuote{饥渴慕义}；赏报是你将\BibleQuote{得饱饫}。诫命是要你\BibleQuote{怜悯人}；赏报是你将\BibleQuote{受怜悯}。同样，诫命是要你洁净心灵；赏报是你将看见天主。\par

9. 但不要这样理解这些诫命和赏报，以致当你听见\BibleQuote{心里洁净的人是有福的，因为他们要看见天主，}就认为神贫的人，或温良的人，或哀恸的人，或饥渴慕义的人，或怜悯人的人，将不会看见祂。不要以为只有那些心里洁净的人才能看见祂，而其他人将被排除在看见祂之外。因为所有这些不同的特征都是同一些人所具有的。他们都将看见；但他们的看见不是由于他们神贫、温良、哀恸、饥渴慕义或怜悯人，而是由于他们心里洁净。正如身体的职能被适当地分派给身体的各个肢体，有人举例说：有脚的人是有福的，因为他们要行走；有手的人是有福的，因为他们要工作；有声音的人是有福的，因为他们要呼喊；有口舌的人是有福的，因为他们要说话；有眼睛的人是有福的，因为他们要看见。同样，我们的主将灵魂的各个肢体按顺序排列，教导了什么是每种肢体所特有的。谦卑使人配得拥有天国；温良使人配得承受土地；哀恸配得安慰；饥渴慕义配得饱饫；怜悯配得受怜悯；洁净的心配得看见天主。\par

10. 因此，如果我们渴望看见天主，我们的眼睛将如何被洁净呢？因为谁不会关心并勤勉地寻求洁净那眼睛的方法，好藉此看见他以全部的热爱所渴望的那一位呢？神圣的记载已明确提到这事，它说：\BibleQuote{藉信德洁净他们的心。}对天主的信德洁净心灵，洁净的心看见天主。但因为这种信德有时被自欺的人这样定义，好像只要相信就够了（因为有些人相信却生活邪恶，他们甚至应许自己看见天主和天国）；宗徒雅各伯似乎被神圣的爱德所激怒和愤慨，在他的书信中说：\BibleQuote{你信只有一个天主。}你为自己的信德而鼓掌，因为你注意到许多不虔敬的人以为有许多神，你因自己相信只有一个天主而沾沾自喜；\BibleQuote{你做得对：魔鬼也信，且颤栗。}他们也能看见天主吗？心里洁净的人将要看见祂。但谁能说不洁的鬼神是心里洁净的呢？然而他们也\BibleQuote{信，且颤栗。}\par

11. 如此，我们的信德必须与魔鬼的信德不同。因为我们的信德洁净心灵，而他们的信德却使他们有罪。因为他们行恶，因此他们对主说：\BibleQuote{我们与祢有什么相干？}当你听见魔鬼这样说时，你以为他们不承认祂吗？他们说：\BibleQuote{我们知道祢是谁：祢是天主子。}这话是伯多禄说的，便受称赞；魔鬼说了，却受谴责。这是为什么？岂不是因为言语虽同，心却不同吗？那么，让我们在我们的信德中有所分别，不要满足于相信。这样的信德并不能洁净心灵。经上记载：\BibleQuote{洁净了他们的心}，又说：\BibleQuote{藉着信德}。但藉着怎样的信德呢？便是保禄宗徒所定义的那一种，他说：\BibleQuote{那藉着爱德行事的信德。}这信德使我们与魔鬼的信德以及那些邪恶败坏之人的行为分别开来。他说：\Quote{信德。}什么信德？\BibleQuote{那藉着爱德行事的信德}，并盼望天主所应许的。没有比这定义更精确完全的了。因此，在信德中有这三样。凡拥有那藉着爱德行事的信德的人，必然盼望天主所应许的。所以望德是信德的伴侣。因为我们尚未看见我们所相信的，望德是必要的，以免因不见而绝望以致跌倒。不见使我们忧伤，但盼望得见却安慰我们。因此望德在此，它是信德的伴侣。还有爱德，藉此我们渴望、努力追求、炽热向往、饥渴慕义。这也被包括在内，于是便有了信德、望德和爱德。因为哪里没有爱德呢？爱德不就是爱吗？而信德本身也被定义为\BibleQuote{那藉着爱德行事的}。除去信德，你所信的一切便消亡；除去爱德，你所行的一切便消亡。因为信德的责任是相信，爱德的责任是实行。你若没有爱而信，就不会致力于善行；即便行善，也是如同仆人而非儿子，出于对刑罚的恐惧，而非对正义的热爱。因此我说，那藉着爱德行事的信德才能洁净心灵。\par

12. 这信德如今产生什么效果呢？它藉着圣经如此多的见证、多样的教训、丰富而充盈的劝诫，除了使我们\BibleQuote{如今藉着镜子观看，模糊不清，到那时则要面对面}之外，还做什么呢？但现在不要让你的思绪再回到你肉身的面容上。只思考心灵的面容。强迫、压服、催促你的心灵默想神圣的事。凡出现在你心中像是什么形体的事物，都把它抛开。你若还不能说\Quote{这就是}，至少要说\Quote{这不是}。因为你何时才能说\Quote{这就是天主}呢？即使你看见祂时，也不能；因为那时你将看见的是不可言传的。正如宗徒所说，他被\BibleQuote{提到第三层天上，听到了不可言传的话语}。如果话语是不可言传的，那么说这些话的那位又是怎样的呢？因此，当你想到天主时，也许你心中浮现出某个伟大非凡的人形，你在心灵的眼前将它设想为极其伟大、庄严、广阔的东西。但你仍然在某处为它设定了界限。如果你设定了界限，那便不是天主。但若你没有设限，面容又在哪里呢？你是在幻想一个庞大的身体，要区分其中的肢体，就必须为它设限。因为除了为这庞大的身体设限，你无法区分肢体。但你这愚昧而肉性的想象啊！你制造了一个大而笨重的身体，并且越大，你就以为越尊崇天主。另一个人给它加了一肘，使它比以前更大。\par

13. 但你会说：\Quote{我读过。}你读了什么？你什么也没有明白！不过告诉我，你读了什么？让我们不要推开这理解力尚幼稚的孩子。告诉我，你读了什么？\BibleQuote{天是我的宝座，地是我的脚凳。}我听见了，我也读过；但也许你以为自己占了优势，因为你既读了又相信了。但我也相信你刚才所说的。那么让我们一同相信吧。我说什么？让我们一同查考。看啊！持守你所读所信的：\Quote{天是我的宝座}（即\Quote{我的座位}，因为希腊文的\OriginalTerm{宝座}{throne}在拉丁文是\OriginalTerm{座位}{seat}），\Quote{地是我的脚凳}。但你没有也读过这句话吗：\BibleQuote{谁曾用手掌量过天呢？}我断定你读过；你承认并相信了；因为在那本书里，我们既读了前者也读了后者，并且两者都相信。但现在请想一想，并教导我。我以你为师，以自己为小学生。请教我，请问：\Quote{谁坐在这手掌上呢？}\par

14. 看，你已从人的身体上画出了天主的肢体轮廓和特征。或许你曾想过，我们是按照天主的肖像造的，这肖像在于身体。我暂且承认这个想法，以便讨论、审视和筛检。那么，若你愿意，请听我说；因为你在乐意发言时，我也听了你。天主坐在天上，以他的手掌量度天。什么！同样的天，当作天主的座位时就变宽，当他量度时就变窄吗？或者天主坐着时，受限于他手掌的度量？若是如此，天主并没有按照他的肖像造我们，因为我们手掌比我们坐的部位窄得多。但如果他手掌和他坐着时一样宽，他就使我们的肢体与他的极不相像。这里毫无相似之处。因此，基督徒该羞愧，不在心中树立这样的偶像。因此，将天理解为所有圣人。因为地也指所有在地上的人，有云：\BibleQuote{愿全地都崇拜你。}若我们论及地上居民时可以说：\BibleQuote{愿全地都崇拜你}，我们同样可以说论及天上居民：\Quote{愿整个天都承载你。}因为连住在世上的圣人，虽身体行走地上，心却住在天上。他们被提醒\BibleQuote{举心向上}，并不徒然；当被提醒时，他们回答\BibleQuote{我们已举心向上}；说\BibleQuote{你们既然与基督一同复活，就该追求天上的事，那里有基督坐在天主的右边。你们要思念天上的事，不要思念地上的事}，也不是徒然的。因此，就他们以那里为居所而言，他们承载天主，他们就是天；因为他们是天主的座位；当他们宣告天主的话语时，\BibleQuote{诸天陈述天主的荣耀。}\par

15. 那么，与我一同回到心的面容，并预备\Emph{它}。天主说话的对象是在内在。耳朵、眼睛和其余一切可见的肢体，要么是内在某物的居所，要么是它的工具。正是内里之人，基督现今因信德住在其中，将来他要以天主性的临在住进去，那时我们将知道\BibleQuote{什么是宽、长、高、深；也将知道基督那超越知识的爱，使我们充满天主的一切丰满。}那么，若你想明白这些话的意义，要集合你所有的能力去领悟宽、长、高、深。不要在你的思想想象中漫游于世界的空间，以及这个巨大身躯的可理解的范围。在你自身中寻找我所说的。\Quote{宽}在于善工；\Quote{长}在于行善时的恒久忍耐和恒心；\Quote{高}在于对天上赏报的期望，为了这高的缘故，你被命令\Quote{举心向上}。行善，并因天主的赏报而恒心行善。将地上的事物视为无有，免得当这地被那位智者用任何鞭打打击时，你说你白白崇拜了天主，白白行了善工，白白在善工中坚持了。因为行善工你仿佛有了\Quote{宽}，恒心行善你仿佛有了\Quote{长}；但由于追求地上的事物，你没有\Quote{高}。现在注意\Quote{深}：它是天主的恩宠，在他旨意的隐秘安排中。\BibleQuote{有谁知道主的心意？或是谁作过他的参谋？}以及，\BibleQuote{你的判断如同深渊。}\par

16. 这番关于行善、在行善上坚持、寄望于上天的赏报、以及天主恩宠的隐秘施与的谈论——乃是出于智慧而非愚妄，也不出于挑剔，因为一人如此，另一人如彼；因为\BibleQuote{天主没有不义}；我说，你若以为有益，也将这番谈论应用到你主的十字架上吧。因为祂选择这种死亡并非没有意义，祂既有能力选择死或不死。现在，如果祂有能力选择死或不死，为什么祂不能选择这样或那样的死法呢？因此，祂选择十字架并非没有意义，为的是藉此将你钉死于这个世界。因为\BibleQuote{宽度}是十字架上的横木，双手被钉之处，象征善工。\BibleQuote{长度}是从这横木向下延伸到地面的那部分木头。因为身体被钉在其上，仿佛站立，这站立象征坚持。而\BibleQuote{高度}则是从同一横木向上延伸到头部的那部分，由此象征对天上之事的期待。至于\BibleQuote{深度}，则莫过于埋入地下的那部分吗？因为恩宠的施与正是如此，隐藏而隐秘。它本身不被看见，但从它投射出所有可见之物。此后，当你领悟了这一切——不仅是头脑的理解，更是在行动中（\BibleQuote{因为凡实行者必获得良好的领悟}），那时，你若能，就尽力伸展自己，以达至对那\BibleQuote{基督的爱超越知识}的认识。当你达到后，你将\BibleQuote{充满天主的一切圆满}。那时，\BibleQuote{面对面}将实现。现在，你将充满天主的一切圆满，并非天主充满你，而是你充满天主。你若有能力，在那里寻找任何有形的面容吧。将这些儿戏从心灵的眼中除去吧。让孩子丢弃他的玩具，学习处理更严肃的事情。在许多事上，我们不过是孩子；当我们比现在更像孩子时，我们的长者包容了我们。\BibleQuote{你们要与众人在和平中，并追求圣德，若无圣德，没有人能见天主。}因为圣德洁净心灵；因为那\BibleQuote{藉着爱德工作}的信德就在心里。因此，\BibleQuote{心里洁净的人是有福的，因为他们要看见天主。}\par

\SourceRecord{Translated by R.G. MacMullen.；Nicene and Post-Nicene Fathers, First Series；Vol. 6.；Edited by Philip Schaff.；Buffalo, NY: Christian Literature Publishing Co.,；1888.；<http://www.newadvent.org/fathers/160303.htm>.}

% structure: p0001 source=h1 target=section reason=page-title

\section{新约讲道词第四篇}

\EditorialNote{本笃版第LIV篇}\par

\Emph{论到福音所记载的，}\BibleRef{玛 5:16}\Emph{，}\BibleQuote{照样，你们的光也当在人前照耀，好使他们看见你们的善行，光荣你们在天之父；}\Emph{相反地，第六章，}\BibleQuote{你们应当心，不要在人前显示你们的义德，为叫他们看见。}\par

1. 亲爱的诸位，常使许多人感到困惑的是，我们的主耶稣基督在他的福音讲道中，首先说了：\BibleQuote{你们的光当在人前照耀，使他们看见你们的善行，光荣你们的天父；}后来又说：\BibleQuote{你们应当心，不要在人前显示你们的义德，为叫他们看见。}因为那理解力薄弱的人的心思就受到困扰，渴望服从这两条诫命，却被各种不同、且矛盾的命令分心。因为一个人若接到矛盾的命令，几乎不能服从一个主人，正如他不能服侍两个主人一样——救主自己在此同一讲道中已证明这是不可能的。那么，处于这种犹豫之中的心思当如何行？它认为自己不能，却又害怕不服从？因为若他将自己的善行放在光中，叫人看见，好履行那命令：\BibleQuote{你们的光当在人前照耀，使他们看见你们的善行，光荣你们的天父；}他就会觉得自己有罪，因为他做了与另一条诫命相反的事，那条诫命说：\BibleQuote{你们应当心，不要在人前显示你们的义德，为叫他们看见。}反之，若他因害怕并避免这一点而隐藏自己的善行，他就会认为没有服从那命说：\BibleQuote{你们的光当在人前照耀，使他们看见你们的善行}的那一位。\par

2. 但是，那具有正确理解的人，会履行两者，并在两者中服从万有的普遍之主；若主命令了那些根本无法做到的事，祂就不会责备懒惰的仆人。请听\BibleQuote{耶稣基督的仆人保禄，蒙召作宗徒，被选拔为传天主的福音}——他既实行又教导这两项职责。请看他的\BibleQuote{光怎样在人前照耀，使人看见他的善行。我们将自己}他说，\BibleQuote{在天主面前推荐给众人的良心。}又说：\BibleQuote{我们不仅在天主面前，也在人面前，预备光明的事。}又说：\BibleQuote{在一切事上要使众人喜悦，就如我在一切事上使众人喜悦。}反之，看他如何谨慎，\BibleQuote{不将他的义德行在人前为叫他们看见。各人}他说，\BibleQuote{应当考验自己的行为，这样，他只在自身夸耀，而不在别人身上。}又说：\BibleQuote{我们的夸耀就是良心的见证。}再没有比这更清楚的了：\BibleQuote{若我}他说，\BibleQuote{仍取悦于人，我就不是基督的仆役。}愿主亲自与我们同在，祂也藉祂的仆人和宗徒说话，向我们启示祂的旨意，并赐给我们遵行的能力。\par

3. 福音的话本身带有解释；它们并不封住饥饿之人的口，因为它们哺育了叩门之人的心。人内心的意向、方向和目的，才是应当注意的。因为若有人希望自己的善行被人看见，他在人前设置自己的光荣和利益，并在此中寻求，他就没有履行主就此事所给的任何一条诫命；因为他立即着眼于\BibleQuote{在人前显示他的义德为叫他们看见}；他的光也没有这样在人前照耀，使人看见他的善行，而光荣他在天之父。他愿意自己受光荣，而不是天主；他寻求自己的利益，不爱主的旨意。关于这样的人，宗徒说：\BibleQuote{众人都寻求自己的事，不寻求耶稣基督的事。}因此，句子并未在\BibleQuote{你们的光当在人前照耀，使他们看见你们的善行}处结束，而是立即附加了为何要如此行：\BibleQuote{好使他们光荣你们在天之父}；好使一个行善的人被人看见时，他在自己的良心中只有善行的意向，而无被认识的意向，除非是为了光荣天主，并为了那些得以认识他的人的益处；因为对他们来说，这益处在于：那位赐予人这种能力的天主开始为他们所喜悦；所以他们不会失望，反而希望如果自己愿意，同样的能力也会赐予他们。同样，祂也没有以别的方式结束另一条诫命：\BibleQuote{你们应当心，不要在人前显示你们的义德}，而是附加了\BibleQuote{为叫他们看见}；祂在这一条中也没有加上\BibleQuote{好使他们光荣你们在天之父}，而是说：\BibleQuote{不然，你们在天父前就没有赏报了。}因为祂借此向我们显示，那些祂不愿祂的忠信者成为的那样的人，正是在被人看见这回事上寻求赏报——他们把自己的好处放在其中，在其中满足自己内心的虚荣——这就是他们的空虚、膨胀、浮肿和衰败。因为为什么说：\BibleQuote{你们应当心，不要在人前显示你们的义德}还不够，祂又加上\BibleQuote{为叫你们看见}？除非因为有些人行他们的\BibleQuote{义德在人前}，不是为了被人看见，而是为了善行本身被看见；并且那在天之父，祂曾乐意用这些恩赐装备那不虔敬而祂已使之成义的人，可以受光荣。\par

4. 这样的人既不将自己的义德视为自己的，而视为祂的义德，他们凭着对祂的信德而生活（因此宗徒也说：\BibleQuote{使我获得基督，并得以在祂内，不是有我来自法律的义德，而是有来自基督信德的义德，即来自天主的义德，基于信德}；在另一处说：\BibleQuote{使我们在祂内成为天主的义德}。因此他也这样责备犹太人：\BibleQuote{他们不认识天主的义德，企图建立自己的义德，便不服从天主的义德}）。所以，凡愿意自己的善行被人看见，好使那赐予他们可见之善的天主受到光荣，并使看见的人因信德之虔诚而被激发效法善行的人，他们的光确实在人前照耀，因为从他们身上放射出爱德之光；他们并非出于骄傲的空虚气焰；而在行动中他们也谨慎，不将自己的义德行在人前为叫人看见，他们不把义德视为自己的，也不为让人看见而行义德；而是为使那称义者得以被人认识，祂在蒙成义的人中被赞美，使赞美者在赞美中成为他自己所赞美的对象，即那赞美者自己成为被赞美者。也要注意宗徒，当他说：\BibleQuote{凡事讨众人喜欢，如同我凡事讨众人喜欢}时，他并没有停在这一点上，好像他把讨人喜欢当作自己意向的终点；否则他就会错误地说：\BibleQuote{我若还讨人的喜欢，我就不是基督的仆人}；但他立刻加上他讨人喜欢的原因：\BibleQuote{不求自己的利益，}他说，\BibleQuote{只求众人的利益，使他们得救}。所以他既不为自己的利益而讨人喜欢，免得他不是\BibleQuote{基督的仆人}；同时他为了救恩而讨人喜欢，为成为基督的忠信仆人；因为他的良心在天主面前已足够；而在人面前，从他身上闪耀出可以让人效法的东西。\par

\SourceRecord{Translated by R.G. MacMullen.；Nicene and Post-Nicene Fathers, First Series；Vol. 6.；Edited by Philip Schaff.；Buffalo, NY: Christian Literature Publishing Co.,；1888.；<http://www.newadvent.org/fathers/160304.htm>.}

% structure: p0001 source=h1 target=section reason=page-title

\section{\WorkTitle{讲道第五篇：论新约}}

\Emph{[LV. Ben.]}\par

\Emph{论福音之言}，\BibleRef{玛 5:22}\Emph{，\BibleQuote{凡向自己的弟兄说\InnerQuote{你是傻瓜}的，就有堕入地狱之火的危险。}}\par

1. 我们刚才听到的福音章节，如果我们有信德，必定会深深惊吓我们；但那些没有信德的人，却不被惊吓。因为不惊吓他们，他们便继续处于虚假的安全中，仿佛不知道如何分辨和区别安全与恐惧的适当时刻。那么，现在过着有限生命的人，应当害怕，以便在无尽的生命中获得安全。因此我们被惊吓了。因为谁不会害怕那说真理者呢？祂说：\BibleQuote{凡向自己的弟兄说\InnerQuote{你是傻瓜}的，就有堕入地狱之火的危险。} 然而\BibleQuote{舌头是人不能制服的。} 人制服野兽，却制服不了自己的舌头；他制服狮子，却不约束自己的言语；他制服一切，却制服不了自己；他制服他所害怕的、并应当害怕的事物，为的是制服他自己，使他不再害怕。但这是怎么回事呢？这是一句真实的句子，出自真理之言：\BibleQuote{但舌头是人不能制服的。}\par

2. 那么，我的弟兄们，我们该怎么办？我确实是在向一大群人讲话，但由于我们在基督内是一体，让我们好像在秘密中商议。没有外人听我们，我们都是一体，因为我们都在同一体内。那么，我们该怎么办？\BibleQuote{凡向自己的弟兄说\InnerQuote{你是傻瓜}的，就有堕入地狱之火的危险。} 但是\BibleQuote{舌头是人不能制服的。} 难道所有人都会进入地狱之火吗？绝对不可！\BibleQuote{主啊，世世代代以来，你是我们的避难所。} 你的愤怒是公义的：你不会不义地将任何人投入地狱。\BibleQuote{我要往哪里去躲避你的圣神？} 我又能逃往何处躲避你，除非到你面前？那么，亲爱的诸位，让我们明白：若没有人能制服舌头，我们就必须投奔天主，让祂来制服它。因为如果你愿意制服它，你却不能，因为你是人。\BibleQuote{舌头是人不能制服的。} 请注意一个类似的例子，关于我们所驯服的野兽。马不能驯服自己；骆驼不能驯服自己；大象不能驯服自己；毒蛇不能驯服自己；狮子不能驯服自己；同样，人也不能驯服自己。但是，要使马、牛、骆驼、大象、狮子、毒蛇被驯服，人就被找来。因此，让天主被寻求，好使人被驯服。\par

3. 因此，\BibleQuote{主啊，你成了我们的避难所。} 我们投奔你，靠着你的帮助，我们就会顺利。因为我们靠自己是不顺的。因为我们离开了你，你便将我们留给了自己。愿我们得以在你内被找到，因为我们在自己内迷失了。\BibleQuote{主啊，你成了我们的避难所。} 那么，弟兄们，如果我们把自己交给他来驯服，我们为什么要怀疑主会使我们温顺呢？你驯服了不是你创造的狮子；难道创造了你的他，不会驯服你吗？因为你从哪里得到力量驯服如此凶猛的野兽？你在体力上与他们相等吗？那么，你凭借什么力量能够驯服大野兽？所谓的役畜，按其本性是野性的。因为在未驯服状态，它们是无用的。但因为习俗从未见过它们不在人的手中、缰绳和权力之下，难道你以为它们生来就是驯服的吗？但无论如何，现在注意那些无疑是野生的野兽。\BibleQuote{狮子吼叫，谁不害怕？} 然而，你发现自己在哪方面比它更强？不在体力，而在理智的内在本性。你比狮子更强，是在你按照天主的肖像被造的方面。什么！天主的肖像竟能驯服野兽，而天主不能驯服自己的肖像？\par

4. 我们的望德在于他；让我们顺服他，恳求他的仁慈。让我们把望德寄托于他，直到我们被驯服，彻底被驯服，即被成全，让我们承受我们的驯服者。因为常常这位驯服者也会拿出他的鞭子。因为如果你拿出鞭子驯服你的野兽，难道天主不会这样做来驯服他的野兽（也就是我们）吗？他正要使他的野兽成为他的儿子。你驯服你的马；当它开始温顺地驮着你，承受你的管教，服从你的规则，成为你忠信、有用的野兽时，你给它什么？你如何回报它？当它死后，你甚至不会埋葬它，而是扔出去被猛禽撕碎。然而，当你被驯服，天主为你保留了一份产业，就是天主自己，尽管短暂死亡，他也会使你复活。他将恢复你的身体，甚至你的每一根头发；他将使你永远与天使在一起，在那里你不再需要他的驯服之手，只需被他的极大仁慈所拥有。因为那时天主将是\BibleQuote{万有中的万有；} 不再有不幸来磨炼我们，只有幸福来滋养我们。我们的天主将亲自作我们的牧者；我们的天主将亲自作我们的杯爵；我们的天主将亲自作我们的荣耀；我们的天主将亲自作我们的财富。无论你在此寻求多少事物，只有他将亲自成为你的一切。\par

5. 人是为着这望德而被驯服的，难道驯服他的天主反倒会被视为不可容忍的吗？人是为着这望德而被驯服的，倘若祂偶然使用鞭子，他能埋怨赐恩的驯服者吗？你们已听到了宗徒的劝勉：\Quote{你们若不受惩戒，便是私生子，而不是儿子；因为哪个儿子，父亲不惩戒他呢？再者，} 他说：\Quote{我们有生身的父亲管教我们，尚且敬畏他们；何况灵性的圣父，我们岂不更当顺服祂而得生吗？} 你的生父为你做了什么？他管教、责打你，拿出鞭子抽你，他能使你永远活着吗？他自己都不能为自己做的事，如何能为你做？为了一点他用高利贷和辛劳积攒的微薄钱财，他用鞭子来训练你，免得他劳苦的成果留给你后，被你放荡的生活挥霍掉。是的，他打儿子，是怕他的劳作付诸东流；因为他留给你的，既不能在这里留住，也不能带走。他并没有留给你任何真正属于他自己的东西；他离去，好让你继承。但你的天主、你的救赎者、你的驯服者、你的惩戒者、你的圣父，祂教导你。为的是什么？为使你承受产业，那时你不再需要将你的父亲带往坟墓，而是你的圣父自己成为你的产业。你为这望德受教导，却埋怨吗？若有什么不幸的事临到你，你（也许）还亵渎吗？你将从祂的圣神那里逃到哪里去？但现在祂任由你，并不鞭打你；或者祂在你亵渎时放弃你；难道你就不经历祂的审判吗？祂鞭打你并接纳你，岂不比祂宽容你却放弃你更好吗？\par

6. 让我们对我们主天主说：\Quote{主，你世世代代作我们的避难所。} 在第一代和第二代中，你作了我们的避难所。我们尚未存在时，你作我们的避难所，使我们得以出生。我们是邪恶的，你作我们的避难所，使我们得以重生。你作避难所，喂养那些离弃你的人。你作避难所，扶起并引导你的子女。\Quote{你作了我们的避难所。} 我们不会离开你，当你救我们脱离一切邪恶，并用你自己的美物充满我们时。现在你赐予美物，温柔地对待我们，免得我们在路上疲乏；你管教、责打、击打并引导我们，免得我们偏离道路。因此，无论你温柔地待我们，免得我们在路上疲乏，还是管教我们，免得我们偏离道路，\Quote{主，你作了我们的避难所。}\par

\SourceRecord{Translated by R.G. MacMullen.；Nicene and Post-Nicene Fathers, First Series；Vol. 6.；Edited by Philip Schaff.；Buffalo, NY: Christian Literature Publishing Co.,；1888.；<http://www.newadvent.org/fathers/160305.htm>.}

% structure: p0001 source=h1 target=section reason=page-title

\section{新约讲道集 第六篇}

[LVI. Ben.]\par

论\BibleRef{玛 6:9}中的主祷文，等，致\OriginalTerm{Competentes}{Competentes}。\par

1. 蒙福的宗徒为显示万民信基督的时候已由先知预言，引用了这见证，其中记着说：\BibleQuote{凡呼号主名的人，必能得救。}因为以前，唯独在以色列人中呼号那创造天地的主的名；其余民族却呼号不能言的聋哑偶像，它们不能垂听，或呼号魔鬼，他们虽得应允，却反被加害。\BibleQuote{但是时期一满}，那预言便应验了：\BibleQuote{凡呼号主名的人，必能得救。}再者，因为犹太人，甚至那些信基督的，对福音心怀嫉妒，不肯传给外邦人，说福音不应传给未受割礼的人；宗徒保禄反对他们，引用了这个见证：\BibleQuote{凡呼号主名的人，必能得救。}他随即又加上这些话，为说服那些不愿让福音传给外邦人的人：\BibleQuote{然而，人若不信他，怎能呼号他呢？人若未听说过他，怎能信他呢？若没有宣讲者，人怎能听到呢？若没有被派遣，人怎能宣讲呢？}因为他既说：\Quote{人若不信他，怎能呼号他呢？}你们便不是先学习主祷文，后学习信经；而是先学习信经，好知道该信什么，然后学习祷文，好知道该呼号谁。那么，信经关乎信德，主祷文关乎祈祷；因为相信的人呼号时，才蒙垂听。\par

2. 但许多人祈求不当求的，不知何为有益。因此，祈祷的人必须提防两件事：不要求不当求的，也不要向不当求的对象求。绝不可向魔鬼、偶像、邪神求任何事。我们必须向我们的主耶稣基督、向先知、宗徒和殉道者的天主父、向我们的主耶稣基督的父、向创造天地、海洋和其中万物的天主，祈求一切当求的。但我们要提防，不可向他求不当求的事。若是我们应当求生命，你却向不能言的聋哑偶像求，这对你有什么益处呢？同样，若你向天上的圣父祈求仇人的死亡，这对你有什么益处呢？难道你没有在圣咏中读过或听过，其中预言了叛徒犹大该死的结局，预言如何论及他说：\BibleQuote{愿他的祈祷变为罪恶吗？}如果这样，你起来为你的仇人求恶，你的\BibleQuote{祈祷便将变为罪恶。}\par

3. 你在圣咏中读到，那在圣咏中说话的人似乎向他的仇人说了许多咒诅的话。诚然，有人会说，在圣咏中说话的人是个义人；那他又为何如此恶愿他的仇人呢？他不是愿，而是预见；这是预言者在讲说将来的事，不是出于诅咒的愿望；因为先知们因圣神知道谁将遭恶，谁将得善；借着预言，他们说话就像愿意他们所预见的事一样。但你怎能知道，你今天为之求恶的人，明天不会变得比你更好呢？但你会说：我知道他是个恶人。好吧，你也必须知道你也是个恶人。纵然你或许自以为能判断他人的心，但对自己，你\Emph{知道}自己是恶人。难道你没有听见宗徒说：\BibleQuote{我从前是亵渎者、迫害者和施暴者；但是我蒙受了怜悯，因为我是在不信中无知而做的。}那时宗徒保禄迫害基督徒，到处捆绑他们，把他们带到司祭长面前受审和受罚；弟兄们，你们认为教会是\Emph{反对}他祈祷，还是\Emph{为他}祈祷呢？的确，天主的教会从她的主那里学到了教训，主在十字架上说：\BibleQuote{父啊，宽恕他们吧，因为他们不知道他们做什么。}同样，她为保禄（那时还是扫禄）祈祷，使那所成就的在他身上得以成就。因他说：\BibleQuote{然而，我在基督内的犹太各教会，只是听闻而已；他们只听说：那个从前迫害我们的人，如今在传扬他曾经破坏的信德；他们便因我而光荣了天主。}他们为何光荣天主呢？岂不是因为他们曾向天主祈求这事，而后它成就了吗？\par

4. 我们的主首先禁止了\Quote{多言}，免得你向天主说许多话，好像你能用许多话去教导他。因此，你祈祷时需要的是虔诚，而不是多话。\BibleQuote{因为你们的父在你们祈求他以前，已知道你们需要什么。}所以不要多说，因为他知道你们需要什么。但恐怕有人会说：既然他知道我们需要什么，我们为什么还要说哪怕几句话？为什么还要祈祷呢？他自己知道；就让他按他所知需要的赐给我们吧。是的，但他愿意你祈祷，好使他赐予你的渴望，使他的恩赐不被轻视；因为他亲自在我们里面形成了这渴望。因此，我们的主耶稣基督在他的祷文中教导我们的那些话，是我们渴望的规则和标准。除了那里所写的，你不可祈求任何事物。\par

5. \BibleQuote{所以，你们要说，}祂说，\BibleQuote{我们的天父，你在天上。}在此，你们看到你们已开始以天主为父。你们将在重生时拥有祂。尽管在你们出生之前，你们已因祂的种子而受孕，仿佛即将在洗礼池——教会之母的子宫中——诞生。\BibleQuote{我们的天父，你在天上。}因此要记住，你们有一位在天上的父。记住你们由你们的父亲亚当出生归于死亡，你们将由天主父重生归于生命。你们所说的，要在心里说。只要祈祷有真挚的热忱，那垂听祈祷者必有有效的回应。\BibleQuote{愿你的名被尊为圣。}你们为何祈求天主的名字被尊为圣？它本身就是神圣的。那么你们为何祈求那已经是神圣的呢？再者，当你们祈求祂的名被尊为圣时，你们不就是在为祂祈祷，而非为自己吗？不。正确理解它，你们是在为自己祈求。因为你们所祈求的，是那本身常是神圣的，得以在你们内被尊为圣。什么是\Quote{被尊为圣}？\BibleQuote{被视为神圣，}不被轻视。可见，你们所愿的善，是愿为自己。因为若你们轻视天主的名字，对你们自己是不好的，对天主则不然。\par

6. \BibleQuote{愿你的国来临。}我们向谁说话？若我们不求，天主的国就不会来吗？因为我们所说的是世界终结后的那个国。因为天主常有一个国；祂从不缺乏一个国，万物都服事祂。但你们所愿的是哪个国呢？是福音中记载的那个：\BibleQuote{你们这些蒙我父祝福的，来承受那从创世以来为你们预备的国吧。}看，这就是我们说\BibleQuote{愿你的国来临}的那个国。我们祈祷它来临\Emph{在我们内}；我们祈祷我们能被发现\Emph{在其中}。因为它必定会来；但若它发现你们在左边，对你们有何益处呢？因此，再次地，你们是在为自己求善；你们为自己祈祷。这是你们所渴望的；在祈祷中渴求这事，使你们如此生活，得以在要赐给众圣人的天主的国有份。所以当你们说\BibleQuote{愿你的国来临}时，你们为自己祈祷，使你们能好好生活。愿我们在你的国有份：愿那要来临到你的圣人和义人身上的国，也来临到我们。\par

7. \BibleQuote{愿你的旨意奉行。}什么！若你们不说这话，天主就不行祂的旨意吗？记住你们在信经中所宣认的：\Quote{我信全能的天主父。}若祂是全能者，你们为何祈求祂的旨意得以奉行？那么这是什么意思呢？\BibleQuote{愿你的旨意奉行}？愿它在我内奉行，使我不抗拒你的旨意。因此，再次地，你们是为自己祈祷，而非为天主。因为天主的旨意将在你们内奉行，即使不由你们奉行。因为对祂将说\BibleQuote{你们这些蒙我父祝福的，来承受那从创世以来为你们预备的国吧}的人，天主的旨意将要奉行，使圣人和义人承受国度；而对祂将说\BibleQuote{你们这些被咒骂的，离开我，到那为魔鬼及其使者预备的永火里去吧}的人，天主的旨意也要奉行，使恶人被定罪于永火。而祂的旨意由你们奉行，是另一回事。那么你们祈求祂的旨意得以在你们内奉行，并非没有原因，而是为了你们的好处。但无论你们是好是坏，它仍将在你们内奉行；但但愿它也由你们奉行。那么我为何说\BibleQuote{愿你的旨意奉行在天上，也在地上}，而不说\Quote{愿你的旨意由天上和地上奉行}呢？因为由你们所行的，是祂自己在你们内所行的。从来没有任何事由你们所行，是祂不在你们内所行的。有时，祂确实在你们内行那不由你们所行的事；但若祂不在你们内行，就绝没有事由你们所行。\par

8. 但\BibleQuote{在天上，也在地上}，或\Quote{如同在天上，也在地上}是什么意思？天使奉行你的旨意；愿我们也如此奉行。\BibleQuote{愿你的旨意奉行在天上，也在地上。}心灵是天，肉体是地。当你们与宗徒一同说（如果你们确实这样说）：\BibleQuote{以心灵我服事天主的法律，但以肉体我服事罪的法律}时，天主的旨意在天上奉行，但尚未在地上。但当肉体与心灵和谐，\BibleQuote{死亡被胜利吞没}时，不再有肉性的欲望使心灵与之冲突，地上的争战止息，内心的战争结束，那话\BibleQuote{肉欲相反圣神，圣神相反肉欲；它们彼此对立，致使你们不能做你们所愿意的事}所表达的成为过去；这场战争，我说，结束时，一切贪欲都变为爱德，就没有什么在肉体里反对心灵，没有什么需要驯服、束缚、践踏；反而一切通过和谐走向正义，天主的旨意将在天上和地上奉行。\BibleQuote{愿你的旨意奉行在天上，也在地上。}我们祈求圆满，当我们为此祈祷时。\BibleQuote{愿你的旨意奉行在天上，如同在地上一样。}在教会中，属神的人是天上，属血气的（carnal）人是地上。因此，\BibleQuote{愿你的旨意奉行在天上，如同在地上一样}，以致正如属神的人服事你，属血气的人经过改造也服事你。\BibleQuote{愿你的旨意奉行在天上，如同在地上一样。}还有一个更属神的意义。因为我们被劝勉要为仇敌祈祷。教会是天，教会的仇敌是地。那么\BibleQuote{愿你的旨意奉行在天上，如同在地上一样}是什么意思？愿我们的仇敌相信，正如我们也相信你！愿他们成为朋友，终结敌意！他们是地，所以反对我们；愿他们成为天，他们将与我们同在。\par

9. \Quote{求祢今天赏给我们日用的食粮。}如今显然，我们是为自己祈祷。当你说：\Quote{愿祢的名被尊为圣，}便需解释为何这是为自己而非为天主祈祷。当你说：\Quote{愿祢的旨意奉行，}又需解释，免得你以为在这祈祷中是为天主求好处，愿祂的旨意成就，而非为自己祈祷。当你说：\Quote{愿祢的国来临，}同样需解释，免得你以为在这祈祷中是祝愿天主为王。但从此处直至祈祷结尾，显而易见的，是我们为自己向天主祈祷。当你说：\Quote{求祢今天赏给我们日用的食粮，}你便承认自己是天主的乞丐。但不要为此羞愧；无论世上何人何等富有，他仍是天主的乞丐。乞丐站在富人家门前，但富人自己却站在大富者的门口。有人向他祈求，他自己也祈求。他若不匮乏，就不会在祈祷中叩击天主的耳朵。富人需要什么？我冒昧地说，富人甚至也需要日用的食粮。他何以拥有一切？岂非因天主赐予了他？若天主收回祂的手，他还能有什么？多少人夜卧于财富，晨起于贫穷？他之缺乏，赖天主仁慈，非己力也。\par

10. 但这食粮，亲爱的诸位，滋养我们的身体，日复一日恢复肉身；这食粮，我说，天主不仅赐给赞美祂的人，也赐给亵渎祂的人；\BibleQuote{祂使太阳升起，光照恶人也光照善人；降雨给义人也给不义的人。}你赞美祂，祂喂养你；你亵渎祂，祂仍喂养你。祂等待你悔改；但你若不改变自己，祂将定你的罪。既然善恶之人都从天主得到这食粮，你以为就没有另一种食粮，就是儿女所祈求的，主在福音中所说的：\Quote{拿儿女的饼扔给小狗，是不对的。}？的确是有。那么那食粮是什么？为何称为日用的？因为这食粮如那食粮一样必要；因为没有它我们不能生活；没有饼我们不能生活。向天主求财富是无耻，求日用食粮却非无耻。助长骄傲的是一回事，维持生命的另一回事。然而，因为这可见可触的饼赐给善恶之人，便有一种日用的食粮，是儿女所祈求的。那就是天主的圣言，日复一日分施给我们。我们的饼是日用之粮；靠它生活的非我们的身体，而是我们的灵魂。对我们这些现今在葡萄园里劳作的工人，它是必要的——它是我们的食物，而非工价。因为雇佣工人进葡萄园的人欠他两样东西：食物，以免他晕倒；工价，使他能喜乐。我们在世上的日用食物就是天主的圣言，常在教会中分施；劳苦之后的工价称为永生。再者，若你以此日用食粮领会信友所领受的，即你们受洗后将要领受的，那么我们祈求\Quote{求祢今天赏给我们日用的食粮}是合理的，好使我们如此生活，不致与圣祭台分离。\par

11. \Quote{求祢宽恕我们的罪债，如同我们也宽恕得罪我们的人。}关于这条祈求，我们无需解释我们是为自己祈祷。因为我们请求我们的罪债获得宽恕。我们是负债者，不是金钱上的，而是罪恶上的。此刻你也许在说：你也是。我们回答：是的，我们也是。什么，你们这些圣主教，也是负债者？是的，我们也是负债者。怎么，我主！万勿如此，不要这样委屈自己。我并未委屈自己，我说的是实话：我们是负债者：\BibleQuote{若我们说我们没有罪，便是欺骗自己，真理也不在我们内。}我们受过洗礼，但仍是负债者。不是说洗礼中有什么未得赦免，而是因为我们生活在不断积累需要每日宽恕的事。那些受洗后立即离世的人，从洗礼池中出来时毫无罪债；毫无罪债离开世界。但那些受洗并仍停留于此生的人，由于有死的软弱而沾染污秽，虽不致沉船，却需要抽水。否则点滴渗入，终致全船沉没。献上这祈祷，就是使用抽水机。但我们不仅应祈祷，还应施舍，因为当抽水以防船沉时，呼声和双手都在工作。我们用呼声工作，便是说：\Quote{求祢宽恕我们的罪债，如同我们也宽恕得罪我们的人。}我们用双手工作，便是践行：\Quote{把你的饼分给饥饿的人，将无家可归的穷人领进你的家。将施舍封在穷人的心中，它必为你向主祈祷。}\par

12. 虽然我们在\BibleQuote{重生之洗}中一切罪过都已获赦免，但若没有主祷文每日的洁净，我们仍将被逼入极大困境。施舍和祈祷涤除罪恶；只是不要犯那些必须使我们与日用食粮分离的罪；避免一切那会导致严厉而确定定罪的债。不要自称义人，好像你们没有理由说：\BibleQuote{宽免我们的罪债，如同我们也宽免我们的债主。}虽然你们避免拜偶像、占星者的安慰、巫师的医治，避免异端者的诱惑、分裂者的分裂，避免凶杀、奸淫和淫乱、偷盗和抢劫、假见证，以及一切其他我不提名且具有毁灭性后果的罪，为此罪人必须与祭台分离，在地上被捆绑，在天上也被捆绑，除非他再次在地上被解开，在天上也被解开——这对他有极大而致命的危险。但在排除这一切之后，人犯罪的机会仍不少。人在看见不应看见的东西而喜悦时犯罪。然而谁能约束眼睛的快速？据说眼睛正是从其快速得名。谁能约束耳朵或眼睛？眼睛可按你的意愿闭上，且瞬间闭上，但耳朵只能费力地关闭；如果有人抓住你的手，它们就保持开着；你也不能关起它们来抵挡辱骂、不洁、谄媚和诱惑的话语。当你听见不应听见的事物时，即使你没有做，你岂不是以耳朵犯罪吗？因为你以喜悦听坏的事物？致死的舌头犯下多大的罪！是的，有时是这种性质的罪，人因此与祭台分离。一切亵渎之事都属于舌头，还说许多无益的空话。但让手不做错事，脚不跑向邪恶，眼不注视不端庄，耳不以喜悦听污秽之言，舌不说不雅之词；然而告诉我，谁能约束思想？我们祈祷时多少次，弟兄们，心思却在他处，好像忘了我们站在谁面前，或俯伏在谁面前！如果所有这些都聚集起来攻击我们，难道它们不会因是小过犯而不压倒我们吗？压垮我们的是铅还是沙子有什么关系？铅是一整块，沙是小颗粒，但数量多时它们压垮你。你的罪过也是如此微小。你难道看不见河流是如何被点滴之水注满，土地是如何被点滴之水荒废吗？它们是小，但数量众多。\par

13. 因此让我们每天说，并真心实意地说，且照我们所说的去做：\BibleQuote{宽免我们的罪债，如同我们也宽免我们的债主。}这是与天主的一个约定、一项盟约、一份协议。你的主天主对你说：宽免，我就宽免。你没有宽免；你保留罪过反对你自己，不是我。我恳求你们，我亲爱的孩子们，因为我深知在主的祷文中什么对你们有益，尤其在那句话中：\BibleQuote{宽免我们的罪债，如同我们也宽免我们的债主；}请听我说。你们即将受洗，宽免一切；无论任何人心中对别人有什么不满，让他从心里宽免它。这样进入，并确信你们所沾染的一切罪过，无论是从父母出生在亚当之后而有的原罪（为此罪你们同婴孩一起奔向救主的恩宠），还是此后在生活中因言语、行为或思想所染的任何罪过，都已全部宽免；你们将走出水，如在你们的主面前，带着一切债务的确实豁免。\par

14. 如今，由于那些我提到的日常罪过，你们有必要在那每日的、如同洁净一般的祈祷中说：\Quote{宽免我们的罪债，如同我们也宽免得罪我们的人}\BibleRef{玛 6:12}；你们将怎么做？你们有仇人。因为谁能在世上没有仇人呢？你们要谨慎，要爱他们。你们的仇人绝不能因他的暴力伤害你们，如同你们因不爱他而伤害自己一样。因为他可能损害你们的产业、牲畜、房屋、男仆、女仆、儿子或妻子；或者，如果他被授予那种权柄，至多伤害你们的身体。但他能像你们自己伤害灵魂那样伤害你们的灵魂吗？亲爱的弟兄们，我恳求你们，努力达到这成全的境界。但我给了你们这种能力吗？只有那一位给了你们，你们对祂说：\Quote{愿祢的旨意奉行在人间，如同在天上}\BibleRef{玛 6:10}。然而，不要让这事在你们看来不可能。我知道，我凭经验知道，有些基督徒确实爱他们的仇人。如果你们认为不可能，你们就不会去做。因此，首先要相信这是可以做到的，并祈求天主的旨意能在你们身上成就。因为你们的邻人作恶对你们有什么好处呢？如果他没有恶，他甚至不会成为你们的仇人。那么，你们要为他祝福，好使他结束他的恶，他就不再是你们的仇人了。因为与你们为敌的不是他里面的人性，而是他的罪。难道他因为是你们的仇人，就因为他有灵魂和身体吗？在这方面，他与你们一样：你们有灵魂，他也有；你们有身体，他也有。他与你们是同一性体；你们都是出于同一尘土，并由同一主所赋予生命。在这一切中，他与你们一样。因此，要在你们里面承认他是你们的弟兄。第一对夫妇，亚当和厄娃，是我们的父母；一个是我们父亲，另一个是我们母亲；因此我们是弟兄。但让我们放下对我们最初起源的考虑。天主是我们的父，教会是我们的母，因此我们是弟兄。但你们会说：我的仇人是外邦人、犹太人、异端者，关于这些人，我从前在\Quote{愿祢的旨意奉行在人间，如同在天上}\BibleRef{玛 6:10}这话上说过。啊，教会，你的仇人是外邦人、犹太人、异端者；他就是 \Emph{尘世}。如果你们是天上，就要呼求你们在天上的父，并为你们的仇人祈祷；因为扫罗也曾是教会的仇敌；就这样为他祈祷，他就成了教会的朋友。他不仅不再迫害教会，而且还努力帮助她。然而，说实话，祈祷是反对他的；但反对的是他的恶毒，而不是他的本性。因此，让你们的祈祷反对你们仇人的恶毒，好使恶毒死去，而他能活着。因为如果你们的仇人死了，你们似乎失去了一个仇人，却并未找到一个朋友。但如果他的恶毒死了，你们就同时失去了一个仇人，找到了一个朋友。\par

15. 但你们仍在说：谁能做到，谁曾做到过？愿天主使这成就在你们心中！我与你们一样知道，只有少数人这样做；这样做的人是伟大的人，是属神的人。难道教会中所有走近祭台、领受基督身体和宝血的信徒都是这样吗？但他们都说：\Quote{宽免我们的罪债，如同我们也宽免得罪我们的人}\BibleRef{玛 6:12}。如果天主回答他们：\Quote{你们为什么求我做我所应许的事，而你们却不做我所命令的事？}我曾应许什么？\Quote{宽免你们的罪债。}我曾命令什么？\Quote{你们也要宽免得罪你们的人。}如果你们不爱你们的仇人，怎能这样做？那么，弟兄们，我们必须做什么？基督的羊群已经减少到如此稀少的数目吗？如果只有那些爱仇人才应当说\Quote{宽免我们的罪债，如同我们也宽免得罪我们的人}\BibleRef{玛 6:12}，我不知道该怎么办，我不知道该说什么。因为我必须对你们说：如果你们不爱仇人，就不要祈祷；我不敢这样说；是的，倒不如祈祷，使你们能爱他们。但我必须对你们说：如果你们不爱仇人，就不要在主祷文中说\Quote{宽免我们的罪债，如同我们也宽免得罪我们的人}\BibleRef{玛 6:12}？假设我说：不要用这些话。如果你们不说，你们的罪债就不得赦免；如果你们说了，却不按此而行，它们也不得赦免。因此，为使它们得以赦免，你们必须既使用这祷词，又按此而行。\par

16. 我看见一个理由，可以安慰的不仅是一些人，而是众多基督徒：我知道你们渴望听到它。基督说过：\Quote{宽恕，好使你们得到宽恕。} 而我们刚才讨论的祷词中你们说什么？\BibleQuote{宽免我们的罪债，犹如我们也宽免得罪我们的人。} 所以，主，宽恕我们，如同我们宽恕人。你们说：\BibleQuote{我们的天父，求祢宽免我们的罪债，犹如我们也宽免得罪我们的人。} 因为你们应当这样做，如果不这样做，你们必将丧亡。当你的仇敌请求宽恕时，立刻宽恕他。这对你们来说难道算多吗？虽然爱你在暴怒中的仇敌是多的，但爱一个在你面前乞求的人算多吗？你有什么可说的？他先前暴怒，那时你恨他。我倒宁愿你那时也不恨他：我宁愿在你遭受他暴怒的时候，你记起了主说：\BibleQuote{父啊，宽恕他们，因为他们不知道他们做的是什么。} 我多么希望，即使在那时你的仇敌向你施暴时，你已留意到你的主天主这样说了。但也许你会说：祂做到了，但祂是主、是基督、是天主子、是独生子、是圣言成了血肉。但我这个软弱有罪的人能做什么呢？如果你的主对你来说是个太高的榜样，就转念想想你的同仆。圣斯德望被人用石头砸死，当他们用石头砸他时，他跪下为仇敌祈祷说：\BibleQuote{主，不要将这罪归于他们。} 他们扔石头，不是请求宽恕，但他仍为他们祈祷。我希望你像他一样；向前迈进。为什么你总是让你的心拖曳在地上？听：\Quote{请举起你们的心，} 向前迈进，爱你们的仇敌。如果你不能爱他在暴怒中，至少在他请求宽恕时爱他。爱那对你说\Quote{兄弟，我犯了罪，请宽恕我。}的人。如果你那时不宽恕他，我不单是说，你把这祷词从你心中抹去，而且你自己也要从天主的生命册上被抹去。\par

17. 但如果那时你至少宽恕了他，或从心中放下了仇恨，我所要你放下的，是从心中来的仇恨，并非正当的惩戒。如果请求我宽恕的人是我理应惩戒的人，那该怎么办？随你处置，因为我想，即使你惩戒你的孩子，你也是爱他的。你不在乎他在棍杖下的哭喊，因为你为他保留着产业。我说的是，当你的仇敌向你请求宽恕时，你要从心中放下一切仇恨。但也许你会说：\Quote{他是在装假，他是在伪装。} 哦，你这评判人心的，告诉我你父亲的想法，告诉你自己昨天的想法。他请求并恳求宽恕；宽恕他，无论如何宽恕他。如果你不宽恕他，你伤害的是你自己，不是他，因为他知道该做什么。你不愿意宽恕你的同仆；他就会到你的主那里，对祂说：\Quote{主，我恳求我的同仆宽恕我，但他不肯；求祢宽恕我。} 难道主没有权能释放祂仆人的罪债吗？因此，他从主那里获得宽恕，得到释放回来，而你却仍被捆绑。如何被捆绑？祈祷的时刻会来到，到那时你不得不说\BibleQuote{宽免我们的罪债，犹如我们也宽免得罪我们的人。} 而主会回答你：你这恶仆，当你欠我这么大一笔债时，你求了我，我就宽免了你；\BibleQuote{难道你不该怜悯你的同伴，如同我怜悯了你一样？} 这些话出自 \WorkTitle{福音}，不是出于我自己的心。但如果在被请求时，你宽恕了那乞求宽恕的人，那么你就可以念这个祈祷。如果你还没有力量爱他在暴怒中，你仍然可以献上这个祈祷：\BibleQuote{宽免我们的罪债，犹如我们也宽免得罪我们的人。} 让我们继续其余部分。\par

18. \BibleQuote{不要让我们陷于诱惑。宽免我们的罪债，犹如我们也宽免得罪我们的人，} 我们这样说，是因为过去的罪，我们不能撤销它们，使它们没有发生。你可以努力不再做从前做过的事，但你如何能使已做的事不发生？关于已经做过的事，这句祈祷对你有帮助：\BibleQuote{宽免我们的罪债，犹如我们也宽免得罪我们的人。} 至于你可能会陷入的，你该怎么办？\BibleQuote{不要让我们陷于诱惑，但救我们脱离凶恶。} \BibleQuote{不要让我们陷于诱惑，但救我们脱离凶恶}，也就是说，脱离诱惑本身。\par

19. 现在，这三个最先的祈求：\BibleQuote{愿祢的名被尊为圣，愿祢的国来临，愿祢的旨意奉行在人间，如同在天上}，这三项关乎永生，因为天主的名字应常在我们内被尊为圣，我们应常在他的国里，我们应常行他的旨意。这将直到永远。但\Quote{每日的食粮}现在是必需的。从这一项我们所祈求的一切，都关乎现世生活的所需。每日的食粮在此生是必需的；宽免罪债在此生是必需的。因为当我们到达另一个生命时，所有罪债都终结了。在此生有诱惑，在此生航行是危险的，在此生总有一些东西通过我们脆弱的缝隙渗进来，必须被抽出去。但是当我们成为与天使平等时，就不再需要向天主说及祈祷宽免我们的罪债，因为那时将没有罪债。所以，这里是\Quote{每日的食粮}；这里是祈祷我们的\Quote{罪债得以宽免}；这里是\Quote{不陷于诱惑}；因为在那个生命中诱惑不会进入；这里是\Quote{脱离凶恶}；因为在那个生命中将没有凶恶，只有永恒而持久的善。\par

\SourceRecord{Translated by R.G. MacMullen.；Nicene and Post-Nicene Fathers, First Series；Vol. 6.；Edited by Philip Schaff.；Buffalo, NY: Christian Literature Publishing Co.,；1888.；<http://www.newadvent.org/fathers/160306.htm>.}

% structure: p0001 source=h1 target=section reason=page-title

\section{关于新约的讲道 第七篇}

[第五十七篇 圣奥斯定]\par

再论\BibleBook{玛窦}{福音}第六章，论主祷文\BibleRef{玛 6:5}。致\OriginalTerm{慕道者}{Competentes}。\par

1. 为造就你们而确立的次序，要求你们先学习当信什么，然后再学习当求什么。因为宗徒说过：\BibleQuote{凡呼号主名的人，必然得救。}\BibleRef{罗 10:13} \Person{圣保禄}从先知引用了这证言；因为先知预言了那些时代，那时所有人将呼号天主；\BibleQuote{凡呼号主名的人，必然得救。} 他又补充说：\BibleQuote{然而，人若不信他，怎能呼号他呢？从未听过他，怎能信他呢？没有宣讲者，又怎能听到呢？除非被派遣，又怎能宣讲呢？}\BibleRef{罗 10:14-15} 因此派遣了宣讲者。他们宣讲基督。他们宣讲时，百姓听了，听了而信，信了而呼号他。因为那时说得最正确、最真实：\Quote{人若不信他，怎能呼号他呢？} 所以你们先学习了当信什么；今天又学习了呼号你们所信的那一位。\par

2. 天主子、我们的主耶稣基督，教给了我们一篇祷文；虽然他自己就是主，正如你们在信经中所听并诵念的，天主的独生子，但他不愿独自一人。他是独生子，却不愿独自一人；他屈尊有了兄弟。因为他向谁说了：\BibleQuote{你们要说：我们的天父。}\BibleRef{玛 6:9} 他要我们称谁为父，难道不是他的父吗？他嫉妒我们吗？父母有时在生了一个、两个或三个孩子后，害怕再生更多，免得使其他孩子穷乏。但是因为他应许我们的产业是这样一种产业，许多人可拥有，且无人会窘迫；所以他召叫万民进入他的兄弟团体；独生子有数不清的弟兄，他们都说：\BibleQuote{我们的天父。} 前人们这样说了；以后的人们也要这样说。看，独生子在他的恩宠中有多少弟兄，他与他为之而死的人分享他的产业。我们在世上曾有父亲和母亲，为使我们生来就劳苦和死亡；但我们找到了另一位父母：天主是我们的父亲，教会是我们的母亲，我们由他们而生，得享永生。所以，亲爱的，让我们想一想我们开始成为谁的孩子；让我们活出配得上有这样一位父亲的生活。看，我们的造物主竟屈尊做了我们的父亲！\par

3. 我们已经听过该呼求谁，以及怀着怎样获得永业的希望，我们开始有一位在天之父；现在让我们听听该向他求什么。这样的一位父亲，我们该求什么呢？难道我们不是求他降雨，今天、昨天和前天吗？向这样的一位父亲求这些算不得什么大事，然而你们看见，当我们害怕死亡、害怕那无人能逃脱的事时，我们是带着多少叹息和多大的渴望求雨。因为人迟早都要死，我们呻吟、祈祷、痛苦挣扎，向天主呼求，使我们晚一点死。那么，我们更应该向他呼求，使我们能到达那永远不死的地方！\par

4. 因此经上说：\BibleQuote{愿你的名被尊为圣。}\BibleRef{玛 6:9} 我们也求他，使他的名在我们内被尊为圣；因为他常是圣的。他的名怎样在我们内被尊为圣呢？除非它使我们成圣。因为我们曾经不是圣洁的，由他的名我们成为圣洁；但他常是圣的，他的名常是圣的。我们是为自己祈祷，不是为天主。因为我们不会希望天主好，他不可能遭遇任何坏事。但我们希望自己得好，使他的圣名被尊为圣，使那常是圣的在我们内被尊为圣。\par

5. \BibleQuote{愿你的国来临。}\BibleRef{玛 6:10} 它必定来临，无论我们是否祈求。的确，天主有一个永恒的国度。因为他何时不统治呢？他何时开始统治呢？他的国无始无终。但为使我们知道在这祷文中我们也是为自己祈祷，而不是为天主（因为我们不说\Quote{愿你的国来临}，好像我们求天主统治）；我们自身将是他国，如果我们信他并在信德中进步。所有忠信者，被他的独生子之血所救赎的，将是他的国。这他的国将在死人复活时来临；因为那时他自己将来临。当死人复活时，他将分开他们，正如他自己所说：\BibleQuote{他把一些人放在右边，一些人放在左边。}\BibleRef{玛 25:32-33} 对右边的人，他将说：\BibleQuote{我父所祝福的，你们来吧！承受国吧。}\BibleRef{玛 25:34} 这就是当我们说\Quote{愿你的国来临}时所希望和祈求的：愿它临于我们。因为如果我们是被弃绝的，那国将临于别人，不临于我们。但如果我们属于那个数目，属于他独生子的肢体，他的国必临于我们，且不会迟延。因为还有多少世代存留，如同已经过去的呢？宗徒\Person{若望}曾说：\BibleQuote{我的孩子们，现在是末时了。}\BibleRef{若一 2:18} 但这是一个长久的时辰，与这长久的白日相称；看，这末时持续了多少年。但尽管如此，你们要像警醒的人那样，睡觉、复活、统治。让我们现在警醒，在死亡中安眠；最终我们将复活，并永远统治。\par

6.\BibleQuote{愿你的旨意奉行在人间，如同在天上。}我们祈求的第三项是，愿祂的旨意奉行在人间，如同在天上。在此我们也为自己求福。因为天主的旨意必定实现。天主的旨意是善人得福，恶人受罚。难道这旨意会不实现吗？但我们说\BibleQuote{愿你的旨意奉行在人间，如同在天上}时，我们为自己求什么福呢？请听。因为这个祈求可以有多种理解，在这个祈求中有许多事应当在我们心中思考，当我们向天主祈祷时说：\BibleQuote{愿你的旨意奉行在人间，如同在天上。}如同祢的天使不得罪祢，愿我们也不得罪祢。再者，如何理解\BibleQuote{愿你的旨意奉行，如同在天上，也在地上}？所有圣祖、所有先知、所有宗徒、所有属神的人，就如天主的天空；而我们与他们相比就是地。\BibleQuote{愿你的旨意奉行，如同在天上，也在地上}：如同在他们身上，也愿在我们身上。再者：\BibleQuote{愿你的旨意奉行，如同在天上，也在地上}：天主的教会是天，祂的仇敌是地。我们祝愿我们的仇敌得好，使他们也相信而成为基督徒，这样天主的旨意就奉行在人间，如同在天上。再者：\BibleQuote{愿你的旨意奉行，如同在天上，也在地上}：我们的心神是天，肉体是地。正如我们的心神因相信而更新，愿我们的肉体也因复活而更新；这样\BibleQuote{天主的旨意奉行，如同在天上，也在地上}。再者，我们的理智，我们藉以看见真理并喜悦这真理，这就是天；如经上说：\BibleQuote{我按内在的人，喜悦天主的法律。}地是什么呢？\BibleQuote{我看出在我的肢体中另有一种法律，与我的理智法律交战？}当这场争战过去，肉体与心神达到完全和谐时，天主的旨意就要奉行如同在天上，也在地上。当我们重复这个祈求时，让我们思想这一切，并向圣父求这一切。亲爱的诸位，我们刚才提到的这一切——这三项祈求——都是关乎永生的。因为如果我们天主的名在我们内被尊为圣，那将是永远的。如果祂的国来临，我们将永远活着，那将是永远的。如果祂的旨意奉行在人间，如同在天上，以我所解释的一切方式，那将是永远的。\par

7. 现在剩下的是我们在这旅途生活中的祈求；因此接着是：\BibleQuote{我们的日用粮，求祢今天赐给我们。}求祢赐给我们永恒的事物，也赐给我们暂时的事物。祢应许了天国，求祢不拒绝我们生活的所需。祢日后将在天上赐给我们永远的荣耀，求祢在此世赐给我们暂时的支持。因此是\BibleQuote{每日}和\BibleQuote{今天}，即在此现世的时间。因为当这生命过去之后，我们还会求日用粮吗？因为那时不再称为\BibleQuote{每日}，而是\BibleQuote{今天}。现在称为\BibleQuote{每日}，因为一天过去，另一天来临。当有一天成为永恒时，还会称为\BibleQuote{每日}吗？这日用粮的祈求无疑有两种理解：既指身体食物的必需供应，也指我们灵性支持的必需。为了维持日常生活，没有它我们就不能生存，这需要身体食物的供应。这就是食物和衣服，但以部分代表全体。当我们求饼时，我们以此理解为一切事物。也有一种灵性食物，是信友所知道的，你们也将知道，当你们在天主的祭台前领受它时。这也是\BibleQuote{日用粮}，仅为此生所需。因为我们到了基督面前，开始与祂永远为王时，我们还会领受圣体吗？所以圣体是我们的日用粮；但我们要如此领受，使我们的肉身不仅得以滋养，灵魂也得滋养。因为在那里所领悟的力量是合一，使我们聚集在祂体内，成为祂的肢体，我们就成为我们所领受的。那才真正是我们的日用粮。再者，我现在在你们面前所讲的也是\Quote{日用粮}；你们在教会中每日听到的读经是日用粮，你们听到并咏唱的圣歌也是日用粮。因为这一切在我们的旅途状态中是必需的。但当我们到了天上，我们还会听圣言吗？我们这些将亲眼看见圣言本身、亲耳听见圣言本身、并像天使现在那样吃喝祂的人？天使需要书籍、讲解者和诵读吗？当然不需要。他们因看见而阅读，因为他们看见真理本身，并从那源泉充分饱足，而我们只得到几滴水珠。所以关于我们的日用粮，已说明这个祈求在此生是必需的。\par

8.\BibleQuote{求你宽恕我们的罪债，如同我们宽恕亏负我们的人。}这求恕除了在此生还有必要吗？因为在来世我们将没有罪债。什么是罪债？不就是罪过吗？看，你快要领洗了，那时你所有的罪都会被抹去，什么都不会留下。无论你曾做过什么恶事，无论是行为、言语、欲望或思想，全都会被抹去。然而，如果在领洗后的生活中没有犯罪的保障，我们就不会学习这样的祈祷：\BibleQuote{求你宽恕我们的罪债。}只是我们务必要做到紧接着的一句：\BibleQuote{如同我们宽恕亏负我们的人。}那么，你们这些即将进入，领受完全彻底罪债赦免的人，尤其要当心，你们心中不要存有对任何人的怨恨，好使你们领洗后如释重负，仿佛一切罪债都免除干净，然后开始计划向仇敌报复，就是那些曾得罪你们的人。宽恕吧，如同你们被宽恕。天主不能得罪任何人，但祂却宽恕，祂不欠任何债。那么，那自身被宽恕的人该当如何宽恕，当祂宽恕所有丝毫不欠祂的人？\par

9. \Quote{不要让我们陷入诱惑，但救我们免于凶恶。}在来世中，这还会是必要的吗？\Quote{不要让我们陷入诱惑，}这句话只有在可能存在诱惑的地方才会被说出来。我们在\BibleBook{约伯}{传}中读到：\Quote{人的生命在世上不就是一种考验吗？}那么我们在祈求什么呢？请听。宗徒雅各伯说：\Quote{当人受诱惑时，不要说：\InnerQuote{我被天主诱惑}。}他指的是那些邪恶的诱惑，人因这些诱惑而被欺骗，并屈服于魔鬼的轭下。这就是他所谈论的那种诱惑。因为还有另一种诱惑，被称为考验；关于这种诱惑，经上写道：\Quote{上主你们的天主考验你们，要知道你们是否爱祂。}\Quote{要知道}是什么意思？\Quote{为使你们知道}，因为祂早已知道。对于那种使我们被欺骗和引入歧途的诱惑，天主不诱惑任何人。但无疑地，在祂深奥而隐藏的审判中，祂会遗弃一些人。当祂遗弃他们时，诱惑者就找到了机会。因为如果天主遗弃他，诱惑者便在他身上找不到任何抵抗其力量的东西，而是立刻将自己呈现为他的占有者。因此，为了祂不遗弃我们，我们才说：\Quote{不要让我们陷入诱惑。}\Quote{因为每个人受诱惑，}同一宗徒雅各伯说：\Quote{都是被自己的私欲所牵引、所诱惑。然后，私欲怀孕后，生出罪恶；罪恶完成后，生出死亡。}那么他借此教导了我们什么呢？要同我们的私欲作战。因为你即将在圣洗圣事中除去你的罪恶；但私欲仍然会存留，在你重生之后，你必须与它们争战。因为与自己作战的战斗仍然存在。不必害怕外部的敌人：征服你自己，整个世界就被征服了。外部的任何诱惑者——无论是魔鬼还是魔鬼的仆役——能对你做什么呢？谁把获利的希望摆在你面前来诱惑你，只要在你身上找不到贪财之心，那么想通过获利来诱惑你的人能有什么效果呢？而如果在你的身上找到了贪财之心，你一见利益就起火，被这腐败食物的饵所捕获。但如果他在你身上找不到贪财之心，网罗就徒然张设。或者诱惑者把一位绝世美女摆在你面前；如果内心有贞洁，外在的邪恶就被战胜了。因此，为了使你不会被陌生女子的美貌所诱，要在内心与自己的私欲作战；你对你的敌人没有直接的知觉，但对你的贪欲却有。你看不见魔鬼，但你看见吸引你的对象。那么，战胜你内心所感觉到的东西吧。勇敢地战斗，因为那重生了你的那一位就是你的审判者；祂布置了战场，祂正在准备冠冕。但是因为你无疑会被击败，如果你没有祂的帮助，如果祂遗弃了你：因此你在祈祷中说：\Quote{不要让我们陷入诱惑。}审判者的愤怒把一些人交给了他们的私欲；宗徒说：\Quote{天主把他们交给了他们心中的私欲。}祂是如何交付他们的？不是通过强迫，而是通过遗弃他们。\par

10. \Quote{救我们免于凶恶，}可能属于同一句话。因此，为使你明白这全是一句话，它如此表述：\Quote{不要让我们陷入诱惑，但救我们免于凶恶。}因此他加上了\Quote{但}，以表明这一切都属于同一句话：\Quote{不要让我们陷入诱惑，但救我们免于凶恶。}这是怎么回事？我分别提出它们。\Quote{不要让我们陷入诱惑，但救我们免于凶恶。}藉着救我们免于凶恶，祂不让我们陷入诱惑；藉着不让我们陷入诱惑，祂救我们免于凶恶。\par

11. 这的确是一个极大的诱惑，亲爱的弟兄们，在今生是一个极大的诱惑，当我们身上那使我们得以获得宽恕的东西本身成为诱惑的对象时，如果我们曾在任何诱惑中跌倒过。这是一个可怕的诱惑，当那可以治愈我们从其他诱惑所受创伤的东西被夺去时。我知道你们还没有明白我的意思。请留意听，好明白。假设贪婪诱惑了一个人，他在某次诱惑中被战胜了（因为有时甚至优秀的摔跤手和战士也可能遭受粗暴对待）：贪婪已经胜过了这个人，尽管他是个好战士，他做了一些贪婪的事。或者有过一时的情欲；它没有使人犯奸淫，也没有达到通奸的地步，因为当这事发生时，人无论如何必须阻止犯罪行为。但\Quote{他看见女人就贪恋她；}他让自己的思想带着比正当更大的愉悦停留在她身上；他接受了攻击；尽管他是优秀的斗士，他受了伤，但他没有同意；他击退了自己情欲的冲动，用忧伤的苦楚惩罚了它，他击退了它；并且得胜了。然而，就他曾经滑倒这一点，他是否有理由说：\Quote{宽免我们的罪债，如同我们也宽免得罪我们的人。}同样，对于所有其他诱惑，难处在于在它们之中没有不应说\Quote{宽免我们的罪债}的机会。那么，我所提到的那可怕的诱惑是什么呢？那严重的、惊人的、必须用我们全部力量、全部决心避免的诱惑是什么？就是当我们想要报复自己时。愤怒被点燃，人渴望报复。可怕的诱惑！你正在失去那可以用来获得其他过犯宽恕的东西。如果你在其他感官和其他情欲上犯了任何罪，你本可以由此得到治愈，因为你可以说：\Quote{宽免我们的罪债，如同我们也宽免得罪我们的人。}但是那煽动你报复的人，将使你失去你说\Quote{如同我们也宽免得罪我们的人}的能力。当那能力失去时，所有的罪都将被保留；完全没有被赦免的。\par

12. 我们的主、导师和救主，知道此生中这危险的诱惑，当他教导我们这祷文中的六或七个祈求时，他没有取其中任何一个亲自讲解并更热切地推荐给我们，超过这一个。我们不是说了：\BibleQuote{我们的天父，你是在天上的，}以及其后的一切吗？为什么在祷文的结语之后，他没有详细向我们讲解，无论是最初提出的，还是结束时或中间的？因为他为何不说：如果天主的名不在你们中被尊为圣，或者如果你们在天主的国度中无份，或者如果天主的旨意不奉行在你们身上如同在天上，或者如果天主不守护你们，使你们不陷于诱惑；为何这一切都不说？但他说什么？\BibleQuote{我实在告诉你们：如果你们宽恕别人的过犯，}这是指那祈求：\BibleQuote{宽免我们的罪债，如同我们也宽免罪债的人。}他略过了他教导我们的所有其他祈求，唯独以特别的力量教导了这一条。对于那些罪，如果人犯了，他知道有方法可以治愈，不必如此强调；有必要强调的，是那种罪，如果你犯了，就没有方法可以治愈其余。因此你应该不断地说：\BibleQuote{宽免我们的罪债。}什么罪债？并不缺乏；因为我们不过是人；我说话稍多于当说的，说了不该说的，笑得过分，吃得过多，愉快地听了我本不该听的，喝得过多，愉快地看了不该看的，愉快地想了不该想的；\BibleQuote{宽免我们的罪债，如同我们也宽免罪债的人。}这一点如果你失去了，你自身就失丧了。\par

13. 注意，我的弟兄们，我的儿子们，天主的儿子们，我恳求你们，注意我向你们所说的。尽你们最大的力量与你们自己的心争战。如果你们看见你们的愤怒在反对你们，就向天主祈祷反对它，使天主使你们战胜你们自己，使天主使你们（我说）战胜的，不是你们外在的仇敌，而是你们内在的灵魂。因为他会赐予你们当下的助佑，并且会这样做。他宁愿我们向他求这个，胜过求雨。因为你们看见，亲爱的，主基督教导了我们多少祈求；几乎找不到其中一条是提到日用食粮的，这是为了使我们的全部思想都按来世的生活塑造吗？因为我们能害怕他不赐给我们什么吗？他曾应许并说：\BibleQuote{你们先寻求天主的国和他的义德，这一切自会加给你们；因为你们的父知道你们需要这一切，在你们求他以前。你们先寻求天主的国和他的义德，这一切自会加给你们。}因为许多人甚至曾受饥饿的考验，却被发现是金子，并没有被天主遗弃。如果那每日内在的食粮离开他们的心，他们早已因饥饿而丧亡了。此后，让我们主要地饥饿。因为：\BibleQuote{饥渴慕义的人是有福的，因为他们要得饱饫。}但是他能仁慈地垂顾我们的软弱，看顾我们，正如经上所说：\BibleQuote{你当记得我们不过是灰尘。}他从尘土造了人并赋予生命，为了他那泥土所作的缘故，赐下他的独生子至死。谁能解释，谁能有相称的概念，他爱我们有多深呢？\par

\SourceRecord{Translated by R.G. MacMullen.；Nicene and Post-Nicene Fathers, First Series；Vol. 6.；Edited by Philip Schaff.；Buffalo, NY: Christian Literature Publishing Co.,；1888.；<http://www.newadvent.org/fathers/160307.htm>.}

% structure: p0001 source=h1 target=section reason=page-title

\section{新约讲道集第八篇}

\Emph{[LVIII. Ben.]}\par

\Emph{再论主祷文}，\Emph{\BibleRef{玛 6:5}} \Emph{。致候洗者}\par

1. 你们刚才重复了信经，其中以简短的概要包含了信德。我先前已经告诉过你们，保禄宗徒说什么：\BibleQuote{他们尚未相信的，怎能呼号祂呢？}因为你们已经听见、学习并重复了如何相信天主；今天要听如何呼号祂。圣子自己，正如你们在宣读福音时所听到的，将这祷文教导了祂的门徒和祂的信徒。我们对于获得我们的案件有良好的希望，因为这样一位辩护人已经口授了我们的诉状。你们所宣认的圣父的同席者，坐在圣父的右边，祂是我们的辩护人，也将是我们的审判者。因为祂将从那里来审判生者死者。那么，也学习这祷文，八天后你们将需要重复它。但你们中凡是没有很好地重复信经的人，还有足够的时间，让他们学习它；因为在安息日，在全体在场者面前，你们将需要重复它：在最后的安息日，当你们在这里领受圣洗时。但从今天起八天后，你们将需要重复这祷文，就是你们今天所听见的。\par

2. 其中的第一句是：\BibleQuote{我们的天父，祢在天上。}我们已在天上找到了父；让我们谨慎生活在地上。因为找到这样一位父的人，理当如此生活，以便配得承受祂的产业。但我们大家一起说：\BibleQuote{我们的天父，祢在天上。}何等大的屈尊！皇帝这样说，乞丐也这样说；奴隶这样说，他的主人也这样说。他们一起说：\BibleQuote{我们的天父，祢在天上。}因此他们明白自己是弟兄，因为同有一位父。现在，主人不要不屑于将他的奴隶当作弟兄，因为主基督已经屈尊将他当作弟兄。\par

3. \BibleQuote{愿祢的名被尊为圣，愿祢的国来临。}天主之名的这种尊圣，正是我们借此成圣的。因为祂的名常是圣的。我们也希望祂的国来临；它必会来临，即使我们不希望；但希望并祈求祂的国来临，无非是向祂祈求，使祂使我们配得祂的国，免得它来临，却不到我们这里来（愿天主禁止）。因为对许多人来说，它永不会来临，尽管它必然来临。因为它要来到那些将听到\BibleQuote{你们这些蒙我父降福的，来吧！承受自创世以来给你们预备了的国度吧}的人身上。但它不会来到那些将听到\BibleQuote{可咒骂的，离开我，到那给魔鬼和他的使者预备的永火里去！}的人身上。因此，当我们说\BibleQuote{愿祢的国来临}时，我们祈求它来临到我们身上。什么是\Quote{来临到我们身上}？就是找到我们是良善的。所以我们祈求的是，祂使我们成为良善的；因为那时祂的国就会来临到我们身上。\par

4. 我们继续：\BibleQuote{愿祢的旨意奉行在人间，如同在天上。}天使在天上侍奉祢，愿我们在地上侍奉祢！天使在天上不冒犯祢，愿我们在地上不冒犯祢！如同他们承行祢的旨意，愿我们也这样承行！而在这里，我们祈求什么，岂不是要我们成为良善的吗？因为当我们承行天主的旨意（因为祂无疑承行祂自己的旨意）时，祂的旨意就在我们身上完成。我们还可以从另一种正确的意义上理解这些话：\BibleQuote{愿祢的旨意奉行在人间，如同在天上。}我们领受了天主的诫命，它令我们喜悦，令我们的心灵喜悦。\BibleQuote{因为我按内心来说，是天主的诫命所喜悦的。}于是祂的旨意成就在天上。因为我们的心灵被比作天，而我们的肉体被比作地。那么\BibleQuote{愿祢的旨意奉行在人间，如同在天上}是什么意思？就是如同祢的命令喜悦我们的心灵，也愿我们的肉体同意它；从而那由宗徒所描述的争斗就结束了，\BibleQuote{因为肉体的欲望相反圣神，圣神的欲望相反肉体。}当圣神相反肉体时，祂的旨意如今已在天上做成；当肉体不相反圣神时，祂的旨意如今已在地上做成。当祂愿意时，将会有完全的和谐；那么现在就让争斗存在，好使以后有胜利。同样，\BibleQuote{愿祢的旨意奉行在人间，如同在天上}，也可以被理解为将\Quote{天}当作教会，因为它是天主的宝座；而\Quote{地}当作不信者，对这些人说：\BibleQuote{你既是土，还要归于土。}因此，当我们为我们的仇敌、教会的仇敌、基督徒圣名的仇敌祈祷时，我们祈求祂的旨意得以奉行\BibleQuote{在人间，如同在天上}，也就是，如同在祢的信徒中，也照样在祢的亵渎者中，好使他们全都成为\Quote{天}。\par

5. 接下来，\BibleQuote{求你今天赏给我们日用的食粮}。可以简单地理解为，我们为每日的养活献上这祈祷，好使我们富足；或者如果不是这样，也使我们不致匮乏。祂说了\Quote{每日}，因为只要\Quote{今天}还在被称呼。我们每日生活，每日起床，每日得喂养，每日饥饿。但愿祂赐给我们日用的食粮。为什么祂没有说\Quote{遮盖}呢？因为支持我们生命的在于饮食，我们的遮盖在于衣服和住处。人不该渴求比这些更多的东西。因为宗徒说：\BibleQuote{我们没有带什么到世界上来，也不能带什么出去；只要有吃有喝有遮盖，就当知足。}贪心消失，本性就富足了。因此，如果这祈祷是指我们的日用养料，既然这是对\BibleQuote{求你今天赏给我们日用的食粮}的良好理解，我们就不要奇怪，在\InnerQuote{食粮}的名下也包括了其他必要的东西。就如当\Person{若瑟}邀请他的兄弟们时，他说：\BibleQuote{这些人今天要与我一同吃饭}。怎么，他们只吃饭吗？不，但唯独提到食粮时，其他的一切都被领会了。同样，当我们为日用食粮祈祷时，我们求的是我们在地上为身体所需的一切。可是\Person{主耶稣}说什么？\BibleQuote{你们先寻求天主的国和它的义德，这一切自会加给你们。}此外，对\BibleQuote{求你今天赏给我们日用的食粮}的另一个很好的理解是，你们的圣体圣事，我们日用的食粮。因为信友们知道他们领受的是什么，对他们而言，领受那为现时必要的日用食粮是好的。他们于是为自己祈祷，好使他们成为好人，好使他们在善、信德和圣善的生活中持守。他们希望这样，他们祈祷这样；因为如果他们不在此善生中持守，他们将被从这食粮中分离。因此，\BibleQuote{求你今天赏给我们日用的食粮}。这是什么？让我们这样生活，好使我们不被从你的祭台分离。同样，向我们展开、并且日复一日如同被擘开的\OriginalTerm{天主的圣言}{Word of God}，就是\Quote{日用食粮}。正如我们的身体渴望那另一食粮，我们的灵魂也渴求这食粮。所以我们既单单祈求这食粮，凡今生为灵魂和身体所需的一切，都包含在\Quote{日用食粮}之内。\par

6. \BibleQuote{求你宽免我们的罪债，}我们说，我们很可以这样说；因为我们说的是真理。因为谁居住在肉躯内而没有罪债呢？有什么人生活得如此，以至于这祈祷对他不是必要的呢？他可能自高自大，却不能为自己辩护。效法税吏，而不是像法利塞人那样自大，是好的，\Quote{他上到圣殿里}，夸耀自己的功德，掩盖了自己的创伤。而那位说\BibleQuote{主啊，可怜我这个罪人吧}的，知道他为什么上去。这祈祷，我的弟兄们，想想，这祈祷是主耶稣教导他的门徒们献上的，那些伟大的首批宗徒，我们羊群的领袖。如果羊群的领袖为自己的罪过祈求赦免，那么羊羔该当如何呢？关于它们经上说：\BibleQuote{把公羊羔献给上主}。你们知道，你们在信经中已经重复了这一点，因为你们在其中提到了\Quote{罪过的赦免}。有一种罪过的赦免是一次而永远赐予的；另一种是每日赐予的。有一种罪过的赦免是在圣洗圣事中一次而永远赐予的；另一种是在我们在此世生活期间，在天主经中赐予的。因此我们说：\BibleQuote{求你宽免我们的罪债}。\par

7. 而且天主把我们带入了一个盟约、一个协定和一个与祂坚固的纽带中，因为我们说：\BibleQuote{如同我们宽免我们的罪债人一样}。凡想有效地说\BibleQuote{求你宽免我们的罪债}的人，必须真诚地说\BibleQuote{如同我们宽免我们的罪债人一样}。如果这最后一句他或不说，或假意地说，那么第一句他便白说了。我们尤其是对你们这些即将领受圣洗圣事的人说：从心里宽恕一切。你们这些信友，利用这机会聆听这祈祷和我们的讲解，你们要完全地、从心里宽恕你们对任何人所有的一切。在那里，在天主看见的地方宽恕。因为有时人用口宽恕，却在心里保留；他为了人的缘故用口宽恕，却在心里保留，因为不怕天主的眼目。但你们要完全宽恕。你们直到这些圣日所保留的一切，在这些圣日，至少宽恕吧。\BibleQuote{太阳不可在你们的愤怒中落下}，但许多太阳已经过去了。那么，现在既然我们正在庆祝伟大太阳的庆日，就是那太阳，经上关于祂说：\BibleQuote{为你们敬畏我名的人，正义的太阳将要升起，以祂的翅膀带来痊愈}，但愿你们的愤怒终于也过去。什么是\Quote{在祂的翅膀中}？是在祂的保护中。因此圣咏说：\BibleQuote{求你把我藏在你翅膀的荫下}。至于其他在审判之日将要后悔，但已太迟，并将哀号却无济于事的人，智慧早已预言他们那时将说什么，他们将在悲愤中后悔和呻吟：\BibleQuote{骄傲给我们带来了什么利益？富贵与夸耀又给我们带来了什么好处？这一切都像阴影一样消逝了。}并且\BibleQuote{因此我们迷失了真理之路，正义的光没有照耀我们，正义的太阳也没有照临我们}。那太阳只照在义人身上；但我们现在所看见的这太阳，天主\Quote{每日}使其\BibleQuote{上升，光照善人，也照恶人}。义人得见那太阳，那太阳现在藉着信德住在我们心里。如果你们发怒，不要让这太阳在你们的心里因你们的愤怒而落下；\BibleQuote{不要让太阳在你们的愤怒中落下}，免得你们发怒，以致正义的太阳在你们身上落下，而你们留在黑暗中。\par

8. 现在不要以为愤怒算不得什么。\BibleQuote{我因愤怒而眼目昏花}\BibleRef{咏 6:7}，先知说。确实，眼目昏花的人不能看见太阳；若他试图去看，只会感到痛苦，毫无乐趣。愤怒是什么？是报复的私欲。人渴望报复，而基督尚未被报复，神圣的殉道者尚未被报复。天主的忍耐仍在等待，使基督的仇敌、殉道者的仇敌能够悔改。我们是谁，竟寻求报复？若天主向我们追讨报复，我们往哪里去？那从未在任何事上伤害过我们的天主，不愿报复我们；而我们这些几乎天天得罪天主的人，却寻求报复？所以，宽恕吧，从心里宽恕。你们若愤怒，却不要犯罪。\BibleQuote{你们愤怒，却不要犯罪}\BibleRef{咏 4:5}。你们作为人，若被愤怒胜过，就当愤怒；却不要犯罪，不要在心中怀怒（因为若怀怒，便是对自己怀怒），免得你们不能进入那光明。所以，要宽恕。那么愤怒是什么？是报复的私欲。仇恨是什么？是根深蒂固的愤怒。若愤怒变得根深蒂固，便称为仇恨。看来先知也承认这点，当他曾说：\BibleQuote{我因愤怒而眼目昏花}之后，又加上说：\BibleQuote{我在所有仇敌中变得根深蒂固}\BibleRef{咏 6:8}。愤怒初起时是\Quote{木屑}，日久天长便成了\Quote{大梁}。我们有时责备一个愤怒的人，自己心中却怀恨于心；于是基督对我们说：\BibleQuote{你看见你兄弟眼中的木屑，却看不见自己眼中的大梁}\BibleRef{玛 7:3-5}。木屑如何长成大梁？因为没有立刻拔除。因为你任凭太阳多次升起又落下，照在你的怒气上，使之根深蒂固；因为你怀着恶意的猜疑，浇灌那木屑，借浇灌滋养它，借滋养使它成为大梁。那么，至少当经上说\BibleQuote{凡恨他兄弟的，就是杀人犯}\BibleRef{若壹 3:15}时，你应当战栗。你没有拔剑，没有造成身体创伤，没有用打击杀死别人；只是心中存有恨意，就此你被视为杀人犯，在天主眼中你有罪。那人还活着，你却已经杀了他。就你而言，你已经杀死了你恨的人。所以，悔改，修正自己。若你的房屋里有蝎子或毒蛇，你岂不竭力清除它们，好能安全居住？然而你们心中有愤怒，甚至有根深蒂固的愤怒，生出那么多仇恨、那么多大梁、那么多蝎子、那么多毒蛇，难道你们还不愿净化天主的殿宇——你们的心吗？那么，要实行经上说的：\BibleQuote{如同我们宽恕我们的罪人一样}；并安心说：\BibleQuote{宽恕我们的罪过}。因为在这世上，你们不能没有罪过；但那些大罪，你们在圣洗中已得赦免，且应当永远摆脱的，是一类；而那些每日的小罪，人在这世上不能没有，因此这每日的祈祷及其盟约是必要的；好使我们能满怀喜乐地说：\BibleQuote{宽恕我们的罪过}，同时能真诚地说：\BibleQuote{如同我们宽恕我们的罪人一样}。关于过去的罪，我们说了这么多；现在关于将来呢？\par

9. \BibleQuote{不要让我们陷入诱惑}\BibleRef{玛 6:13}：宽恕我们已经做的，并恩赐我们不再犯罪。因为凡被诱惑所胜的，就是犯罪。因此，宗徒雅各伯说：\BibleQuote{当人受诱惑时，不要说：\InnerQuote{我是由天主受诱惑}，因为天主不会受邪恶诱惑，也不诱惑任何人。每个人受诱惑，是被自己的私欲所勾引、所饵诱。然后，私欲怀孕，产生罪；罪完成后，产生死亡。}\BibleRef{雅 1:13-15}。为此，你们不要被自己的私欲勾引；不要同意它。它无法怀孕，除非你同意。你一旦同意，就像在心中接纳了它的拥抱。私欲已升起，你要拒绝它，不要跟随它。它是一种违法、不洁、可耻的私欲，会使你远离天主。因此，不要给它同意的拥抱，免得你要哀叹那出生；因为如果你同意，即当你拥抱它时，它就怀孕，\Quote{私欲怀孕，产生罪}。你还不害怕吗？\Quote{罪产生死亡}；至少，害怕死亡吧。如果你不害怕罪，至少害怕它所导致的下场。罪是甘甜的，但死亡是苦涩的。这就是人的不幸：他们因之犯罪的东西，死时都留在地上，而罪本身却随身携带。你为了金钱犯罪，它必须留在地上；或为了乡间别墅；它必须留在地上；或为了某个女人；她必须留在地上；无论你为何犯罪，当你闭上眼睛死去时，你都必须把它留在地上；然而你所犯的罪本身，你却随身携带。\par

10. 愿罪获得宽恕；过去的被宽恕，未来的停止。但在此世，你们不能没有罪；无论是大罪、小罪，或是微罪。然而，连这些小罪和微罪也不可轻视。江河是由点滴之水汇聚而成的。连较小的罪也不可轻视。水从船上的狭缝渗入，船舱不断积水，若不加理会，船就会沉没。但水手并非闲懒；他们的手在活动——活动着，好让水一天天被排出。同样，你们的手也要活动，好让你们天天把水舀出去。\Quote{你们的手要活动}是什么意思？让他们施舍，行善工，这样你们的手就忙碌起来。\Quote{把你的饼分给饥饿的人，将贫穷无家者领到你家中；若见赤身露体者，就给他衣服穿。}尽你所能，用你能支配的方法，怀着喜乐去做，然后满怀信心地献上你的祈祷。这祈祷将有两只翅膀，一份双重的施舍。什么是\Quote{双重的施舍}？\Quote{宽恕，你们也必蒙宽恕；施舍，你们也必得施舍。}一种施舍是发自内心的，当你宽恕你的兄弟的罪过时。另一种施舍是从你的财物中行出的，当你把饼分给穷人时。要献上两者，否则你的祈祷就会因缺少翅膀而无法起飞。\par

11. 因此，当我们说了\Quote{不要让我们陷入诱惑}之后，接着就是\Quote{但救我们免于凶恶}。凡希望从凶恶中被拯救出来的人，就证明他正处于凶恶之中。宗徒也这样说：\Quote{赎回光阴，因为日子是邪恶的。}但谁\Quote{愿得享生命，愿看见好日子}？既然所有人在此肉身中只有邪恶的日子，谁不渴望它呢？要这样做：\Quote{使你的舌头远离邪恶，使你的嘴唇不说欺诈的话；远离邪恶，行善，寻求和平，并追随它；}这样你就摆脱了邪恶的日子，你的祈祷\Quote{救我们免于凶恶}就得以成就。\par

12. 因此，前三个祈求：\Quote{愿你的名被尊为圣，愿你的国来临，愿你的旨意奉行在人间，如同在天上，}是为永恒而求。但接下来的四个与此生有关：\Quote{求你今天赏给我们日用的食粮。}当我们到达那圆满福乐时，我们还会日复一日地求日用食粮吗？\Quote{求你宽恕我们的罪过。}在那没有罪过的国度里，我们还会这样说吗？\Quote{不要让我们陷入诱惑。}那时没有诱惑，我们还能这样说吗？\Quote{但救我们免于凶恶。}那时没有需要被拯救的凶恶，我们还会这样说吗？因此，这四个祈求因我们的日常生活而必要，而前三个则是关乎永生。但我们要祈求一切，以达致那生命，并在此祈祷，好使我们不致与此生命分离。你们受洗后，每天必须诵念此祈祷。因为天主经每天在教会天主的祭台前诵念，信众都听得到。所以我们不用担心你们学不好，因为即使有人不能完全记住，他也会因天天聆听而学会。\par

13. 因此，在星期六，当靠天主的恩宠你们将举行守夜时，你们要背诵的不是祈祷文，而是信经。因为如果你们现在还不懂信经，你们就不会每天在教会中和在人群中听到它。但当你们学会后，为了不忘记，你们要每天早晨起床时念；准备睡觉时，重温你们的信经，向主重温它，提醒自己，不要厌倦重复。因为重复是有益的，免得遗忘悄然侵入。不要说：\Quote{我昨天念了，今天念了，每天念，我完全知道了。}要使你的信德存于心中，反省自己，让你的信经如同一面镜子摆在你面前。从中看见你自己，是否相信你所宣誓的一切，并因而日日因你的信德而喜乐。让它成为你的财富，让它成为你灵魂每日的衣物。你不是每天起床时都穿衣服吗？同样，借着每天重复信经来穿戴你的灵魂，免得遗忘使它赤裸，使你赤身露体，并发生宗徒所说的（愿它远离你们！）：\Quote{如果我们脱下衣服，不会被发现是赤身露体的。}因为我们将被我们的信德所穿戴：这信德同时是衣服和护心镜；衣服抵挡羞耻，护心镜抵挡逆境。但当我们到达那我们将统治的地方时，就不再需要念信经了。我们将看见天主；天主本身将是我们的异象；看见天主将是我们现今信德的赏报。\par

\SourceRecord{Translated by R.G. MacMullen.；Nicene and Post-Nicene Fathers, First Series；Vol. 6.；Edited by Philip Schaff.；Buffalo, NY: Christian Literature Publishing Co.,；1888.；<http://www.newadvent.org/fathers/160308.htm>.}

% structure: p0001 source=h1 target=section reason=page-title

\section{\WorkTitle{论新约第九篇讲道}}

\Emph{[第五十九篇 圣奥斯定]}\par

\Emph{再论主祷文}，\BibleRef{玛 6:5}。\Emph{致候洗者\OriginalTerm{候洗者}{Competentes}}。\par

1. 你们已经重温了你们所信的，现在要听你们当求的。因为你们不能呼求你们尚未相信的那一位；正如宗徒所说：\BibleQuote{从未相信祂的，怎能呼求祂呢？}因此，你们首先学习了《信经》，其中藏着你们信德简短而崇高的准则：言辞简短，内涵崇高。而你们今天领受并要在八天后背诵的这篇祈祷，是主亲自（正如你们在宣读福音时所听到的）教导祂的门徒，并由他们传给我们，因为\BibleQuote{他们的声音传遍了普世}。\par

2. 你们既已找到天上的父，切莫贪恋地上的事物。因为你们即将说：\Quote{我们的天父，祢在天上}。你们已开始属于一个伟大的家族。在这位父之下，主人与奴隶是弟兄；在这位父之下，将军与普通士兵是弟兄；在这位父之下，富人与穷人是弟兄。所有基督信徒在世上各有不同的父亲，有的显贵，有的卑微；但他们同呼一位天上的父。若我们的父在那里，那里就为我们预备了产业。但祂是这样一位父：我们可以与祂一同拥有祂所赐的产业。祂赐予产业，却不是因死亡而留给我们。因为祂自己并不离去，而是永远存留，好让我们来到祂面前。既然我们知道了该向谁祈求，也让我们知道该求什么，免得我们因妄求而冒犯这样一位父。\par

3. 那么，主耶稣基督教导我们向天上的父求什么呢？\Quote{愿祢的名被尊为圣}。我们向天主求的这是什么样的恩惠呢？天主的圣名永远是圣的；为何我们祈求它被尊为圣？除非是我们藉以成为圣的。因此，我们祈求那永远圣的，在我们内被尊为圣。当你们受洗时，天主的圣名在你们内被尊为圣。那么，你们领洗之后为何还要献上这祈祷？岂不是为使你们那时所领受的，永远存留在你们内？\par

4. 接着是另一祈求：\Quote{愿祢的国来临}。天主的国必然来临，无论我们求与不求。那么为何这样祈求？岂不是为使那将临于诸圣的国，也临于我们；使天主也把我们算在祂诸圣的行列中，好让祂的国降临到我们身上？\par

5. 我们在第三个祈求中说：\Quote{愿祢的旨意奉行在人间，如同在天上}。这是什么意思？是愿天使们在天上怎样事奉祢，我们也在世上怎样事奉祢。因为祂的圣天使服从祂；他们不冒犯祂；他们因爱祂而遵行祂的命令。我们就是这样祈求的：愿我们也能以爱遵行天主的命令。此外，这句话还有另一层理解：\Quote{愿祢的旨意奉行在人间，如同在天上}。我们内的天上是指灵魂，我们内的地上是指肉身。那么\Quote{愿祢的旨意奉行在人间，如同在天上}是什么意思？是愿我们听从祢的训诫时，我们的肉身也与我们同心；免得肉身与灵性相争，使我们无法完成天主的命令。\par

6. 接着祈祷说：\Quote{求祢今天赏给我们日用的食粮}。我们在此祈求父赐予肉身所需之粮，以\Quote{食粮}指代我们所需的一切；或者，我们理解为那日用的食粮，就是你们即将从祭台领受的圣体——我们祈求祂赐予我们，这很好。我们祈求的是什么呢？无非是我们不犯任何邪恶，以免与那圣洁之粮分离。再者，每日宣讲的天主圣言也是日用之粮。它虽非肉身之粮，也并非就不是灵魂之粮。但当此生终结后，我们既不再寻求那为解饥饿的食粮，也无需再领受祭台上的圣事，因为那时我们已与基督同在，现今我们领受的正是祂的身体；此刻我们向你们所说的这些话，也无需再说，圣卷也无须再读，因为我们将亲眼看见那本身就是天主圣言的那一位，万物由祂而造，天使由祂而饱饫，天使由祂而光照，天使由祂而智慧；他们不需要迂回的话语，而是畅饮那唯一的圣言，满怀祂而欢呼咏叹，永不止息。因为圣咏说：\BibleQuote{住在祢殿中的人，是有福的；他们不断地赞美祢。}\par

7. 因此，在今世生活中，我们祈求接下来的话：\Quote{宽免我们的罪债，如同我们也宽免得罪我们的人}。在圣洗中，一切债，即一切罪，都被完全宽免。但因为在世无人能无罪生活，即使没有犯下须与祭台分离的重大罪行，也绝无人能在世上全然无罪；我们只能一次领受唯一洗礼；因此，我们在这篇祈祷中得知如何日日被洗涤，使我们的罪日日被宽免；但条件是我们要做到紧随其后的话：\Quote{如同我们也宽免得罪我们的人}。因此，我的诸位弟兄，我劝告你们——你们是在天主恩宠中的我的儿女，但同在天父之下是我的弟兄——我劝告你们：每当有人得罪你们、对你们犯罪，前来认错并请求宽恕时，你们要宽恕他，并立刻从心里原谅他；免得你们对自己关闭了来自天主的宽恕。因为你们若不宽恕，祂也不宽恕你们。因此，我们是在今世献上这祈求，因为只有在今世，罪过才能被赦免，因为它们能被犯下；而在来世，罪过不被赦免，因为不再被犯下。\par

8. 在此之后我们祈祷说：\BibleQuote{不要让我们陷于诱惑，但救我们免于凶恶。} 这也意味着，我们不被引入诱惑，在今生有必要祈求，因为今生有诱惑；并且祈求\BibleQuote{我们得以脱离凶恶，}因为这里有凶恶。因此，在这七个祈求中，三个关乎永生，四个关乎现世生命。\BibleQuote{愿你的名受显扬。} 这将是永远的。\BibleQuote{愿你的国来临。} 这国将是永远的。\BibleQuote{愿你的旨意奉行在人间，如同在天上。} 这将是永远的。\BibleQuote{求你今天赏给我们日用的食粮。} 这将不是永远的。\BibleQuote{求你宽恕我们的罪过。} 这将不是永远的。\BibleQuote{不要让我们陷于诱惑。} 这将不是永远的。\BibleQuote{但救我们免于凶恶。} 这将不是永远的；但哪里有诱惑，哪里有凶恶，就有必要献上这个祈祷。\par

\SourceRecord{Translated by R.G. MacMullen.；Nicene and Post-Nicene Fathers, First Series；Vol. 6.；Edited by Philip Schaff.；Buffalo, NY: Christian Literature Publishing Co.,；1888.；<http://www.newadvent.org/fathers/160309.htm>.}

% structure: p0001 source=h1 target=section reason=page-title

\section{论新约讲道集 第十篇}

\Emph{[LX. Ben.]}\par

\Emph{论福音之言，\BibleRef{玛 6:19}，\BibleQuote{不要在地上为自己积蓄财宝}，等。劝勉人施舍。}\par

凡身处患难、自身资源不足的人，都会寻找一位明智之人向他请教，从而知道该做什么。那么，让我们假设整个世界如同一个人。他寻求逃避邪恶，却迟迟不行善；随着患难增多，自身资源枯竭，他还能找到比基督更明智的咨询者吗？无论如何，至少让他找到一个更好的，然后按自己的意愿行事。但如果他找不到更好的，就让他来到那位无处不在的主面前：让他咨询，从他那里接受劝告，遵守善诫，逃避大恶。因为人们如此惧怕、如此抱怨的暂时现世祸患，并因抱怨而冒犯那纠正他们的主，以致找不到他的拯救帮助；这些现世的祸患，我再说，无疑是短暂的；它们要么经过我们，要么我们经过它们；要么在我们活着时过去，要么在我们死亡时留下。在患难中，时间短促的算不得重大。无论你是谁，思考明天的事，却不回想昨天的事。后天来临，这个明天也将成为昨天。但如果人们因忧虑逃避那些短暂、甚至飞逝的暂时患难而如此不安，那么他们应当怎样思考才能逃避那些永无止境的患难呢？\par

人生是艰难的境况。诞生不就是进入劳苦的生活吗？婴儿的啼哭正是我们即将面临的劳苦的见证。没有人能免饮这忧苦之杯。亚当所饮的杯，必须饮尽。诚然，我们是由真理之手所造，但因罪恶被放逐到虚妄的日子。\BibleQuote{我们是按天主的肖像被造的}，但我们因罪恶的过犯而损毁了它。因此，圣咏提醒我们如何被造，以及我们沦落到了何种境地。因为它说：\BibleQuote{人纵然行走在天主的肖像中}。看，他受造为何物。他沦落到了何处？请听随后的话：\BibleQuote{却要徒然烦恼}。他行走在真理的肖像中，却在虚妄的计谋中烦恼。最后，看他的烦恼，看吧，如同照镜子一般，对自己不满。\BibleQuote{纵然}，他说，\BibleQuote{人行走在天主的肖像中}，因此是某种伟大的存在，\BibleQuote{却要徒然烦恼}；好像我们要问：请问，人如何徒然烦恼？\BibleQuote{他积蓄财宝}，他说，\BibleQuote{不知为谁积聚}。看，这个人，即全人类代表为一个人，自身毫无资源，失去劝告，迷失了健全心智的道路；\BibleQuote{他积蓄财宝，不知为谁积聚}。还有比这更疯狂、更不幸的吗？但他真的为己而积吗？并非如此。为何不是为己？因为他必须死，因为人生短暂，因为财宝存留，而积累者迅速逝去。因此，他怜悯那行走在天主肖像中的人，他承认真实的事，却追随虚妄的事，他说：\BibleQuote{他将徒然烦恼}。我为他悲伤；\BibleQuote{他积蓄财宝，不知为谁积聚}。他是为自己积聚吗？不。因为人死了，财宝却存留。那么为谁呢？如果你有任何好建议，请告诉我。但你没有任何建议给我，因此你也没有给自己。既然我们两人都没有，就让我们一同寻求，一同领受，一同思考此事。他烦恼，他积蓄财宝，他思考，劳苦，因焦虑而失眠。你终日受劳作折磨，整夜被恐惧搅动。为使你的钱柜装满金钱，你的灵魂却在焦虑中发热。\par

3. 我看见了，我为你们忧伤；你们烦乱不安，而那位不能欺骗的保证我们说，\Quote{你们是徒然烦乱不安。}因为你们在积蓄财宝：假设你们一切事业都顺利，且不说损失，不说在追求每一种利益时所遭遇的如此巨大的危险和死亡（我不说身体的死亡，而是邪恶思想的死亡，为了黄金能进来，正直就出去；为了你们在外表上得以穿衣，内在却赤身露体），但且搁下这些和其他类似的事，且对一切不利于你们的事保持缄默，让我们只考虑顺利的情况。看，你们在积蓄财宝，利益从四面八方流入你们手中，你们的金钱像泉水般涌流；凡有缺乏的地方，丰盛就涌来。你们难道没有听过，\Quote{如果财富增加，不要将心寄托其上？}看，你们在获取，你们烦乱不安，并非没有结果，但仍是徒然。\Quote{怎么，}你们要问，\Quote{我是怎样徒然烦乱不安？我装满我的钱柜，我的墙壁几乎装不下我所获得的，那我怎么是徒然烦乱不安呢？}\Quote{你们在积蓄财宝，却不知道为谁收集。}或者如果你们知道，我请你们告诉我。我会听你们。为谁呢？如果你们不是徒然烦乱不安，告诉我你们为谁积蓄财宝？\Quote{为我自己，}你们说。你们敢这样说吗？你们这么快就要死的人？\Quote{为我的子女。}你们敢这样说他们吗？他们也要这么快就死。据说，为人父亲为儿子积蓄是很大的亲情责任；更确切地说，是很大的虚妄，一个快死的人为同样快死的人积蓄。如果为你们自己，为什么聚集，既然你们死时都要留下一切？你们的子女也是这样；他们会继承你们，但不会长久停留。我不说他们会是什么样的子女，也许放荡会挥霍掉贪婪所累积的。所以另一个人因放荡浪费了你们用许多辛劳聚集的。但我搁下这个。也许他们会是好子女，不会放荡，会保存你们留下的，会增加你们保存的，不会散尽你们集聚的。那么你们的子女会和你们一样虚妄，如果他们这样做，如果他们在这方面效法你们做父亲的。我要对他们说刚才对你们说的话。我要对你们的儿子，对你们为之节省的那个人说，\Quote{你们在积蓄财宝，却不知道为谁收集。}因为你们不知道，他也不知道。如果虚妄在他身上继续存在，真理对他会失去效力吗？\par

4. 我不再强调，甚至可能在你们有生之年，你们不过是在为盗贼积蓄。在一个夜晚，他们可能来，发现这么多日日夜夜的积聚都已齐备。可能你们是在为强盗或拦路贼积蓄。我对此不再多说，免得我回想并重新揭开过去苦难的伤口。空虚无谓的虚荣所积聚的许多东西，敌人的残忍已发现随手可得。这不是我所希望的事；但大家都关心要害怕它。愿天主免去！愿祂自己的鞭笞足够。愿我们所祈祷的那位宽恕我们！但如果祂问我们为谁积蓄，我们怎么回答？那么，人啊，无论你是谁，你徒然积蓄财宝，当我与你处理这件事，与你共同寻求一个共同原因中的意见时，你如何回答我？因为你发言并回答：\Quote{我为我自己、为我的子女、为我的后代积蓄。}我已经说过，对于这些子女本身，有许多畏惧的理由。但我搁下这样的考虑：你的子女可能这样生活，以致成为你的咒诅，像你的敌人所愿的那样；假设他们像父亲自己所愿的那样生活。然而，有多少人陷入这些不幸，我已经宣布并提醒过你们了。你们对此战栗，虽然你们没有改正自己。因为你们除了回答\Quote{也许不会如此}之外，还有什么可回答？好吧，我也这样说了；也许我说你们不过是在为盗贼、或强盗、或拦路贼积蓄。我没有肯定地说，而是也许。哪里有一个\Quote{也许}，哪里就有一个\Quote{也许不}；所以你们不知道将要发生什么，因此你们\Quote{是徒然烦乱不安}。你们现在看见真理说得多么真，虚荣如何徒然地烦乱不安。你们听见了，终于学会了智慧，因为当你们说\Quote{也许是为我的子女}，却不敢说\Quote{我肯定是为我的子女}，你们实际上不知道为谁积蓄财富。所以，正如我所看见并已经说过的，你们自己毫无办法；你们找不到什么来回答我，我也不能来回答你们。\par

5. 为此，我们两人都应寻求并请求忠告。我们有机会咨询的，不是任何智者，而是智慧本身。让我们两人都听从耶稣基督，\BibleQuote{对于犹太人来说是绊脚石，对于外邦人是愚妄，但对于蒙召的，无论是犹太人还是希腊人，基督却是天主的德能和天主的智慧。}你为什么为你的财富准备强有力的辩护？请听天主的德能，没有比祂更强有力的。你为什么为保护你的财富准备明智的忠告？请听天主的智慧，没有比祂更智慧的。或许当我说出我要说的话时，你会感到冒犯，这样你就成了犹太人，\BibleQuote{因为基督对犹太人是冒犯。}或者，也许当我说完后，你会觉得那是愚妄，这样你就成了外邦人，\BibleQuote{因为基督对外邦人是愚妄。}然而你是基督徒，你已被召叫。\BibleQuote{但对于蒙召的，无论是犹太人还是希腊人，基督是天主的德能和天主的智慧。}那么，当我说完我要说的话时，不要悲伤；不要被冒犯；不要以轻蔑的神情嘲笑我的愚妄（如你所认为的）。让我们倾听。因为我要说的话，基督已经说过了。如果你轻视传报者，仍要畏惧审判者。那么我要说什么呢？福音的诵读刚刚解除了我的困境。我不再读任何新的内容，只是唤起你们对刚才所读之事的回忆。你们曾因自身能力不足而寻求忠告；那么请看这正确忠告的泉源吧，从它的溪流中无需恐惧毒药，你能从中汲取多少就汲取多少。\par

6. \BibleQuote{不要在地上为自己积蓄财宝，那里有虫蛀和锈蚀，也有贼挖洞偷窃；但要在天堂为自己积蓄财宝，那里没有贼靠近，也没有虫蛀；因为你的财宝在哪里，你的心也必在哪里。}你还要等什么？事情是清楚的。忠告是公开的，但邪恶的贪欲隐藏着；不，不止如此，更糟的是，它也是公开的。因为掠夺没有停止它的蹂躏；贪婪没有停止欺诈；恶意没有停止发假誓。这一切都是为了什么？为了积聚财宝。积聚在\Emph{哪里}？在地上，确实，\Emph{来自}尘土\Emph{归于}尘土。因为对那个犯罪并为我们（如我所说）预定了劳苦之杯的人，说过：\BibleQuote{你是尘土，你将归于尘土。}财宝在地上是合理的，因为心也在那里。那么，\Quote{我们举心向上，归于主？}又在哪里呢？为你自己悲伤吧，你这理解我的人；如果你真心悲伤，就改过自新。你们鼓掌却不去实行要到几时呢？你们所听的是真的，再真不过了。那么就让这真事得以实行吧。我们赞美唯一的天主，但我们不改变，以免在这赞美中徒然烦乱。\par

7. 因此，\BibleQuote{不要在地上为自己积蓄财宝；}无论你已从经验中发现地上所积蓄的如何丧失，或你尚未有这种经验，仍要惧怕以免如此。经验要改造那言语所不能改造的人。现在人不能起来，不能出去，却全都同声呼喊：\Quote{我们有祸了，世界正在倒塌。}如果它在倒塌，你为什么不迁移？如果一位建筑师告诉你，你的房子即将倒塌，你难道不会在你徒然的哀叹之前就迁移吗？世界的建造者告诉你们世界即将倒塌，你们还不相信吗？请听预言者的声音，听那警告你的忠告。预言的声音是：\BibleQuote{天地要过去。}警告的声音是：\BibleQuote{不要在地上为自己积蓄财宝。}那么，如果你相信天主的预言；如果你不轻视祂的警告，就让祂所说的话得以实行。那给你如此忠告的天主不会欺骗你。你所施舍的不会失去，反而要跟随那已先行送走的。因此我的忠告是：\BibleQuote{施舍给穷人，你就在天堂有了财宝。}你不会没有财宝；但你在焦虑中所拥有的地上财宝，你将无忧无虑地拥有天堂的财宝。那么，转移你的财物吧。我是在给你保守的忠告，而非丧失的忠告。祂说：\BibleQuote{你将有财宝在天堂，并且来，跟随我，}好让我带你去到你的财宝那里。这不是浪费，而是储蓄。为什么人保持沉默？让他们听见，并终于因经验而知道该惧怕什么，让他们去做那不会带来惧怕的事，让他们把财物转移到天堂。你把小麦放在低洼处；你的朋友来了，他知道谷粒和土地的性质，指示你的无知，对你说：\Quote{你做了什么？}你把谷粒放在平坦的土壤、低洼的地里；土壤潮湿，它会全烂掉，你会白费力气。你回答：那么我该怎么办？他说，把它移到高处。那么，你听从一位朋友关于你谷物的忠告，却轻视天主关于你心的忠告吗？你惧怕把谷粒放在低洼的地里，却愿意把你的心失在地里吗？看，你的主天主在给你有关你心的忠告时说：\BibleQuote{因为你的财宝在哪里，你的心也必在哪里。}祂说，把你的心举向天堂，使它不致在地里腐烂。这是那想要保全你的心而非毁灭它的天主的忠告。\par

8. 若是如此，那些后来没有这么做的人，该怎样悔改呢？现在他们如何自责啊！我们本可以在天堂我们现在在地上失去的。仇敌拆毁了我们的房屋；但他能打破天堂吗？他杀死了被派来守护的仆人；但他能杀死那本会保护他们的主吗？那地方\BibleQuote{没有贼接近，也没有虫蛀}。现在有多少人在说：\Quote{我们本可以在那里拥有，并安全地藏起我们的宝藏，在那里稍后我们就可以安全地跟随它们。为什么我们没有听从我们的主？为什么我们轻视了圣父的劝诫，以致经历了仇敌的侵犯？}如果这是一个好劝告，我们就不应迟缓地留意它；如果我们拥有的必须被转移，就让我们把它转移到那个地方，从那里我们不能失去它。我们所施舍的穷人，不正是我们的搬运者吗？我们通过他们将我们的财物从地上运到天上。那么，施舍吧：你只是给予你的搬运者，他把你所给的带到天上。你怎么说，他如何把它带到天上？因为我看见他通过吃来结束它。毫无疑问，他搬运它，不是通过保存它，而是通过把它当作食物。什么？你们忘记了：\BibleQuote{我父所祝福的，来，承受天国；因为我饿了，你们给了我吃的}，并且，\BibleQuote{你们既然对我这些最小兄弟中的一个做了，就是对我做了}。如果你们没有轻视站在你们面前的乞丐，就请考虑你们给他的来到了谁那里。\BibleQuote{你们既然}祂说，\BibleQuote{对我这些最小兄弟中的一个做了，就是对我做了}。那给了你们施舍之资的，已经接受了它。那在末日将祂自己赐予你们的，已经接受了它。\par

9. 为此我曾多次提醒你们，亲爱的，我向你们承认，在天主的圣经中，有一件事使我非常震惊，我应当反复提醒你们注意。我请你们想想我们的主耶稣基督自己所说的：在世界末日，当祂来审判时，祂将把万民聚集在祂面前，把人分成两部分；祂要把一些人放在右边，另一些人在左边；祂要对右边的人说：\BibleQuote{我父所祝福的，来，承受自创世以来为你们预备的天国吧}。但对左边的人说：\BibleQuote{可咒骂的，离开我，到那给魔鬼和他的使者预备的永火里去吧}。探究这如此大的赏报或如此大的惩罚的原因。\BibleQuote{承受天国}和\BibleQuote{到永火里去}。为什么前者要承受天国？\BibleQuote{因为我饿了，你们给了我吃的}。为什么后者要进入永火？\BibleQuote{因为我饿了，你们没有给我吃的}。我问这意味着什么？我看到关于那些要承受天国的人，他们作为善良忠信的基督徒施舍了，不轻视主的话，并且以确实的信赖盼望祂的应许，他们如此行了；因为如果他们不这样做，这种贫瘠不会与他们的善行相符。因为他们可能是贞洁的，不是骗子，不是酒徒，并远离恶事。但如果他们没有加上善行，他们就会仍然贫瘠。因为他们会遵守\BibleQuote{远离邪恶}，但他们不会遵守\BibleQuote{并力行善事}。尽管如此，甚至对他们，祂也不说：\BibleQuote{来，承受天国}，因为你生活贞洁；你没有欺骗任何人，你没有欺压任何穷人，你没有侵犯他人的地界，你没有因发誓而欺骗任何人。祂没有说这个，而是说：\BibleQuote{承受天国，因为我饿了，你们给了我吃的}。这是多么超越一切！当主不提其余，只提这个时。又对其他人说：离开我到永火里去吧，那是为魔鬼和他的使者预备的。祂能对不虔敬的人提出多少指控，如果他们问：\Quote{为什么我们要进入永火？}为什么？你们问，你们这些奸淫者、杀人者、骗子、亵渎神明者、不信者。然而祂没有提这些中的任何一个，而是说：\BibleQuote{因为我饿了，你们没有给我吃的}。\par

10. 我看得出你们和我一样感到惊奇。这确实是一件奇妙的事。但我尽我所能揣摩这件如此奇异之事的原因，并不向你们隐瞒。经上记载：\BibleQuote{水能熄灭火焰，施舍也能熄灭罪。}又记载：\BibleQuote{你要把施舍存在穷人的心中，它必会为你向主求情。}又记载：\BibleQuote{大王，请听我的劝告，用施舍来赎你的罪罢。}还有许多其他神圣言语的见证，表明施舍对于熄灭和消除罪过大有功效。因此，对那些将要被定罪——不，毋宁说将要被加冕——的人，祂只举出施舍一事，仿佛祂要说：我若仔细检视和衡量你们，以极大的精确性审察你们的行为，我就很难找不到定你们罪的把柄；但是，\BibleQuote{你们进入天国罢，因为我饿了，你们给了我吃的。}所以你们要进入天国，不是因为你们没有犯罪，而是因为你们用施舍赎了你们的罪。又对另一些人说：\BibleQuote{你们进入那为魔鬼及其使者预备的永火里去罢。}他们同样有罪，在罪中已久，迟迟才害怕自己的罪，当他们心中思量自己的罪时，他们岂敢说自己是冤枉被定罪，如此公义的审判者对他们宣判这样的判决是冤枉的呢？他们省察自己的良心和灵魂的所有创伤，岂敢说：我们是被冤枉定罪的。关于他们，曾在\BibleBook{}{智慧篇}中预言说：\BibleQuote{他们自己的罪要当面指证他们。}毫无疑问，他们会看到自己因罪和邪行而被定罪是公义的；然而，祂好像对他们说：\BibleQuote{不是因为你们所想的那样，而是\InnerQuote{因为我饿了，你们没有给我吃的。}}因为如果你们远离所有这些行为，转向我，并用施舍赎了所有这些罪过和罪，那么那些施舍现在就会拯救你们，并使你们从如此重大过犯的罪责中解脱出来；因为\BibleQuote{怜悯人的人是有福的，因为他们要受怜悯。}但现在你们进入永火里去罢。\BibleQuote{那不怜悯人的，也要受无怜悯的审判。}\par

11. 但愿我已说服你们，我的弟兄们，施舍你们地上的食粮，并敲打天上的食粮！主就是那食粮。祂说：\BibleQuote{我是生命的食粮。}但如果你不给那有需要的人，祂怎能赐给你呢？一个人在你面前有需要，你在另一位面前有需要；既然你在另一位面前有需要，另一人在你面前有需要，那另一人在他本身有需要的人面前有需要。因为你在祂面前有需要的那一位，一无所缺。所以，你们愿意人怎样待你们，你们也要怎样待人。因为这与那些朋友的情况不同——他们习惯于用各自的善举彼此责备，例如：\Quote{我为你做了这事}，对方回答：\Quote{我也为你做了这事}——祂希望我们为祂做些好事，因为祂先为我们做了好事。祂一无所缺，因此祂就是主。我曾对主说：\BibleQuote{祢是我的天主，因为祢不需要我的财物。}尽管如此，尽管祂是主，是真正的主，不需要我们的财物，但为了让我们也能为祂做些什么，祂竟屈尊在祂的穷人身上饥饿。\Quote{我饿了，}祂说，\BibleQuote{你们给了我吃的。主啊，我们几时见了祢饿了？凡你们对我这些最小的兄弟中的一个所做的，就是对我做的。}简而言之，让人们听见并恰当地思考：当基督饥饿时给祂食物是何等大的功德，而当基督饥饿时轻视祂是何等大的罪恶。\par

12. 对罪的悔改确实能使人向善转变；但若缺乏慈悲的工作，似乎连悔改也毫无益处。这真理通过\Person{若翰}的口证实了，他对那些来到他那里的人说：\BibleQuote{毒蛇的种类！谁指示你们逃避那将要来临的忿怒？那么，结出与悔改相称的果子来罢！不要心里说：我们有\Person{亚巴郎}为父。我告诉你们：天主能从这些石头中给\Person{亚巴郎}兴起子孙来。现在斧子已经放在树根上；凡不结好果子的树，就砍下来，丢在火里。}关于这果子，他上面说：\BibleQuote{结出与悔改相称的果子来。}所以，凡不结这些果子的人，不要以为他能借着不结果子的悔改获得罪的赦免。这些果子是什么，他后来亲自说明了。因为在他这些话之后，群众问他说：\BibleQuote{那么我们该做什么呢？}也就是说，你们用如此可怕的力量劝我们要结出的果子是什么？\BibleQuote{有两件衣裳的，就分给那没有的；有食物的，也当这样行。}我的弟兄们，还有什么比这更清楚、更确定、更明确的呢？那么，他上面说的\BibleQuote{凡不结好果子的树，就砍下来，丢在火里}还有什么别的意思呢？不就是左边的人将听到的：\BibleQuote{你们进入永火里去罢，因为我饿了，你们没有给我吃的}吗？所以，离开罪恶只是一件小事，如果你忽略了医治过去的罪，正如经上所记：\BibleQuote{孩子，你犯了罪，不要再犯了。}但为免他仅凭这一点就自以为稳妥，经文又说：\BibleQuote{并为你从前的罪祈祷，求它们得以赦免。}但如果你祈祷赦免，却不结出与悔改相称的果子，使自己配得垂听，以致像一棵不结果的树被砍下来丢在火里，那么祈祷对你有什么益处呢？所以，如果你们希望祈祷赦罪时蒙垂听，就要\BibleQuote{赦免，你们也必蒙赦免；施舍，你们也必得施舍。}\par

\SourceRecord{Translated by R.G. MacMullen.；Nicene and Post-Nicene Fathers, First Series；Vol. 6.；Edited by Philip Schaff.；Buffalo, NY: Christian Literature Publishing Co.,；1888.；<http://www.newadvent.org/fathers/160310.htm>.}

% structure: p0001 source=h1 target=section reason=page-title

\section{\WorkTitle{论新约·第11讲}}

\EditorialNote{LXI. Ben.}\par

\Emph{论福音经文}，\BibleRef{玛 7:7}\Emph{，\BibleQuote{你们求，必要给你们；}等等。劝勉行哀矜。}\par

1. 在圣福音的教训中，主劝勉我们祈祷。祂说：\BibleQuote{你们求，必要给你们；你们找，必要找着；你们敲，必要给你们开门。因为凡求的，就必得到；找的，就必找着；敲的，就必给他开门。或者，你们中间有哪个人，儿子向他求饼，反而给他石头？或者求鱼，反而给他蛇？或者求鸡蛋，反而给他蝎子？你们纵然不善，尚且知道把好东西给你们的儿女，何况你们在天之父，岂不更要把好东西赐给求祂的人？}祂说：\BibleQuote{你们纵然不善，尚且知道把好东西给你们的儿女。}弟兄们，真是奇事！我们是不善的，却有一位好父亲。还有什么比这更明显的？我们听到了我们适当的名字：\BibleQuote{你们纵然不善，尚且知道把好东西给你们的儿女。}现在请看，祂对祂称之为不善的人显示了什么样的父亲。\BibleQuote{何况你们的父？}谁的父？无疑是不善者的父。那是怎样的父？\BibleQuote{除天主一个外，没有谁是善的。}\par

2. 为此，我们这些不善的人有一位好父亲，为的是使我们不再常是不善的。没有一个不善的人能使别人成为善的。如果不善的人不能使别人成为善的，那么不善的人怎能使自己成为善的？只有那永远为善的那一位，才能使一个不善的人成为善的。\BibleQuote{求祢医治我，我必痊愈；求祢拯救我，我必得救。}那么，为什么那些虚妄的人用他们自己虚妄的话对我说：\Quote{你愿意就能拯救自己}？\BibleQuote{上主，求祢医治我，我必痊愈。}我们是由善者造为善的；因为\BibleQuote{天主造人原是正直的}，但藉着我们的自由意志，我们变成了不善的。我们从善变为不善是有能力的，我们从不善变为善也将有能力。但那总是善的那一位，使善从不善中出来；因为人凭自己的意志无力医治自己。你受伤时，不会去找医生，反而自己弄伤自己；但你弄伤自己后，就会找人来医治你。于是，现世暂时的善，那些涉及身体和肉身的善，我们纵使不善，也懂得如何给我们的儿女。因为连这些也是善的，谁能否认？鱼、鸡蛋、饼、果实、麦子、我们所看见的光、我们所呼吸的空气，这一切都是善的。甚至那些使人自高自大、不愿承认别人与自己平等的财富；那些使人更爱炫目的衣饰而不思想共同本性的财富；我再说，甚至这些财富也是善的。但我刚才所提的所有这些善，善人与不善人都能拥有；虽然它们本身是善的，却不能使它们的拥有者成为善的。\par

3. 有一种善能使人成为善的，又有另一种善你能藉此而行善。使人成为善的善是天主。因为除那永远善的天主外，没有人能使人为善。因此，为叫你能成为善的，当呼求天主。但另有一种善，你能藉此而行善，那就是你所拥有的一切。有金子，有银子；它们是善的，不是那种能使你成为善的善，而是你能藉此而行善的善。你有金和银，你却想要更多的金银。你既拥有，又渴望拥有；你既满足，又口渴。这是病，不是富足。人患水肿时，充满了水，却总是口渴。他们充满了水，却仍渴求水。那么，你既有了这种渴求的病，又怎能以富足为乐？你有金子，它是善的；然而你并没有能藉此成为善的善，却有能藉此而行善的善。你问：我能用金子行什么善？难道你没有在圣咏中听见：\BibleQuote{他广施钱财，周济穷人；他的仁义永远常存。}这就是善，这就是能使你成为善的善：仁义。如果你有能使你成为善的善，就用那不能使你成为善的善去行善吧。你有钱财，慷慨地施舍。藉着慷慨施舍，你增加了仁义。\BibleQuote{因为他广施钱财，周济穷人；他的仁义永远常存。}请看，什么减少了，什么增加了。你的钱财减少了，你的仁义增加了。那减少的是你不久必会失去的，那减少的是你不久必会留下的；那增加的是你将永远拥有的。\par

4. 那么，这是一个赢利的秘诀，我正在传授；要学着这样交易。你称赞那卖掉铅而得到金的商人，难道你不称赞那施舍钱财而得到仁义的商人吗？但你会说：我不施舍我的钱财，因为我还没有仁义。那有仁义的人，让他施舍钱财吧；我没有仁义，那么至少让我保有我的钱财。那么，你因为还没有仁义，就不想施舍你的钱财吗？不如还是施舍吧，好叫你能获得仁义。因为除了从天主这仁义之源那里，你从何处能得到仁义呢？所以，如果你愿意有仁义，就成为天主的乞丐吧。祂刚才藉着福音劝你求、找、敲。祂认识祂的乞丐，看哪，那家主，那大富者，即富有属灵和永久的财富者，劝勉你说：\BibleQuote{求、找、敲；求的，就得到；找的，就找着；敲的，就给他开门。}祂劝你去求，难道祂会拒绝你所求的吗？\par

5. 试想一个取自相反情形（如同那不义的审判官）的比喻或比较，这是对我们祈祷的鼓励。\BibleQuote{有一个，}主说，\BibleQuote{在城里有一个法官，不敬畏天主，也不尊重人。}一个寡妇天天烦扰他，并说：\BibleQuote{给我伸冤。}他长久不肯；但她不停地恳求，终于因她的强求，他做了他不愿凭自己善心去做的事。因为祂正是通过一个相反的情形，教导我们要祈祷。\par

6. 再者，祂说：\Quote{有一个人，有客人来到他那里，就到他的朋友那里，开始敲门说：\InnerQuote{有客人来到我这里，借给我三个饼。}}那人回答说：\Quote{我已经上了床，我的仆人同我在一起。}那人不肯离开，而是站着继续请求，敲门并恳求，如同一个朋友对另一个朋友一样。祂又说什么？\Quote{我告诉你们，他起来，不是因他友谊的缘故，}而是\Quote{因那人的强求，就给他所需之量。不是因为他们的友谊，}虽然他是他的朋友，但是\Quote{因着他的强求。}\Quote{因着他的强求}是什么意思呢？因为他不停地敲门；即使他的请求被拒绝，他也没有转身离开。那不愿给予的人，因请求者不灰心而给予所求。那么，那位劝我们祈求、不祈求就感到不悦的善者，岂不更要赐予吗？但有时祂赐得较慢，是为了显示祂恩赐的价值，并非拒绝。长久渴望之物更能喜乐地获得，而速得之物则被轻视。所以，祈求，寻找，坚持。借着祈求与寻找，你们自身长大，以能容纳更多。天主为你们保存着祂不愿速速赐下的东西，为使你们学会以极大的渴望来渴望大事。因此\BibleQuote{所以我们应当常常祈祷，不可灰心。}\par

7. 如果天主藉劝诫、鼓励和命令使我们成为祂的乞丐，要我们祈求、寻找和敲门，那么我们也要留意那些向我们祈求的人。我们祈求，我们向谁祈求？我们祈求的是什么人？我们求什么？向谁求？我们是什么人？我们求什么？我们向善的天主祈求，而祈求者是恶人；但我们求的是义德，借此我们得以成为善人。我们所求的，是那可以永远拥有的事物，当我们充满它时，将不再缺乏。但为使我们被充满，让我们饥渴；饥渴着，让我们祈求、寻找、敲门。\BibleQuote{因为饥渴慕义的人是有福的。}他们为何有福？他们饥渴，他们有福吗？缺乏岂是福？他们之有福，不在于饥渴，而在于将被充满。福乐在于满全，而不在于饥饿。但饥饿必须在满全之前，以免对食粮生厌。\par

8. 我们已经说了我们向谁求、我们是什么人、以及我们求什么。但我们自己也会被人求。因为我们是天主的乞丐；为使祂承认祂的乞丐，我们也要承认我们自己的乞丐。但在此情形中再想想，当有人向我们祈求时，他们是谁，向谁求，求什么？那么，祈求的是什么人？是人。他们向谁求？向人。他们是什么人？会朽坏的人。向谁？向会朽坏的人。他们是什么人？软弱的人。向谁？向软弱的人。他们是什么人？可怜人。向谁？向可怜人。除了财富方面，祈求者与被祈求者是一样的。你在你的主人面前怎能厚颜祈求，却不承认你的同类？\Quote{我不是，}他会说，\Quote{像他那样，}我绝不像他那样。一个穿着丝绸、骄傲自大的人就是这样谈论一个衣衫褴褛的人。但我问你，当你们俩都赤身裸体时。我不问你现在穿着衣服的样子，而问你刚出生时的样子。两人都是赤身、软弱，开始一个痛苦的生命，因此以啼哭开始。\par

9. 那么，请看，富人啊，回想你最初的开端；看你是否带了什么到世界上来。你确实来了，并发现了如此丰盛。但请教我，你带来了什么？告诉我，你若羞于说，就听宗徒的话。\BibleQuote{我们没有带什么到世界上来。}他说：\BibleQuote{我们没有带什么到世界上来。}但也许因为你没有带进来，却在这里找到了许多，你会带些什么离去吗？这话也许因着爱财之心，你不敢承认。也听听这个，让不会阿谀的宗徒告诉你。\BibleQuote{我们没有带什么到世界上来，}即我们出生时；\BibleQuote{也不能带什么出去，}即我们离开世界时。你没有带进来，也不能带出去。那么你为什么在穷人面前自高自大呢？当婴儿初生时，只让父母、仆人、从属和谄媚的成群随从退开；然后让富家子女以他们的哭声被认出。让富有的妇女和穷妇一同生产；让她们不留意自己的孩子，离开一会儿；然后回来，若可能，认出他们。那么，富人啊，\BibleQuote{你没有带什么到世界上来，也不能带什么出去。}我对他们出生时所说的话，也可在他们的死亡时说。如果不是这样，当古老的坟墓偶然被打开时，让富人的骨骸可以被认出吧。因此，富人啊，要听从宗徒：\BibleQuote{我们没有带什么到世界上来。}承认吧，这是真的。\BibleQuote{也不能带什么出去。}承认吧，这也是真的。\par

10. 那么接下来是什么呢？\BibleQuote{只要有吃有穿，就当知足；因为那些想发财的人，陷于诱惑，陷入许多有害的欲望，这欲望把人淹没在败坏和灭亡中。因为贪饕是万恶的根源，有些人随从它，便偏离了信德。} 现在要思考他们所舍弃的。你们为他们舍弃了这些而悲伤，但请看他们陷入了什么。请听；\BibleQuote{他们偏离了信德，使自己陷入许多痛苦。} 但谁呢？\BibleQuote{那些想发财的人。} 成为富人与想成为富人是两回事。富人是指出身富裕家庭的人，他富有不是因为他想要，而是因为许多人给他留下了遗产。我看到他的财富，我不问他对财富的喜爱。在这段圣经中，受谴责的是贪饕，而不是金子、银子或财富，而是贪饕。因为那些不想成为富人，或不在意财富，不燃烧着贪欲，也不受贪饕之火灼烧，却仍然是富人的，让他们听宗徒的话（今天已诵读过）：\BibleQuote{你要嘱咐那些在今世富有的人。} 嘱咐他们什么？首先嘱咐他们不要心高气傲，因为财富最容易滋生骄傲。每一种果实、每一粒谷、每一棵树都有其特有的害虫，苹果的害虫是一种，梨的另一种，豆的另一种，麦的另一种。财富的害虫就是骄傲。\par

11. \BibleQuote{所以你要嘱咐那些今世富有的人，不要心高气傲。} 他已摒除滥用，现在要教导正当的使用。\BibleQuote{不要心高气傲。} 但抵御骄傲的屏障从何而来？来自接下来的话：\BibleQuote{也不要寄望于无常的财富。} 那些不寄望于无常财富的人，就不会心高气傲。如果他们不心高气傲，就让他们敬畏。如果他们敬畏，他们就不会心高气傲。有多少人昨天还是富人，今天却成了穷人？有多少人睡觉时富有，因盗贼来临抢走一切，醒来时却成了穷人？因此，\BibleQuote{嘱咐他们不要寄望于无常的财富，而要寄望于永生的天主，他厚赐百物给我们享受。} 这是暂时的和永恒的事物。但永恒的事物更值得享受，暂时的事物则为使用。暂时的事物如同旅人所需，永恒的事物如同家乡所需。暂时的事物，藉以行善；永恒的事物，使我们得以成善。因此，让富人这样做：\BibleQuote{不要心高气傲，也不要寄望于无常的财富，而要寄望于永生的天主，他厚赐百物给我们享受。} 让他们这样做。但他们能用所拥有的做什么呢？请听：\BibleQuote{让他们在善工上富足，乐于施舍。} 因为他们有资本。为什么他们不这样做呢？贫穷是艰难的状况。但他们可以轻易给予，因为他们有办法。\BibleQuote{让他们分享，} 也就是说，承认他们的同胞与自己平等。\BibleQuote{让他们分享，为自己积存美好的根基，预备将来。} 因为，他说，当我说 \BibleQuote{让他们乐于施舍，让他们分享，} 我不是要夺去或剥夺他们，或使他们一无所有。我所教导的是痛苦的教训；我给他们指明存放财货的地方，\BibleQuote{让他们为自己积蓄。} 因为我并不希望他们一直贫穷。\BibleQuote{让他们为自己积蓄。} 我不是要他们损失自己的财物，而是指示他们转移到何处。\BibleQuote{让他们为自己积蓄美好的根基，预备将来，使他们能抓住真正的生命。} 那么，现在的生命是虚假的生命；让他们抓住真正的生命。\BibleQuote{因为虚而又虚，万事皆虚。人在太阳下辛勤劳苦，究竟有什么益处？} 因此，必须抓住真正的生命，我们的财富必须转移到真正生命的地方，使我们能在这里给出的一切在那里找到。那交换我们财物的，也是改变我们自己的。\par

12. 那么，我的弟兄们，请施舍给穷人，\Quote{只要有食物和遮盖，我们就当知足。} 富人从自己的财富中并没有得到什么，除了穷人所乞求于他的食物和遮盖。你从你所拥有的一切中，还得到了什么呢？你得到了食物和必需的遮盖。我说的是必需的，不是无用的，也不是多余的。从你的财富中你还得到了什么？告诉我。的确，你所多余的一切都将成为多余的。那么，让你的多余成为穷人的必需吧。但你会说，我享用昂贵的宴席，我吃昂贵的肉食。而穷人吃什么？他们吃廉价的食物；穷人吃廉价的东西，而我，他说，吃昂贵的肉。好，我问你，当你们两个都吃饱时，昂贵的食物进入你体内，但一旦进入，它变成了什么？如果我们体内有镜子，难道我们不会为你所吃下的一切昂贵肉食而感到羞愧吗？穷人饥饿，富人也饥饿；穷人寻求饱足，富人也寻求饱足。穷人用廉价的食物吃饱，富人用昂贵的肉吃饱。两者都同样饱足，他们所追求的目标是同一个，只是一个人通过捷径达到，另一个人通过迂回的路达到。但你会说，我更喜欢我的昂贵食物。不错，但你很难满足，因为你如此挑剔。你不知道饥饿所调味的美味。我这样说，并不是要强迫富人吃穷人的食物和饮料。让富人使用他们的软弱所习惯的东西；但让他们感到遗憾，因为他们不能做别的事。因为如果他们能的话，对他们会更好。如果穷人因贫穷而不自高，你为什么因你的软弱而自高呢？那么，就使用精选昂贵的肉吧，因为你已经习惯了，因为你不能做别的事，因为如果你改变习惯，你就会生病。我允许你这一点，使用多余的东西，但给穷人必需的东西；使用昂贵的肉，但给穷人廉价的食物。穷人期待从你那里得到，而你也期待从天主那里得到；穷人仰望着和你一样被造的手，而你仰望着那创造你、并且不仅创造你、也创造穷人同你的那双手。祂为你们俩安排了同一个旅程，即今生；你们发现自己是旅伴，你们走在同一条路上：他什么也不带，而你负担过重；他什么也不携带，而你带着比你需要的更多。你负载了；把你所有的分给他一些；这样你既喂养了他，也减轻了你自己的负担。\par

13. 那么，请施舍给穷人；我恳求、我劝告、我命令、我吩咐你们。请随意施舍给穷人。因为我不会向你们隐藏，亲爱的诸位，为什么我认为有必要向你们讲这篇道理。当我往来于教会时，穷人们不断地请求我，让我向你们说话，以便他们能从你们那里得到一些东西。他们催促我向你们说话；当他们看到从你们那里什么也得不到时，他们就以为我在你们中间的所有劳动都是徒劳的。他们也期待从我这里得到一些东西。我尽我所能给他们；但我有足够的资源来满足他们所有的需要吗？既然我没有足够的资源来满足他们所有的需要，我至少是他们的使者，来到你们这里。你们听了并且鼓掌了；感谢天主。你们接受了种子，你们给出了回响。但你们的这些称赞反而压垮了我，使我处于危险之中。我忍受着，并且在忍受时颤抖。然而，我的弟兄们，你们的这些称赞不过是树上的叶子；我所寻求的是果实。\par

\SourceRecord{Translated by R.G. MacMullen.；Nicene and Post-Nicene Fathers, First Series；Vol. 6.；Edited by Philip Schaff.；Buffalo, NY: Christian Literature Publishing Co.,；1888.；<http://www.newadvent.org/fathers/160311.htm>.}

% structure: p0001 source=h1 target=section reason=page-title

\section{《新约讲道集 第十二篇》}

\Emph{[LXII. Ben.]}\par

\Emph{论福音的话}，\BibleRef{玛 8:8} \Emph{，\BibleQuote{主，我不堪当祢到我舍下来}等等，以及宗徒的话}，\BibleRef{格前 8:10} \Emph{，\BibleQuote{因为如果有人看见你这有知识的人，在偶像的庙里坐席}等等。}\par

1. 我们听到了，当福音宣读时，那在谦卑中彰显的信德受到了赞美。因为当主耶稣应许要往百夫长的家里去医治他的仆人时，他回答说：\BibleQuote{主，我不堪当祢到我舍下来，只要祢说一句话，我的仆人就会痊愈。}他自称不配，反倒表明自己配得基督进入他的心中，而非他的屋内。若他没有在心中怀有基督，他不敢怀着如此大的信德与谦卑说这样的话；因他害怕基督进入他的屋子，这并非什么大幸福。因为这位以言以行教导谦卑的夫子，甚至曾坐在一个名叫西满的骄傲法利塞人的家中；虽然他坐在他家里，但在那心中却没有\BibleQuote{人子可以枕头的地方}。\par

2. 因为如此，我们从主自己的话可以明白，祂曾阻止一个骄傲的人跟随祂，那人自愿要与祂同去：\BibleQuote{主，祢无论往哪里去，我都要跟随祢。}主看透他心中那不可见的事，便说：\BibleQuote{狐狸有洞，天空的飞鸟有窝，但人子却没有枕头的地方。}意思是：在你心里，有如狐狸的诡诈，以及如天空飞鸟的骄傲。但人子，单纯以对诡诈，卑微以对骄傲，没有枕头的地方；而这\InnerQuote{枕下}而非\InnerQuote{抬头}的姿势，正教导谦卑。因此，祂把那个自愿跟随的人拉回来，而另一个拒绝的人祂却吸引前去。因为同样在那一处，祂对某个人说：\BibleQuote{跟随我。}那人说：\BibleQuote{主，我要跟随祢，但请允许我先去埋葬我的父亲。}他的推辞实在是出于孝道；故此他更值得被除去推辞，并坚定他的蒙召。他想做的事是孝道的行为；但夫子教导他什么是更优先的。因祂愿意他成为活圣言的宣讲者，使别人得以存活。但另有其他人可以完成那首要的职责。祂说：\BibleQuote{让死人埋葬他们的死人。}当不信者埋葬一具死尸时，是死人埋葬死人。一人的身体失去了灵魂，另一人的灵魂失去了天主。因为灵魂是身体的生命，同样，天主的生命是灵魂的生命。身体失去灵魂时便消逝，同样，灵魂失去天主时便消逝。失去天主是灵魂的死亡；失去灵魂是身体的死亡。身体的死亡是必然的；灵魂的死亡是自愿的。\par

3. 主曾坐在一个骄傲的法利塞人家里。就如我所说的，祂在他家里，却不在他心里。然而，祂没有进入这位百夫长的家，却占据了他的心。至于匝凯，他既在家中也在心中接待了主。然而，百夫长的信德因谦卑受到赞美。因为他说：\BibleQuote{主，我不堪当祢到我舍下来。}主便说：\BibleQuote{我实在告诉你们，在以色列我从未见过这样大的信德。}这是按肉身说的。因为他按精神无疑也是以色列人。主曾来到肉身的以色列，即犹太人那里，首先在这民族中寻找迷失的羊，并且祂也自这民族取了祂的身体。祂说：\BibleQuote{在那里我没有见过这样大的信德。}我们只能以人能判断的方式衡量人的信德；但那位鉴察人心、无人能欺骗的主，祂为此人的心作证，听到了谦卑的话，并宣布了医治的判决。\par

4. 但他从何获得这种确信呢？他说：\BibleQuote{我也是个在权力下的人，手下有兵士；我对这个人说：去，他就去；对另一个说：来，他就来；对我的仆人说：你做这事，他就做。}我对我属下的一些人有权威，而我本人也在某个更高的权威之下。如果我这在权力下的人有命令的能力，那么祢，所有权力都事奉的主，该有何等的能力呢？这个人是外邦人，因为他是个百夫长。那时，犹太民族中有罗马帝国的兵士驻扎。他在那里过军旅生活，按百夫长的权限，既在权力之下，又对他人拥有权威；作为顺服的属下，又管辖其下的人。但是主（尤其请留意这一点，亲爱的各位，因为你们需要知道），虽然祂仅在犹太民族中，却早已预告教会将遍布全世界，为建立教会祂将派遣宗徒；祂自己虽不被看见，却被外邦人因信德而接受；被犹太人看见并处死。因为正如主并没有亲身进入此人的家，但虽然身体不在，却以威严临在，治愈了他的信德和他的家；同样，这位主也曾在身体上仅临在于犹太民族中；在其他民族中，祂既未由童贞女诞生，未受苦难，未行走，未忍受祂的人间苦难，也未施行祂的神迹奇事。这一切在其他民族中一件也未发生，然而论到祂的话却应验了：\BibleQuote{我不认识的人民，竟事奉了我。}如果他们不认识祂，又如何呢？\BibleQuote{他们因耳闻而服从了我。}犹太民族认识祂并钉死了祂；普世万民听闻而相信了。\par

5. 祂身体的这种缺席，可说是如此，而祂的能力却临在于万民之中，祂也在那触摸祂衣边的妇人身上表明了这一点，当时祂问说：\BibleQuote{谁摸了我？}祂询问，仿佛祂是缺席的；又如亲临般，祂治愈了。门徒们说：\BibleQuote{群众拥挤着祢，祢还说：谁摸了我？}因为好像祂行走时无人能触摸祂一样，祂说：\BibleQuote{谁摸了我？}他们回答：\BibleQuote{群众拥挤着祢。}主似乎要说：我问的是那触摸我的人，不是拥挤我的人。如今祂的身体也是如此，也就是祂的教会。少数人的信德\Quote{触摸}它，多数人的群众\Quote{拥挤}它。因为你们作为她的子女已经听过，基督的身体就是教会，如果你们愿意，你们自己就是。宗徒在许多地方这样说：\BibleQuote{为了祂的身体，就是教会；}又说：\BibleQuote{你们却是基督的身体，各自做肢体。}如果我们是祂的身体，那么祂的身体当年在人群中所受的，如今祂的教会也在经受。它被许多人拥挤，被少数人触摸。肉体拥挤它，信德触摸它。所以，请抬起你们的眼睛，我恳求你们，那些有眼可看的人。因为你们面前有可看之物。抬起信德的眼睛，只触摸祂衣边的边缘，这足以带来救恩的健康。\par

6. 你们看，你们从福音中所听到的当初将要发生的事，如今已经临现。因此，在称赞那位百夫长的信德时——他按肉身是外邦人，但心中却是家人——祂说：\BibleQuote{因此，我告诉你们：将有许多人从东从西而来。}不是所有的人，而是\BibleQuote{许多人}；然而他们要从\BibleQuote{东从西而来}；这两个部分代表整个世界。\BibleQuote{将有许多人从东从西而来，与亚巴郎、依撒格、雅各伯一同坐席于天国；但本国的子民要被驱逐到外面的黑暗中。}\BibleQuote{但本国的子民}，即犹太人。为何是\BibleQuote{本国的子民}？因为他们领受了法律，先知被派遣到他们那里，圣殿和司祭职在他们中间；他们庆祝一切未来之事的预像。然而，他们庆祝了预像，却不承认实体的临在。所以，祂说：\BibleQuote{因此，本国的子民}要进入外面的黑暗，那里有哀号和切齿。于是我们看到犹太人被弃绝，而基督徒从东从西被召叫到天国的筵席，与亚巴郎、依撒格和雅各伯同坐，那里的食粮是正义，杯爵是智慧。\par

7. 那么，弟兄们，请想一想，因为你们就是这些人；你们就是这民族，那时被预言，如今被彰显。是的，确实，你们就是那些从东从西被召叫，在天国坐席的人，不是在偶像的庙宇中。所以，你们要成为基督的身体，而不是拥挤祂身体的人。你们有祂的衣边可以触摸，好从血漏中得到痊愈，即从肉体的快乐中得到痊愈。我说，你们有衣边可以触摸。请将宗徒们视为衣裳，他们借着合一的织工紧密地依附在基督身边。在这些宗徒中，保禄就像衣边，是最小和最后的，正如他自己所说：\BibleQuote{我是宗徒中最小的。}在衣裳中，最后和最小的东西是衣边。衣边在外表上可轻视，但触摸它却有救恩的功效。\BibleQuote{直到此时此刻，我们又饥又渴，又赤身露体，又挨打。}还有什么境况比这更低微、更可轻视的呢？那么，如果你患有血漏，就去触摸吧。那衣裳的主人会发出能力，治愈你。刚才已向你们提出要触摸的衣边，当时从同一位宗徒的读经中读到：\BibleQuote{若有人看见你这有知识的人在偶像的庙里坐席，这软弱人的良心岂不壮胆去吃那祭偶像之物吗？你这知识就使软弱的弟兄丧亡，基督为他而死的弟兄丧亡！}你们想，人会被偶像欺骗到什么程度？他们以为偶像受到基督徒的尊敬。有人会说：\Quote{天主知道我的心。}是的，但你的弟兄不知道你的心。如果你软弱，要提防更大的软弱；如果你强壮，要关心你弟兄的软弱。他们看见你所做的，就壮胆去做更多，以致不仅想要吃，还想要在那里献祭。看哪，\BibleQuote{因你的知识，软弱的弟兄就丧亡。}那么，我的弟兄，请听；如果你轻视软弱者，你也要轻视兄弟吗？醒醒。如果你这样犯罪，难道不是直接得罪基督吗？因为你要注意你绝不能轻视的事。他说：\BibleQuote{但你们这样得罪了弟兄们，伤了他们软弱的良心，就是得罪基督。}那些轻视这些话的人，现在就去吧，在偶像的庙里坐席；他们岂不是那些拥挤而不触摸的人吗？当他们已在偶像的庙里坐席后，让他们来充满教会——不是来领受救恩，而是来造成拥挤。\par

8. 但你会说：\Quote{我害怕冒犯在我之上的人。}你当然应该害怕冒犯他们，这样你就不会冒犯天主。因为你害怕冒犯在你之上的人，看看是否没有一位高于你所害怕冒犯的那一位。那么你当然应该不愿冒犯在你之上的人。这是你的一条确定规则。但难道不明显，那高于一切者绝不可被冒犯吗？现在数一数在你之上的人。首先是你的父亲和母亲，如果他们正确地教育你，如果他们为基督养育你；在一切事上要听从他们，在他们的一切命令上要服从；只要他们不命令任何违背他们之上者的事，这样就应该服从他们。你会说：\Quote{谁高于生我的人呢？}那创造你的那位。因为人生育，但天主创造。人如何生育，他不知道；他将生育什么，他不知道。但是那在祂所造者存在之前就看见你以便造你的那位，当然高于你的父亲。你的祖国又应高于你的父母；因此，但凡你的父母命令任何反对祖国的事，就不应听从他们。但凡你的祖国命令任何反对天主的事，也不应听从。因为如果你愿意被治愈，如果你在血漏之后，如果你在患病十二年后，如果你在将一切花在医生上而未得健康之后，终于希望痊愈；啊，女人，我把你当作教会的预像，你的父亲命令你这样，你的民族那样。但你的主对你说：\BibleQuote{忘记你的民族和你的父家。}有什么好处？有什么益处？有什么有用的结果？\BibleQuote{因为君王恋慕你的美丽。}祂恋慕祂所造的，因为祂在你丑陋时就爱上了你，为能把你造得美丽。因为不信、丑陋的你，祂倾流了祂的血，使你成为忠信而美丽，祂在你们身上爱上了祂自己的恩赐。因为你给你新郎带来了什么？你从你从前的父亲和从前的民族那里得到了什么嫁妆？岂不是罪过的放纵和破烂衣裳？祂丢弃了你的破烂，撕裂了你污秽的长袍。祂怜悯你，为能装饰你。祂装饰你，为能爱你。\par

9. 弟兄们，还需要更多吗？你们是基督徒，并且听见了：\BibleQuote{如果你们得罪弟兄，并伤及他们软弱的良心，就是得罪基督。}不要不理会此事，如果你们不愿被从生命册上涂抹。我还要多久用光明悦耳的话向你们讲述，而我的悲伤却迫使我用某种方式说出，并不许我隐瞒呢？凡有意忽视这些事并得罪基督的人，让他们想想自己在做什么。我们希望其余的外邦人被聚集进来；而你们却是他们路上的绊脚石；他们有进来之意，却跌倒了，于是回头。因为他们心里说：\Quote{我们为什么要离弃那些连基督徒自己也同样敬拜的神明呢？}天主不许，你会说，\Quote{我怎么会敬拜外邦人的神明呢？}我知道，我明白，我相信你。但你在如何对待你正在伤害的软弱良心呢？你如何对待他们的价值，如果你轻视了救赎的代价呢？想一想救赎的代价是何等巨大。\BibleQuote{因着你的知识，}宗徒说，\BibleQuote{软弱的弟兄将丧亡；}你所自称的知识，即你知道偶像是虚无，你在心里只想到天主，于是便坐在偶像的庙里。在这知识中，软弱的弟兄丧亡了。并且免得你不关心软弱的弟兄，他补充道：\BibleQuote{基督为他死了。}如果你不屑于他，仍要思量他的代价，用基督的血来权衡整个世界。并且免得你仍然认为你得罪的是软弱的弟兄，因而当他听说那是\Person{伯多禄}时，就认为这是微不足道的过犯，他说：\BibleQuote{你得罪了基督。}因为人们习惯说：\Quote{我得罪了人，难道我也得罪了天主吗？}那么，你否认基督是天主吧。你敢否认基督是天主吗？你是在偶像庙里坐席时学到了这别的道理吗？基督的学校不接受那道理。我问：你在哪里学得基督不是天主？外邦人通常这样说。你看到坏的交往会带来什么吗？你看到\BibleQuote{不良的交往会败坏良好的品行}吗？在那里你不能谈论福音，而你会听到别人谈论偶像。在那里你失去了基督是天主的真理；你在那里所吸取的，却在教会里呕吐出来。也许你在这里大胆发言，大胆在人群中嘟囔：\Quote{基督不也是人吗？祂不也被钉死了吗？}这是你从外邦人学来的。你失去了你灵魂的健康，你没有触摸繸子。那么在这事上再触摸繸子，获得健康。正如我教导你在所记载的这事上触摸它：\BibleQuote{凡看见弟兄在偶像庙里坐席的}；也要触摸基督天主性的繸子。同一繸子论到犹太人说：\BibleQuote{圣祖是他们的，并且按照肉身，基督也是从他们来的，祂是在万有之上，永远受赞美的天主。}看，你犯罪反对的是谁？正是至高的天主，当你与假神同坐时。\par

10. 你可能会说：\Quote{这不是神，因为它是\Place{迦太基}的守护神灵。}就好像如果它是\Person{玛尔斯}或\Person{墨丘利}，它就会是神一样。但请考虑他们如何看待它，而不是它本身是什么。因为我和你一样知道，它不过是一块石头。如果这个\Quote{神灵}是一种装饰，那么让\Place{迦太基}的市民生活正直，他们自己就会成为\Place{迦太基}的这个\Quote{神灵}。但如果这个\Quote{神灵}是一个魔鬼，你已在同一部圣经中听到：\BibleQuote{外邦人所祭献的，是祭献给魔鬼，而不是给天主；我不愿意你们与魔鬼有来往。}我们很清楚它不是神；但愿他们也知道这一点！但因为那些不知道的软弱者，他们的良心不应受伤。这是宗徒告诫我们的。因为他们把那雕像视为某种神圣之物，并把它当作神，祭台就是见证。如果它不被视为神，祭台在那里做什么？不要有人对我说：它不是神，它不是天主。我已经说过：\Quote{但愿他们也知道这一点，如同我们所有人一样。}但他们如何看待它，把它当作什么，以及他们关于它做什么，那个祭台就是见证。它针对所有在那里敬拜者的意图定他们的罪，但愿它不会也定那些与他们一同坐席的人的罪！\par

11. 是的，如果外邦人压迫教会，不要让基督徒也压迫她。她是基督的身体。我们不是说过，基督的身体曾被压迫，却没有被触摸吗？祂忍受了那些压迫祂的人，并寻找那些\Quote{触摸}祂的人。并且，弟兄们，但愿如果基督的身体被外邦人压迫（他们惯于压迫她），至少基督徒不要压迫基督的身体。弟兄们，我的职责是对你们说话，我的职责是对基督徒说话；\BibleQuote{关於外人，我有什麼资格判断他们呢？}宗徒自己说。我们以另一种方式对待他们，如同软弱者。我们必须温柔地对待他们，使他们能听到真理；在你们身上，必须割除腐败。如果你们问如何赢得外邦人，如何光照他们，召叫他们得救恩：放弃他们的庆典，放弃他们虚浮的表演；然后，如果他们不认同我们的真理，就让他们为自己的贫乏而羞愧。\par

12. 如果那管辖你的人是个好人，他就是你的养育者；如果是坏人，他就是你的试探者。对前者，你要欣喜地接受养育；在试探中，你要证明自己是可嘉许的。成为金子。把这个世界视为金匠的熔炉；在一个狭窄的地方，有金子、糠秕、火。火作用于前两者，糠秕被烧毁，金子被炼净。一个人向威胁屈服，被带到偶像的庙里：唉！我哀悼糠秕；我看到了灰烬。另一个人尚未向威胁或恐惧屈服，被带到审判官面前，在宣认中站立稳固，没有向偶像屈膝：火焰对他做了什么？它岂不炼净金子吗？所以，弟兄们，你们要站稳在主内，召叫你们的那位更有能力。不要害怕不敬虔之人的威胁。忍受你们的仇敌；在他们身上，你们有可以为你们祈祷的对象；他们绝不可使你们恐惧。这是救恩的健康，在此盛宴中从这里汲取；在这里饮那能使你们满足的，而不是在其他那些盛宴中，只饮那能使你们癫狂的。站稳在主内。你们是银子，你们将成为金子。这个比喻不是我们自己的，是出自圣经。你们读过也听过：\BibleQuote{祂在熔炉中试炼了他们，如同黄金，并悦纳了他们如同全燔祭。}看看你们在天主宝藏中将成为什么。你们要在天主面前富足，不是使祂富足，而是从祂得到富足。让祂充满你们；不要在心中容纳别的任何事物。\par

13. 我们是否自高自大，或告诉你们要藐视那些被设立的权柄？不是这样。你们这些在此事上软弱的，是否也想摸一下衣边？宗徒自己说：\BibleQuote{每人要服从上级有权柄的人，因为没有权柄不是从天主来的，现有的权柄都是由天主规定的。所以，反抗权柄的就是反抗天主的规定。}但如果它命令你去做你本不该做的事呢？在这种情况下，绝对要因畏惧那更高的权柄而忽略这权柄。考虑这些不同级别的人间权柄。如果地方官命令某事，难道不应当执行吗？但如果他的命令与总督相悖，你肯定不是在藐视权柄，而是选择服从更大的权柄。在这种情况下，较小的权柄不应当生气，如果更大的被优先服从。同样，如果总督本人命令某事，而皇帝命令另一事，难道有任何疑问，我们应忽略前者而服从后者吗？所以，如果皇帝命令一件事，而天主命令另一件事，你们如何判断？\Quote{给我纳税，服从我的统治。}对，但不是在偶像的庙里。在天主的庙里，祂禁止它。谁禁止它？一个更大的权柄。那么请原谅我：你威胁监狱，祂威胁地狱。在此你必须立刻拿起你们的\BibleQuote{信德作盾牌，使你们能扑灭恶者的一切火箭。}\par

14. 但这些权柄中的一个正在密谋，对你图谋恶计。好吧：他只是在磨快剃刀，用来剃掉头发，而不是砍掉头。你们刚才在圣咏中听到我所说的：\BibleQuote{你图谋欺诈，有如锋利的剃刀。}为什么祂把恶人的欺诈比作剃刀？因为它只触及我们多余的部分。正如我们身体上的毛发似乎多余，剃掉它们对肉身毫无损失；同样，一个有权柄的人在愤怒中从你那里所能夺走的，只能算作你的多余之物。他夺走你的贫乏？他能夺走你的财富吗？你的贫乏是你心中的财富。他只能夺走你多余的东西，即使他被许可伤害身体，也只能伤害这些。甚至这现世的生命，对于那些思念另一种生命的人而言，这现世的生命，我说，也可算作多余之物。因为殉道者们就是这样轻视了它。他们没有失去生命，反而获得了生命。\par

15. 弟兄们，务要确信：仇敌不能加害于忠信者，除非是为使他们受试炼并获得益处。弟兄们，对此务必确信，不容任何人提出异议。将你们的一切挂虑都投靠于主，将你们自己完全彻底地交付于祂。祂必不撒手使你跌倒。那位创造了我们的天主，甚至就我们的头发也赐予了保证。\BibleQuote{我实在告诉你们：就是你们的头发，也都一一数过了。}我们的头发被天主数算过；那么，祂岂不更认识我们的行为吗？祂既如此深知我们的头发，岂不更明了我们的行径？请看，天主是如何不轻忽我们最微小的事。因为祂若轻忽，就不会创造它们。祂确实创造了我们的头发，并仍一一数点。但你或许会说：\Quote{它们现在虽被保存，将来却可能朽坏。}关于此，再听祂的话：\BibleQuote{我实在告诉你们：连你们的一根头发，也不会失落。}人啊，你既在天主的怀抱中，为何还惧怕人呢？不要从祂的怀抱中跌落；在那里你所受的一切，都将有助于你的救恩，而非毁灭。殉道者曾忍受肢体被撕裂的苦痛，难道基督徒反要惧怕基督时代的伤害吗？如今想要加害你的人，只能战战兢兢地行事。他不会公然说：\Quote{来参加偶像的宴席吧}；他不会公然说：\Quote{来我的祭坛前赴宴吧}。即便他这样说了，而你拒绝了，让他以此控告你，将之作为对你的指控和罪状：\Quote{他不肯来我的祭坛，不肯来我敬拜的神庙。}让他这样说吧。他不敢这么做；却以诡计另谋攻击。预备好你的头发吧；他正磨利剃刀，要剃去你多余之物，剃去那些你不久必将遗弃的东西。若有能力，让他剃去那永恒之物吧。这强敌究竟夺走了什么？他夺走了什么伟大的事物？那不过是一个窃贼或破门盗贼所能夺走的东西；在他极度狂怒之下，他所能夺的，也无非是强盗所能夺的。即便他获得许可杀害肉身本身，他所夺走的不过是强盗所能夺的吗？我称他为\Quote{强盗}，实在是抬举了他。因为无论那强盗是谁，他终是个人。他从你身上取走的，无非是一阵热病、一条毒蛇或一朵毒菌所能取走的。这便是人类狂怒的全部能力——所能做的事，不过是一朵毒菌所为！人吃了毒菌便死。看啊！人的生命何其脆弱；迟早你总要弃它而去。因此，不要如此为生命挣扎，以至于你自己反被遗弃。\par

16. 基督是我们的生命；那么，要思念基督。祂来，是为受苦，也是为得荣耀；祂受轻贱，也为受尊崇；祂死，也为复活。倘若劳苦使你恐惧，要看到它的赏报。你为何想不经艰辛就抵达那唯有艰苦劳作才能抵达之处呢？现在，你惧怕失去你的钱财，因为你以极大的辛劳赚得它。若你不经辛劳就获得那总有一天要失去（至少在你死时）的钱财，难道你竟想不经辛劳就获得永生吗？要让那经过一切劳苦后所获得、永不再失去的东西，在你眼中更为贵重。若这钱财（你历经劳苦所得，却终有失去之日）在你眼中如此贵重，那么，对那些永恒的事物，我们岂不更当渴慕？\par

17. 不要轻信他们的话，也不要惧怕他们。他们说我们是他们偶像的仇敌。愿天主恩准，并赐予我们能力，如同祂已使我们有能力砸碎的那些偶像一样。我这样说，亲爱的诸位，是使你们不要在不合法拥有能力时贸然行事；因为那些不守规矩的人及狂热的\OriginalTerm{寻衅派}{Circumcelliones}，惯于在没有能力时仍施暴，且常无端渴望死亡。你们曾听我们诵读给你们（凡当时在\Place{玛帕里亚}的人）的话：\BibleQuote{当那片土地被交到你们手中时（祂首先说\Quote{交到你们手中}，然后才吩咐当行之事），那时，你们要拆毁他们的祭坛，打碎他们的神柱，砍倒他们所有的偶像。}当我们得了能力，就这样做。若能力尚未赐予我们，我们便不做；一旦赐予，我们绝不怠慢。许多外教人在自己的田产上仍有这些可憎之物，难道我们去砸碎它们吗？不，因为我们首要的努力是使他们心中的偶像被砸碎。当他们自己也成为基督徒时，他们或是邀请我们行此善工，或是抢先自行毁除。目前，我们当为他们祈祷，而非向他们发怒。若我们心生痛楚，那更是针对基督徒，针对我们的弟兄——他们怀着一种心态进入教会：身体在这里，心却在别处。一切都应存于内在。若人所见的已是内在，那为什么天主所见的却是外在呢？\par

18. 亲爱的诸位，你们能够知道，这些人将自己的怨言与异端者和犹太人联合在一起。异端者、犹太人和异教徒已经联合起来反对合一。因为有时在某些地方，犹太人因自己的恶行而受到惩罚，他们便指责并怀疑我们，或假装相信我们总在寻求对他们同样的处置。又因为有时在某些地方，异端者因自己暴行的不敬和狂怒而受到法律的刑罚，他们立刻就说我们千方百计寻求某种毁灭他们的祸害。又因为已决定应制定法律反对异教徒——不，其实是为他们，如果他们明智的话。（因为正如愚蠢的孩童玩弄泥巴，弄脏了双手，严厉的教师走来，抖落他们手中的泥巴，并递给他们书本；同样，天主通过祂的臣仆——君王的手，来惊动他们幼稚愚昧的心，使他们扔掉手中的污秽，开始做些有用的事。这手中有用的事若不是\BibleQuote{把你的饼分给饥饿的人，将无家可归的穷人领到你家里}，又是什么呢？然而这些孩童却躲开教师的视线，偷偷溜回泥巴那里，被发现时就藏起双手不让人看见。）既然这样取悦了天主，他们便以为我们在各处搜寻偶像，并在所有发现偶像的地方将其捣毁。怎么会呢？难道没有一些偶像就在我们眼前的地方？或者我们真的不知道它们在何处？然而我们并未捣毁它们，因为天主没有把它们交到我们手中。天主何时将偶像交到我们手中？当这些事物的主人成为基督徒时。最近某地的主人正希望这样做。如果他本无意将那块地方本身交给教会，却只下令在他的产业上不得有偶像，我想应当以极大的热忱去执行，以帮助那位不在场的基督徒兄弟的灵魂，他愿在自己的土地上感谢天主，不愿那里有任何东西羞辱天主，他的同会基督徒应助他一臂之力。此外，在这次事件中，他将那块地方本身交给了教会。难道教会的产业上可以有偶像吗？弟兄们，请看，异教徒究竟不满什么？他们不满我们不去没收他们的产业，不捣毁偶像；他们甚至希望我们自己的地方也保留偶像。我们宣讲反对偶像，从人心中除去偶像；我们是偶像的迫害者，我们公开承认。难道我们要成为偶像的保护者吗？若我没有权力，我就不碰它们；若产业的主人抱怨，我就不碰它们；但若他愿意并为此感谢，我若不去做，就会招致罪过。\par

\SourceRecord{Translated by R.G. MacMullen.；Nicene and Post-Nicene Fathers, First Series；Vol. 6.；Edited by Philip Schaff.；Buffalo, NY: Christian Literature Publishing Co.,；1888.；<http://www.newadvent.org/fathers/160312.htm>.}

% structure: p0001 source=h1 target=section reason=page-title

\section{新约讲道集 第十三篇}

\Emph{[第六十三篇 本笃]}\par

\Emph{论福音的经文，\BibleRef{玛 8:23}，\BibleQuote{当祂上了船}等。}\par

1. 赖主的降福，我要以刚才宣读的福音教训来向你们讲话，并藉此劝勉你们：在这世界的风浪与波涛中，愿信德不要沉睡在你们心里。\Quote{因为主基督并非没有死亡或睡眠的能力，也许当祂航行时，睡眠违背祂的意愿胜过了全能者？}如果你们相信这事，祂就在你们里面睡着了；但如果基督在你们里面醒着，你们的信德就是醒着的。宗徒说：\BibleQuote{愿基督因信德住在你们心里。}\BibleRef{弗 3:17}因此，基督的睡眠是一个深奥奥秘的标记。水手是那些在木头上渡过世界的灵魂。那只船也是教会的预象。的确，每一个人都是天主的圣殿，他自己的心就是各人航行的船；只要思想是良善的，他就不会遭受船难。\par

2. 你听到了侮辱，那是风；你动怒，那是浪。因此，当风吹起，浪涌高涨时，船就危险了，心就危险了，心被颠簸摇荡。当你听到侮辱时，你渴望报复；看，你报复了，于是因着别人的损害而欢喜，你便遭受了船难。为什么呢？因为基督在你里面睡着了。这是什么意思：基督在你里面睡着了？你忘记了基督。那么，唤醒祂，记起基督，让基督在你里面醒着，留意祂。你曾渴望什么？报复。你是否忘记了，当祂被钉十字架时，祂说：\BibleQuote{父啊，宽恕他们，因为他们不知道自己做的是什么。}\BibleRef{路 23:34}那在你心中睡着了的祂，并不想要报复。那么，唤醒祂，记起祂。对祂的记念就是祂的话，对祂的记念就是祂的命令。然后，如果基督在你里面醒着，你会说：我是什么人，竟想报复！我是什么人，竟向别人发出威胁！也许在我报复之前我就会死去。而在我断气之时，充满愤怒、渴望报复而离开这身体时，那不愿报复的祂不会接纳我；那说\BibleQuote{你们给，就会给你们；你们宽恕，就会宽恕你们}\BibleRef{路 6:38}的祂不会接纳我。因此，我要抑制愤怒，回到内心的安宁。基督命令了海，平静就恢复了。\par

3. 现在我所讲的关于愤怒的话，你们要牢记，作为所有诱惑中的准则。诱惑发生，那是风；你心乱，那是浪。那么就唤醒基督，让祂对你们说话。\BibleQuote{这是谁？连风和海都听从祂。}\BibleRef{玛 8:27}这是谁？海也听从祂？\BibleQuote{海洋是祂的，祂创造了它。}\BibleRef{咏 95:5}\BibleQuote{万物是藉着祂造成的。}\BibleRef{若 1:3}那么，效法风和海，听从造物主。海听从基督的命令，而你却聋吗？海听从，风平息，而你还在吹吗？什么！我说，我做，我计划，这不就是在吹，而不愿服从基督的话吗？不要让你在这样动荡的心中屈服于浪。然而，既然我们是人，如果风吹向我们，激起我们灵魂的情感，不要绝望；让我们唤醒基督，好让我们在平静的海洋中航行，到达我们的家乡。\Quote{让我们转向主，等等。}\par

\SourceRecord{Translated by R.G. MacMullen.；Nicene and Post-Nicene Fathers, First Series；Vol. 6.；Edited by Philip Schaff.；Buffalo, NY: Christian Literature Publishing Co.,；1888.；<http://www.newadvent.org/fathers/160313.htm>.}

% structure: p0001 source=h1 target=section reason=page-title

\section{新约讲道第十四篇}

\Emph{[LXIV. Ben.]}\par

\Emph{论福音 \BibleRef{玛 10:16} 之言：\BibleQuote{看，我派遣你们好像羊进入狼群中}等等。在殉道者庆节上宣讲。}\par

当圣福音被诵读时，弟兄们，你们听到了我们的主耶稣基督如何以祂的教导坚固祂的殉道者，说：\BibleQuote{看，我派遣你们好像羔羊在狼群中}。现在，我的弟兄们，请思考祂做了什么。若只有一只狼进入许多羊中间，即使羊有千万只，也会因一只狼在它们中间而惊惶；尽管并非全部都会被撕裂，但全部都会恐惧。这是什么计划、什么谋略、什么能力，不是让狼进入羊群，而是派遣羊去攻击狼！祂说：\BibleQuote{我派遣你们好像羊在狼群中}；不是到狼群的附近，而是\BibleQuote{在狼群中}。那时曾有一群狼，却只有几只羊。因为当许多狼杀死了那几只羊时，狼就被改变而成了羊。\par

那么，让我们听听这位赐下冠冕却先指定战斗、作为战斗的观看者并协助他们劳累的那一位，给了什么劝告。祂规定了怎样的战斗？祂说：\BibleQuote{你们要像蛇一样机警，像鸽子一样纯朴}。凡明白并持守这点的，可以确信他不会真正死去而死。因为没有人应该在这种确信中死去，除非他知道他将要这样死去，以至只有死亡在他内死去，而生命得到冠冕。\par

因此，亲爱的诸位，尽管我已多次谈论过这个主题，但我仍要向你们解释，什么是\BibleQuote{像鸽子一样纯朴，像蛇一样机警}\BibleRef{玛 10:16}。如果鸽子般的纯朴被命令给我们，那么蛇的机警在鸽子的纯朴中有什么用呢？在鸽子中我所爱的是它没有苦胆；在蛇中我所怕的是它有毒液。但现在不要完全害怕蛇；它有你可憎恨的东西，也有你可模仿的东西。因为当蛇因年老而沉重，感到多年岁月的重负时，它便收缩并挤入一个洞中，脱去旧的皮衣，从而焕发新生。你这基督徒，听到基督说\BibleQuote{你们要从窄门进去}\BibleRef{玛 7:13}，要在这一点上效法它。宗徒保禄也对你说：\BibleQuote{脱去旧人和他的行为，穿上新人}\BibleRef{哥 3:9}。这样你就有蛇可模仿之处。不要为\Quote{旧人}而死，而是为真理而死。凡为任何暂时利益而死的，是死于\Quote{旧人}。但当你脱去了所有\Quote{那旧人}时，你就模仿了蛇的机警。再效法它这一点：\Quote{保护你的头}。这是什么意思，保护你的头？就是让基督与你同在。你们中或许有人观察到，当你想杀死一条蝰蛇时，它如何为了保护头而将整个身体暴露在攻击者的打击下？它不愿它的那部分被击中，因为它知道生命所在。我们的生命就是基督，因为祂亲口说：\BibleQuote{我是道路、真理、生命}\BibleRef{若 14:6}。这里宗徒也说：\BibleQuote{男人的头是基督}\BibleRef{格前 11:3}。所以，谁把基督保守在自己内，就是保护自己的头。\par

请看鸽子的争执，正如宗徒所说的：\BibleQuote{如果有人不听从我们这封信上的话，要记下他，不要与他交往}\BibleRef{得后 3:14}。请看这争执；但要注意，这是鸽子的争执，而不是狼的争执。他立刻加上：\BibleQuote{但不要以他为仇敌，而要劝戒他如弟兄}\BibleRef{得后 3:15}。因此，拥有鸽子的纯朴和蛇的机警，以清醒的心灵，而不是以肉体的放纵，来庆祝殉道者的庆节，向天主唱赞歌。因为那位是殉道者的天主，也是我们的主天主，祂将给我们加冕。如果我们善斗，我们将被祂加冕，祂已经加冕了那些我们所愿意效法的人。\par

\SourceRecord{Translated by R.G. MacMullen.；Nicene and Post-Nicene Fathers, First Series；Vol. 6.；Edited by Philip Schaff.；Buffalo, NY: Christian Literature Publishing Co.,；1888.；<http://www.newadvent.org/fathers/160314.htm>.}

% structure: p0001 source=h1 target=section reason=page-title

\section{论新约第十五篇讲道}

[LXV. Ben.]\par

\Emph{论福音之言}，\BibleRef{玛 10:28} \Emph{，\Quote{不要怕那些杀害肉身的人。}在殉道者庆节上宣讲。}\par

1. 刚才宣读的神圣训谕教导我们，在敬畏中不畏惧，在不畏惧中敬畏。你们在听圣福音时注意到，我们的主天主在为我们而死之前，愿意我们坚定；祂以此劝诫我们\Quote{不}要畏惧，同时又叫我们畏惧。因为祂说：\Quote{不要怕那杀害肉身，而不能杀害灵魂的。}看，祂在哪里劝我们不要畏惧。现在看祂在哪里劝我们要畏惧。\Quote{但}祂说，\Quote{更要怕那能把灵魂和肉身都投入地狱中的。}因此，让我们畏惧，以免我们畏惧。畏惧似乎与怯懦相关；似乎是弱者的特质，而非强者的。但请看圣经怎么说：\Quote{敬畏上主是力量的希望。}那么，让我们畏惧，以免我们畏惧；也就是说，让我们明智地畏惧，以免我们徒然畏惧。在圣殉道者的庆节上，为了他们的隆重典礼，读了这段福音；他们因畏惧而不畏惧，因为他们敬畏天主，就不顾念人。\par

2. 因为人何必畏惧人？一个人何以使另一个人畏惧呢？他们二者不都是人吗？一人威胁说：\Quote{我要杀你；}他并不害怕，恐怕他在威胁之后，未及实现就先死了。\Quote{我要杀你，}他说。谁说这话，对谁说？我听见两个人，一个威胁，一个惊慌；其中一个有权势，另一个软弱，但二者都是有死的。那么，他为何如此膨胀呢？他在尊荣上稍显膨胀的权势，在身体上却是同等的软弱。让那不怕死亡的人安然威胁死亡吧。但如果他害怕那使他害怕的事，就让他想想自己，把自己同他所威胁的人比较一下。让他看见他所威胁的人身上有一个相似的状况，于是同他一起向主祈求同样的怜悯。因为他不过是一个人，威胁另一个人，一个受造物，另一个受造物；只不过一个在天主眼中自高自大，另一个投奔同一天主。\par

3. 因此，愿刚毅的殉道者，作为一个人站在另一个人面前，说：\Quote{我不畏惧，因为我敬畏。}除非祂愿意，你不能做你所威胁的事；但祂所威胁的，无人能阻止祂做。再者，你威胁什么？你若被允许，又能做什么？你的暴力只及于肉身，灵魂是安全的。你不能杀害你看不见的东西：你是有形的，你威胁我身上有形之物。但我们有一个共同的有形无形的创造者，我们都应敬畏祂；祂从既可见又不可见之物创造了人。祂用地上的尘土造了有形的人，用祂的气息吹入无形之神。因此，无形的实体，即灵魂，它使躺卧的尘土升起，当你攻击尘土时，它并不害怕。你能打击居所，但你能打击住在那里的人吗？当锁链断裂时，原先被捆绑的人逃脱了，他将在暗中被加冕。那么，你为何威胁我呢？你丝毫不能加害于我的灵魂。因着那你所不能伤害之物的功德，你的权力所及之物将要复活。因为靠着灵魂的功德，肉身也要复活，并要归还给它的居住者，不再朽坏，而存留到永远。看（我用殉道者的话说），看，我说，甚至为了我的肉身，我也不怕你的威胁。我的肉身的确受你的权力支配，但我头发的数目都被我的创造者数过了。我何必害怕失去我的肉身呢？我连一根头发也不会失去。那连我最微小之事都如此清楚知道的一位，怎会不顾惜我的肉身呢？这肉身可能被创伤、被杀，暂时成为灰烬，但将来要永远不死。但这是为谁而设的呢？肉身即使被杀、被毁灭、被风吹散，仍要复活得永生，这是为谁呢？是为那不怕舍弃自己性命的人，因为他并不害怕自己的肉身被杀。\par

4. 因为，弟兄们，灵魂被称为不死的，按照它自己的某种方式，它是不死的；因为它是一种生命，能借其临在使身体获得生命。因为肉身因灵魂而活。这生命不能死，因此灵魂是不死的。那么，我为何说按照它自己的某种方式呢？请听原因。因为有一种真正的不死性，一种完全不变的不死性；关于此，宗徒论及天主时说：\Quote{只有祂是不死不灭的，住在不可接近的光中，没有人看见过，也不能看见；愿尊荣和永远的权能归于祂。阿们。}如果只有天主有不死性，那么灵魂必定是有死的。那么请看，我为何说灵魂按它自己的某种方式是不死的。因为事实上它也可能死。亲爱的，请理解这一点，就不会再有困难了。所以我敢说，灵魂能死，也能被杀。然而它无疑是不死的。看，我敢说，它既是不死的，又可能被杀；因此我说有一种不死性，一种完全的不变性，唯独天主有，经上说：\Quote{只有祂是不死不灭的；}因为如果灵魂不能被杀死，那么主自己为使我们要畏惧而说：\Quote{要怕那能在地狱里把灵魂和肉身都杀死的那一位。}\par

5. 迄今我证实了困难，却未解决。我已证明灵魂可以被杀。福音不可被反驳，除非被不敬的灵魂。看，此时我想到一件事，进入我的脑海要说出。生命不可被反驳，除非被死亡的灵魂。福音是生命，不敬与不信是灵魂的死亡。那么看，它能死，却仍是不朽的。它如何是不朽的？因为总有一种生命在它内永不被熄灭。它又如何死？并非因停止成为生命，而是因失去其生命。因为灵魂既是其他事物的生命，也有其自身的生命。考虑受造物的秩序。灵魂是身体的生命：天主是灵魂的生命。正如生命，即灵魂，与身体同在，以使身体不死；同样，灵魂的生命，即天主，也当与灵魂同在，以使灵魂不死。身体如何死？因灵魂离开它。我说，因灵魂离开，身体死了；它躺卧，只是一具尸体，先前是可欲的，如今是可鄙的对象。其肢体仍在：眼睛、耳朵；但这些只是房屋的窗户，其住客已去。哀悼死者的人，徒然在房屋的窗户呼喊，里面无人听见。哀悼者痴情的爱发出多少言语，列举并追忆多少事情；并以何等疯狂的悲伤，可以说，对着一仿佛感知所发生之事的人说话，而实际是对着一个已不在的人说话？他述说他的优点，以及他对自己的善意标记。是你给了我这个；为我做了这个那个；是你这样做了，如此深爱我。但若你肯思考并理解，克制你悲伤的疯狂，那曾爱你的人已去；你的敲击徒然落在房屋上，再找不到居住者。\par

6. 让我们回到刚才谈论的主题。身体是死的。为什么？因为它的生命，即灵魂，已离去。再者，身体活着，那人却不敬、不信、顽固、不可救药；在这种情况下，身体活着，而使身体生活的灵魂却死了。灵魂是如此卓越之物，即使死了也有能力给身体生命。我说，灵魂是如此卓越，如此卓越的受造物，即使自身死了，也有能力使身体活过来。因为不敬、不信、无度之人的灵魂是死的，但藉着这死的灵魂，身体却活着。因此它在身体内：它使手动工，使脚行走；指导眼睛看，安排耳朵听，辨别滋味，逃避痛苦，寻求快乐。所有这些是身体生命的标记；但来自灵魂的临在。若我问一个身体它是否活着，它会回答我：你看我行走，看我工作，听我说话，你察觉我有某些追求和厌恶，难道你不明白身体活着吗？藉着内在灵魂的这些作为，我明白身体活着。我也问灵魂它是否活着？它也有其固有的作为，藉以显示其生命。脚行走。我由此明白身体活着，因灵魂的临在。我现在问：灵魂活着吗？这些脚行走（仅就这一动作而言）。我在质问身体和灵魂，关乎它们的生命。脚行走，我明白身体活着。但它们走向何处？走向奸淫，据说。那么灵魂是死的。因为无误的圣经曾说：\Quote{那生活在宴乐中的寡妇是死的。}既然\Quote{宴乐}与奸淫之间差别甚大，那被称为死在宴乐中的灵魂，如何能在奸淫中活着？它肯定是死的。但即使不在这种情况下，它也是死的。我听见一个人说话；那么身体活着。因为舌头不能在口中自己移动，并通过其各种动作发出清晰的声音，除非内有居住者，且仿佛有乐师弹奏这乐器来使用他的舌头。我完全明白。如此身体说话；那么身体活着。但我问：灵魂也活着吗？看，身体说话，因此活着。但它在说什么？如我所说关于脚：它们行走，身体活着，然后我问：它们走向何处？好让我明白灵魂是否也活着。同样，当我听见一个人说话，我明白身体活着；我问他说什么，好让我知道灵魂是否也活着。他说谎。若如此，灵魂便死了。我们如何证明这一点？让我们请教真理本身，它说：\Quote{说谎的口杀害灵魂。}我问，为什么灵魂死了？如同刚才我问：为什么身体死了？因为灵魂，它的生命，已离去。为什么灵魂死了？因为天主，它的生命，已离开了它。\par

7. 经过这番简短考察后，要知道并确信：身体无灵魂是死的，灵魂无天主是死的。每个没有天主的人都有死的灵魂。你哀悼死者：倒不如哀悼罪人，倒不如哀悼不敬之人，倒不如哀悼不信之人。经上记载：\Quote{哀悼死者七日；为愚昧和不敬之人，则一生之久。}什么？难道你内没有基督徒的怜悯之心，以致你为灵魂离去的身体哀哭，却不哀哭灵魂，因天主已离开它？让殉道者记念此事，回答威胁他的人：\Quote{你为何强迫我否认基督？}那么你要强迫我否认真理？若我不肯，你要做什么？你要攻击我的身体，使我的灵魂离开它；但这同样的灵魂在我内，它拥有身体只是为了灵魂的缘故。它不是如此愚昧或不智。你会伤害我的身体；但你是否要使我因害怕你伤害我的身体、我的灵魂离开它，而伤害我自己的灵魂，使我的天主离开它？那么，殉道者，不要惧怕执行者的刀剑，只当惧怕你自己的舌头，免得你对自己行刑，杀害的不是你的身体，而是你的灵魂。要为你的灵魂惧怕，免得它死于地狱之火。\par

8. 因此，主说：\Quote{有权力把身体和灵魂都投入地狱之火的}。如何？当不虔敬的人被投入地狱之火时，他的身体和灵魂会在那里燃烧吗？永远的惩罚将是身体的死亡；天主的缺席将是灵魂的死亡。你想知道灵魂的死亡是什么吗？请听那位先知说：\Quote{让不虔敬的人被带走，使他不再看见主的荣耀。}那么，让灵魂害怕它自己的死亡，而不害怕身体的死亡。因为如果它害怕自己的死亡，从而活在天主内，不冒犯祂、不把祂从自己身边推开，它将被发现配得在末时再次接受自己的肉体；不是像不虔敬的人那样为了永远惩罚，而是像义人那样为了永生。由于害怕这种死亡，并热爱那种生命，殉道者们怀着天主许诺的希望，并蔑视迫害者的威胁，使自己获得与天主同享荣冠，并为我们留下了这些庆节的庆祝。\par

\SourceRecord{Translated by R.G. MacMullen.；Nicene and Post-Nicene Fathers, First Series；Vol. 6.；Edited by Philip Schaff.；Buffalo, NY: Christian Literature Publishing Co.,；1888.；<http://www.newadvent.org/fathers/160315.htm>.}

% structure: p0001 source=h1 target=section reason=page-title

\section{新约讲道集 第16篇}

\Emph{[LXVI. 本笃版]}\par

\Emph{关于福音的话}，\BibleRef{玛 11:2}\Emph{，\Quote{若翰在狱中听到基督的作为，便派门徒去问他说：\InnerQuote{你就是要来的那一位，还是我们该期待另一位？}}等等。}\par

1. 圣福音的读经给我们提出了一个关于洗者若翰的问题。愿主助我向你们解释清楚，如同祂已向我们解释了一样。若翰，如你们所听说的，曾受基督的称赞，并且受到这样的称赞：在妇女所生的人中，没有兴起比若翰更大的。但有一位由童贞女所生的，比他更大。大多少呢？让那位前驱自己宣告，他与他的审判者——他是其前驱——之间的差别何等巨大。因为若翰在出生和宣讲上都先于基督；但他是以服从的身份先祂而行，并非自认为先于祂。因为随从的整个行列都走在审判官前面；然而那些走在前面的人，实际上是在他之后。那么，若翰给了基督何等著名的见证呢？他甚至说：\Quote{不配解祂的鞋带}。还有什么？\Quote{从祂的满盈中，我们都领受了}，他说。他承认自己不过是靠着祂的光点燃的灯，于是他便投奔在祂的脚下，以免因高傲而登高，被骄傲的风吹灭。他确实如此伟大，竟被人当作基督；如果他本人没有作证自己不是基督，误会就会继续，他就会被认为是基督。何等惊人的谦卑！人民向他表示敬意，他却自己拒绝了。人们对他的伟大产生了错误，他却贬抑了自己。他不愿因人的言语而增长，因为他已领悟了天主的圣言。\par

2. 这就是若翰论及基督的话。那么基督论及若翰的话又是什么呢？我们刚才听到了。耶稣就对群众讲论若翰说：\Quote{祂开始对群众讲论若翰说：\InnerQuote{你们出去到旷野里看什么呢？被风吹动的芦苇吗？}}当然不是；因为若翰不被每一种教义之风吹来吹去。\Quote{你们出去到底要看什么？穿细软衣服的人吗？}不，因为若翰穿着粗糙的衣服；他穿的是骆驼毛的衣服，不是绒毛。\Quote{你们出去究竟要看什么？一位先知吗？是的，我告诉你们，而且比先知还大。}为什么\Quote{比先知还大}？先知们预言主将要来，他们渴望看见祂，却没有看见；但若翰却获得了他们所寻求的。若翰看见了主；他看见祂，用手指指着祂说：\Quote{看，天主的羔羊，除免世罪者；看，祂就在这里。}现在祂来了，却没有被认出来；所以连若翰本人也出现了误会。看哪，这就是族长们渴望看见、先知们所预言、法律所预表的。\Quote{看，天主的羔羊，除免世罪者。}他给主作了美好的见证，主也给他作了见证。\Quote{在妇女所生的人中，}主说，\Quote{没有兴起比洗者若翰更大的；但在天国里较小的，也比他大。}在时间上较小，但在尊威上更大。祂说这话，意指自己。洗者若翰在人间确实极其伟大，唯有基督比凡人更大。也可以这样陈述和解释：\Quote{在妇女所生的人中，没有兴起比洗者若翰更大的；但在天国里最小的，也比他大。}并不像我之前解释的那样。\Quote{但在天国里最小的，也比他大；}祂所说的天国是指天使们所在的地方；那么天使中最小的，也比若翰大。这样祂向我们显明了那应渴望之天国的卓越；将一座城摆在我们面前，我们应渴望作其公民。那里的公民是怎样的？他们有多大！谁在那里最小，也比若翰大。比什么若翰？\Quote{在妇女所生的人中，没有兴起比他更大的那位。}\par

3. 这样我们听到了若翰论基督和基督论若翰的真实而美好的见证。那么，当若翰被囚在监狱里，临死之前，派他的门徒到祂那里，对他们说：\Quote{去对祂说：\InnerQuote{你就是要来的那一位，或是我们该期待另一位？}}这又是什么意思呢？难道这一切赞美变成了怀疑吗？若翰啊，你说什么？你在对谁说话？你说什么？你对你的审判者说话，你自己却是前驱。你伸出手指指向祂，你说：\Quote{看，天主的羔羊，看，除免世罪者。}你说：\Quote{从祂的满盈中，我们都领受了。}你说：\Quote{我不配解祂的鞋带。}现在你却说：\Quote{你就是要来的那一位，或是我们该期待另一位？}这不是同一位基督吗？你又是谁？你不是祂的前驱吗？你不是那预言之所说的：\Quote{看，我派遣我的使者在你面前，预备你的道路。}你如何预备道路，而你自己却偏离了道路呢？于是若翰的门徒来了；主对他们说：\Quote{你们去告诉若翰：\InnerQuote{瞎子看见，聋子听见，瘸子行走，癞病人洁净，死人复活，穷人得听福音；凡不因我而跌倒的，是有福的。}}不要怀疑若翰在基督身上跌倒了。然而他的话听来的确如此：\Quote{你就是要来的那一位？}问我的作为吧：\Quote{瞎子看见，聋子听见，瘸子行走，癞病人洁净，死人复活，穷人得听福音；}你还问我是否就是那一位吗？我的作为，祂说，就是我的言语。\Quote{你们再去告诉他。当他们离开的时候。}恐怕有人会说：若翰起初是好的，但天主圣神离弃了他；所以在他派去的人离开之后，祂说了这些话；在若翰所派的人离开之后，基督称赞了若翰。\par

4. 那么这个隐晦的问题究竟是什么意思？愿那盏灯从它取得火焰的太阳照耀我们。所以这问题的解答是完全清楚的。\Person{若翰}有自己的门徒，并非与\Person{基督}分离，而是预备作祂的见证。因为如此一个人应当为\Person{基督}作证，他自己也在聚集门徒，并且可能对祂怀有嫉妒，因为他不能看见祂。因此，\Person{若翰}的门徒非常尊敬他们的师傅，他们从\Person{若翰}那里听到他关于\Person{基督}的见证，就感到惊奇；当他快要死时，他希望能通过\Person{基督}使他们得到坚定。毫无疑问，他们私下议论：他说了如此伟大的事关于祂，却对自己只字不提。\Quote{那么去吧，问祂；}不是因为我怀疑，而是为了使你们得到教导。\Quote{去吧，问祂，}从祂自己那里听到我惯常告诉你们的话；你们已经听到了报信者，让审判者来坚定你们。\Quote{去吧，问祂：\InnerQuote{祢就是那要来的，还是我们期待另一位？}}他们于是去了并询问；不是为\Person{若翰}的缘故，而是为他们自己。\Person{基督}为了他们的缘故说：\BibleQuote{瞎子看见，瘸子行走，聋子听见，癞病人洁净，死人复活，贫穷人听到福音。}你们看见我，就认出我来；你们看见这些工程，就认出行工程者。\BibleQuote{凡不因我而绊倒的，是有福的。}但我指的是你们，不是\Person{若翰}。因为我们要知道祂说这话不是关于\Person{若翰}，当他们离开时，\Quote{祂开始对群众讲论\Person{若翰}；}那真实的，真理本身，宣扬了他真实的赞美。\par

5. 我想这个问题已经解释得足够清楚了。那么就此延长我的讲道便已足够。现在要记住穷人。你们尚未施舍的人，请施舍吧；相信我，你们不会失去它。是的，确实，你们似乎只会失去那些你们没有带到竞技场去的东西。现在我们必须把你们中任何人的奉献交给穷人，而我们现有的数额比你们平日的奉献少得多。摆脱这种懒惰。我已成了乞丐的乞丐；这与我何干？我愿做乞丐的乞丐，使你们能被列入子女的行列。\par

\SourceRecord{Translated by R.G. MacMullen.；Nicene and Post-Nicene Fathers, First Series；Vol. 6.；Edited by Philip Schaff.；Buffalo, NY: Christian Literature Publishing Co.,；1888.；<http://www.newadvent.org/fathers/160316.htm>.}

% structure: p0001 source=h1 target=section reason=page-title

\section{新约讲道 第十七篇}

\Emph{[LXVII. Ben.]}\par

\Emph{论福音之言}，\BibleRef{玛 11:25}\Emph{，\BibleQuote{父啊，天地的主宰，我称谢祢，因为祢将这些事瞒住了智慧及明达的人，而启示给了小孩子。}等}\par

1. 当圣福音被宣读时，我们听到主耶稣在圣神内欢跃，并说：\BibleQuote{父啊，天地的主宰，我称谢祢，因为祢将这些事瞒住了智慧及明达的人，而启示给了小孩子。}首先说到这一点，在我们继续深入之前，如果我们以应有的注意、勤勉，尤其是以虔诚来思想主的话，就会发现：当我们读到圣经中的\Quote{宣认}时，我们不应一成不变地理解为罪人的宣认。此刻有必要说这番话，并提醒你们，亲爱的诸位，因为当诵读者的声音一说出这个词，随之而来的是你们捶胸的声音。我是说，当你们听到主所说的\BibleQuote{我称谢祢，父啊}时，你们捶胸。在说出\Quote{我称谢}这些词时，你们捶胸。但这捶胸意味着什么？岂不是要显露那隐藏在胸中的事物，并以这可见的捶打惩罚那隐秘的罪？你们为什么这样做？岂不是因为你们听到\BibleQuote{我称谢祢，父啊}？你们听到了\Quote{我称谢}这些词，却没有思考是谁在宣认。但现在请思考：如果基督，那与一切罪毫不相干的，说了\Quote{我称谢}，那么宣认就不只属于罪人，有时也属于那赞美天主的人。所以，我们宣认，无论是在赞美天主，还是在责备自己。无论哪种情况，都是虔敬的宣认：或是当你责备自己——你并非无罪，或是当你赞美祂——祂不可能有罪。\par

2. 但如果我们好好思考：你们自己的责备，就是祂的赞美。因为为什么你现在宣认，责备自己的罪？为什么你在宣认？岂不是因为你已从死亡中复活？因为经上记着：\BibleQuote{宣认从死人中消失，如同从不存在的人中消失。}如果宣认从死人中消失，那么宣认的人必是活的；如果他宣认罪，那他无疑已从死亡中复活。那么，如果那宣认罪的人已从死者中复活，是谁复活了他？没有死人能复活自己。唯有那一位能复活自己，祂的身体虽死，祂自己却没有死。祂复活了那已死的。祂复活了自己，祂在自己内是活的，但在那将要被复活的肉体内是死的。因为不仅圣父（关于祂，宗徒曾说：\BibleQuote{因此天主也高举了祂}）复活了圣子，而且主也复活了自己，即祂的肉体。因此祂说：\BibleQuote{拆毁这座圣殿，三天内我要把它重建起来。}但罪人是死的，尤其是那被罪恶习惯的重担压住的人，他就像拉匝禄一样被埋葬。因为他不仅死了，而且被埋葬了。那么，凡是被邪恶习惯、邪恶生活、世俗欲望的重担所压的人（我是指，在他身上应验了某篇圣咏中悲惨描述的：\BibleQuote{愚妄人心里说：没有天主}），他就是这样的人，正如经上所说：\BibleQuote{宣认从死人中消失，如同从不存在的人中消失。}除了那移开石头、大声呼喊说：\BibleQuote{拉匝禄，出来！}的那一位，谁能复活他呢？所谓\Quote{出来}，不就是把隐藏的东西带出来吗？那么，凡宣认的，就是\Quote{出来}。\Quote{出来}他不能，若非已活着；他不能活着，若非已复活。因此，在宣认中，责备自己，就是对天主的赞美。\par

3. 有人会说：如果那宣认的人立刻被主的声音复活而出来，那么教会有什么益处呢？那宣认的人从教会得到什么益处呢？主曾对教会说：\BibleQuote{凡你们在地上捆绑的，在天上也要被捆绑。}想想拉匝禄的例子：他出来了，却还带着缠布。他通过宣认已经活了，但还没有自由行走，因为他被缠布捆着。那么，教会——主曾对她说\BibleQuote{凡你们释放的，必要被释放}——做了什么？不就是主立刻对他的门徒说的\BibleQuote{解开他，让他走吧}吗？\par

4. 所以，无论我们责备自己，还是直接赞美天主，两种方式我们都赞美天主。如果我们怀着虔敬的意向责备自己，我们这样做就是赞美天主。当我们直接赞美天主时，我们就是在庆祝祂的圣洁，祂是无罪的。但当我们责备自己时，我们将光荣归于那使我们复活的天主。如果你们这样做，仇敌在审判者面前就找不到任何机会来欺骗你们。因为当你们成为自己的控告者，而主成为你们的解救者时，他除了是个纯粹的诬告者之外，还能是什么？基督徒借此有充分的理由为自己提供了保护，以对抗他的仇敌——不是那些可见的、血肉的（他们是应该可怜而不是可怕），而是对抗那些宗徒劝勉我们要武装自己去对付的仇敌。\BibleQuote{我们不是对抗血肉，}而是对抗那些看不见的仇敌。\BibleQuote{魔鬼进入了犹达斯的心，使他出卖主。}有人或许会说：我做了什么？请听宗徒的话：\BibleQuote{不要给魔鬼留余地。}\par

5. 因此，祂警告并说：\Quote{我们战斗不是对抗血肉，而是对抗率领者和掌权者。}有人可能以为这是指对抗地上的君王、对抗这世界的权势。怎么会呢？他们难道不是血肉吗？并且已经明确说：\Quote{不是对抗血肉。}请将你们的注意力从所有人身上移开。那么还剩下什么敌人呢？\Quote{对抗属灵邪恶的率领者和掌权者，这世界的统治者。}这似乎让魔鬼和他的使者得到了超出实际的地位。的确，祂称他们为\Quote{这世界的统治者。}但为了避免误解，祂解释了这个世界的含义，他们是这个世界的统治者。\Quote{这世界的统治者，这黑暗的统治者。}什么是\Quote{这世界的，这黑暗的}？世界充满了爱世界的人和不信者，魔鬼是他们的统治者。宗徒称此为黑暗。这黑暗是魔鬼和他的使者所统治的。这不是天然的、不可改变的黑暗：这黑暗会改变，变成光明；它相信，并因相信而被光照。当这发生时，它将听到这句话：\Quote{因为你们从前是黑暗，但现在在主内却是光明。}因为当你们是黑暗时，你们不在主内；同样，当你们是光明时，你们的光明不是出于自己，而是出于主。\Quote{你有什么不是领受的呢？}既然他们是看不见的敌人，就必须用看不见的方法来征服。可见的敌人你确实可以用打击来战胜；看不见的敌人你则依靠信仰来征服。人是可见的敌人；打击也是可见的。魔鬼是不可见的敌人；信仰也是不可见的。因此，对抗看不见的敌人，是一场看不见的战斗。\par

6. 面对这些敌人，有谁可以说自己是安全的呢？因为这是我开始谈到的，但我认为有必要稍为详细地论述这些敌人。现在我们知道我们的敌人，让我们看看如何抵御他们。\Quote{在赞美中我要呼求主，这样我便能从我的敌人那里得救。}你们知道该做什么。\Quote{在赞美中呼求；}即\Quote{在赞美主时呼求。}因为如果你们赞美自己，就不会从敌人那里得救。\Quote{在赞美中呼求主，你们将从敌人那里得救。}因为主自己说什么？\Quote{赞美的祭献将荣耀我，那里有一条路，我将在其中显示我的救恩。}路在哪里？在赞美的祭献中。那么，你们的脚不要偏离这条路。守在这条路上，不要离开它；从对主的赞美中，连一毫都不要偏离。因为如果你们偏离这条路，赞美自己而不赞美主，你们就不会从敌人那里得救；关于敌人，经上说：\Quote{他们在路上为我设下了绊脚石。}因此，无论你们认为自己有多少善是出于自己，你们已经偏离了对天主的赞美。那么，如果敌人诱骗你们，你们自己就是自己的诱骗者，这有什么奇怪呢？听宗徒的话：\Quote{如果有人自以为什么，其实什么都不是，他便是欺骗自己。}\par

7. 那么，请注意主承认说：\Quote{我称谢祢，圣父，天地的主宰。}我称谢祢，即我赞美祢。我赞美祢，并非控告自己。就取人性而言，一切都是恩宠，特殊而完美的恩宠。那成为基督的人有什么功德呢？如果除去这恩宠，甚至那卓越的恩宠——由此必须有一个基督，并且我们所承认的那一位就是祂——那么除去这恩宠，基督除了是一个纯粹的人，还能是什么呢？除了像你们自己一样，还能是什么呢？祂取了灵魂，取了身体，取了完美的人；祂将之与自己结合，主与仆人成为同一性体。这是何等卓越的恩宠！基督在天上，基督在地上；基督同时在天上地上；不是两个基督，而是同一个基督，既在天上又在地上。基督与圣父同在，基督在童贞女的胎中；基督在十字架上，基督在阴府里救助一些灵魂；同一日，基督与那悔悟的强盗同在乐园中。那强盗如何得到这福乐，难道不是因为他坚守了那条路，在那条路上\Quote{祂显示祂的救恩}吗？那是一条路，你们的脚不可偏离。因为他控告自己，赞美天主，并使自己的生命有福。他满怀希望地向主期待这事，并对祂说：\Quote{主，当祢进入祢的王国时，请记念我。}因为他考虑了自己的恶行，并认为即使最后蒙受怜悯也是极大的恩惠。但主在祂说\Quote{记念我}——何时？\Quote{当祢进入祢的王国时}——之后，立即说：\Quote{我实在告诉你们，今天你们要与我同在乐园里。}怜悯立刻赐予，而苦难的拖延却还在。\par

8. 那么，请听主称谢说：\BibleQuote{我称谢祢，父，天地的主宰。} 我称谢什么？我在什么事上赞美祢？因为这种称谢，如我先前所说，表示赞美。\BibleQuote{因为祢将这些事隐瞒了智慧和明达的人，而启示给了小孩子。} 这是怎么回事，弟兄们？借着与它们对立的事来理解。祂说：\BibleQuote{祢将这些事隐瞒了}智慧和明达的人，祂没有说：祢启示给了愚昧和不明达的人，而是\BibleQuote{祢将这些事隐瞒了}智慧和明达的人，\BibleQuote{而启示给了小孩子。} 对这些智慧和明达的人——他们实在是可笑的，是骄傲的人，虚假地自以为伟大，而实际上只是膨胀——祂对立的不是愚昧人，也不是不明达的人，而是小孩子。谁是小孩？谦卑的人。因此\BibleQuote{祢将这些事隐瞒了智慧和明达的人。} 以智慧和明达之人的名号，祂亲自解释了骄傲的人被理解，当祂说\BibleQuote{祢启示给了小孩子}时。因此，从那些不是小孩子的人那里，祢隐藏了它们。从那些不是小孩子的人那里是什么意思？从那些不谦卑的人那里。他们是谁，不就是骄傲的人吗？主的道路啊！要么没有，要么隐藏着，好启示给我们。主为什么欢喜？\BibleQuote{因为祢启示给了小孩子。} 我们必须成为小孩子；因为如果我们愿意成为伟大的，也就是所谓的\Quote{智慧和明达的}，它就启示给我们了。这些伟大的人是谁？智慧和明达的人。\BibleQuote{他们自称是智慧的，却成了愚笨的。} \BibleRef{罗 1:22} 这里你有一个由对立面提出的补救方法。因为如果由于\Quote{宣称自己是智慧的，你成了愚笨的；宣称自己是愚笨的，你将成为智慧的。} 但要真实地宣称，发自内心地宣称，因为事实正如你所宣称的。如果你宣称，不要在人前宣称，而在天主面前不要停止宣称。至于你自己和你的一切，你完全是黑暗的。因为成为愚笨的，除了是心中黑暗，还能是什么？祂最后论到他们说：\BibleQuote{他们自称是智慧的，却成了愚笨的。} \BibleRef{罗 1:22} 在他们宣称这一点之前，我们看到什么？\BibleQuote{他们愚笨的心就黑暗了。} \BibleRef{罗 1:21} 因此承认你自己不是自己的光。你至多是一只眼睛，而不是光。即使是一只睁开而健康的眼睛，若缺少光，又有什么益处？因此承认你自己并非是自己的光；并如经上记载的那样呼喊：\BibleQuote{上主，祢必点燃我的灯；上主，我的天主，祢必以祢的光照亮我的黑暗。} \BibleRef{咏 18:29} 就我自己而言，我完全是黑暗；但祢是驱散黑暗并照亮我的光；我自己不是自己的光，是的，我只有在祢内才有光的份。\par

9. 同样，\Person{若翰}，新郎的朋友，被认为是基督，被认为是那光。\BibleQuote{他不是那光，而是来为光作证的。} \BibleRef{若 1:8} 但那光是什么？它是真光。什么是真光？\BibleQuote{那照亮每个人的光。} \BibleRef{若 1:9} 如果那是照亮每个人的真光，那么它也照亮了\Person{若翰}，他正确宣告并承认说：\BibleQuote{从祂的圆满中我们都领受了。} \BibleRef{若 1:16} 你看出他是否说了别的话，而只是\BibleQuote{上主，祢必点燃我的灯}？最后，既已受光照，他作了见证。为了盲人的益处，这灯为白日光照作证。请看，他是一盏灯；祂说：\BibleQuote{你们曾派人到若翰那里去，他曾为真理作证。他是一盏点着而发亮的灯。} \BibleRef{若 5:35} 他，这灯，即一个被光照之物，被点亮以发光。能被点亮的也能被熄灭。为了不被熄灭，不要让它暴露在骄傲之风中。因此，\BibleQuote{我称谢祢，父，天地的主宰，因为祢将这些事隐瞒了智慧和明达的人，} 向那些自以为光是黑暗的人隐瞒；因为他们既是黑暗，又自以为光，甚至不能被光照。而那些是黑暗并承认自己是黑暗的人，是小孩子，不是伟大的；是谦卑的，不是骄傲的。因此他们正当地说：\BibleQuote{上主，祢必点燃我的灯。} \BibleRef{咏 18:29} 他们认识自己，他们赞美主。他们没有偏离救恩之路；\Quote{他们在赞美中呼求主，他们就从仇敌手中得救。}\par

10. 然后，转向主我们的天主，全能之父，以纯洁的心，我们愿尽我们软弱之所能，向祂献上我们最高而充足的感恩，以整个心神恳求祂的独一美善，愿祂按祂的善意，垂听我们的祈祷，愿祂以大能驱除我们行为和思想中的仇敌，扩大我们的信德，引导我们的心智，赐予我们属神的思想，并安全地将我们引向祂无尽的真福，藉着祂的圣子\Person{耶稣基督}。阿们。\par

\SourceRecord{Translated by R.G. MacMullen.；Nicene and Post-Nicene Fathers, First Series；Vol. 6.；Edited by Philip Schaff.；Buffalo, NY: Christian Literature Publishing Co.,；1888.；<http://www.newadvent.org/fathers/160317.htm>.}

% structure: p0001 source=h1 target=section reason=page-title

\section{新约讲道第十八篇}

\EditorialNote{[第六十八篇 本笃版]}\par

\Emph{再论福音之言，\BibleRef{玛 11:25}，\BibleQuote{父啊，天地的主宰，我感谢祢}等等。}\par

我们听到天主子说：\BibleQuote{父啊，天地的主宰，我感谢祢。}祂感谢祂什么？祂在哪里赞美祂？\BibleQuote{因为祢将这些事瞒住了智慧通达的人，而启示给了小孩子。}谁是\Quote{智慧通达的人}？谁是\Quote{小孩子}？祂瞒住了智慧通达的人什么，又启示给了小孩子什么？通过\Quote{智慧通达的人}，祂指的是圣保禄所说的那些人；\BibleQuote{智者在哪里？经师在哪里？今世的辩士在哪里？天主岂不是使今世的智慧变成了愚妄吗？}也许你还要问他们是谁。他们可能是那些对天主多有争论，却妄论祂的人；他们被自己的学说所膨胀，完全无法找到并认识天主；对于那本质不可理解、不可看见的天主，他们却以为空气和天空是天主，或太阳是天主，或任何在受造物中居于高位的东西是天主。因为他们观察受造物的伟大、美丽和能力，就停留在它们身上，而没有找到造物主。\par

\WorkTitle{智慧篇}斥责这些人，其中说：\BibleQuote{如果他们能知道这么多，以致能推测世界，怎么没有更早地发现它的主呢？}他们被控浪费时间和忙碌的争论，在研究和仿佛测量受造物；他们探求星辰的轨道、行星的间距、天体的运动，以便通过某种计算达到那样的知识，以致能预言日月食；并且正如他们所预言的，事件应按日、时和天体被蚀的部分应验。极大的勤奋，极大的心智活动。但在这些事中，他们寻求造物主，祂离他们不远，他们却没有找到祂。如果他们能找到祂，他们本可能拥有祂在自己内。因此，他们最有理由、非常正确地被控，他们能研究星辰的数目和它们多样的运动，并知道并预言发光体的亏蚀：我再说，他们被正确控告，因为他们没有找到创造和安排这些的那一位，因为他们忽视了寻求祂。但如果你不知道星辰的轨道和天体与地体的比例，不必过分不安。看世界的美，赞美造物主的智慧。看祂所造的，爱那造它的：这是你最应关心的事。爱那造它的，因为祂也按自己的肖像造了你，使你能够爱祂。\par

如果这很奇怪，那些基督说\Quote{祢将这些事瞒住了智慧通达的人}的事对这样的智者隐藏了，他们完全专注于受造物，选择漫不经心地寻求造物主，不能找到祂；那么更奇怪的是，竟有一些\Quote{智慧通达的人}能够认识祂。\BibleQuote{原来天主的忿怒，从天上发显，攻破一切不虔敬和人的不义，就是那些以不义抑制真理的人。}也许你问，他们以不义抑制什么真理？\BibleQuote{因为认识天主为他们是很明显的事。}如何明显？他接着说：\BibleQuote{原来天主已将自己显示给他们了。}你还要问，祂如何显示给他们，那些没有领受法律的人？如何？\BibleQuote{因为祂那看不见的事，从创造世界以来，藉着所造之物，就可以明白，都被看见了，就是祂永远的大能和祂的天主性。}那么，曾有一些这样的人，不像天主的仆人梅瑟，也不像许多先知，他们对这些事有洞见和知识，并由信德吸入、以虔诚的喉咙饮下、并由内在之人的口再涌出的天主之神所助。他们不像这样的人；但远远不如他们，他们藉这可见的受造物能够达到认识造物主，并对天主所造的这些事说：看祂所造的，祂也治理并容纳。祂造了它们，祂自己以祂的临在充满祂所造的。他们能够说这么多。因为保禄在\WorkTitle{宗徒大事录}中也提到了他们，当他论到天主说：\BibleQuote{因为我们生活、行动、存在都在祂内}（因为他是在对雅典人说话，那些有学问的人曾在那里），他立刻补充说：\BibleQuote{正如你们中有些人也说过。}他们所说的并非小事：\BibleQuote{我们在祂内生活、行动、存在。}\par

听我开始引用的宗徒的话：\BibleQuote{天主的忿怒}，他说，\BibleQuote{从天上发显，攻破一切不虔敬}（即，甚至对那些没有领受法律的人）；\BibleQuote{攻破一切不虔敬和人的不义，就是那些以不义抑制真理的人。}什么真理？\BibleQuote{因为认识天主为他们是很明显的事。}是谁使这明显？\BibleQuote{原来天主已将自己显示给他们了。}如何？\BibleQuote{因为祂那看不见的事，从创造世界以来，藉着所造之物，就可以明白，都被看见了，就是祂永远的大能和祂的天主性。}祂为什么显示？\BibleQuote{叫他们无可推诿。}那么，他们在什么上被责备？\BibleQuote{因为他们虽然认识天主，却没有以祂为天主而光荣祂。}\par

5. 这些话是什么意思？\BibleQuote{他们不当作天主而光荣祂？}他们没有感谢祂。那么，光荣天主就是感谢天主吗？是的，确是如此。因为如果你是按天主的肖像受造的，又认识了天主，却不感谢祂，还有什么比这更糟的呢？的确，光荣天主就是感谢天主。信友们知道在何时何处说：\Quote{让我们向主天主献上感谢。}但是除了那\Quote{向上主举起自己心}的人，谁能感谢天主呢？因此他们是应受责备、无可推诿的，\BibleQuote{因为他们虽然认识了天主，却没有当作天主而光荣祂，也没有感谢祂。但是}——什么？\BibleQuote{但是他们在自己的推理中变成虚妄的。}他们怎样变成虚妄的，还不是因为骄傲？烟因升腾而消散，火焰则越压低越明亮炽热；\BibleQuote{他们在自己的推理中变成虚妄的，他们愚昧的心就昏暗了。}同样，烟虽比火焰升得更高，却是黑暗的。\par

6. 最后，注意下文，看清整个问题的关键。\BibleQuote{他们自称为聪明，反而成了愚拙。}由于把天主所赐的归给自己，天主就收回了祂所赐的。因此，祂向骄傲者隐藏了自己，只向那些透过受造物勤恳寻求造物主的人传达对自己的认识。我们的主说得很好：\BibleQuote{你把这些事向聪明和通达的人隐藏起来；}无论是对那些在极多辩论和忙碌探究中完全查考了受造物，却不认识造物主的人，还是对那些认识了天主却不当作天主光荣祂、感谢祂，且因骄傲不能完美或健康地看见的人。\BibleQuote{因此，你把这些事向聪明和通达的人隐藏起来，却向婴孩启示出来。}什么样的婴孩？向谦卑的人。请说：我的圣神安息在谁身上？\BibleQuote{在谦卑、安静、对我话语战栗的人身上。}听了这些话，\Person{伯多禄}战栗；\Person{柏拉图}却不战栗。让渔夫紧握那位最著名的哲学家所失去的。\BibleQuote{你把这些事向聪明和通达的人隐藏起来，却向婴孩启示出来。}你向骄傲者隐藏，向谦卑者启示。这些事是什么？因为当祂说这话时，祂并不是指天地，也不是像说话时用手指着它们。因为这些谁看不见呢？善人看见，恶人也看见；因为祂\BibleQuote{叫太阳上升，光照恶人，也光照善人}。那么这些事是什么？\BibleQuote{我父把一切都交给了我。}\par

\SourceRecord{Translated by R.G. MacMullen.；Nicene and Post-Nicene Fathers, First Series；Vol. 6.；Edited by Philip Schaff.；Buffalo, NY: Christian Literature Publishing Co.,；1888.；<http://www.newadvent.org/fathers/160318.htm>.}

% structure: p0001 source=h1 target=section reason=page-title

\section{新约讲道集 第十九篇}

\Emph{[LXIX. Ben.]}\par

\Emph{关于福音中的话，\BibleRef{玛 19:28}，\Quote{凡劳苦和负重担的，你们都到我这里来}，等等。}\par

1. 我们在福音中听见：主满心喜乐于圣神，向天主圣父说：\BibleQuote{父啊！天地的主！我称谢你，因为你将这些事瞒住了智慧聪明的人，而启示给了小孩子。是的，父！因为你喜欢这样。我父把一切交给了我；除了父外，没有人认识子；除了子和子所愿意启示的人外，也没有人认识父。} 我说话劳苦，你们听也劳苦；愿我们一同听从那继续说话的主：\BibleQuote{凡劳苦的，你们都到我这里来。} 因为我们劳苦，岂不是因我们是必死的人，脆弱的受造物，软弱无力，携带着拥挤和互相挤压的瓦器吗？但如果这些血肉之器被挤压，愿爱德的广阔空间得以扩展。那么，他说\Quote{凡劳苦的，你们都到我这里来}，除非是说你们不再劳苦，还有什么意思呢？总之，他的应许十分清楚：既然他召唤那些劳苦的人，他们或许会问，被召有什么益处：他说，\Quote{我要使你们安息。}\par

2. \Quote{你们背起我的轭，跟我学吧！} 不是要建造世界的架构，不是要创造有形无形的万物，不是在已经创造的世界中行奇迹、复活死人；而是\BibleQuote{我是良善心谦的。} 你愿意成为伟大的，就从小处开始。你想建造高大的建筑，首先要思想谦卑的基础。无论你想要建造多么大的建筑，基础越深，建筑才能越高。建筑在建造过程中向高处升起，但挖基础的人必须先降到很低。因此，你看，建筑在升高之前是低的，顶部只在低后升高。\par

3. 在我们正在建造的那座建筑中，顶部是什么？这座建筑的最高点将到达哪里？我立刻回答：甚至达到看见天主。你看，看见天主是多么崇高、多么伟大的事。谁渴望它，就懂得我所说的和他所听的。看见天主的应许赐给了我们，就是看见那至高无上的天主。因为这是好事：看见那位看见者。因为那些崇拜假神的人，容易看见它们；但他们所看见的，是\BibleQuote{有眼却看不见}的。但应许给我们的是看见生活而看见的天主，好使我们热切渴望看见那位天主，经上论及他说：\BibleQuote{栽培耳朵的，难道自己听不到吗？形成眼睛的，难道自己看不见吗？} 那么，为你造了耳朵使你听见的，难道他不听吗？为你造了眼睛使你看的，难道他看不见吗？因此，在同一篇圣咏的前面的话中，他说：\BibleQuote{民间的愚人，你们要明白；愚昧的人，你们何时才能明智？} 因为许多人犯罪作恶，而以为天主看不见他们。事实上，他们很难相信他\Emph{不能}看见他们；但他们认为他\Emph{不愿}看见。很少有人如此不虔敬，以至于应验了经上所记：\BibleQuote{愚人在心里说：没有天主。} 这只是少数人的疯狂。因为正如大虔敬只属于少数人，大不虔敬也不例外。但大多数人说：什么！天主现在在想这个吗？他要知道我在家里做什么吗？天主关心我在床上选择做什么吗？谁说这话？\BibleQuote{民间的愚人，你们要明白；愚昧的人，你们何时才能明智？} 因为作为人，你要知道家里发生的一切，所有仆人的行为和言语都要传到你的耳中，这是劳苦的；你认为天主观察你，同样劳苦吗？他创造你的时候没有劳苦。难道他不注视你吗？他造了你的眼睛。你未曾存在，他创造了你，给了你存在；现在你存在了，难道他不关心你吗？他\BibleQuote{叫那不存在的成为存在的}。那么，你就不要对自己许诺这个。不管你愿意与否，他看见你，你无处可以躲避他的眼睛。\BibleQuote{如果我升上高天，他在那里；如果我下到地狱，他也在那里。} 你若不愿离开恶行，劳苦就大了；却又希望不被天主看见。真是艰难的劳苦！你天天想作恶，却怀疑自己不被看见吗？请听经上所说：\BibleQuote{栽培耳朵的，难道自己听不到吗？形成眼睛的，难道自己看不见吗？} 你能在哪里躲避天主的眼睛，隐藏你的恶行？如果你不愿离开它们，你的劳苦实在很大。\par

4. 那么请听祂说：\BibleQuote{凡劳苦的，到我这里来。}你不能通过逃避来解决你的劳苦。你愿意逃离祂，而不是到祂那里去吗？那么找出你能逃到哪里，然后逃吧。但如果你不能逃离祂，因为祂无所不在；那就逃向天主吧，祂就在你站立的地方。逃吧。看，你在逃跑中已经越过了诸天，祂在那里；你下到地狱，祂也在那里；无论你选择大地上的什么荒僻之处，祂都在那里，因为祂说过：\BibleQuote{我充满天地。}如果祂充满天地，没有地方你能逃离祂，那么停止你这劳苦，逃到祂的面前，免得你感到祂的来临。从希望中获得勇气：你将通过善生见到祂，而即使在你的恶生中，祂也看到你。因为在恶生中，你能被看见，却不能看见；但通过善生，你既被看见，又能看见。因为祂在怜悯中看到那不配的你，呼唤你，那么当祂加冕那配得的人时，将以何等更温柔的亲近来看顾你呢？\Person{纳塔乃耳}对尚不认识的祂说：\BibleQuote{你从哪里认识我？}主对他说：\BibleQuote{当你还在无花果树下时，我就看见了你。}基督在你自己的荫蔽下看见了你；难道祂不会在祂的光明中看见你吗？因为\BibleQuote{当你还在无花果树下时，我就看见了你}是什么意思？这意味着什么？要记起亚当的原罪，我们都在他内死去。当他初次犯罪时，他用无花果树叶给自己做了围裙，这些叶子象征着因犯罪而陷入的欲念的刺激。我们就是这样出生的；我们生在这种状况中；生在有罪的肉身内，只有\BibleQuote{有罪肉身的形状}才能治愈这肉身。因此\BibleQuote{天主派遣了自己的儿子，带着有罪肉身的形状。}祂来自这肉身，但祂不像其他人那样来临。因为童贞女因信德而非因欲念怀了祂。祂来到童贞女内，祂在童贞女之前就已存在。祂拣选了她，祂创造了她；祂创造了她，祂决意拣选她。祂给童贞女带来了生育，却没有夺走她完整无缺的纯洁。因此，祂来到你这里，没有无花果树树叶的刺激，祂在\BibleQuote{你还在无花果树下时}就看见了你。那么，预备好在你因怜悯而被看见的那一位的荣耀高峰中看见祂吧。但因为高峰是高的，要思想根基。什么根基？你说？\BibleQuote{你们跟我学习，因为我是良善心谦的。}在你内深深挖掘谦卑的根基，这样你才能达到爱德的加冕高顶。\Quote{转向主，}等等。\par

\SourceRecord{Translated by R.G. MacMullen.；Nicene and Post-Nicene Fathers, First Series；Vol. 6.；Edited by Philip Schaff.；Buffalo, NY: Christian Literature Publishing Co.,；1888.；<http://www.newadvent.org/fathers/160319.htm>.}

% structure: p0001 source=h1 target=section reason=page-title

\section{论新约讲道第二十篇}

[LXX. Ben.]\par

\Emph{再次论福音的话}，\BibleRef{玛 11:28} \Emph{，\Quote{凡劳苦和负重担的，你们都到我这里来，我要使你们安息。}等等}。\par

1. 弟兄们，有些人听到主说：\Quote{凡劳苦和负重担的，你们都到我这里来，我要使你们安息。你们背起我的轭，向我学习，因为我是良善心谦的，这样你们必要找到你们灵魂的安息。因为我的轭是柔和的，我的担子是轻松的。}他们觉得奇怪，认为那些无所畏惧地向这轭低下颈项，并极其顺从地把这担子放在肩上的人，在世上却遭受如此大的困难，以致他们似乎不是从劳苦被召到安息，而是从安息被召到劳苦；因为宗徒也说：\Quote{凡愿意在基督耶稣内虔敬生活的，都必遭受迫害。}所以有人会说：\Quote{轭如何是柔和的，担子如何是轻松的？}既然背这轭、挑这担子，无非就是在基督内虔敬生活？又如何说：\Quote{凡劳苦和负重担的，你们都到我这里来，我要使你们安息。}而不是说：\Quote{你们安逸闲散的人，来劳苦吧。}因为祂发现那些闲散安逸的人，就是祂雇来进葡萄园，好让他们承担一天炎热的人。我们听到宗徒在这柔和的轭和轻松的担子下说：\Quote{我们在一切事上，表明自己是天主的仆役：以极大的坚忍，在患难、贫乏、困苦、鞭打}等等，在同一封书信的另一处说：\Quote{被犹太人鞭打了五次，每次四十减去一下；被棍打了三次，被石头砸了一次；遭遇船难三次，在深海里度过了一夜一天}等等，以及其余的危险，这些危险诚然可以列举，但除非藉圣神的助佑，是无法忍受的。\par

2. 他所提到的所有这些痛苦和沉重的考验，他都极其频繁而丰富地承受了；但实际上，圣神与他同在，在外人日渐衰败的同时，内心却日日更新，并藉着天主的喜悦，在灵性安息的品尝中，以对将来福乐的希望软化一切现世的艰难，减轻一切沉重的考验。看，他背负的基督之轭是多么甘饴，担子是多么轻松；以至于他能说所有那些一听到就让人战栗的艰难和痛苦，不过是\Quote{轻微的苦难}；因为他以内在的眼目，信德的目光，看到永生的价值如此之高，必须以现世如此大的代价来换取，逃脱恶人永久的痛苦，完全享受义人永福而无任何焦虑。人们忍受切割和灼烧，以便通过更剧烈的痛苦来赎买，不是永恒的，而是某种比平常更持久的病痛。为了一个短暂而不确定的短暂休息时期，而且是在生命的尽头，士兵在战争的所有艰难考验中耗尽，或许在劳苦中度过了比他将安享休息更多的年头。忙碌的商人们将自己暴露在何种风暴和狂风、何种可怕而惊人的天空和海洋的狂暴中，以便获得像风一样无常、充满甚至比获得它们时所经历的更大的危险和风暴的财富！猎人们忍受着何种炎热和寒冷，何种来自马匹、沟渠、悬崖、河流、野兽的危险，何种饥渴的痛苦，何种低廉粗劣饮食的严格限制，只为捕获一只野兽！有时，经过这一切，他们为此忍受一切的野兽的肉对餐桌毫无用处。尽管捕获了一头野猪或一只雄鹿，对猎人来说，捕获的喜悦比食客品尝烹调后的美味更甜美。男孩们的稚嫩年龄被几乎每日鞭打的严厉纠正所驯服！在学校里，他们为了财富和虚荣荣誉，为了学习算术、其他文学和诡辩的欺骗，而不是学习真正的智慧，忍受着何等巨大的警醒和斋戒的痛苦！\par

3. 现在，在所有这些情形中，不爱这些事物的人感到它们是极大的严厉；而爱它们的人虽然也忍受同样的东西，却似乎不觉得严酷。因为爱使一切最艰难、最痛苦的事变得完全容易，几乎不算什么。那么，爱德为了真福，岂不更能确信而容易地成就那仅凭欲望在痛苦中勉强做到的事吗？若可避免永罚并获得永息，暂时的逆境是多么容易忍受！那\OriginalTerm{特选的器皿}{vessel of election}怀着极大的喜乐说：\BibleQuote{现时的苦楚，与将来在我们身上要显示的光荣，是不能较量的。}你看，那\BibleQuote{轭是柔和的，担子是轻的。}是如何实现的。虽然它对少数选择它的人是狭窄的，但对所有爱它的人却是容易的。圣咏作者说：\BibleQuote{因了祢口中的言语，我谨守了艰险的道路。}但那些对劳苦的人艰难的事，当这些人爱上它们时，就失去了粗糙。因此，天主良善的安排是这样制定的：使那\BibleQuote{内在的人却日日更新}，不再处于法律之下，而是在恩宠之下，并从无数仪规的重担中释放出来——这些仪规诚然是沉重的轭，但曾恰当地加在顽梗的颈项上——这样，那被逐的王子从外部加于外在之人的一切惨痛磨难，都应借着单纯的信德、良好的望德和神圣的爱德之轻松，因内在的喜乐而变为轻省。因为对于一个善意，没有什么比这善意本身更容易，而这对天主已经足够了。所以，无论这个世界如何狂怒，当主降生成人时，天使们最真实地高呼：\BibleQuote{天主受享光荣于高天，主爱的人在世享平安。}因为那时诞生的那一位，\BibleQuote{祂的轭是柔和的，祂的担子是轻的。}正如宗徒所说：\BibleQuote{天主是忠信的，祂不许你们受那超过你们能力的试探，而且随着试探也要开一条出路，叫你们能够承担。}\par

\SourceRecord{Translated by R.G. MacMullen.；Nicene and Post-Nicene Fathers, First Series；Vol. 6.；Edited by Philip Schaff.；Buffalo, NY: Christian Literature Publishing Co.,；1888.；<http://www.newadvent.org/fathers/160320.htm>.}

% structure: p0001 source=h1 target=section reason=page-title

\section{新约讲道集第二十一篇}

\Emph{[LXXI. Ben.]}\par

\Emph{论福音之言，\BibleRef{玛 12:32}，\BibleQuote{凡说话干犯圣神的，不论在今世或来世，决不得赦免。} 或 \Quote{论\OriginalTerm{亵渎圣神}{blasphemy against the Holy Ghost}。}}\par

1. 关于刚读过的福音篇章，曾有一个重大问题被提出，凭我自己的力量是无法解答的；但\Quote{我们的充足乃出于天主}，无论我们能在多大程度上接受祂的助佑。首先，请思考这问题的重大；如此你们看到它重压在我的肩上，就会为我的辛劳祈祷，而在赐予我的助佑中，为你们自己的灵魂找到益处。当\Quote{一个附魔的人被带到主前，又瞎又哑，主治好了他，使他能说话能看见，众人都惊讶说：这不是达味之子吗？法利塞人听见了，就说：这个人驱魔，除非仗赖魔王贝尔则步。耶稣知道他们的思念，就对他们说：凡一国自相纷争，必成荒芜；一城一家自相纷争，必站立不住。如果撒殚驱逐撒殚，便是自相纷争，他的国如何能站立得住？}在这些话中，祂愿意从他们自己的承认中让他们明白，由于他们不信祂，他们已选择属于魔鬼的国，而魔鬼的国自相纷争，因此不能站立。那么，让法利塞人任选其一吧。如果撒殚不能驱逐撒殚，他们就无从指摘主；但如果他能，那么他们更要反省自己，离开他的国，因为自相纷争的国不能站立。\par

2. 但为使不让他们以为主耶稣基督是靠魔王驱魔，请听下文：\Quote{如果我是仗赖贝尔则步驱魔，你们的子弟是靠谁驱魔？为此，他们将是你们的审判者。}祂这话无疑是指祂的门徒，那百姓的\Quote{子弟}；他们作为主耶稣基督的门徒，深知从善师那里没有学到任何邪术，以致能靠魔王驱魔。\Quote{为此，}祂说，\Quote{他们将是你们的审判者。}祂说，这世上一无所有的卑微之人，在他们身上看不到这种诡诈的恶意，只看到我权能的神圣纯朴；他们将是我的见证人，他们将是你们的审判者。然后祂继续：\Quote{如果我是仗赖天主圣神驱魔，那么天主的国已来到你们中间了。}这是什么意思？祂说：\Quote{如果我是仗赖天主圣神驱魔，而你们的子弟，我并未给他们有害诡诈的道理，只给了纯朴的信德，他们也只能如此驱魔；那么无疑天主的国已来到你们中间；由此魔鬼的国被颠覆，而你们也与它一同被颠覆。}\par

3. 在祂说了\Quote{你们的子弟是靠谁驱魔？}之后，为显示在他们身上是祂的恩宠，而非他们自己的功德，祂说：\Quote{或者，一个人怎能进入壮士的家，抢夺他的家具，除非先把那壮士捆住，然后才抢夺他的家？}祂说：你们的子弟，或已信了我，或将要信我，并驱魔，不是靠魔王，而是靠圣德的纯朴，他们从前或现在与你们一样，也是罪人及不虔敬的人；因此，在魔鬼的家里，作为魔鬼的器皿，他们怎能从他那被解救出来，因为那在他们身上作王的邪恶曾牢牢抓住他们，除非我被我的正义的锁链捆住那壮士，好从他那里夺走那些曾是愤怒的器皿，并使他们成为怜悯的器皿？这也是蒙福的宗徒在责备骄傲之人及那些自夸功德之人所说的：\Quote{谁使你与众不同？}也就是说，谁使你与出自亚当的毁灭之堆及愤怒的器皿分别出来？为使人不说\Quote{我的义德，}他说：\Quote{你有什么不是领受的呢？}关于这一点，论及自己他也说：\Quote{我们从前也是本性为义怒之子，如同其他人一样。}所以当他曾是教会的迫害者，\Quote{亵渎者、凌辱者、生活在恶毒与嫉妒中}（正如他承认的）时，他自己也是那壮士家中的一个器皿，在邪恶中强壮。但那捆绑壮士的，从他那里夺走了这毁灭的器皿，并使之成为蒙拣选的器皿。\par

4. 以后，为使不信主、不敬虔、即基督徒名号的敌人，不致因那些在基督徒名号下聚集迷失羊群的各色异端与分裂，而以为基督的国也是自相纷争的，祂接着补充说：\BibleQuote{那不与我的相合的，就是反对我的；不同我收集的，就是分散的。} 祂不是说，凡外表承认我名号或领受我圣事形式的；而是说：\BibleQuote{那不与我的相合的，就是反对我的。} 祂也不是说，凡不以外表承认我名号而聚集的；而是说：\BibleQuote{那不同我收集的，就是分散的。} 如此，基督的国并非自相纷争；但人却试图分裂那以基督宝血的代价所买来的。\BibleQuote{主认识属于祂的人。并且，凡称呼基督之名的，要离开不义。} 因为人若不离弃不义，即使他称呼基督之名，他也不属于基督的国。为举例说明，贪婪之灵和纵欲之灵，因一个聚敛，一个挥霍，是自相纷争的；然而两者都属于魔鬼的国。在偶像崇拜者中，\Person{朱诺}之灵和\Person{赫丘利}之灵是自相纷争的；两者都属于魔鬼的国。外邦人——基督的敌人，与犹太人——基督的敌人，是自相纷争的；两者都属于魔鬼的国。\OriginalTerm{亚略派}{Arianus}与\OriginalTerm{弗提诺派}{Photinianus}两者都是异端者，且两者都自相纷争。\OriginalTerm{多纳特派}{Donatist}与\OriginalTerm{马克西米安派}{Maximianist}两者都是异端者，且两者都自相纷争。人类一切彼此相反的恶习与错谬，都是自相纷争的，且全部属于魔鬼的国；因此他的国必站立不住。但义人与不虔敬者、信徒与不信者、\OriginalTerm{天主教徒}{Catholic}与异端者，的确自相纷争，但他们并不都属于基督的国。\BibleQuote{主认识属于祂的人。} 没有人当因一个名号而自夸。若他愿主的名号有益于他，就让他\BibleQuote{凡呼求主名的人要离开不义。}\par

5. 但这些福音的话，虽然有些隐晦（我想靠主的助佑已经解释了），却不若接下来似乎更困难的话那样难懂。\BibleQuote{为此，我告诉你们：一切的罪和亵渎，人都可得赦免；但亵渎圣神的罪，却不得赦免。凡说话得罪人子的，可得赦免；但凡说话得罪圣神的，今世及来世都不得赦免。} 那么，教会所想赢得的人将如何呢？当他们从任何错误中悔改并进入教会时，所应许他们赦免一切罪的希望，难道是虚假的吗？因为谁没有在成为基督徒或天主教徒之前就被定罪说过得罪圣神的话呢？首先，那些被称为外教人、崇拜众多假神、敬拜偶像的人，不是因他们说主基督藉魔法行奇迹，而像那些说祂藉魔王驱魔的人吗？再者，当天天亵渎我们的圣化时，他们不是亵渎圣神吗？怎么？犹太人——他们所说关于我们主的话引起了这场议论——他们不是直到今日还说亵渎圣神的话，否认祂如今在基督徒内，正如那些人否认祂在基督内一样吗？因为他们甚至没有这样不配地或类似地谈论圣神，如断言祂不存在，或虽然存在却不是天主而是受造物，或没有能力驱魔。其实，\Person{撒杜塞人}否认圣神；但\Person{法利塞人}反对他们的异端而主张圣神存在，却否认圣神在主耶稣基督内；他们认为主耶稣是藉魔王驱魔，而实际上祂是藉圣神驱魔。因此，犹太人以及所有承认圣神却否认圣神在基督身体内的异端者——这身体就是祂独一无二的教会，即唯一的天主教会——无疑如同当时的法利塞人，他们虽然承认圣神的存在，却否认圣神在基督内，将基督驱魔的作为归于魔王。我不提有些异端者胆大妄为地主张圣神不是创造者而是受造物，如\OriginalTerm{亚略派}{Arians}、\OriginalTerm{欧诺米派}{Eunomians}和\OriginalTerm{马其顿派}{Macedonians}，或者全然否认祂的存在，否认天主是圣三，而断言只有圣父是天主，有时称为圣子，有时称为圣神；如\OriginalTerm{撒伯流派}{Sabellians}，有人称他们为\OriginalTerm{父苦论派}{Patripassians}，因为他们认为圣父受苦；并且由于他们否认有任何圣子，无疑也否认圣神。\OriginalTerm{弗提诺派}{Photinians}又说只有父是天主，圣子仅是凡人，他们完全否认有圣神的第三位。\par

6. 显然，圣神也被外邦人、犹太人或异端者亵渎。那么，他们是否该被放弃，被认为毫无希望，因为判决已定：\BibleQuote{凡说话干犯圣神的，在今世及来世都不得赦免}？只有自幼就是天主教徒的人才能被认定无此最重罪过的罪债吗？因为所有相信天主圣言而成为天主教徒的人，必然是从外邦人、犹太人或异端者中皈依，来到基督的恩宠与平安中；若他们因所说干犯圣神的话不得赦免，那我们向人应许并宣讲归向天主、藉圣洗或在教会中获得平安与罪赦，便是徒然。因为经上并没有说：\Quote{除非藉着圣洗，不得赦免}；而是说：\BibleQuote{不得赦免，既不在此世，也不在来世}。\par

7. 有些人认为，只有那些在教会中藉重生的水泉受洗并领受圣神后，仿佛对救主如此伟大的恩赐忘恩负义，此后陷入任何大罪，如奸淫、杀人或完全背道，或完全背弃基督徒名号，或背弃天主教会者，才是干犯圣神。但我不知道这种理解如何证明；因为任何罪过在教会中并不否认悔改之处；而且宗徒说，连异端者自己也要受责备，为的是\BibleQuote{如果天主或许赐给他们悔改，得以认识真理；并且他们可以摆脱魔鬼的罗网，他们已被魔鬼俘虏，随从他的意愿}。因为不接受任何赦免希望的悔改有什么益处呢？最后，主并没有说：\Quote{受洗的天主教徒若说话干犯圣神}；而是说：\Quote{他，} 即凡说话者，无论他是谁，\BibleQuote{不得赦免，既不在此世，也不在来世}。所以，不论他是外邦人、犹太人、基督徒，或出自犹太人或基督徒的异端者，或任何其他错误称号，经上没有说这个人或那个人；而是说：\BibleQuote{凡说话干犯圣神的}，即亵渎圣神者，\BibleQuote{不得赦免，既不在此世，也不在来世}。此外，如果与真理相反、与基督徒平安为敌的每一种错误，如我们先前所表明的，都是\BibleQuote{说话干犯圣神}；然而教会并没有停止从各种错误中改造并聚集那些将要领受罪赦的人，以及他们所亵渎的圣神本身；那么，我认为我已经发现了一个重要的奥秘，可以用来阐明这个如此重大的问题。就让我们从主那里寻求解释之光吧。\par

8. 那么，兄弟姐妹们，请向我竖起你们的耳朵，并向主举起你们的心。我告诉你们，我亲爱的；或许在整个圣经中找不到一个更重要或更困难的问题。因此（为向你们坦白一下自己），我总是在向民众的讲道中回避这个问题的困难和窘迫；并非我对这个问题毫无想法，因为对于如此重要的事，我不曾疏忽去\BibleQuote{求}、\BibleQuote{寻找}、\BibleQuote{敲门}；而是因为我认为自己无法对所得到的某种程度的理解，用当时想到的言语恰当地表达出来。但当我听着今天的读经（我有责任向你们讲论它），当福音被宣读时，我的心跳动得如此厉害，以至于我相信是天主的旨意，要你们通过我的服务听到一些关于这个问题的道理。\par

9. 首先，我请你们思考并理解：主并没有说：\Quote{不得赦免任何对圣神的亵渎}，或\Quote{凡说话干犯圣神的，无论说什么，都不得赦免}；而是说\BibleQuote{凡说话干犯圣神}；因为若他说了前者，我们就没有任何辩论的余地了。因为如果对圣神没有任何亵渎、任何所说的话都不得赦免，那么教会就无法从所有不虔敬的罪人中——无论犹太人、外邦人、任何种类的异端者，甚至在天主教会内部一些知识浅薄的人——赢得任何人，这些人反对基督的恩赐和教会的圣化。但主绝不可这样说：天主禁止！我说，真理绝不可说：任何对圣神所说的亵渎和任何言语，既不在此世，也不在来世，都不得赦免。\par

10. 祂的旨意确是要藉问题的困难来锻炼我们，而非藉错误的判断来欺骗我们。因此，无人需要认为，凡是对圣神的亵渎或所说的每句话都不得赦免；但显然必须存在某种特定的亵渎，某种话若是对圣神说的，便永远无法获得宽恕和赦免。因为若我们将它理解为\Quote{每句话}，那谁能得救呢？但若我们认为没有这样的\Quote{话}，我们又与救主相矛盾。因此，无疑存在某种特定的亵渎和某种话，若是对圣神说的，便不得赦免。现在这究竟是什么话，是主愿意我们探究的；因此祂没有明说。我说，祂的意愿是它应当被探究，而非被否认。因为圣经的文风常常如此：当某事被表达得既不限定于普遍意义也不限定于特殊意义时，我们不必将其理解为普遍，而非特殊。那么，这个命题若要完全表达，即普遍地，就会说：\BibleQuote{凡对圣神的亵渎都不得赦免}；或者：\BibleQuote{无论谁说了任何话反对圣神，必不得赦免，无论在今世，或在来世。}但若是部分地，即特殊地表达，就会说：\Quote{某种对圣神的亵渎不得赦免。}然而，由于这个命题既非以普遍形式也非以特殊形式提出（因为并未说\Quote{凡亵渎}，或某种对圣神的亵渎；而只是不确定地说\Quote{对圣神的亵渎不得赦免}；同样也没有说\Quote{无论谁说了任何话}，或\Quote{无论谁说了某句话}，而是不确定地说\Quote{无论谁说了话}），所以我们不必理解\Quote{凡亵渎和每句话}；但显然主有意让我们理解某种亵渎和某种话，尽管祂不愿明说，好使我们若藉着祈求、寻找、叩门而获得正确的理解，就不致轻视它。\par

11. 为更清楚看到这一点，请思考同一位主论及犹太人的话：\BibleQuote{若我没有来向他们讲话，他们便没有罪。}因为这句话并非意指：若主没有来向他们讲话，犹太人就会完全没有罪。因为事实上，祂发现他们充满了并被罪重压着。因此祂说：\BibleQuote{凡劳苦和负重担的，你们都到我这里来。}负重担！背负什么？岂不是罪孽和违犯法律的负担？\BibleQuote{因为法律进入，是为叫过犯增多。}既然祂在另一处说：\BibleQuote{我不是来召义人，而是召罪人悔改。}那么，若祂没有来，他们\Quote{便没有罪}是什么意思呢？岂不是因为这个命题既非普遍地也非特殊地表达，而是不确定地，并不强迫我们理解为所有的罪？但除非我们理解为有一种罪是犹太人若基督没有来向他们讲话就不会有的，否则我们必须说这个命题是假的，这绝不可能。因此祂没有说：\Quote{若我没有来向他们讲话，他们就没有罪}，以免真理说谎。祂也没有明确地说：\Quote{若我没有来向他们讲话，他们就没有某种罪}，以免我们虔诚的热忱得不到锻炼。因为在圣经的丰富中，我们以明白的部分为食，以模糊的部分受锻炼；前者驱除饥饿，后者增添美味。既然没有说\Quote{他们没有罪}，我们无需不安，即使我们承认即或主没有来，犹太人也是罪人。但既然说\BibleQuote{若我没有来，他们便没有罪}，就必须是他们因主的来临而染上了某种（尽管不是全部）先前没有的罪。这实在就是：他们不相信临在并向他们讲话的祂，并且因祂说真理而视祂为仇敌，进而将祂置于死地。这一罪如此重大而可怕，显然若祂没有来向他们讲话，他们就不会有。因此，当我们听到\Quote{他们没有罪}这句话时，我们理解为并非所有罪，而是某些罪；同样，当我们今天所读的经文中听到\BibleQuote{对圣神的亵渎不得赦免}时，我们理解为并非所有亵渎，而是一种特定的亵渎；当我们听到\BibleQuote{无论谁说了话反对圣神，必不得赦免他}时，我们不应理解为每句话，而是某种特定的话。\par

12. 因为在此同一段经文中祂也说：\BibleQuote{但对圣神的亵渎不得赦免}；我们当然必须理解为并非对任何神的亵渎，而是对圣神。尽管祂在其他地方没有更明确地说明这一点，但谁会愚钝到以其他方式理解呢？根据同一说话规则，这句表达也被理解：\BibleQuote{除非人由水和圣神而生}。因为祂在那处没有说\Quote{和圣神}，但这一点仍被理解。也正因为祂说了\Quote{由水和圣神}，无人被迫理解为任何神。因此，当你听到\Quote{但对圣神的亵渎不得赦免}时，正如你不可理解为任何神，同样也不可理解为对圣神的任何亵渎。\par

13. 我看你们现在渴望知道，既然不是对圣神的每一种亵渎，那么那不得赦免的亵渎是什么，以及那句言语是什么，既然不是每一句言语若得罪圣神，就不得赦免，也不在这世上，也不在来世。至于我，我倒愿意立刻告诉你们你们如此专注等待听的事；但请忍耐片刻这更谨慎勤勉所要求的延迟，直到藉着主的助佑，我阐明眼前这段经文的全部意义。另外两位圣史，\Person{马尔谷}和\Person{路加}，在论及同一件事时，并没有说\Quote{亵渎}或\Quote{言语}，好让我们明白这不是指每一种亵渎，而是某一种亵渎；不是每一句言语，而是某一句确定的言语。那么他们说了什么？在\BibleBook{马尔谷}{福音}中这样记载：\BibleQuote{我实在告诉你们：人的一切罪和亵渎，无论他们怎样亵渎，都可获得赦免；但谁若亵渎了圣神，永远不得赦免，而要承担永久的罪责。}在\BibleBook{路加}{福音}中是这样：\BibleQuote{凡说话得罪人子的，尚可获得赦免；但那亵渎圣神的，决不能得赦免。}是否因为表达方式的不同，而有违同一命题的真理性？事实上，圣史们之所以不用同样的方式叙述同一件事，别无其他原因，而是让我们借此学会更重视事物而非言辞，不把言辞置于事物之上，并且从说话者那里只寻求他的意图，而言语仅用来传达这意图。说\Quote{圣神的亵渎不得赦免}，或说\Quote{那亵渎圣神的，必不得赦免}，究竟有什么实质区别呢？或许只是后者比前者更明确地宣告同一件事；因此一位圣史并不是推翻另一位，而是解释另一位。\Quote{圣神的亵渎}这个表述并不明确；因为没有直接说明是哪个神；因为并非每个神都是圣神。因此，当人用神魂亵渎时，可以称为\Quote{神魂的亵渎}；正如当人用神魂祈祷时，可以称为\Quote{神魂的祈祷}。因此宗徒说：\BibleQuote{我要用神魂祈祷，也要用理智祈祷。}但当说\Quote{那亵渎圣神的}时，这些歧义就消除了。因此，\Quote{永远不得赦免，而要承担永久的罪责}这个说法，若不是按\BibleBook{玛窦}{福音}所表达的\BibleQuote{不得赦免，也不在此世，也不在来世}，又是什么呢？完全相同的思想用不同的言语和不同的表达方式表达出来。至于玛窦所说的\Quote{凡说话得罪圣神的}，好让我们明白这不是指别的，而是亵渎，其他人则更清楚地表达为\Quote{那亵渎圣神的}。然而所有人都说了同样的事；没有一人偏离了发言者的意图，而正是为了理解这意图，言语才被说出、写下、阅读和听闻。\par

14. 但有人可能会说：你看，我承认并理解，当使用\Quote{亵渎}一词，且既未表明确指所有的亵渎，也未表明确指某一种亵渎时，它可以理解为所有的，或某一种的，但不一定是所有的；但另一方面，如果它不被理解为某一种，那么所说的话就不真实：同样，如果没有说每一句言语或某一句言语，就不必理解为每一句言语，但除非理解为某一句言语，否则所说的话绝不可能为真。但当我们读到\Quote{那亵渎的}，当\Quote{亵渎}一词未被使用时，我怎能理解某一种亵渎呢？或者当\Quote{言语}一词未被使用时，我怎能理解某一句言语呢？相反，它似乎是笼统地说\Quote{那亵渎的}。对此我如此回应：如果在这段经文中也说\Quote{那用任何亵渎亵渎圣神的}，那么我们就没有理由认为要寻找某一种特定的亵渎，而应当理解所有的亵渎；但之所以不能意指所有的亵渎，是为了不使外邦人、犹太人、异端者以及所有因各自的错误和矛盾而亵渎圣神的人，失去悔改后获得赦免的希望；由此毫无疑问，在记载着\Quote{那亵渎圣神的，永远不得赦免}的这段经文中，所指的人不是任何方式亵渎了的人，而是以那种特别的方式亵渎了，以致永不能得到宽恕的人。\par

15. 因为正如其中所说\BibleQuote{天主不诱惑人}，这不能被理解为天主绝不进行任何诱惑，而只是不进行某种特定的诱惑；否则那所记载的\BibleQuote{上主你们的天主试探你们}就成虚假，而我们读到祂问祂的门徒\BibleQuote{试探他；但他自己原知道要怎样行}时，也不能承认基督是天主，或说福音是假的。因为有一种诱惑引人犯罪，天主不进行这种诱惑；另有一种诱惑只考验我们的信德，甚至天主也屈尊用这种诱惑来试探人。因此，当我们听到\Quote{那亵渎圣神的}时，我们不可理解为每一种亵渎，正如在其他地方不可理解为每一种诱惑一样。\par

16. 同樣，當我們聽到：\BibleQuote{信而受洗的，必要得救}，我們當然不會理解為那些像\BibleQuote{魔鬼也信，且戰慄}那樣相信的人，也不會理解為那些像西滿魔術士那樣領受洗禮的人——他雖能受洗，卻不能得救。因此，當主說：\BibleQuote{信而受洗的，必要得救}時，祂所指的並非所有相信並受洗的人，而僅是某些人；即那些有份於那按宗徒的區別而\BibleQuote{以愛德行事}的信德的人。同樣，當祂說：\BibleQuote{凡褻瀆聖神的，永遠不得赦免}時，祂並非指任何一種褻瀆，而是指一種特定的褻瀆聖神的罪，凡被這罪所束縛的人，絕不會因任何赦免而被釋放。\par

17. 祂的另一句話：\BibleQuote{吃我的肉，喝我的血，便住在我內，我也住在他內}，我們該如何理解？我們能否將那些連宗徒也說\BibleQuote{他們吃喝自己的定案}的人——當他們吃這肉、喝這血時——也包括在這些話語中？怎麼！難道不虔誠的猶達斯，那位出賣並背叛其師傅的人（儘管按路加福音所更明確記載的，他與其他門徒一同吃了並喝了由主的手祝聖的祂體血的第一件聖事），他\BibleQuote{住在基督內，基督也住在他內}嗎？總之，那麼多要麼偽善地吃那肉、喝那血，要麼在吃了喝了之後成為背道者的人，他們\BibleQuote{住在基督內，或基督住在他們內}嗎？然而，確有一種吃那肉、喝那血的方式，凡如此吃和喝的人，\BibleQuote{他住在基督內，基督也住在他內}。因此，正如並非任何方式\BibleQuote{吃基督的肉、喝基督的血}的人都\BibleQuote{住在基督內，基督也住在他內}，而只有某種特定方式，主說這話時無疑所指的正是這種方式。同樣，在這句話中：\BibleQuote{凡褻瀆聖神的，永遠不得赦免}，並非任何方式褻瀆的人都犯下這不可赦免的罪，而是以那特定的方式，說出這真實而可怕判決的那位願意我們去尋求和理解的。\par

18. 現在，關於這種褻瀆的方式——或者不如說是無節制——是什麼，那種特定的褻瀆是什麼，以及那針對聖神的話語是什麼，我的講話次序要求我說出我的想法，不要再拖延你們那早已被推遲、但卻是有必要地推遲的期待。親愛的弟兄們，你們知道，在我們信德和至公教會所持守並宣講的那不可見、不朽壞的\OriginalTerm{天主聖三}{Trinity}中，天主聖父不是聖神之父，而是聖子之父；天主聖子不是聖神之子，而是聖父之子；而天主聖神不僅是聖父的聖神，也不僅是聖子的聖神，而是聖父和聖子的聖神。並且，這\OriginalTerm{天主聖三}{Trinity}，儘管每一位的\OriginalTerm{屬性}{Property}和\OriginalTerm{具體實存}{Subsistence}得以保存，但由於永恆、真理與美善的不可分割、不可分離的本質或本性，並非三位神，而是一位天主。以此方式，按我們的能力，以及按我們被允許\BibleQuote{藉著鏡子觀看，模糊不清}地看見這些事（尤其是我們現在這樣的人），我們被傳達了聖父中的\OriginalTerm{本源}{Origination}、聖子中的\OriginalTerm{受生}{Nativity}、以及聖神內聖父與聖子的\OriginalTerm{共融}{Communion}，以及在三位中的平等。因此，藉著那作為聖父與聖子之間\OriginalTerm{共融的紐帶}{Bond of communion}的那一位，他們樂意我們既在彼此之間、也與他們共融，並藉著那同一個恩賜（他們兩位所共同擁有的，即聖神，既是天主又是天主的恩賜）將我們聚集為一。因為在祂內，我們與天主性和好，並以之為樂。因為我們所認識的任何美善的知識，若非我們也愛它，對我們有何益處？但正如我們藉著真理學習，同樣我們藉著愛德去愛，以便我們也能達到更完全的認識，並在真福中享受我們所認識的。\BibleQuote{愛被傾注在我們心中，是因著所賜予我們的聖神。}並且，因為是藉著罪我們疏遠了對真正美善的擁有，\BibleQuote{愛遮蓋了許多罪過。}因此，聖父自己是聖子的真正本源，聖子就是真理，源出於真正的聖父；而聖神是美善，從美善的聖父和美善的聖子傾注而出；但在三位中，神性是同等，並且\OriginalTerm{不可分的合一}{Unity Inseparable}。\par

19. 首先，为使我们获得将于末后赐下的永生，从我们信仰的开端，天主的慈爱就赐给我们一项恩惠，即罪赦。因为只要这些罪尚存，在某种意义上，我们与天主的敌意和疏远就仍存留，这是由我们内的邪恶而来的；因为圣经不说谎话，它说：\BibleQuote{你们的罪孽使你们与你们的天主隔绝。}除非祂除去我们的邪恶，祂就不会将其美善赐给我们。前者随着后者的减少而增加；前者未达圆满，后者也未被终结。但现在主耶稣藉着圣神赦免罪过，正如祂藉着圣神驱逐魔鬼一样，这可以由以下事实理解：祂从死者中复活后，对门徒说：\BibleQuote{你们领受圣神吧，}随即加上：\BibleQuote{你们赦免谁的罪，就给谁赦免；你们保留谁的罪，就给谁保留。}因为那重生（其中一切过去的罪都被赦免）也是由圣神成就的，正如主说：\BibleQuote{人除非由水和圣神重生，不能进入天主的国。}但由圣神而生是一回事，由圣神滋养是另一回事；正如由肉身而生是一回事（发生在母亲分娩时），由肉身滋养是另一回事（发生在母亲哺育婴儿时，婴儿转动身子，从那里吸吮乳汁以获得生命，他从此获得生命支持，正如他从此获得出生的起源）。因此，我们必须相信，天主慈爱的第一项恩惠，在圣神内，就是罪赦。由此，作为主的前驱被派遣的若翰洗者的宣讲也以它开始。因为经上记载：\BibleQuote{那时，若翰洗者出现在犹太旷野宣讲说：你们悔改吧！因为天国临近了。}同样，我们读到主的宣讲的开始：\BibleQuote{从那时候，耶稣开始宣讲说：你们悔改吧！因为天国临近了。}若翰在向那些来受他洗礼的人所说的话中，也这样讲：\BibleQuote{我固然用水洗你们，为使你们悔改；但在我以后要来的那一位，比我更强，我连给祂提鞋也不配，祂要以圣神和火洗你们。}主也说：\BibleQuote{若翰固然用水施了洗，但不多几天以后，你们要因圣神受洗。}那是在五旬节。至于若翰所说的\Quote{火}，虽然也可能理解为信徒为基督之名将要遭受的苦难，但我们有理由认为，同一位圣神也被\Quote{火}这名所表达。因此，当祂来临时，经上说：\BibleQuote{有散开的火，像舌头似的，出现在他们面前，停在每一个人的头上。}所以主自己也说：\BibleQuote{我来是为把火投在地上。}所以宗徒也说：\BibleQuote{在圣神内，常要热情。}因为爱的热忱来自祂。\BibleQuote{因为天主的爱，藉着所赐与我们的圣神，已倾注在我们心中了。}与此热忱相反的，是主所说的：\BibleQuote{许多人的爱必要冷淡。}如今圆满的爱是圣神的圆满恩赐。但第一项\Quote{恩赐}与罪赦有关；藉着这项恩惠，\BibleQuote{我们被救脱离了黑暗的权势}；这世界的王，他以罪的共融和束缚在悖逆之子身上工作，却因我们的信德被\BibleQuote{驱逐出去}。因为藉着圣神（天主的子民藉祂被聚拢为一），那分裂的、反对自己的不洁之神被驱逐出去。\par

20. 不悔改的心反对这项无偿的恩赐，反对天主的恩宠。因此，这种不悔改就是\BibleQuote{亵渎圣神，今世及来世都不得赦免。}因为反对圣神（藉着祂，那些罪得赦免的人受洗，教会也领受了祂，好使\BibleQuote{凡她赦免的罪，可得赦免}），凡是在思想中或也在言语上说极其邪恶、极不虔敬的话的人——当他经历天主的耐心引导他悔改时，却依恃顽固和不悔改的心，为自己积蓄忿怒，直到义怒及显示天主公义审判的那天，祂要照各人的行为报应各人——都是在亵渎圣神。那么，这种不悔改（我们可以用一个总名称来称这种亵渎和那永远不得赦免的反对圣神的话），这种不悔改，我再说，是前驱和审判者都大声呼喊反对的：\BibleQuote{你们悔改吧！因为天国临近了。}主也是反对它公开了福音宣讲的口，并且预言了同一福音要传遍普世，当祂从死者中复活后对门徒说：\BibleQuote{默西亚必须受苦，第三天从死者中复活，并且必须从耶路撒冷开始，因祂的名向万邦宣讲悔改，以得罪之赦。}总而言之，这种不悔改\BibleQuote{今世及来世都不得赦免}；因为只有悔改在今世获得赦免，才能在来世有效。\par

21. 但是，只要人还生活在肉身之内，就不能断定这种不知悔改或不悔改的心。因为我们不应绝望于任何人，只要\BibleQuote{天主的耐心引导恶人悔改}，并不急忙将他逐出此生；\BibleQuote{天主，祂不愿罪人死亡，而愿他离开自己的道路而生活。}今天他是外邦人；但你怎么知道他明天不会成为基督徒？今天他是异端者；但明天他若追随公教真理又如何？今天他是裂教徒；但明天他若拥抱公教和平又如何？那些你现在看到在任何可能的错误中，并且你断定他们处于最绝望境地的人，如果他们在此生结束之前悔改，并在未来的生命中寻得真生命，又该如何？因此，弟兄们，也让宗徒的话激励你们：\BibleQuote{不到时候，什么也不要判断。}因为这种对圣神的亵渎，是没有赦免的（我理解这不是每一种亵渎，而是一种特定的亵渎，并且如我所说或发现的，甚至如我认为清楚表明的情况，是顽固不悔改之心的持续坚硬），只要人还在此生，就无法在任何人身上把握住它，我再说一遍。\par

22. 请不要觉得荒谬：一个人若在顽固不悔中持续到底，甚至直到此生终结，长久大量地反对圣神的这种恩宠；然而福音却将这如此长久的不悔改之心的反对，视作短暂之事，称为\Quote{一句话}，说：\BibleQuote{凡说话干犯人子的，还可赦免他；但凡说话干犯圣神的，今世来世总不得赦免。}因为虽然这种亵渎持续长久，由很多话语组成并延展，但圣经的习惯是将许多话语也称为\Quote{一句话}。因为没有任何一位先知只说一句话；然而我们读到：\BibleQuote{那\Emph{话}临到某位先知。}宗徒说：\BibleQuote{长老们应受加倍的尊敬，尤其是那些在圣言和教导上劳苦的人。}他并没有说\Quote{在言语上}，而是说\Quote{在圣言上}。圣雅各伯说：\BibleQuote{你们应按照圣言行事，不要只听从。}他也没有说\Quote{听言语}，而是说\Quote{听圣言}；尽管在教会庆祝和隆重的礼仪中，会诵读、宣讲和聆听许多圣经的话语。因此，无论我们中任何人宣讲福音多久，他不是被称为言语的宣讲者，而是圣言的宣讲者；无论你们中任何人如何专心致志地聆听我们的宣讲，他被称为最认真的\Quote{听者}，不是言语的听者，而是圣言的听者。同样地，按照圣经的风格和教会的习惯，凡人在整个肉身生命期间，无论时间多长，无论说了多少话，无论是用口说还是仅以不悔改之心思想，反对教会中所赐予的罪的赦免，他就是说了一句干犯圣神的\Quote{话}。\par

23. 因此，不仅一切说话干犯人子的，而且事实上，一切罪过和亵渎都将得到赦免；因为如果在一个人身上没有这种以不悔改之心干犯圣神的罪（借由圣神，罪过在教会中获得赦免），那么所有其他的罪都会得到赦免。但是，那种阻碍其他罪过赦免的罪，又如何能得到赦免呢？因此，一切罪过都得到赦免的，是那些没有这种永不得赦免之罪的人；但若有人有这种罪，既然这罪永不赦免，其他罪也不会赦免；因为一切赦免都被这一罪的束缚所阻碍。由此可见，并不是\BibleQuote{凡说话干犯人子的，必得赦免}，而是\BibleQuote{凡说话干犯圣神的，总不得赦免}；这并不是说在天主圣三中，圣神比圣子更大——这是连异端者也从未主张过的——而是因为，无论何人，即使在天主圣言——祂是子、人子和真正的基督——这样在人间显现，以致\BibleQuote{成了血肉，住在我们中间}之后，若抗拒真理并亵渎真理，如果他没有也说出那种以不悔改之心干犯圣神的话（关于圣神曾说过：\BibleQuote{除非由水和圣神而生}，又说：\BibleQuote{你们领受圣神吧！你们赦免谁的罪，就给谁赦免}），也就是说，如果他悔改，他就会借此获得赦免所有罪过的恩赐，也包括他\BibleQuote{曾说过干犯人子的话}的罪，因为他在无知、固执或无论何种亵渎的罪上，没有再加上对天主的恩赐以及由圣神在教会中所赐予的重生或和好的恩宠不悔改的罪。\par

24. 因此，我们也不可像某些人所想象的那样，以为说话干犯人子的，必得赦免，但说话干犯圣神的，却不得赦免，因为基督由于取了肉身而成为人子，在这方面圣神当然更大，祂在自己的本体中与圣父和独生子按神性同等，而独生子本身也按此与圣父和圣神同等。若果真如此，那么关于其他任何亵渎都必不会说什么，唯有那说话干犯人子（仅被视为人）的才可能获得赦免。然而，首先如此说：\BibleQuote{凡各样的罪和亵渎，都必赦免给人；} 在另一位福音作者那里也这样表达：\BibleQuote{凡人的一切罪和一切亵渎，无论怎样亵渎，都必得赦免；} 毫无疑问，那干犯圣父的亵渎也包括在那总括的表述中；然而，唯独那干犯圣神的被定为不得赦免。什么！难道圣父也取了仆人的形状，以致在这方面圣神比祂更大吗？决不是：但在普遍提及一切罪和一切亵渎之后，祂要更突出地表达干犯人子的亵渎，理由如下：虽然人甚至可能被捆绑在祂所说的那罪中，那时祂说：\BibleQuote{我若没有来对他们说话，他们就没有罪；} 在若望福音中，祂也表明这罪是非常严重的，当祂应许将要派遣圣神时，论到圣神祂说：\BibleQuote{祂要责备世界关于罪、关于义、关于审判：关于罪，是因为他们不信我；} 然而，若那不知悔改之心硬没有说出干犯圣神的话，就连这干犯人子的也必得赦免。\par

25. 也许有人会问：\Quote{是否只有圣神赦免罪过，而圣父和圣子却不赦免？} 我回答：圣父和圣子也都赦免。因为圣子论到圣父曾说：\BibleQuote{如果你们宽免别人的过犯，你们的圣父也必宽免你们。} 我们在主祷文中向祂说：\BibleQuote{我们的圣父，你在天上。} 在其它祈求中，我们也这样求：\BibleQuote{宽免我们的罪债。} 祂论到自己又说：\BibleQuote{为叫你们知道人子在地上有权赦罪。} 那么，\Quote{你会说：圣父、圣子和圣神都赦免罪过，为什么那永不蒙赦免的不知悔改只与亵渎圣神有关，好像那被捆绑在这不知悔改之罪中的人，似乎是在抵抗圣神的恩赐，因为罪的赦免正是藉着那恩赐完成的？} 关于这一点，我也要问：难道只有基督驱逐魔鬼，而圣父和圣神却不驱逐吗？因为如果只有基督，那么祂说的 \BibleQuote{住在我内的圣父，祂在做这些事} 又是什么意思？因为经上说 \BibleQuote{祂在做这些事}，好像圣子不做，而是住在圣子内的圣父在做。那么为什么在另一处祂又说：\BibleQuote{我的圣父一直工作，我也工作。} 过了一会儿：\BibleQuote{凡祂所做的，圣子也同样做。} 但在另一处祂说：\BibleQuote{我若没有在他们中间做过别人没有做过的事，} 祂说话好像只有祂独自在做。如果这些事情是这样表达的，而圣父和圣子的工作仍然是不可分的，那么对于圣神，我们该相信什么？岂不也相信祂同样与他们一同工作？因为就在那引发我们讨论的经文中，当圣子驱逐魔鬼时，祂却说：\BibleQuote{我若以圣神驱逐魔鬼，那么天主的国已临到你们了。}\par

26. 在此或许有人会说：圣神是由圣父和圣子所赐，而非凭自己的意志行任何事；这话的含义即是：\BibleQuote{在圣神内我驱魔}，因为不是圣神本身，而是基督在圣神内行了此事；因此\BibleQuote{我在圣神内驱魔}这句话可理解为\BibleQuote{我藉圣神驱魔}。因为这是圣经惯用的说法：\Quote{他们用刀剑杀害}，即藉刀剑杀害；\Quote{他们用火焚烧}，即藉火焚烧；\BibleQuote{若苏厄取了火石刀，用以给以色列子民行割损}，即用以行割损。但那些因此剥夺圣神本有德能的人，当注意我们所读的主所说的：\BibleQuote{风随意向哪里吹}。至于宗徒所说：\BibleQuote{这一切工作都是同一圣神所行，随祂的意愿分给各人}；恐怕有人会以为圣父和圣子不行这些事：然而在这些工作中，他明确提到了\BibleQuote{治病的奇恩}和\BibleQuote{行奇迹的能力}，其中无疑也包括驱魔。但当他加上\BibleQuote{随祂的意愿分给各人}时，岂不也清楚显示了圣神的德能，然而又同样明确地与圣父和圣子不可分？既然如此，这些事如此表述，并不妨碍我们理解天主圣三的行动是不可分的：因此当说到圣父的行动时，应理解为祂并非没有圣子和圣神而行动；当说到圣子的行动时，并非没有圣父和圣神；当说到圣神的行动时，并非没有圣父和圣子。这对那些有健全信德的人，或那些尽己所能理解的人，都足够清楚：既知\Quote{祂行这些工程}是就圣父而言，因为万有从祂而来工程的\OriginalTerm{第一首因}{first principle}，也是由祂而来共同行动者的存在；因为圣子由祂所生，圣神由祂所发，祂是\OriginalTerm{原始首因}{First Beginning}，圣子由祂所生，祂与圣子共有同一圣神；再者，当主说：\BibleQuote{我若没有在他们中间行过别人未曾行的事}，祂并非指圣父和圣神不参与那些工作；而是指世人，据我们所知，许多奇迹由他们行过，但没有任何奇迹如同圣子所行的。至于宗徒论及圣神所说：\BibleQuote{这一切工作都是同一圣神所行，随祂的意愿分给各人}，并非因为圣父和圣子不与祂合作；而是因为在这些工作中并非有众多神，而是独一圣神，在祂各种不同的行动中祂并不自相矛盾。\par

27. 然而，说\BibleQuote{你是我的爱子，我因你而喜悦}这句话的是圣父，而非圣子和圣神，这并非没有原因，而是有理有据的。尽管如此，我们并不否认圣子和圣神合作完成了这从天上发声的奇迹，尽管我们知道这仅属于圣父位格。因为虽然圣子带着肉体在地上与人交谈，但当那声音从云中发出时，祂同时也在圣父怀中作为\OriginalTerm{独生子圣言}{Only-Begotten Word}；也不能明智地且藉圣神相信，圣父将这些可闻且转瞬即逝的话语的行动与祂的智慧及圣神的合作分离。同样，当我们最正确地说，不是圣父，也不是圣神，而是圣子在海面上行走，因为只有祂拥有那肉体及其踏在波涛上的双脚；然而谁会否认圣父和圣神合作完成了如此伟大的奇迹呢？同样，我们最真实地说，只有圣子取了我们的血肉，而非圣父或圣神，然而若有人否认圣父或圣神在唯独属于圣子的\OriginalTerm{降生成人}{Incarnation}工作中合作，那便是没有真正的智慧。同样，我们也说，既非圣父，也非圣子，唯独圣神以\BibleQuote{鸽子的形状}和\BibleQuote{火舌}显现；并赐给祂所到之人能力，以各种不同的方言宣讲\BibleQuote{天主的奇事}；然而对于这唯独关乎圣神的奇迹，我们不能将其与圣父和\OriginalTerm{独生子圣言}{Only-Begotten Word}的合作分离。同样，整个天主圣三共同完成天主圣三中每位位格的工作，另外两位在另一位的工作中合作，通过天主圣三之间圆满的和睦行动，而非因其中一位有效行动能力的不足。既然如此，主耶稣在圣神内驱魔便是由此而来。并非祂不能单独完成此事，或祂借助那援助乃因不足以胜任；而是因为那分裂的魔鬼理当被那圣神所驱逐，这圣神是圣父和圣子所共有的，而圣父和圣子本身并不分裂。\par

28. 因此，罪恶在教会之外不得赦免，而必须由那使教会合而为一的圣神来赦免。事实上，若有人在教会之外为自己的罪忏悔，却对他与天主的教会隔绝这一如此重大的罪怀有不悔的心，那么那另一种忏悔对他有何益处呢？就凭这一点，他说的就是反对圣神的话，因为他与那领受此恩赐、即在其内藉圣神赐予罪之赦免的教会隔绝了。此赦免虽是整个天主圣三的工程，却被特别理解为属于圣神。因为祂是义子之恩的神，我们藉祂\BibleQuote{呼号：阿爸，父啊！}，使我们能向祂说\BibleQuote{赦免我们的罪债}。并且，\Person{若望}宗徒说\BibleQuote{我们由此知道}\BibleQuote{基督藉祂所赐给我们的圣神住在我们内}。\BibleQuote{圣神亲自与我们的心神一同作证：我们是天主的子女。}因为共融属于祂，我们藉此共融成为天主独生子的同一身体。因此经上记载：\BibleQuote{所以，若在基督内有任何鼓励，有任何爱的安慰，有任何圣神的共融。}为了这共融，祂最初降临于谁，谁就说万国的方言。因为正如藉着方言，人类的共融更为紧密地联合；同样，天主的子女和基督的肢体、这要在万民中实现的共融，也应藉万国的方言来表示；使得正如那时，凡说万国方言的人就被认为领受了圣神；同样现在，凡为那遍及万国的教会的和平联系所维系的人，就当承认他领受了圣神。因此宗徒说：\BibleQuote{竭力保持圣神的合一，以和平的联系。}\par

29. 如今，祂是圣父之神，圣子亲自说：\BibleQuote{祂由父而发。}在另一处：\BibleQuote{因为不是你们说话，而是你们父的圣神在你们内说话。}而且祂也是圣子之神，宗徒也说：\BibleQuote{天主派遣了祂圣子的圣神到你们心中，呼喊着：阿爸，父啊！}也就是说，使你们呼喊。因为是我们呼喊；但在祂内，即因祂将爱倾注在我们心中，没有祂，谁呼喊也是徒然。因此他又说：\BibleQuote{谁若没有基督的圣神，就不属于祂。}那么，在天主圣三中，这共融的相通还能特别属于哪一位呢？岂不属于那为圣父和圣子所共有的圣神吗？\par

30. \Person{犹达}宗徒已极其清楚地声明，那与教会分离的人没有这圣神，他说：\BibleQuote{那些自己分离的人，属血气，没有圣神。}因此，\Person{保禄}宗徒责备那些甚至在教会内、藉着人的名号、虽在教会合一中拥有地位却引起某种分裂的人时，在其他话中说道：\BibleQuote{然而属血气的人不能领受天主圣神的事，因为为他是愚妄；他也不能认识，因为这些事是藉圣神分辨的。}这显示他的意思是\BibleQuote{不能领受}，即不接受知识的言语。这些人在教会中有地位，他称他们为婴孩，尚未属神，仍是属肉体的，应当吃奶，不能吃饭。他说：\BibleQuote{甚至，}\BibleQuote{我如同对基督内的婴孩，给予你们奶，而非饭；因为直到现在你们还不能吃，就是现在你们也不能。}当我们说\BibleQuote{还不能}时，我们不可绝望，如果那\BibleQuote{还不能}的去向是成为。因为他说：\BibleQuote{你们仍是属肉体的。}并指出他们如何是属肉体的，他说：\BibleQuote{因为你们中间有嫉妒、纷争、分裂，你们岂不是属肉体的、按照人的方式行事吗？}并且更明确地说：\BibleQuote{因为有人说：我是属\Person{保禄}的，另有人说：我是属\Person{阿波罗}的，你们岂不是属肉体的人吗？\Person{保禄}算什么？\Person{阿波罗}算什么？不过是仆役，藉着你们信了。}这些人，即\Person{保禄}和\Person{阿波罗}，在圣神的合一与和平的联系中同心合意；然而因为格林多人开始在他们之间制造分裂，并\BibleQuote{互相自高自大}，他们被称为人——属肉体和属血气的人，不能领受天主圣神的事；然而因为他们没有与教会分离，他们被称为\BibleQuote{基督内的婴孩}；事实上，他希望他们成为天使，甚至成为神，他责备他们因为是人，即在那些争执中，\BibleQuote{他们不体味天主的事，只体味人的事。}但对于那些与教会分离的人，不仅仅说\BibleQuote{不能领受天主圣神的事}，以免这被指为知识的领受；而是说\BibleQuote{没有圣神}。因为拥有圣神的人未必也藉知识认识他所拥有的。\par

31. 那么，那些仍在教会有地位的\Quote{基督内的婴孩}——他们仍是属血气和属肉体的，不能\Quote{领受}，即理解和认识他们所拥有的——却有这圣神。因为除非他们由圣神重生，怎能成为基督内的婴孩？一个人可以拥有某物却不知道他拥有什么，这也不足为奇。因为且不提全能者的神性和不变圣三的合一，谁能轻易藉知识认识灵魂是什么？然而谁又没有灵魂呢？最后，为使我们确知，\Quote{基督内的婴孩}虽不能\Quote{领受天主圣神的事}，却仍拥有天主之神；让我们看\Person{保禄}宗徒，在稍后责备他们时，如何说：\BibleQuote{你们不知道你们是天主的殿，天主之神住在你们内吗？}这话他决不会对那些与教会分离、被描述为\Quote{没有圣神}的人说。\par

32. 但在心中怀有诡诈、仅凭肉身的交往而与基督的羊群混在一起的人，也不能说他在教会内，并属于那圣神的共融。因为\BibleQuote{智慧之神必远避奸诈}\BibleRef{智1:5}。所以，凡是在分裂者或异端者的聚会或分离中受洗的人，虽然他们未由圣神重生，如同依市玛耳一样——他是亚巴郎按肉身所生的儿子，而非像依撒格那样是因恩许按圣神所生的儿子——然而，当他们来到天主教会，并加入那教会之外肯定没有的圣神的共融时，他们肉身的洗不必重复。因为\BibleQuote{这敬虔的外形}\BibleRef{弟后3:5}即使在他们在外时也不欠缺；但加给他们的，却是\BibleQuote{在和平的联系中，圣神的合一}\BibleRef{弗4:3}，这只能在教会内赐予。事实上，在他们成为天主教徒之前，他们就如宗徒所说的那些人：\BibleQuote{虽有敬虔的外形，却否认了它的能力}\BibleRef{弟后3:5}。因为树枝有形的外形即使在离开葡萄树后仍然存在，但树根无形的生命只有在葡萄树内才能获得。因此，那些与基督身体的合一分离的人所施行和庆祝的肉身圣事，固然能给予\BibleQuote{敬虔的外形}\BibleRef{弟后3:5}；但敬虔无形而属神的能力绝不能在他们内，正如人的肢体从身体上被切除后，感觉也就不再伴随着它。\par

33. 既然如此，罪的赦免——既然它只能由圣神赐予——也只能在那拥有圣神的教会内赐予。因为罪的赦免的效果是：罪恶之君、那自相分裂的邪神不再在我们内为王，并且我们从那邪神的权下获得释放后，从此成为圣神的宫殿，并接受祂——借着接受宽恕我们由祂洁净——住在我们内，以施行、增长并成全义德。因为在祂初次降临之时，那些接受了祂的人说万国的方言，宗徒伯多禄向那些惊奇在场的人发言，他们心中刺痛，对伯多禄和其余宗徒说：\BibleQuote{诸位仁人弟兄，我们该作什么？}\BibleRef{宗2:37}指示我们。\BibleQuote{伯多禄对他们说：你们悔改，并要因耶稣基督之名受洗，好赦免你们的罪，并领受圣神的恩惠。}\BibleRef{宗2:38} 在那确实拥有圣神的教会内，这两件事都得以实现，即罪的赦免和这恩惠的领受。因此是\BibleQuote{因耶稣基督之名}\BibleRef{宗2:38}，因为当主应许同一圣神时，祂说：\BibleQuote{就是父因我的名所要派遣来的}\BibleRef{若14:26}。因为圣神不能没有父和子而住在任何人内；同样，子也不能没有父和圣神，父也不能没有他们。他们的内住是不可分的，正如他们的行动不可分；但有时他们藉由受造物的象征以分离的方式表现自己，而不是按其本体。就像他们在话语中以音节被分别宣示，这些音节各自占据自己的空间，然而在时间上没有间隔或片刻的分离。因为它们永远不能被同时宣诵，而它们却永远只能同时存在。但正如我已经说过，且不止一次，罪的赦免——借此那自相分裂的邪神的国度被推翻并驱除——以及天主教会合一的共融——在这共融之外没有罪的赦免——被视为圣神的特殊工作，当然有父和子的合作，因为圣神本身在某种意义上是父和子的共融。因为子与圣神不共同拥有父为父，因为祂不是两位的父。而父与圣神也不共同拥有子为子，因为祂不是两位的子。但是父与子共同拥有圣神为神，因为祂是两位的唯一圣神。\par

34. 因此，凡对抗圣神——教会的共融合一是藉着圣神聚集的——犯有\OriginalTerm{怙恶不悛}{impenitence}之罪的人，将永不得赦免；因为他为自己堵塞了宽恕的泉源，且理当与那分裂自身的灵一同被定罪，而他自己也因对抗那不分裂自身的圣神而分裂。福音本身的见证，只要我们留心查考，也警告我们这一点。因为按\BibleBook{路加}{福音}，主并非说：\BibleQuote{那亵渎圣神的，不得赦免}；而主是在回应那些说祂藉魔王驱魔的人时说的。由此看来，主并非仅一次说这话；但我们也不可轻忽地略过这最后一次说这话的场合的考量。因为当祂论及那些在人前承认或否认祂的人时，祂说：\BibleQuote{我又告诉你们，凡在人前承认我的，人子也要在天主的众天使前承认他；但在人前否认我的，在天主的众天使前也要被否认。}并且，免得因此而对宗徒伯多禄的得救绝望，祂随即加上：\BibleQuote{凡说话干犯人子的，还可赦免；但那亵渎圣神的，必不得赦免}——也就是说，以一颗怙恶不悛的心所行的亵渎，由此抗拒那藉圣神在教会中赐予的罪的赦免。伯多禄没有犯这亵渎之罪，因为他当时就悔改了，\BibleQuote{痛哭流泪}，而且他在胜过了那分裂自身、想要\BibleQuote{筛他如同麦子}的灵之后，——主曾\BibleQuote{为他祈求，使他的信德不至丧失}——他甚至领受了至圣的圣神，并未抗拒祂，好使他的罪不仅得到赦免，而且借着他，罪的赦免得以宣讲和分施。\par

35. 而在其他两位福音作者的叙述中，主说出这句关于亵渎圣神之判词的场合，乃因提及那分裂自身的不洁之灵。因为有人论及主说：\BibleQuote{祂是藉魔王驱魔}。在那里，主说：\BibleQuote{祂是藉圣神驱魔}，好使那不分裂自身的圣神得以胜过并驱逐那分裂自身的灵；但那因怙恶不悛而拒绝进入那不分裂自身者的平安中的人，将永留在他的丧亡中。因为\BibleBook{马尔谷}{福音}的叙述如此：\BibleQuote{我实在告诉你们：人的一切罪过和亵渎，无论怎样亵渎，都可赦免；但谁若亵渎了圣神，永远不得赦免，而要承担永罚的罪债。}主说了这些话之后，马尔谷又加上自己的话：\BibleQuote{因为他们说祂附有不洁之灵}；为表明主说这话的原因，是因为他们曾说\BibleQuote{祂是藉贝耳则步魔王驱魔}。这并不是说这亵渎是不可赦免的，因为即便是这亵渎，若随后有真诚的悔改，也可得到赦免；而是因为，如我所说，此处由此引发了主说出那判词的原因，因为提到了那不洁之灵，主表明他是分裂自身的；这是因为圣神不仅不分裂自身，反而使祂所聚集的人不分裂，赦免那些分裂自身的罪，并居住在那被洁净的人内，好使他们如\BibleBook{宗徒}{大事录}所记载：\BibleQuote{众信徒都是一心一意}。这赦免的恩赐，无人抗拒，除非是那心硬而怙恶不悛的人。因为另一处经文中，犹太人也说主附有魔鬼，但主在那里并未提及亵渎圣神的事；因为他们并未如此提及那不洁之灵，以致能从他们自己的口中显出他是分裂自身的，如同贝耳则步，他们说魔鬼是藉他而被驱逐的。\par

36. 但在\BibleBook{玛窦}{福音}的这段经文中，主更明白地解释了祂在此所要表达的意思：即，那以怙恶不悛的心抗拒\OriginalTerm{教会的合一共融}{Unity of the Church}的人，就是说话干犯圣神的人，因为在那合一共融中，罪的赦免是藉圣神赐予的。因为那些即使持守并施行基督的圣事，却与祂的集会分离的人，并不拥有这圣神，如前所述。因为当主论及撒殚自相分裂，以及祂自己是藉圣神——即那不分裂自身的圣神——驱魔时，惟恐有人因那些以基督之名、但不在祂羊栈内聚集非法集会的人，而以为基督的国也自相分裂，主随即加上：\BibleQuote{不与我相合的，就是反对我；不同我收集的，就是分散。}为表明那些在\Quote{外面}聚集、不愿\Quote{收集}、反而\Quote{分散}的人，不属于祂。之后祂又加上：\BibleQuote{所以我告诉你们：一切罪过和亵渎，人都可得赦免；但亵渎圣神的罪，必不得赦免。}这\Quote{所以}是什么意思呢？难道唯独亵渎圣神不得赦免，因为\Quote{不与我相合的就是反对我，不同我收集的就是分散}吗？确实如此。因为那不同祂收集的人，无论他以祂的名义如何聚集，都不拥有圣神。\par

37. 这样，祂就彻底迫使我们理解：任何罪过或亵渎的赦免，只能在基督的集会上生效，这集会不会四散分离。因为这集会是在圣神内聚集的，圣神不象那不洁之神，自相分裂。因此，所有自称是基督教会，却自相分裂、彼此敌对、与合一集会（即祂的真教会）为敌的集会，或不如说分散团体，并不属于祂的集会，因为它们似乎拥有祂的名号。但它们本可以属于它，如果圣神——这集会藉祂而结合——是自相分裂的话。但事实并非如此（\BibleQuote{因为不站在基督一边的，就是反对祂；不与祂聚集的，就是分散}）；因此，在这集会中，一切罪过和一切亵渎都将被赦免，这集会是基督在圣神内聚集的，圣神不自相分裂。但对圣神本身的亵渎，即藉一颗不知悔改的心，甚至到了今生的尽头，仍抗拒天主如此伟大的恩赐，将不得赦免。因为即使一个人如此反对真理，以致抗拒天主说话——不是在先知们中，而是在祂的独生子内（因为祂为了我们，乐意祂成为人子，以便藉祂向我们说话），当他悔改并投奔天主的仁慈时，仍将获得赦免；因为天主\BibleQuote{不乐意恶人丧亡，却乐意他离开自己的道路而生活}，已将圣神赐予祂的教会，使凡在圣神内赦免罪过的，都得赦免。但谁若以敌人的姿态对待这恩赐，以致不悔改寻求它，反而因不知悔改而反驳它，他的罪就成为不可赦免的；这不是某一特定种类的罪，而是对罪之赦免本身的蔑视，甚至反对。因此，当人们从不曾从分散中来到那接受圣神以赦免罪过的集会时，就是说了一句反对圣神的话。任何人只要以真诚来到这集会，即使是通过一个邪恶的圣职人员——一个被弃绝的和伪善的人——的服务，只要他是天主教司铎，他仍将在圣神内获得罪之赦免。因为圣神在圣教会中的运作就是如此，即使在现在，当麦子如同与糠秕一起被扬打时，祂不轻视任何人的真诚告解，也不被任何人的虚假伪装所欺骗，并因此远离被弃绝之人，却仍藉着他们的服务来聚集那些被悦纳的人。因此，对抗不可赦免之亵渎的唯一避难所，就是我们要提防一颗不知悔改的心；并且不要以为悔改能有什么益处，除非保持在教会内，那里赐予罪之赦免，并在和平的联系中保持圣神的共融。\par

38. 我藉着主的怜悯和助佑，尽我所能地处理了这个最困难的问题——如果我真的多少能做到的话。然而，无论我在这困难中未能领会什么，都不要归咎于真理本身——真理即使隐藏，对虔诚人也是有益的操练——而要归咎于我的虚弱，我或许既不能看到别人可能理解的，也不能解释我所理解的。但至于我可能通过默想的力量有所发现并用言语阐明的，必须感谢祂，我从祂那里寻求、祈求、敲门，使我能在默想中得到滋养，并在讲论中为你们服务。\par

\SourceRecord{Translated by R.G. MacMullen.；Nicene and Post-Nicene Fathers, First Series；Vol. 6.；Edited by Philip Schaff.；Buffalo, NY: Christian Literature Publishing Co.,；1888.；<http://www.newadvent.org/fathers/160321.htm>.}

% structure: p0001 source=h1 target=section reason=page-title

\section{新约讲道集第22篇}

\Emph{[LXXII. Ben.]}\par

\Emph{论福音的话，\BibleRef{玛 12:33}，\BibleQuote{或使树好，它的果子也好}，等等。}\par

主耶稣曾告诫我们，要我们成为好树，这样我们才能结出好果子。因为祂说：\BibleQuote{或使树好，它的果子也好；或使树坏，它的果子也坏，因为树是由它的果子被认出的。}当祂说：\BibleQuote{使树好，它的果子也好}时，这当然不是一句劝诫，而是一个有益的诫命，必须服从。但当祂说：\BibleQuote{使树坏，它的果子也坏}时，这不是要你做的诫命，而是一个劝诫，要你提防它。因为祂是反对那些认为虽然自己是恶的，却能说好话或行善事的人。主耶稣说这是不可能的。因为人自己必须先被改变，他的行为才能改变。因为如果一个人停留在恶的状态中，他不能有善行；如果他停留在善的状态中，他不能有恶行。\par

但谁被主看为好呢，既然\BibleQuote{基督为不虔敬的人死了}？祂发现他们全是坏树，但那些\BibleQuote{信祂名的人，祂赐他们权能成为天主的子女}。因此，现在无论谁是个好人，即一棵好树，原是发现自己坏，而被改好的。如果祂来时选择拔除那些坏树，还有什么树不配被拔除而得以存留呢？但祂先来施行仁慈，以便后来执行审判，正如对祂所说的：\BibleQuote{上主，我要歌唱仁慈与审判。}于是祂把罪的赦免赐给了信祂的人，祂甚至不与他们清算了过去的账目。祂赐予罪的赦免，使他们成为好树。祂推迟了斧头，祂给予了安全。\par

关于这斧头，\Person{若翰}说：\BibleQuote{现在斧头已放在树根上；凡不结好果子的树，必被砍下，投入火中。}福音中的家主用这斧头威吓说：\BibleQuote{看，这三年来我来到这棵树前，在这树上找不到果子。}现在我必须清理土地；因此，让它被砍下吧。园丁求情说：\BibleQuote{主，再留它这一年吧！等我在它周围掘土，加上肥料；如果结果子，便罢；不然的话，祢就来把它砍下。}所以，主好像三年访问了人类，即三次不同的时期。第一次是在法律之前；第二次是在法律之下；第三次是现在，即恩宠时期。因为如果祂没有在法律之前访问人类，那么\Person{亚伯尔}、\Person{哈诺客}、\Person{诺厄}、\Person{亚巴郎}、\Person{依撒格}、\Person{雅各伯}又从何而来呢？祂乐意被称为他们的主。祂，万邦都属于祂，却好像只是三个人的天主，说：\BibleQuote{我是亚巴郎的天主、依撒格的天主、雅各伯的天主。}但如果祂没有在法律之下访问，祂就不会赐下法律本身。在法律之后，家主亲自来了；祂受苦、死亡、复活了；祂赐下了圣神，使福音被传遍全世界，然而有一棵树仍然不结果子。仍然有一部分人类至今没有改过自新。园丁求情；宗徒为人民祈祷；他说：\BibleQuote{我屈膝……在天父面前，为使你们在爱德上根深蒂固，奠定基础，能同众圣人一同领悟什么是宽度、长度、高度和深度，并知道基督的爱超越知识，使你们充满天主的一切圆满。}通过屈膝，他为我们向家主求情，使我们不被拔除。因此，既然祂必然要来，我们当注意，使祂发现我们结果子。在树周围掘土是忏悔者的谦卑。因为每一条沟都是低的。加上肥料是忏悔的肮脏衣服。因为有什么比粪土更肮脏呢？然而如果善加利用，什么更有益呢？\par

那么，让每一个人都成为好树；不要以为如果他仍然是坏树，就能结出好果子。只有好树才能结出好果子。改变心，行为就会改变。拔除欲望，栽种爱德。因为\Quote{欲望是万恶的根源}，同样爱德是万善的根源。那么，为什么人们烦恼争斗，说：\Quote{什么是善？}哦，但愿你知道什么是善！你所希望拥有的并不太善；你所不愿成为的，才是善。因为你希望拥有身体的健康；这确是好的；然而你不能认为这是多么重大的善，因为恶人也同样拥有。你希望拥有金银；我承认这些也是善物，但惟有当你善用时才是；而如果你自己是恶的，你就不会善用它们。因此，金银对恶人是恶，对善人是善，不是因为金银使他们善，而是因为他们发现它们是善的，就善用它们。同样，你希望拥有荣誉，这是好的；但这也只有在你善用时才是。荣誉对多少人成了毁灭的时机！而荣誉对多少人又成了善事的工具！\par

5. 那么，我们若能，就对这些美善作一区分；因为我们所谈论的是好树。这里每人都最应思考的，莫过于将目光转向自己，在自己内学习、省察、检查、探索，并发现自己的真相；杀死那不喜悦的，渴望并种植那中悦（天主）的。因为人若发现自己如此缺乏更好的美善，为何还贪图外在的美善？装满美善的箱子，良心却是空虚的，又有何益？你愿拥有美善之物，难道不愿自己成为善人吗？你看不见，如果你的房屋满是美善之物，而你这主人却是邪恶的，你岂不更应为你的美善之物羞愧？告诉我，你愿拥有的是什么东西？我敢肯定，没有一样是恶的；既不是妻子，也不是儿子、女儿、仆役、婢女、乡间别墅，也不是外衣，甚至不是鞋子；但你却愿意拥有邪恶的生活。我求你，将你的生活方式置于你的鞋子之上。环绕你视线的一切优雅美丽之物，你都珍宝一般重视；但你对自己却如此轻视，如此缺乏美丽？若你房屋中充满的美善之物，你渴望拥有、害怕失去的，能够回答你，它们岂不向你呼喊：正如你愿我们为善，我们也愿拥有善的主人？如今它们以无声的话语向你的主控告你：\Quote{看！你给了他如此多的美善之物，他本人却是邪恶的。他拥有这一切，却不拥有那赐给他一切的主，于他有何益处？}\par

6. 那么，一个人若曾受劝告，或许被这些话触动悔改，可能会问：什么是美善？美善的本性是什么？它从何而来？你明白询问这事是你的本分，这很好。我将回答你的疑问，并说：\Quote{不能违背你的意愿而失去的，那才是美善。}因为金子你可能违背意愿而失去；房屋、荣誉，甚至身体的健康，也是如此；但那使你真正成为善人的美善，你既不会违背意愿而接受，也不会违背意愿而失去。于是我追问：\Quote{这美善的本性是什么？}一篇圣咏教导我们一件重要的事，或许正是我们寻求的。因为经上说：\Quote{人子们，你们心重要到何时？}那棵树要在这三年中不结果实多久？\Quote{人子们，你们心重要到何时？}什么是\Quote{心重}？\Quote{你们为何喜爱虚浮，寻求谎言？}接着它继续说我们必须真正寻求的：\Quote{你们要知道，主已显扬了他的圣者？}如今基督已来临，已被显扬，已复活升天，他的名传扬于普世：\Quote{你们心重要到何时？}过去的岁月已足够；如今那位圣者已被显扬，\Quote{你们心重要到何时？}三年之后，剩下的不就是斧头吗？\Quote{你们心重要到何时？你们为何喜爱虚浮，寻求谎言？}这些虚妄、无用、轻浮、转瞬即逝的东西仍在被寻求，而基督、那位圣者已被如此显扬？真理如今在大声呼喊，虚浮仍被寻求吗？\Quote{你们心重要到何时？}\par

7. 这个世界受严重的鞭打，是理所当然的；因为世界如今已经知道它主人的话。他说：\Quote{那仆人}，他说，\Quote{不知道主人的旨意，却做了该受鞭打的事，要受少些鞭打。}为什么？好使他寻求主人的旨意。那么，那不知道他旨意的仆人，就是世界，在\Quote{他显扬了他的圣者}之前；它是\Quote{那不知道主人旨意的仆人}，因此\Quote{要受少些鞭打。}但那如今知道主人旨意的仆人，就是现在，自从天主\Quote{祝圣了他的圣者}，却\Quote{不行他的旨意，要受多些鞭打。}那么，若世界如今大受鞭打，又何足为奇？\Quote{这就是那知道主人旨意，却做了该受鞭打之事的仆人。}因此，让他不要拒绝受多些鞭打；因为若他在不义中不愿听从他的老师，他必在正义中感受到他的报应者。至少，当他看到自己该受鞭打时，不要抱怨那责打他的主，好使他能获得怜悯；因我们的主耶稣基督，他和天主圣父及圣神，永生永王，万世无疆。阿们。\par

\SourceRecord{Translated by R.G. MacMullen.；Nicene and Post-Nicene Fathers, First Series；Vol. 6.；Edited by Philip Schaff.；Buffalo, NY: Christian Literature Publishing Co.,；1888.；<http://www.newadvent.org/fathers/160322.htm>.}

% structure: p0001 source=h1 target=section reason=page-title

\section{\WorkTitle{论新约 第二十三篇讲道}}

[伦.第七十三篇]\par

\Emph{论福音之言}，\BibleRef{玛 13:19} \Emph{等，其中主耶稣解释撒种的比喻。}\par

1. 昨天和今天，你们都听到了我们主耶稣基督关于撒种者的比喻。昨天在场的人，今天请回忆一下。昨天我们读到那个撒种者，当他撒种时，\Quote{有的落在路旁}，被飞鸟啄食；\Quote{有的落在石头地上}，因日晒而枯干；\Quote{有的落在荆棘中，被窒息了}，不能结实；\Quote{有的落在好地里，就结了果实，有一百倍的、六十倍的、三十倍的。}但今天主又讲了另一个关于撒种者的比喻，\Quote{他撒好种子在田里。人们睡觉时，仇敌来了，在麦子中间撒上莠子。}只要还在苗期，就看不出来；但当好种子的果实开始显现时，\Quote{莠子也显出来了。}家主仆人们看见好麦子中长了许多莠子，便恼怒，想要拔掉它们，却不被允许；而是对他们说：\Quote{让两者一起生长，直到收割的时候。}现在，主耶稣基督也解释了这比喻；祂说祂是撒好种子的人，并指出撒莠子的仇敌是魔鬼；收割的时候是世界的终结；祂的田地是全世界。祂还说：\Quote{在收割的时候，我要对收割的人说：你们先把莠子收集起来，捆成捆去烧，但把麦子收进我的谷仓。}祂说：你们这些充满热忱的仆人，为什么这样急躁？你们在麦子中看见莠子，在好人中看见坏基督徒，就想拔掉坏人；安静吧，现在还不是收割的时候。那时刻会到来，但愿那时你们是麦子！你们为什么自寻烦恼？为什么不能容忍恶人与好人混在一起？在田野里，他们可能和你们在一起，但在谷仓里就不会了。\par

2. 现在你们知道，昨天提到的种子不生长的三种地方，\Quote{路旁}、\Quote{石头地}和\Quote{荆棘地}，与这些\Quote{莠子}是一样的。它们只是在不同比喻下取了不同的名字。因为当使用比喻时，术语的字面意义并未表达，而是通过它们传达真理的相似性。我注意到只有少数人理解了我的意思；但我说话是为了所有人的益处。在有形事物中，路旁就是路旁，石头地就是石头地，荆棘地就是荆棘地；它们就是它们自己，因为名称是按字面意义使用的。但在比喻和象征中，同一事物可以用许多名称来称呼；因此，我告诉你们，那\Quote{路旁}、那\Quote{石头地}、那些\Quote{荆棘地}是坏基督徒，他们也是\Quote{莠子}，这并不可矛盾。基督不是被称为\Quote{羔羊}吗？基督不也是\Quote{狮子}吗？在野兽和牲畜中，羔羊就是羔羊，狮子就是狮子；但基督两者都是。前者按字面表达各自是什么；后者在象征意义上一起是。不仅如此，在象征下，非常不同的事物也可以用同一个名称来称呼。有什么比基督和魔鬼更不同呢？但基督和魔鬼都被称为\Quote{狮子}。基督被称为\Quote{狮子}：\Quote{犹大支派的狮子已得胜}；魔鬼也被称为狮子：\Quote{你们不知道你们的仇敌魔鬼如同吼叫的狮子，巡游寻找可吞吃的人吗？}那么，两者都是狮子；一个是因其力量而称狮子，另一个是因凶残；一个是为\Quote{得胜}的狮子，另一个是为伤害的狮子。魔鬼又是蛇，\Quote{那古蛇}；那么，当我们的牧者告诉我们\Quote{你们要像蛇一样机警，像鸽子一样纯朴}时，我们是被命令去效法魔鬼吗？\par

3. 因此，昨天我向\Quote{路旁}、向\Quote{石头地}、向\Quote{荆棘地}讲话，我说：趁你们还能改变时改变吧！用犁翻开硬地，把石头从田里拣出，拔掉其中的荆棘。不要保留那刚硬的心，免得天主的话很快过去而失落。不要有那样浅薄的土壤，使爱德的根不能扎深。不要以世俗的贪欲和忧虑窒息我劳苦撒在你们里面的好种子。因为播种的是主，我们只是祂的工人。但你们要成为\Quote{好地}。我昨天说过，今天也对所有人再说一次：让一个人结出\Quote{一百倍，另一个六十倍，另一个三十倍}。有人果实多，有人少；但所有人都会在谷仓里有一席之地。昨天我说了这一切，今天我要对莠子讲话；但羊群自己就是莠子。哦，坏基督徒，哦，你们这些只通过恶行压迫教会的人，在收割之前改过自新吧！\Quote{不要说：我犯了罪，又有什么事临到我呢？}天主并未失去大能；祂正要求你们悔改。我向那些虽然是基督徒的恶人说话；我向莠子说话。因为他们是在田里；有可能今天他们是莠子，明天就变成麦子。所以我也要向麦子说话。\par

4. 你们这些生活良好的基督徒，你们叹息呻吟，因为你们是少数，在众多人中是少数，在非常众多的人中是少数。冬天将过去，夏天将来到；看！收割很快就要到了。天使们将要来，他们能进行分离，并且不会犯错误。我们现在这个时候，就像那些仆人们，关于他们曾说过：\BibleQuote{你愿意我们去把莠子收集起来吗？}因为我们曾经希望，如果可能的话，没有一个恶人留在善人中间。但有人告诉了我们：\BibleQuote{让两者一起生长直到收割。}为什么？因为你们这样的人可能会受骗。最后请听：\BibleQuote{免得你们收集莠子时，连麦子也一起拔出来。}你们在做什么好事？你们的热忱会使我的收成浪费吗？收割者们将要来到，而收割者是谁，祂已经解释了：\BibleQuote{收割者是天使。}我们不过是人，收割者是天使。我们确实，如果我们完成我们的路程，也将成为与天主的天使同等；但现在当我们对恶人感到恼怒时，我们仍然不过是人。我们现在应该听从这些话：\BibleQuote{所以，那自以为站得稳的，要小心，免得跌倒。}因为你们以为，我的弟兄们，我们所读的这些莠子不会登上这个座位吗？你们以为它们都在下面，而上面这里没有吗？愿天主赐恩，使我们不至如此。\BibleQuote{至于我被你们审断，或被人审判，这为我是一件极小的事。}我真实地告诉你们，我亲爱的诸位，即使在这些高高的座位上，也有麦子和莠子，而在平信徒中，也有麦子和莠子。让善人容忍恶人；让恶人改变自己，效法善人。让我们所有人，如果可能的话，到达天主那里；让我们所有人通过祂的仁慈，逃离这个世界的邪恶。让我们寻求好日子，因为我们如今在坏日子中；但在坏日子中，让我们不要亵渎，这样我们才能到达好日子。\par

\SourceRecord{Translated by R.G. MacMullen.；Nicene and Post-Nicene Fathers, First Series；Vol. 6.；Edited by Philip Schaff.；Buffalo, NY: Christian Literature Publishing Co.,；1888.；<http://www.newadvent.org/fathers/160323.htm>.}

% structure: p0001 source=h1 target=section reason=page-title

\section{论新约讲道 第廿四篇}

\Emph{[LXXIV. Ben.]}\par

\Emph{论福音之言}，\BibleRef{玛 13:52} \Emph{，\BibleQuote{因此，凡成为天国门徒的经师，}等等。}\par

1. 福音读经提醒我去找出并解释给你们，亲爱的，正如主赐给我能力，那位受教于天国的经师是谁，他是\BibleQuote{如同一个家主，从他的宝库里拿出新旧的东西}。因为福音读经在这里结束了。\Quote{受教的经师的新旧东西是什么？}现在众所周知他们是谁，就是古人按照我们圣经的习惯称为经师的人，即那些声称精通法律的人。因为在犹太人中，这样的人被称为经师，不像现在在法官的职务或国家的习惯中如此称呼的那样。因为我们进学校不能徒劳无益，但我们必须知道在什么意义上理解圣经的词语，以免当提及圣经中的某个词时，若该词在世俗用语中另有通常含义，听众就会误解，并因想到其惯常意义，而无法理解所听到的内容。当时的经师就是那些声称精通法律的人，保存、研习、抄写和注释法律书都属于他们的职责。\par

2. 这些人就是我们的主耶稣基督所斥责的，因为他们有天国的钥匙，却\BibleQuote{自己不进去，也不让别人进去}；用这些话指责法利塞人和经师，即犹太法律的教师。关于他们，主在另一处说：\BibleQuote{凡他们所说的，你们要遵行，但不要效法他们的行为，因为他们说而不做。}为什么对你们说：\BibleQuote{因为他们说而不做？}岂不是因为有一些人，宗徒所说的在他们身上清楚显明：\BibleQuote{你教导人不可偷盗，你自己还偷盗吗？你教导人不可奸淫，你自己还奸淫吗？你憎恶偶像，你自己还行亵渎之事吗？你以法律自夸，却因违犯法律而羞辱天主？因为天主的名因你们在外邦人中被亵渎。}显然主所说的就是指这些人：\BibleQuote{因为他们说而不做。}那么他们是经师，但不是\BibleQuote{受教于天国的}。\par

3. 或许你们中有人会说：\Quote{一个恶人怎能说好话呢？因为经上记载主自己的话说：\InnerQuote{善人从心里善的宝库发出善事，恶人从恶的宝库发出恶事。你们这些伪君子，你们既是恶人，怎能说出好话来呢？}}主在一处说：\BibleQuote{你们既是恶人，怎能说出好话来呢？}在另一处说：\BibleQuote{凡他们所说的，你们要遵行，但不要效法他们的行为，因为他们说而不做。}如果他们\BibleQuote{说而不做}，他们就是恶人；如果他们是恶人，他们就不能\BibleQuote{说好话}；那么，我们怎能去做我们从他们那里听到的事，既然我们不能从他们那里听到好事？现在请注意，圣洁而亲爱的，这个问题如何解决。恶人从自己发出的，都是恶的；恶人从自己心里发出的，都是恶的；因为那里有恶的宝库。但善人从自己心里发出的，都是善的；因为那里有善的宝库。那么，那些恶人从哪里发出善事呢？\BibleQuote{因为他们坐在梅瑟的座位上。}如果主没有首先说\BibleQuote{他们坐在梅瑟的座位上}，他就绝不会命令应当听从恶人。因为他们从自己心里恶的宝库里发出的是一回事；他们从梅瑟的座位上发出的，即作为法官的传令官，是另一回事。传令官所说的话，如果他在法官面前说话，永远不会归到他身上。传令官在自己家里说的是一回事，传令官从法官那里听到而说出的又是另一回事。因为不管他愿意与否，传令官必须宣告对他自己朋友的刑罚判决。同样，不管他愿意与否，他必须宣告对他自己敌人的无罪判决。假设他凭自己的心意说话：他宣告朋友无罪，惩罚敌人。假设他坐在法官的席位上说话：他惩罚朋友，宣告敌人无罪。那么，经师也是如此；假设他们凭自己的心意说话：你会听到\BibleQuote{我们吃喝吧，因为明天我们要死了。}假设他们从梅瑟的座位上说话：你会听到\BibleQuote{不可杀人，不可奸淫，不可偷盗，不可作假见证。应孝敬你的父亲和母亲；应爱人如己。}那么，你们要履行经师之口所宣告的职位之言，而不是他们心里发出的。因为这样，你既遵行了主的这两项判断，就不会在一个方面顺从，在另一方面违背；而会明白两者是一致的，并且会认为这两点是真实的：即\BibleQuote{善人从心里善的宝库发出善事，恶人从恶的宝库发出恶事}；以及另一点，即那些经师不是从自己心里恶的宝库发出善事，而是能够从梅瑟座位的宝库发出善事。\par

4. 那么，当主说\Quote{每棵树凭它自己的果实就能被认出。人从荆棘上岂能摘到葡萄？从蒺藜上岂能摘到无花果？}时，这些话不会使你困扰。犹太人的经师和法利塞人原是荆棘和蒺藜，然而他却说：\Quote{他们所说的，你们要遵行，但不要效法他们的行为。}因此，从荆棘上能摘到葡萄，从蒺藜上能摘到无花果，正如他刚刚讲述的方法所启示你的那样。因为有时在葡萄园的荆棘篱笆中，葡萄藤会缠绕其中，一串串葡萄会从荆条上垂挂下来。你刚一听到荆棘的名称，就几乎要忽视葡萄。但你要寻找荆棘的根，就会看到它在哪里。也要寻找垂挂葡萄串的根，就会看到它在哪里。因此，你们要明白：一个指的是法利塞人的心，另一个指的是梅瑟的座位。\par

5. 但他们为何如此？\Quote{因为}，圣保禄说，\Quote{帕子还在他们心上。他们看不到旧事已经过去，一切都成了新的。}所以他们如此，所有现在与他们一样的人也如此。为何是旧事？因为它们早已被公布。为何是新事？因为它们关乎天主国。那么，帕子如何被除去？宗徒亲自告诉我们：\Quote{但你们何时转向主，帕子就会被除去。}因此，那不转向主的犹太人心思不朝向终向。正如那时以色列子民在这预象中不将目光凝视\Quote{终向}，即梅瑟的面容。因为梅瑟发光的容貌含有真理的预象；帕子放在中间，因为以色列子民还不能看见他面容的光荣。\Quote{这个预象已被废除。}因为宗徒这样说：\Quote{它已被废除。}为何废除？因为当皇帝来到时，他的肖像就被除去。当皇帝不在时，人们观看肖像；但肖像所代表的皇帝在哪里，肖像就被除去。因此，在我们的皇帝主耶稣基督来到之前，有肖像在他面前行进。当肖像被除去时，皇帝亲临的光荣就显现了。因此，\Quote{任何人转向主，帕子就被除去。}因为梅瑟的声音透过帕子传出，但梅瑟的面容却没有被看见。如今，基督的声音通过旧经书向犹太人发出：他们听见声音，却看不见说话者的面容。那么，他们愿意帕子被除去吗？\Quote{让他们转向主。}因为那时旧事不是被除去，而是被存放在宝库里，使经师从此成为\Quote{在天主国受教的经师，从他的宝库里拿出}不只是\Quote{新事}，也不只是\Quote{旧事}。因为如果他只拿出\Quote{新事}或只拿出\Quote{旧事}，他就不是\Quote{在天主国受教的经师，从他的宝库里拿出新和旧的东西}。如果他说而不做，那么他是从教导的座位上拿出，而不是从他心中的宝库拿出。并且（我们说实话，圣洁的弟兄们），凡从旧事中拿出的，都由新事所光照。因此，\Quote{我们要转向主，好使帕子被除去。}\par

\SourceRecord{Translated by R.G. MacMullen.；Nicene and Post-Nicene Fathers, First Series；Vol. 6.；Edited by Philip Schaff.；Buffalo, NY: Christian Literature Publishing Co.,；1888.；<http://www.newadvent.org/fathers/160324.htm>.}

% structure: p0001 source=h1 target=section reason=page-title

\section{新约讲道集 第25篇}

\Emph{[LXXV. Ben.]}\par

\Emph{论福音经文}，\BibleRef{玛 14:24}，\Emph{\BibleQuote{那时船已在海中，被波浪颠簸。}}\par

1. 我们刚才听到的福音教训，是关于谦逊的教训，为叫我们看见并知道我们在哪里，以及我们应该朝向和奔赴何方。因为那条载着门徒、被逆风抛掷于波浪中的船，并非没有含义。主离开群众，独自上山祈祷，然后来到门徒那里，发现他们处于危险之中，行走在海面上，上船后坚固他们，平息风浪，这也并非没有含义。但若祂能平息万事，又有何奇？祂创造了万物。然而，祂上船后，船上的人前来说道：\Quote{你确实是天主子。}但在祂清楚显现自己之前，他们惊慌失措，说道：\Quote{是个鬼怪。}祂上船后，从他们心中消除了内心的摇摆，因为那时他们由于怀疑而在灵魂中比先前在身体上因风浪而更危险。\par

2. 然而，在主所做的这一切中，祂教导我们关于我们在此世生活的本质。在这世界上，没有一个人不是旅客，尽管并非所有人都渴望返回自己的家乡。正是由于这旅途，我们遭受风浪和风暴；但我们必须至少待在船上。因为即使船上有危险，但离开船则必定灭亡。因为无论一个人在开阔的海中游泳的臂力有多强，最终仍会被巨浪冲走而沉没，被波浪的威力征服。因此，我们必须待在船里，即我们必须被木头承载，才能渡过这海。承载我们软弱的这木头就是主的十字架，我们因它被标记，并从这世界的危险风暴中得救。我们遭受波浪的冲击，但帮助我们的是天主。\par

3. 因为当主离开群众后，\Quote{他独自上山祈祷}；那山象征天的高度。因为离开群众后，主复活后独自升天，并且\Quote{在那里}，如宗徒所说，\Quote{他为我们转求。}那么，他\Quote{离开群众，独自上山祈祷}是有含义的。因为他独自一人，从死者中复活的身体之后，成为首生者，坐在天父的右边，是我们祈祷的大司祭和辩护者。教会的元首在上，叫其余肢体在末后跟随。既然如此，他\Quote{为我们转求}，超越所有受造物的高度，如同在山顶上，\Quote{他独自祈祷。}\par

4. 与此同时，载着门徒的船，即教会，被诱惑的风暴抛掷和摇撼；逆风，即她的仇敌魔鬼，不停息，竭力阻止她到达安息。但\Quote{为我们转求的那位}更伟大。因为在这使我们劳苦的颠簸中，他赐给我们信心，来到我们中间并坚固我们；只是我们不要在困境中将自己抛出船外，投进海里。因为即使船在困境中，它仍是船。只有她载着门徒，并接纳基督。海上确实有危险，但离开她则立刻丧亡。因此，你当待在船里，向天主祈祷。因为当所有计谋失败，甚至连舵也失灵，张帆反而危险而无用时，当所有人力与力量都消失时，留给水手的只有恳切的祈求呼声，向天主倾心祈祷。那么，那位赐予水手到达港口的，岂会舍弃自己的教会，不带她到达安息吗？\par

5. 然而，弟兄们，这极大的困境在这船中，仅当主不在时才会发生。什么！在教会里的人，能使主远离他吗？他何时使主远离他？当他被某种情欲战胜时。因为在某处经文中以比喻说：\Quote{不可让太阳落在你们的愤怒上；也不可给魔鬼留地步。}这并不指那可见的太阳，它在可见受造物中享有荣耀的顶峰，我们与野兽都能同样看到；而是指那只有内心纯洁的信徒才能看到的光，如经上所记：\Quote{那是真光，照亮每个来到世界的人。}因为可见太阳的光也照亮最微小的动物。义德和智慧就是那真光，当人被愤怒的混乱如同云遮蔽时，心灵就停止看见它；然后，就好像太阳落在人的愤怒上。同样，在这船中，当基督不在时，每个人都被自己的风暴、不义和邪恶欲望摇撼。例如，法律告诉你：\Quote{不可作假见证。}如果你遵守见证的真理，你就在灵魂中有光；但如果你被不洁之财的欲望战胜，决意要说假见证，你将立刻因基督的缺席开始被风暴扰乱，被你贪婪的波浪抛来抛去，被你情欲的暴风雨危及，并且因基督的缺席仿佛即将沉没。\par

6. 有何畏惧之由，恐船偏离航向，退向后方？此乃发生，当人放弃天上赏报之望，情欲把持舵柄，转向可见易逝之物！凡为情欲诱惑所扰，却仍反观内心者，其处境尚非完全绝望，若能恳求宽恕己过，克己奋发，战胜怒海汹涛。然而，有人背离本来之我，心中自语：\Quote{天主不见我；因祂不思我，也不关心我是否犯罪。}此人已转舵，为风暴所携，返归原处。因人心多念；基督不在时，船便为世俗波浪与诸般风暴所颠簸。\par

7. 夜间的第四更，即是夜之终末；因每一更含三小时。此象征：如今在世界终末，主已来助，被见行于海上。虽然此船为诱惑风暴所颠簸，但她看见祂荣耀的天主行走于一切海涛之上，即此世一切掌权者之上。因先前，根据祂受难之时，按肉体祂展示谦卑的榜样，曾以配合时势的话说过：海涛徒然向祂发怒，祂为我们的缘故自愿屈服，为应验那预言：\BibleQuote{我来到深海中，洪流淹没了我。}祂并未斥退假证人，也未斥退那些喊叫\BibleQuote{钉死祂！}之人的野蛮呼声。祂并未以大能压制那些狂怒之徒的野蛮心灵和言语，却以忍耐承受了一切。他们任意待祂，因为祂\BibleQuote{服从至死，且死在十字架上。}但当祂从死中复活后，好独自为祂的门徒祈祷——他们被安置在教会中，如在一艘船上，并凭信德承载于祂的十字架之木上，在此世诱惑中航行，如在海浪中——祂的名号开始在这曾轻蔑、控告、杀害祂的世上受到尊敬；如此，祂在肉体受苦安排中\BibleQuote{曾来到深海中，洪流曾经淹没祂}的那位，如今借着祂名号的荣耀，践踏骄傲者的颈项，如踏汹涌波涛。正如我们现在所见的，主仿佛行于海上，我们看到此世的一切疯狂尽伏于祂脚下。\par

8. 除了风暴之险，还有异端者的谬误；且不乏有人这样试炼船上之人的心思，竟说基督非为童贞女所生，亦无真实身体，只是眼睛看来如此，实则不然。这些异端者的意见现在兴起，正当基督的名号已受万民荣耀之时；也就是说，基督如今仿佛行于海上。门徒在考验中说：\BibleQuote{是个鬼怪。}但祂以自己的声音赐予我们力量抵御这些有害的意见：\BibleQuote{放心！是我，不要怕！}因为人们基于虚妄的恐惧而对基督产生这些意见，瞻仰祂的尊荣与威严，以为祂既堪受如此荣耀，就不能如此降生，从而畏惧祂\Quote{在海面上行走。}因这行动象征祂尊荣的卓越；所以他们以为祂是个幻影。但祂说：\BibleQuote{是我}——这岂非说祂身上无一不是真实？因此，若祂显示祂的肉体，就是肉体；若显示骨，就是骨；若显示伤痕，就是伤痕。因为\BibleQuote{在祂并没有\InnerQuote{是}而又\InnerQuote{非}，在祂只有\InnerQuote{是}}，如宗徒所言。故此祂说：\BibleQuote{放心！是我，不要怕！}即是说：不要因畏惧我的威严，就想要剥夺我存在的真实性。虽然我行于海上，虽然我将此世的傲慢与骄傲如怒涛踩在脚下，但我仍以真人显现，我的福音确实传报关于我的真理：我为童贞女所生，我是降生成人的圣言；我真实地说过：\BibleQuote{你们摸摸我，看看！因为神魂没有骨肉，如同你们看见我有的。}怀疑的宗徒亲手抚摸的，是我伤痕的真迹。因此，\BibleQuote{是我；不要怕。}\par

9. 但门徒们以为祂是幻影，这不仅指那些否认主具有人性肉身的人，也不仅指那些有时以盲目的乖戾甚至动摇船上的人；也指那些认为主在任何事上说了谎，并且不信祂对恶人所威胁的事必将实现的人。仿佛祂部分真实、部分虚假，在言语上如幻影出现，如同祂是\Quote{是与否}。但那些正确理解祂声音的人，祂说：\BibleQuote{是我，不要害怕}；他们立刻相信主的一切言语，如此，他们既盼望祂所应许的赏报，也畏惧祂所警告的惩罚。因为祂对右边的人所说的\BibleQuote{我父所祝福的，你们来吧！承受自创世以来为你们预备的国度}是真实的；同样，左边的人所听见的\BibleQuote{你们到那为魔鬼及其使者预备的永火里去吧}也是真实的。因为这种见解——即认为基督对不义和堕落之人的威胁不真实——源于他们看到无数民族和群众臣服于祂的名；由此，基督在他们看来如同幻影，因为祂行走在海面上；也就是说，祂在惩罚的威胁上似乎说谎，因为祂似乎不能消灭这些臣服于祂名号和尊荣的无数人民。但让他们听祂说：\BibleQuote{是我}；因此，那些相信基督在一切事上真实的人，不仅寻求祂所应许的，也避免祂所警告的，他们就\BibleQuote{不要害怕}；因为虽然祂行走在海面上，即世上万民都臣服于祂，但祂并非幻影，因此祂说话不虚假，当祂说：\BibleQuote{不是凡向我说\InnerQuote{主啊，主啊！}的人，就能进入天国}。\par

10. 那么，伯多禄敢于在水面上走向祂，又预示什么呢？因为伯多禄通常代表教会的形象。那么，\BibleQuote{主，如果是祢，就叫我在水面上走到祢那里去}又意味着什么，难道不是：主，若祢是真实的，且凡事不说谎，就让祢的教会也在世上得荣耀，因为预言已论及祢说：\BibleQuote{人民的富足人必恳求祢的面}。但因世人的赞美对主不是诱惑，而人在教会中常因人的赞美和尊荣而混乱，几乎被它们淹没；因此伯多禄在海中颤抖，因暴风的威力而恐惧。因为谁不害怕那些话：\BibleQuote{那称你们有福的人，正是迷惑你们，扰乱你们脚下的路}？既然灵魂与渴求人类赞美的欲望搏斗，那么在此危险中，投身祈祷和恳切哀求是好的：免得那被赞美迷惑的人，被责难淹没而沉没。让伯多禄在水中动摇时呼喊说：\BibleQuote{主，救我}。因为主要伸出祂的手，虽责备说：\BibleQuote{小信德的人哪，为什么怀疑？}——为什么你不直视你所要走向的那一位，只以主为夸耀？但祂必将他从波浪中救出，不许那承认自己软弱、恳求祂帮助的人丧亡。当他们将主接入船中，信德得以坚固，一切疑虑消除，海上的风暴平息，以至于到达稳固平安的岸边时，他们都朝拜祂说：\BibleQuote{祢真是天主子}。因为这就是那永恒的喜乐：真理显明，天主的圣言，创造万有的智慧，以及祂慈爱的崇高都得以认知和爱慕。\par

\SourceRecord{Translated by R.G. MacMullen.；Nicene and Post-Nicene Fathers, First Series；Vol. 6.；Edited by Philip Schaff.；Buffalo, NY: Christian Literature Publishing Co.,；1888.；<http://www.newadvent.org/fathers/160325.htm>.}

% structure: p0001 source=h1 target=section reason=page-title

\section{论新约讲道集 第26篇}

\Emph{[LXXVI. Ben.]}\par

\Emph{再论 \BibleRef{玛 14:25}：主行走在海浪之上，以及伯多禄的摇晃。}\par

1. 刚才所读的福音书论及主基督行走在海面上；宗徒伯多禄行走时因恐惧而摇晃，因缺乏信德而沉没，又因宣认而重新站起，这使我们明白：大海象征现世，而宗徒伯多禄代表唯一教会。因为在宗徒中，伯多禄居于首位，在爱基督方面最为热切，常常代表众人回答。再者，当主耶稣基督问门徒们人们说祂是谁，门徒们报告了众人的各种意见后，主又问他们说：\Quote{但你们说我是谁？} 伯多禄回答说：\Quote{祢是基督，永生天主之子。} 一人为众人回答，合一在多人中。于是主对他说：\Quote{你是有福的，西满巴约纳，因为不是血肉启示了你，而是我在天之父。} 然后祂加说：\Quote{我也对你说。} 仿佛祂在说：\Quote{因为你对我说：\InnerQuote{祢是基督，永生天主之子。} 所以我也对你说：\InnerQuote{你是伯多禄。}} 因为他先前被称为西满。这\Quote{伯多禄}的名字是主赐给他的，寓意深远，应象征教会。因为基督是\OriginalTerm{磐石}{Petra}，伯多禄就是基督徒子民。因为磐石（Petra）是原初之名。因此，伯多禄由磐石得名，而非磐石由伯多禄得名；正如基督不是由基督徒而被称为基督，而是基督徒由基督而得名。祂说：\Quote{因此，你是伯多禄；在这磐石上——就是你所宣认的磐石，你所承认的磐石，说\InnerQuote{祢是基督，永生天主之子。} ——我要建立我的教会。} 即在我自己，永生天主之子身上，\Quote{我要建立我的教会。} 我要把你建立在我身上，而非把我建立在你身上。\par

2. 因为那些愿意建立在人身上的人说：\Quote{我是属保禄的；我是属阿颇罗的；我是属\OriginalTerm{刻法}{Cephas}的}，刻法就是伯多禄。但那些不愿意建立在伯多禄身上，而愿意建立在磐石上的人说：\Quote{我是属基督的。} 当宗徒保禄清楚自己是被拣选的，而基督被轻视时，他说：\Quote{基督被分裂了吗？保禄为你们被钉在十字架上了吗？或者你们受洗是归于保禄的名下吗？} 既然不是归于保禄的名下，也不是归于伯多禄的名下，而是归于基督的名下：这样伯多禄就建立在磐石上，而非磐石建立在伯多禄上。\par

3. 正是这位被磐石称为\Quote{有福的}伯多禄，他承载着教会的形象，在宗徒职位中居首位，但在他听到自己被称为\Quote{有福的}之后不久，听到自己被称为\Quote{伯多禄}之后不久，听到自己将被\Quote{建立在磐石上}之后不久，当他听到主预言即将来临的苦难时，却使主不悦。因为主曾预告门徒们说苦难即将来临。他害怕会因死亡而失去他所宣认为生命源泉的那一位。他慌乱地说：\Quote{主，这事决不可临到祢身上。愿祢怜悯自己，天主，我不愿祢死。} 伯多禄对基督说：我不愿祢死；但基督远更美好地说：我愿为你而死。然后祂立刻责备了那位不久之前还称赞的人；称他为\Quote{撒殚}（\OriginalTerm{撒殚}{Satan}），而先前曾称他为\Quote{有福的}。\Quote{撒殚，退到我后面去！} 祂说：\Quote{你是我的绊脚石，因为你所体会的，不是天主的事，而是人的事。} 既然祂责备我们只因我们是人，祂要我们在这现世做什么呢？你想知道祂要我们做什么吗？请听圣咏：\Quote{我曾说：你们是神，你们都是至高者的儿子。} 但由于体会人的事，\Quote{你们要像世人一样死亡。} 正是同一位伯多禄，不久之前是\Quote{有福的}，之后却成为\Quote{撒殚}，转瞬之间，短短几句话之内！你惊讶于称呼的差异，就请注意理由的差异。为什么他不久之前是\Quote{有福的}，之后却成为\Quote{撒殚}？请注意他被称为\Quote{有福的}的理由。\Quote{因为不是血肉启示了你，而是我在天之父。} 因此，他是有福的，因为不是血肉启示了你。因为如果是血肉启示了你，那是属于你自己的；但因为不是血肉启示了你，而是我在天之父，那是属于我的，不是属于你自己的。为什么是属于我的？\Quote{因为父所有的一切都是我的。} 所以，你听到了为什么他是\Quote{有福的}，为什么他是\Quote{伯多禄}的理由。但为什么他成了我们战栗而不愿重复的那位？为什么？是因为出于你自己的。\Quote{因为你所体会的，不是天主的事，而是人的事。}\par

4. 让我们在这教会的肢体中审视自己，分辨什么是属于天主的，什么是属于我们自己的。因为这样我们就不会摇晃，就能建立在磐石上，稳固坚定，抵挡风、洪水、激流——即现世的诱惑。然而请看这位伯多禄，他曾是我们的预像：时而信赖，时而摇晃；时而宣认不死的主，时而害怕祂会死。为什么？因为基督的教会既有强壮者，也有软弱者；不能没有强壮者或软弱者；因此宗徒保禄说：\Quote{我们强壮者应担待软弱者的脆弱。} 当伯多禄说：\Quote{祢是基督，永生天主之子} 时，他代表强壮者；但当他摇晃，不愿基督受苦，因害怕为基督而死，而不承认生命时，他代表教会的软弱者。因此，在那一位宗徒，即伯多禄，宗徒中为首和居首位的、教会预表在他身上的那位，两种类型都必须被代表，即强壮者和软弱者；因为教会缺少两者都不能存在。\par

5. 因此，刚才读到的也是这个道理：\Quote{主，如果是祢，请吩咐我走在水面上到祢那里去。}因为这事我不能凭自己，只能凭祢。他承认自己有什么，以及从祂那里有什么，他相信凭着祂的旨意，他能做到人性软弱所不能做的事。因此，\Quote{如果是祢，请吩咐我}；因为当祢吩咐时，这事就必成就。我所不能凭自己承担的事，祢能藉着吩咐我而成就。主说：\Quote{来吧。}伯多禄毫不迟疑地，因着那吩咐他者的言语、那扶持他者的临在、那引导他者的临在，立刻跳进水里，开始行走。他能做主所做的事，不是在自己内，而是在主内。\Quote{因为你们从前曾是黑暗，但现在在主内却是光明。}没有人在保禄内、没有人在伯多禄内、也没有人在任何其他宗徒内所能做的事，可以在主内做到。因此，保禄藉着对自己有益的自谦并赞美祂，说得很好：\Quote{保禄为你们被钉死了吗？或者你们受洗是归于保禄名下吗？}所以，你们不是\Emph{在}我内，而是\Emph{与}我一起；不是在我之下，而是在祂之下。\par

6. 因此，伯多禄遵照主的吩咐在水面上行走，知道他自己不能有这能力。因着信德，他获得了力量去做人性软弱不能做的事。这就是教会中的强者。你们要留意、聆听、明白并照样去做。因为我们对强者必须以这样的原则对待，好使他们成为弱者；但对弱者我们却要以这样的原则对待，好使他们成为强者。但自恃己力使许多人远离了力量。除非人感到自己软弱，否则不会从天主得到力量。\Quote{天主为祂的产业降下了甘霖。}你们这些知道我要说什么的人，为什么抢在我前面呢？愿你们的敏捷有所收敛，好让其余的人能跟上。我说过，并且再说一次：你们要聆听、接受并按此原则行事。除非人感到自己软弱，否则不会被天主造就成强者。因此，正如圣咏所说，\Quote{甘霖}——\Quote{甘霖}不是出于我们的功劳，而是\Quote{甘霖}。所以，\Quote{天主为祂的产业降下了甘霖}；因为\Quote{它原是软弱的，但祢使它坚强了。}因为祢\Quote{为它预备了甘霖}，不看重人的功劳，而看重祢自己的恩宠和怜悯。这产业于是被削弱，并在自身内承认自己的软弱，好能在祢内坚强。如果它不软弱，就不会被坚强，以便藉着祢而在祢内得以\Quote{成全}。\par

7. 请看保禄，他是这产业的一小部分，看他在软弱中说：\Quote{我不配称为宗徒，因为我曾迫害过天主的教会。}那么你为何成了宗徒呢？\Quote{因天主的恩宠，我成了今日的我。我不配，但因天主的恩宠，我成了今日的我。}保禄原是\Quote{软弱的}，但祢\Quote{成全}了他。但现在因\Quote{天主的恩宠，他成了今日的他}，请看下文：\Quote{祂的恩宠在我身上没有落空，反而我比他们众人更劳苦。}你要小心，免得因骄矜而失去那曾因软弱而获得的一切。这是好的，非常好：\Quote{我不配称为宗徒。因祂的恩宠，我成了今日的我，祂的恩宠在我身上没有落空。}这一切都极好。但\Quote{我比他们众人更劳苦}——你似乎开始把方才归给天主的据为己有。请留意并继续看：\Quote{其实不是我，而是天主的恩宠偕同我。}好！你这软弱的人，你将在极大的力量中被高举，因为你没有忘恩。你就是这同一位保禄，在自己内微小，在主内伟大。你曾三次求主，使\Quote{那加在你身上的刺，即撒殚的使者，那击打你的，愿它离开你。}当你这样祈求时，对你说的是什么？你听到什么？\Quote{我的恩宠为你够了，因为我的德能在软弱中才全显出来。}因为他原是\Quote{软弱的}，但祢\Quote{成全}了他。\par

8. 同样，伯多禄也说：\Quote{请吩咐我走在水面上到祢那里去。}我敢这样做的不过是一个人，但我恳求的不是人。愿那神人命令，使人能做人所不能做的事。祂说：\Quote{来吧。}他便下去，开始在水面上行走；伯多禄能够做到，因为磐石吩咐了他。看啊，伯多禄在主内是什么？他自己内是什么？\Quote{他一见风势猛烈，就害怕；并开始下沉，便喊说：\InnerQuote{主，我丧命了，救救我！}}当他从主那里寻求力量时，他便从主那里得了力量；作为人，他摇动，但他回到主那里。\Quote{我若说我的脚滑了}（这是圣咏的话，是圣歌的曲调；我们若承认，它们也是我们的话；是的，我们若愿意，它们也是我们的）。\Quote{我若说我的脚滑了}——怎么滑了？除非是因为那是我的。下文是什么？\Quote{主，是祢的仁慈扶持了我。}不是我的力量，而是祢的仁慈。因为天主岂会遗弃那摇动、那呼求祂的人呢？那么，\Quote{谁呼求过天主，而被祂遗弃了呢？}又在哪里说：\Quote{凡呼号主名字的人，必获救}？祂立刻伸出右手的援助，把下沉的他扶起来，并责备他的不信：\Quote{小信德的人哪，你为什么怀疑？}你曾信靠我，现在对我起了疑心吗？\par

9. 诸位弟兄，我的讲道必须结束了。当将世界视为大海；风浪汹涌，狂风大作。各人独特的情欲便是他的风暴。你爱天主，便能行走于海上，脚下是世界的汹涌。你若爱世界，它必将你吞没。它只懂得吞吃爱慕它的人，而非承载他们。但当你的心被情欲颠簸时，为使你可战胜自己的情欲，当呼求基督的天主性。你们以为逆风便是今生的逆境吗？因为当有战争、骚乱、饥荒、瘟疫，甚至当个人遭遇私有的祸患时，人便以为这是逆风，以为必须呼求天主。但当世界以暂时的幸福微笑时，仿佛就没有逆风了。然而，不要以时世的平静来询问此事：只问你自己的情欲。看看你内里是否平静：看看有没有内在的风将你倾覆；要注意这个。要有极大的德行才能与幸福搏斗，免得这幸福诱惑你、败坏你、推翻你。我说，要有极大的德行才能与幸福搏斗，并要有极大的幸福才能不被幸福征服。那么，学习践踏世界吧；记住要信靠基督。\Quote{若你的脚滑了；}如果你摇摇欲坠，如果有某些你无法克服的事，如果你开始下沉，就说：\Quote{主啊，我丧亡了，救我。}要说：\Quote{我丧亡了，}免得你真正丧亡。因为只有那在身体内为你而死的那一位，才能救你脱离身体的死亡。让我们转向主，等等。\par

\SourceRecord{Translated by R.G. MacMullen.；Nicene and Post-Nicene Fathers, First Series；Vol. 6.；Edited by Philip Schaff.；Buffalo, NY: Christian Literature Publishing Co.,；1888.；<http://www.newadvent.org/fathers/160326.htm>.}

% structure: p0001 source=h1 target=section reason=page-title

\section{论新约讲道 27}

\Emph{[LXXVII. Ben.]}\par

\Emph{论福音之言，\BibleRef{玛 15:21}，\Quote{耶稣从那里出去，退到提洛和漆冬一带。看，有一个客纳罕妇人，}等等。}\par

1. 这位客纳罕妇人，刚刚在福音的课文中被带到我们面前，给我们展示了谦逊的榜样以及虔诚的道路；她教导我们如何从谦逊升高。\newline{}现在，很明显，她不是以色列子民，以色列出了圣祖、先知以及按照血肉是主耶稣基督的父母；童贞玛利亚本身也是这子民中的一位，她是基督之母。\newline{}那么这位妇人不是这子民中的，而是外邦人。\newline{}因为，正如我们所听到的，主\Quote{退到提洛和漆冬一带，看，有一个客纳罕妇人从那一带出来}，极其恳切地祈求祂怜悯，好医治她的女儿，\Quote{她被魔鬼痛苦地折磨}。\newline{}提洛和漆冬不是以色列子民的城市，而是外邦人的，尽管它们与该子民接壤。\newline{}因此，由于渴望得到怜悯，她呼叫并大胆地敲门；祂却装作没有听见她，不是为了拒绝怜悯她，而是为了激发她的渴望；并且不仅激发她的渴望，而且如我之前所说，为了彰显她的谦逊。\newline{}所以她呼叫，而主装作没听见，却在静默中安排祂将要作的事。\newline{}门徒们为主代求，说：\Quote{打发她走吧，因为她在我们后面喊叫。} \newline{}祂说：\Quote{我被派遣，只是为了以色列家迷失的羊。}\par

2. 从这些话中产生了一个问题：\Quote{如果祂没有被派遣，除非到以色列家迷失的羊那里，我们外邦人如何进入基督的羊栈？这奥秘如此深刻的安排是什么意思？主知道祂来的目的——祂要在万民中有一个教会，祂却说\InnerQuote{祂没有被派遣，除非到以色列家迷失的羊那里}？} 我们由此明白，祂必须在那个民族中彰显祂的肉身显现、祂的诞生、祂的奇迹的展示以及祂复活的能力：这样规定，这样从一开始就预定了，这样预言了，这样应验了；基督耶稣要来到犹太人的民族，被看见并被杀害，并从他们中赢得祂所预知的人。\newline{}因为那个民族并非完全被定罪，而是被筛过。他们中间有大量的糠秕，但也有隐藏的麦粒的价值；他们中间有要被焚烧的，也有要装满谷仓的。\newline{}因为宗徒们从何而来？伯多禄从何而来？其他人从何而来？\par

3. 保禄本人从何而来？他起先被称为撒乌耳。那是指起初骄傲，后来谦逊吗？因为当他还是撒乌耳时，他的名字出自撒乌耳：撒乌耳是个骄傲的君王；他在位时迫害谦逊的达味。所以当后来成为保禄的人还是撒乌耳时，他是骄傲的，那时是义人的迫害者，那时是教会的摧残者。\newline{}因为他从大司祭那里领了书信（他因对会堂的热心而焚烧，并迫害基督徒的名号），以便他可以搜出他找到的任何基督徒，加以惩罚。\newline{}他正走在路上，正喘着杀气，正渴慕血腥，却被基督从天上来的声音击倒在地——这个迫害者，被提拔起来成为宣道者。\newline{}在他身上应验了先知所记载的：\Quote{我打伤，我也医治。} 因为天主打伤的，只是人身上那自高反抗天主的部分。\newline{}那位切开肿胀、切割或烧灼腐烂部位的外科医生并非不仁慈。他确实带来痛苦；但他带来痛苦，只为使病人恢复健康。他带来痛苦；但如果他不这样做，他就无法行善。\newline{}基督于是用一句话将撒乌耳击倒，将保禄扶起；也就是说，祂击倒骄傲的，扶起谦逊的。\newline{}他改名的原因是什么？为什么他从前叫撒乌耳，后来却选择叫保禄？不就是因为他承认自己身上那作为迫害者时的名字\Quote{撒乌耳}是骄傲之名吗？\newline{}所以他选择了一个谦逊的名字：被称为保禄，也就是最小的。因为保禄就是\Quote{最小的}。保禄无非是\Quote{微小}。\newline{}如今他以这个名字为荣，并给我们上了一堂谦逊的课，他说：\Quote{我原是宗徒中最小的。} \newline{}那么，他又是从哪里来的呢？不就是从犹太人那里吗？其他宗徒也是从他们中来的，保禄也是从他们中来的，同一位保禄提到的那些见过主复活的人也是从他们中来的。\newline{}因为他说：\Quote{祂曾同时显现给五百多位弟兄；其中大部分至今还在，有些已经长眠。}\par

4. 这个民族，即犹太民族中，也有这样一些人：当\Person{伯多禄}讲述\Person{基督}的受难、复活和天主性时（在领受圣神之后，所有领受圣神的人都说起万国的方言），他们听了他的话，心中刺痛，为自己的得救寻求劝告，因为他们明白自己犯了流\Person{基督}血的罪过；因为他们钉死了祂，杀死了祂，而他们却看到以祂的名（尽管被他们所杀）行了如此伟大的奇迹，并看到圣神的临在。于是他们寻求劝告，得到答复说：\BibleQuote{你们悔改，并要以我们的主\Person{耶稣基督}的名受洗，使你们的罪得赦。}当杀害\Person{基督}的罪也被赦免时，谁还该绝望于自己的罪得赦呢？他们从犹太民族中悔改，悔改并受洗了。他们来到主的祭台前，凭着信德饮下了他们曾在愤怒中倾流的血。他们悔改的程度、决心和完全性，\WorkTitle{宗徒大事录}都有显示。\BibleQuote{因为他们变卖了所有的一切，把卖得的钱放在宗徒们脚前；然后照各人所需用的分给各人；没有一个人说有什么东西是自己的，一切都公有公用。}并且，如经上所载：\BibleQuote{他们同心合意。}看哪，这就是祂所说的羊：\BibleQuote{我被派遣，只是到以色列家迷失的羊那里去。}因为祂向他们显示了祂的临在，当祂被钉十字架时，祂为他们祈祷说：\BibleQuote{父啊，宽恕他们，因为他们不知道他们做了什么。}这位医师知道，那些狂怒的人正在疯狂中杀害医师，在杀害他们的医师时，虽然他们不知道，却为自己准备了药物。因为借着被杀害的主，我们所有人都被治愈了，被祂的血救赎了，靠着圣体之饼脱离了饥荒。\Person{基督}就是这样向犹太人显示了祂的临在。所以祂说：\BibleQuote{我被派遣，只是到以色列家迷失的羊那里去。}以便向他们显示祂身体的临在；并非祂不顾念和越过祂在外邦人中的羊。\par

5. 因为祂自己没有到外邦人那里去，而是派遣了祂的门徒。先知的话便应验了：\BibleQuote{一个我不认识的民族服事了我。}请看，这预言是多么深奥、清晰、明确：\InnerQuote{一个我不认识的民族}（即我未曾向他们显示我的临在）\InnerQuote{服事了我。}如何呢？接着又说：\BibleQuote{他们凭着耳听服从了我。}即他们不是凭眼见，而是凭耳听而相信了。因此外邦人更值得称赞。因为犹太人看见并杀了祂；外邦人听见并相信了。现在，为了召集和聚集外邦人，使那句我们刚才咏唱的经文\BibleQuote{从外邦人中聚集我们，使我们称谢你的名，以你的赞美为荣耀}得以应验，宗徒\Person{保禄}被派遣了。他，最小的，变得伟大，不是靠自己，而是靠祂曾迫害的那位，被派往外邦人，从强盗变成牧人，从狼变成羊。这位最小的宗徒被派往外邦人，在外邦人中辛勤劳作，通过他，外邦人相信了。他的书信就是见证。\par

6. 在\WorkTitle{福音}中，你们也看到一幅非常神圣的预象。有一个会堂长的女儿真的死了，她父亲恳求主到她那里去；他离开她时，她正在生病，情况极度危险。主出发去探望和医治病人；正在这时，有人来报告说女孩已经死了，并告诉父亲：\BibleQuote{你的女儿死了，不要麻烦老师了。}但主知道自己能复活死人，并没有使绝望的父亲失去希望，对他说：\BibleQuote{不要怕，只要信。}于是祂继续向女孩走去；在路途中，有一个患血漏的女人，她在长期的病痛中已经把所有财产都花在医生身上，但毫无效果，她挤入人群中，尽她所能。当她触摸到祂衣穗时，就痊愈了。主说：\BibleQuote{谁摸了我？}门徒们不知道发生了什么事，他们看到主被众人拥挤，而祂却为那个轻轻触摸祂的女人操心，惊讶地回答说：\BibleQuote{群众拥挤着您，您还说谁摸了我？}祂说：\BibleQuote{有人摸了我？因为那些人在拥挤，她却触摸了。}接着又说：\BibleQuote{有人摸了我，因为我感觉到有德能从身上出去了。}那女人见自己隐藏不住，就俯伏在祂脚前，承认所发生的事。之后，祂又出发了，到达了目的地，将那位会堂长死去的女儿复活了。\par

7. 这是一个按字面解释的事实，并且如所述那样实现了；然而主所做的这些事本身还具有更深的意义，它们（如果可以这样说）是一种可见且有意义的言语。这一点尤其明显在：祂在不在果子的时节在树上找果子，因为找不到，就用诅咒使树枯干。除非将这个行动视为一个象征，否则其中没有任何好的意义；首先，在不是任何树结果子的时节，在那棵树上找果子；其次，即使那时正是结果子的时节，树没有果子又是什么过错呢？但因为这意味着，祂寻找的不只是叶子，也是果子，也就是说，不只是人的言语，也是人的行为，祂使那棵只找到叶子的树枯干，象征着那些能说好话却不愿行善之人的惩罚。在此处也是如此。因为其中确实有一个奥秘。祂预知一切，却说：\Quote{谁摸了我？}造物主使自己像一个无知的人；祂询问不仅知道此事，连其他一切也都预知的那一位。无疑，基督要通过这个意味深长的奥秘向我们说话。\par

8. 那会堂主管的女儿是犹太人民的象征，基督是为他们而来的，祂说：\BibleQuote{我奉差遣，不过是到以色列家迷失的羊那里去。}\BibleRef{玛 15:24}但那位患血漏的妇女，象征着来自外邦人的教会，基督并未以肉身临在于此。祂正前往前者，专注于她的康复；与此同时，后者跑向祂，仿佛在祂不知情的情况下触摸了祂的衣边；也就是说，她得到了那位在某种意义上缺席者的医治。祂说：\Quote{谁摸了我？}好像祂要说：我不认识这人民；\BibleQuote{我素不认识的百姓要侍奉我。}\BibleRef{咏 18:43}\Quote{有人摸了我，因为我觉得有能力从我身上出去了。}那就是说，我的福音已经传出去，充满了全世界。现在触摸的是衣边，是衣服上细小和外在的部分。把宗徒们看作基督的衣服。在他们当中，保禄是衣边，即最后和最小的。因为他自称两者都是：\Quote{我是宗徒中最小的。}因为他是在他们之后蒙召的，在他们之后相信的，但他的医治比他们所有人都多。主并非被派遣，除非是\BibleQuote{到以色列家迷失的羊那里去。}\BibleRef{玛 15:24}但是，因为\BibleQuote{我素不认识的百姓也要侍奉祂，}\BibleRef{咏 18:43}并在耳中听到就顺从，祂在其他人中间也提到了他们。因为同一主在某处说：\BibleQuote{我还有别的羊，不属于这一羊栈；我也必须领他们来，这样将只有一个羊群，一个牧人。}\BibleRef{若 10:16}\par

9. 这妇女就是其中之一；因此她没有被拒绝，只是被延缓。祂说：\BibleQuote{我奉差遣，不过是到以色列家迷失的羊那里去。}\BibleRef{玛 15:24}而她坚持呼喊：她恒心，她敲门，好像她已经听到了：\BibleQuote{你们求，就要得到；你们找，就要找到；你们敲，就要给你们开。}\BibleRef{玛 7:7}她坚持，她敲门。因为当主说这些话时：\BibleQuote{你们求，就要得到；你们找，就要找到；你们敲，就要给你们开；}祂先前也曾说过：\BibleQuote{不要把圣物给狗，也不要把你们的珍珠丢在猪前，免得它们用脚践踏，并且转过来撕裂你们。}\BibleRef{玛 7:6}也就是说，免得轻看你们的珍珠，甚至加害你们。因此，不要将他们所轻视的东西丢在他们面前。\par

10. 然而我们如何区分（有人可能会问）谁是\Quote{猪}、谁是\Quote{狗}？这在那个妇女的个案中已经显明了。因为祂只回应她的恳求说：\BibleQuote{拿儿女的饼丢给狗吃，是不合适的。}\BibleRef{玛 15:26}\Quote{你是狗，你是外邦人之一，你崇拜偶像。}但狗最合适做什么呢？不是舔石头吗？\BibleQuote{因此，拿儿女的饼丢给狗吃，是不合适的。}\BibleRef{玛 15:26}如果她在这些话后离开，她就像来时一样是狗；但通过敲门，她从狗变成了人。因为她恒心求告，并且通过那似乎是一种责备，她显出了谦卑，获得了怜悯。因为她没有被激怒、没有发怒，虽然被称为狗，她仍然求祝福，求怜悯，却说：\Quote{主啊，是的；\BibleRef{玛 15:27}你称我为狗，我确实是狗，我承认我的名字：是真理在说话；但我不应该因此而拒绝这个祝福。我确实是狗；\InnerQuote{可是狗也吃主人桌子上掉下来的碎渣。}我所渴望的不过是适度的、小小的祝福；我不挤到桌子前，我只找一些碎渣。}\par

11. 请看，弟兄们，谦卑的价值是如何呈现在我们面前的！主曾称她为狗；她并没有说\Quote{我不是}，反而说\Quote{我是}。因为她承认自己是狗，主立刻说：\BibleQuote{妇人，你的信德真大！就照你所求的给你成就吧。}你已经承认自己是狗，我现在承认你是人。\Quote{妇人，你的信德真大；}你已经求了，找了，敲了；你必得到，找到，门必为你打开。请看，弟兄们，在这个客纳罕妇人身上——她来自外邦人，是教会的预象，即象征——谦卑的恩宠如何卓越地呈现在我们面前。因为犹太民族，为了被剥夺福音的恩宠，被骄傲所膨胀；因为他们曾蒙恩典领受法律，从这个民族中产生了圣祖、先知，天主的仆人\Person{梅瑟}曾在埃及行了我们在圣咏中听过的大奇迹，带领百姓从分开的水中穿过红海，领受了法律，并传给了这个民族。这正是犹太民族骄傲的原因，而正是这种骄傲使他们不愿向谦卑的创始者、骄傲膨胀的抑制者基督，向医生天主谦卑自己；天主本是天主，却因此成了人，好让人知道自己不过是人。伟大的药方！如果这药方不能治愈骄傲，我不知道什么能治愈。祂是天主，却成了人；祂放下祂的天主性，即某种程度地隐藏了祂原有的东西，只显示祂所取的人性。祂是天主，却成了人；而人却不愿承认自己是人，即不愿承认自己是必死的、软弱的、有罪的、患病的，好至少作为病人去寻求医生；但更危险的是，他自以为健康。\par

12. 因此，那个民族没有来到祂面前，就是因为骄傲；而且本来的枝条被描述从橄榄树上折下来，即从圣祖所建立的民族中被折下；换句话说，犹太人因骄傲的精神受罚而公正地不结果实；野橄榄被接在橄榄树上。野橄榄树就是外邦人的民族。宗徒这样说：\BibleQuote{野橄榄被接在好橄榄树上，但本来的枝条被折下来。}因骄傲他们被折下；野橄榄因谦卑而被接上。那妇人彰显了这种谦卑，当她说：\BibleQuote{主，是的，我是狗，我只希望得到碎屑。}百夫长也因这种谦卑而中悦祂；他请求主治愈他的仆人，主说：\BibleQuote{我去治好他。}百夫长回答说：\BibleQuote{主，我不堪当祢到我的屋檐下来，但只要祢说一句话，我的仆人就会痊愈。我不堪当祢到我的屋檐下来。}他没有接待主到他的房屋，却已在心中接待了祂。越是谦卑，就越能容纳，也越满溢。因为山岭阻挡水流，而山谷却被水充满。那么，当主说了\BibleQuote{我不堪当祢到我的屋檐下来}之后，祂对那些跟随祂的人说了什么呢？\BibleQuote{我实在告诉你们，在以色列我没有见过这么大的信德。}即在我所来的那个民族中，\BibleQuote{我没有见过这么大的信德。}这信德从何而来？大是因着极小，即因着谦卑。\BibleQuote{我没有见过这么大的信德；}如同芥菜子，它越小，就越火热。因此主立刻将野橄榄接在好橄榄树上。祂这样做是当祂说：\BibleQuote{我实在告诉你们，在以色列我没有见过这么大的信德。}\par

13. 最后，请注意下面的话。\Quote{因此}——即因为\BibleQuote{我在以色列没有见过这么大的信德}，即如此伟大的谦卑与信德——\BibleQuote{因此我告诉你们，许多人将从东方和西方来，与亚巴辣罕、依撒格和雅各伯在天国里同席。}\Quote{同席}，即\Quote{安息}。因为我们不应在那里想象肉体的筵席，也不该贪求那国度中的任何此类事物，以致不是以美德代替恶习，而只是交换恶习。因为渴慕天国是为了智慧和永生是一回事；为了世俗的福乐，以为在那里我们会更丰富、更大量地拥有它，是另一回事。如果你以为在那国度中会富有，你并没有断绝欲望，只是改变了欲望；然而你确实会富有，而且除了那里之外别无他处你会富有；因为在这里，你的匮乏聚集了事物的丰裕。为什么富人拥有许多？因为他们需要许多。更大的匮乏堆砌了更大的财富；在那里，匮乏本身将死亡。那时你将真正富有，当你一无所缺时。因为现在你确实不富有，而天使贫穷，他没有马、车和仆人。为什么？因为他不需要这些：因为他的能力越大，匮乏就越少。因此那里有财富，真正的财富。所以不要为那地方想象这世上的筵席。因为这世上的筵席是每日的药；它们对于我们所拥有的某种疾病是必要的，我们与生俱来的疾病。这种疾病，每个人在过了用餐的时刻时都会感觉到。你愿意看见这疾病有多大吗？就像急性热病在七天内会致命一样？那么不要自以为健康。不朽将是健康。因为现在只是一场漫长的疾病。因为你用每日的药物支持你的疾病；你自以为健康；拿掉药物，然后看看你能做什么。\par

14. 因为我们从出生那一刻起，就必然在走向死亡。这种病必然将我们带向死亡。医生在诊察病人时确实这样说。例如：\Quote{这个人患了水肿，他快要死了；这病无法医治。这个人患了癞病：这病也无法医治。他得了肺痨。谁能医治呢？他必死，他必定灭亡。}看，医生已经断定他得了肺痨，他必死无疑；然而有时水肿病人并没有死于他的病，癞病病人也没有，肺痨病人也没有；但如今，每一个出生的人都必然死于此。他死于此，别无选择。医生和没有学识的人都能断定这一点；即使他死得稍慢一些，难道就不死了吗？那么，哪里才有真正的健康呢？除非在真正的永生那里。但如果那是真正的永生，没有腐朽，没有消耗，那里哪里还需要营养呢？因此，当你听见说：\BibleQuote{他们要与\Person{亚巴郎}、\Person{依撒格}、\Person{雅各伯}一同坐席。}要整理的不是你的身体，而是你的灵魂。在那里，你将被饱饫；这个内在的人有它合适的食粮。与此相关，经上说：\BibleQuote{饥渴慕义的人是有福的，因为他们将得饱饫。}他们将如此真实地饱饫，以致不再饥饿。\par

15. 因此，主立刻嫁接野橄榄枝，当他说：\BibleQuote{将有许多人从东方和西方来，同\Person{亚巴郎}、\Person{依撒格}、\Person{雅各伯}在天国里一起坐席；}也就是说，他们将被嫁接在好橄榄树上。因为\Person{亚巴郎}、\Person{依撒格}、\Person{雅各伯}是这橄榄树的根；但\BibleQuote{但天国之子}，即不信的犹太人，\BibleQuote{要到外面的黑暗里去。}\BibleQuote{天然的枝条将被折下，好让野橄榄枝嫁接上去。}那么，天然的枝条为何该被剪除呢？岂非因为骄傲？野橄榄枝为何被嫁接呢？岂非因为谦逊？正因如此，那妇人也说：\BibleQuote{主！是的，可是小狗也吃主人桌子上掉下来的碎屑。}于是她听见：\BibleQuote{妇人，你的信德真大！}同样，那位百夫长也说：\BibleQuote{我当不起你到我的舍下来。}\BibleQuote{我实在告诉你们：在\Place{以色列}，我从未见过这样大的信德。}让我们学习谦逊，或持守谦逊。如果我们尚未拥有，让我们学习它；如果我们已经拥有，让我们不要失去它。如果我们尚未拥有，让我们拥有它，好使我们被嫁接进去；如果我们已经拥有，让我们持守它，好使我们不被剪除。\par

\SourceRecord{Translated by R.G. MacMullen.；Nicene and Post-Nicene Fathers, First Series；Vol. 6.；Edited by Philip Schaff.；Buffalo, NY: Christian Literature Publishing Co.,；1888.；<http://www.newadvent.org/fathers/160327.htm>.}

% structure: p0001 source=h1 target=section reason=page-title

\section{新约讲道第28篇}

\Emph{[LXXVIII. Ben.]}\par

\Emph{论福音之言}，见：\BibleRef{玛 17:1}，\Emph{\Quote{六天后，耶稣带着伯多禄、雅各伯和他的兄弟若望，}等。}\par

我们现在必须考察并论述主在山中所显示的那个神视。因为正是这个神视，祂曾说过：\Quote{我实在告诉你们，站在这里的有些人还没有尝到死味，就必看见人子来到祂的国里。}见：\BibleRef{玛 16:28} 接着便开始了刚才所读的经文。\Quote{祂说完这些话，六天后，带着三个门徒——伯多禄、雅各伯和若望——上了一座山。}见：\BibleRef{玛 17:1} 这三人便是祂所说的\Quote{有些人}，即\Quote{站在这里的有些人还没有尝到死味，就必看见人子来到祂的国里}的人。这里有一个不小的难题：因为那座山并非祂王国的全部；对于拥有诸天的那位，一座山算什么呢？我们不仅读到祂拥有诸天，而且在某种程度上用心灵的眼睛看见它。祂称那为祂的王国，在许多地方祂也称之为\Quote{天国}。天国就是圣人们的王国。\Quote{因为诸天述说天主的荣耀。}见：\BibleRef{咏 19:2} 并且在这诗篇中紧接着说：\Quote{没有语言，没有话语，也听不到他们的声音。它们的声音传遍普世，它们的言语传到地极。}见：\BibleRef{咏 19:3-5} 这是谁的话语？正是诸天的话语，也就是宗徒们和一切忠信宣讲天主圣言的人。因此，这些诸天将与被造诸天的那一位一同为王。现在思考所成就的事，好使这事显明出来。\par

2. 主耶稣自身明亮如太阳，祂的衣服洁白如雪，梅瑟和厄里亚与祂交谈。耶稣自身像太阳一样发光，表明\Quote{祂是照亮每一个来到世上的人的真光。}见：\BibleRef{若 1:9} 这太阳对肉体的眼睛如何，祂对心灵的眼睛也如何；那太阳对肉身的照射如何，祂对心灵的照射也如何。至于祂的衣服，那就是祂的教会。因为衣服若不靠着穿着它的人，就会脱落。在这件衣服上，保禄如同一道最后的衣边。因为他自己说：\Quote{我是宗徒中最小的一个。}见：\BibleRef{格前 15:9} 在另一处又说：\Quote{我是宗徒中最后的一个。}如今在衣服上，衣边是最后最小的部分。因此，正如那位患血漏的妇人，在触摸主的衣边后得了痊愈，同样，由外邦人中出来的教会，因保禄的宣讲得了痊愈。当你听见先知依撒意亚说：\Quote{你们的罪虽似朱红，我将使它们白如雪}见：\BibleRef{依 1:18} 时，有什么可奇怪的呢？梅瑟和厄里亚，即法律和先知，若不同主交谈，有什么益处呢？除非它们为主作证，否则谁会读法律或先知呢？请注意宗徒如何简练地表达这一点：\Quote{因为法律带来对罪的认识；但如今，天主的正义在法律之外已彰显出来}——看，太阳；\Quote{有法律和先知作证}——看，太阳的光辉。见：\BibleRef{罗 3:20-21}\par

3. 伯多禄看见这情景，作为思虑属世之事的人，他说：\Quote{主，我们在这里真好！}见：\BibleRef{玛 17:4} 他因人群而感到疲惫，如今找到了山间的宁静；在那里他有了基督——灵魂的食粮。难道他还要从那里再回到劳苦和痛苦中去吗？他怀有对天主的圣爱，因而有良好的品行？他为自己着想，于是补充说：\Quote{你若愿意，我就在这里搭三个帐棚：一个为你，一个为梅瑟，一个为厄里亚。}主没有回答；但伯多禄仍然得到了答复。\Quote{因为他还在说话的时候，忽然有一片光耀的云彩笼罩了他们。}见：\BibleRef{玛 17:5} 他想要三个帐棚；天上的答复告诉他，我们只有一个，而人的判断却想把它分开。基督——天主的圣言，在法律中的圣言，在先知中的圣言。伯多禄，你为什么想要分开他们？你更应该把他们联合起来。你寻求三个；要知道它们本是一个。\par

4. 当云彩笼罩了他们，并仿佛为他们搭了一个帐棚时，\Quote{从云中有声音说：这是我的爱子。}见：\BibleRef{玛 17:5} 梅瑟在那里，厄里亚也在那里；但并没有说：\Quote{这些都是我的爱子。}因为独生子是一回事，义子又是另一回事。祂被单独选出，法律和先知都荣耀祂。\Quote{这是我的爱子，我所喜悦的，你们要听祂！}因为你们已在先知中听见了祂，在法律中听见了祂。你们在哪里没有听见祂呢？\Quote{门徒听见了，就俯伏在地。}因此，在教会中，天主之国已显示给我们。这里有主，这里有法律和先知；但主是主，法律在梅瑟身上，预言在厄里亚身上；他们只是仆人和事奉者。他们是器皿，祂是泉源。梅瑟和先知们说话、写作；但他们所倾注的，都是从祂那里汲取的。\par

5. 但主伸出祂的手，扶起倒地的他们。然后\Quote{他们不见一人，只见\Person{耶稣}独自一个。}这是什么意思？当宗徒书信被宣读时，你们听到了：\BibleQuote{因为我们现在是借着镜子观看，模糊不清，到那时就要面对面了。}以及\BibleQuote{语言之恩终必停止，}那时我们所盼望和相信的将要实现。他们仆倒在地，象征我们死亡，因为曾对肉体说：\BibleQuote{你原是土，还要归于土。}但当主扶起他们时，祂象征了复活。复活后，法律对你们还有什么？先知呢？因此，既不见\Person{梅瑟}，也不见\Person{厄里亚}。唯有祂留给你们：\BibleQuote{在起初已有圣言，圣言与天主同在，圣言就是天主。}祂留给你们：\BibleQuote{好叫天主成为万物之中的万有。}\Person{梅瑟}将在那里；但法律不再有了。我们也将看见\Person{厄里亚}在那里；但先知不再有了。因为法律和先知只是为基督作证：祂必须受难，第三日从死者中复活，并进入祂的光荣。在这光荣中，祂对所爱之人所作的应许得以实现：\BibleQuote{爱我的，我父必爱他，我也要爱他。}好像有人说：祢既爱他，要给祢什么？\BibleQuote{我要亲自向他显现。}伟大的恩赐！伟大的应许！天主不把任何自己的东西留给你作赏报，而是赐给你自己。哦，你这贪得无厌的人；基督所应许的还不够吗？你自以为富足；但若没有天主，你有什么呢？别人贫穷，但有天主，他又缺什么呢？\par

6. 下来，\Person{伯多禄}：你曾想留在山上安息；下来，\BibleQuote{宣讲圣言，不论顺境逆境，务必坚持，以全般的坚忍和训导，指责、警戒、劝勉。}忍受、劳苦，承担你的苦难分量；好使你获得主的白衣所象征的，即因在爱德中正直劳苦而有的光明与美丽。因为当宗徒书信被宣读时，我们听到了赞美爱德的话：\BibleQuote{她不求自己的利益。她不求自己的利益；}因为她施与自己所拥有的。在另一处，这句话若理解不当，则更加危险。因为宗徒按爱德的原则命令基督的忠实肢体，说：\BibleQuote{不可寻求自己的利益，而应寻求他人的利益。}听到这话，贪欲便准备好欺骗，在生意场合，以寻求他人利益为借口，去诈骗别人，所以是\BibleQuote{不求自己的利益，而求他人的利益。}但让贪欲克制自己，让正义显现；这样我们才能听而明白。对爱德所说的正是：\BibleQuote{不可寻求自己的利益，而应寻求他人的利益。}现在，哦，你这贪得无厌的人，如果你仍要抗拒，并将这诫命曲解为你应贪求他人的东西，那么就会失去自己的东西。但据我所知，你既想拥有自己的，又想拥有他人的。你将行欺诈来获得他人之物；然后忍受被抢而失去自己之物。你不愿寻求自己的利益，却去夺取他人的。你这样做是不好的。听着，你这贪得无厌的人：宗徒在另一处更清楚地解释了他所说的：\BibleQuote{不可寻求自己的利益，而应寻求他人的利益。}他论到自己说：\BibleQuote{不求自己的利益，但求众人的利益，为使他们得救。}这事\Person{伯多禄}还不明白，当他渴望与基督在山上同住时。这事祂为你保留，\Person{伯多禄}，在死后。但现在主亲自说：\Quote{下来，在地上劳苦；在地上服事，被轻视，被钉死在地上。生命下来了，为要被杀；食粮下来了，为要饥饿；道路下来了，为要在路上疲倦；泉源下来了，为要口渴；而你不愿劳苦吗？\InnerQuote{不求自己的利益。}拥有爱德，宣讲真理；这样你将进入永恒，在那里你必得安稳。}\par

\SourceRecord{Translated by R.G. MacMullen.；Nicene and Post-Nicene Fathers, First Series；Vol. 6.；Edited by Philip Schaff.；Buffalo, NY: Christian Literature Publishing Co.,；1888.；<http://www.newadvent.org/fathers/160328.htm>.}

% structure: p0001 source=h1 target=section reason=page-title

\section{新约讲道集 第29篇}

\EditorialNote{[LXXIX. Ben.] 本尼狄克版本编号}\par

\Emph{再论福音之言}，\BibleRef{玛 17}，其中\Person{耶稣}在山上向三位门徒显示自己。\par

1. 我们聆听圣福音宣读时，听到了山上伟大的神视，\Person{耶稣}在山上向三位门徒——\Person{伯多禄}、\Person{雅各伯}和\Person{若望}——显示了自己。\BibleQuote{祂的脸面发光如太阳}，这是福音光辉的图像。\BibleQuote{祂的衣服洁白如雪}，这是教会纯洁的图像，关于这一点，先知曾说过：\BibleQuote{你们的罪虽似朱红，我必使之洁白如雪}。\Person{厄里亚}和\Person{梅瑟}在与祂交谈，因为福音的恩宠得到法律和先知的见证。法律由\Person{梅瑟}代表，先知由\Person{厄里亚}代表，简单地说。因为有一位圣殉道者，天主通过他要叙述祂所赐的仁慈。让我们倾听。\Person{伯多禄}愿意搭三个帐棚，一个为\Person{梅瑟}，一个为\Person{厄里亚}，一个为\Person{基督}。山上的孤寂吸引了他；他因世务的喧嚣而疲惫。但他为什么寻求三个帐棚，岂不是因为他还不认识法律、先知和福音的统一？最后，他被云彩纠正了：\BibleQuote{他还在说话的时候，看，有一朵光明的云彩遮蔽了他们}。看，云彩成了一个帐棚；你为什么寻求三个呢？\BibleQuote{从云彩里有声音说：这是我的爱子，我所喜悦的，你们要听祂}。\Person{厄里亚}说话，但\BibleQuote{你们要听祂}；\Person{梅瑟}说话，但\BibleQuote{你们要听祂}。先知说话，法律说话，但\BibleQuote{你们要听祂}。祂是法律的声音，先知的舌头。祂在他们内说话，当祂俯允这样做时，祂亲自显现。\BibleQuote{你们要听祂}：那么让我们听祂。当福音说话时，要认为那是云彩：从那里发出声音传到我们。让我们听祂；就是说，让我们做祂所说的，让我们盼望祂所应许的。\par

\SourceRecord{Translated by R.G. MacMullen.；Nicene and Post-Nicene Fathers, First Series；Vol. 6.；Edited by Philip Schaff.；Buffalo, NY: Christian Literature Publishing Co.,；1888.；<http://www.newadvent.org/fathers/160329.htm>.}

% structure: p0001 source=h1 target=section reason=page-title

\section{新约讲道集 第30篇}

\Emph{[LXXX. Ben.]}\par

\Emph{论福音的话，\BibleRef{玛 17:19}，\BibleQuote{我们为什么不能赶出它呢？}等等，以及论祈祷。}\par

1. 我们的主耶稣基督斥责祂自己的门徒的不信，正如我们刚在读福音时所听到的。因为他们曾说：\BibleQuote{我们为什么不能赶出他呢？}祂回答说：\BibleQuote{因为你们的不信。}如果宗徒是不信者，谁又是信徒呢？如果公羊摇摆不定，羔羊当如何行？然而主的仁慈并未藐视他们的不信，而是斥责、滋养、成全、给他们加冕。因为他们自己也意识到自己的软弱，便向祂说，正如我们在福音某处所读到的：\BibleQuote{主，增加我们的信德。}他们说：\BibleQuote{主，}\BibleQuote{增加我们的信德。}知道自己有所欠缺，是第一个益处；更大的幸福是知道他们向谁祈求。\BibleQuote{主，}他们说\BibleQuote{增加我们的信德。}看，他们岂不是仿佛将自己的心带到泉源，叩门以求打开，好让他们能充满吗？因为祂愿意人向祂叩门，并非要拒斥叩门者，而是要操练那些渴望的人。\par

2. 弟兄们，你们以为天主不知道你们需要什么吗？祂知道且先于我们的渴望，祂了解我们的匮乏。因此，当祂教导门徒祈祷，并警告他们不要用许多言辞时，祂说：\BibleQuote{不要用许多言辞；因为你们的父在你们祈求祂以前，已知道你们需要什么。}然而主所说的与此有所不同。这是什么意思？因为祂不喜悦我们在祈祷中用许多言辞，便对我们说：\BibleQuote{当你们祈祷时，不要用许多言辞；因为你们的父在你们祈求祂以前，已知道你们需要什么。}如果我们的父在你们祈求祂以前，已知道你们需要什么，那么为什么我们甚至用很少的言辞呢？如果我们的父已经知道我们需要什么，那么祈祷有什么用处呢？祂对一个人说：不要向我作冗长的祈祷；因为我知道你需要什么。若是这样，主啊，我为什么还要祈祷呢？祢不愿意我用长祷，祢甚至几乎命令我完全不要祈祷。那么，另一处的诫命又是什么意思？因为那位说\BibleQuote{在祈祷中不要用许多言辞}的，在另一处说：\BibleQuote{求，就必得到。}并且，免得你以为这第一个求的诫命是随口而出的，祂又加上：\BibleQuote{寻找，就必找到。}并且，免得你以为这也是随口而出的，请看祂又加上什么，请看祂以什么结束。\BibleQuote{叩门，就必给你们开门；}请看祂加上什么。祂愿意你们求，以便领受；寻找，以便找到；叩门，以便进入。既然我们的父已经知道我们需要什么，我们为何又如何祈求？为何寻找？为何叩门？为何在祈求、寻找、叩门中劳苦自己，去教导那位已经知道的？在另一处，主的话是：\BibleQuote{人应当时常祈祷，不可灰心。}如果人应当时常祈祷，那么祂怎么说\BibleQuote{不要用许多言辞}？若是我很快就结束，我怎能时常祈祷？这里祢命我快快结束，那里\BibleQuote{时常祈祷，不可灰心}：这是什么意思？为使你们理解这一点，\Quote{求，寻找，叩门。}因为门关闭，并非为了将你们关在外面，而是为了操练你们。因此，弟兄们，我们应当劝勉祈祷，既劝勉自己，也劝勉你们。因为在今世的多重邪恶中，我们别无希望，只有借着祈祷叩门，相信并坚定地持守心中的信念：你们的父只赐给你们祂知道对你们有益的事物。因为你们知道你们所愿的；祂知道什么对你们有益。设想你们在一位医生手下，身体虚弱，正如实际情况一样；因为我们整个生命都是虚弱；长寿不过是延长的虚弱。那么设想你们在医生手中患病。你们渴望请求医生允许喝一杯新酒。你们并未被禁止请求，因为或许这对你们无害，甚至接受它有益。所以不要犹豫去请求；请求，不要犹豫；但如果你们未得到，不要放在心上。如果你们在人的手中，在身体医生的手中会如此行，那么在天主手中，祂是你们身心的医生、创造者和恢复者，又该如何呢？\par

3. 因此，请看主在这段经文中如何劝勉祂的门徒祈祷，祂说：\BibleQuote{你们不能因你们的不信赶出这魔鬼。}因为那时祂劝勉他们祈祷，便这样结束：\BibleQuote{这类非藉祈祷和禁食不能赶出。}若人必须祈祷，才能从他人身上赶出魔鬼，更何况赶出自己的贪婪？更何况赶出自己的酗酒？更何况赶出自己的奢侈？更何况赶出自己的不洁？人身上有多少事物，若坚持不舍，便不容许进入天国！弟兄们，请想想，为保存暂时的健康，人是如何恳求医生；若有谁病入膏肓，他岂会羞怯或迟疑俯伏在人脚前？以泪水沾湿某位极能干的首席医生的脚步？若医生对他说：\Quote{你除非让我捆绑你，并使用火和刀，否则不能痊愈。}他会回答：你愿做什么就做什么，只求你治好我。他以何等热切渴望那如云雾般短暂的数日健康，为了它，他甘愿被捆绑，忍受火与刀，受监视，不得随意吃喝！这一切他都忍受，为能晚一点死；然而他却不肯忍受一丝一毫，为能永不死亡。若天主，就是那在我们之上的天上医生，对你说：\Quote{你愿意痊愈吗？}你除了说\Quote{愿意}之外还能说什么？或者你也许不会这样说，因为你自以为健康，也就是说，因为你病得更重。\par

4. 因为我们若设想两个病人，一个含泪恳求医生，另一个则病中狂乱地嘲笑医生；医生会对那哭泣者怀有希望，而哀叹那发笑者的境况。为什么？因为人自认越健康，其病就越危险！犹太人正是如此。基督来到他们这些病人中间；祂发现他们全都病了。因此，谁也不可因自己貌似健康而自诩，免得医生放弃他。祂发现全都病了；这是宗徒的判断：\BibleQuote{因为众人都犯了罪，缺乏了天主的光荣。}虽然祂发现他们都病了，但其中有两种病人。一种来到医生面前，紧依基督，聆听、尊敬、跟随祂，并归依了。祂接纳所有人，不鄙视任何人，为要医治他们——祂无偿地医治，以全能之力治愈。当祂接纳他们，并使他们与祂结合以得医治时，他们便喜乐了。但另一种病人，已因不义之病而变得狂乱，竟不知自己有病；他们嘲笑祂，因祂接纳病人，并对祂的门徒说：\BibleQuote{看，你们的师傅和税吏及罪人一起吃饭。}祂知道他们是什么人，便回答说：\BibleQuote{健康的人不需要医生，有病的人才需要。}祂向他们指明谁是\Quote{健康的人}，谁是\Quote{病人}。祂说：\BibleQuote{我不是来召义人，}而是\BibleQuote{来召罪人。}如果罪人不来就我，那么我为何而来？我为谁而来？若众人都健康，何以如此伟大的医生从天降下？祂为何为我们预备了药——不是出自祂的仓库，而是出自祂自己的血？因此，那种病情较轻、自觉有病的病人，紧依医生，为得痊愈。但那些病更危险的人，嘲弄医生，辱骂病人。他们的狂乱最终发展到何种地步？抓住医生，捆绑、鞭打、给祂戴上荆棘冠冕、将祂悬于木架、在十字架上杀死祂！你为何惊奇？病人杀死了医生；但医生藉着被杀，治愈了那狂暴的病人。\par

5. 首先，祂在十字架上不忘祂自己的身份，向我们显示祂的忍耐，并给我们爱仇敌的榜样；当祂看见他们围在祂身边狂怒时——祂这医生知道他们的病，知道使他们迷狂的疯狂——祂立刻对父说：\BibleQuote{父啊，宽恕他们，因为他们不知道自己做的是什么。}现在，你们以为那些犹太人不是恶毒、残忍、嗜血、骚乱、敌视天主之子的吗？你们以为那呼声\BibleQuote{父啊，宽恕他们，因为他们不知道自己做的是什么}是无效而徒然的吗？祂看见他们所有人，但祂知道他们中那些终将成为祂的人。简言之，祂死了，因为如此合宜，好借着祂的死来杀死死亡。天主死了，好通过一种天上的交换来实现一种交易，使人不再见到死亡。因为基督是天主，但祂不是在祂是天主的那个性体内死了。因为同一为既是天主又是人；因为天主和人是一个基督。人性被取了，好使我们能变得更好；祂没有将神性贬低到较低的层次。因为祂取了祂所不是的，没有失去祂所是的。因为祂既是天主又是人，乐意我们因祂所有的而活，祂就在我们所是的里面死了。因为祂自己没有什么可使其死；我们也没有什么可使其活。因为那没有什么可使其死的祂是谁？\BibleQuote{在起初已有圣言，圣言与天主同在，圣言就是天主。}如果你在天主内寻找什么可使祂死，你不会找到。但我们这些属于血肉的都要死；人带着有罪的血肉。找出那可使罪活着的东西；它没有。因此，祂不能在自己所有的之内有死亡，我们也不能在自己所有的之内有生命；但我们从祂所有的得到生命，祂从我们所是得到死亡。何等交换！祂给了什么，又收了什么？经商的人为了交换物品而进行商业往来。因为古代商业只是物品的交换。一个人给出他有的，接受他没有的。例如，他有小麦，但没有大麦；另一个人有大麦，但没有小麦；前者给出他有的小麦，接受他没有的大麦。多么简单，用较大的量来补偿较差的种类！于是另一个人给出大麦，来收小麦；最后，另一个人给出铅，来收银子，只是他给出许多铅以换取少量银子；另一个人给出羊毛，来收一件成衣。谁能数尽所有这些交换？但没有人给出生命来接受死亡。因此，医生挂在树上时发出的声音并非徒然。因为为使我们能死，圣言不能死，\BibleQuote{圣言成了血肉，居住在我们中间。}祂挂在十字架上，但在肉体中。有那卑贱，犹太人鄙视它；有那宝贵，犹太人借此得救。因为正是为他们说了\BibleQuote{父啊，宽恕他们，因为他们不知道自己做的是什么。}那声音并非徒然。祂死了，被埋葬，复活了，与门徒一起过了四十天，升了天，祂派遣圣神到等候应许的他们身上。他们充满了所领受的圣神，并开始说万国的言语。那时在场的犹太人惊奇：没有学问无知的人（他们曾知道这些人是在他们中间用一种语言长大的）竟然因基督的名说各种语言，他们惊讶，并从伯多禄的话得知这恩赐从何而来。那挂在树上的祂赐予了。那悬在树上被讥讽的祂赐予了，好从祂在天上的宝座赐予圣神。祂曾为他们说\BibleQuote{父啊，宽恕他们，因为他们不知道自己做的是什么}的那些人听了，信了。他们信了，受了洗，他们的悔改生效了。什么悔改？因信德他们喝了基督的血，就是他们在狂怒中所流的。\par

6. 因此，为了以我们开始时的同一话题来结束这篇讲道，让我们祈祷，让我们依靠天主；让我们按祂的命令生活；当我们在这尘世中摇摇欲坠时，让我们像门徒那样呼求祂，说：\BibleQuote{主，请增加我们的信德。}伯多禄既信赖祂，又动摇；然而他没有被忽视而沉下去，而是被扶起并提起来。因为他的信赖从何而来？不是出于他自己，而是出于主的。如何？\BibleQuote{主，如果是祢，请命令我走到祢那里，行走在水上。}因为主是在水上行走。\BibleQuote{如果是祢，请命令我走到祢那里，行走在水上。}因为我知道，如果是祢，祢一命令，事就成了。\BibleQuote{祂说：来吧。}他按照祂的命令下来，但因自己的软弱而害怕。然而，当他害怕时，他喊说：\BibleQuote{主，救我！}于是主抓住他的手，说：\BibleQuote{小信德的人哪，你为什么怀疑？}祂先召叫他，祂拯救了他，当他动摇、失足之时；好应验圣咏上所说的：\BibleQuote{我若说我的脚滑了一下，上主，祢的仁慈就扶持了我。}\par

7. 福分有两种：暂时的和永恒的。暂时的福分是健康、资财、尊荣、朋友、家园、儿女、妻子，以及我们在此寄居的生命中的其他事物。那么，我们在这生命的客栈中歇宿，如同过路的旅客，而非打算长住的业主。但永恒的福分首先是永生本身，灵魂与身体的不朽与不死，天使的团体，天上的城，永不衰残的光荣，圣父与天乡：前者无死亡，后者无敌对。让我们以极其切望的心渴慕这些福分，以恒久不懈的心祈求它们，不是用冗长的言语，而是以叹息为证。渴望的愿望常祈祷，即使舌头沉默。你若常渴望，你就常祈祷。祈祷何时沉睡？当渴望冷淡时。所以，让我们以全部渴望之情恳求这些永恒的福分，以完全的热诚寻求那些好事，以全然的确信祈求那些好事。因为那些好事有益于拥有它们的人，不能伤害它们。但那些暂时的好事有时有益，有时有害。贫困有益于许多人，财富有害于许多人；卑微的生活有益于许多人，高升的尊荣有害于许多人。再者，金钱有益于有些人，尊贵的地位有益于有些人：有益于那些善用它们的人；但对于那些滥用它们的人，不收回它们反而更伤害了他们。所以，弟兄们，让我们也祈求那些暂时的福分，但要适中，确信如果我们真得到了它们，是那知道什么对我们有益的一位所赐的。你求了，而你所求的没有赐给你？信赖你的圣父，若它对你有益，祂必会赐给你。看！你以自身判断此事吧。因为你的儿子，他不认识人的方法，在你面前如何，你自身对主也是如此，你不认识天主的事。看，你的儿子在你面前哭喊一整天，要你给他一把刀或一把剑；你拒绝给他，你不会给，你不顾他的眼泪，免得你将来哀悼他的死亡。任凭他哭喊，捶打自己，或扑倒在地，要你抱他上马；你也不会做，因为他不知道如何驾驭马匹，可能会摔倒摔死他。你拒绝给他一部分，是为他保全全部。但为使他长大，安全地拥有全部，你不给他那对他充满危险的小东西。\par

8. 所以，弟兄们，我们说：要尽力祈祷。邪恶泛滥，天主愿意邪恶泛滥。但愿邪恶的人不泛滥，那样邪恶就不会泛滥。恶时！多难之时！人们这样说。让我们的生命良善，那么时代就是良善的。我们造就自己的时代；我们如何，时代也如何。但我们能做什么？我们或许无法感化大众归于良善的生活。但让那少数听从的人善度生活；让那少数善度生活的人容忍那许多生活邪恶的人。他们是谷粒，他们在打谷场上；在打谷场上他们能与糠秕同在，但在谷仓里他们不会再有糠秕。让他们忍受所不愿的，为得到所愿的。为何我们忧愁，责怪天主？邪恶在世上泛滥，为使世界不能吸引我们的爱。那些轻视世界及其一切诱惑的人，是伟大的人物、忠信的圣人；我们甚至无法轻视如今已面目全非的世界。世界是邪恶的，看，它是邪恶的，竟仍被爱好像它是良善的。但这邪恶的世界是什么？因为诸天、大地、水泽，以及其中的万物——鱼、鸟、树木——都不是邪恶的。这一切都是良善的；但邪恶的人造就了这邪恶的世界。既然我们不能没有邪恶的人，让我们，如我所说，在活着时，在上主、我们的天主面前倾诉我们的叹息，忍受邪恶，为得到良善的事物。我们不要责怪家主；因为祂对我们慈爱。祂容忍我们，而不是我们容忍祂。祂知道如何管理祂所造的；做祂所吩咐的，盼望祂所应许的。\par

\SourceRecord{Translated by R.G. MacMullen.；Nicene and Post-Nicene Fathers, First Series；Vol. 6.；Edited by Philip Schaff.；Buffalo, NY: Christian Literature Publishing Co.,；1888.；<http://www.newadvent.org/fathers/160330.htm>.}

% structure: p0001 source=h1 target=section reason=page-title

\section{论新约·第三十一篇讲道}

[LXXXI. Ben.]\par

\Emph{论福音的话语}，\BibleRef{玛 18:7}，\Emph{其中劝诫我们提防世上的绊脚石。}\par

1. 我们刚才听到的圣经教训，在宣读之时就警告我们要积蓄美德，坚固基督徒的心，以抵御那预言将要来临的绊脚石，这乃是出于主的仁慈。圣经说：\BibleQuote{人算什么，竟蒙祢眷顾？}主说：\BibleQuote{祸哉，这世界因了绊脚石！}真理如此说；祂警告并提醒我们，不愿我们松懈；当然祂不会使我们绝望。针对这\BibleQuote{祸哉}，即那应受恐惧、战栗并防备的邪恶，圣经在同一处劝勉、激励并教导我们，那里说：\BibleQuote{爱慕祢法律的人，必得大平安，他们没有任何绊脚石。}祂向我们指明了要防备的仇敌，但也没有忘记给我们展示一座防御的墙。当你听到\BibleQuote{祸哉，这世界因了绊脚石}时，你曾想：要躲避绊脚石，你能到世界之外去，以免受其害吗？因此，为了躲避绊脚石，除非你逃向创造世界的那一位，你还能去哪里呢？我们又如何能逃向创造世界的那一位呢？除非我们听从那在各地宣讲的祂的法律。而仅仅听从还不够，除非我们爱慕它。因为圣经为使你免受绊脚石之害，并没有说：\BibleQuote{爱慕祢法律的人，必有大平安}。原来\BibleQuote{听法律的并不在天主面前成义；}但是，因为\BibleQuote{行法律的才要成义，}并且\BibleQuote{信德以爱德行事}，所以圣经说：\BibleQuote{爱慕祢法律的人，必得大平安，他们没有任何绊脚石。}与我们依次吟唱的经文也相合：\BibleQuote{但温良的人要继承大地，并喜爱丰富的和平。}因为\BibleQuote{爱慕祢法律的人，必得大平安。}这些\BibleQuote{温良}的人就是\BibleQuote{爱慕天主法律}的人。因为\BibleQuote{上主，祢所教训并教以法律的人，是有福的，好使他从灾难的日子得享平安，直到为恶人挖了陷阱。}这些经文看似多么不同，却融入并汇合为一个意思，使你从那最丰富的源泉中所听到的一切，只要你在其中安息，并与真理和谐相爱，你立刻充满了平安；因爱德而火热，并因之而防御了绊脚石。\par

2. 那么，我们应当去认识、寻求或学习，如何成为\BibleQuote{温良}的人；藉我刚才从圣经中引用的经文，我们就被引导去寻找我们所追求的。因此，亲爱的诸位，请稍加留意；这是一件重要的事，正在处理，就是我们要成为温良的人；在人生的逆境中是必需的。但人生的逆境并不被称为绊脚石；但请注意什么是\BibleQuote{绊脚石}。例如，一个人因某种严峻的困境而被重压所困。他被重压所困，这并非绊脚石。殉道者们也曾受过这样的压力，但并未被压倒。要提防绊脚石，但不必过分担忧重压。后者压你，绊脚石却压迫你。那么两者有什么区别呢？在重压之下，你准备保持忍耐，持守坚忍，不放弃信德，不屈从于罪恶。你若保持或将要保持这一点，那么压迫你的困苦就不会使你跌倒；相反，那重压会如同油压机一样，不是为了毁灭橄榄，而是为了压出油来。总之，如果在这压迫你的困苦中，你将赞美归于天主，那么这重压对你将何其有用，竟压出如此宝贵的油！宗徒们在这种压力下坐监，并在那种压力下向天主咏唱了赞美诗。这是多么宝贵的油被压榨出来！在沉重的压力下，约伯坐在粪堆上，无依无靠，无人帮助，没有财物，没有子女；只有满身的蛆虫——就外在的人而言——但他内心充满了天主，他赞美天主，那压力对他并非\BibleQuote{绊脚石}。那么\BibleQuote{绊脚石}在哪里呢？当他的妻子来对他说：\BibleQuote{你说一句亵渎天主的话，然后死去吧。}当魔鬼夺走了他所有的一切时，一位厄娃被留给了这位受试炼者，不是为了安慰，而是为了试探她的丈夫。请看绊脚石在哪里。她夸大他的痛苦，也夸大她自己的痛苦与他的一起，并开始劝他亵渎天主。但他，因为\BibleQuote{温良}，因为\BibleQuote{天主教训他，并以法律赐他平安，使他脱离灾难的日子}，他心中\BibleQuote{得享大平安，因他爱慕天主的法律，没有什么绊脚石能绊倒他}。她是绊脚石，但未能绊倒他。总之，请看这位温良的人，请看这位受教于天主法律的人——我指的是永恒的天主法律。因为在约伯时代，那写在石版上的法律尚未赐给犹太人，但在虔诚人心中，永恒的法律仍然存在，而那赐予百姓的法律是据此抄录的。因此，藉天主的法律，他\BibleQuote{从灾难的日子得享平安}，并\BibleQuote{因爱慕天主的法律而享有大平安}，请看他是如何的\BibleQuote{温良}，以及他所回答的。由此你要学习我所要探讨的：谁是温良的人。他说：\BibleQuote{你说话像愚昧的妇人一样。我们从主手里得福，难道不也受苦吗？}\par

3. 我们已经藉着一个例子听到谁是温良的人：如果可能，让我们用言语给他们下定义。温良的人是这样的人：在他们的一切善行中，在他们所做的一切好事中，除了天主，没有什么是令他们喜悦的；在他们所遭受的一切恶事中，天主并未令他们不悦。弟兄们，现在请留意这个规则，这个榜样；让我们向它伸展，让我们寻求增长，好能充满它。因为，除非天主赐给增长，我们栽种、浇灌又有何益？\BibleQuote{可见，栽种的不算什么，浇灌的也不算什么，只在那使它生长的天主。} 请听，无论你是谁，你愿意成为\Quote{温良}的人，你愿意\Quote{从患难之日获得安息，你爱慕天主的法律}，好使\Quote{没有绊脚石到你那里}，并使你\Quote{有大平安}，使你\Quote{承受大地，喜爱和平的丰盛}；请听，无论你是谁，你愿意成为\Quote{温良}的人。无论你做什么善事，不要自满。\BibleQuote{因为天主拒绝骄傲人，却赐恩宠给谦卑的人。} 所以，无论你做什么善事，除了天主以外，不要让任何事使你喜悦；无论你遭受什么恶事，不要让天主使你不悦。你还需要什么？这样做，你必将生存。患难之日不会压倒你，你将逃脱那所说的：\BibleQuote{世界因了恶表是有祸的。} 因为哪个世界因恶表而有祸呢？不就是那曾说的\BibleQuote{世界却不认识祂}的世界吗？而不是那曾说的\BibleQuote{天主在基督内使世界与自己和好}的世界。有一个邪恶的世界，也有一个良善的世界；邪恶的世界是这个世界中所有邪恶的人，良善的世界是这个世界中所有良善的人。正如我们常见一块田地：这块田地满了什么？满了麦子。但我们也会说，而且说得正确，这块田满了糠秕。同样，一棵树，它满了果实。另一个人说，它满了叶子。说它满了果实的人说得对，说它满了叶子的人也说得对。叶子的繁盛并没有夺去果实的空间，果实的丰硕也没有驱走叶子的众多。它两者都满，但叶子被风吹散，果实被农夫收聚。所以，当你听到\BibleQuote{世界因了恶表是有祸的}时，不要害怕；\Quote{爱慕天主的法律，没什么会成为你的绊脚石。}\par

4. 但是你的妻子来到你面前，劝你做某件恶事。你爱她，如同妻子应当被爱的那样；她是你的一个肢体。而你刚才从福音中听到：\BibleQuote{倘若你的眼绊倒你，倘若你的手绊倒你，倘若你的脚绊倒你，砍下来，丢掉它们。} 无论谁是你所亲爱的，无论谁是你所敬重的，只要他还没有开始绊倒你，也就是说，劝你做任何恶事，你就当一直敬重他，一直把他当作你亲爱的肢体。现在请听，这就是\Quote{绊脚石}的意思。我已经举了\Person{约伯}和他妻子的例子；但在那里，\Quote{绊脚石}这个词并未出现。请听福音：当主预言祂的受难时，\Person{伯多禄}开始劝祂不要受苦。\BibleQuote{退到我后面去，\Person{撒殚}，你是我的绊脚石。} 这里无疑地，主给了你生活的榜样，教训你什么是\Quote{绊脚石}，以及如何避免绊脚石。祂刚才还对他说：\BibleQuote{你是有福的，\Person{西满巴尔多纳}，}表明他是祂的肢体。但当他开始成为绊脚石时，祂就砍掉那个肢体；只是祂后来恢复了那肢体，又把它放回原处。那么，凡是开始劝你行恶的人，就会成为你的\Quote{绊脚石}。亲爱的，在这里要注意；这种事情多半不是出于恶意，而是出于一种错误的善意。你的朋友，他爱你，你也爱他，你的父亲、兄弟、孩子、妻子，见你处于困境，就想让你做恶事。我所谓\Quote{见你处于困境}是什么意思？就是见你被某种患难所压迫。这种压迫也许你是因为义德而受的；你之所以受它，是因为你不肯作假见证。我只是举例说明；类似的例子很多，因为\BibleQuote{世界因了恶表是有祸的。} 例如，某个有权势的人，为了掩饰他的掠夺和抢劫，要你作假见证。你拒绝了：拒绝发虚誓，免得你否认那真实的祂。我不必多说，他发怒，他有权势，他压迫你；这时一个朋友来了，他不愿你受这股患难的压迫、处于这种困境，说：\Quote{我求你，照他说的做吧；这有什么大不了呢？} 然后，也许就像撒殚对主所说：\BibleQuote{经上记载：祂为你吩咐了祂的天使，保护你，免得你的脚碰在石头上。} 也许你这个朋友，因为看到你是基督徒，想要凭法律说服你去做他认为你该做的事。\Quote{照那人说的做。} \Quote{什么？照那人所愿的做。} \Quote{但那是谎言，是假的。} \Quote{好吧，难道你没有读过：\InnerQuote{众人都是虚谎的。}？} 现在，他就成了一个\Quote{绊脚石}。他是朋友，你该怎么办？他是眼，他是手：\BibleQuote{砍下来，丢掉。} 什么叫做\Quote{砍下来，丢掉}？就是不要同意他。因为我们身体上的肢体是因同意而构成统一，因同意而生存，因同意而彼此结合。哪里有不和，那里就有疾病或溃烂。他既是你的一个肢体，你会爱他。但他成了你的绊脚石；\BibleQuote{砍掉他，丢掉他。} 不要同意他；把他从你的耳际赶走，或许他会回头改正。\par

5. 你要怎么做我所说的这事：\Quote{把他砍掉，把他从你身上扔掉}，这样一来，或许能改正他？回答我，你打算怎么做？他想在法律之外说服你说谎。因为他说：\Quote{说话。}也许他不敢说：\Quote{说句谎话；}而是说：\Quote{说别人希望听的话。}你说：\Quote{可这是一句谎话。}他为自己开脱说：\Quote{所有人都是说谎者。}那么，我的兄弟，你要反驳说：\Quote{说谎的口杀害灵魂。}注意，你所听见的不是轻描淡写的话：\Quote{说谎的口杀害灵魂。}那个压迫我的强大仇敌能对我做什么，以至于你怜悯我和我的处境，不愿我处于这种邪恶境况，而你却愿意我作恶？那个强大的人能对我做什么？他能压迫什么？肉身。他会压迫你的身体，你会说：我承认他可能压迫它以致毁灭。然而他对待我，比我若说谎对自己还温和得多！他杀我的肉身；我杀我的灵魂。他凭其权势和愤怒杀死身体；\Quote{说谎的口杀害灵魂。}他杀身体，身体即使不被杀也必死；但不义所不能杀的灵魂，真理却永远接纳。那么，保全你能保全的；让那迟早必朽坏的东西朽坏吧。那么，你已经作了回答，但你没有解决\Quote{所有人都是说谎者}的问题。对此也向他回答，免得他认为自己引用法律中的见证就说服了人说谎，从而用法律来反对法律。因为法律上写着：\Quote{不可作假见证；}法律上也写着：\Quote{所有人都是说谎者。}那么，回到我刚刚提示的话，就是我尽量用言语为\Quote{温良}的人下的定义。他是\Quote{温良}的，在他一切善行中，只有天主令他喜悦，在他所受的一切恶事中，天主并不令他不悦。那么，向那说\Quote{说谎吧，因为经上记载：\InnerQuote{所有人都是说谎者}}的人回答：我不说谎，因为经上记载：\Quote{说谎的口杀害灵魂。}我不说谎，因为经上记载：\Quote{你必消灭说谎言的人。}我不说谎，因为经上记载：\Quote{不可作假见证。}即使因真理而不悦的人用压迫折磨我的身体，我也要听从我主的话：\Quote{不要怕那些杀身体的人。}\par

6. \Quote{那么怎么所有人都是说谎者呢？什么！你难道不是人吗？我想是。}快而真实地回答。\Quote{但愿我不是人，这样我就不会是撒谎者。}因为请看：\Quote{天主从天上俯视世人，要看有没有明智的，有没有寻求天主的。他们都离弃了正道，一同成了无用的；没有行善的，连一个也没有。}为什么？因为他们愿作世人。但为了将他们从这些不义中解救出来，医治、治愈、改变这些世人，\Quote{祂赐给他们权能，成为天主的子女。}那么有何奇妙！你们曾是世人，如果我们曾是世人；你们所有人都是人，也都是说谎者，因为\Quote{所有人都是说谎者。}天主的恩宠临到你们，并\Quote{赐给你们权能，成为天主的子女。}听我父的声音说：\Quote{我曾说，你们是神；你们都是至高者的子女。}既然他们是人，是世人，如果他们不是至高者的子女，他们就是说谎者，因为\Quote{所有人都是说谎者。}如果他们是天主的子女，如果他们被救主的恩宠赎回，被祂的宝血买回，由水和圣神重生，预定得天上的产业，那么他们真是天主的子女，因此也就是神。那么，谎言与你们何干？因为亚当只是一个凡人，基督是人而天主；天主是万物的创造者。亚当只是一个凡人，那人基督是天人中保，圣父的独生子，天主而人。看，你，人啊，远离天主，天主远高于人；在两者之间，天主而人将自己置于中间。承认基督，借着祂这位人升到天主那里。\par

7. 既然如今已改过自新，并且，若我的言语蒙受祝福，成了\Quote{温良的}，就让我们\BibleQuote{坚定不移地持守所宣认的信仰。}让我们爱天主的法律，好能逃避那所记载的：\BibleQuote{世界因了绊脚石是有祸的。}现在我要略略谈一谈\Quote{绊脚石}，世界充满了它，并谈谈绊脚石怎样增多，患难怎样层出不穷。世界荒废了，榨酒池被践踏。啊！基督徒，属天的嫩芽，你们是地上的旅客，寻求天上的城邑，渴望与圣天使为伴；要明白你们来到此地只有一个条件，就是应当速速离开。你们正经过世界，努力到达创造它的那一位那里。不要让那些爱世界、愿留于世界、然而不论愿意与否都不得不离开世界的人扰乱你们，欺骗你们，诱惑你们。这些层出不穷的患难并不是绊脚石。你们要成为义人，它们就只是考验。患难来了；它将成为你所选择的那样，或为考验，或为定罪。它发现你是什么，它就是什么。患难是火；它若发现你是金子，它就除去渣滓；它若发现你是糠秕，它就将其烧成灰烬。因此，层出不穷的患难并不是\Quote{绊脚石}。但什么是\Quote{绊脚石}呢？就是那些表达，那些对我们这样说的言语：\Quote{看，基督徒时代带来了什么}；看，这些才是真正的绊脚石。因为这话对你说，为的是这个目的：你若爱世界，就会亵渎基督。说这话的人是你的朋友、你的谋士，所以是你的\Quote{眼}。说这话的人服事你，与你同劳，所以是你的\Quote{手}。说这话的人也许支持你，把你从低微的地上地位抬举起来，所以是你的\Quote{脚}。全都抛开，全都砍断，全都丢弃；不要赞同他们。你要回答这样的人，如同那被劝去作假见证的人所回答的。你也要这样回答；对那向你说\Quote{看，正是在基督徒时代才有这样层出不穷的患难；整个世界都荒废了}的人，你要回答他：\Quote{这事在发生之前，基督已预先告诉了我。}\par

8. 那么你为何烦乱？你的心灵因世界的层出不穷的患难而烦乱，正如那船一样，基督在里面睡着了。啊！刚毅的人，你的心灵为何烦乱？那基督睡在里面的船，就是信德睡在里面的心灵。因为有什么新事，有什么新事，我问你，基督徒，告诉了你呢？\Quote{在基督徒时代，世界荒废了，世界正在衰亡。}难道你的主没有告诉你，世界将要荒废吗？难道你的主没有告诉你，世界将要衰亡吗？为什么当初应许之时你相信，如今应验之时你却烦乱呢？因此狂风猛烈地冲击你的心灵；当心船破，唤醒基督吧。宗徒说：\BibleQuote{为使基督因着信德住在你们心中。}基督因信德住在你们内。信德临在，就是基督临在；信德清醒，就是基督清醒；信德昏睡，就是基督昏睡。起来吧，振作自己；说：\BibleQuote{主啊，我们丧亡了。}请看外教人对我们说什么；更糟的是，恶毒的基督徒说什么！醒起来吧，主啊，我们丧亡了。让你的信德醒起来，基督就开始对你说话。\Quote{你\InnerQuote{为什么烦乱？}这一切事我早已预先告诉了你。我预言了它们，为的是当邪恶来临时，你仍能期望美善，而不在邪恶中昏厥。}你奇怪世界正在衰亡吗？奇怪世界已经衰老了。它如同一个人，出生，长大，变老。老年有许多痛苦：咳嗽、流涕、眼花、烦躁和疲惫。所以，正如一个人老了，满是痛苦；世界也老了，满是患难。天主为你所做的难道是小事吗？在世界的暮年，祂派遣基督到你这里来，好使祂在万有衰败之时更新你。你难道不知道祂在亚巴郎的后裔中预示了这事？宗徒说：\BibleQuote{\InnerQuote{亚巴郎的后裔}就是基督。祂并不是说\InnerQuote{后裔们}，如同许多；而是单指一位：\InnerQuote{你的后裔}，就是基督。}因此，亚巴郎老年时生了一个儿子，因为在这世界的暮年，基督将要来临。祂来的时候，万有都已衰老，祂使它们更新。世界作为一个受造、生成、必朽之物，正在趋于败亡。它必然充满患难；祂来，既要安慰你在当前的患难中，又要应许你永远的安息。那么，不要选择依附这个衰老的世界，而不愿在基督内恢复青春，祂告诉你：世界正在消逝，世界正在变老，世界正在衰亡；因老年的沉重喘息而痛苦。但不要害怕，\BibleQuote{你的青春将如鹰一般被更新。}\par

9. 看，他们说，在基督徒的时代，罗马正在毁灭。也许罗马并非正在毁灭；也许她只是在受鞭打，并未完全摧毁；也许她是在受惩戒，而非化为乌有。可能如此；罗马不会毁灭，如果罗马人不毁灭。他们也不会毁灭，如果他们赞美天主；如果他们亵渎祂，他们就会毁灭。因为什么是罗马？不就是罗马人吗？因为问题不在于她的木石，不在于她高耸的宫殿和她所有宽阔的城墙。这一切的建造，只是建立在有朝一日会倒塌的条件上。当人建造它时，他把石头叠在石头上；当人毁灭它时，他把石头一块块拆掉。人建造了它，人毁灭了它。罗马受到了什么伤害吗，因为有人说：\Quote{她正在倒塌}？不，不是罗马，也许是她的建造者。那么，我们是否伤害了她的建造者，因为我们说罗马正在倒塌，而罗马是\Person{罗慕路斯}建造的？这个世界本身也将被火烧毁，而这是天主建造的。但人的工程不会倒塌，除非天主愿意；天主的工程也不会倒塌，除非祂愿意。因为如果人的工程没有天主的旨意就不会倒塌，那么天主的工程怎能因人的意愿而倒塌呢？然而，天主创造了这个世界，它终有一日会为你而倒塌；因此祂也创造了你，作为一个终有一死的人。人本身，城市的装饰，人本身，城市的居民、统治者、治理者，以此条件而来——他终将离去，以此条件而生——他终将死亡，以此条件进入世界——他终将消逝；\Quote{天地要过去}，那么，如果一座城市某时有个尽头，有什么奇怪呢？然而也许城市的尽头现在还没来到；但终有一天会来到。但为什么罗马在基督徒的祭献中毁灭？为什么她的母城\Place{特洛伊}在外教人的祭献中被焚毁？罗马人寄托全部希望的神祇，是的，外教罗马人寄托希望的神祇，从特洛伊的火焰中解救出来，去建立罗马。这些罗马的神祇原本就是特洛伊的神祇。特洛伊被焚毁，\Person{埃涅阿斯}带走了逃亡的神祇；或者更确切地说，他本人是一个逃亡者，带走了这些无感觉的神祇。因为他们能被逃亡者携带，但他们自己不能逃跑。他带着这些神祇进入意大利，带着这些假神，建立了罗马。要详述整个故事太长了；但我愿意简短地提及他们自己的著作所包含的内容。他们的一位广为人知的作者这样说：\Quote{据我所知，特洛伊人在埃涅阿斯的带领下，像逃亡者一样漂泊，没有固定的居所，最初建造并居住在罗马。}所以，他们带着他们的神祇，在\Place{拉丁姆}建立了罗马，并在那里安置了他们先前在特洛伊崇拜的神祇。他们的诗人描写\Person{尤诺}，对着埃涅阿斯和逃亡的特洛伊人发怒，说道：\par

\Quote{我被厌恶的一群流浪奴隶，\newline{}乘着顺风驶过第勒尼安海，\newline{}驶向富饶的意大利，\newline{}为他们战败的神祇，设计新的庙宇\newline{}在那里。}\par

如今，当这些战败的神祇被带到意大利时，是作为保护神，还是作为他们未来倒下的预兆？\Quote{因此，你们要爱慕天主的法律，凡事就不会使你们跌倒。}我们恳求你们，我们哀求你们，我们劝勉你们；要温良，同情受苦的人，扶持软弱的人；借此众多陌生人、穷困者和受苦者聚集的机会，愿你们的款待和善行更加丰盛。只要基督徒做基督所吩咐的，那么外教人只会徒然诋毁而自受其害。\par

\SourceRecord{Translated by R.G. MacMullen.；Nicene and Post-Nicene Fathers, First Series；Vol. 6.；Edited by Philip Schaff.；Buffalo, NY: Christian Literature Publishing Co.,；1888.；<http://www.newadvent.org/fathers/160331.htm>.}

% structure: p0001 source=h1 target=section reason=page-title

\section{新约讲道第三十二篇}

\Emph{[LXXXII. Ben.]}\par

\Emph{论福音的话}，\BibleRef{玛 18:15} \Emph{，\BibleQuote{如果你的弟兄得罪了你，去，趁只有你和他在一起的时候，指出他的过错。}以及撒罗满的话：用眼睛行诡诈的，使人愁苦；但公开责备的，带来和平。}\par

一、我们的主警告我们不要忽略彼此的罪，不是要去找可以指责的事，而是要看有什么可以改进的。因为祂说，那在自己眼中有梁木的人，怎能看出弟兄眼中的木屑呢？现在，亲爱的，我将简要地告诉你们这是什么意思。眼中的木屑是愤怒；眼中的梁木是憎恨。因此，当怀恨的人指责愤怒的人时，他想要取出弟兄眼中的木屑，却被自己眼中的梁木所阻碍。木屑是梁木的开始。因为梁木在成长过程中，首先是一根木屑。浇水使木屑长成梁木；用恶意的猜疑滋养愤怒，便使之变成憎恨。\par

二、愤怒者的罪与怀恨者的残忍之间有巨大的差别。因为我们甚至会对我们的孩子发怒；但谁曾见过恨自己的孩子呢？甚至在牲畜中，母牛有时会因某种厌烦而在愤怒中赶走正在吮乳的牛犊；但不久它就怀着母亲的一切柔情拥抱它。当它用角顶撞时，似乎厌恶它；但当它失去它时，就会寻找它。我们管教孩子，也无不是带着某种愤怒和义愤；但我们管教他们，完全是出于对他们的爱。因此，每个发怒的人离憎恨如此之远；以至于有时人若是不发怒，反而更会被定为憎恨。假设一个孩子想在某条河流中玩耍，而水流的力量会让他丧命；如果你看到这事，并忍耐地任凭他去，这就是憎恨；你忍耐地任凭他，便是他的死亡。发怒并纠正他，远胜于因不发怒而任凭他灭亡！因此，最要紧的是避免憎恨，并拔除眼中的梁木。确实，一个人在愤怒中说话越过了应有的限度，后来通过悔恨而洗刷掉，这与在心里隐藏着险恶的意图，两者之间有很大的差别。最后，这两句圣经的话之间也有大的差别：\BibleQuote{我的眼睛因愤怒而昏乱。} 而另一句则说：\BibleQuote{凡恨他弟兄的，就是杀人犯。} 眼睛昏乱与完全失明，有很大的差别。木屑使人昏乱，梁木使人完全失明。\par

三、那么，为了我们能妥善地实行并完成今天所告诫我们的事，首先让我们说服自己：最重要的是没有憎恨。因为当你自己的眼中没有梁木时，你就能正确地看见你弟兄眼中有什么东西，并且感到不安，直到你从他眼中取出你所看见的伤害他的东西。你里面的光不允许你忽略你弟兄的光。然而，如果你怀恨在心，却要纠正他，当你自己失去了光时，你如何改善他的光呢？因为同一圣经记载说：\BibleQuote{凡恨他弟兄的，就是杀人犯。} 也明确地告诉我们这一点：\BibleQuote{凡恨他弟兄的，至今仍在黑暗中。} 因此，憎恨就是黑暗。现在，恨别人的人，必先伤害自己，这是必然的。因为他试图在外表上伤害别人，却在内心毁灭自己。由于我们的灵魂比身体更珍贵，所以我们更应当照顾灵魂，不让它受到伤害。但恨别人的人，正伤害自己的灵魂。他对所恨的人能做什么呢？他能做什么呢？他拿走他的钱财，难道能拿走他的信德吗？他伤害他的名声，难道能伤害他的良心吗？他造成的任何伤害，都只是外在的；现在看看他给自己的伤害是什么？因为恨别人的人，内心是自己的敌人。但由于他感觉不到自己在伤害自己，所以他对别人施暴，而且更加危险，因为他感觉不到自己在作恶；因为这种暴力本身使他失去了感知的能力。你对你的敌人施暴；通过你的暴力，他被掠夺，而你将陷入邪恶。两者之间有很大的差别。他失去了钱财，你失去了纯洁。请问哪一个损失更重？他失去了一件必定会朽坏的东西，而你却成了必须自己灭亡的人。\par

4. 因此，我们应当怀着爱心去责备；不是出于任何伤害人的热望，而是出于真诚渴望改正的关怀。若我们如此存心，便能最卓越地实践今天所劝勉我们的：\BibleQuote{如果你的弟兄得罪了你，你要在独处时，只在你和他之间责备他。}为什么你要责备他？因为你因他得罪你而忧伤吗？决不是。若你出于自爱而这样做，你便一无所成。若你出于对他的爱而这样做，你便做得极好。事实上，请在这些话中观察，你究竟应为爱谁而做这事，是为自己还是为他。\BibleQuote{如果他听你，你便赚得了你的弟兄。}那么，要为他而行，好使你\Quote{赚得}他。若你这样做\Quote{赚得}了他，那么你若不做，他早已失丧了。既然如此，许多人为何轻视这些罪过，说：\Quote{我做了什么大事？我只是得罪了人而已。}切勿轻视。你得罪了人，但你要知道，得罪人时，你正失丧自己。若那被得罪的人\BibleQuote{在独处时，只在你和他之间责备了你}，而你听从了他，他便\Quote{赚得}了你。\Quote{赚得了你}是什么意思？岂不是说，若他不赚得你，你早已失丧了？因为若你本不会失丧，他又怎能赚得你呢？因此，人不可轻视得罪弟兄的事。因为宗徒在某处说：\Quote{但你们这样得罪弟兄，并伤及他们软弱的良心，便是得罪基督了。}这是因为我们全都成了基督的肢体。你怎能不因得罪基督的肢体而得罪基督呢？\newline{}\par

5. 因此，谁也不要说：\Quote{我没有得罪天主，只得罪了弟兄。我得罪了人，这是小罪，或根本不算罪。}或许你说这是小罪，因为它很快就能治愈。你得罪了弟兄：给他赔补，你就痊愈了。你做了件致命的事，但很快也找到了补救。我的弟兄们，我们中谁能指望天国呢？当福音说：\BibleQuote{凡向自己的弟兄说\InnerQuote{傻子}的，就要受地狱的火刑。}这是极大的恐惧！但请看同一处经文的补救：\BibleQuote{若你献礼物在祭台上，在那里想起你的弟兄有什么怨你的事，就把礼物留在祭台前。}天主并不因为你迟延献礼而发怒。天主更看重的是你本人，而非你的礼物。因为若你带着礼物来到天主面前，却对弟兄怀着恶意，祂必回答你：\Quote{你已失丧，你带给我的是什么？你带了礼物，而你自己却不成全给天主的礼物。基督更珍视祂以自己的血所救赎的人，而非你从仓库里找到的东西。}故此，\BibleQuote{把你的礼物留在祭台前，先去与你的弟兄和好，然后再来献你的礼物。}看，那\Quote{地狱的火刑}多么迅速地化解了！当你尚未和好时，你\Quote{处于地狱火刑的危险中}；一旦和好，你便可安然在祭台前献礼。\newline{}\par

6. 然而，人加害于人很容易，寻求和好却很困难。有人说，要向你所得罪、所伤害的人请求宽恕。他回答：\Quote{我不愿这样卑躬屈节。}但如今你若轻视你的弟兄，至少听从你的天主吧。\BibleQuote{凡自谦的，必被高举。}你已跌倒，却还不愿自卑吗？自卑者与伏地者之间大有区别。你已伏在地上，还不愿自卑吗？你大可以说：\Quote{我不愿下降}——若你起初就不愿跌倒。\newline{}\par

7. 那么，行不义的人应当这样做。而遭受不义的人又当如何做呢？我们今日所听到的：\BibleQuote{如果你的弟兄得罪了你，你要单独地规劝他。}如果你忽略此事，你就比他更坏。他行了不义，并且因行不义，给自己造成了重伤；你岂能不顾你弟兄的伤口？你岂能看着他丧亡，或已经丧亡，而不管他的情况？你保持沉默比他的辱骂更坏。因此，当有人得罪我们时，我们应当倍加谨慎，不是为了我们自己，因为忘记伤害是光荣的事；只忘记你自己的伤害，而不要忘记你弟兄的伤口。因此，\BibleQuote{你要单独地规劝他}，目的在于使他改正，但顾及他的羞耻。因为可能由于羞愧，他会开始辩护自己的罪，这样你会使你想要改正的人变得更糟。所以，\BibleQuote{规劝他，只在你们两人之间。如果他听从你，你便赚得了你的弟兄}，因为如果你没有这样做，他本会丧亡。但如果\BibleQuote{他不听从你}，也就是说，如果他辩护自己的罪，仿佛那是正义的行为，\BibleQuote{你就带一两个人同去，以便凭两三个证人的口，句句都可定准；如果他仍不听从他们，你就告诉教会；但如果他也不听从教会，你就把他看作外邦人和税吏。}不再将他列在弟兄之数中。然而也绝不可因此忽略他的救恩。因为即使外邦人，即异教徒，我们也不列在弟兄之中；但我们仍不断寻求他们的救恩。这就是我们所听到的主如此劝诫，并如此谨慎地命令的，以致他立刻附加了这句：\BibleQuote{我实在告诉你们：凡你们在地上所捆绑的，在天上也要被捆绑；凡你们在地上所释放的，在天上也要被释放。}你已开始把你的弟兄看作税吏；\BibleQuote{你在地上捆绑了他}；但要当心，你捆绑得公正。因为不义的捆绑，正义会使之破裂。但当你已改正，并与你的弟兄和好时，\BibleQuote{你在地上释放了他}。而\BibleQuote{当你在地上释放了他，他在天上也要被释放。}因此，你行了一件大事，不是为你自己，而是为他；因为他所行的大不义，不是对你，而是对他自己。\par

8. 但既然是这样，那么\Person{撒罗满}所说的——我们今日从另一篇读经中首先听到的——是什么呢？\BibleQuote{以眼示意奸诈的人，使忧愁加增于人；但公开规劝的，缔造和平。}那么，如果\BibleQuote{公开规劝的，缔造和平}，那么\BibleQuote{你要单独地规劝他}又怎样呢？我们必须担心，以免天主的诫命彼此矛盾。但绝非如此：我们应明白它们之间有最完美的和谐。我们不要追随某些虚妄之人的幻想，他们错误地认为旧约和新约两卷书中的两个圣约彼此对立；以至于我们会认为这里有任何矛盾，因为一处出自\Person{撒罗满}的书，另一处出自福音书。因为如果有任何不精通并诋毁圣经的人说：\Quote{看，这两部圣约彼此矛盾。主说：\InnerQuote{你要单独地规劝他。}\Person{撒罗满}说：\InnerQuote{公开规劝的缔造和平。}}难道主不知道他所命令的吗？\Person{撒罗满}要使罪人的硬脑门受伤；基督则顾及那为自己的罪而羞愧者的羞耻。因为一处写着：\BibleQuote{公开规劝的缔造和平}；但在另一处：\BibleQuote{你要单独地规劝他}，不是\BibleQuote{公开的}，而是私下且隐秘的。但无论你是谁，你若这样想，你要知道这两部圣约并非彼此对立，因为前一段经文见于\Person{撒罗满}的书，后一段见于福音书吗？请听宗徒的话。当然，宗徒是新约的仆人。那么请听宗徒\Person{保禄}训诫\Person{弟茂德}时所说的话：\BibleQuote{那些犯罪的人，你要当众规劝，使其余的人也有所惧怕。}这样看来，似乎不是\Person{撒罗满}的书，而是宗徒\Person{保禄}的书信与福音书相冲突。那么，让我们暂时不凭私见放下\Person{撒罗满}；让我们听主\Person{基督}和他的仆人\Person{保禄}的话。主啊，祢说什么？\BibleQuote{如果你的弟兄得罪了你，你要单独地规劝他。}宗徒啊，你说什么？\BibleQuote{那些犯罪的人，你要当众规劝，使其余的人也有所惧怕。}我们在做什么？我们像法官一样听取这场争论吗？决不可。相反，作为在审判者之下的人，让我们敲门，好使我们得到开启；让我们飞到我主天主的翼下。因为他并没有与他的宗徒说话矛盾，因为他亲自也\Quote{在}他内说话，正如他所说：\BibleQuote{你们既然寻求基督在我内说话的凭证吗？}基督在福音书中，基督在宗徒内：因此基督两者都说了；一个亲口所说，另一个借他传令官的口所说。因为当传令官从法庭上宣布任何事时，记录上不会写\Quote{传令官说了它}，而是写下令传令官说话的那一位说了它。\par

9. 那么，弟兄们，让我们如此听从这两条诫命，好使我们能理解它们，并在二者之间的平安中安定自己。只要我们与自己的心一致，圣经就丝毫不会自相矛盾。两条诫命全是真实的；但我们必须加以区分，有时当行其一，有时当行另一；有时应当\Quote{在单独和他之间规劝兄弟}，有时应当\Quote{在众人面前规劝兄弟，以使其他人也有所畏惧}。我们若有时行这一样，有时行那一样，就能持守圣经的和谐，并在成全和服从它们时不会出错。但有人会对我说：\Quote{我何时该行这一样，何时该行那一样？免得我在应当\InnerQuote{在众人面前规劝}时，却\InnerQuote{在单独和他之间规劝}；或是在应当暗中规劝时，却在众人面前规劝。}\par

10. 亲爱的诸位，你们很快会明白，我们应当做什么，以及何时做；只愿我们不要懈怠去实践它。请留意听着：\Quote{如果你的兄弟得罪了你，要在单独和他之间规劝他。}为什么呢？因为他是得罪了你。什么是\Quote{得罪了你}？你晓得他犯了罪。因为他得罪你时是隐密的，你改正他的罪时也要寻求隐密。因为若是只有你知道他得罪了你，而你却要\Quote{在众人面前规劝他}，你便不是规劝者，而是告密者。请思考那位\Quote{义人}若瑟是如何以极大的仁慈对待他的妻子，尽管他怀疑她犯了那么大的罪，在他知道她因谁怀孕之前；因为他察觉她怀了孕，而他知道自己并没有与她同房。那么，剩下一种不可避免的奸淫嫌疑，但只因他独自察觉，独自知道，福音论到他时说了什么？\Quote{她的丈夫若瑟是个义人，不愿公开羞辱她。}丈夫的悲伤并未寻求报复；他愿意使罪人得益，而非惩罚。\Quote{不愿公开羞辱她，便有意暗暗地休退她。}正当他思虑这事时，\Quote{看，上主的天使在梦中显现给他}，告诉他事情的原委：她没有玷污丈夫的床，而是因圣神怀了孕，这圣神是他们二人的主。那么，你的兄弟得罪了你；若只有你知道，那么他确实是单独得罪了你。因为如果他在许多人面前伤害了你，他也得罪了那些人，就是那些他使之成为其恶行之见证的人。因为我告诉你们，我极亲爱的弟兄们，你们能在自己身上认出这一点。当任何人在我面前伤害我的兄弟时，愿天主禁止我认为那伤害与我无关。他确实也伤害了我；甚至更伤害了我，因为他以为我所做的是讨我喜悦的。因此，那些在众人面前犯的罪，要在众人面前规劝；那些更隐密犯的罪，要更隐密地规劝。区分时间，圣经就自相和谐了。\par

11. 让我们如此行事；而且我们不仅当那罪是得罪我们自身时这样行，当任何人所犯的罪不为他人所知时，也当如此行。我们应当在暗中规劝，在暗中责备他；免得我们若要公开责备他，反而出卖了那人。\Emph{我们}愿意规劝并使他改过；但如果他的仇敌正窥探着要听到什么以便惩罚他呢？例如，一位主教知道有人杀了另一个人，而除此之外无人知晓。我愿意公开责备他；但你却要追诉他。那么，我决不会出卖他，也不会忽视他；我将在暗中规劝他；我将天主的审判置于他眼前；我将警醒他那沾满血污的良心；我将劝他悔改。我们应当以这种爱德来装备自己。因此，人们有时指责我们，好像我们不规劝；或者他们以为我们知道我们所不知道的事，或者以为我们隐瞒我们所知道的事。而你所知道的，我或许也知道，但我不愿在你面前规劝，因为我愿意医治，而不是做告发者。有些人在自己家中犯奸淫，他们暗中犯罪，有时通过他们自己的妻子——通常是出于嫉妒，有时是寻求丈夫的得救——而被我们发现；在这种情况下，我们不公开出卖他们，而是暗中规劝他们。恶事发生在哪里，就让恶事在那里消亡。然而，我们并不忽略那个创伤；最重要的是向那处于如此罪恶状态、背负如此受伤良心的人表明，这是一个致命的创伤，而遭受它的人有时因难以理解的乖僻而轻视它，并为自己寻求支持——不知从何而来，但肯定是无效和虚妄的——说：\Quote{天主不关心肉身的罪。}那么，我们今天所听见的\Quote{淫乱和奸淫的人，天主必要审判}又在哪里呢？看！无论你是谁，如果你受此疾病之苦，请听。听天主所说的话；而不是听你的心出于放纵你自己的罪所说的话，也不是听你的朋友——实际上他是你的仇敌，也是他自己的仇敌，被同一罪恶的锁链捆绑——所说的话。那么，请听宗徒所说的话：\Quote{婚姻在众人中当受尊重，床第也是无玷的；但淫乱和奸淫的人，天主必要审判。}\par

12. 来罢，弟兄，你应当改过。你害怕仇人控告你，难道不怕天主审判你么？你的信德在哪里？当惧怕的时候，就该惧怕。审判的日子确实遥远，但每个人的末日不会遥远；因为生命短暂。既然这短暂不可知，你不知你的末日何时来到。你今天就要改过，为了明天。现在秘密的责备对你有益。因为我公开说话，却是秘密地责备。我敲着所有人的耳朵，但触及某些人的良心。若我说：\Quote{你这个奸夫，改过自新罢；}或许首先我可能说了我所不知道的，或许是轻率听信谣言而猜测。所以我不说：\Quote{你这个奸夫，改过自新罢；}而是：\Quote{你们中无论谁是奸夫，改过自新罢。}这样责备是公开的，改过是隐秘的。我知道，凡惧怕的人，必要改过自新。\par

13. 谁也不要在心里说：\Quote{天主不关心肉身的罪。} \Quote{你们不知道，}宗徒说，\BibleQuote{你们是天主的殿，天主圣神住在你们内吗？若有人毁坏天主的殿，天主必要毁坏他。} \Quote{谁也不可自欺。}但或许有人说：\Quote{我的灵魂是天主的殿，不是我的身体，}并附加这段经文：\BibleQuote{凡有血肉的尽都如草，凡人的一切光荣都如草上的花。}不幸的解释！该受惩罚的狂妄！血肉被称为草，因为它会死；但要当心，那暂时死去的不带着罪过复活。你愿意也在这点上辨别明确的判断吗？\Quote{你们不知道，}同一位宗徒说，\BibleQuote{你们的身体是在你们内的圣神的殿，这圣神是你们从天主得来的？}所以不要再轻忽肉身的罪，因为你们的\BibleQuote{身体是在你们内的圣神的殿，这圣神是你们从天主得来的。}若你轻忽肉身的罪，你岂能轻忽你得罪圣殿的罪呢？你的身体就是在你内之天主圣神的殿。现在你要注意你如何对待天主的殿。若你选择在这墙壁内的圣堂中行奸淫，还有什么比这更邪恶？但现在你自身就是天主的殿。在你出、在你入，在你居家中，在你起身时，你无时无刻不是一座殿。所以小心你所行的，当心不要冒犯殿的居住者，免得他离弃你，使你倒塌。\Quote{你们不知道，}他说，\BibleQuote{你们的身体是在你们内的圣神的殿，这圣神是你们从天主得来的，你们不是属于自己？}因为\BibleQuote{你们是用重价买来的。}若你这样轻看自己的身体，至少想一想你的代价。\par

14. 我知道，但凡稍加思虑的人都知道，但凡敬畏天主的人，没有不因他的话而改过自新的，除非他以为自己还有更长的寿数。这就是杀死许多人的原因，他们总说：\Quote{明天，明天；}突然门就关了。他留在外面，像乌鸦一样呱呱叫，因为他没有鸽子的哀鸣。\Quote{明天，明天}是乌鸦的叫声。你要像鸽子那样哀鸣，捶胸顿足；但当你捶打胸膛时，要因这捶打而变得更好，免得你似乎不是捶打良心，反而用捶打使它变硬，使恶良心更加顽固而非改善。不要徒然哀鸣。因你可能对自己说：天主已应许我赦免，只要我改过，我就能安全；我读圣经：\BibleQuote{当恶人转离他的恶行，遵行法律和正义时，我将不再记念他的一切过犯。}所以我安全了，只要我改过，天主必要赦免我的恶行。对此我能说什么？我岂能高声反对天主？我岂能对天主说：不要赦免他？我岂能说：这没有记载，天主没有这样应许？若我说这样的话，便是说谎。你说得好、说得对；天主应许在你改过后赦免，我不能否认；但请告诉我：看，我同意、我允许、我承认天主应许了赦免，但谁应许了你一个明天？你读到的地方说我若改过必得赦免，在那里也读到我能活多久？你承认：\Quote{我在那里读不到。}那么你不知道你能活多久。改过吧，因此常做准备。不要怕末日，如同贼在你睡觉时挖窟窿，但要醒起，今天就改过。为何拖延到明天？若你的生命很长，就让它又长又好。没有人因为一顿好饭时间长就推迟它；而你却希望过又长又恶的生活？当然，若它是长的，若它好岂不更好？若它是短的，它的好尽可能延长岂不善？但人们如此忽略生命，以至于除了生命本身，他们不愿有任何不好的东西。你买田地，就要找好的；娶妻，要选好的；生子，你盼望好的；买鞋，你不愿意要坏的；然而你爱过坏的生活。你的生活怎么得罪了你，使你只愿它坏，使你拥有一切好东西，唯独你自己是坏的？\par

15. 所以，各位弟兄，倘若我愿私下逐一责备你们中的某人，或许他会听从我。如今我在公开场合责备你们许多人；众人都赞美我；但愿有人留心听我！我并非爱那以口舌赞美我、心中却轻视我的人。因为当你赞美，却不改正自己时，你便是反对自己的见证。如果你邪恶，却对我所言感到喜悦，就当对自己不悦；因你若因自己的邪恶而自省，待你改正时，便会对自己十分满意。若我没记错，这是前天我说过的。在我所有的言语中，我将一面镜子摆在你们面前。这些并非我自己的话，而是我奉主的命令说话；因祂的威严，我不敢缄默不言。谁不愿选择缄默，而不必为你们交账呢？然而如今我已担负这重担，我不能、也不该将它从肩上卸下。各位弟兄，当诵读\BibleBook{}{希伯来书}时，你们听见了：\BibleQuote{你们要服从那些带领你们的人，且要顺服，因为他们为你们的灵魂警醒，如同将要交账的人，使他们能怀着喜乐、而非忧苦去做这事，因为这对你们并无益处。}我们何时怀着喜乐做这事？当我们看见人在天主的言语上长进的时候。田间的工人何时怀着喜乐劳作？当他看见树木，见其果实；当他观看庄稼，见场上谷物丰沛的远景；当他看见自己没有白费劳力，没有弯腰或伤手，没有徒然忍受寒暑的时候。这就是他所说的：\BibleQuote{使他们能怀着喜乐、而非忧苦去做这事；因为这对你们并无益处。}他是否说\Quote{对他们无益}？不。他说\Quote{对你们无益}。因为当那些管辖你们的人因你们的恶行而忧伤时，这对他们有益；他们的忧伤本身对他们有益；但对你们却无益。但我们不愿任何对我们有益的事，却对你们无益。所以，各位弟兄，让我们一同在主的田地里行善，好叫我们在赏报时能一同喜乐。\par

\SourceRecord{Translated by R.G. MacMullen.；Nicene and Post-Nicene Fathers, First Series；Vol. 6.；Edited by Philip Schaff.；Buffalo, NY: Christian Literature Publishing Co.,；1888.；<http://www.newadvent.org/fathers/160332.htm>.}

% structure: p0001 source=h1 target=section reason=page-title

\section{新约讲道集 第33篇}

\Emph{[LXXXIII. Ben.]}\par

\Emph{论福音之言}\BibleRef{玛 17:21}\Emph{，\Quote{我的兄弟得罪我，该多少次……}等等。}\par

1. 昨日的圣福音劝告我们不要忽视弟兄的罪过：\Quote{但若你的弟兄得罪了你，去，要在你和他独处的时候，规劝他；如果他听从了你，你便赚得了你的弟兄；但他如果不听，你就另带一两个人，为叫任何事，凭两个或三个见证人的口供，得以成立。若是他仍不听从他们，你要告诉教会；如果他连教会也不听从，你就将他看作外邦人或税吏。}今日接着读的段落，涉及同一主题。因为主耶稣向伯多禄说了这话后，伯多禄便问他的师傅，弟兄得罪自己，该宽恕他多少次；他问七次是否足够。\Quote{主回答他说：\InnerQuote{不仅七次，而是七十个七次。}}接着祂讲了一个非常可怕的比喻：\Quote{天国好比一个家主，同他的仆人算账；其中有一个欠他一万塔冷通。当主人命将他本人、他的一切和全家都卖掉来还债时，他俯伏在主人脚下，恳求延期，并获得了全部豁免。因为我们听到：\InnerQuote{主人动了怜悯之心，就赦免了他全部的债。}于是那免了债的人，却成了不义的奴隶，他从主人面前出来，遇见一个欠他债的同僚，那同僚欠他，不是像他被宽免的一万塔冷通，而是一百德纳；他就抓住那同僚的脖子，说：\InnerQuote{还清你欠的债！}那同僚像他求主人那样求他，但他找着的同僚，并不像他找着的那样主人。他不仅不宽免那同僚的债，连延期也不答应。他刚刚被主人免了债，却用暴力催促那同僚还债。其他同僚很不高兴，\InnerQuote{就去把所行的事报告了他们的主人}；主人就把那仆人召来，对他说：\InnerQuote{恶仆！你欠我这么大一笔债，我怜悯了你，全部赦免了。难道你不该怜悯你的同僚，如同我怜悯了你一样吗？}于是主人命令将所赦免的一切全部还清。}\par

2. 因此，祂讲这个比喻是为了教导我们，通过这个警告，祂要救我们免于丧亡。祂说：\Quote{如果你们每人不是从心里宽恕自己的弟兄，我的天上的父也必要这样对待你们。}看，弟兄们，这事很清楚，劝告有益，必须完全听从，好使所吩咐的得以实行。因为每个人都是天主的债务人，同时也有某个弟兄对自己负债。因为谁不是天主的债务人呢？除非他找不到任何罪过。谁没有一个弟兄对自己负债呢？除非没有人得罪过他。你以为在人类中能找到一个人，他自己不受某种罪的束缚，同时也有他自己的债务人吗？所以每个人都是债务人，但也有他自己的债务人。因此，公义的天主为你对待你的债务人定下一个规矩，祂自己也会遵守。因为有两样仁慈的工作能解救我们，主自己在福音中简短地阐述了：\Quote{宽恕，你们也必被宽恕；给予，你们也必被给予。}\Quote{宽恕，你们也必被宽恕}涉及赦免。\Quote{给予，你们也必被给予}涉及行善。关于祂所说的赦免，你既希望自己的罪被赦免，你也有另一个你可以赦免的人。同样，关于行善：乞丐向你求乞，而你是天主的乞丐。因为我们祈祷时都是天主的乞丐；我们站在大宅主的门前，甚至还俯伏在地，叹息哀求，希望得到一些东西；而这东西就是天主自己。乞丐向你求什么？饼。你向天主求什么？基督，祂说：\Quote{我是从天上降下的生命之粮。}你愿意被宽恕吗？宽恕吧。\Quote{宽恕，你们也必被宽恕。}你愿意领受吗？\Quote{给予，你们也必被给予。}\par

3. 但现在请听，在如此明确的诫命中，我可能会引起一个难题。关于宽恕的问题，当请求宽恕时，它应当由应宽恕者给予，对我们来说可能是一个难题，正如对伯多禄一样。\BibleQuote{我该宽恕多少次？七次够吗？}\BibleQuote{不够，}主说，\BibleQuote{我不对你说：直到七次，而是直到七十个七次。}现在请计算你的弟兄得罪了你多少次。如果你能数到第七十八次，从而超出了七十个七次，那么就去报复吧。难道主真是这个意思吗？果真如此，如果他犯了\BibleQuote{七十个七次}，你就应该宽恕他；但如果他犯了七十八次，你就可以不宽恕了吗？不，我敢说，即使他犯了七十八次，你也必须宽恕。是的，如我所说，如果他犯了七十八次，宽恕他。如果他犯了一百次，也宽恕。我何必说如此这般多少次呢？一言以蔽之，无论他犯多少次罪，都宽恕他。那么，我是否僭越了我主的尺度？祂将宽恕的限度定在七十七之数；我岂敢逾越这个限度？并非如此，我并没有丝毫僭越。我听见主亲自通过祂的宗徒说话，那里没有确定的尺度或数目。因为祂说：\BibleQuote{彼此宽恕，若有人对谁有怨，如同天主在基督内宽恕了你们一样。}这里就是规则。如果基督只宽恕了你的罪\BibleQuote{七十个七次}，只赦免到这个限度，而拒绝宽恕更多；那么你也定下这个限度，不愿宽恕超出此限。但如果基督发现了成千上万的罪加于罪，却仍宽恕了全部；那么你不要收回你的仁慈，而要祈求那庞大数目的宽恕。因为主说\BibleQuote{七十个七次}并非没有意义；因为没有任何过犯是你所不应该宽恕的。请看比喻中的这个仆人，他身为债务人，却发现自己有一个债务人，欠他\OriginalTerm{一万塔兰特}{ten thousand talents}。我想一万塔兰特至少代表一万个罪。因为我不会说一个塔兰特如何包括所有的罪。但另一个仆人欠他多少？他欠了\OriginalTerm{一百德纳}{hundred denarii}。这难道不是超过了\BibleQuote{七十七}吗？然而主却发怒，因为他没有宽恕他。不仅是一百超过了\BibleQuote{七十七}，而且一百德纳或许相当于一千个阿司。但这与一万塔兰特相比又算什么呢？\par

4. 因此，让我们准备好宽恕所有得罪我们的过犯，如果我们想要被宽恕。因为如果我们考虑自己的罪，并计算我们在行为上、在眼目上、在耳朵上、在思想上、在无数行动中所犯的罪；我不知道我们是否睡觉时也不欠一个塔兰特。因此我们每日乞求，每日用祈祷叩击天主的耳朵，每日俯伏在地说：\BibleQuote{宽免我们的罪债，如同我们宽免我们的债务人一样。}你的罪债是什么？全部，还是一部分？你会回答：全部。那么你也这样对待你的债务人。这就是你定下的规则，你所说的条件；这是你祈祷时所提及的盟约和协议，说：\BibleQuote{宽免我们，如同我们宽免我们的债务人一样。}\par

5. 那么，弟兄们，\BibleQuote{七十个七次}是什么意思？请听，因为这是一个伟大的奥迹，一个奇妙的圣事。当主受洗时，福音圣史路加在该处按常规的次序、系统和血统纪念了祂的世代，直到基督降生的那一代。玛窦从亚巴郎开始，以降序下到若瑟；而路加却开始以升序计算。为什么一个以降序，另一个以升序计算？因为玛窦展示了基督降生到我们中间的血统；因此他在基督降生时开始以降序计算。而路加则是在基督受洗时开始计算；这是升序的开始，他开始以升序计算，在他的计算中，他完成了七十七代。他从谁开始计算？请注意从谁开始？他从基督开始一直上溯到亚当本人，亚当是第一个罪人，他以罪的锁链生了我们。他上溯到亚当，因此计算出了七十七代；就是说，从基督到亚当，再从亚当到基督，就是前面所说的七十七代。所以，如果没有省略任何一代，那么就没有任何过犯是可以不宽恕的。因此，他计算了这七十七代，这个数字正是主在谈到宽恕罪过时所提到的；因为他从洗礼开始计算，在洗礼中所有的罪都得赦免。\par

6. 而且，弟兄们，请在此看到一个更伟大的奥迹。在七十七这个数字中，有罪过得赦免的奥迹。从基督到亚当恰好有这么多代。那么，现在就更加谨慎努力地探求这个数字的秘密含义，并探究其隐藏的意义；更加谨慎努力地敲门，好让它为你打开。义德在于遵守天主的法律：确实如此。因为法律体现在十条诫命中。因此比喻中的那个仆人\BibleQuote{欠了一万塔兰特}。这就是那由天主的指头所写、并由天主的仆人梅瑟交给百姓的可纪念的十诫。祂\BibleQuote{欠了一万塔兰特}；这象征所有的罪，与法律的数目有关。而另一个\BibleQuote{欠了一百德纳}；同样衍生自同一数目。因为一百乘一百等于一万；十乘十等于一百。一个\BibleQuote{欠了一万塔兰特}，另一个欠了十乘十德纳。因为都没有离开法律的数目，在这两个数目中，你会发现包括各种罪。两人都是债务人，都恳求并乞求宽恕；但那个邪恶、忘恩负义的仆人不愿偿还他所得到的，不愿施予那本不配得的仁慈。\par

7. 那么，诸位弟兄，请思考；人人自圣洗开始；他自由地出去，那\BibleQuote{一万塔冷通}被赦免了；当他出去时，很快会找到某个同仆欠他的债。那么，让他注意罪本身是什么；因为十一这个数字是法律的逾越。因为法律是十，罪是十一。因为法律以十表示，罪以十一表示。为什么罪以十一表示？因为要达到十一，便是逾越了十。但正当的限度定在法律中；逾越它便是罪。当你越过了十，你就到了十一。这高深的奥迹在命建造会幕时被预表出来。那里提到许多数字上的事物，都是极大的奥迹。其中，命制做山羊毛的帐幕，不是十幅，而是十一幅；因为山羊毛象征罪的告解。那么，你还需要什么？你愿知道所有罪都包含在这\BibleQuote{七十七}这个数字中吗？七通常被用来表示整体；因为七天内时间的循环完成，当第七天结束时，它又回到第一天，以便继续同样的循环。在这样的循环中，整个世代过去；然而并没有离开七这个数字。因为祂论到所有的罪，当祂说：\BibleQuote{七十个七次；}因为将十一乘以七，就得到七十七。因此祂愿所有罪都被赦免，因为祂以七十七这个数字标示了它们。那么，不要有人因拒绝宽恕而保留怨恨，免得当他祈祷时，这怨恨被保留在他身上。因为天主说：\BibleQuote{宽恕，你们也将被宽恕。}因为我已经先宽恕了你们；你们至少以后要宽恕。因为如果你们不宽恕，我将召回你们，并把所有我已豁免的再度加在你们身上。因为真理不会说谎；基督既不欺骗，也不受欺骗，祂在比喻的结尾说：\BibleQuote{同样，你们的圣父也要这样对待你们。}你找到了圣父，效法你的圣父。因为如果你不效法祂，你就是在谋划成为弃子。\BibleQuote{同样}，然后\BibleQuote{我的圣父也要这样对待你们，如果你们不各自从心里宽恕自己的弟兄。}不要用口舌说：\Quote{我宽恕，}却延迟在心里宽恕；因为天主以报复的威胁向你显示惩罚。天主知道你在哪里说话。人能听见你的声音；天主洞察你的良心。如果你说：我宽恕，就宽恕吧。你在言语上激烈，却在心里宽恕，总胜于言语柔和，心里却毫无怜悯。\par

8. 那么，不守规矩的孩童会哀求，并很难接受被打，当我们想要以这种方式惩戒他们时，他们便对我们提出异议。\Quote{我犯了罪，但请宽恕我。}好，我宽恕了，他又犯罪。\Quote{请宽恕我，}他喊叫，我宽恕了他。他第三次犯罪。\Quote{请宽恕我，}他喊叫，我第三次宽恕了他。那么，第四次让他挨打。他会说：\Quote{什么！我已经使你厌倦到七十七次了吗？}如果因这样的异议，纪律的严厉沉睡，纪律被压制，邪恶将不受惩罚地肆虐。那么该做什么？让我们以言语责备，如果需要，以鞭打；但同时让我们宽恕罪，并从心里抛弃对它的记忆。因此主补充说：\BibleQuote{从心里，}为的是虽然出于爱心执行纪律，但温柔不离开心。因为什么像外科医生持刀那样仁慈温柔？要被割的人哭喊，却还是被割；要被烧灼的人哭喊，却还是被烧灼。这不是残忍；绝不可称那外科医生的治疗为残忍。他对伤口残忍，为使病人得痊愈；因为如果伤口被温柔对待，这人就完了。因此，我的诸位弟兄，我如此建议：我们应爱我们的弟兄，无论他们怎样得罪了我们；不让我们对他们的爱离开我们的心，并在需要时，对他们执行纪律；免得因纪律的松弛，邪恶增长，我们开始在天主面前被控告，因为经上已向我们宣读：\BibleQuote{你在众人前斥责犯罪的人，为使其余的人也有所惧怕。}当然，如果有人以唯一正确的方式分辨时机，从而解决这个问题，一切就都是真的。如果罪是隐秘的，就秘密地斥责它。如果罪是公开明显的，就公开斥责它，使罪人得以改正；并且\BibleQuote{为使其余的人也有所惧怕。}\par

\SourceRecord{Translated by R.G. MacMullen.；Nicene and Post-Nicene Fathers, First Series；Vol. 6.；Edited by Philip Schaff.；Buffalo, NY: Christian Literature Publishing Co.,；1888.；<http://www.newadvent.org/fathers/160333.htm>.}

% structure: p0001 source=h1 target=section reason=page-title

\section{论新约第三十四篇讲道}

\Emph{[LXXXIV. Ben.]}\par

\Emph{论福音的话}，\BibleRef{玛 19:17}，\Emph{，\Quote{你若愿意进入生命，就该遵守诫命。}}\par

1. 主曾对一位青年人说：\BibleQuote{你若愿意进入生命，就该遵守诫命。}祂没有说：\Quote{你若愿意进入永生，}而是说：\Quote{你若愿意进入生命；}以此将生命定为永生。让我们先陈述一下对今世生活的爱的价值。因为即使是现世生活，无论环境如何，也是被爱的；人们害怕并畏惧结束它，无论它是什么样子；即使充满烦恼和痛苦。因此我们可以看出，可以思考，永生应该如何被爱，当这如此悲惨且必将结束的生命被如此深爱时。弟兄们，请思考，那永不结束的生命，当被何等深爱？你似乎爱这现世生活，在其中你如此劳苦，奔波忙碌，承受疲惫；是的，这悲惨生活的需求几乎数不胜数：播种、犁地、清理土地、航海、磨粉、烹饪、纺织；在这一切之后，你必须结束你的生命。看看你所爱的这悲惨生活中所遭受的恶事；你以为你不会死，永远活着吗？圣殿、石头、大理石，用铁和铅如此牢固地连接在一起，尽管坚固，却会倒塌；人岂能以为自己永不死？那么，弟兄们，学习寻求永生吧，在那里你将不会承受这一切，而是与天主永远为王。\BibleQuote{因为谁愿生活，}正如先知所说，\BibleQuote{就爱享见幸福的日子。}因为在恶日，人宁愿求死而非求生。我们不是常听见并看见，当人们陷入某些磨难和困境，诉讼或疾病，且看到自己受苦时，我们岂不是听见他们只说：\Quote{天主啊，赐我一死，加速我的日子吧}？然而当疾病来临时，他们四处奔走，请医生，许诺金钱和报酬。死亡亲自对你说：\Quote{看，我在这里，刚才你还向主求我，为什么现在要逃避我？我发现你自欺欺人，是这悲惨生活的爱好者。}\par

2. 但关于我们现在所度过的这些日子，宗徒说：\BibleQuote{要赎回时间，因为日子是邪恶的。}这些日子不正是邪恶的吗？我们在这可朽的肉身上，在如此沉重的可朽身体的重担下，在如此大的诱惑中，在如此大的困难中，那里只有虚假的快乐，没有安全的喜乐，折磨人的恐惧，贪婪的贪欲，令人憔悴的悲伤。看，何等的恶日！然而没有人愿意结束这些恶日，因此人们热切祈求天主让他们长寿。然而长寿是什么，不就是长期受折磨吗？长寿是什么，不就是将恶日加到恶日上吗？当男孩长大时，似乎日子在增加；然而他们不知道日子在减少；他们的计算是错误的。因为当我们成长时，我们的天数不是增加，而是减少。例如，为任何人在他出生时指定八十年；他每活一天，就从这个总数中减少一些。然而愚蠢的人却为自己和孩子的频繁生日而欢喜。明智的人啊！如果你瓶子里的酒减少了，你会难过；你正在失去日子，你却高兴吗？那么这些日子是邪恶的；而且因为它们被爱，它们就更邪恶。这个世界是如此诱人，以至于没有人愿意结束悲伤的生命。因为真正有福的生活是当我们复活并与基督一起为王时。不虔诚的人也会复活，但要去火中。然后生命又在那里，但那是有福的。而有福的生命只能是永恒的生命，那里有\BibleQuote{幸福的日子}；而且这些不是许多日子，而是一天。它们被称为\Quote{日子}是沿用今世生活的习惯。那一天没有升起，也没有落下。那一天之后没有明天；因为之前没有昨天。这一天，或这些日子，和这个生命，这个真正的生命，是我们所应许的。它是某种工作的报酬。所以如果我们爱报酬，就不要在工作中失败；这样我们将与基督永远为王。\par

\SourceRecord{Translated by R.G. MacMullen.；Nicene and Post-Nicene Fathers, First Series；Vol. 6.；Edited by Philip Schaff.；Buffalo, NY: Christian Literature Publishing Co.,；1888.；<http://www.newadvent.org/fathers/160334.htm>.}

% structure: p0001 source=h1 target=section reason=page-title

\section{新约讲道35}

\Emph{[第八十五篇，本笃版]}\par

\Emph{论福音之言}，\BibleRef{玛 19:17}，\Emph{\BibleQuote{你若愿意进入生命，就该遵守诫命。}}\par

弟兄们，此刻在我们耳中回响的福音讲道，与其说需要一位讲解者，不如说需要一位留心听和遵行的人。还有什么比这光更清楚的呢？\BibleQuote{你若愿意进入生命，就该遵守诫命。}那么我还有什么可说的呢？只有说：\BibleQuote{你若愿意进入生命，就该遵守诫命。}谁不渴望生命呢？然而谁又真的愿意遵守诫命呢？你若不愿遵守诫命，为何寻求生命？你若对工作懈怠，为何冲向奖赏？福音中的那位富少年说他已遵守了诫命；然后他听到了更大的诫命：\BibleQuote{你若愿意是成全的，你还缺少一件：去变卖你所有的一切，施舍给穷人；}你不会失去它们，反而\BibleQuote{你将有宝藏在天堂；然后来跟随我。}你若做到其余一切，却不跟随我，这对你有什么益处呢？但是正如你所听见的，他忧闷地\BibleQuote{走了}，且\BibleQuote{悲伤；因为他拥有许多财产。}他所听见的，我们也听见了。福音是基督的声音。他坐在天堂，却不停地在世上说话。我们不要耳聋，因为他正在呼喊。我们不要死亡，因为他正在打雷。若你不愿做更大的事，至少做那些较小的。若更大的负担对你太重，至少担起那些较小的。你为什么对两者都迟钝？为什么反对两者？更大的事是：\BibleQuote{变卖你所有的一切，施舍给穷人，然后跟随我。}较小的事是：\BibleQuote{不可杀人，不可奸淫，不可偷盗，不可作假见证。孝敬你的父母；并且，你当爱你的邻人如同你自己。}做这些事吧；我为何要呼吁你变卖你的财产，既然我不能从你那里得到什么，只希望你不再抢夺他人的东西？你已听见\BibleQuote{不可偷盗}，你却抢夺。在如此伟大的审判者眼前，我发现你不只是小偷，更是强盗。宽恕你自己，怜悯你自己吧。今生还给你喘息的机会，不要拒绝改正。昨天你是小偷；今天不要再这样。或者，如果你今天已经这样做了，明天不要再做。在某时停止你的恶行，这样你才能要求善的奖赏。你渴望善的事物，却不愿成为善人；你的生命与你的愿望相矛盾。若拥有一座好的庄园是极大的善，那么拥有邪恶的灵魂又是多大的恶！\par

2. 富少年\BibleQuote{忧闷地走了}；主说：\BibleQuote{有钱财的人进天国是多么难啊！}他用一个比喻显示出困难程度以至于完全不可能。因为不可能的事都是困难的，但不是所有困难的事都是不可能的。至于多么困难，请留意这个比喻：\BibleQuote{我实在告诉你们：骆驼穿过针眼，比有钱人进天国还容易。}骆驼穿过针眼！若他说是蚊子，那是不可能的事。当他的门徒们听见了，就忧闷地说：\BibleQuote{若是这样，谁能得救呢？}什么有钱人？那么请听基督的话吧，你们穷人，我是对天主子民说话。你们中穷人比富人更多，所以你们至少接受我所说的，留心听。你们中凡夸耀自己贫穷的，要小心骄傲，免得谦卑的富人胜过你们；要小心不敬，免得虔诚的富人胜过你们；要小心醉酒，免得节制的富人胜过你们。不要因你的贫穷而自夸，如果他们不能因他们的财富而自夸。\par

3. 让富人们听吧，如果他们确实是富人的话；让他们听宗徒的话：\BibleQuote{你要嘱咐今世的富人，}因为另有一些人是另一个世界的富人。另一个世界的富人是宗徒们，他们说：\BibleQuote{似乎一无所有，却拥有万有。}这样你就知道他所说的是哪种穷人了：他加上了\BibleQuote{今世}。那么让\BibleQuote{今世的富人}听宗徒的话吧，他说：\BibleQuote{你要嘱咐今世的富人，不要心高气傲。}骄傲是财富的第一条蛀虫。这是一种吞噬一切，甚至使之化为尘土的病态蛀虫。\BibleQuote{因此，你要嘱咐他们，不要心高气傲，也不要寄望于无定的钱财}（这是宗徒的话），\BibleQuote{而要寄望于永生的天主。}贼可能偷走你的金子；谁能偷走你的天主？如果富人没有天主，他有什么？如果穷人有了天主，他缺什么？因此他说：\BibleQuote{也不要寄望于钱财，而要寄望于永生的天主，他厚赐百物供我们享用；}赐予这一切的同时，他也赐予了自己。\par

4. 如果他们不应该\BibleQuote{信赖钱财}，也不该\BibleQuote{依靠它们，而应依靠永生的天主}，那他们该如何处理自己的钱财呢？请听：\BibleQuote{让他们在善工上富有。}这是什么意思？解释一下，啊，宗徒。因为许多人不愿理解他们不愿实践的事。解释一下，啊，宗徒；不要因你言辞的晦涩而给恶行以机会。告诉我们你所说的\BibleQuote{让他们在善工上富有}是什么意思。让他们听到并理解；不要让他们被容许为自己辩护；反而让他们开始控告自己，并说出我们刚才在圣咏中所听到的：\BibleQuote{因为我承认我的罪恶。}告诉我们这是什么意思：\BibleQuote{让他们在善工上富有。让他们容易施舍。}而\BibleQuote{让他们容易施舍}又是什么？怎么！连这个也不明白吗？\BibleQuote{让他们容易施舍，让他们分享。}你有，别人没有：分享吧，这样天主才能与你分享。在这里分享，你将在那里分享。在这里分享你的饼，你将在那里领受食粮。什么饼在这里？就是你汗流满面、辛苦劳作所收获的饼，这是对第一个人的诅咒。什么食粮在那里？就是那说过\BibleQuote{我是从天上降下的永生之粮}的那一位。在这里你是富有的，但在那里你是贫穷的。你有金子，但你还没有基督的临在。把你所有的拿出来，好领受你所没有的。\BibleQuote{让他们在善工上富有，让他们容易施舍，让他们分享。}\par

5. 那么他们必须放弃所有吗？他说：\BibleQuote{让他们分享}，而不是\Quote{让他们全部施舍。}让他们为自己保留足够的部分，让他们保留比足够更多的。让我们捐出其中一部分。什么部分？十分之一？经师和法利塞人捐献了十分之一，而基督尚未为他们倾流鲜血。经师和法利塞人捐献了十分之一；以免你或许认为自己将饼分给穷人是什么大事；而这几乎仅是你财富的千分之一。然而，我并非对此有所指责；就这样做吧。我是如此饥渴，以至于连这些碎屑也高兴。但我还是不能隐瞒为我们而死的那位在世时所说的话：\BibleQuote{除非你们的义德超过经师和法利塞人的义德，你们决进不了天国。}他没有对我们温柔；因为他是医生，他切中要害。\BibleQuote{除非你们的义德超过经师和法利塞人的义德，你们决进不了天国。}经师和法利塞人捐献了十分之一。你怎么样？问问自己。想想你做了什么，用什么做的；你给了多少，给自己留了多少；你在慈悲上花费了多少，在奢侈上留了多少。因此，\BibleQuote{让他们容易施舍，让他们分享，让他们为自己积攒美好的根基，以备将来，好能得着永生。}\par

6. 我已劝告了富人；现在请听，穷人。你们富人，施舍你们的钱财；你们穷人，不要抢夺。你们富人，分施你们的财物；你们穷人，节制你们的欲望。请听，穷人，这位同一宗徒的话：\BibleQuote{虔敬与知足，乃是大利。}利益就是获得收益。世界是你们与富人共有的；你们没有与富人共有的房屋，但你们共有天空，共有光。只寻求知足，寻求够用，不要奢求更多。其余的一切，都是负担而非帮助；是累赘而非荣耀。\BibleQuote{虔敬与知足，乃是大利。}首先是虔敬。虔敬就是敬拜天主。\BibleQuote{虔敬与知足。因为我们没有带什么到这世上来。}你带什么来了吗？不，即使是你们富人也没有带来任何东西。你们在这里发现了一切，你们出生时赤身露体，如同穷人。两者都有同样的身体软弱；同样的婴儿啼哭，这是我们悲惨的见证。\BibleQuote{因为我们没有带什么到这世上来，也不能带什么出去。只要有衣有食，就当知足。}\BibleQuote{因为他们想要发财。}\BibleQuote{他们想要发财}，而不是已经发财。因为那些已经发财的人，很好。他们已经听到了他们的教训，要\BibleQuote{在善工上富有，容易施舍，分享。}他们已经听到了。现在你们这些尚未富有的人，请听。\BibleQuote{他们想要发财，就陷于诱惑和罗网，以及许多有害而愚昧的欲望。}你们不害怕吗？请听下面的话：\BibleQuote{这使人沉溺在败坏和灭亡中。}你们现在还不害怕吗？\BibleQuote{因为贪婪是万恶之源}？贪婪就是想要发财，而非已经富有。这就是贪婪。你们不害怕\BibleQuote{沉溺在败坏和灭亡中}吗？你们不害怕\BibleQuote{贪婪是万恶之根}吗？你从你的田里拔出荆棘的根，却不肯从你的心里拔出恶欲的根吗？你洁净你的田地，好使你的身体得到果实，却不肯洁净你的心灵，好使你的天主居住其中吗？\BibleQuote{因为贪婪是万恶之源，有些人因贪求它，就离弃了信仰，并用许多愁苦把自己刺透了。}\par

7. 你们现在已听到了你们必须做什么，已听到了你们必须害怕什么，已听到了天国如何可以买得，已听到了什么会阻碍天国。愿你们全体一心一意，服从天主的圣言。天主创造了富人和穷人。圣经说：\BibleQuote{富人和穷人彼此相遇，上主是他们的造主。}富人和穷人彼此相遇。在什么方面？只在此生？富人和穷人出生一样。你们在同行时相遇。不要压迫或欺骗。一个有所缺乏，另一个绰绰有余。但\BibleQuote{上主是他们的造主。}通过那有的，他帮助那缺乏的；通过那没有的，他检验那有的。我们已经听到，我们已经说过；让我们敬畏，让我们留心，让我们祈祷，让我们达到。\par

\SourceRecord{Translated by R.G. MacMullen.；Nicene and Post-Nicene Fathers, First Series；Vol. 6.；Edited by Philip Schaff.；Buffalo, NY: Christian Literature Publishing Co.,；1888.；<http://www.newadvent.org/fathers/160335.htm>.}

% structure: p0001 source=h1 target=section reason=page-title

\section{新约讲道集 第三十六篇}

\Emph{[LXXXVI. Ben.]}\par

\Emph{论《福音》中的话，\BibleRef{玛 19:21}，\Quote{你去，变卖你所有的，分给穷人}等等。}\par

1. 今日的福音读经提醒我，亲爱的诸位，要向你们谈论天上的宝藏。因为我们的天主并非如不信的贪婪之人所想，要我们失去所有；若我们正确理解、虔诚相信、热心领受所吩咐的，祂并非要我们失去，反而指示了一个可以积蓄的地方。因为没有人能不想着自己的财宝，他的心总在某种旅程中追随他的财富。若财宝藏于地下，他的心必寻觅最底层之处；若财宝存于天上，他的心便居高位。因此，若基督徒愿意行他们所知并公开承认的（并非所有听见者都知道；但愿知道者不致徒然知道），若他们愿意\Quote{举心向上}，就应将所爱的积蓄在天上；虽仍处于血肉的尘世，却与基督同心居住；如同教会之首已先行，基督徒的心也当先行。正如肢体将追随已经先行的元首基督，同样，每人复活时也将追随其心现已先行之处。让我们用那可先行的一部分从此世先行；我们的整个自我将追随已先行的部分。我们地上的房屋必将倾颓；我们天上的房屋是永恒的。让我们将财物先行搬往我们准备前去的所在。\par

2. 我们刚听见一位富人来向\Quote{善师}求问得永生的方法。他所爱的是大事，他不愿舍弃的却无价值。于是，因着对无价值之物的过度贪爱，他心怀悖逆，听了方才称为\Quote{善师}的那位的话，却失去了宝贵之物的拥有。他若不愿得永生，就不会求问如何得永生。那么，弟兄们，他为何拒绝那被他称为\Quote{善师}者从信德道理中为他提出的话语呢？什么？难道在教导之前是\Quote{善师}，教导之后却成了恶师？在教导之前，他被称作\Quote{善}。他没有听见他所愿听的，却听见了对他有益的话；他带着渴望而来，却忧伤而去。若祂对他说\Quote{失去你所有的}，他又将如何？当祂说\Quote{你的财宝安然无恙}时，他却忧伤而去。祂说：\BibleQuote{你去，变卖你所有的，分给穷人。}你或许害怕失去。请看下文：\BibleQuote{你将有宝藏在天上。}也许你曾雇一个年轻奴仆看守你的财宝；你的天主要作你黄金的守护者。那在地上赐予你的，祂要在天上亲自为你保存。或许他本不犹豫将财物托付基督，却因听到\Quote{分给穷人}而忧伤，仿佛在内心说：\Quote{你若说\InnerQuote{交给我，我在天上为你保存}，我必不犹豫交给我的主、\InnerQuote{善师}；但如今你却说了\InnerQuote{分给穷人}}\par

3. 没有人该害怕施舍穷人；没有人以为他所看见的手就是接受者。那吩咐你施舍的，祂接受。我说这话并非出于己意或人的臆测；请听祂亲自说的话，祂既劝勉你，又给你安全的保证。祂说：\BibleQuote{我饿了，你们给了我吃的。}当他们列举所有善事后回答：\BibleQuote{我们几时见了你饿了？}祂答道：\BibleQuote{你们对我这些最小兄弟中的一个所做的，就是对我做的。}是穷人在乞讨，但接受者却是富有的。你给那会浪费的人，接受者却要归还。祂不仅归还所接受的，还乐意付利息，应许的比你给的更多。现在放纵你的贪欲，想象自己是个放高利贷者。你若真是放高利贷者，必受教会责斥，被天主的话语驳倒，所有弟兄都唾弃你，如同残酷的高利贷者，想从别人的眼泪中榨取利益。但现在，做个放高利贷者吧，无人阻拦你。你愿意借给穷人，他偿还时总带愁苦；但如今借给那有能力偿还的债主，祂甚至劝你接受祂所应许的。\par

4. 向天主施舍，并向天主索要报偿。不，更确切地说，向天主施舍，你将被迫接受报偿。的确，在世上你必须寻找你的债务人；而他也寻找，但只是为了找到何处可以躲避你的面。你曾到法官面前说：\Quote{命我的债务人被传来；}他听到这话就逃走了，甚至不屑于向你问安，尽管你或许在他困乏时通过借贷给了他财富。那么你有一个可以放心投入钱财的对象。向基督施舍；祂会主动催迫你接受，而你甚至会惊奇祂究竟从你这里领受了什么。因为对那些站在祂右边的人，祂首先要说：\BibleQuote{你们蒙我父降福的，来吧。}\Quote{来}到哪里？\BibleQuote{承受自创世以来为你们预备的国度。}为什么？\BibleQuote{因为我饿了，你们给了我吃的；我渴了，你们给了我喝的；我作客，你们收留了我；我赤身露体，你们给了我穿的；我患病，你们看顾了我；我在监里，你们来探望了我。}他们要说：\BibleQuote{主啊，我们几时见了祢？}这是什么意思？债务人催迫着要偿还，而债权人却推辞。但那位可靠的债务人不会让他们因此受损。\Quote{你们犹豫着不肯接受吗？我已经接受了，你们还不知道吗？}祂便回答如何接受的：\BibleQuote{你们对我这些最小兄弟中的一个所做的，就是对我做的。}\Quote{不是我自己接受的，而是通过我的弟兄。给他们的就归到我身上；放心，你们没有失去。你们在世上指望那些偿还能力有限的人；你们在天上有一位偿还能力无限的主。我，}祂说，\Quote{已经接受了，我将偿还。}\par

5. 而我接受了什么，我又偿还什么？\Quote{\InnerQuote{我饿了，}祂说，\InnerQuote{你们给了我吃的；}等等。}我接受了尘世之物，我将赐予天堂；我接受了暂时之物，我将恢复永恒；我接受了饼，我将赐予生命。是的，我们甚至可这样说：\Quote{我接受了饼，我将赐予饼；我接受了饮料，我将赐予饮料；我接受了住处，我将赐予房屋；我在病中被探望，我将赐予健康；我在监中被探望，我将赐予自由。你们给我的穷人的饼已被消耗；而我将赐予的饼既恢复软弱，又不衰竭。}愿祂赐予我们饼，祂是从天降下的生命之饼。当祂赐予饼时，祂将赐予自己。因为你们放债时意图什么？给钱，收钱；但少给多收。\Quote{我，}天主说，\Quote{会为你给我的一切给你更佳的交换。}因为如果你给一磅银子，收一磅金子，你会有多大的喜乐？考察并质问贪婪。\Quote{我给了一磅银子，收到一磅金子！} 银子与金子之间有何比例！更有甚者，尘世与天堂之间有何比例！你的金银你要留在此世；而你自己不会永远留在此处。\Quote{我要给你另一样、更多一样、更好一样；我要给你甚至那永存之物。}所以，弟兄们，愿我们的贪婪被遏制，好使另一种神圣的贪婪被点燃。她的劝告完全是邪恶的，她阻止你们行善。你们宁愿服侍邪恶的主母，却不承认良善之主。有时两个主母占据人心，将理应承受双重轭的奴隶撕碎。\par

6. 是的，有时两个对立的主母占据一个人：贪婪和奢靡。贪婪说：\Quote{守着；}奢靡说：\Quote{花掉。}在两个发号施令、索取不同东西的主母之下，你能做什么？她们各有其言语方式。当你开始不愿服从她们，朝自由迈出一步时，由于她们无权命令，便改用谄媚。她们的谄媚比命令更需警惕。贪婪说什么？\Quote{为自己留着，为你的孩子留着。如果你变得匮乏，没有人会给你。不要只为当下而活；为将来考虑。}另一方面是奢靡。\Quote{趁你活着就享受吧。善待你自己的灵魂。你必定会死，且不知何时；你不知道你将把你的财产留给谁，谁会拥有它。你在夺走你口中的食粮，也许你死后，你的继承人甚至不会在你的坟墓上放一杯酒；或者即使放了一杯，他也会自己喝醉，一滴也不会到你的嘴里。因此，趁你还能的时候，善待你自己的灵魂吧。}所以贪婪命令一件事：\Quote{为自己留着，为将来考虑。}奢靡命令另一件：\Quote{善待你自己的灵魂。}\par

7. 然而，你这蒙召得自由的人，你当厌烦，你当厌烦你侍奉这类女主人的奴役。你要承认你的救赎者、你的解救者。你当事奉祂，祂所命令的事更为容易，祂所命令的事并非彼此矛盾。我敢进一步说：贪婪和奢侈向你命令了相反的事，使你不能同时服从二者；一个说：\Quote{为你自己存留，为未来谋划}；另一个说：\Quote{任意挥霍，善待你的灵魂。}如今，让你的主和你的救赎者出来，祂也要说同样的话，却没有任何矛盾。你若不愿，祂的家不需要不情愿的仆人。你要想想你的救赎者，想想你的赎价。祂来救赎你，祂倾流了祂的宝血。祂看你是宝贵的，因你付出了如此宝贵的代价。你既承认买你的那一位，就要想想祂是从什么中解救你。其他那些骄傲地主宰你的罪，我暂且不提，因为你曾服事过无数的主人。我只说这两个：奢侈和贪婪，它们向你发出相反的指令，把你拖向不同的事。你要从它们中解救自己，来到你的天主面前。你若曾是不义的奴仆，如今要做正义的奴仆。它们曾对你说的话，以及它们给你的相反命令，如今你从你的主那里听到完全相同的话，然而祂的命令并不矛盾。祂没有夺去它们的话语，却夺去了它们的权势。贪婪曾对你说什么？\Quote{为你自己存留，为未来谋划。}话语并未改变，但人改变了。如今，你若愿意，就请比较这些劝谕者：一个是贪婪，另一个是正义。\par

8. 检查这些相反的诫命。\Quote{为你自己存留}，贪婪说。假设你愿意服从它，问它要存留在何处？它会指给你一个防守严密的地方，有墙的密室，或铁柜。好吧，你尽可采取一切防范措施；然而，也许家里有个贼会撬开秘处；你一面为钱财防范，一面要为自己的性命担忧。或许当你看守你的积蓄时，那存心要劫掠它的人，甚至已在图谋杀害你。最后，即使你通过各种防范，使你的财宝和衣物免于盗贼，却仍要提防锈蚀和蛀虫。那时你能做什么呢？外界没有敌人夺去你的货物，但内里却有一个消耗它们的。\par

9. 因此，贪婪没有给出良策。看，它命令你存留，却未找到一个可以存留的地方。让它给出另一个建议：\Quote{为未来谋划。}为哪一个未来？为那些不确定的短暂日子。它对一个也许活不到明天的人说：\Quote{为未来谋划。}但假设他活得像贪婪所想的那样长久——不是它能证明、保证或确信的那样，而是像它所想的那样，他渐渐老去，最终死去：当他已老迈佝偻，拄着拐杖时，他仍求利，仍听见贪婪说：\Quote{为未来谋划。}为哪一个未来？当他已经奄奄一息时，它仍说话。它说：\Quote{为你的儿女起见。}但愿我们至少没有发现那些无儿女的老人也是贪婪的。然而对这些人，甚至对那些连用任何虚情假意来掩盖其不义都不能的人，它也喋喋不休地说：\Quote{为未来谋划。}但这些人或许很快就会为自己羞愧；所以我们来看那些有儿女的人，他们是否肯定自己的儿女会拥有他们留下的东西？让他们在生前观察别人的儿女：有些因他人的不义暴力而失去所有，有些因自己的邪恶而消耗他们所拥有的；那些曾是富家子女的人，却落得贫困境地。所以，不要再作贪婪的家生奴隶了。但有人会说：\Quote{我的儿女将会拥有这些。}这是不确定的；我不说它是假的，但充其量是不确定的。但即使它是确定的，你愿意留给他们什么呢？你自己所挣得的。确实，你所挣得的并没有人留给你，而你却拥有了它。如果你能拥有不是留给你的东西，那么他们也能拥有你不留给他们的东西。\par

10. 如此，贪婪的劝谕已被驳斥；但现在让主来说同样的话，让正义发言：话语相同，但含义不同。\Quote{为你自己存留}，主说，\Quote{为未来谋划。}现在问祂：\Quote{我该存留在何处？} \BibleQuote{你必有财宝在天上，那里没有贼接近，也没有蛀虫腐蚀。}为怎样一个持久的未来，你当存留它！\BibleQuote{来吧，我父所祝福的，承受从创世以来为你们预备的国。}这个国是多少天的呢？这节经文的结尾指明了。因为祂对左边的人说了\BibleQuote{这些人要进入永火里去}之后，对右边的人说：\BibleQuote{但义人却要进入永生。}这就是\Quote{为未来谋划。}一个不再有未来的未来。那些无尽的日子既被称为\Quote{日子}，又被称为\Quote{一日}。有一个人论到那些日子时说：\BibleQuote{使我得以住在上主的殿中，直至悠长岁月。}它们也被称为一日：\BibleQuote{你是我的儿子，我今日生了你。}那些日子是一日；因为其中没有时间；那一日既无昨日在先，也无明日在后。所以，让我们\Quote{为\Emph{那}未来谋划：}从前贪婪对你说的话语，其实与此并无不同；然而却藉着它们，贪婪被推翻了。\par

11. 或许还可说一句：\Quote{但我该怎样对待我的儿女呢？} 关于这一点，也当听从你主的劝告。倘若你的主对你说：\Quote{创造他们的我，比生养他们的你，更关心他们，} 你或许无话可说。然而你将会看见那福音中忧愁而去的富人，并受责备，你或许会对自己说：那富人作恶，没有变卖一切周济穷人，因他没有儿女；但我有儿女，我该为他们留存一些。在这软弱中，主也预备与你商议。我愿因祂的仁慈而大胆发言；我愿说些话，并非出于我的想象，而是出于祂的怜悯。那么，也为你的儿女留存吧，但请听我说。假定（按人的处境）有人失去一个儿女；请注意，弟兄们，请注意，那贪吝在此世或来世都没有藉口。我说，这是人的处境；并非我愿如此，但我们眼见实例。某个基督徒孩子去世了：你失去了一个基督徒孩子；其实你并未失去他，而是送他先你而去。他并非完全离去，而是先行一步。请问你自己的信德：你不久也将到那里，到他所先去的地方。我只问一个简短的问题，我想无人能回答。你的儿子活着吗？请问你的信德。若他活着，为何他的那份被弟兄们抢占？但你必会说：怎么，他会回来再拥有它吗？那就把它送到他先去的地方吧；他不能回到他的财物那里，但他的财物可以到他那里去。只须想想他与谁同在。倘若有个儿子在宫廷供职，成了皇帝的朋友，对你说：\Quote{卖掉我在那里的那份，寄给我；} 你还能回答他吗？你的儿子现今与万皇之皇、万王之王、万主之主同在；寄给祂吧。我不是说你的儿子自身有所欠缺；而是与他同在的主，在地上有所欠缺。祂俯允在此收取，而在天上赐予。去做那些贪婪之人常做的事吧：立契据，赐予那些在旅途中的人，使你在本国得以收取。\par

12. 但现在我根本不谈你自己，而是谈你的孩子。你犹豫不决，不肯给出属于你自己的，不，是归还属于别人的；由此你无疑被定罪：你并非为儿女积蓄。看，你不给儿女，甚至要从儿女那里夺走。至少从这个孩子你要夺走。为何他不配领受他的一份，因为他与一位比所有人更尊贵者同住？这理由若成立，除非你儿子与之同住的那位不愿接收。如今你要为你的家，却不为天主的家而富足。因此，我要对你说：\Quote{将你的所有给出，} 远非如此；我是在对你说：\Quote{偿还你所欠的。} 但你必会说：\Quote{他的弟兄们将得到它。} 恶毒的准则，会教导你的儿女盼望兄弟死亡。若他们因已故兄弟的财产而致富，当心他们在你家中彼此守望。那么你将做什么？你会分他的遗产，从而教导弑亲吗？\par

13. 但我不愿谈及失去孩子的事，免得我似乎以灾祸相威胁，而灾祸的确会临到人。让我们用更愉快和吉利的口吻来说话。那么我不说你会少一个；不如算你多了一个。给基督在你儿女中留一个位置，让你的主加入你的家庭；让你的创造者加入你的后代，让你的弟兄加入你儿女的数目。因为虽有如此大的鸿沟，祂仍屈尊成为弟兄。尽管祂是圣父的独生子，祂却俯允有共同继承人。看，祂赐予得何其丰盛！你为何要如此吝啬地给予？你有两个孩子，算祂为第三个；你有三个，让祂算为第四个；你有五个，让祂称为第六个；你有十个，让祂为第十一个。我不再多说；为主保留一个孩子的位置。因为凡你给主的，将有益于你和你的儿女；反之，你为儿女不义地保留的，将伤害你与他们。如今你将给出一份，算作一个孩子的一份。就算你多了一个孩子。\par

14. 弟兄们，这是何等重大的要求？我只是给你们劝告，难道我使用暴力吗？正如宗徒所说：\BibleQuote{我说这话，是为你们自己的益处，并非要陷害你们。} 我设想，弟兄们，对于为人父者，以为自己多了一个孩子，从而获得那永久的产业，即你和你的儿女都可永远拥有的，这实在是轻松且容易的念头。贪吝在这事上无话可说。你们听了这些话已欢呼。倒要转你们的言语攻击她；别让她胜过你们，别让她在你们心中有比你们的救赎主更大的权能。别让她在你们心中有比那劝勉我们\Quote{举心向上}的主更大的权能。现在，让我们辞别她。\par

15. 奢靡说甚么？甚么？\Quote{善待你的灵魂。}请看，主也说了同样的话：\Quote{善待你的灵魂。}奢靡对你说的话，正义也对你说同样的话。但在这里，要再思量这些话语是在何种意义上用的。你若愿意善待自己的灵魂，便想想那个听从奢靡和贪婪的劝告、想要善待自己灵魂的富人罢。他的\Quote{田地出产丰盈，他没有地方存放他的果实；他说：我该怎么办？}我没有地方存放我的果实；我找到了办法：\Quote{我要拆掉我的}旧\Quote{仓房，建造新的，}将其装满，\Quote{然后对我的灵魂说：你拥有许多财物，尽情享受罢。}请听反对奢靡的劝告：\Quote{愚昧人哪，今夜就要索回你的灵魂；你所预备的一切，将归谁呢？}那将被索回的靈魂要往哪里去？今夜就要索回，而他不知往何处去。\par

16. 再想想那个奢侈、傲慢的富人。他\Quote{天天奢华宴乐，身穿紫袍和细麻衣；}而\Quote{穷人躺在门前，满身疮痍，}徒然\Quote{渴望从富人桌上掉下的碎屑；}他用疮痍喂了狗，却没有得到富人的喂养。两人都死了；其中一人被埋葬了；至于另一人，说了甚么？\Quote{他被天使带到亚巴郎的怀中。}富人看见了穷人；不，现在是穷人看见了富人；他渴望从他的指尖得到一滴水落在舌头上，而那人曾渴望从他桌上得到一块碎屑。他们的命运确实颠倒了。死去的富人在徒然祈求此事：但愿我们这些活着的人不要徒然听闻。因为他想再回到世上，却不被允许；他想从死者中派一人到他的弟兄那里去，这也未被允准。但对他说的甚么？\Quote{他们有梅瑟和先知。}他说：\Quote{除非有人从死者中复活，他们不会听从。}亚巴郎对他说：\Quote{如果他们不听从梅瑟和先知，即使有人从死者中复活，他们也不会信服。}\par

17. 因此，奢靡在施舍和为灵魂预备未来的安息——以便我们\Quote{善待自己的灵魂}——这一点上，以歪曲的方式所说的，梅瑟和先知也说过。让我们活着的时候倾听。因为在那里，谁若在这里轻视了这些话语，必将徒然渴望听见。我们是否期待有人甚至从死者中复活，告诉我们善待自己的灵魂？这事已经发生了：你的父亲没有复活，但你的主复活了。听从祂，接受善意的劝告。不要吝惜你的财物，尽可能地慷慨使用。这是奢靡的声音：它已成为主的声音。尽可能地慷慨使用，善待你的灵魂，免得今夜你的灵魂被索回。那么，在这里你以基督之名得到了我认为关于施舍义务的一篇讲道。你们现在鼓掌赞许的声音，只有在主看见你们的手在行慈悲工作时，才真正令祂喜悦。\par

\SourceRecord{Translated by R.G. MacMullen.；Nicene and Post-Nicene Fathers, First Series；Vol. 6.；Edited by Philip Schaff.；Buffalo, NY: Christian Literature Publishing Co.,；1888.；<http://www.newadvent.org/fathers/160336.htm>.}

% structure: p0001 source=h1 target=section reason=page-title

\section{论新约·第三十七篇讲道}

\EditorialNote{LXXXVII. Ben.}\par

\Emph{主日宣讲，论福音所载}，\BibleRef{玛 20:1}\Emph{，\BibleQuote{天国好像一个家主，清晨出去为自己的葡萄园雇工人。}}\par

1. 你们已从圣福音中听到一个非常适合当前时节的比喻，论葡萄园的工人。因为现在是物质收获的时节。也有一种属灵的收获，天主在其中因祂葡萄园的果实而喜乐。因为我们耕种天主，天主也耕种我们。但我们并非以致使祂变得更好的方式耕种祂。因为我们的耕种是心灵的劳动，而非双手的。祂耕种我们，如同农夫耕种田地。因此，当祂耕种我们时，祂使我们变得更好；因为农夫通过耕种使田地变得更好，祂在我们身上寻求的果实正是我们耕种祂。祂施于我们的耕种是：祂不断用祂的圣言从我们心田里拔除恶种，用祂圣言的犁打开我们的心，种下祂诫命的种子，等待虔敬的果实。因为当我们如此接受那耕种进入内心，以至善于耕种祂时，我们就不是忘恩负义对待我们的农夫，而是结出祂所喜悦的果实。我们的果实并不使祂更富有，而是使我们更幸福。\par

2. 那么请看，听我如何说：\Quote{天主耕种我们。}因为至于我们耕种天主，无须向你们证明。因为所有人都把这挂在嘴上，说人耕种天主；但听众听到天主耕种人时，会感到某种敬畏，因为按照人的惯常用语，不说天主耕种人，而说人耕种天主。因此我们必须向你们证明，天主也耕种人；免得我们被认为说了违背纯正道理的话，而人们在心里与我们争辩，因不知我们的意思而责难我们。所以我决意要给你们指明，天主也耕种我们；但正如我已经说过的，如同田地，为使祂使我们变得更好。主在福音中这样说：\BibleQuote{我是葡萄树，你们是枝条；我父是农夫。}农夫做什么？我问你们这些农夫。我想他耕种自己的田地。如果天主父是农夫，祂就有田地；祂耕种祂的田地，并期望从中得到果实。\par

3. 再者，祂\BibleQuote{栽种了一个葡萄园}，正如主耶稣基督亲自说的，\BibleQuote{租给了园户，他们应在适当的时候把果实交给祂。祂派仆人到他们那里去收取葡萄园的租金。但他们虐待了他们，杀了一些人，}并轻蔑地拒绝交出果实。\BibleQuote{祂又派了别的仆人，}他们也遭受了同样的对待。于是那家主，即田地的耕种者、葡萄园的栽种者和出租者说：\BibleQuote{我要派我的独生子去，也许他们会尊敬他。}因此祂说，\BibleQuote{祂也派了自己的儿子。他们彼此说：这是继承人，来，我们杀了他，产业就归我们了。于是他们杀了他，把他扔出葡萄园。葡萄园的主人来了，他会怎样处置那些恶园户呢？他们回答说：他要凶恶地消灭那些恶人，并把葡萄园租给别的园户，他们会在季节到来时把果实交给他。}葡萄园是在犹太人心中赐下法律时栽种的。先知们被派去寻求果实，即他们的良善生活；先知们被他们虐待、杀害。基督也被派来，即家主的独生子；他们杀了那位继承人，从而失去了产业。他们的恶谋适得其反。他们杀了他，以为可以占有产业；但因为杀了他，他们失去了产业。\par

4. 你们刚才也听到了出自神圣福音的比喻：\BibleQuote{天国好像一个家主，清早出去雇工人进他的葡萄园。他在早晨出去}，雇了他所遇见的人，与他们讲定一个银币作为工钱。他\BibleQuote{又在第三时辰出去，看见另有些人}，就带他们到葡萄园里去工作。\BibleQuote{约在第六和第九时辰，他又照样做了。他也在第十一时辰出去}，接近一天的结束，\BibleQuote{看见有些人闲着站在那里，就对他们说：你们为什么站在这里？} 你们为什么不在葡萄园里工作？他们回答说：\BibleQuote{因为没有人雇我们。} \BibleQuote{你们也去吧，} 他说，\BibleQuote{我给你们合理的报酬。} 他乐意把工钱定为一个银币。那些只工作了一个时辰的人怎敢希望得到一个银币呢？然而他们庆幸自己有望得到一些报酬。于是这些人连仅一个时辰也被带进去了。到了一天结束时，他吩咐从最后来的到最先来的给所有人付工钱。于是他从那些在第十一时辰进来的人开始付，命令给他们一个银币。那些在第一个时辰来的人看见别人得到了一个银币，而那是他与他们讲定的工钱，\BibleQuote{他们以为会多得一些；} 轮到他们时，他们也得到了一个银币。\BibleQuote{他们就对家主抱怨说：你看，我们这些承担了白天的炎热和劳苦的，你竟使他们与只工作了一个时辰的人平等一样。} 而\BibleQuote{家主} 很公正地回答其中一个说：\BibleQuote{朋友，我没有亏负你；} 就是说，我没有欺骗你，我给了你与你讲定的工钱。\BibleQuote{我没有亏负你，} 因为我给了你讲定的。对于这些后来的人，我愿意不偿还工钱，而是施予礼物。\BibleQuote{难道不许我拿我的财物做我所愿意的事吗？因为我善良，你就眼红吗？} 如果我拿走了不属于任何人的东西，我确实该被责备为欺诈和不义；如果我没有付清某人的工钱，我确实该被责备为欺诈和扣留他人的财物；但我付清了应付的，并另外给我所愿意的人，那么我欠他工钱的人不能挑剔，我给予的人更应该高兴。他们无话可答；所有人都平等了；\BibleQuote{后来的成了在先的，在先的成了后来的；} 这是通过待遇平等，而非颠倒顺序。因为\BibleQuote{后来的成了在先的，在先的成了后来的} 是什么意思呢？即先来的和后来的都得到了相同的。\par

5. 为什么他从最后的人开始付钱？难道不是所有人如我们所读的那样一起领受吗？因为我们在福音的另一处读到，祂将对那些放在右边的人说：\BibleQuote{你们蒙我父降福的，来承受自创世以来为你们预备的国度吧。} 既然所有人都要一起领受，我们如何在此处理解：那些在第十一时辰开始工作的人先领受，而最先被雇的人后领受呢？如果我能如此讲话，以致达到你们的理解，愿天主受赞美。因为你们应当感谢祂，祂藉我分施给你们；因为我分施的不是我自己的。例如，若你们问我，两人中哪一个先领受了——一个在一小时后领受，另一个在十二小时后？每个人都会回答说：在一小时后领受的人先于在十二小时后领受的人。因此，虽然他们在同一时辰领受，但因为有些人一小时后领受，另一些人十二小时后领受，那些在这么短时间内领受的人被说成是先领受的。最初的义人，如亚伯尔和诺厄，仿佛在第一个时辰蒙召，他们将与我们一起领受复活的福乐。在他们之后的义人，亚巴郎、依撒格、雅各伯以及他们同时代的所有人，仿佛在第三时辰蒙召，将与我们一起领受复活的福乐。其他义人，如梅瑟和亚郎，以及与他们一同的人，仿佛在第六时辰蒙召，将与我们一起领受复活的福乐。在他们之后，神圣的先知们，仿佛在第九时辰蒙召，将与我们一起领受同样的福乐。在世界末日，所有的基督徒，仿佛在第十一时辰蒙召，将与其余的人一起领受那复活的福乐。所有人都将一起领受；但想想那些最初的人，他们是在多久之后才领受的？如果那些最初的人在很久之后才领受，我们则在短暂之后领受；虽然我们都同时领受，但我们似乎先领受了，因为我们的工钱不会迟迟不来。\par

6. 在那工钱中，我们都将平等，在先的如同在后的，在后的如同在先的；因为那个银币是永生，而在永生中所有人都平等。因为虽然圣人因德行的差异而闪耀，有的更亮，有的稍逊；但就永生之恩赐而言，对所有人都是一样的。因为永生对这人不会更长，对那人更短，它同样是永恒的；那没有终点的，对你我没有终点。在那生命中，婚配的贞洁是一种样式，童贞的纯洁是另一种；善工的果实是一种样式，殉道的荣冠是另一种。这人是一种，那人是另一种；然而就永远活着而言，这人不会比那人活得更久，那人也不会比这人活得更久。因为他们同样无终地活着，尽管每个人都以自己的光辉活着；比喻中的银币就是永生。那么，在许久之后领受的人，不要抱怨在短暂之后领受的人。对前者，是偿还；对后者，是白白施予；然而同样的事物平等地给予两者。\par

7. 在此生也有相似的情形。除了比喻的那种解释，即被召于第一时辰的是指\Person{亚伯尔}和他那时代的义人，第三时辰的是指\Person{亚巴郎}和他那时代的义人，第六时辰的是指\Person{梅瑟}和\Person{亚郎}以及他们那个时代的义人，第十一时辰的则如同在世界末日，指所有基督徒；除了这个比喻的解释之外，这个比喻也可以在今生得到解释。因为那些刚出母胎就开始作基督徒的人，就好像是在第一时辰被召；孩童好像在第三时辰被召，青年在第六时辰，将近老年的人在第九时辰，而那些完全衰老的人就好像是在第十一时辰被召；然而所有这些人都要领取同一个永生的\OriginalTerm{德纳}{denarius}。\par

8. 但是，弟兄们，你们要听并且要明白，免得有人拖延进入葡萄园，因为他确信，无论何时来，他都必得到这个德纳。他确实确信德纳已应许给他，但这并非拖延的借口。因为那些被雇进葡萄园的人，当家主出来雇他所找到的人时——比如在第三时辰雇了他们——他们有没有对他说：\Quote{等等，我们要到第六时辰才去那里}？或者那些他在第六时辰找到的人，有没有说：\Quote{我们要到第九时辰才去}？或者那些他在第九时辰找到的人，有没有说：\Quote{我们要到第十一时辰才去？因为祂要同等赐予所有人，我们何必过度劳累自己呢？}祂要赐予什么、祂要做什么，都在祂自己旨意的奥秘中。你在被召时就来。因为同等的赏报应许给所有人；但关于这指定的工作时辰，却有一个重要的问题。例如，如果那些在第六时辰——即在那生命阶段，如同正午的酷热，感到壮年岁月的炽热——如果这些在壮年被召的人说：\Quote{等等，因为我们从福音中听到所有的人都将得到同样的赏报，我们要到第十一时辰，等我们老了再来，仍会得到同样的赏报。我们何必增加我们的劳苦呢？}那么他们必这样被回答：\Quote{你现在不愿劳力，却不知道是否能活到老年？你在第六时辰被召，来吧。家主确实应许你，如果你在第十一时辰来，必得一个德纳，但是否能活到第七时辰，却无人应许你。我且不说第十一时辰，甚至第七时辰也没有人应许你。那么，你为什么拖延那召你的人呢？你固然肯定赏报，却不确定时日。因此务要小心，以免祂因应许要赐给你的，你却因拖延而自己夺去。}如今，如果这话可以正当地对属于第一时辰的婴孩说，可以正当地对属于第三时辰的孩童说，可以正当地对处于生命强盛时期、如同第六时辰正午烈日的人说；那么，对衰老之人更应如何说呢？看，已经是第十一时辰了，你还站着，还迟迟不来吗？\par

9. 也许家主没有出去召你？如果祂没有出去，那么我们对你们的讲道有何意义？因为我们是祂家的仆人，奉命去雇工人。那么你为什么还站着？你已经满了你的年岁，赶快追求那德纳吧。因为家主的\Quote{出去}就是祂使自己为人所知；因为在屋里的人隐藏着，不被外面的人看见；但当祂\Quote{出了}屋子，就被外面的人看见。同样，基督在隐藏中，只要祂不被认识和承认；但当祂被承认时，祂就出去雇工人了。因为如今祂已从隐藏处出来，为使人认识祂：基督到处被认识，基督被宣讲；普天之下所有地方都大声宣扬基督的光荣。祂在犹太人中间几乎成了嘲笑和藐视的对象，祂显现为卑微并被轻视。因为祂隐藏了自己的威严，显示了自己的软弱。祂内那显现出来的被轻视，那隐藏的却不被认识。\BibleQuote{因为他们若认识祂，就不会钉死光荣的主了。}难道如今祂坐在天上还要被轻视吗？如果祂曾被轻视当祂悬挂在木架上时？那些钉死祂的人摇头，站在祂的十字架前，好像达到了他们残忍愤怒的果实，他们嘲笑说：\Quote{如果祂是天主子，就让祂从十字架上下来吧。祂救了别人，却不能救自己。}祂没有下来，因为祂隐藏着。因为对于祂来说，从十字架上下来远比从坟墓中复活更容易，因为祂有能力从坟墓中复活。祂显现了忍耐的榜样，为教导我们。祂延缓了自己的能力，而不被承认。因为那时祂还没有出去雇工人，祂出去了，但没有使自己被人认识。第三天祂复活了，显现给祂的门徒，升了天，并在复活后第五十天、升天后第十天派遣了圣神。所派遣的圣神充满了所有在一间屋子里的一百二十人。他们\BibleQuote{充满了圣神，并开始说起万国的方言}；如今呼召是明显的，如今祂出去雇工了。因为如今真理的能力开始被众人知晓。因为那时一个人领受了圣神，就能独自说起万国的方言。但如今在教会中，合一本身就如同一个人说起万国的方言。因为基督信仰没有到达什么语言呢？它没有延伸到什么边界呢？如今没有一个人\BibleQuote{能躲避它的热气}，而延迟仍被那些在第十一时辰站着的人冒险。\par

10. 那么，我亲爱的弟兄们，这是显而易见的，对所有人都是显而易见的，你们要牢牢记住并确信：每当有人从无益或放荡的生活方式，转向我主耶稣基督的信仰时，他过去的一切都被宽恕了，就像他所有的债务都被取消了一样，与他开始了新的账目。一切都完全被宽恕了。不要让任何人因担忧还有什么未被宽恕而焦虑。但另一方面，也不要让任何人安于一种错误的安稳。因为这两样是灵魂的死亡：绝望和错误的希望。因为正如正确而正直的希望拯救人，同样，错误的希望欺骗人。首先，思考绝望如何欺骗人。有些人，当他们开始反省自己所行的恶事时，以为自己不能被宽恕；而他们以为不能被宽恕时，立刻就把灵魂交付给毁灭，因绝望而丧亡，心中这样说：\Quote{如今我们没有希望了；因为我们所犯的如此大罪不能被赦免或宽恕；那么，我们何不满足我们的情欲呢？至少让我们充满现时的快乐吧，既然我们在将来没有赏报。让我们做我们喜欢的事，尽管不合法；这样我们至少可以享受暂时的欢乐，因为我们不能达到接受永生的赏报。} 说这些话时，他们因绝望而丧亡，要么在尚未相信之前，要么在已经成为基督徒后，因邪恶的生活而陷入各类罪与恶中。葡萄园的主走向他们，借厄则克耳先知敲门，在他们绝望中呼唤他们，而他们却背对着呼唤他们的那一位。\Quote{在任何人悔改之日，我要忘记他的一切过犯。} 如果他们听见并相信这话，他们就从绝望中被挽救，从那非常深且无底的深渊中再次升起，他们曾沉陷其中。\par

11. 但这些必须害怕，免得他们跌入另一个深渊，并因错误的希望而死，他们本不能死于绝望。因为他们改变了想法，这些想法虽然与以前大不相同，但同样有害，又开始在心里说：\Quote{如果在任何一天我离开了最邪恶的道路，仁慈的天主，正如祂借先知真实应许的，将忘记我的一切过犯，那么我为何要在今天回转，而不等到明天呢？让这一天像昨天一样在罪恶的享乐中、在放纵的洪流中度过，让它沉溺于致命的欢愉；明天我将\InnerQuote{回转}，一切就了结了。} 有人可能回答你：什么了结？你可能会说：我的过犯了结。很好，确实，喜乐吧，因为明天你的过犯将了结。但如果在明天之前你自己的结局就到来呢？那么你确实做得很好，为天主应许了你悔改时赦免你的过犯而喜乐；但没有人应许你明天。或者，如果碰巧某个占星家应许了明天，那与天主的应许大不相同。这些占星家欺骗了许多人，因为他们曾应许自己好处，却只找到了损失。因此，为了这些希望错误的人，家主再次出去。正如祂走向那些错误地绝望并迷失在绝望中的人，把他们呼唤回希望中；同样，祂也走向那些因邪恶希望而丧亡的人；并通过另一部书对他们说：\Quote{不要迟延归向上主。} 正如祂曾对其他人说过：\Quote{在任何人悔改之日，我要忘记他的一切过犯。} 并除去了他们的绝望，因为他们已将灵魂交付丧亡，绝望于任何方式的宽恕；同样，祂也走向那些有意因希望和迟延而丧亡的人；并对他们说话，斥责他们说：\Quote{不要迟延归向上主，不要一日一日地拖延；因为主的怒气会突然降临，在报复之日祂将毁灭你。} 因此，不要迟延，不要对现在敞开的门自己关闭。看，宽恕的施与者为你开门；你为什么迟延？你本该喜乐，即使祂在你叩门许久之后才开门；你还没有叩门，祂就开了门，而你却留在外面？不要迟延。圣经在某处论及慈悲的善工时说：\Quote{不要说：去吧，明天再来，我明天才给你；既然你能立即行善，因为你不知道明天会发生什么。} 那么，这里有一条不可迟延对他人慈悲的诫命，而你却要用迟延来残忍地对待自己吗？你不该迟延给面包，而你要迟延接受宽恕吗？如果你在怜悯他人时不迟延，\Quote{也要在取悦天主的事上怜悯自己的灵魂。} 也要施舍给你的灵魂。不，我不是说，给它施舍，而是不要推开那要赐予你的手。\par

12. 但人在害怕得罪他人时，常常极其严重地伤害自己。因为好友对行善大有影响，而恶友对行恶亦然。因此，主不愿首先拣选元老，却拣选渔夫，为教训我们为获得自己的救恩而藐视有权势者的友谊。啊，造物主何其显著的慈悲！因祂知道，若祂拣选了元老，元老会说：\Quote{我的品级被拣选了。}若祂首先拣选了富人，富人会说：\Quote{我的财富被拣选了。}若祂首先拣选了皇帝，皇帝会说：\Quote{我的权力被拣选了。}若拣选了演说家，他会说：\Quote{我的口才被拣选了。}若拣选了哲学家，他会说：\Quote{我的智慧被拣选了。}同时祂说：让这些骄傲的人暂且搁置吧，他们膨胀得太厉害了。实在的，实质的丰盈与肿胀之间有很大的区别；两者固然都大，但两者并非都是健康的。因此，祂说：让这些骄傲的人暂且搁置吧，他们必须以坚实之物来医治。首先，祂说：把这位渔夫给我。\Quote{来吧，你这贫穷者，跟随我；你一无所有，一无所知，跟随我。你这贫穷而无知的人，跟随我。你身上没有什么可令人敬畏，却有许多可被充满。}如此丰盈的源泉，应携带空器皿前来。于是渔夫舍弃了他的网，渔夫领受了恩宠，成为神圣的演说家。请看主做了什么，关于祂，宗徒说：\BibleQuote{天主拣选了世上软弱的，以叫那强壮的羞愧；又拣选了世上卑贱的，并且那些无有的，如同它们存在一样，为要废掉那有的。}因此，如今渔夫的话被诵读，演说家的颈项被降卑。愿一切虚风都被除去，愿那升起即散的烟雾被除去；当涉及这救恩时，让他们完全被藐视。\par

13. 倘若在某城中有人患了某种身体疾病，而那里有一位极高明的医生，却是病人有权势的朋友们的仇敌；倘若我所说的，在某城中有人正受某种危险身体疾病的折磨；而同一城中有一位极高明的医生，如我所说，是病人有权势的朋友们的仇敌，他们会对他们的朋友说：\Quote{不要叫他来，他什么都不知道；}他们说这话并非出于理智的判断，而是出于对他的厌恶；难道他不会为了自己的安全，摒除他有权势的朋友们的无端断言，而且无论冒犯他们多大，为能多活几天，去请那位医生来，就是众人口传为最擅长驱除他身体疾病的那位吗？看，整个人类都患病了，不是身体的疾病，而是罪恶。那里躺着一位巨大的病人，从东到西布满全世界。为了医治这位巨大的病人，全能的医生降临了。祂谦卑自下，甚至取了会死的肉身，犹如到了病人的床边。祂赐下健康的规诫，却被藐视；那些遵守的人获得拯救。祂被藐视，当有权势的朋友们说：\Quote{祂什么都不知道。}若祂什么都不知道，祂的能力怎能充满万民？若祂什么都不知道，祂怎能在与我们同在之前就已存在？若祂什么都不知道，祂怎会先派遣先知？那些古时预言的事，如今不是应验了吗？这位医生岂不是通过祂应许的实现，证明了祂医道的大能吗？致命的错谬不是被推翻了吗？靠打谷般的世界，情欲不是被制伏了吗？没有人可说：\Quote{从前的世界比现在好；自从那位医生开始施展祂的医术，我们在这里目睹许多可怕的事。}对此不必惊异？在任何人接受治疗之前，医生的住所似乎没有血迹；但现在，因着看见你所做的，抖落一切虚妄的快乐，来就医生吧，这是医治的时候，不是享乐的时候。\par

14. 那么，弟兄们，让我们思想得到医治吧。若我们还不认识这位医生，却至少不要像狂乱的人那样对祂施暴，或像嗜睡的人那样远离祂。因为许多人在这种狂暴中丧亡，许多人在沉睡中丧亡。狂乱的人就是因缺乏睡眠而疯狂的人。嗜睡的人就是被过多睡眠压垮的人。这两种人都可以找到。有些人存心对这位医生施暴，而由于祂自己坐在天上，他们就在地上迫害祂的忠信者。但就连这样的人，祂也医治。他们中许多人从仇敌转变为朋友，从迫害者转变为宣讲者。犹太人就是这样的人，虽然他们在祂在世时像狂乱的人一样对祂施暴，祂却治愈了他们，并当祂悬在十字架上时为他们祈祷。因为祂说：\BibleQuote{父啊，宽赦他们吧，因为他们不知道他们做的是什么。}然而，他们中许多人在暴怒平息后，仿佛狂乱被制服，认识了天主和基督。当圣神在升天后被派遣时，他们归向了他们所钉死的那一位，并作为信者饮下了他们在暴怒中所流的祂的血的圣事。\par

15. 我们在此有实例。扫禄迫害耶稣基督的肢体——基督如今坐在天上；他在狂乱中、在失去理智中、在疯狂的冲动中，严重地迫害了他们。但主从天上以一句话召唤他：\BibleQuote{扫禄，扫禄，你为什么迫害我？}击倒了狂乱者，扶起了痊愈者，杀死了迫害者，使宣讲者复活。同样，许多昏睡者也得了痊愈。因为那些不暴力反对基督、不恶意敌对基督徒、却只是在延误中因昏沉言语而迟钝沉重、迟迟不愿睁眼看光、并恼怒唤醒他们的人，正与这类人相似。沉重的昏睡者说：\Quote{离开我，我求你，离开我。}为什么？\Quote{我想睡觉。}但你会因此而死。他因贪睡而回答：\Quote{我想死。}但来自高天的爱呼喊道：\Quote{我不愿这样。}儿子常对年迈的父亲表现这种孝爱之情，尽管父亲数日后必离世，且已至风烛残年。若儿子见他昏睡，并从医生得知他患了昏睡病，医生告诉他：\Quote{叫醒你父亲，不要让他睡觉，否则你无法救他的命！}那时，儿子会走到老人身边，拍打、挤压、拧、刺，或让他感到任何不适，这一切都是出于对他的孝敬之情；他不让父亲立即死去——尽管他因年迈不久必将离世——若生命因此得救，儿子便因能与即将离世、为他让路的父亲多活几天而欢喜。那么，我们更应当以多大的爱德恳切劝勉我们的朋友？我们与他们一起，不只在今世活几天，而是在天主面前永远生活！让他们爱我们，听从我们的话，敬拜我们所敬拜的那一位，好使他们能领受我们所希望的一切。\Quote{让我们转向主，}等等。\par

\SourceRecord{Translated by R.G. MacMullen.；Nicene and Post-Nicene Fathers, First Series；Vol. 6.；Edited by Philip Schaff.；Buffalo, NY: Christian Literature Publishing Co.,；1888.；<http://www.newadvent.org/fathers/160337.htm>.}

% structure: p0001 source=h1 target=section reason=page-title

\section{新约讲道集 38}

\Emph{[LXXXVIII. Ben.]}\par

\Emph{论福音的话}，\BibleRef{玛 20:30}\Emph{，关于两个盲人坐在路旁，喊说：}\Quote{主，达味之子，可怜我们罢。}\par

圣洁的弟兄们，你们和我们同样都知道，我们的主和救主耶稣基督是我们永恒健康的医师；为此祂取了我们的本性的软弱，好使我们的软弱不再持久。因为祂接受了可死的身体，在其中杀死死亡。并且，\BibleQuote{祂虽然由于软弱而被钉在十字架上}，正如宗徒说，\BibleQuote{却因天主的德能而生活}。这也是同一宗徒的话；\BibleQuote{祂不再死，死亡不再统治祂}。这些事，我说，你们的信德都熟知。还有由此而来的这一点，即我们应该知道，祂在身体上所做的一切奇迹，都有益于我们的教导，使我们能从中领悟那不会过去、没有终结的事。祂使盲人重见那双死亡早晚要闭上的眼睛；祂使拉匝禄复活，而拉匝禄将来还要再死。祂为身体健康所做的任何事，都不是为了使它们永远健康；但在末日，祂将连身体本身也赐予永恒的健康。但因为那些看不见的事不被相信，祂便借这些看得见的事，建立起对那看不见之事的信德。\par

2. 所以，弟兄们，谁也不要说我们的主耶稣基督现在不行这些事，并因此将教会以前的世代置于现世之上。的确，在同一处，主更喜欢那些\BibleQuote{没有看见而相信}的人，胜过那些看见因而相信的人。因为那时祂门徒的软弱如此犹豫不决，以致他们认为他们所看见已复活的主必须被触摸，好使他们相信。他们的眼睛看见祂还不够，除非他们的手也触摸祂的肢体，并触摸祂最近伤口的疤痕；好使那位怀疑的门徒在触摸并认出疤痕时突然喊出：\BibleQuote{我的主，我的天主！}这些疤痕彰显了那曾治愈他人所有伤口的那位。难道主不能没有疤痕而复活吗？能，但祂知道门徒心中的伤口，为了医治他们，祂保留了自己身上的疤痕。主对那现在承认并说出\BibleQuote{我的主，我的天主！}的人说了什么？祂说：\BibleQuote{因为你看见了我，你才相信；那些没有看见而相信的人是有福的。}祂说的是谁，弟兄们，难道不是我们吗？并非只有我们，还有那些以后要来的人。因为不久之后，当祂从人的眼前离去，为要信德在他们心中建立起来，凡相信的，虽然未看见，却相信了，他们的信德功劳甚大；因为他们取得这信德，只凭一颗虔诚的心，而非手的触摸。\par

3. 这些事，主做了，为邀请我们进入信德。这信德如今在遍布全世界的教会中为王。现在祂施行更大的医治，为此祂那时也不屑于显示那些较小的。因为灵魂比身体更好，所以灵魂的救恩比身体的健康更好。失明的身体现在不再因主的奇迹而睁眼，但蒙蔽的心却因主的话而睁眼。可死的尸体现在不再复活，但灵魂却复活了，它曾活在活的身体里却已死去。身体的聋耳现在不再开启；但有多少人心中的耳朵关闭，却因天主穿透的话而敞开，以致那不信的人相信了，那恶生的人善生了，那不服从的人服从了；我们说：\Quote{这样的人成了信徒！}我们惊奇地听说那些我们曾经以为是硬心的人。那么，你为什么惊奇于一个现在相信、生活无罪并事奉天主的人？难道不是因为你看见他看见了，而你以为他是瞎眼的？看见他活着，而你以为他是死的？看见他听见了，而你以为他是聋的？因为要想到，有另一种意义上的死人，就是主曾对那个因想埋葬父亲而迟延跟随祂的人所说的：\BibleQuote{让死人去埋葬他们的死人。}当然这些埋葬死人的人身体并不死；因为若是那样，他们就不能埋葬死尸。但祂称他们为死人，何处？岂非在他们内的灵魂？正如我们常在一户人家中看到，房子本身健康完好，主人却躺卧死了；同样，在健康的身体里，许多人携带一个死去的灵魂；宗徒唤醒他们时说：\BibleQuote{你这睡眠的，醒来，从死者中起来，基督必要光照你。}那正是那给瞎子光明、唤醒死人的同一位。因为正是用祂的声音，宗徒向死人呼喊：\BibleQuote{你这睡眠的，醒来。}当瞎子复活时，他将被光照。主说：\BibleQuote{有耳听的，就听吧！}当时祂眼前有多少聋子？因为谁站在祂面前却没有肉体的耳朵？祂所求的，岂非是内心者的耳朵吗？\par

4. 再者，当祂对那些确实看见、却只用肉眼看见的人说话时，祂寻求的是怎样的眼睛？因为当斐理伯对祂说：\Quote{主，请将圣父显示给我们，我们就满足了；}他确实明白，如果将圣父显示给他，那足以使他满足；但那与圣父同等的一位都不能使他满足，圣父又怎能使他满足呢？为什么祂不能使他满足？因为祂没有被看见。为什么祂没有被看见？因为那藉以看见祂的眼睛尚不健全。因为主在肉身中被肉眼看见，这不仅尊敬祂的门徒看见了，连钉死祂的犹太人也看见了。因此，那愿意以另一种方式被看见的，就寻求其他的眼睛。所以，当那人对祂说：\Quote{请将圣父显示给我们，我们就满足了；}祂回答说：\BibleQuote{我跟你们在一起这么长久，你们还不认识我吗？斐理伯！谁看见了我，就是看见了圣父。}同时，祂为了医治信德的眼睛，先给了他关于信德的教导，以便他达到看见。免得斐理伯以为他应当以当时在肉身中看见主耶稣基督的同样形状来想象天主，祂立刻接着说：\BibleQuote{你不信我在圣父内，圣父在我内吗？}祂早已说过：\BibleQuote{谁看见了我，就是看见了圣父。}但斐理伯的眼睛还不够健全，不能看见圣父，因而也不能看见与圣父同等的圣子。所以耶稣基督着手医治，并用信德的药物和药膏来坚固他心灵的眼睛，这眼睛尚微弱，不能注视如此伟大的光明，祂说：\BibleQuote{你不信我在圣父内，圣父在我内吗？}那么，那还不能看见主有一天要向他显示之事的人，不要先寻求看见他应当相信的；而要首先相信，以便那藉以看见的眼睛得到医治。因为只有奴仆的形状显示给奴仆的眼睛；因为如果\BibleQuote{祂虽具有天主的形体，却没有以自己与天主同等为应当把持不舍的}，现在能够被祂愿意医治的人看见与天主同等，祂就不需要\BibleQuote{使自己空虚，取了奴仆的形状}。但因为除了人能够被看见的方式外，没有方式能使天主被看见，所以那本是天主的成了人，以便那被看见的医治那藉以看见祂的。因为祂在别处亲自说：\BibleQuote{心里洁净的人是有福的，因为他们要看见天主。}斐理伯当然可以回答说：\Quote{主，看，我看见了你；圣父就像我看见你那样吗？因为你曾说：\InnerQuote{谁看见了我，就是看见了圣父}？}但在斐理伯这样回答之前，或者甚至在他这样想之前，当主说：\BibleQuote{谁看见了我，就是看见了圣父；}祂立刻加上：\BibleQuote{你不信我在圣父内，圣父在我内吗？}因为用那只眼睛，他既不能看见圣父，也不能看见与圣父同等的圣子；但为了他的眼睛得到医治而能看见，他需要被傅油以相信。所以，在你看见你现在不能看见的事之前，要相信你尚未看见的事。\BibleQuote{凭信德往来}，好使你达到看见。看见不会使那人在家中喜悦，如果信德在路上不给他安慰。因为宗徒这样说：\BibleQuote{我们几时住在这肉身内，就是与主远离。}他立刻加上为什么我们仍然\Quote{远离}，虽然我们已经相信；他说：\BibleQuote{因为我们是凭信德往来，}他说，\BibleQuote{并非凭目睹。}\par

5. 我们的全部事业，弟兄们，在此生就是医治这心灵的眼睛，藉以看见天主。为此举行神圣奥迹；为此宣讲天主的圣言；为此有教会的道德劝诫，即那些关乎矫正品行、克制肉欲、弃绝世俗的劝诫，不仅是在言语上，更是在生活的改变上：为此一切神圣圣经的旨趣都指向那内在的人得以净化，除去那妨碍我们看见天主的障碍。因为正如那被造来看见这暂时光明的眼睛——虽然这光是天上的、物质的、明显的，不仅对人，甚至对最卑贱的动物也是如此（因为眼睛被造就是为了看见这光）——如果有东西丢入或落入其中，使它混乱，它就被拒于这光之外；虽然光以其临在包裹它，眼睛却转离它，远离它；并且因其混乱状态，它不但对临在的光视而不见，甚至那为看见而被造的光也成为它的痛苦。同样，心灵的眼睛，当它混乱受伤时，也转离义德之光，不敢也不能注视它。\par

6. 是什么扰乱了心灵之眼？是邪恶的欲望、贪婪、不义、世俗的情欲，这些扰乱了、封闭了、蒙蔽了心灵之眼。然而，当身体之眼出了毛病时，人是怎样求医，毫不迟延地去打开、洗净它，以便那能看见外界光明的眼得以痊愈！人们来回奔走，没有一个人停歇，没有一个人耽搁，就算最细小的灰尘落入眼中也是如此。而天主必定创造了太阳，我们渴望用健康的眼睛去看它。创造太阳的那一位无疑更加光明；而且心灵之眼所关注的光明根本不属于这一类。那光明是永恒智慧。人啊，天主按自己的肖像造了你。祂给了你看见祂所造的太阳的能力，难道不给你看见创造你的那一位的能力吗？祂按自己的肖像造了你，也给了你这种能力。两者祂都给了你。但你太爱这外在的眼睛，太轻视那内在的眼睛；你带着它，它却受了伤、有了创伤。是的，如果你的造物主愿意向你显示自己，那对你将是一种惩罚；在它被治愈和痊愈之前，对你的眼睛是一种惩罚。因为亚当在乐园中犯了罪，就躲避天主的面容。因此，当他拥有纯洁良心的健康之心时，他因天主的临在而喜乐；当那眼睛因罪而受伤时，他开始惧怕神光，逃回黑暗和茂密的树丛中，逃避真理，寻求荫庇。\par

7. 因此，我的弟兄们，既然我们也从他而生，正如宗徒所说：\BibleQuote{在亚当内众人都死了}\BibleRef{格前 15:22}；因为如果我们当初愿意听从医生，我们本可以不生病；那么现在，让我们听从祂，好从疾病中得救。医生在我们健康时给了我们诫命，祂给了我们诫命，好使我们不需要医生。祂说：\BibleQuote{健康的人不需要医生，而是有病的人}\BibleRef{玛 9:12}。当我们健康时，我们藐视了这些诫命，并且亲身经验到，藐视祂的诫命是如何导致我们的毁灭。现在我们病了，我们痛苦，我们躺在软弱的床上；但让我们不要绝望。因为我们不能到医生那里去，祂却屈尊亲自来到我们这里。尽管人在健康时藐视了祂，祂却没有在人跌倒时藐视他。祂没有停止给软弱的人其他的诫命，这些人当初不肯遵守最初的诫命以免软弱；好像祂在说：\Quote{你确实从经验中感受到，当我说\InnerQuote{不要碰这个}时，我说的是真话。那么现在，终于痊愈吧，恢复你失去的生命。看，我正承担你的软弱；喝下这苦杯。因为你凭自己使我当初在健康时赐给你的那些甜蜜的诫命变得如此艰难。它们被藐视了，于是你的痛苦开始；除非你喝下苦杯，否则你无法痊愈。这苦杯就是诱惑之杯，就是今生充满的苦难之杯，就是患难、痛苦和折磨之杯。那么喝吧，}祂说，\Quote{喝吧，好使你活下去。}为了使病人不回答：\Quote{我不能，我不能忍受，我不喝；}这位医生，尽管完全健康，却先喝了，好使病人不再犹豫去喝。因为在这杯里有什么苦味是祂没有喝过的呢？如果是侮辱，祂在驱魔时首先听到了：\BibleQuote{祂附有魔鬼，赖贝耳则步驱魔}\BibleRef{玛 12:24}。因此，为了安慰病人，祂说：\BibleQuote{如果人家称家主为贝耳则步，何况他的家人呢？}\BibleRef{玛 10:25}。如果痛苦是这苦杯，祂被捆绑、鞭打、钉在十字架上。如果死亡是这苦杯，祂也死了。如果软弱对某种死亡方式特别恐惧，那么当时没有比十字架之死更耻辱的了。因为宗徒在陈述祂的服从时并非无故加上：\BibleQuote{服从至死，且死在十字架上}\BibleRef{斐 2:8}。\par

8. 但是，因为祂有意在世界的终末荣耀祂的信徒，祂首先在这世上荣耀了十字架；以至于相信祂的世上君王禁止任何罪犯被钉十字架；而犹太迫害者以极大的嘲笑为主预备的十字架，如今祂的仆人君王们却以极大的信心佩戴在额上。唯独我们的主屈尊为我们所受的死的耻辱性质如今不那么明显了，祂正如宗徒所说：\BibleQuote{为我们成了诅咒}\BibleRef{迦 3:13}。当祂被悬挂时，犹太人的盲目嘲笑祂，当然祂本可以从十字架上下来，因为如果祂不愿意，祂就不会在十字架上；但从坟墓中复活比从十字架上下来更伟大。因此，我们的主在行这些神性的事和受这些人性的事时，借着祂身体的奇迹和身体的忍耐教导我们，使我们相信，并得以痊愈，去看见那些身体之眼无法知道的不可见之事。正是为此，祂治好了这些盲人，福音中刚才读到了关于他们的记述。请思考，借着他们的痊愈，祂给内在有病的人带来了怎样的教导。\par

9. 请思考此事的结果与状况的顺序。那两个坐在路旁的盲人，当主经过时，喊叫说，求主怜悯他们。但与主同行的群众却阻止他们喊叫。切莫以为这一细节没有奥秘的含义。但他们凭着呼喊的极大恒心，克服了阻拦他们的群众，使自己的声音能达于主耳；仿佛主并未预知他们的意念。于是那两个盲人喊叫，好让主听见，且不能被群众阻止。主\BibleQuote{经过}，他们喊叫；主\BibleQuote{站住}，他们便痊愈了。因为\BibleQuote{主耶稣站住，叫了他们来，说：\InnerQuote{你们要我为你们做什么？}他们对他说：\InnerQuote{主啊，要我们的眼睛能看见。}}主照他们的信德行事，恢复了他们的视力。若我们现在明白了病者、聋者、死者，内在的病者、聋者、死者，那么也让我们在此处寻找内在的瞎子。心灵的眼睛是闭着的；\BibleQuote{耶稣经过}，好让我们喊叫。什么是\BibleQuote{耶稣经过}？耶稣所行的事只持续一时。什么是\BibleQuote{耶稣经过}？耶稣所行的事是会过去的。请看，祂有多少事已经\BibleQuote{过去了}。祂由童贞玛利亚诞生，难道祂永远在诞生吗？祂曾如婴儿被哺乳，难道祂永远被哺乳吗？祂经历了生命的各个阶段直到成年，难道祂的身体永远在成长吗？童年继婴儿之后，少年继童年之后，成年继少年之后，经过多次的过渡。就连祂所行的奇迹也已经\BibleQuote{过去了}，它们被阅读并相信。因为这些奇迹被记载下来以便阅读，当它们发生时就已\BibleQuote{过去了}。总之，不再赘言，祂曾被钉十字架：难道祂永远挂在十字架上吗？祂曾被埋葬，祂复活了，祂升到天上；\BibleQuote{如今祂不再死，死亡再也不能主宰祂了。}祂的天主性永存，是的，祂身体的不朽如今永不失效。但尽管如此，祂在时间中所行的一切事都已经\BibleQuote{过去了}；它们被记载下来供阅读，被宣讲以供相信。在这一切事中，\BibleQuote{耶稣经过}。\par

10. 那\BibleQuote{两个坐在路旁的盲人}，不就是耶稣来医治的那两个民族吗？让我们在圣经中指明这两个民族。福音中记载：\BibleQuote{我另外有羊，不是这圈里的；我必须领他们来，他们也要听我的声音，并且要合成一群，归一个牧人。}那么这两个民族是谁？一个是犹太人，另一个是外邦人。\BibleQuote{我奉差遣，}祂说，\BibleQuote{不过是到以色列家迷失的羊那里去。}祂对谁说了这话？对门徒；当时那自认是狗的客纳罕妇人喊叫，求主使她配得从主人桌子上掉下来的碎屑。因她被认为配得，如今祂所来的两个民族便显明了：即犹太民族，就是祂所说的\BibleQuote{我奉差遣，不过是到以色列家迷失的羊那里去}的那个民族；以及外邦民族，这妇人正是其预表，祂起初拒绝她说：\BibleQuote{不好拿儿女的饼丢给狗吃。}而当她说：\BibleQuote{主啊，不错，但是狗也吃它们主人桌子上掉下来的碎屑。}祂回答：\BibleQuote{妇人，你的信德真大！照你所要的，给你成就了吧。}同样，那百夫长也属于这个民族，主曾论到他：\BibleQuote{我实在告诉你们，这么大的信德，就是在以色列中，我也没有遇见过。}因为他说过：\BibleQuote{主啊，你到我舍下，我不敢当；只要你说一句话，我的仆人就必好了。}因此，主在受难和受光荣之前，就已经指出了两个民族：一个是祂因对祖先的应许而来到的民族；另一个是祂因慈悲而不拒绝的民族；为要实现给亚巴郎的应许：\BibleQuote{地上万国都要因你的后裔蒙福。}因此，宗徒在主复活升天之后，被犹太人藐视时，便转向外邦人。但这并不意味着他向由犹太信徒组成的教会缄默；他说：\BibleQuote{我本来在犹太的各基督教会内是未曾谋面的。不过他们只是听说：那从前迫害我们的，如今却传扬他原先所破坏的信仰，他们就因我光荣天主。}同样，基督被称为\BibleQuote{使两者合一的角石。}因为角石连接两面从不同方向来的墙。还有什么比割损与未割损更不同呢？一面墙来自犹太，另一面来自外邦。但她们被角石联合起来。\BibleQuote{匠人所弃的石头，已成了屋角的基石。}在建筑中，除非两面从不同方向来的墙相遇并以某种合一相连接，否则就没有角石。那么，那\BibleQuote{两个盲人}向主呼喊，按照预表，就是这两面墙。\par

11. 现在，亲爱的诸位，请留意。主正\Quote{路过}，瞎子们便\Quote{喊叫}。\Quote{路过}是什么意思？正如我们已说过的，祂所行的事工都是\Quote{路过}的。如今，我们的信德正是建立在这些路过的事工之上。因为我们相信天主之子，不仅相信祂是天主的圣言，万物借祂而造；因为如果祂一直保持\BibleQuote{具有天主的形体，与天主同等}，而没有\BibleQuote{空虚自己，取了奴仆的形体}，瞎子们甚至不会看见祂，好能喊叫。但祂行路过的事工时，即\BibleQuote{祂贬抑自己，服从至死，且死在十字架上}，那\BibleQuote{两个瞎子就喊说：达味之子，可怜我们吧！}因为正是这件事——祂，达味的主和造物主，也愿意成为达味之子——祂在时间中行了，祂\Quote{路过}行了。\par

12. 弟兄们，所谓\Quote{向基督喊叫}，岂不是以善行回应基督的恩宠吗？我说这话，弟兄们，免得我们徒然用声音高喊，在生活中却默不作声。谁向基督喊叫，好使他的内心盲瞎能被\Quote{路过}的基督驱除呢？即祂正将那暂时的圣事分施给我们，借此我们受教以领受永恒之事。谁向基督喊叫？凡轻视世俗的，就是向基督喊叫。凡轻视世俗享乐的，就是向基督喊叫。凡不用舌头而用生活说\BibleQuote{世界于我已被钉在十字架上，我于世界也被钉在十字架上}的，就是向基督喊叫。凡\BibleQuote{散财而施舍，他的仁义永存}的，就是向基督喊叫。因为：让那听见且不聋于声音的人，以\BibleQuote{变卖你们所有的，施舍给穷人；为自己准备不朽的钱囊，在天上备下取不尽的宝藏}为诫命；让他听见这如同基督\Quote{路过}的脚步声时，便在其瞎盲中以喊叫响应，也就是说，让他行这些事。让他的声音在于行动。让他开始轻视世俗，将财产分施给穷人，视他人所贪爱的为无物，不介意伤害，不求报复，让他\BibleQuote{连左颊也转给}打他的人，让他为他的仇人祈祷；若有人\BibleQuote{夺去了他的东西，让他不要再索回}；若他从任何人取过什么东西，\BibleQuote{让他偿还四倍}。\par

13. 当他开始行这一切时，他的亲属、亲戚和朋友都会骚动。那些爱现世的人，会反对他。这是何等疯狂！你太过极端了：什么！难道其他的人不是基督徒吗？这是愚拙，这是疯狂。群众喊出诸如此类的话，阻止瞎子喊叫。群众斥责他们喊叫，却未能压制他们的喊声。凡渴望得医治的人，要明白他们该做什么。耶稣如今也正\Quote{路过}；让那些在路旁的人喊叫。这些就是那些\BibleQuote{用嘴唇尊敬天主，心却远离祂}的人。这些人是在路旁，耶稣像给心里盲瞎的人一样赐予他们诫命。因为每当叙述耶稣所行的那些路过之事时，耶稣常被呈现给我们为\Quote{路过}的。因为直到世界终末，不会缺少\Quote{坐在路旁的瞎子}。所以，坐在路旁的他们有必要喊叫。与主同在的群众会压制那些寻求痊愈者的喊声。弟兄们，你们明白我的意思吗？因为我不知道如何说话，但更不知道如何缄默。那么我就要说话，且说得明明白白。因为我害怕\Quote{耶稣路过}和\Quote{耶稣站住}，因此我不能保持沉默。恶且冷漠的基督徒阻碍那些真正热忱、愿意遵行福音所载天主诫命的善良基督徒。这同主同在的群众阻碍那些喊叫的人，即阻碍那些行善的人，以免他们因持守而得医治。但要让他们喊叫，不要气馁；不要像被多数人的权威所误导；不要效法那些比他们先成为基督徒却过着恶行且嫉妒他人善行的人。不要让他们说：\Quote{让我们像这许多人那样生活。}为何不更按福音的规定生活呢？为什么你们愿意依从那阻碍你们的群众的劝诫而活，却不跟随\Quote{路过}的主的脚步？他们会嘲笑、辱骂、把你们叫回来；你们要喊叫，直到耶稣的耳朵听见。因为那些坚持做基督所命之事，不理会阻碍他们的群众，也不看重自己似乎是跟随基督的（即被称为基督徒），却爱基督即将恢复给他们的光明胜于畏惧阻碍者的喧嚷的人，他们绝不会与祂分离，耶稣会\Quote{站住}，并使他们痊愈。\par

14. 我们的眼睛如何得以痊愈？当我们藉信德在时间中的救恩计划里领悟基督\Quote{经过}时，我们就能在祂不变永恒中认识祂\Quote{停住}。因为当对基督神性的认识达到时，眼睛就痊愈了。愿你们的爱德领悟这一点；请留意我要讲述的伟大奥秘。主耶稣基督在时间中所行的一切，都在我们内培植信德。我们信从天主子，不单是信那\BibleQuote{万物藉祂而造的}圣言，更是信这\BibleQuote{成了血肉、居住在我们中间的}圣言，祂由童贞玛利亚诞生，以及信经所包含的其他内容，这些向我们启示，为使基督\Quote{经过}，使瞎子听到祂\Quote{经过}的脚步声，藉他们的作为\Quote{呼喊}，以他们的生活彰显信德的宣认。但如今，为使那些呼喊的人得以痊愈，\Quote{耶稣停住了}。因为那说\BibleQuote{我们纵然曾按血肉认识基督，但如今不再这样认识祂了}的人，看见了耶稣如今\Quote{停住}。因为他看见了基督的神性，在现世生活中尽其可能。所以在基督内有神性和人性。神性\Quote{停住}，人性\Quote{经过}。神性\Quote{停住}是什么意思？它不改变、不被动摇、不离去。因为祂来到我们这里，并没有离开圣父；祂升天，也没有改变地方。当祂取了肉身时，改变了地方；但取了肉身的天主，既然不在地方内，就不改变地方。让我们被\Quote{停住}的基督触摸，我们的眼睛就痊愈了。但谁的眼睛？那些在祂\Quote{经过}时\Quote{呼喊}的人的眼睛；即那些藉着在时间中分施、为教导我们幼年时期的信德，而行善工的人。\par

15. 还有什么比痊愈的眼睛更珍贵呢？他们欢喜看见这从天上照耀的受造之光，甚至灯所发出的光。那些不能看见这光的人，显得多么不幸！但我为何说这些，谈论这一切，无非是劝勉你们在耶稣\Quote{经过}时都要\Quote{呼喊}。我高举这光——也许你们尚未看见——作为你们爱慕的对象，圣善的弟兄们。你们虽然还未看见，但要相信；并要\Quote{呼喊}，好使你们看见。那些看不见这肉身之光的人，他们的不幸被认为是多么大！有人失明了，立刻就说：\Quote{天主向他发怒，他做了恶事。}多俾亚的妻子就这样对她的丈夫说。他因那只山羊羔而呼喊，怕它是偷来的；他不愿家里有任何偷窃之物的声音；而她坚持自己所做的，辱骂她的丈夫；当他说：\BibleQuote{若是偷来的，就还回去}时，她侮辱地回答说：\BibleQuote{你的义行在哪里？}她坚持偷窃，她的盲目是多么大！而命令归还偷窃之物的人，看见了多么明亮的光！她在外表上因太阳的光而喜乐，他在内里因正义的光而喜乐。他们谁处在更好的光中呢？\par

16. 亲爱的诸位，我要劝勉你们爱慕这光：当主\Quote{经过}时，你们要藉着作为\Quote{呼喊}；让信德的声音发出，使\Quote{停住的耶稣}，即不变、永存的天主智慧，以及\BibleQuote{万物藉祂而造的}天主圣言的尊威，开启你们的眼睛。那位多俾亚在劝告他儿子时，教导他要呼喊；即教导他行善工。他告诉他要施舍穷人，嘱咐他周济贫苦的人，并教训他说：\BibleQuote{我儿，施舍救人不至于进入黑暗。}这位盲人劝告人接受并获得光明。他说：\BibleQuote{施舍救人不至于进入黑暗。}倘若他儿子惊讶地问他：父亲，你难道没有施舍吗，如今你在眼瞎中对我说话？你岂不在黑暗中，却对我说\BibleQuote{施舍救人不至于进入黑暗}？不，他深知那光是什么，他就是按此教导他儿子的；他深知在内里看到了什么。儿子向父亲伸出手，使他在世上行走；而父亲向儿子伸出手，使他能居住在天上。\par

17. 简而言之，弟兄们，让我以一个与我息息相关、令我极为痛心的事来结束这篇讲道。请看，有多少群众在\Quote{斥责那呼喊的瞎子}。但愿你们中间凡渴望获得治愈的人，不要被他们吓倒；因为有许多人徒有基督徒之名，行为却是不敬天主的；不要让他们阻止你们行善。在那些阻拦你们、叫你们回头、辱骂你们、生活邪恶的群众中间，你们要呼喊。因为坏基督徒不仅用言语，也用恶行压制善人。好基督徒不愿参加公开表演。正是因为他克制自己去看戏的欲望，他就是呼求基督，呼求得到治愈。别人都涌向那里，但也许他们是外教人或犹太人？啊！的确，如果基督徒不去看戏，那里的人就会少得可怜，他们会因羞愧而离去。所以基督徒也往那里跑，只是带着圣名走向自己的定罪。那么，你们要藉着不去看戏、在内心抑制这种世俗的贪欲来呼喊；以坚强而持久的呼喊，紧抓住救主的耳朵，使耶稣\Quote{站住}并治愈你们。就在群众中间呼喊，不要绝望于达到主的耳朵。因为福音中的瞎子并非在没有人群的地方呼喊，好让他们能在没有阻碍的方向被听见。他们就在群众中间呼喊，主仍然听见了他们。同样，你们也要在罪人、纵欲者、贪爱世俗虚荣的人中间呼喊，使主能治愈你们。不要到别处去向主呼喊，不要到异端者那里去向他呼喊。弟兄们，试想，在那阻拦他们呼喊的群众中，那呼喊的人正是得了痊愈。\par

18. 圣洁的弟兄们，也要注意这一点：坚持呼喊意味着什么。我要讲一讲许多人和我自己在基督名下所经历的事；因为教会不断产生这样的人。当一个基督徒开始善度生活、热心行善、轻视世俗时，在这新生活中，他会受到冷淡基督徒的非议和反对。但如果他坚持，并以忍耐胜过他们，在善行上不松懈，那么先前阻拦他的人现在反而会尊敬他。因为他们责备、阻拦、反对他，只要他们还有希望他会屈服。但如果进步的人以他们的坚持胜过了他们，他们就转变态度，开始说：\Quote{他是个伟人，圣人，天主赐给他如此恩宠，他是幸福的。}现在他们尊敬他，祝贺他，祝福并赞美他；正如那与主同在的群众所做的那样。他们起初阻拦瞎子，不让他们呼喊；但当瞎子继续呼喊直到被听见、获得主的怜悯时，那群众就说：\Quote{耶稣叫你呢。}而刚才还\Quote{斥责他们，叫他们不作声}的人，现在却用劝勉的话了。只有那不在世上劳苦的人才不被主呼召。但在今生，谁不是因自己的罪恶和不义而劳苦呢？但如果众人都劳苦，那就对众人说：\Quote{凡劳苦和负重担的，你们都到我这里来。}既然是对众人说的，你为什么把失败归咎于那这样邀请你的主呢？来吧。他的家并不狭窄；天主的国度为众人平等地拥有，也为每人全部拥有；它不因拥有者增多而减少，因为它不分。多人同心所拥有的，为每人都是完整无缺的。\par

19. 然而，弟兄们，从这段经文的神秘意义中，我们认出了在圣经其他部分最明确表达的内容：教会内有善有恶，正如我常说的，麦子和糠秕。没有人应在时候未到之前离开禾场，他应在打谷时容忍糠秕，在禾场上容忍它。因为在仓里，他就无需容忍了。那簸扬者将要来，他将把坏人与好人分开。那时也将有身体的分离，而现在的精神分离已先发生。在内心要常与坏人分离，在身体上暂时与他们联合，但须谨慎。然而，不要疏忽纠正那些属于你的人、那些以某种方式属于你照管的人，要以劝告、教训、鼓励或警告来纠正。尽你所能去做。因为你在圣经和圣人的榜样中（无论他们是在主降临今生之前或之后生活的）发现，坏人与好人合一并不玷污好人，所以不要因此变得在纠正坏人上松懈。坏人在两方面不会玷污你：如果你不赞同他，并且如果你责备他；这就是不与他共融、不赞同他。因为当意志或赞同的同意与他的行为结合时，就有共融。这是宗徒教导我们的，他说：\Quote{不要与黑暗无益的作为共融}。又因不赞同是小事，如果伴随以纠正上的疏忽，他就说：\Quote{但更要指摘他们}。看，他如何同时包含了两者：\Quote{不要共融，但要指摘他们}。什么是\Quote{不要共融}？不赞同他们，不称赞他们，不认可他们。什么是\Quote{但更要指摘他们}？指摘、斥责、抑制他们。\par

20. 但在纠正和制止他人的罪过时，必须谨慎，以免在责备他人时抬高自己；应牢记宗徒的那句话：\BibleQuote{所以，凡自以为站得稳的，务要小心，免得跌倒。}让责备的声音在外表上带着威吓，但内心保持仁爱和温柔的精神。\BibleQuote{如果一个人偶然被过犯所胜，}正如同一宗徒所说，\BibleQuote{你们属灵的人当用温柔的心把他挽回过来；又当自己小心，恐怕也被引诱。你们各人的重担要互相担当，如此，就完成了基督的法律。}又在另一处说：\BibleQuote{主的仆人不可争竞，只要温温和和地待众人，善于教导，存心忍耐，用温柔劝戒那抵挡的人；或者天主给他们悔改的心，可以明白真道，叫他们这已经被魔鬼任意掳去的，可以醒悟，脱离他的网罗。}因此，既不可赞同恶事而认可它，也不可疏忽而不责备；更不可骄傲，以侮辱的口气责备。\par

21. 但谁若离弃合一，便破坏了爱德；谁若破坏了爱德，无论他拥有多大的恩赐，都算不得什么。\BibleQuote{我若能说人间的语言，并能说天使的语言；我若有先知之恩，又明白一切奥秘和各种知识；我若有全备的信德，甚至能移山；我若舍己身叫人焚烧，却没有爱德，于我毫无益处。}他拥有一切却无益，因为他没有那能使一切善用的一件东西。因此，让我们拥抱爱德，\BibleQuote{以和平的联系，努力保持心神的合一。}不要被那些按肉体理解圣经的人欺骗，他们制造身体上的分裂，却因精神上的亵渎而与散布全世界的教会好种子分离。因为好种子是撒在全世界的。那位好撒种者，人子，不仅在非洲，而且到处撒下好种子。但仇敌却在其中撒了莠子。然而家主说什么？\BibleQuote{让两样一起长到收割的时候。}在哪里长？当然是在田里。田是什么？是非洲吗？不！那么是什么？我们不要自己解释，让主来说；我们不许任何人凭己意猜测。因为门徒对老师说：\BibleQuote{请把莠子的比喻讲给我们听。}主就讲明了：\BibleQuote{好种子}，祂说，\BibleQuote{就是天国之子。莠子就是那恶者之子。}谁撒的？\BibleQuote{那撒莠子的仇敌}，祂说，\BibleQuote{就是魔鬼。}田是什么？\BibleQuote{田}，祂说，\BibleQuote{就是这世界。}收割是什么？\BibleQuote{收割}，祂说，\BibleQuote{就是世界的终结。}收割的人是谁？\BibleQuote{收割的人}，祂说，\BibleQuote{就是天使。}非洲就是世界吗？现在是收割的时候吗？\Person{多纳图}是收割的人吗？那么，要在全世界寻找收割，要在全世界长到收割，要在全世界容忍莠子直到收割。不要让那些邪僻的人欺骗你们，那些轻飘飘的糠秕，在簸扬者来临之前就从禾场上飞走；不要让他们欺骗你们。你们要紧紧抓住这个关于莠子的单一比喻，不容他们谈论别的。这个人，有人会说，他交出了圣经；不，不是这样；但另一个人交出了圣经。无论谁交出了圣经，难道他们的不信就废掉了天主的信实吗？什么是\Quote{天主的信实}？就是祂向\Person{亚巴郎}所应许的：\BibleQuote{地上万国都要因你的后裔而蒙福。}什么是天主的信实？\BibleQuote{让两样一起长到收割的时候。}在哪里长？在整个田里。什么叫整个田里？就是整个世界。\par

22. 他们在这里说：\Quote{诚然，两种麦子曾一度生长于全世界，但好麦子减少了，并局限于我们的这个国家和我们小小的团体。} 但主不允许你按自己的意愿解释。那亲自解释这比喻的主，堵住了你的口——你这亵渎、不敬、邪恶的口，与你自身的利益相悖，而你悖逆遗嘱者，如同祂召唤你承受产业。祂如何堵住你的口？祂说：\BibleQuote{让两者一起生长直到收获。} 如果收获已经来到，让我们相信麦子减少了。然而即使在那时，它也不会减少，而是被收进仓里。因为祂如此说：\BibleQuote{你们先把莠子捆成捆，留着焚烧，却把麦子收入我的仓里。} 如果它们一直生长直到收获，并且在收获后被收进去，它们怎么会减少呢，你这邪恶、不敬的人？我承认，与莠子和糠秕相比，麦子在数量上较少；但\BibleQuote{两者一起生长直到收获。} 因为\BibleQuote{由于罪恶的增多，许多人的爱渐渐冷淡了；} 莠子和糠秕增多了。但在全世界，麦子不会缺乏，那些\BibleQuote{坚持到底的，必然得救}，两者一起生长直到收获。若因恶人之多，经上说\BibleQuote{当人子来临时，你以为祂会在世上找到信德吗？} 此名称指代所有因违反法律而效仿那被说\BibleQuote{你是土，还要归于土} 之人的；但亦因善人之多，并因那被说\BibleQuote{你的后裔必像天上的星和海边的沙} 的，经上也记着说：\BibleQuote{将有许多人从东方和西方来，同亚巴郎、依撒格在天国里坐席。} 因此\BibleQuote{两者一起生长直到收获}，莠子和糠秕在圣经中有它们的段落，麦子也有。那些不理解它们的人，混淆了它们并使自己困惑；在盲目欲望中，他们制造如此大的喧嚷，以致真理的明确彰显也不能使他们安静。\par

23. 看，他们说，先知说：\BibleQuote{你们离开吧，从那里出去，不要触摸不洁之物；} 那么，为了和平，我们怎能容忍恶人呢？我们不是被命令\BibleQuote{出去、离开，以免触摸不洁之物} 吗？我们灵性地理解那\Quote{离开}，他们却物质地理解。因为我也与先知一同呼喊（无论我是何等卑微的器皿，天主仍用我来服务你们）；我也呼喊并对你们说：\BibleQuote{你们离开吧，从那里出去，不要触摸不洁之物；} 但这是心的触摸，不是身体的。因为\Quote{触摸不洁之物} 是什么？不就是同意犯罪吗？\Quote{从那里出去}是什么？不就是去做那属于指责恶人的事，尽可能根据各人的等级和身份，在保持和平的情况下进行？你对一个人的罪感到不满，你就没有\Quote{触摸不洁之物}。你责备、斥责、劝诫了他，若情况需要，你施行了适当的管教，且不破坏合一；这样你就\Quote{从那里出去了}。现在思想圣人们的作为，免得这似乎是我自己的解释。正如圣人们理解了这些话，我们也应当如此理解。先知说：\BibleQuote{你们从他们中间出去。} 我要先根据惯常用法维护这话的意思，然后表明这意思并非我个人独创。常有人被控告；当他们被控告时，他们为自己辩护；当被控告者以合理正义的理由辩护时，听众便说：\Quote{他摆脱了。} 摆脱了？他去了哪里？他还留在原地，却\Quote{摆脱了}。他如何摆脱的？通过他做出的良好陈述和最令人满意的辩护。这正是圣宗徒们所做的，当他们\BibleQuote{跺掉脚上的尘土}，反对那些不接待所传和平信息的人。那位守望者\Quote{从那里出去了}，就是那被说\BibleQuote{我立了你作以色列家的守望者} 的人。因为曾告诉他说：\BibleQuote{你若警告恶人，他仍不离开他的邪恶和他的道，那恶人必因自己的罪而死，你却救了你的灵魂。} 他若这样做，就\Quote{从他那里出去了}，不是通过身体的分离，而是通过为自己工作的辩护。因为他做了他该做的事；尽管那人有责任服从，却未服从。这就是\BibleQuote{你们从那里出去}。\par

24. 梅瑟、依撒意亚、耶肋米亚和厄则克耳就这样呼喊。那么，我们要看看他们是否这样做了，是否离开了天主的子民，投奔了外邦人。耶肋米亚对他的百姓中的罪人和恶人发出了多少严厉的责备啊！然而他生活在他们中间，与他们进入同一座圣殿，举行同样的奥迹；他住在那一群恶人里，却通过呼喊\Quote{出离他们}。这正是\Quote{出离他们}；这正是\Quote{不触摸不洁之物}，即不认同他们的意愿，并且在言语上不宽容他们。我该怎么说耶肋米亚、依撒意亚、达尼尔、厄则克耳以及其余的先知呢？他们没有离开恶民，以免抛弃混在那些百姓中的善人，而他们自己正是在那些人中得以成为他们那样的人。当梅瑟自己在山上领受法律时，下面的百姓造了一个偶像。天主的子民，曾从为他们分开的红海波涛中走过，并淹没了追赶他们的敌人，经历了如此多显示在埃及人身上的征兆和奇迹直到死亡，以及为保护他们直到拯救，却要求一个偶像，强求了一个偶像，制造了一个偶像，朝拜了一个偶像，向偶像献祭。天主向祂的仆人显明了百姓所做的事，并说祂要消灭他们。梅瑟为他们求情，因为他正要回到这百姓中去；然而他本有机会离开并\Quote{出离他们}，如这些人所理解的那样，以便他\Quote{不触摸不洁之物}，不生活在他们中间，但他没有这样做。而且，为了不让人以为他是出于不得已而非出于爱才这样做，天主向他提出了另一个民族，以便消灭这些百姓。天主说：\BibleQuote{我要使你成为一个大民族。}但他没有接受；他紧贴罪人，为罪人祈祷。他是怎样祈祷的？噢，爱的显著证明，我的弟兄们！他是怎样祈祷的？留意那种我曾说过的母性的柔情。当天主威胁那亵渎的百姓时，梅瑟慈爱的心颤抖了，他替他们抵挡天主的义怒。他说：\BibleQuote{吾主，如果祢宽恕他们的罪，就宽恕；但若不，求祢从祢所写的册上涂抹我的名字。}这是何等父性母性的柔情，又是何等确信地说出这话，因为他同时考虑到天主的公义和慈悲：由于祂是公义的，祂不会毁灭义人；由于祂是慈悲的，祂会宽恕罪人。\par

25. 如今你们的明辨必定清楚，应当如何接受圣经的所有这类见证；因此当圣经说我们必须离开恶人时，我们被命令只能理解为在内心离开；以免因与善人分离而犯下比因与恶人联合而退缩更严重的恶，正如这些\OriginalTerm{多纳徒派}{Donatists}所做的那样。但如果他们真的善良，并且这样责备了恶人，而不是他们自己邪恶而诽谤善人，他们就会为了和平容忍任何人，无论他们是谁，因为他们已经接纳了\OriginalTerm{马克西米安派}{Maximianists}为健全的，而他们先前曾谴责他们为失丧的。先知无疑明确地说过：\BibleQuote{离开，从那里出去，不要触摸不洁之物。}但为了理解他说的话，我留意他所做的事。他以自己的行动解释了他的话语。他说：\BibleQuote{离开。}他对谁说的？当然是义人。他命令他们离开谁？当然是罪人和恶人。那么我问，他本人是否离开了这样的人？我发现他没有。所以，他以另一种意义理解。因为他肯定会首先实行他所吩咐的。他在内心离开了他们，责备并斥责了他们。由于他避免认同他们，他\Quote{不触摸不洁之物}；但通过斥责他们，他在天主面前\Quote{出去}为自由；在天主面前，他既不被归咎自己的罪，因为他没有犯罪；也不被归咎别人的罪，因为他没有赞同；也不被归咎疏忽，因为他没有缄默；也不被归咎骄傲，因为他保持了合一。所以，我的弟兄们，无论你们当中有多少仍然被对世界的爱所重压的人，贪婪的，或作伪证的，奸淫的，看戏的，咨询占星家的，走火入魔的，占卜的，占兆的，预言的，醉酒的，纵欲的，你们知道自己当中有多少恶人；尽你们所能表示不赞成这一切，以便在内心\BibleQuote{离开}；并责备他们，以便你们\BibleQuote{出离他们}；不要认同他们，以便你们\BibleQuote{不触摸不洁之物}。\par

\SourceRecord{Translated by R.G. MacMullen.；Nicene and Post-Nicene Fathers, First Series；Vol. 6.；Edited by Philip Schaff.；Buffalo, NY: Christian Literature Publishing Co.,；1888.；<http://www.newadvent.org/fathers/160338.htm>.}

% structure: p0001 source=h1 target=section reason=page-title

\section{新约讲道集 第39篇}

\Emph{[LXXXIX. Ben.]}\par

\Emph{论福音经文}，\BibleRef{玛 21:19}，\Emph{那里耶稣使无花果树枯干；以及论经文}，\BibleRef{路 24:28}，\Emph{那里祂装作还要往前行。}\par

1. 刚刚读过的神圣福音的教训给我们一个警醒的警告，免得我们只有叶子，而没有果实。简言之，免得只有言语而无行为。非常可怕！当在这读经中，用心眼看见那枯干的树，因听到那句话而枯干：\BibleQuote{从今以后，你永远不再结果子}时，谁不害怕呢？愿恐惧带来改正，改正带来果实。因为毫无疑问，主基督预见了某棵树会理应枯干，因为它有叶子，却没有果实。那树就是会堂，不是那被召的，而是那被弃的。因为从其中也召出了天主的人民，他们以真诚和真理，在先知中等待天主的救恩，耶稣基督。由于他们在信德中等待，他们被认为配得认识祂，当祂显现时。因为宗徒们从其中出来，一大群走在主的驴前面的人也从其中出来，他们喊道：\BibleQuote{贺三纳于达味之子！奉主名而来的，当受赞颂。}当时有一大群信主的犹太人，一大群在基督为他们倾流祂的血之前就信了祂的人。因为主自己来到世上不是为别的，\BibleQuote{只为以色列家迷失的羊}，这不是没有意义的。但在其他的人身上，当祂被钉十字架、并已被高举升天后，祂找到了悔改的果实；祂没有使这些枯干，而是在祂的田地里栽培他们，用祂的话浇灌他们。属于这个数目的是那四千犹太人，他们在宗徒们和与他们一起的人充满了圣神、说万国方言之后信了主，并在那种方言的多样性中预先宣告了教会将遍布万国。他们那时信了，\BibleQuote{他们是以色列家迷失的羊}；但是因为\BibleQuote{人子来，是为寻找并拯救迷失了的}，祂也找到了这些。但他们隐藏在荆棘中，各处分散，好像被狼撕碎；因为他们隐藏在荆棘中，祂没有来寻找他们，除非被祂苦难的荆棘刺伤；然而祂来了，找到了，救赎了他们。他们杀死的，与其说祂，不如说是他们自己。他们被为他们而死的那位所救。因为当宗徒们说话时，他们被刺痛了；他们被良心刺痛，他们曾用长枪刺祂；被刺痛后，他们寻求建议，得到后听从了，悔改了，获得了恩宠，信了并喝下了他们狂怒时所流的血。但那些留在这恶劣而不结果子的种族中，直到今日，并要留到终结的，在那棵树上得到了预表。你们今天到他们那里，发现他们拥有全部先知著作。但这些只是叶子；基督饿了，祂寻找果实；但在他们中间找不到果实，因为祂在他们中间找不到自己。因为谁没有基督，谁就没有果实。谁不持守基督的合一，谁没有爱德，谁就没有基督。因此，通过这个链条，谁没有爱德，谁就没有果实。请听宗徒的话：\BibleQuote{圣神的效果是爱德}；这样他称赞这簇果实，即这果子：\BibleQuote{圣神的效果}，他说，\BibleQuote{是爱德、喜乐、平安、忍耐。}当爱德领先时，不要对后面的话感到惊奇。\par

因此，当门徒对无花果树的枯萎感到惊奇时，祂便向他们阐明\(\OriginalTerm{信德}{fides}\)的价值，对他们说：\(\BibleQuote{如果你们有信德，毫不怀疑；}\) 就是说，如果你们在一切事上信赖天主，而不说：\(\Quote{天主能作这事，这事祂不能作；}\) 而是依靠全能者的全能；\(\BibleQuote{你们不只能作这事，而且如果你们对这座山说：起来，投到海里去！也必成就。凡你们在祈祷中祈求的，只要相信，就必获得。}\) 现在我们读到奇迹由门徒们所行，更确切地说，是由主藉门徒而行；因为祂说：\(\BibleQuote{离了我，你们什么也不能作。}\) 主没有门徒能行许多事，但门徒没有主却不能行任何事。那甚至能创造门徒本身的主，当然不需要他们协助来创造他们。我们读到宗徒们的奇迹，但从未读到他们使一棵树枯萎，或将一座山移入海中。因此，让我们探究这事是在何处成就的。因为主的话不会没有效果。如果你按字面通常的意义思考\(\Quote{树木}\)和\(\Quote{山岭}\)，这事并未发生。但如果你思想祂所说的那棵树，以及先知所说的主的山：\(\BibleQuote{在末日，上主殿宇的山岭必要矗立；}\) 如果你这样思想并领悟，这事就发生了，且是由宗徒们成就的。那棵树就是犹太民族，但我再说一遍，是其中被遗弃的部分，而非被召叫的部分；我们所说的那棵树就是犹太民族。那座山，一如先知见证所教导我们的，就是主自己。枯萎的树是失去基督荣耀的犹太民族；海洋是这世界及万民。现在请看宗徒们对即将枯萎的这棵树说话，并将山移入海中。在\(\BibleBook{宗徒}{大事录}\)中，他们对那些抗拒真理之言、即只有叶子却没有果实的犹太人说话：\(\BibleQuote{天主的话本来应先讲给你们听，但因你们拒而不受}\)（你们引用先知的话，却不承认先知所预言的那位，即只有叶子），\(\BibleQuote{看，我们转向外邦人。}\) 这也曾被先知预言：\(\BibleQuote{看，我已立你作外邦人的光明，使你成为我的救恩，直到地极。}\) 看，那棵树已枯萎；基督已被移向外邦人，山被移入海中。那栽种在葡萄园中的树怎能不枯萎呢？关于这葡萄园曾说过：\(\BibleQuote{我要命令云不再降雨在上面。}\)\par

主为传达这一真理而行事预言性，我的意思是，就这棵树而言，祂的旨意并不仅仅要展示一个奇迹，而是藉着奇迹预示将来要发生的事，有许多事教导并说服我们，甚至违背我们的意愿强使我们相信。首先，这棵树没有果子算什么过错？因为即使它在适当季节，即结果子的季节没有果子，这也绝不是树的过错；树没有感觉和理性，不能受责备。但这里还加上另一事实，正如我们在另一位明确提到此事的福音作者的叙述中读到：\(\BibleQuote{还不是无花果的季节。}\) 因为那时正是无花果树发出嫩叶的时候，我们知道叶子先于果实；这我们可以证明，因为主受难的日子临近了，我们知道祂是在何时受难；即使我们不知道，我们当然也应该相信那位说\(\BibleQuote{还不是无花果的时候}\)的福音作者。因此，如果只打算展示一个奇迹，而不预示什么，那么主以祂的慈悲和怜悯，使任何祂发现的枯萎之树再次变绿，不是更符合祂的作风吗？正如祂治愈病人、洁净癞病人、复活死人一样。然而相反，祂似乎违反祂通常的仁慈规则，找到一棵绿叶树，虽尚未在正确季节结果，但尚未使栽培者失望，祂却使它枯萎；仿佛祂会对我们说：\(\Quote{我并不喜悦这棵树的枯萎，但我这样向你们传达，我并非无缘无故这样做，而是因为我想借此给你们一个你们更应重视的教训。我所咒骂的不是这棵树，我并非对这无感觉的树施加惩罚，而是使你们——无论你是谁，若思考这事——心生畏惧，免得当基督饿了时你们轻看祂，宁愿以果实为富足，而不以叶子为遮盖。}\)\par

4. 主藉祂所行的这一件事所暗示的，正是祂有意要表明的这一点。还有别的吗？祂饿了，来到树前，寻找果子。难道祂不知道那时不是结果子的时候吗？知道这棵树情况的农夫，难道它的创造主不知道吗？祂在树上寻找那尚未结出的果子。祂是真的在寻找，还是\OriginalTerm{装作}{pretence}寻找呢？因为如果祂真的寻找，祂就是错了。但愿祂绝不会错！所以祂是装作寻找。你害怕承认祂装作寻找，就承认祂错了。你又拒绝承认祂错了，于是又陷入祂装作寻找的想法。我们在这两者之间焦灼不堪。既然焦灼，就让我们求雨，好使我们得以青翠，免得我们说了任何对主不敬的话，反而枯萎了。圣史确实说：\BibleQuote{祂来到树前，在树上没有找到果子。} 除非祂或者真的寻找，或者明知那里没有果子却装作寻找，否则不会说祂\Quote{没有找到}。因此，我们毫不迟疑地绝不说基督错了。那么怎样？我们说祂装作寻找吗？我们这样说吗？我们如何摆脱这个困境？让我们说那些若不是圣史在另一处论及主的话，我们不敢自行说出的。让我们说圣史所记载的，说了以后，让我们理解它。但为了理解，我们先要相信。因为先知说：\BibleQuote{除非你们相信，你们不会明白。} 主基督复活后，在路上与两位门徒同行，他们尚未认出祂，祂作为第三位旅伴加入他们。他们到了要去的地方，圣史说：\BibleQuote{祂却装作还要前行。} 但他们以亲切的善意挽留祂，说天色已晚，求祂留住；他们接待了祂，祂擘饼时，在祝福和擘饼中他们认出了祂。因此，现在我们不必害怕说祂装作寻找，如果祂曾装作前行。但这里又出现了另一个问题。昨天我详细论述了宗徒们中的真理；那么，我们如何在主自己身上找到任何\Quote{装作}呢？因此，弟兄们，我必须告诉你们，并按主为你们的益处赐给我的微薄才能教导你们，并将你们可以作为解释全部圣经的准则传达给你们。一切所言所行，都应按其字面意义理解，否则便有所象征；或者至少两者兼备，既有其字面解释，也有象征意义。我如此提出了三点，现在必须给出例子；从何处引出呢？当然是从圣经。字面意义上说，主受难、复活、升天；我们将在世界末日复活，如果我们不轻慢祂，我们将与祂永远为王。这一切都按字面理解，不必寻求象征；正如所说，实际也是如此。各种行动也是如此。宗徒上耶路撒冷去见伯多禄，宗徒确实这样做了，确实发生了，这是他个人的行动。这是一个事实，他告诉你；一个简单的事实，按字面意义。\BibleQuote{匠人弃而不用的石头，反而成了屋角的基石} 是比喻的说法。如果按字面理解\Quote{石头}，哪块\Quote{石头曾被匠人弃用，却成了屋角的基石}？如果按字面理解\Quote{石头}，这\Quote{石头}成了哪个\Quote{屋角}的基石？如果我们承认这是比喻的说法，并按其比喻理解，那么\OriginalTerm{基石}{cornerstone}就是基督；屋角之头就是教会的头。为什么教会是屋角？因为她从一边召叫了犹太人，从另一边召叫了外邦人，这两道墙如同从不同方向而来，相遇结合为一，藉着和平的恩宠将她束缚在一起。因为\BibleQuote{祂是我们的和平，使双方合而为一。}\par

5. 你们已经听到了字面表达、字面行为和比喻表达的例子；你们正在等待一个比喻行为的例子。这样的例子很多，但与此同时，正如刚才提到\OriginalTerm{角石}{corner-stone}所提示的，当\Person{雅各伯}用油傅抹他睡觉时枕在头下的石头时，他在梦中看到了一个奥秘的景象：从地上到天上的梯子，\Person{天使}上去下来，\Person{主}站在梯子上。他明白了这景象所要象征的，就把那石头当作\Person{基督}的预像，以此向我们证明，他对那异象和启示的理解并不陌生。因此，他傅抹那石头并不奇怪，因为\Person{基督}的名字就来自\Quote{\OriginalTerm{受傅}{the anointing}}。后来，\Person{雅各伯}在圣经中被记载为\Quote{一个无诡诈的人}。而你们知道，这\Person{雅各伯}又名以色列。因此在福音中，\Person{主}看见\Person{纳塔乃耳}时，就说：\BibleQuote{看，这确是一个以色列人，在他内毫无诡诈。}那个以色列人还不知道与他说话的是谁，就回答说：\BibleQuote{祢从哪里认识我？}\Person{主}对他说：\BibleQuote{当你还在无花果树下时，我就看见了你；}好像说：当你还在罪的荫蔽下时，我已预定了你。\Person{纳塔乃耳}因为想起自己曾在那无花果树下，而\Person{主}并不在那里，就承认了祂的天主性，回答说：\BibleQuote{祢是天主子，祢是以色列的君王。}那个曾在无花果树下的人并没有成为一棵枯萎的无花果树；他认出了\Person{基督}。\Person{主}就对他说：\BibleQuote{因为我说：当你还在无花果树下时，我就看见了你，你就相信吗？你将看到比这些更大的事。}这些\Quote{更大的事}是什么？\Quote{我实实在在告诉你们}（因为他是一个\BibleQuote{无诡诈的以色列人}；请记住无诡诈的\Person{雅各伯}；并回想他所说的话：他头下的石头、梦中的景象、从地上到天上的梯子、上去下来的天使；这样你就会明白\Person{主}想要对那个\BibleQuote{无诡诈的以色列人}说的是什么）；\BibleQuote{我实实在在告诉你们：你们将要看见天开了}（听啊，无诡诈的\Person{纳塔乃耳}，无诡诈的\Person{雅各伯}所看见的）；\BibleQuote{你们将要看见天开了，天主的天使上去下来}（向谁？）\BibleQuote{在人子身上。}因此，祂作为人子，头上受了傅油；因为\BibleQuote{女人的头是男人，而男人的头是基督。}现在请注意，祂并没有说：\Quote{从人子那里上去，并下降到人子身上}，好像祂只在上面；而是说\BibleQuote{在人子身上上去下来}。请听人子从上面呼喊：\BibleQuote{扫禄，扫禄。}请听人子从下面说：\BibleQuote{你为什么迫害我？}\par

6. 你们已经听到了字面表达的例子，如\Quote{我们将要复活}；字面行为的例子，如经上所说：\BibleQuote{保禄上耶路撒冷去见伯多禄。}\BibleQuote{匠人弃而不用的石头}是比喻表达；\Person{雅各伯}头下的\Quote{受傅的石头}是比喻行为。现在应你们期望，有一个两者结合的例子：一件既是字面事实，同时又象征其他事物的东西。\BibleQuote{我们知道亚巴郎有两个儿子，一个由婢女所生，一个由自由妇人所生}；这不仅是故事，更是字面事实；你们在寻找其中所象征的吗？\BibleQuote{这两个是两约。}所以，比喻性的说法是一种虚构。但由于它有真实事件作为预表，且比喻本身也有真理根基，所以它摆脱了虚假的指控。\BibleQuote{播种的人出去撒种；撒种时，有的落在路旁，有的落在石地上，有的落在荆棘中，有的落在好地上。}谁出去\Quote{撒种}？他何时出去？落在什么\Quote{荆棘}、\Quote{石头}、\Quote{路旁}或什么田里？如果我们把它当作虚构故事，我们以比喻来理解；它是虚构的。因为如果真有播种的人出去，真的把种子撒在这些不同地方，正如我们所听到的，那就不是虚构，也就没有虚假。但现在尽管是虚构，却不是虚假。为什么？因为虚构有更深层的含义，不欺骗人。它只需要人理解，不会引人误入歧途。因此，基督想向我们传达这个教训，就寻找果实，并以此向我们展示了一个比喻性的、不骗人的虚构；因此这是一个值得称赞、而非责备的虚构；不是通过考查它而陷入虚假，而是通过勤勉探究而发现真理。\par

7. 我看到有人可能会说：请给我解释；\BibleQuote{祂假装还要前行}是什么意思？因为如果它没有更深的意义，那就是欺骗，是谎言。那么，我们必须按照我们的解释规则和区分，告诉你们这个\BibleQuote{祂假装还要前行}意味着什么；\BibleQuote{祂假装还要前行}，却被留住了。因此，主基督既然在身体临在方面被认为缺席，被人以为真的不在，祂就好像要\BibleQuote{前行}。但是，你们要用信德紧紧抓住祂，在擘饼时紧紧抓住祂。我还能说什么呢？你们认出了祂吗？如果是，那么你们就找到了基督。我不再多谈这\OriginalTerm{圣事}{Sacrament}。那些推迟认识这\OriginalTerm{圣事}{Sacrament}的人，基督就离他们更远。那么，让他们紧紧抓住它，不要让祂离开；让他们邀请祂到自己的家，这样他们就被邀请到天堂。\par

\SourceRecord{Translated by R.G. MacMullen.；Nicene and Post-Nicene Fathers, First Series；Vol. 6.；Edited by Philip Schaff.；Buffalo, NY: Christian Literature Publishing Co.,；1888.；<http://www.newadvent.org/fathers/160339.htm>.}

% structure: p0001 source=h1 target=section reason=page-title

\section{新约讲道集 第四十篇}

\EditorialNote{[XC. Ben.]}\par

\Emph{论福音话语，\BibleRef{玛 22:2}等，关于王子的婚宴；反对\OriginalTerm{多纳徒派}{Donatists}，论爱德。在\Place{迦太基}的\Place{瑞斯提图塔}宣讲。}\par

1. 所有信徒都知道王子的婚宴和他的盛宴，主的筵席铺开，凡愿意的都可以来。但重要的是，每个人都要注意自己如何走近，即使他未被禁止走近。因为圣经教导我们，主的宴会有两次：一次是善恶都来，另一次是恶人不来。所以在刚才福音宣读时我们听到的那场盛宴，有善客也有恶客。凡是推脱不来赴宴的，都是恶人；但所有进来的人不都是善人。所以你们这些赴宴的善客，我对你们说，你们心中记着这话：\BibleQuote{那不相称地吃和喝的人，就是吃喝自己的罪案。}我向你们所有这样的人说话，叫你们不要在外面寻找善人，要在里面容忍恶人。\par

2. 我不怀疑你们愿意听，亲爱的，我所说的人是谁，就是他们不要在外面寻找善人，要在里面容忍恶人。如果里面全是恶人，我对谁说话呢？如果里面全是善人，我又劝他们容忍谁为恶人呢？所以让我先靠主的帮助，尽力从这个困难中脱身。如果严格完美地说善，除天主独一外，没有一个是善的。你们有主最清楚的话：\BibleQuote{你为什么称我为善？除天主一个外，没有善的。}那么，那婚宴怎么会有善客和恶客呢，既然\BibleQuote{除天主独一外，没有善的}？首先你们要知道，在某种意义上，我们都是恶的。是的，无疑在某种意义上我们都是恶的；但在任何意义上我们都不是善的。因为难道我们能和宗徒们相比吗？主亲自对他们说：\BibleQuote{你们纵然不善，尚且知道把好东西给你们的儿女。}如果我们查考圣经，十二宗徒中只有一个恶人，主在某处论到他说：\BibleQuote{你们是洁净的，但不是全部。}但当他一起对他们说话时，他说：\BibleQuote{你们纵然不善。}伯多禄听到了，若望听到了，安德肋听到了，其余十一位宗徒都听到了。他们听到了什么？\BibleQuote{你们纵然不善，尚且知道把好东西给你们的儿女，何况你们在天之父，岂不更将好东西赐给求他的人？}当他们听到他们不善时，他们绝望了；但当他们听到天主在天上是他们的父时，他们又振奋了。\BibleQuote{你们纵然不善}；那么恶人应得什么？只有惩罚。\BibleQuote{何况你们在天之父？}子女应得什么？继承的希望。所以，在某个方面，他们是恶的，在另一个方面，他们是善的。因为对他们说了\BibleQuote{你们纵然不善，尚且知道把好东西给你们的儿女}之后，立刻加上\BibleQuote{何况你们在天之父？}\par

3. 所以，根据某一角度，他们是恶的，而根据另一角度，他们是善的。因为对他们说了\BibleQuote{你们纵然不善，尚且知道把好东西给你们的儿女}之后，立刻加上\BibleQuote{何况你们的天父？}所以他是那些恶人的父，但不是让他们一直如此，因为他要医治他们。根据某一角度，他们是恶的。然而，我想那家主的婚宴客人并不属于那一类，即关于他们所说的\BibleQuote{邀请了善人和恶人}，以致他们被算在恶人之中，因为我们听见那个没有婚宴礼服的人被赶出去了。根据某一角度，我再重复，他们是恶的，然而也是善的；根据某一角度，他们是善的，然而也是恶的。听若望按着什么角度他们是恶的：\BibleQuote{如果我们说我们没有罪，就是自欺，真理不在我们内。}看，按着什么角度他们是恶的：因为他们有罪。按着什么角度他们是善的？\BibleQuote{如果我们承认我们的罪，他是忠信正义的，必赦免我们的罪，并洗净我们一切的不义。}那么，如果我们根据这个解释原则（我想你们已经听我刚刚从圣经中引出的），即同一些人在某一方面是善，在另一方面是恶；如果我们愿意按照这个意义接受\BibleQuote{邀请了善人和恶人}这句话，即同一些人同时是善又是恶；如果我们愿意这样接受，我们却不能这样做，因为那个人被发现有\BibleQuote{没有穿婚宴礼服}，他不仅被\BibleQuote{赶出去}，失去了那盛宴，而且被定罪，受永远黑暗的惩罚。\par

4. 但有人会说：一个人算什么？如果一个人群中\Quote{没有穿婚宴礼服}的人偷偷溜进来，没有被家主仆人们发现，这算什么稀奇、什么大事？难道就因为那个人，就能说\Quote{他们邀请了好的和坏的}吗？因此，我的弟兄们，请留心并明白。那个人代表了一类人，因为这样的人很多。这里，一位勤勉的听众可能会回答我说：\Quote{我不希望你告诉我你的猜测；我希望你向我证明，那个人代表了许多人。}靠着主的当前帮助，我将清楚证明这一点；我也不必远求，就能证明。天主会在这段经文中用祂自己的话帮助我，并藉着我的服务为你提供一个明白的证明。\Quote{家主进来观看坐席的客人。}\newline{}看，我的弟兄们，仆人的工作仅仅是邀请并带进好人和坏人；请注意，经上并没有说仆人们留心观察客人，并在他们中间发现一个没有穿婚宴礼服的人，便对他说话。这没有记载。是家主看见了他，是家主发现了，是家主检查了，是家主把他分别出来。不能忽略这一点。但我已承担了确立另一点的任务，即那个人如何代表了许多人。\Quote{家主}然后\Quote{进来观看坐席的客人，他在那里看见一个人没有穿婚宴礼服。就对他说：朋友，你没有穿婚宴礼服，怎么进到这里来？那人无言可答。}\newline{}因为询问他的那一位，是他无法虚伪造答的。所寻找的礼服是在心里，不在身上；因为如果是穿在外面的，就连仆人也无法隐藏。那婚宴礼服必须穿在哪里，请听这话：\Quote{愿你的司铎披上正义。}关于那礼服，宗徒说：\Quote{如果我们被发现是穿着的，而不是赤身的。}\newline{}因此，那个逃脱了仆人注意的人，被主发现了。被质问时，他无言可答；他被捆绑、被抛弃、被定罪，一个人代表了许多人。我说：主啊，祢教导我们，在这件事上祢是在警告所有人。那么，我的弟兄们，请与我一同回想你们所听过的话，你们就会立刻发现、立刻断定，那一个人就是许多人。的确，主质问的是一个人，对一个人说：\Quote{朋友，你怎么进到这里来？}那一个人无言可答，关于那同一个人，经上说：\Quote{捆起他的手脚来，把他丢在外面的黑暗中；在那里要有哀号和切齿。}为什么？\Quote{因为被召的人多，被选的人少。}谁能否认这真理的显明呢？\Quote{把他丢在外面黑暗中，}祂说。\Quote{他}，无疑就是那个人，主关于他说：\Quote{因为被召的人多，被选的人少。}那么，不被丢弃的是少数。的确，只有一个人\Quote{没有穿婚宴礼服。把他丢出去。}但为什么他被丢出去？\Quote{因为被召的人多，被选的人少。}保留少数，丢出多数。的确，那个人只是一个。然而，那一个人无疑不仅是许多人，而且那些许多人在数量上远远超过了好人的数量。因为好人也是多的；但与坏人相比，他们是少的。在收成中有许多麦子；但与糠秕相比，麦粒是少的。同样的人，在他们自己看来是多的，与坏人相比是少的。我们如何证明在他们自己看来是多的？\Quote{将有许多人从东方和西方来。}他们来哪里？来到那个好人和坏人都进去的宴席。但谈到另一个宴席时，祂接着说：\Quote{并同亚巴郎、依撒格和雅各伯在天国里坐席。}那是坏人不能接近的宴席。愿我们相称地接受那个现在的宴席，好能到达另一个。那么，同一些人，既是多的，也是少的；在他们自己看来是多的，与坏人相比是少的。因此，主怎么说？祂发现了一个，便说：\Quote{让多人被丢出去，少数留下。}因为说\Quote{被召的人多，被选的人少}，无非就是清楚指明，在现在的宴席中，哪些人被认为是可被带到另一个宴席的，那里没有坏人会来。\par

5. 那又是什么呢？我不愿意你们所有在今生接近主的圣筵的人，同那许多将被关在门外的人在一起，而是同那少数蒙拣选的人在一起。但你们怎能达到这一点呢？拿起\Quote{婚宴礼服}。你们会说：\Quote{请解释这\InnerQuote{婚宴礼服}给我们听。}毫无疑问，那就是只有善人所拥有的礼服，他们将在筵席中被留下，被保留至那另一个筵席——没有恶人会接近的筵席——他们将因主的恩宠安全地到达那里；这些人拥有\Quote{婚宴礼服}。那么，我的弟兄们，让我们在信徒中寻找那些拥有恶人所没有的东西，这就是\Quote{婚宴礼服}。如果我们谈论圣事，你们看这些是善恶共有的。是圣洗圣事吗？没有圣洗圣事确实无人能接近天主，但并非每个领受圣洗圣事的人都达到祂。因此我不能认为圣洗圣事，即圣事本身，是\Quote{婚宴礼服}；因为这礼服我既在善人身上看到，也在恶人身上看到。或许是祭台，或在祭台上所领受的。不；我们看到许多人吃，且\BibleQuote{吃与喝自己的审判}。那么它是什么呢？是禁食吗？恶人也禁食。是涌向教会吗？恶人也涌向那里。最后，是神迹吗？不仅善恶都能行神迹，而且有时善人也不行神迹。看，在古代民族中，法老的术士行了神迹，以色列子民却没有；在以色列子民中，只有梅瑟和亚郎行了神迹；其余的人没有行，但他们看见了，就畏惧而相信。行神迹的法老术士，难道比不能行神迹的天主子民更优秀吗？然而那子民是天主的子民。在教会本身，听宗徒说：\BibleQuote{众人岂都是先知？众人岂都有治病的恩赐？众人岂都说语言？}\par

6. 那么什么是\Quote{婚宴礼服}？这就是婚宴礼服：宗徒说：\BibleQuote{诫命的终向，是出自纯洁的心、善良的良心和无伪信德的爱德。}这就是\Quote{婚宴礼服}。不是任何一种爱德；因为常有同伙作恶的人似乎彼此相爱。那些一起抢劫的人，一起喜爱可恶的巫术和舞台的人，一起为战车赛或斗兽场呐喊的人，这些人常常彼此相爱；但在他们内没有\Quote{出自纯洁的心、善良的良心和无伪信德的\OriginalTerm{爱德}{charity}}这样的\Quote{婚宴礼服}。\BibleQuote{我若能说人间的语言和天使的语言，但我若没有爱德，我就成了鸣的锣、响的钹一般。}唯独语言来了，却对他们说：\Quote{你们没有穿婚宴礼服，怎么进到这里来？}\BibleQuote{我若有先知之恩，又明白一切奥秘和一切知识；并且有全备的信德，甚至能移山，但我若没有爱德，我什么也不是。}看，这些是常没有\Quote{婚宴礼服}之人的奇迹。\Quote{我若有这一切，而没有基督，我什么也不是。}那么\Quote{先知之恩}什么也不是吗？\Quote{明白奥秘的知识}什么也不是吗？并不是这些什么也不是；而是\Quote{我}若拥有它们，\Quote{却没有爱德，就什么也不是}。没有这唯一的好事，多少好事毫无益处！因此，我若没有爱德，即使我慷慨地把所有施舍给穷人，即使我到流血和赴火来宣认基督之名，这些事也可能出于求光荣的动机，因而落空。因为它们可能出于求光荣的动机，而不是出于虔诚情感洋溢的爱德，所以他也明确地一一列举，你们要听：\BibleQuote{我若把我所有的财产全施舍给人，我若交出我的身体被火烧，但我若没有爱德，为我毫无益处。}这就是\Quote{婚宴礼服}。你们问问自己：若你们拥有它，你们就可以在主的筵席中无所畏惧。在同一个人内，存在两种东西：爱德和\OriginalTerm{欲望}{desire}。若爱德还未出生，愿它在你内出生；若已经出生，愿它被养育、被培养、被增长。至于那欲望，尽管在今生无法完全熄灭，\BibleQuote{因为如果我们说我们没有罪，就是欺骗自己，真理也不在我们内}；但罪恶在多大程度上在我们内，我们就在多大程度上没有罪：愿爱德增长，欲望减少；这样，那一个——即爱德——终有一天得以圆满，而欲望被销毁。穿上\Quote{婚宴礼服}：我正对你们这些尚未拥有它的人说话。你们已经进来了，已经接近了筵席，却还没有那件尊荣新郎的礼服；\BibleQuote{你们还在寻求自己的事，而不是耶稣基督的事。}因为\Quote{婚宴礼服}是为尊荣这结合，即新郎与新妇的结合而领受的。你们知道新郎：就是基督。你们知道新妇：就是教会。尊重新妇，尊重新郎。若你们给他们应有的尊荣，你们就是他们的子女。因此，在这事上要长进。爱主，并这样学会爱自己；这样，当你们因爱主而爱了你们自己之后，就可以安心地爱人如己。因为当我找到一个不爱自己的人，我怎能将他的近人托付给他，要他爱人如己呢？你们会说，有谁不爱自己呢？谁呢？看，\Quote{凡喜爱不义的人，就是恨自己的灵魂。}那爱自己的身体却恨自己的灵魂，以致伤害自己和自己的灵魂身体的人，算是爱自己吗？谁爱自己的灵魂？那全心、全意爱天主的人。对这样的人，我立刻将他的近人托付给他。\BibleQuote{爱你的近人如你自己。}\par

7. 有人会说：\Quote{谁是我的邻人？} 每个人都是你的邻人。我们不是都有同一对父母吗？每种动物彼此都是邻人：鸽子对鸽子，豹对豹，蝮蛇对蝮蛇，绵羊对绵羊；难道人不是人的邻人吗？请回想受造的秩序。天主发言，水就产生了游动的生物、巨大的鲸鱼、鱼类、鸟类以及诸如此类的东西。难道所有的鸟都来自一只鸟吗？所有的兀鹰都来自一只兀鹰吗？所有的鸽子都来自一只鸽子吗？所有的蛇都来自一条蛇吗？所有的金鲷都来自一条金鲷吗？所有的绵羊都来自一只绵羊吗？不，大地确实是同时产生了所有这些种类。但说到人，大地并没有产生人。我们被造了一位父亲，甚至不是两位——父亲和母亲：我们被造了一位父亲，我说，甚至不是两位——父亲和母亲；但来自这位父亲的一位母亲；这位父亲不从任何人而来，而是由天主所造，这位母亲则从他而来。那么，请注意我们种族的本性：我们从一个源泉涌出；因为那源泉变成了苦毒，我们全部从一根好橄榄树变成了野橄榄树。于是恩宠也来了。那一位生我们归于罪与死，然而作为一个种族，然而作为彼此的邻人，然而不仅相似，而且彼此相关。那一位来对抗一位；对抗那分散的一位，那聚集的一位来了。这样，对抗那杀戮的一位，就是那使人活着的一位。\BibleQuote{因为正如在亚当内众人都死了，照样在基督内众人都要复活。} 现在，正如凡生于第一个的都死，同样凡相信基督的都活过来。当然，前提是他有\Quote{婚宴礼服}，并作为将要留下、而非被赶出的人而被邀请。\par

8. 所以，我的弟兄们，要有爱德。我已说明这衣裳就是婚宴礼服。信德受到赞美，这是显而易见的，它受到赞美；但宗徒辨别这是怎样的信德。因为有些人夸耀信德，却没有良好的品行，宗徒雅各伯斥责他们说：\Quote{你信只有一个天主，你做得对；魔鬼也信，且战栗害怕。} 请与我一同回想，伯多禄是因什么受到赞美，因什么被称为有福的。是因他说：\Quote{你是基督，永生天主之子} 吗？那称他有福的，并非看言语的声音，而是看内心的情感。因为你们要知道吗，伯多禄的福分并不在于这些话？魔鬼也说了同样的话。\Quote{我们知道你是谁，你是天主子。} 伯多禄宣认\Quote{他是天主子}；魔鬼也宣认\Quote{他是天主子}。\Quote{请分辨，我主，请分辨两者。} 我明确加以分辨。伯多禄出于爱德而说，魔鬼出于恐惧而说。伯多禄又说：\Quote{我与你同在，至死不渝。} 魔鬼说：\Quote{我们与你有何相干？} 所以，你们这些来赴宴席的人，不要只以信德夸耀。好好辨别这信德的性质；这样，在你身上就会认出\Quote{婚宴礼服}。让宗徒来辨别，让他教导我们：\Quote{割损毫无益处，未受割损也毫无益处，唯有信德。} 告诉我们，什么信德？难道魔鬼不也相信且战栗吗？\Quote{我告诉你们，} 他说，\Quote{请听，我现在要辨别：} \Quote{但信德借爱德行事。} \Quote{我若有全备的知识，} 他说，\Quote{和全备的信德，甚至能移山，却没有爱德，我什么也不算。} 要有信德与爱德；因为你们不可能有爱德而没有信德。我警告你们，我劝勉你们，我以主的名教导你们，亲爱的，你们要有信德与爱德；因为你们可能有信德而没有爱德。我并非劝你们要有信德，而是爱德。因为你们不可能有爱德而没有信德；我指的是对天主和对邻人的爱；没有信德，这爱从何而来呢？那不相信天主的人，怎能爱天主呢？那愚妄人怎能爱天主，他\Quote{心中说：没有天主}？可能你们相信基督已来到，却不爱基督。但你们不可能爱基督，却还说基督已来到。\par

9. 那么，要有信德并怀有爱德。这就是\Quote{婚宴礼服}。你们这些爱基督的人，要彼此相爱，爱你们的朋友，爱你们的仇敌。不要让这事对你们显得困难。那么，你们因此会失去什么呢，当你们获得如此之多？什么？你们向天主祈求，让你们的仇敌死去，作为某种大恩惠吗？这不是\Quote{婚宴礼服}。请将你们的思绪转向那位新郎自己，祂为你们悬挂在十字架上，并为祂的仇敌向天父祈祷；祂说：\BibleQuote{父啊，宽恕他们，因为他们不知道他们在做什么。}你们已经看到新郎这样说；也请看新郎的朋友，一位穿着\Quote{婚宴礼服}的宾客。看真福斯德望，他如何仿佛在愤怒和怨恨中斥责犹太人：\BibleQuote{你们这些执拗而心耳未受割损的人，你们时常抗拒圣神。你们的祖先杀害了哪位先知呢？}你们已听到他口舌多么严厉。于是你们立刻准备攻击任何人；而我希望那是攻击得罪天主的人，而非得罪你们的人。有人得罪了天主，你们却不指责他；他得罪了你们，你们就喊叫；那\Quote{婚宴礼服}在哪里呢？你们已听到斯德望如何严厉；现在请听他是如何爱的。他冒犯了他所斥责的人，并被他们用石头砸死。当他被四面狂暴迫害者的手和石块的打击所压倒、伤重至死时，他首先说：\BibleQuote{主耶稣基督，接纳我的灵魂。}然后，在他为自己站立的祈祷之后，他跪下为那些用石头砸他的人祈祷，说：\BibleQuote{主，不要将这罪归于他们；让我在身体里死去，但不要让这些人在灵魂里死去。说了这话以后，他就安眠了。}这些话之后他再没有添加什么；他说了，便离去了；他最后的祈祷是为了他的仇敌。你们由此学会拥有\Quote{婚宴礼服}。所以你们也要屈膝，以额头叩地，当你们即将接近主的祭台、圣经的盛筵时，不要说：\Quote{但愿我的仇敌死去！主，如果我曾对您有任何功德，请杀死我的仇敌。}因为如果你们这样说，难道不怕祂回答你们：\Quote{如果我选择杀死你的仇敌，我首先会杀死你。什么！你因为现在被邀请来到这里而夸耀？想想你不久前是什么。你难道没有亵读过我？你难道没有嘲笑过我？你难道没有希望从地上抹去我的名字？然而现在你因为被邀请来到这里而自夸！如果我曾在你作我仇敌时杀死你，我怎能让你成为我的朋友？相反，你用邪恶的祈祷来教导我做我未曾在你身上做过的事？}是的，相反，天主对你们说：\Quote{让我教导你效法我。当我悬挂在十字架上时，我说：\InnerQuote{宽恕他们，因为他们不知道他们在做什么。}我把这功课教给了我的勇敢士兵。做我反抗魔鬼的新兵吧。除非你为你的仇敌祈祷，你绝不可能以任何方式无敌战斗。然而，你确实可以祈求这件事，是的，祈求这件事，祈求你可以逼迫你的仇敌；但要明辨地祈求，清楚区分你所求的是什么。看，一个人是你的仇敌；回答我，他里面什么与你为敌？是在他作为人这一点上与你为敌吗？不。那是什么呢？是他邪恶。就他是人，是我所造的他而言，他并不与你为敌。}祂对你们说：\Quote{我没有造人邪恶；他因不顺从而成为邪恶，他顺从魔鬼而非天主。他所造成的自己，与你为敌；就他邪恶而言，他是你的仇敌；而非就他是人而言。因为我听到\InnerQuote{人}和\InnerQuote{邪恶}这两个词；一个是本性的名字，另一个是罪的名字；我治愈罪，而保存本性。所以，你们的天主对你们说：\InnerQuote{看，我为你伸冤，我确实杀死你的仇敌；我除去使他邪恶的东西，我保存使他成为人的东西：现在，如果我使他成为一个好人，我岂不是杀了你的仇敌，并使他成为你的朋友？}因此，要祈求你所祈求的，不是让人灭亡，而是让这些敌意消亡。因为如果你为此祈祷，希望那人死去，那是一个恶人针对另一个恶人的祈祷；当你说\InnerQuote{杀死那恶人}时，天主回答你：\InnerQuote{你们中哪一个？}}\par

10. 那么，扩展你的爱，不要把它限制在你的妻子和儿女身上。这样的爱甚至在走兽和麻雀中也能找到。你知道麻雀和燕子如何爱它们的伴侣，如何一同孵蛋，一同养育雏鸟，出于一种自由而自然的善意，而且不图回报。因为麻雀不会说：\Quote{我要养育我的幼雏，这样当我年老时，它们可以喂养我。}它没有这样的想法；它爱它们并喂养它们，出于对它们的爱；展现了父母的情感，并不期待回报。同样，我知道，我确信，你们也爱你们的子女。\BibleQuote{因为儿女不该为父母积蓄，而是父母该为儿女积蓄。}是的，基于这个理由，你们中许多人以此为自己的贪婪开脱，说你们在为子女积攒，为他们储存。但我告诉你们，扩展你们的爱，让这爱增长；因为爱妻子和儿女，还不是那\Quote{婚宴礼服}。要怀着对\Person{天主}的信德。首先爱\Person{天主}。将你们自己延伸到\Person{天主}那里；尽你们所能，将每一个人引向\Person{天主}。这里有你的仇敌：让他被引向\Person{天主}。这里有儿子、妻子、仆人：让他们都被引向\Person{天主}。这里有陌生人：让他被引向\Person{天主}。这里有仇敌：让他被引向\Person{天主}。拉，拉你的仇敌；通过拉他，他就不再是你的仇敌。因此，让爱德得到推进，这样被滋养，得到滋养后得以成全；这样穿上\Quote{婚宴礼服}；这样，通过我们的前进，\Person{天主}的肖像——我们正是按此肖像受造的——在我们内重新被刻印。因为罪使这肖像受损、磨损。如何受损？如何磨损？当它被擦在地上时？什么是\Quote{当它被擦在地上时}？当它因世俗的贪欲而磨损时。因为\BibleQuote{人虽带着这肖像行走，却徒然烦扰不安。}在\Person{天主}的肖像中寻找的是真理，而非虚幻。因此，借着对真理的爱，那按此肖像受造的肖像重新被刻印，并将他应缴的税赋归还给我们的\Person{凯撒}。因为你们已从\Person{主}的回答中听到，当犹太人试探祂时，祂说：\BibleQuote{你们为什么试探我，你们这些假善人？拿\OriginalTerm{税钱}{tribute money}给我看，}也就是肖像的\OriginalTerm{印记与刻文}{impress and superscription}。告诉我你们付的是什么，你们准备的是什么，你们被征收的是什么。他们给祂看了一个\OriginalTerm{德纳}{denarius}；祂问这上面是谁的\OriginalTerm{肖像与刻文}{image and superscription}。他们回答说：\Quote{\Person{凯撒}的。}所以\Person{凯撒}寻找他自己的肖像。\Person{凯撒}不愿他所命令制造的东西损失，\Person{天主}当然也不愿祂所创造的东西损失。我的弟兄们，\Person{凯撒}并没有制造钱币；铸币厂的匠师制造它；工人按命令行事，他向他的臣仆下达指令。他的肖像被印在钱币上；钱币上刻着\Person{凯撒}的肖像。然而，他要求别人所印的归给他；他将它存入他的宝库；他不会让人拒绝给他。\Person{基督}的钱币就是人。在他身上有\Person{基督}的肖像，有\Person{基督}的名号，\Person{基督}的恩赐，\Person{基督}的职责规条。\par

\SourceRecord{Translated by R.G. MacMullen.；Nicene and Post-Nicene Fathers, First Series；Vol. 6.；Edited by Philip Schaff.；Buffalo, NY: Christian Literature Publishing Co.,；1888.；<http://www.newadvent.org/fathers/160340.htm>.}

% structure: p0001 source=h1 target=section reason=page-title

\section{讲道41：论新约}

\Emph{[XCI. Ben.]}\par

\Emph{论福音之言}，\BibleRef{玛 22:42} \Emph{，主问犹太人，他们以为达味是谁的儿子。}\par

1. 当犹太人被问（正如我们刚才从所读的福音中听到的），我们的主耶稣基督，达味自己称祂为主的那位，怎么是达味之子，他们无法回答。因为他们在主身上所看见的，他们知道。因为祂以人子的身份向他们显现；但作为天主之子，祂是隐藏的。因此，他们以为祂可以被战胜，并在祂悬在木头上时嘲笑祂，说：\Quote{祂若是天主之子，就从十字架上下来，我们就相信祂。} 他们看见了祂的一方面，却不知道另一方面，\BibleQuote{因为他们若认识祂，就不会钉死光荣的主了。} 然而他们知道基督应该是达味之子。因为直到现在他们还盼望祂会来。他们不知道祂已经来了，但他们的无知是自愿的。因为即使他们在木头上没有认出祂，他们也不应该不认识祂坐在宝座上的身份。因为万民是因谁的名字而被呼召和祝福呢？不就是他们认为不是基督的那一位吗？因为这位达味之子，即 \BibleQuote{按肉身是达味的后裔}，是亚巴郎之子。既然曾对亚巴郎说过：\BibleQuote{地上万民都要因你的后裔蒙受祝福}；而他们现在看见在我们的基督里万民都受了祝福，为什么他们还等待已经来了的，而不害怕那将要来的呢？因为我们的主耶稣基督，使用了一个先知性的证词来维护祂的权威，称自己为 \Quote{磐石}。是的，这样一块磐石，\BibleQuote{凡绊倒在这磐石上的，必会跌碎；那磐石落在谁身上，就会把他压得粉碎。} 因为当这磐石被绊倒时，它是低下的；由于低下，它 \Quote{跌碎} 那绊倒它的人；被高举后，当它降下时，它把骄傲的人 \Quote{压得粉碎}。因此，犹太人已经因那次绊倒而 \Quote{跌碎} 了；还剩下的是，在祂光荣的来临中，他们也要被 \Quote{压得粉碎}，除非在他们还活着的时候，他们就承认祂，以免死亡。因为天主是忍耐的，天天邀请他们来信德。\par

2. 但当犹太人不能回答主提出的问题，问他们 \Quote{他们说基督是谁的儿子；} 他们回答 \Quote{达味之子；} 祂继续向他们提出进一步的问题，\Quote{那么，达味怎么在圣神里称祂为主，说：上主对我主说：你坐在我右边，等我使你的仇敌做你的脚凳。如果达味}，祂说，\Quote{在圣神里称祂为主，祂怎么是他的儿子？} 祂没有说 \Quote{祂不是他的儿子，}但说 \Quote{祂怎么是他的儿子？} 当祂说 \Quote{怎么，}这并非否定的字，而是询问，好像祂对他们说：\Quote{你们确实说得好，基督是达味之子，但达味自己称祂为主；他所称为主的那位，怎么是他的儿子？} 如果犹太人曾受过我们持有的基督信仰的教导；如果他们没有封闭自己的心而不接受福音，如果他们希望自己里面有属灵的生命，他们就会，如同受过教会信仰教导的人，回答这个问题并说：\BibleQuote{因为在起初已有圣言，圣言与天主同在，圣言就是天主：}看，祂怎么是达味的主。但因为\BibleQuote{圣言成了血肉，居住在我们中间；}看，祂怎么是达味之子。但因无知，他们沉默了，也没有在闭口时打开耳朵，以便他们被问到时不能回答的，可以在接受教导后知道。\par

3. 但看到这是一件大事，要知道这奥秘：祂怎么既是达味之子，又是达味的主；同一人怎么既是人又是天主；在人的形式下祂比父小，在天主的形式下祂与父同等；祂又怎么说，一方面 \BibleQuote{父比我大，}另一方面 \BibleQuote{我与父原是一体；}看到这是一个伟大的奥秘，我们的行为必须被塑造，以便能理解它。因为向不配的人是封闭的，向那些适合的人则是敞开的。我们不是用石头、棍棒、拳头或脚跟来叩击主。是生命在叩击，是生命被打开。寻求在心中，求问在心中，叩门在心中，敞开心扉。现在那颗正确求问、正确叩门和正确寻求的心，必须是虔诚的。必须先为天主本身而爱天主（因为这就是虔诚）；并且不为自己设定任何从天主那里期待的回报，除天主本身以外。因为没有什么比祂更好。一个人能向天主求什么宝贵的东西，如果在他眼中天主本身被轻看？祂赐予地，你就欢喜，你这爱地的人，你自己成了地。如果当祂赐予地上的好处时，你欢喜，那么当祂赐予你祂自己——祂造了天地——时，你更应当欢喜！因此，必须为天主本身而爱天主。因为魔鬼不知道圣若伯心中所想的，就把这作为一个重大的指控加在他身上，说：\BibleQuote{若伯岂是无缘无故敬畏天主呢？}\par

4. 那么，如果那控告者提出了这项指控，我们就应当害怕这指控会落到我们头上。因为我们所面对的是一个极其恶毒的控告者。如果他要捏造并不存在的事，那么对于真实存在的事，他又会如何竭力地指控呢？然而，让我们喜乐吧，因为我们的审判者是这样的一位，祂不会被我们的控告者所欺骗。因为如果我们有一位人做审判者，那仇敌就可以随意为他捏造。因为没有谁比魔鬼更精于捏造了。就是他在如今也捏造一切针对圣人的虚假指控。他知道他的指控在天主面前毫无效力，所以他将这些指控散布在人间。但这对他有什么益处呢？正如宗徒所说：\BibleQuote{我们的夸耀是这一点：我们良心的见证。}然而，你们以为他捏造这些虚假指控是毫无诡计的吗？是的，他很清楚，如果信德的警觉不抵抗他，他会借此造成怎样的邪恶。因为正是为此，他把对善人的恶毒指控散布出去，好让软弱的人以为世上没有善人，于是他们便任由自己被欲望冲昏头脑，成为欲望的掠物，同时心里说：\Quote{有谁遵守天主的诫命？有谁保持贞洁？}而当一个人认为没有人做到时，他自己就成了那个没有人。这就是魔鬼的伎俩。但\Person{约伯}是这样一个人，魔鬼无法对他捏造任何这样的指控；因为他的生活太过众所周知、显而易见。但因为\Person{约伯}拥有大量财富，魔鬼就把那一点拿来指控他；假如那一点存在，它可能隐藏在内心深处，而不表现在行为上。\Person{约伯}敬拜天主，施舍财物；他怀着怎样的心做这些事，没有人知道，连魔鬼自己也不知道；但天主知道。天主为自己的仆人作证；魔鬼毁谤天主的仆人。魔鬼被允许试炼他，\Person{约伯}被证明是义人，魔鬼被挫败。\Person{约伯}被发现是为了天主本身而敬拜天主，为了天主本身而爱天主；不是因为天主给了他什么，而是因为天主没有将祂自己收回。因为他说：\BibleQuote{主赐予，主取去；主看着好，就如此成就了，愿主的名受赞美。}诱惑的火临近了他；但它发现他是金子，不是碎秸；它除去了渣滓，却没有将他烧成灰烬。\par

5. 那么，为了理解天主的奥秘——基督如何既是人又是天主——心必须被洁净；而心是通过善行、纯洁的生活、贞洁、圣洁和爱，以及通过\BibleQuote{以爱德行事的信德}来洁净的（现在我所讲的这一切，就好比一棵树，它的根在心内；因为只有从心的根才能生出行动；如果你在其中种下欲望，就会长出荆棘；如果你种下爱德，就会结出好果子）。主在向犹太人提出那个问题之后，当他们无法回答时，祂立刻就转而谈论善行，为要表明他们为何不配理解祂所问的。因为当那些骄傲而可怜的人无法回答时，他们当然应该说：\Quote{我们不知道；老师，请告诉我们。}但他们没有：他们在问题面前哑口无言，没有开口求教。因此主针对他们的骄傲立刻说：\BibleQuote{你们要提防经师，他们喜爱会堂里的首位和筵席上的首座。}不是因为他们的职位，而是因为他们喜爱这些。因为这些话指控了他们的心。没有人能指控心，除非那能洞察心的人。因为天主的仆人在教会中担任某种尊贵职务时，应当被给予首位；因为如果不给他，拒绝给予的人就有祸了；但对接受者而言，这本身并不是好事。因此，在基督徒的集会中，主教应当坐在显眼的位置，好让他们因自己的座位而被识别出来，并使他们的职分得到应有的体现；但并非因此他们就可以因座位而骄傲，而应将其视为一项负担，他们要为这负担交账。但谁知道他们是喜爱这一点还是不喜爱？这是内心的事，只有天主才能审判。主亲自警告祂的门徒，不要陷入这酵母中，正如祂在另一处所说的：\BibleQuote{你们要提防法利塞人和撒杜塞人的酵母。}当他们以为祂说这话是因为他们没有带饼时，祂回答他们说：\BibleQuote{你们忘了那五个饼使多少人吃饱了吗？}然后经上说\BibleQuote{他们明白祂称他们的教训为酵母。}因为他们贪爱这些现世的暂时之物，既不惧怕永恒的痛苦，也不爱慕永恒的美善。因此他们的心封闭了，不能理解主所问的。\par

6. 那么，天主的教会该如何行，才能明白它首先藉恩宠所相信的呢？它必须使心智有容量，以接受将要赐予它的。而为了做到这一点，为了使心智——也就是说——有容量，主天主暂停了祂的应许，并没有收回它们。祂暂停它们，好使我们伸展自己；而我们伸展自己，为能成长；我们成长，为能达到它们。请看宗徒\Person{保禄}如何向这些暂停的应许伸展自己：\BibleQuote{这并不是说我已经达到了，或已经是成全的了；弟兄们，我不认为我已经获得了；我只做一件事：忘记背后，奋力向前，向着目标直跑，为得天主在基督耶稣内从高天召叫的奖品。}他在地上奔跑，奖品悬在天上。他在地上奔跑，但精神上升。请看他这样伸展自己，看他向着悬着的奖品伸出：\BibleQuote{我向着目标直跑，为得天主在基督耶稣内从高天召叫的奖品。}\par

7. 我们必须继续前行，然而为此无需涂抹脚油、寻找坐骑或预备船只。以心灵的挚爱奔跑，以爱情前行，以爱德攀登。你为何寻找道路？当紧贴基督，祂藉着下降与上升亲自成为了道路。你愿意上升吗？紧抓住那上升者。因为你凭自己无法上升。\BibleQuote{没有人上过天，除了那自天降下的人子，祂在天上。}若无人上升，除了那下降者，即人子、我们的主耶稣，你也愿意上升吗？那么，就当成为那唯一上升者的肢体。因为祂是头，与众肢体合成一人。既然无人能上升，除非在祂的身体内成为祂的肢体，如此便应验了：\BibleQuote{没有人上过天，除了那下降者。}因为你不能这样说：\Quote{看，为何伯多禄上升了？为何保禄上升了？为何宗徒们上升了？既然无人上升，除了那下降者？}对此的答案是：\Quote{伯多禄、保禄和其余宗徒及所有信友，他们从宗徒那里听到什么？\InnerQuote{你们便是基督的身体，各自都是肢体。}如果基督的身体和祂的肢体属于一个整体，就不要把他们分开。因为祂\InnerQuote{离开了父亲和母亲，依附于自己的妻子，两人成为一体。}祂离开了父亲，因为在此祂没有显示与父同等，反而\InnerQuote{空虚了自己，取了奴仆的形象。}祂也离开了母亲，即按血肉从祂而生的会堂。祂依附于祂的妻子，即祂的教会。在基督亲自提出这证据的地方，祂表明婚姻的约束不可解除：\InnerQuote{你们没有读过吗？那起初造人的天主，造了他们一男一女，并说：两人要成为一体？所以，天主所结合的，人不可拆散。}\InnerQuote{两人要成为一体}是什么意思？祂接着说：\InnerQuote{因此，他们不再是两个，而是一体。}这样，\InnerQuote{没有人上过天，除了那下降者。}}\par

8. 为使你们知道，新郎与新妇依照基督的肉身是一体，而非依照祂的天主性（因为依照祂的天主性，我们不能成为祂；祂是造物主，我们是受造物；祂是创造者，我们是祂的化工；祂是塑造者，我们被祂塑造；但为使我们能在祂内与祂合一，祂屈尊成为我们的头，从我们取了血肉，在其中为我们而死）；为使你们知道这整个就是一位基督，祂藉依撒意亚说：\BibleQuote{祂给我束上冠冕如同新郎，又给我披上盛装如同新娘。}因此祂同时是新郎和新妇。也就是说，在祂自己作为头时是新郎，在身体中则是新娘。\BibleQuote{因为二人，}祂说，\BibleQuote{将成为一体；所以如今不再是两个，而是一体。}\par

9. 既然我们是祂的肢体，为使我们理解这奥迹，正如我所说的，弟兄们，让我们圣洁地生活，为天主本身而爱天主。如今，祂在我们旅途中向我们显示奴仆的形象，却为那些到达家乡的人保留了天主的形象。祂以奴仆的形象铺设了道路，以天主的形象预备了家园。既然我们难以理解这事，却不难相信它；因为依撒意亚说：\BibleQuote{除非你们相信，你们不会明白。}让我们\BibleQuote{凭信德行走，因为我们还在主内作客旅，直到我们来到面对面看见的境地。}当我们凭信德行走时，让我们行善工。在这些善工中，要有为天主本身而发的自由之爱，以及活跃的近人之爱。因为我们不能为天主做什么；但因我们能给近人做什么，我们将藉着对穷人的善行，获得那一切丰盛之源者的恩宠。所以，让每个人尽其所能为他人做事；让他将多余的慷慨施予穷人。有人有钱；让他喂养穷人，给赤身者衣服，建造教堂，用他的钱财尽一切可能行善。有人有善意的劝告；让他引导近人，以圣洁的光明驱散疑虑的黑暗。有人有学问；让他从这主的宝库中取出，分给同仆们食物，坚固信友，召回迷途者，寻找失丧者，尽一切可能行善。即使穷人也彼此有东西可施予；有人借脚给跛子，有人用自己的眼睛引导盲人；有人探访病人，有人埋葬死者。这些都是所有人都能做的事，总之，很难找到一个人没有任何方式行善。最后，还有宗徒所说的那重要责任：\BibleQuote{你们要彼此承担重担，如此便满全了基督的法律。}\par

\SourceRecord{Translated by R.G. MacMullen.；Nicene and Post-Nicene Fathers, First Series；Vol. 6.；Edited by Philip Schaff.；Buffalo, NY: Christian Literature Publishing Co.,；1888.；<http://www.newadvent.org/fathers/160341.htm>.}

% structure: p0001 source=h1 target=section reason=page-title

\section{新约讲道集 第42篇}

\Emph{[XCII. Ben.]}\par

\Emph{论同一福音的话}，\BibleRef{玛 22:42}\par

1. 应向犹太人提出的问题，基督徒应该解答。因为主耶稣基督向犹太人提出了这一问题，却没有亲自向他们解答——我是指向犹太人，祂没有解答，但为我们解答了。亲爱的诸位，我要提醒你们，你们将会发现祂已解答了。但首先请考虑问题的症结。祂问犹太人，他们\Quote{以为基督是谁？祂该是谁的儿子？}因为犹太人也期待基督。他们在先知书中读到祂，期待祂来临，祂来了他们却杀害了祂；因为他们读到基督将降临的地方，也读到他们将杀害基督。然而，他们在先知书中盼望祂未来的降临；因为他们没有预见到自己未来的罪行。因此，祂这样问他们关于基督的事，并非仿佛是一位他们不认识、或从未听过其名、或从未盼望其来临的对象。因为他们仍盼望着祂，甚至在此点上犯错。而我们也确实盼望祂；但我们盼望祂是作为审判者降临，而非被审判。因为圣先知们预言了这两点：祂将先来临，不公地被审判；而后将带着正义降临审判。于是祂说：\Quote{你们以为基督怎样？祂是谁的儿子？}他们回答说：\Quote{达味的儿子。}这完全符合圣经。但祂说：\Quote{那么，达味怎么在圣神内称祂为主，说：\InnerQuote{上主对我主说：你坐在我右边，等我把你的仇敌放在你的脚下。}如果达味在圣神内称祂为主，祂怎么又是他的儿子呢？}\par

2. 此处需要谨慎，免得有人认为基督否认了自己是达味之子。祂并未否认自己是达味之子，而是在探究方式。你们说基督是达味之子，我不否认；但达味称祂为主；告诉我，祂怎么既是达味之子，又是他的主？告诉我吧？他们没有回答，而是沉默。那么，让我们借着基督自己的解释来说明。在哪里？通过祂的宗徒。但首先，我们如何证明基督亲自解释了？宗徒说：\Quote{你们愿意接受在我内说话的基督的证据吗？}因此，在宗徒内，祂屈尊解答了这个问题。首先，基督通过宗徒对弟茂德说了什么？\Quote{你要记念耶稣基督，按我的福音，从达味的后裔中复活了。}看，基督是达味之子。祂怎么又是达味的主？请告诉我们，宗徒：\Quote{祂虽具有天主的形体，却没有以自己与天主同等为应当把持不舍的。}请承认达味的主。如果你承认达味的主、我们的主、天地的主、天使的主、与天主同等、具有天主的形体，祂怎么又是达味的儿子？注意下文。宗徒通过说\Quote{祂虽具有天主的形体，却没有以自己与天主同等为应当把持不舍的}向你显示了达味的主。祂怎么又是达味的儿子？\Quote{祂却使自己空虚，取了奴仆的形体，成了人的模样；以人的形状出现，祂谦卑了自己，听命至死，且死在十字架上。因此，天主极其举扬了祂。}基督\Quote{从达味的后裔}，达味之子，复活了，因为\Quote{祂使自己空虚}。祂如何\Quote{使自己空虚}？通过取了祂所不是的，而非失去祂所是的。祂\Quote{使自己空虚}，祂\Quote{谦卑了自己}。祂虽是天主，却显现为人。祂在地上行走时被轻视，祂，那造了天的主。被轻视，仿佛只是一个凡人，仿佛毫无能力。是的，不仅被轻视，还被杀害。祂就是那块躺在地上的石头，犹太人绊在上面，动摇。祂自己说什么？\Quote{凡跌在这石头上的，必被摔碎；这石头落在谁身上，就要把谁砸得粉碎。}起初，祂低卧，他们绊倒了；祂将从上面降临，将那些已经动摇的\Quote{砸得粉碎}。\par

3. 这样你们就听到了，基督既是达味之子，又是达味的主：永远是达味的主，时间是达味之子；达味的主，由圣父的本体所生；达味之子，由童贞玛利亚所生，因圣神受孕。让我们紧握两者。其中之一将是我们永久的居所，另一个是我们从现世流放中的解救。因为如果我们的主耶稣基督没有屈尊成为人，人就灭亡了。祂成了祂所造的，好使祂所造的不致灭亡。真天主，真人；天主和人，整个基督。这就是公教信德。谁否认基督是天主，就是斐提尼派；谁否认基督是人，就是摩尼教徒。谁明认基督是与圣父同等的天主和真人，祂真实受苦，真实流了血（因为真理若为我们付了虚假的赎价，就不会解救我们）；谁明认两者，就是公教徒。他有家乡，他有道路。他有家乡：\Quote{在起初已有圣言}；他有家乡：\Quote{祂虽具有天主的形体，却没有以自己与天主同等为应当把持不舍的}。他有道路：\Quote{圣言成了血肉}；他有道路：\Quote{祂使自己空虚，取了奴仆的形体}。祂是我们前往的家乡，祂是我们行走的道路。让我们借着祂走向祂，我们就不会走错。\par

\SourceRecord{Translated by R.G. MacMullen.；Nicene and Post-Nicene Fathers, First Series；Vol. 6.；Edited by Philip Schaff.；Buffalo, NY: Christian Literature Publishing Co.,；1888.；<http://www.newadvent.org/fathers/160342.htm>.}

% structure: p0001 source=h1 target=section reason=page-title

\section{新约讲道集 第四十三篇}

\Emph{[第九十三篇，本笃会辑]}\par

\Emph{论福音的话}，\BibleRef{玛 25:1}，\Emph{\BibleQuote{那时，天国好比十个童女。}}\par

昨天在场的人记得我的许诺；在主助佑下，今天要向你们，也要向许多其他聚在一起的人实现。这不是一个容易的问题：十个童女是谁，其中五个明智，五个愚昧。然而，按这段经文的上下文——今天我希望你们再听一次——亲爱的诸位，据主恩赐我的理解，我以为这个比喻或譬喻并非只关乎那些因独特而更卓越的圣德而被称为教会中贞女的女子——我们通常也更常用一个词称她们为\Quote{修道者}；但若我没有弄错，这个比喻关乎整个教会。但即使我们只理解为那些被称为\Quote{修道者}的人，难道他们只有十个吗？天主不许！那么一大群贞女竟然减少到如此小的数目！但也许有人会说：\InnerQuote{可是，如果他们在外表上虽如此众多，实际上却如此稀少，连十个都难找到呢！}并非如此。因为如果祂的意思只是让那十个指善贞女，祂就不会在其中安排五个愚昧的。因为如果这就是蒙召的贞女数目，为什么大宅的门对五个关闭呢？\par

2. 因此，亲爱的诸位，让我们明白，这个比喻关乎我们所有人，即整个教会，不仅关乎我们昨天所说的圣职人员，也不仅关乎平信徒，而是一般地关乎所有人。那么，为什么童女是五个和五个？这五个和五个童女就是所有基督徒灵魂。但让我告诉你们，在主启示下我所认为的：它们不是每一种灵魂，而是那些拥有大公信仰，且似乎在教会中拥有善工的灵魂；可是即使在他们中，\Quote{五个是明智的，五个是愚昧的。}那么首先让我们看看为什么她们被称为\Quote{五个}，又为什么被称为\Quote{童女}，然后再考虑其余部分。每一个在身体中的灵魂因此以数字五来象征，因为灵魂运用五种感官。因为我们对身体的一切感知，都通过这五重门户：或视觉、或听觉、或嗅觉、或味觉、或触觉。凡因此戒绝非法的看见、非法的听见、非法的嗅闻、非法的品尝、非法的触摸，由于他的无玷，便获得了童女之名。\par

3. 但如果戒绝感官的非法兴奋是好的，并且为此每个基督徒灵魂都得了童女之名，为什么五个被接纳，五个被拒绝？她们都是童女，却仍被拒绝。她们是童女，这还不够；她们有灯，这也不够。她们是童女，因戒绝了感官的非法放纵；她们有灯，因有善工。关于这些善工，主说：\BibleQuote{你们的光当在人前照耀，好叫他们看见你们的善工，光荣你们在天之父。}祂又对门徒说：\BibleQuote{你们腰间当束上带，灯也当点着。}在\Quote{束上的腰}上是童贞；在\Quote{点着的灯}上是善工。\par

4. 童贞的称号通常不用于已婚者；但即使在已婚者中，也有一种信德的童贞，产生婚配的贞洁。为使你们明白，圣洁的弟兄们，每一个人、每一个灵魂，就灵魂本身而言，那因信仰的无玷——由此实践戒绝非法之事，并由此行善工——未始不可称为\Quote{童女}；整个教会，由童女、男孩、已婚男女组成，以一个名字被称为贞女。我们从何处证明这一点？请听宗徒的话，他不是只对献身的妇女，而是对整个教会说：\BibleQuote{我曾把你们许配给一个丈夫，要把你们当作贞洁的童女献给基督。}因为要提防魔鬼——这童贞的败坏者，宗徒说了\BibleQuote{我曾把你们许配给一个丈夫，要把你们当作贞洁的童女献给基督}之后，又加上：\BibleQuote{但我担心，正如蛇用它的诡诈诱惑了厄娃，你们的心意也会败坏了，失去那在基督内的纯朴。}在身体上持守童贞的人少；在心中，所有人都应当持守。既然如此，如果戒绝非法之事是好的——由此获得童贞之名，而善工是可称赞的——由灯所象征；为什么五个被接纳，五个被拒绝？如果有一个童女，一个提灯的人，却未被接纳；那么那个既未保持戒绝非法之事的童贞，又因不愿行善工而行在黑暗中的人，将置身何处呢？\par

5. 我诸位弟兄，是的，我们应当更详细地讨论这些事。凡不愿看见罪恶、不愿听见罪恶、对违法的熏香和祭献的非法食物掩鼻而去、拒绝拥抱他人妻子、将自己的饼分给饥饿的人、接待旅客、赤身露体者给他衣穿、使争论者和好、探访病人、埋葬死者的人，他必定是童贞女，他必定有灯。我们还要求什么？有人会说：你还在求什么？我还在求一件事；是福音使我追寻。它说就连这些童贞女，拿着灯，也有明智和愚昧的。我们如何看出这点？我们如何区分？藉着油。这油象征某种伟大而极其伟大的事物。你们以为它不是爱德吗？我们这是探索它是什么而说的；我们不作轻率的判断。我告诉你们为什么爱德似乎由油来象征。宗徒说：\BibleQuote{我指示你们一条更高超的道路。}即使我说人言和天使的话语，若没有爱德，我便成了鸣的锣、响的钹一般。这\Quote{爱德}就是\Quote{那条更高超的道路}，它合理地由油来象征。因为油浮于所有液体之上。倒进水，再倒上油，油会浮在上面。倒进油，再倒进水在它上面，油仍会浮在上面。你若保留平常的顺序，它在最上面；你若颠倒顺序，它还是最上面。\Quote{爱德永不坠落。}\par

6. 那么，诸位弟兄，这又是什么意思呢？我们现在来讨论那五个明智和五个愚昧的童贞女。她们想去迎接新郎。所谓\Quote{出去迎接新郎}是什么意思？就是用心灵出去，等候他的来临。但他延迟了。\Quote{当他延迟时，她们都睡了。}\Quote{都}是什么意思？就是愚昧的和明智的，\Quote{都打盹睡着了。}我们认为这睡眠是好的吗？这睡眠是什么？是不是因为新郎延迟，以致\BibleQuote{由于罪恶增多，许多人的爱德冷淡了}？我们是否这样理解这睡眠？我不以为然。我告诉你们为什么。因为在她们中间也有明智的童贞女；的确，当主说：\BibleQuote{由于罪恶增多，许多人的爱德冷淡了}时，祂接着又说：\BibleQuote{唯独坚持到底的，才可得救。}你要把这些明智的童贞女放在哪里？她们不是属于那些\Quote{坚持到底}的人吗？诸位弟兄，她们得以进入里面，没有别的原因，就是因为她们\Quote{坚持到底}。因此，爱德的冷气没有侵入她们，她们的爱德没有冷淡；反而保持着炽热直到终局。因为她们直到终局仍保持炽热，所以新郎的门为她们敞开；所以她们被吩咐进去，如同那良善的仆人：\BibleQuote{进入你主人的福乐吧！}那么，她们\Quote{都睡了}是什么意思？有另一种睡眠，没有人能避免。你们岂不记得宗徒说：\BibleQuote{弟兄们，关于安眠的人，我不愿你们无知，}即关于死去的人？为什么他们被称为\Quote{安眠的人}，岂不是因为在他们的日子里？因此\Quote{她们都睡了}。你以为因为一个人明智，就不必死吗？无论是愚昧的童贞女还是明智的童贞女，都同样遭受死亡的睡眠。\par

7. 但人们常常自言自语说：\Quote{看，审判的日子现在就要来了，发生了如此多的灾祸，如此多的艰难在加剧；看，先知们所说的一切，几乎都实现了；审判的日子已经近在眼前了。}凡这样说，且怀着信德这样说的人，就等于是以这样的心思出去\Quote{迎接新郎}。但是，看！战争不断，艰难不断，地震不断，饥荒不断，民族互相攻击，而新郎仍未到来。就在人们期待祂来临的时候，所有那些说：\Quote{看，祂要来了，审判的日子会在这里找到我们，}的人，都睡着了。当他们说这话时，他们就睡着了。因此，每个人都要留意自己的睡眠，并坚持在爱德中直到睡眠；让睡眠发现他如此等候。因为假如他睡着了，\Quote{那睡了的人岂不能再起来吗？}因此，比喻中明智和愚昧的童贞女，\Quote{都睡了}，经上说：\Quote{她们都睡了。}\par

8. \Quote{看，在半夜有人喊叫。}\Quote{半夜}是什么意思？就是没有任何期待，丝毫也不相信的时候。夜晚代表无知。一个人在内心中计算：\Quote{看，从\Person{亚当}以来已经过了这么多年，六千年即将结束，然后按照某些诠释者的计算，审判的日子立即就会来到；}然而这些计算来了又过了，新郎的来临仍然延迟，那些出去迎接祂的童贞女都睡着了。而看，当祂不被期待时，当人们说：\Quote{六千年被等待，看，已经过去了，那么我们将如何知道祂何时来呢？}祂要在半夜来临。\Quote{要在半夜来临}是什么意思？就是当你不知道的时候来临。为什么祂要在你不知道的时候来临？听主自己说：\BibleQuote{父以自己的权柄所定的时候和日期，你们是没有必要知道的。}宗徒说：\BibleQuote{主的日子要像夜间的盗贼一样来到。}因此，你要在夜间警醒，免得被盗贼所乘。至于死亡的睡眠——不管你愿不愿意——它总会来到。\par

9. \BibleQuote{但当那喊声在半夜发出之际。} 这喊声是什么？岂非宗徒所说的：\BibleQuote{一眨眼，在末次号角之时}？\BibleQuote{因为号角将响，死人将要复活，成为不朽的，我们也要被改变}。因此当那喊声在半夜发出时：\BibleQuote{看，新郎来了}；接着是什么？\BibleQuote{于是那些童女都起来}。\Quote{她们} 全都起来是什么意思？主亲自说：\BibleQuote{时候将到，凡在坟墓里的都要听见他的声音，并要出来}。所以在末次号角时，她们全都起来了。\BibleQuote{那些明智的童女带着灯，也在器皿里预备了油；但愚拙的却没有预备油}。\BibleQuote{没有在器皿里预备油} 是什么意思？\BibleQuote{在器皿里} 是什么意思？在他们的心里。因此宗徒说：\BibleQuote{我们的夸耀就是良心的见证}。这就是油，珍贵的油；这油是天主的恩赐。人可以把油放在器皿里，却不能创造橄榄。看，我有油，但你造了油吗？这是天主的恩赐。你有油，就带着它。\Quote{带着它} 是什么意思？就是持守在内心，在那里取悦天主。\par

10. 因为，看，那些\BibleQuote{没有预备油的愚拙童女}，想藉着她们的贞洁（因此她们被称为童女）和善行来取悦于人，当她们似乎提着灯的时候。如果她们想取悦于人，并为此行一切可称赞的事，她们就没有带着油。那么，你要带着它，带在内心，在天主看见的地方；在那里带着良心的见证。因为谁若为了得到别人的见证而行，他就没有带着油。如果你禁戒不法之事，并因受人的称赞而行善，那么里面就没有油。所以，当人们停止称赞时，灯就熄灭了。所以，亲爱的，在这些童女睡觉之前，经上没有说她们的灯熄灭了。明智童女的灯因内在的油、因良好良心的确证、因内在的荣耀、因最深处的爱德而燃烧。然而，愚拙童女的灯也燃烧着。为什么那时它们燃烧呢？因为那时还没有缺少人的称赞。但在她们起来以后——即在从死者中复活时——她们开始整理她们的灯，即准备向天主交出自己的作为。因为那时没有人称赞，每个人都完全忙于自己的事，那时没有人不为自己着想，因此没有人卖油给他们；于是她们的灯开始熄灭，愚拙的便转向那五个明智的，说：\BibleQuote{请把你们的油分些给我们，因为我们的灯快要灭了}。她们寻求她们素常寻求的：即靠着别人的油发光，靠别人的称赞行事。\BibleQuote{请把你们的油分些给我们，因为我们的灯快要灭了。}\par

11. 但她们说：\BibleQuote{不必这样，恐怕不够我们和你们用的；你们不如到卖油的那里去买吧}。这不是给建议者的回答，而是嘲笑者的回答。为什么她们嘲笑呢？因为她们是明智的，因为智慧在她们内。因为她们的明智不是出于自己；而是那智慧在她们内，正如某卷书上所写：当那些藐视智慧的人陷入智慧所警告的灾祸时，智慧要说：\BibleQuote{我也必笑你们的灾祸}。那么，明智的童女嘲笑愚拙的童女有什么奇怪呢？这嘲笑又是什么意思呢？\par

12. \BibleQuote{你们到卖油的那里去买吧}：你们这些人，从来不是因为人称赞你们（这些卖油的人）才生活得好。这话何意？\Quote{卖油给你们}？\Quote{出卖称赞}。谁出卖称赞？奉承者。你们若没有顺从奉承者，并带着内心的油，为良心的缘故而行一切善事，那该多好！那么你们就能说：\BibleQuote{义人用仁慈纠正我、责备我，罪人的油却不可抹在我的头上}。他倒说：让义人纠正我，让义人责备我，让义人打击我，让义人纠正我，而不要叫\BibleQuote{罪人的油抹在我的头上}。罪人的油是什么？岂非奉承者的甜言蜜语？\par

13. \BibleQuote{你们到卖油的那里去}，这是你们素常所做的。但我们不会给你们。为什么？\BibleQuote{恐怕不够我们和你们用的}。\BibleQuote{恐怕不够} 是什么意思？这话并非出于缺乏希望，而是出于清醒而虔诚的谦卑。因为善人虽有良好的良心，但他怎能知道那位不受任何人欺骗的主将如何审判呢？他有良好的良心，没有罪恶在内心搅扰他；然而，尽管良心良好，但因日常生活中的罪过，他向天主说：\BibleQuote{宽免我们的罪债}，因为他做到了下一句：\BibleQuote{如同我们也宽免得罪我们的人}。他从内心深处把食物分给饥饿的人，从内心深处给赤身者衣服穿；借着内在的油，他行了善工，然而在那审判中，就连他的良好良心也要战栗。\par

14. 那么请看，这\BibleQuote{给我们油吧}是什么意思。他们被告知：\BibleQuote{你们到卖油的那里去吧}。因为你们习惯于靠人的表扬而活，你们没有随身带油；但我们不能给你们；\BibleQuote{恐怕不够我们和你们用的}。因为我们自己还难以判断自己，何况判断你们呢？\Quote{我们自己还难以判断自己}是什么意思？因为，\BibleQuote{当正义的君王坐在宝座上时，谁能夸耀自己的心是纯洁的呢？}也许你不在自己的良心里发现什么；但那位看得更清楚者，祂的圣目洞察更深，或许会发现什么，祂或许会看见什么，祂发现什么。你向祂说\BibleQuote{求你不要与你的仆人进入审判}是多么好啊！是的，\BibleQuote{宽免我们的罪债}是多么好啊！因为也因那些火炬、那些灯，将对你们说：\BibleQuote{我饿了，你们给了我吃的}。那么怎样？那些糊涂的童女不是也这样做了吗？是的，但她们不是在祂面前做的。那么她们是如何做的呢？正如主所禁止的，祂说：\BibleQuote{你们应当小心，不要在人前行你们的正义，为叫他们看见；否则，你们在天之父就没有赏报了。当你们祈祷时，不要像伪君子那样，因为他们喜欢站在街角祈祷，为叫人看见。我实在告诉你们，他们已经获得了他们的赏报。}她们买了油，付了代价；她们买了，没有被人表扬所欺骗，她们寻求人的表扬，并且得到了。这些人的表扬在审判之日对她们无益。但其余的童女是怎么做的呢？\BibleQuote{你们的光当在人前照耀，使他们看见你们的善行，光荣你们在天之父。}祂并没有说\Quote{光荣你们}。因为你们自己没有油。夸耀你自己说，我有它；但来自祂，\BibleQuote{你有什么不是领受的呢？}因此，一方是这样做的，另一方是那样做的。\par

15. 难怪，在\BibleQuote{她们去买的时候}，同时她们寻找可受表扬的人，却找不到；同时她们寻找可受安慰的人，却找不到；那时门开了，\BibleQuote{新郎来了}，新娘——教会——那时与基督一起受光荣，好使各个肢体被聚集到整体之中。\BibleQuote{她们同他进去赴婚宴，门就关上了。}后来那些糊涂的童女来了；但她们买了油吗？或者找到了可以向之买油的人吗？因此她们发现门关了；她们开始敲门，但已经太迟了。\par

16. 经上说，这话是真的，不是骗人的话：\BibleQuote{敲门，就给你们开门}；但此时是仁慈的时候，不是审判的时候。因为这两个时期不能混淆，因为教会向她的主咏唱\BibleQuote{仁慈与审判}。现在是仁慈的时候；悔改吧。你在审判的时候能悔改吗？那时你就像那些童女，门对着她们关上了。\BibleQuote{主啊，主啊，给我们开门吧}。怎么！她们没有悔改吗？她们没有带油吗？是的，但当真正的智慧嘲笑她们时，她们迟来的悔改有什么益处呢？因此\BibleQuote{门就关上了}。又对她们说了什么？\BibleQuote{我不认识你们}。难道祂不认识他们吗？祂知道一切。那么\BibleQuote{我不认识你们}是什么意思？我拒绝，我舍弃你们。在我的技艺中，我不承认你们，我的技艺不认识恶；现在这是一件奇妙的事，它不认识恶，却审判恶。它在实践中不认识恶；它通过斥责来审判。因此，\BibleQuote{我不认识你们}。\par

17. 五个明智的童女来了，\BibleQuote{进去了}。你们有多少人呢，我的弟兄们，在基督圣名的宣认中！愿你们中间有五个明智的，但不要只有这样五个人。愿你们中间有五个明智的，属于这数字五的智慧。因为那个时辰要来临，而且在我们不知道的时候来临。它要在半夜来临。你们要警醒。福音就这样结束：\BibleQuote{所以你们要警醒，因为你们不知道那日子，也不知道那时辰}。但如果我们都将睡觉，我们怎么能警醒呢？用心灵警醒，用信德警醒，用望德警醒，用爱德警醒，用善行警醒；然后，当你们在身体内睡觉的时候，时辰来到，你们就要复活。当你们复活后，预备好你们的灯。那时它们不再熄灭，那时它们会被良心的内在之油更新；那时那位新郎将在祂属灵的拥抱中怀抱你们，那时祂将把你们带入祂的家中，在那里你们永不睡觉，你们的灯永不熄灭。但现在我们正在劳苦，我们的灯在这生命的风和诱惑中闪烁；但只愿我们的火焰猛烈燃烧，以致诱惑的风反而增加火焰，而不是熄灭它。\par

\SourceRecord{Translated by R.G. MacMullen.；Nicene and Post-Nicene Fathers, First Series；Vol. 6.；Edited by Philip Schaff.；Buffalo, NY: Christian Literature Publishing Co.,；1888.；<http://www.newadvent.org/fathers/160343.htm>.}

% structure: p0001 source=h1 target=section reason=page-title

\section{关于新约的讲道集 第44讲}

\Emph{[XCIV. Ben.]}\par

\Emph{论福音的话，\BibleRef{玛 25:24}等，即懒惰的仆人因不肯将他所领受的\OriginalTerm{塔冷通}{talent}拿出来放贷而被定罪者。}\par

诸位大人、我的弟兄、诸位同仁主教屈尊来探望我们，以他们的临在使我们喜乐；但我不知为何他们不愿在我疲惫时帮助我。亲爱的，我在他们面前对你们说这些话，为使你们的听闻能在某种程度上为我向他们求情，好使我请求他们时，他们也可轮流对你们讲论。让他们分施他们所领受的，让他们屈尊工作而非推辞。然而，请你们听我这疲惫难言的人只说几句话。因为我们还有一份记录，关于天主通过一位神圣殉道者赐予的仁慈，我们必须全心聆听。那么，是什么呢？我该对你们说什么呢？你们已在福音中听到了善仆的应得赏报和恶仆的惩罚。而那被弃绝并受严厉谴责的仆人的全部邪恶，就在于他不肯用他的金钱去放贷。他保留了他所领受的全部款项；但主期望从中获利。天主在我们的救恩上是贪婪的。如果那不放贷者尚且如此受罚，那些丢失所领受者的人又将面临什么呢？那么，我们是分施者，我们放贷，你们收受。我们期望利润；你们要生活得好。因为这就是我们与你们交往中的利润。但你们不要认为这放贷的职责也不属于你们。你们确实不能从这高座上执行它，但你们可以在你们所到之处执行。无论何处基督受攻击，你们要保卫祂；答复抱怨者，斥责亵渎者，与他们保持距离。这样，你们若使任何人获益，你们便是放贷了。在你们自己的家中履行我们的职务。主教之得名即由此而来，因为他监督、照顾并关注他人。因此，每个人，若他是自己家庭的首领，便应当拥有\OriginalTerm{主教职}{Episcopate}的职责，照顾他的家眷如何信仰，使他们中无一人陷入异端，无论是妻子、儿子、女儿，甚至他的奴隶，因为他是以如此巨大的代价被赎买的。宗徒的教导将主人置于奴隶之上，并将奴隶置于主人之下；然而基督为两者都付出了同样的代价。那么，不要忽视你们中最小的成员，要全神贯注地照顾你们全家的救恩。你们若这样做，就是放贷；你们将不会是懒惰的仆人，你们将不必恐惧如此可怕的惩罚。\par

\SourceRecord{Translated by R.G. MacMullen.；Nicene and Post-Nicene Fathers, First Series；Vol. 6.；Edited by Philip Schaff.；Buffalo, NY: Christian Literature Publishing Co.,；1888.；<http://www.newadvent.org/fathers/160344.htm>.}

% structure: p0001 source=h1 target=section reason=page-title

\section{新约讲道集 第四十五篇}

\Emph{[XCV. Ben.]}\par

\Emph{论\WorkTitle{福音}的话，\BibleRef{谷 8:5}等，其中记述了七个饼的奇迹。}\par

1. 我给你们讲解圣经，就如同为你们擘饼。你们当怀着饥饿领受，从心田涌出丰满的言辞；你们这些在宴席上富足的人，莫在善行善事上匮乏。我所分给你们的，原非我所有。你们所吃的，我也吃；你们赖以生活的，我也赖以生活。我们在天上有一个共同的仓库，因为天主圣言从那里而来。\par

2. \Quote{七个饼}象征圣神的七重化工；\Quote{四千人}象征建立在四福音上的教会；\Quote{七筐碎块}象征教会的圆满。因为七这个数字极常象征圆满。\Quote{我一日七次赞美你}这句话从何而来？难道一个人若不如此多次赞美主，便有罪吗？那么\Quote{七次赞美}是什么意思？岂不是\Quote{我永不止息地赞美}？因为说\Quote{七次}的人是指一切时间。故此在这世界上有连续七日的循环。那么\Quote{我一日七次赞美你}是什么意思？不就是另一处所说的\Quote{赞美祂常在我口中}吗？为这圆满的缘故，\Person{若望}写信给七个教会。\WorkTitle{默示录}是圣史\Person{若望}的书，他写信\Quote{给七个教会}。你们当饥饿，领受这些筐。因为那些碎块并未失落；既然你们也属于教会，它们必然有益于你们。我向你们解释这些，是在服侍基督；你们平安地听，你们\Quote{坐下}。我身体坐，心中却站着，且战战兢兢服侍你们；恐怕食物或器皿冒犯你们中任何人。你们知道天主的宴席，你们常听见，那是为心，并非为腹。\par

3. 果然，四千人用七个饼吃饱了；还有什么比这更奇妙呢！然而这还不够，七筐碎块也装满了。奥义啊！这些是行动，行动在说话。你们若明白这些作为，它们就是言语。你们也属于那四千人，因为你们生活在四福音之下。孩童和妇女不包括在这数目中。因为经上记载：\Quote{吃的人，除了妇女孩子，约有四千。}仿佛无理解和软弱的都不算在数目内。然而让这些也吃吧。让他们吃；也许孩子会长大，不再做孩子；也许软弱者会改正，成为贞洁。让他们吃；我们分施，我们分发给他们。但这些是谁，天主察看他自己的宴席，如果他们不改过，那位知道如何邀请他们来的，也知道如何把他们从众人中分别出来。\par

4. 你们知道，亲爱的诸位；请回忆福音的比喻，主怎样来察看祂某次宴席上的宾客。那位邀请他们的家主，正如经上记载，\Quote{在那里发现一个没有穿婚宴礼服的人}。因为那位新郎曾邀请他们赴婚宴，祂是\Quote{在世人中最美丽的}。那新郎为了祂那不端庄的新娘，变得不端庄，为使她能美丽。那美丽的一位如何变得不端庄？如果我不证明，就是亵渎。那位美丽的见证，先知给我说：\Quote{你在世人中最美丽。}那丑陋的见证，另一位先知给我说：\Quote{我们看见祂，没有威仪，也没有美丽；祂的面容被毁伤，容貌全无光彩。}先知啊，你说\Quote{你在世人中最美丽}；你被反驳了；另一位先知出来反对你，说：\Quote{你说谎。我们见过祂。你怎么说\InnerQuote{你在世人中最美丽}？}难道这两位先知在和平的基石上不一致吗？二人都论及基督，都论及基石。在角上，墙壁连接。若不连接，就不是建筑，而是废墟。不是的，先知们是一致的，不要让他们争吵。是的，我们应明白他们的和平；因为他们不懂得争吵。先知啊，你说\Quote{你在世人中最美丽}；你在哪里看见祂？回答我，你在哪里看见祂？\Quote{祂虽具有天主的形体，却不以自己与天主同等为僭越。}我在那里看见祂。难道你怀疑那\Quote{与天主同等}的，是\Quote{在世人中最美丽的}吗？你已回答；现在让那说\Quote{我们看见祂，祂没有威仪，也没有美丽}的来回答。你这么说；告诉我们你在哪里看见祂？他从另一人的话开始；另一人结束之处，他起始。另一人在哪里结束？\Quote{祂虽具有天主的形体，却不以自己与天主同等为僭越。}看，他在那里看见了那\Quote{在世人中最美丽的}；你告诉我们，你在哪里看见祂\Quote{没有威仪，也没有美丽}。\Quote{祂却使自己空虚，取了奴仆的形体，成为人的模样，且以人的形态出现。}关于祂的丑陋，他更说：\Quote{祂贬抑自己，听命至死，且死在十字架上。}看，我在那里看见了祂。因此，二人都在和平的和谐中，一同和平。有什么比天主更美丽？有什么比被钉者更丑陋？\par

5. 如此，这位新郎，\BibleQuote{美丽超过世人子}，却变得丑陋，为要使祂的新娘美丽，就是那被吩咐说\BibleQuote{女子中美丽的}，关于她也曾说\BibleQuote{那上升的、被漂白的是谁}，因光明而非虚假的色彩而被漂白！祂既召唤他们赴婚宴，遇见一个没有穿婚宴礼服的人，就对他说：\BibleQuote{朋友，你没有穿婚宴礼服，怎么进到这里来？他哑口无言。}因为他找不出话来回答。那邀请他的家主说：\BibleQuote{捆起他的手脚，丢到外面的黑暗中去；在那里要有哀号和切齿。}为这么小的过失，这么大的惩罚？它确实是大的。不穿婚宴礼服，被称为小过失，但那只是对不明白的人而言。若这不是极其严重的过失，祂何至于如此愤怒，如此审判，因他没有穿婚宴礼服，就\BibleQuote{捆起手脚丢到外面的黑暗中，在那里有哀号和切齿}呢？我说这话，是因为你们是借着我被邀请的；虽然祂邀请你们，却是借着我的职务邀请。你们都在宴席上，要有婚宴礼服。我要解释那是什么，使你们都能拥有它；若现在有人听我说话却没有它，愿他在家主来视察宾客之前改过自新，领受\Quote{婚宴礼服}，然后满怀信心地坐下。\par

6. 事实上，亲爱的诸位，那从宴席中被丢出去的人，并非指一个人，远非如此。他们有很多。主自己设了这个比喻，这位新郎亲自召集宾客赴宴，并赐生命给祂所召叫的人，祂亲自向我们解释了，那人并非代表一个人，而是许多人，就在那里，同一个比喻中。我不需要远寻，我在那里找到了解释，我在那里擘饼，摆在你们面前吃。因为祂说，当那没有\Quote{婚宴礼服}的人被丢到外面黑暗中去时，祂立刻接着说道：\BibleQuote{因为被召的人多，被选的人少。}你从那里丢出一个人，却说\BibleQuote{因为被召的人多，被选的人少。}无疑，被选的人不会被丢出去，他们就是少数留下的宾客；而\Quote{许多人}由那一个人代表，因为那没有\Quote{婚宴礼服}的人乃是恶人的身体。\par

7. 什么是\Quote{婚宴礼服}？让我们在圣经中寻找。什么是\Quote{婚宴礼服}？无疑，它是善人与恶人所不共有之物；让我们发现它，我们就发现了\Quote{婚宴礼服}。在天主的恩赐中，善人与恶人不共有的是什么？我们是人而非野兽，这是天主的恩赐，但善人与恶人共有之。天上的光照耀我们，云中的雨降下，泉水涌流，田地出产果实；这些是恩赐，但善人与恶人共有之。让我们进入婚宴，让那些被召而不来的人留在外面。让我们考虑那些宾客本身，即基督徒。圣洗是天主的恩赐，善人与恶人都领受之。祭台的圣事，善人与恶人一同领受。撒乌耳尽管作恶，在他大发烈怒迫害一位圣洁至义的人时，仍然说预言。难道只有善人才相信？\BibleQuote{魔鬼也信，且战栗。}我当如何？我筛过一切，尚未找到\Quote{婚宴礼服}。我展开了我的包裹，考虑了所有或几乎所有，仍未找到那件礼服。宗徒保禄在某处给我带来了许多卓越之物的集合；他将其展开在我面前，我对他说：\Quote{请给我看，若你从中找到了那件\InnerQuote{婚宴礼服}。}他开始逐一展开，说：\BibleQuote{我若能说人间的语言和天使的语言，有全备的知识和先知之恩，有全备的信德以致能移山，我若将一切财产分施给穷人，并舍身投火被焚。}珍贵的衣物！然而，其中还没有\Quote{婚宴礼服}。现在请向我们展示\Quote{婚宴礼服}。宗徒啊，你为何让我们悬疑？也许先知之恩是天主的恩赐，但善人与恶人都不拥有它？\BibleQuote{如果我没有爱德，毫无益处。}看，这就是\Quote{婚宴礼服}；穿上它吧，宾客们，好使你们安稳地坐下。不要说：\Quote{我们太穷，没有那件礼服。}你穿给别人，自己就被穿上了。现在是冬天，给裸体者衣服穿。基督是赤裸的；祂将赐予你们那件\Quote{婚宴礼服}，凡没有它的人。跑到祂那里，恳求祂；祂知道如何圣化祂的信徒，祂知道如何给祂的裸体者穿上衣服。为使你们能拥有\Quote{婚宴礼服}，免于对外面黑暗以及手足被捆的恐惧，你们不要使自己的行为落空。若它们落空，手被捆着，你能做什么？脚被捆着，你能逃往何处？那么，持守那件\Quote{婚宴礼服}，穿上它，当祂来视察时，好安稳地坐下。审判之日将要来到；祂现在赐予长时间的宽限，让先前裸体的人现在穿上衣服吧。\par

\SourceRecord{Translated by R.G. MacMullen.；Nicene and Post-Nicene Fathers, First Series；Vol. 6.；Edited by Philip Schaff.；Buffalo, NY: Christian Literature Publishing Co.,；1888.；<http://www.newadvent.org/fathers/160345.htm>.}

% structure: p0001 source=h1 target=section reason=page-title

\section{新约讲道第四十六篇}

\Emph{[第九十六篇 圣奥斯定]}\par

\Emph{讲论福音中\BibleRef{谷 8:34}的话：\Quote{谁若愿意跟随我，该弃绝自己}等等；并讲论\BibleRef{若一 2:15}的话：\Quote{谁若爱世界，父的爱就不在他内。}}\par

1. 主所命令的，即\BibleQuote{谁若愿意跟随我，该弃绝自己}，似乎艰难而沉重。但那命令我们的主帮助我们，使祂的命令得以实现，因此祂的命令并不艰难或沉重。因为圣咏中向祂所说的这句话是真实的：\BibleQuote{因祢嘴唇的话语，我谨守了艰难的道路。}而祂自己所说的这句话也是真实的：\BibleQuote{我的轭是柔和的，我的担子是轻的。}因为凡所命令的艰难之事，爱德能使它变得容易。我们知道爱情本身能成就何等大事。这爱情常常甚至是可憎与不洁的；但人为达到他们所爱的对象，忍受了多少艰难，承受了多少侮辱与不可忍受之事？无论是爱钱财者被称为贪财者，或爱荣誉者被称为野心家，或爱美女者被称为纵欲者。谁能列举所有种类的爱？然而，试想所有爱者所受的劳苦，他们却不自觉其劳苦；而当他们被阻碍不能劳苦时，反而最感劳苦。既然多数人如何，他们的爱也如何，并且我们生活的规范除了选择我们所当爱的之外，别无他虑；那么，你为何惊奇，如果有人爱基督并愿意跟随基督，因爱祂而弃绝自己？因为人若因爱自己而丧亡，无疑因弃绝自己而得以寻获。\par

2. 人的首要毁灭是爱自己。因为假如他不爱自己，假如他将天主置于自己之上，他就愿意永远服从天主；就不会转向忽略天主的旨意而随从自己的私意。因为爱自己就是愿意随从自己的私意。将天主的旨意置于此之上；学习通过不爱自己而爱自己。为使你知道爱自己是一种恶习，宗徒这样说：\BibleQuote{因为人将成为爱自己者。}那爱自己的人怎能对自己有任何可靠的信赖？不能；因为他开始爱自己而离弃天主，并被驱离自己，去爱那些在他之外的事物，以致当上述宗徒说了\BibleQuote{人们将成为爱自己者}之后，立即加上\BibleQuote{爱钱财者。}你已经看见你在外面。你开始爱自己：你如果能，就站在自己里面。为什么你走到外面？因为你富有钱财，你成了爱钱财者？你开始爱在你之外的事物，你丧失了自己。当人的爱甚至从自己转向那些外在事物时，他开始分享其虚妄欲望的虚空，并且仿佛浪子一般耗尽自己的力量。他分散、枯竭、无援、无力，他放猪；厌倦了这放猪的职务，他终于记起自己是什么，并说：\BibleQuote{我父亲有多少雇工都口粮充裕，我却在这里饿死！}但当比喻中的儿子这样说时，关于那曾挥霍一切于娼妓、希望自己掌握那在他父亲那里为他妥善保存的财产的人，他希望能随意处理它，挥霍一切，陷入赤贫，对他怎样说？\BibleQuote{当他回到自己时。}如果\BibleQuote{他回到自己}，他就曾离开过自己。因为他曾从自己堕落，离开了自己，他首先回到自己，以便回到那因离开自己而堕落的境地。因为正如因离开自己而留在自己内，同样因回到自己，他不应留在自己内，免得再离开自己。他回到自己，为了不留在自己内，他说了什么？\BibleQuote{我要起来，到我父亲那里去。}看，他从哪里离开自己？他从父亲那里离开自己；他离开自己，走到外面的事物中。他回到自己，并到他父亲那里去，在那里他可以稳妥地保存自己。如果他曾离开自己，那么在他回到自己时，也离开他曾离开的那位，以便\BibleQuote{到他父亲那里去}，弃绝自己。什么是\BibleQuote{弃绝自己}？让他不信任自己，让他感觉自己是人，并尊重先知的话：\BibleQuote{凡信赖世人的人，是可咒骂的。}让他从自己退出，但不走向低处。让他从自己退出，好与天主结合。凡他所拥有的善，让他交付于造他的主；凡他所拥有的恶，是他自己造成的。天主没有造成他内的恶；让他摧毁自己所造成的，而他自己也因此被摧毁。\BibleQuote{让他弃绝自己}，祂说，\BibleQuote{并背起自己的十字架，跟随我。}\par

3. 那时主必须被跟随到何处？我们知道祂去了哪里；不过几天前我们刚庆祝了那庄严的纪念。因为祂已复活，升了天堂；必须在那里跟随祂。无疑，我们不应对此绝望，因为祂亲自向我们许诺了，并非因人能做什么。天堂离我们很远，在我们的头进入天堂之前。但现在，如果我们那头的肢体，为何要绝望呢？那么，必须在那里跟随祂。谁不愿意跟随祂到那样的居所呢？尤其当我们在地上因恐惧与痛苦如此劳苦时。谁不愿意跟随基督到那里，那里有至高的福乐、至高的平安、永久的安稳？跟随祂到那里是好的：但我们必须看我们该走什么道路。因为主耶稣说我们正在思考的这些话时，祂尚未从死者中复活。祂尚未受难，祂还必须走向十字架，走向祂的羞辱、凌辱、鞭打、荆棘、创伤、嘲笑、侮辱、死亡。这道路似乎是崎岖的；它使你迟缓；你无心跟随。但继续跟随。人自造的崎岖之路，却被基督经过而踏平了。因为谁不希望走向高升呢？升高令众人喜悦；但谦卑是通向它的阶梯。为何你要将脚伸过自己？你意在跌倒，而非攀登。从阶梯开始，于是你已上升。这两个门徒不愿看到这谦卑的阶梯，他们说：\BibleQuote{主，请吩咐，在你王国里，一个坐在你右边，一个坐在你左边。}他们寻求升高，却未看到阶梯。但主指示了阶梯。因为祂如何回答他们？\BibleQuote{你们寻求升高之山，能喝谦卑之杯吗？}因此祂不只是说：\BibleQuote{让那人弃绝自己，跟随我}；而是进一步说：\BibleQuote{让他背起自己的十字架，跟随我。}\par

4. 什么是\BibleQuote{让他背起自己的十字架}？让他承受他所有的任何困难；这样让他跟随我。因为当他开始根据我的生活和诫命跟随我时，会有许多人反对他，许多人阻止他，许多人劝阻他，甚至从那些似乎是基督同伴的人中。那些阻止瞎子呼喊的人是和基督一起行走的。因此，无论是威吓还是安抚，或任何阻碍，如果你想跟随，就将它们变成你的十字架，背负它，扛着它，不要屈服。主这些话似乎是对殉道的劝勉。如果有迫害，难道不应当为基督之故而藐视一切吗？世物被爱；但应更爱创造世物的那一位。世界是伟大的；但创造世界的那一位更伟大。世界是美的；但创造世界的那一位更美。世界是甘饴的；但创造世界的那一位更甘饴。世界是恶的；而创造世界的那一位是善的。我怎能解释并解开我说的这些？愿天主助我？因为我说的什么？你们鼓掌的什么？看，这只是一个问题，而你们已鼓掌。如果创造世界的那一位是善的，那么世界如何是恶的？天主岂没有创造一切，\BibleQuote{看，都很好}？圣经岂不在每一个创造工程中见证天主创造它是好的，说\BibleQuote{天主看了认为好}，最后又将其全部总结，即天主创造了它们，\BibleQuote{看，都很好}？\par

5. 那么，世界如何是恶的，而创造世界的那一位是善的？如何？\BibleQuote{世界是藉着祂造的，世界却不认识祂。}世界是藉着祂造的——天、地及其中万物；\BibleQuote{世界却不认识祂}——爱世界的人；爱世界而轻蔑天主的人；这\BibleQuote{世界却不认识祂}。因此世界是恶的，因为那些喜爱世界胜过喜爱天主的人是恶的。而创造世界——天、地、海及那些爱世界者本身的那一位是善的。因为只有这一点——他们爱世界而不爱天主——祂没有在他们身上创造。但他们自身，凡属于其本性的一切，祂都创造了；凡属于罪过的一切，祂没有创造。这就是我刚才说的：\Quote{让人抹去他所造的，这样他才能取悦那创造他的主。}\par

6. 因为在人类自身中也有一个善的世界；但那是由恶变为善的。因为整个<世界>——若你取\Quote{世界}一词指人，撇开我们所称的世界（天、地及其中万物）；若你取世界指人——那首先犯罪者使整个世界变为恶。整个质体在根源中被败坏。天主造人本是善的；经上如此说：\BibleQuote{天主造人原为正；但他们自己找出许多计谋。}从这些\Quote{许多}逃向独一，将你分散之物收束为一；一起流出，围护自己，与独一者同在；不要走向许多事物。那里有福乐。但我们流散了，走向毁灭；我们都带着罪而生，且因恶行加增了生来的罪，于是整个世界成了恶的。但基督来了，祂拣选了祂所造的，而非祂所找到的；因祂找到的全是恶的，却藉祂的恩宠使它们成了善的。于是形成了另一个\Quote{世界}；现在这\Quote{世界}逼迫那\Quote{世界}。\par

7. 那迫害人的\Quote{世界}是什么？就是那对我们所说的：\Quote{你们不要爱世界，也不要爱世界上的事物。谁若爱世界，圣父的爱就不在他内。因为凡在世界上的一切，就是肉身的贪欲、眼目的贪欲和生活的骄傲，都不是出于圣父，而是出于世界。世界和它的贪欲都要过去；但那遵行天主旨意的，却永远存留，}如同天主永远存留一样。看！我讲了两个\Quote{世界}：迫害人的\Quote{世界}和受迫害的\Quote{世界}。那迫害人的\Quote{世界}是什么？\Quote{凡在世界上的一切，就是肉身的贪欲、眼目的贪欲和生活的骄傲，都不是出于圣父，而是出于世界；}并且\Quote{世界都要过去。}看，这就是迫害人的\Quote{世界}。那受它迫害的\Quote{世界}是什么？\Quote{那遵行天主旨意的，却永远存留，}如同天主永远存留一样。\par

8. 但请看，那迫害人的被称为\Quote{世界}；让我们证明那受迫害的是否也被称为\Quote{世界}。怎么！你们难道聋于基督所说的话，或更确切地说，聋于\WorkTitle{圣经}所见证的话：\Quote{天主在基督内使世界与自己和好。}\Quote{如果世界恨你们，你们要知道，它先恨了我。}看，\Quote{世界}在恨。它所恨的不就是\Quote{世界}吗？哪一个\Quote{世界}？\Quote{天主在基督内使世界与自己和好。}那被定罪的\Quote{世界}迫害；那被和好的\Quote{世界}受迫害。那被定罪的\Quote{世界}是在教会之外的一切；那被和好的\Quote{世界}就是教会。因为他说：\Quote{人子来不是为审判世界，而是为叫世界借着他而得救。}\par

9. 如今，在这个世界——圣洁、良善、和好、得救，更确切地说，将要得救，且如今在望德中得救，因为\Quote{我们得救是在于望德}——在这个世界，我说，就是在完全跟随基督的教会内，他曾普遍地说：\Quote{谁若愿意跟随我，该弃绝自己。}因为不是只有童贞女该听从这话，而已婚妇女不必；或只有寡妇该听从，而仍有丈夫的妇女不必；或只有隐修士该听从，而已婚男子不必；或只有圣职人员该听从，而平信徒不必；而是让整个教会、整个身体、所有成员，在各种职务上区分并分布，都跟随基督。让整个教会跟随他，那唯一的教会，让鸽子跟随他，让新娘跟随他，让她——被赎且以新郎的血为嫁妆的——跟随他。那里，童贞的贞洁有其位置；那里，寡妇的节欲有其位置；那里，已婚的贞洁有其位置；但奸淫在那里没有自己的位置；那里也没有邪淫——那非法而该受惩罚的——的位置。然而，让这些各按其类、按其位置和尺度拥有位置的成员\Quote{跟随基督}；让他们\Quote{弃绝自己}；就是说，不要自恃什么；让他们\Quote{背起自己的十字架}；就是说，让他们为了基督在世界忍受世界可能加给他们的一切。让他们爱他，唯有他不欺骗，唯有他不被欺骗，唯有他不骗人；让他们爱他，因为他所应许的是真实的。但因为他没有立刻赐予，信德就动摇了。坚持，坚韧，忍受，承担迟延，你就背起了十字架。\par

10. 不要让童贞女说：\Quote{我将独自在那里。}因为玛利亚不会独自在那里，寡妇亚纳也会在那里。不要让有丈夫的妇女说：\Quote{寡妇会在那里，我不会；}因为并非亚纳会在那里，而苏撒纳却不在。但无论如何，凡愿意在那里的人，要以此证明自己：他们在此处居于较低位置的不嫉妒，反而爱他人更好的位置。例如，我的弟兄们，为使你们明白：一人选择了婚姻生活，另一人选择了节欲生活。若那选择婚姻生活的人有奸淫的贪欲，他就算\Quote{回头看了}；他贪图了非法的事。那后来愿意从节欲生活返回婚姻生活的人，也\Quote{回头看了}；他选择了本身合法的事，却\Quote{回头看了}。那么，婚姻该受谴责吗？不。婚姻不该受谴责；但请看那选择它的人走到了哪里。他已经走到了它前面。当他年轻时生活在纵情享乐中，婚姻在他前面；他正向它走去；但当他选择了节欲，婚姻就在他后面。主说：\Quote{你们要记得罗特的妻子。}罗特的妻子因回头一看，就停住了。因此，无论一个人已到达什么地步，都要惧怕从那里\Quote{回头}；并要行在路上，要\Quote{跟随基督}。\Quote{忘记背后，努力面前的，向着目标直奔，要得天主在基督耶稣内召我向上获得的奖品。}让已婚者视未婚者在自己之上；让他们承认他们更好；让他们在他们身上爱自己所没有的；并在他们身上爱基督。\par

\SourceRecord{Translated by R.G. MacMullen.；Nicene and Post-Nicene Fathers, First Series；Vol. 6.；Edited by Philip Schaff.；Buffalo, NY: Christian Literature Publishing Co.,；1888.；<http://www.newadvent.org/fathers/160346.htm>.}

% structure: p0001 source=h1 target=section reason=page-title

\section{\WorkTitle{新约讲道第47篇}}

\Emph{[XCVII. Ben.]}\par

\Emph{论福音的话}，\BibleRef{谷 13:32} \Emph{，\Quote{至于那日子或那时刻，没有人知道，连天上的天使也不知道，连子也不知道，惟独父知道。}}\par

1. 诸位弟兄，你们刚才听到圣经所给的劝告，它叫我们要警醒等候末日。每个人都应当把这视为关乎他自己的末日，免得当你判断或以为世界的末日还很遥远时，你却在自己的末日里沉睡了。你们听到了耶稣论及这世界末日所说的话：\Quote{连天上的天使也不知道，连子也不知道，惟独父知道。} 这里确实有一个很大的难点，免得我们用属肉体的方式去理解，认为父知道什么子却不知道。因为当祂说\Quote{父知道它}时，祂这样说是因为在父内，子也知道它。因为有什么日子不是由圣言所造，而日子正是借圣言而造的呢？因此，没有人应当追究末日何时来临；但愿我们都藉着良好的生活警醒，免得我们任何一个人的末日发现我们毫无准备。无论一个人在他最后的日子以什么状态离开此世，他在世界的末日也要被发现是那样的。那时，你在今世没有做过的事，将丝毫不会帮助你。各人的作品将帮助他，或各人的作品将压倒他。\par

2. 我们在圣咏中是怎样向主歌唱的：\Quote{上主，求你怜悯我，因为人把我踏倒了}？\Quote{人}是指那按人的方式生活的人。因为对那些按天主生活的人，经上说：\Quote{你们是神，你们都是至高者的子女。} 但对于那些被召成为天主的子女却宁愿做人（即按人的方式生活）的败类，祂说：\Quote{你们却要像人一样死亡，像一位首领一样倒毙。} 因为人是会死的，这理当有助于他的教化，而不是自夸。明天就要死的小虫有什么可夸的呢？我向你们的爱德说，诸位弟兄，骄傲的凡人应当被魔鬼羞愧。因为魔鬼虽是骄傲的，却是不朽的；他是一个灵，尽管是恶灵。末日为他存留，作为他最终的惩罚；然而他不受我们所受的死亡之约束。但人听到了判决：\Quote{你必定要死。} 让他善用他的惩罚吧。我所说的\Quote{让他善用他的惩罚}是什么意思？让他不要因为他所受的惩罚而陷入骄傲；让他承认他是会死的，并以此打破他的高傲。让他听这话对他说的：\Quote{为什么尘土和灰烬自高自大？} 即使魔鬼是骄傲的，他也不是\Quote{尘土和灰烬}。因此经上说：\Quote{你们却要像人一样死亡，像一位首领一样倒毙。} 你们不考虑自己是会死的人，却像魔鬼一样骄傲。因此人应当善用他的惩罚，诸位弟兄；让他善用他的恶，以便转向他的善。谁不知道，我们必须死亡这一事实是一种惩罚，而且更严重的是，我们不知道它何时来临？惩罚是确定的，时辰是不确定的；并且只有这个惩罚在人类事务的普通进程中是我们确定的。\par

3. 我们所有其他的事，无论是好是坏，都不确定；只有死亡是确定的。我说的是什么意思？一个孩子被怀上，也许他会出生，也许会成为不足月的胎儿。所以这不确定：也许他会长大，也许他不会长大；也许他会变老，也许他不会变老；也许他会富有，也许贫穷；也许他会显赫，也许卑贱；也许他会有子女，也许没有；也许他会结婚，也许不会；如此类推，凡你所能举出的好事都是这样。现在再看看生活的祸患：也许他会生病，也许不会；也许会被蛇咬，也许不会；也许会被野兽吞噬，也许不会。所以看看所有的祸患：处处都有\Quote{也许会}和\Quote{也许不会}。但你能说\Quote{也许他会死}和\Quote{也许他不会死}吗？就像医生检查疾病，确定是致命的，他们这样宣布：\Quote{他会死，他不会好过来的。} 因此，从人出生的那一刻起，就可以说：\Quote{他不会好过来的。} 当他出生时，他就开始患病了。当他死时，他确实结束了这种病：但他不知道他是否陷入更糟的境地。福音中的那个富人结束了他纵欲的病，却来到了折磨的病。那个穷人结束了他的病，却得到了完全的健康。但他在今生选择了将来所要拥有的；他在那里收割的，是他在这里播种的。因此，当我们活着时，我们应该警醒，并选择我们在将来世界所要拥有的。\par

4. 我们不要爱世界。它淹没那些爱它的人，引导他们无益。我们更应在其中劳苦，免得它诱惑我们，而不要惧怕它会倾倒。看，世界倾倒，基督徒却屹立不倒，因为基督不倾倒。为什么主说：\BibleQuote{你们欢喜吧，因为我已经战胜了世界}？我们若愿意，可以回答祂说：\Quote{\InnerQuote{欢喜}，是的，祢欢喜。祢若\InnerQuote{战胜了}，祢欢喜。为什么我们应当欢喜呢？}为什么祂对我们说：\BibleQuote{欢喜}，岂不是因为祂为我们战胜了，为我们战斗了？祂在何事上战斗呢？在于祂取了人性。除去祂的童贞诞生，除去祂自甘空虚，\BibleQuote{取了奴仆的形象，成为人的样式，且以人的形态出现}；除去这些，那么战斗在哪里？争战在哪里？考验在哪里？胜利在哪里？那胜利没有战斗作先导。\BibleQuote{在起初已有圣言，圣言与天主同在，圣言就是天主。万物是藉着祂造成的，凡受造的，没有一样不是由祂造成的。}犹太人能钉死这圣言吗？那些不敬的人能嘲弄这圣言吗？这圣言能受掌掴吗？这圣言能戴荆棘冠冕吗？但为能承受这一切，\BibleQuote{圣言成了肉体}，在承受这一切之后，祂藉着复活\Quote{战胜了}。所以祂为我们\Quote{战胜了}，祂向我们显示了祂复活的确实凭据。于是你向天主说：\BibleQuote{主，求你怜悯我，因为人将我践踏。}你自己不要\Quote{践踏}自己，人就不能战胜你。因为看，一个有权势的人恐吓你。他凭什麼恐吓你？\Quote{我要掠夺你，定你的罪，折磨你，杀死你。}于是你呼喊：\BibleQuote{主，求你怜悯我，因为人将我践踏。}如果你说的是实话，并反省自己，一个必死的人\Quote{践踏你，}因为你害怕人的威胁；而人\Quote{践踏你，}因为若不是你也是人，你就不会害怕。那么补救之法是什么？人啊，你要依附于天主，祂使你成为人；紧紧依附祂，信赖祂，呼求祂，让祂做你的力量。你对祂说：\BibleQuote{主，我的力量在于祢。}那时你面对人的威胁就要歌唱；而你要歌唱什么，主亲自告诉你：\BibleQuote{我要寄望于天主，不惧怕人能对我做什么。}\par

\SourceRecord{Translated by R.G. MacMullen.；Nicene and Post-Nicene Fathers, First Series；Vol. 6.；Edited by Philip Schaff.；Buffalo, NY: Christian Literature Publishing Co.,；1888.；<http://www.newadvent.org/fathers/160347.htm>.}

% structure: p0001 source=h1 target=section reason=page-title

\section{新约讲道集第48篇}

[第九十八篇 圣奥思定]\par

\Emph{论福音之言}——\BibleRef{路 7:2}\Emph{，等；论主所复活的三位死者。}\par

1. 吾主耶稣基督救主所行的奇迹，凡听见并相信的人，无不深受感动，但因人而异。有些人惊异于祂在人的身体上所行的奇迹，却不知分辨那更大的；另有些人则赞叹在现今人的灵魂上所实现的更圆满的满全，他们听说这些事曾在人的身体上发生过。主亲自说：\BibleQuote{就如父复活死者，使之生存；同样，子也随意使谁生存。}\BibleRef{若 5:21}当然，不是子\Quote{使}一些人生存，父使另一些人生存；而是父与子\Quote{使}同一些人；因为父藉着子行一切事。所以，凡是基督徒，不可怀疑现今仍有死者复活。所有人都有眼睛，能看见死者像那位寡妇的儿子一样复活，我们刚从福音中读到；但看见内心死了的人复活的那些眼睛，并非所有人都有，除非他们自己已在内心复活了。复活一个将永远活下去的人，比复活一个必将再死的人，是一个更大的奇迹。因为主若没有来复活死者，宗徒就不会说：\BibleQuote{醒来，你这睡着的人，从死者中起来，基督必光照你。}\BibleRef{弗 5:14}\par

2. 那位寡妇母亲因那青年的复活而喜乐；而慈母教会则因每日在心灵上复活的人而喜乐。他固然肉体死了，他们却灵魂死了。他那可见的死被人可见地哀悼；他们那不可见的死既未被探究，也未被察觉。祂寻找那些祂知道已死的人；唯独祂知道他们是死的，也只有祂能叫他们活。因为主若没有来复活死者，宗徒就不会说：\BibleQuote{醒来，你这睡着的人，从死者中起来，基督必光照你。}\BibleRef{弗 5:14}你听见睡着的人时，在话中：\BibleQuote{醒来，你这睡着的人}；但当你听见\BibleQuote{从死者中起来}时，你明白这是指死者。所以，那些甚至身体已死的人经常被称为睡着。的确，在祂能唤醒他们这一点上，他们都只是睡着。因为就你而论，死人确实是死的，因为无论你怎么打他、刺他、撕他，他都不会醒来。但就基督而论，那个听到\BibleQuote{起来}就立刻起来的人，只是睡着了。没有人能像基督在坟墓中那样轻易地唤醒他人。\par

3. 现在我们发现，主显然复活了三位死者，\Quote{有形可见地}复活了数千人\Quote{无形不可见地}。不，谁知道祂复活了多少可见的死者呢？因为祂所行的一切并没有都记下来。若望告诉我们说：\BibleQuote{耶稣所行的还有许多别的事，假使把它们一一写出来，我想世界上的书也容不下。}\BibleRef{若 21:25}所以无疑还有许多其他人被复活了；但三位的明确记载并非没有意义。因为吾主耶稣基督愿意祂在身体上所行的事也得到属灵的理解。因为祂行奇迹不仅是为奇迹本身，而是使祂所行的事在看见的人心中引起惊叹，并向理解的人传达真理。正如一个人看到一份优美的写本，却不认识字，他固然赞美抄写者的手，惊叹字形的优美；但他不知道那些字的意思或含义；他眼之所见赞美这作品，但心中并不理解；而另一个人既赞美作品，又能理解它；我是指那种既能看见人人共见的事物，又能阅读的人；从未学过的人则不能。所以那些看见基督的奇迹却不懂其含义、不懂它们在某种意义上传达给有理解力的人什么的人，只惊叹奇迹本身；而另一些人既惊叹奇迹，又领悟其含义。我们在基督的学校里应当如此。因为有人说基督行奇迹只是为了奇迹本身，那么他也可以说，当祂在无花果树下找果子时，祂不知道那不是结果子的时候。因为正如福音记载，那不是结果子的时候；但祂饿了，仍在树上找果子。难道基督不知道任何一个农夫都知道的事吗？难道树的主人不知道，树的创造者反而不知道吗？所以当祂饿了在树上找果子时，祂表示祂饿了，是在寻求别的什么东西；祂发现那棵树没有果子，却长满了叶子，就诅咒它，它立刻枯干了。那棵树不结果子有什么过错？不结果子难道是树的罪过吗？不是；而是有些人因自己的意愿不能结果子。不结果子是\Quote{他们}的过错，因为他们的意愿本可结果子。因此，犹太人拥有法律的言语，却没有实行，他们满是叶子，却不结果子。我说这些是要说服你们，吾主耶稣基督行奇迹是带着这样的意图：藉着这些奇迹，祂要表示更深的意义，除了奇迹本身奇妙、伟大、神圣之外，我们也能从它们学到些什么。\par

4. 让我们看看祂愿我们从祂所复活的三位死者中学到什么。祂复活了会堂长已死的女儿，她在患病时曾向祂请求，求祂救她脱离病痛。祂正前行时，有人报告说她已经死了；似乎祂现在只是徒然劳累，有人给她的父亲传话说：\BibleQuote{你的女儿死了，何必再烦劳师傅呢？} 但祂继续前行，对女孩的父亲说：\BibleQuote{不要怕，只管信。} 祂来到家中，发现已准备好了惯例的丧葬事宜，便对他们说：\BibleQuote{不要哭，这女孩不是死了，而是睡着了。} 祂说的是真理；她睡着了，就是说，对祂而言，她能因祂而醒。于是唤醒她，将她活生生地交还给她的父母。同样，祂也唤醒了那位寡妇的儿子；这个事例提醒我现在要对你们讲论这个主题，亲爱的诸位，正如祂亲自赐给我能力一样。你们刚才听到他是怎样被唤醒的。主\BibleQuote{走近城门，看，正抬出一个死人}，已出了城门。祂动了怜悯之心，因为那母亲是个寡妇，且失去了独子，正在哭泣，祂便行了你们所听到的事，说：\BibleQuote{青年人，我对你说：起来。那死者便坐起来，并开始说话，耶稣便把他交给了他的母亲。} 祂也同样从坟墓中唤醒了拉匝禄。在那件事中，当与祂谈话的门徒们知道他病了时，祂说（\BibleQuote{耶稣爱他}）：\BibleQuote{我们的朋友拉匝禄睡着了。} 他们想着病人的健康睡眠，便说：\BibleQuote{主，他若是睡着了，就必好了。} 于是耶稣便更明白地说：\BibleQuote{我告诉你们，我们的朋友拉匝禄死了。} 在这两件事上祂都说了真话：\Quote{对你们来说，他是死了；对我而言，他是睡着了。}\newline{}\par

5. 这三种死者，是三种罪人，基督至今仍复活他们。因为会堂长的已死女儿是在屋内，尚未从墙内的隐秘处被抬到公众眼前。她在那里被复活，活生生地交还给她的父母。但第二个死者现在已不在屋内，却尚未在坟墓中；已被抬出墙外，但尚未被安葬。那位复活了尚未被抬出的少女的，也复活了现在已被抬出但尚未埋葬的死者。还有第三种情形，就是复活一个已被埋葬的人；这祂在拉匝禄身上行了。因此，有些人在内心暗中怀有罪恶，却尚未在公开行为上表现出来。例如，一个人被任何情欲扰乱。因为主亲自说：\BibleQuote{凡注视妇女，有意贪恋她的，他已在心里奸淫了她。} 他尚未在身体上接近她，但在心里已同意了；他在屋内有一个死者，尚未将他抬出去。而时常发生的是，正如我们所知，正如人每日在自己身上所经验到的，当他们听到天主的话语时，仿佛主说\BibleQuote{起来}；他们谴责对罪的同意，重新呼吸于救恩和正义之中。屋内的死人起来了，心在思想的隐秘处重获生命。这死灵魂的复活发生在内部，在良心的隐秘处，仿佛在屋内的墙壁之间。另一些人在同意之后走向公开行为，仿佛将死者抬出去，使那隐藏于隐秘的显现于公众。这些已走向外在行为的人，就没有希望了吗？福音中对那青年不也说了：\BibleQuote{我对你说：起来}？他不是也被交还给了他的母亲吗？因此，那犯了公开行为的人，若偶然因真理之言被劝诫和唤醒，在基督的声音中起来，就活生生地复原了。他能走到那么远，却不能永远丧亡。但那些藉作恶而将自己陷入邪恶习惯的人，以至于这恶习本身不让他们看见它是恶的，成为自己恶行的辩护者；当被指责时便发怒；甚至如同古时索多玛人对指责他们可憎计谋的义人所说：\BibleQuote{你来是寄居的，不是来做官判案的。} 在那地方，可憎污秽的习惯如此强大，以致放荡被当作正义，而阻止者反比行恶者更受责难。这样的人被恶习压住，仿佛被埋葬了。是的，弟兄们，我该说什么呢？他们被埋葬到那种程度，正如论及拉匝禄所说的：\BibleQuote{现在已经臭了。} 那堆在坟墓上的土堆，就是这顽固的习惯势力，灵魂被压住，既不能起来，也不能重新呼吸。\newline{}\par

6. 如今曾说过：\BibleQuote{他已死了四天。}事实上，灵魂达到我所说的那种习惯，是经由一种四重进展。第一，可以说是内心的愉悦的\OriginalTerm{诱惑}{provocatio}；第二，\OriginalTerm{同意}{consent}；第三，\OriginalTerm{外在行为}{actus}；第四，习惯（\OriginalTerm{习惯}{habitus}。因为有些人如此彻底地将不法之事从他们的思想中排除，甚至对他们毫无愉悦之感。有些人感到愉悦，却不给予同意；死亡尚未完成，但以某种方式开始了。愉悦之感加上同意，便立刻招致谴责。同意之后，进展至公开行为；行为变成习惯；进而产生一种绝望的状态，以致可以说：\BibleQuote{他已死了四天，如今已经臭了。}因此，主来了——当然，在主一切都是容易的——祂却仿佛在此事上遇到一种困难。祂\BibleQuote{心神感伤}，表明需要大量而强烈的规劝，才能唤醒那些因习惯而变得刚硬的人。然而，在主呼喊的声音中，束缚的枷锁被冲破。地狱的势力战栗，拉匝禄重获生命。因为主甚至从恶习中拯救那些\Quote{已死四天}的人；福音中这个\Quote{已死四天}的人，在基督看来只是睡了，因为祂定意要使他复活。但主说了什么？请观察他复活的方式。他从坟墓中活着出来，却无法行走。主对门徒说：\BibleQuote{解开他，让他行走。}\Quote{主}使他从死亡中复活，\Quote{门徒}解开他的束缚。请注意，有些事情属于天主独特的尊威，祂使人复活。一个陷入恶习的人受到真理之言的谴责。有多少人被谴责，却听而不闻！那么，是谁在内里感动那听信的人？是谁在内里赋予他生命？是谁驱散那看不见的死亡，赐予那看不见的生命？在谴责和规劝之后，人们不是独自被留在自己的思想中，难道他们不是开始思考自己过着多么邪恶的生活，被多么糟糕的习惯所压迫吗？然后，他们对自己不满，决心改变生活。这样的人已经复活了；那些对自己过去的生活感到不满的人，已经重获新生；但尽管重获新生，他们却不能行走。这些是他们罪恶的束缚。因此，那重获生命的人需要被解开，被释放。主将这项任务交给了门徒，祂对他们说：\BibleQuote{凡你们在地上所束缚的，在天上也要被束缚。}\par

7. 因此，亲爱的诸位，我们要这样聆听这些事：使活着的人继续活着；使死了的人重新活过来。无论是罪过仅在心中酝酿，尚未成为公开行为，就应悔改，加以纠正，让那死在良心之屋内的人复活。或者，他实际行了所思想的事，即使如此，也不该绝望。那死在屋内的人尚未复活，当\BibleQuote{他被抬出来}时，他就应复活。他应悔改自己的行为，立刻重获生命；不要堕入坟墓深处，不要承受习惯的重负。但或许我现在正对一个被自己习惯的硬石所压迫的人说话，一个已被习俗的重量所压倒的人，一个\Quote{已在坟墓里四天，并且已经发臭}的人。然而，连这样的人也不该绝望；他在深处死了，但基督在高处被高举。祂知道如何以祂的呼喊冲破尘世的重担，祂知道如何亲自在内里恢复生命，并将他交给门徒解开。连这样的人也应当悔改。因为当拉匝禄在四天后复活时，在他活着的时候，恶臭不再停留。所以，让活着的人仍然活着；让死了的人，无论他们是谁，无论他们发现自己处于这三种死亡中的哪一种，都赶快立刻复活。\par

\SourceRecord{Translated by R.G. MacMullen.；Nicene and Post-Nicene Fathers, First Series；Vol. 6.；Edited by Philip Schaff.；Buffalo, NY: Christian Literature Publishing Co.,；1888.；<http://www.newadvent.org/fathers/160348.htm>.}

% structure: p0001 source=h1 target=section reason=page-title

\section{新约讲道集 第四十九篇}

\Emph{[XCIX. Ben.]}\par

\Emph{关于福音的话，\BibleRef{路 7:37}，\Quote{看，城里有一个罪妇，}等。关于罪过的赦免，驳斥多纳徒派。}\par

既然我相信这是天主的旨意，我现在要依照主的话出自圣经所提醒我们的主题，靠祂的助佑，向你们，蒙爱的，讲一篇关于罪过的赦免的讲道。因为当福音被宣读时，你们极其认真地倾听，故事被报道，并在你们心灵的眼前呈现。因为你们不是用身体，而是用理智，看见了主耶稣基督\Quote{在法利塞人家里坐席，}并被祂邀请时，并不不屑于前往。你们也看见了一个城里著名的\Quote{妇人，}确实在恶名中著名，\Quote{是个罪妇，}未受邀请却强行闯入宴会，在那里她的医师正坐席，并以一种神圣的无耻寻求健康。于是她闯了进去，就宴会而言似乎不合时宜，但就她所期望的祝福而言却是合时的；因为她深知自己正受着何等沉重的疾病，她知道她所来的那位能够治愈她；于是她走近了，不是到主的头前，而是到祂的脚前；她既长久行在邪恶中，现在寻求正直的足迹。首先她流泪，心头的血；并以告解的责任洗了主的脚。用她的头发擦干，她亲吻，她抹油：她用沉默说话；她没有发出一句话，但表达了自己的虔敬。\par

所以，因为她摸了主，用眼泪浇灌、亲吻、清洗、抹油祂的脚；那邀请主耶稣基督的法利塞人，看到祂属于那种骄傲人，依撒意亚先知论他们说：\Quote{他们说：\InnerQuote{离我远点，不要摸我，因为我是洁净的。}} 认为主不认识那妇人。他心里这样想，心里说：\Quote{这人如果是先知，就会知道触摸祂脚的那个妇人是怎样的。} 他认为主不认识她，因为主没有拒绝她，没有禁止她靠近，容许她被那罪妇触摸。因为他从哪里知道主不认识她呢？可是，哦，你这法利塞人，邀请主却又嘲笑主的人！你喂养主，却不知道你自己要被谁喂养。你从哪里知道主不知道那妇人过去是怎样的，除非因为她被允许靠近祂，因为因祂的容许她亲吻了祂的脚，清洗了，抹了油？因为一个不洁的妇人不应该被允许对这些洁净的脚做这些事吗？那么，如果这样的妇人接近那法利塞人的脚，他肯定会说出依撒意亚论及这类人所说的话：\Quote{离开我，不要摸我，因为我是洁净的。} 但她带着不洁走近主，为能洁净而归；她带病走近，为能痊愈而归；她凭告解走近祂，为能宣认祂而回归。\par

因为主听见了法利塞人的心思。现在让法利塞人由此明白，那能听见他心思的，难道不能看见她的罪吗？于是主向那人讲了一个关于两个人欠同一个债主债的比喻。因为祂也渴望治好法利塞人，免得白白在他家吃饭。祂渴慕那喂养祂的人，祂愿意改造他，杀死他，吃掉他，将他转入祂自己的身体；正如祂对撒玛黎雅妇人说的：\Quote{我渴。} \Quote{我渴}是什么意思？我渴望你的信德。因此，在这个比喻中说了主的话；其中有一个双重目的，既使那邀请者与同席者一起被治愈，他们同样看见了主耶稣基督，也同样不认识祂，也使那妇人得到她告解所应得的保证，不再被良心的刺痛所折磨。\Quote{一个，}祂说，\Quote{欠了五百个\OriginalTerm{德纳}{denarii}，另一个欠了五十个；祂都赦免了他们：哪一个更爱他？} 那被问比喻的人回答了他当然按常理必须回答的话：\Quote{主，我想是那被赦免更多的。} 然后转身对那妇人，祂对西满说：\Quote{你看见这妇人吗？我进了你的家，你没有给我水洗脚；她却用眼泪洗了我的脚，并用她的头发擦干。你没有给我亲吻；这妇人自从进来，就不住地亲吻我的脚。你没有用油抹我的头；但这妇人却用香膏抹了我的脚。因此，我告诉你，她许多的罪都赦免了，因为她爱得多。但那被赦免少的，爱得也少。}\par

4. 这里出现一个困难，必须真实地加以解决，并需要你们专注的注意，亲爱的朋友们，免得我的言辞因时间紧迫，无法完全消除和澄清其全部晦暗；尤其是因为这副被暑热耗尽的血肉之躯，现在渴望得到恢复，索求它应得的，并阻碍灵魂的渴慕，证明了经上所说的：\BibleQuote{心神固然切愿，但肉体却软弱。} 这里有理由害怕，是的，大理由害怕，免得因主的这些话，在那些不正确理解它们、放纵肉体情欲、又不愿被带离它们进入自由的人心中，潜入了那种甚至在宗徒们宣讲时也出现在毁谤者口中的情绪，关于这些毁谤者，宗徒保禄说：\BibleQuote{也有人毁谤我们，说我们主张\InnerQuote{作恶以成善}。} 因为有人会说：\Quote{如果\BibleQuote{那得赦免少的，爱得也少}；那得赦免多的，爱得也多；爱多比爱少好；那么，我们正应该多犯罪，多欠债，以便渴望被赦免，这样我们就会更加爱那位赦免我们巨债的主。因为在福音中，那位有罪的妇人，正是因为她欠得越多，就越爱那赦免她债的主，正如主自己说的：\BibleQuote{她的许多罪赦了，因为她爱得多。} 如今她为什么爱得多呢？岂不是因为她欠得多吗？后来主又加上说：\BibleQuote{但那得赦免少的，爱得也少。} 岂不是更好吗，} 他可能说，\Quote{赦免我的多而不是少，以便我能更爱我的主？} 你们无疑看到了这个困难的深奥；是的，我肯定你们看到了。你们也看到了我时间的紧迫；是的，你们也看到并感觉到这一点。\par

5. 那么，请听几句话。如果我不能公正地处理这个问题的重大性，请暂且存留我现在所说的，并把我当作欠债者，以待将来。假设有两个人，好借着更清楚的例证使你们思考我所提出的问题。其中一人罪孽深重，长久以来生活极其邪恶；另一人只犯了些小罪。两人都来到恩宠前，都接受了圣洗，他们作为负债者进入，作为自由者离开；一人被赦免得多，另一人被赦免得少。我问，每人爱了多少？如果我发现在赦罪最多的人中，他爱得最多，那么他多犯罪对他更有利，他的大罪对他更有利，这样他的爱才不至于冷淡。我问另一人爱多少，发现他爱得少；因为如果我发觉他也爱了，与被赦免得多的人同样多，那么我如何回答主的话呢？真理所说的\BibleQuote{那得赦免少的，爱得也少}又怎能成立呢？\Quote{看，}有人说，\Quote{我被赦免的少，我没有多犯罪；然而我爱得与被赦免得多的人一样多。} 你说的是真话，还是基督？你的谎言被赦免，难道是为了让你去指责那赦免你的人说谎吗？如果赦免你的少，你爱得也少。因为如果赦免你的虽少，你爱得却很多，你就与那说\BibleQuote{那得赦免少的，爱得也少。}的主相矛盾了。因此，我更相信那位比你更了解你的主。如果你认为赦免你的少，那么可以肯定你爱得也少。\Quote{那么，}他说，\Quote{我该怎么办呢？我应当多犯罪，好让他能够赦免我许多罪，使我能够爱得更多吗？} 这使我为难，但愿那向我们说出这真理之言的主，救我脱离这困境。\par

6. 这话是为那法利塞人说的，他以为自己没有罪，或只有少许罪。其实，他若没有一些爱，就不会邀请主。但这份爱何等渺小！他没有给主行亲吻礼，连洗脚的水也没有，更不用说眼泪；他没有以那妇人所行的一切尊敬之礼来尊崇主；那妇人深知自己何等需要被治愈，也知道谁能治愈她。啊，法利塞人，所以你爱得少，因为你妄自认为你被赦免的少；并非真的赦免得少，而是你以为被赦免的只是少许。\Quote{那么，}他说：\Quote{难道我从未杀人，就被算作杀人犯吗？难道我从未犯奸淫，就要为奸淫受罚吗？或者这些我从未犯过的罪，也要被赦免吗？}看：再假设两个人，让我们对他们说话。一人前来哀求，一个罪人，满身荆棘如刺猬，又极其胆怯如野兔。但磐石是刺猬和野兔的避难所。他于是来到磐石前，找到避难所，得到援助。另一人没有犯许多罪；我们该为他做什么，使他能爱得多？我们该劝他什么？我们要违背主的话吗？\Quote{那赦免少的，爱得也少。}是的，千真万确，那真正被赦免少的。可是，啊，你说你没有犯许多罪：你为何没有？因谁的引导？感谢天主，因你的动作和声音表明你已明白我。那么，我想，难题已经解决了。一人犯了许多罪，因此欠下许多债；另一人因天主的引导只犯了少许罪。那被赦免的，一人归功于主，另一人也把他未犯的罪归功于主。在你过去无知的生命中，你尚未被光照，尚未分辨善恶，尚未信靠那在你不知时引导你的主，你未曾犯奸淫。你的天主这样对你说：\Quote{我为你引导我自己，我为你保守我自己。为使你不犯奸淫，没有诱惑者靠近你；没有诱惑者靠近你，是我的作为。时机和地方都缺乏；它们缺乏，又是我的作为。或者诱惑者靠近你，时机和地方都不缺；为使你不赞同，是我使你警醒。那么，承认他的恩宠吧，你也亏欠他，因为你没有犯罪。另一人亏欠我他所行的，你看见他被赦免了；你也亏欠我你所未行的。}因为一个人所犯的任何罪，若缺了那造人的引导者，另一个人也可能犯。\par

7. 那么，既然我已尽我所能在这短短时间内解决了这个深奥的难题（或者若我尚未解决，我仍如之前所说，是余下部分的债务人）；现在让我们简要思考一下赦罪的问题。基督曾被认为只是一个人，被邀请他的法利塞人和与他同席的客人视为如此。但那个有罪的妇人却在主身上看见了更多。因为她为何行这一切，岂不是为叫她的罪得赦吗？那么她知道他能赦罪；而他们知道没有人能赦罪。我们必须相信，他们所有人——即那些坐席的人，以及那靠近主双足的妇人——都知道没有人能赦罪。既然他们都知晓这一点，那相信他能赦罪的妇人，就认出他超过人。所以当他对妇人说：\Quote{你的罪赦了，}他们立刻说：\Quote{这人是谁，竟然也赦罪？}这人是谁，那有罪的妇人早已认识？你坐席如同健康人，却不认识你的医生；因为或许因更重的热病，你甚至失去了理智。因那狂乱的患者欢笑时，健康的人为他哀哭。然而，你确实知道并坚守这个真理；是的，坚守它：没有人能赦罪。这个妇人相信她能由基督得赦，就相信基督不仅是人，也是天主。\Quote{是谁，}他们说，\Quote{这人是谁，竟然也赦罪？}主并没有在他们问\Quote{这是谁？}时告诉他们：\Quote{他是天主圣子，天主的圣言；}他没有这样告诉他们，而是容许他们暂时仍停留在先前的意见中，却确实解决了令他们激动的问题。因为他看见他们坐席时，听见了他们的心思，就转身对妇人说：\Quote{你的信德救了你。}让这些说\Quote{这人是谁，竟然也赦罪？}、认为我只是人的人，仍认为我只是人吧。因为你，\Quote{你的信德救了你。}\par

8. 这位良医不但治愈了当时在场的病人，也为日后的人作了预备。日后将有人要说：\Quote{是我赦罪，是我使人成义，是我圣化，是我治好我所施洗的人。}这些人便是说\Quote{不要摸我}的那等人。是的，他们如此彻底地属于这一党，以致近日在我们的会议中——正如你们可在会议记录中所见——当专员为他们提供席位，请他们与我们同坐时，他们竟以为理应回答：\Quote{圣经告诉我们不可与这样的人同坐。}以免通过座位的接触，我们的传染（如他们所想）会传到他们那里。看，这岂不正是\Quote{不要摸我，因为我是洁净的}？但在另一日，我有更好的机会，当有人就教会提出疑问，说其中的恶人不会污损善人时，我便向他们指出这最可悲的虚荣。我回答他们，因为他们因此不愿与我们同坐，并说他们曾受天主圣经如此劝告，因为经上确实写着：\BibleQuote{我未曾坐在虚妄者的集会中；}我说：\Quote{你们若因经上写着\BibleQuote{我未曾坐在虚妄者的集会中}而不愿与我们同坐，那么，你们为何与我们同进这地方呢？因为紧接着经上又写着\BibleQuote{我决不与作恶的人同席}？}因此，他们所说的\Quote{不要摸我，因为我是洁净的}，正如那个邀请主的法利塞人，他以为主不认识那妇人，只因主没有阻止她触摸自己的脚。但就另一方面而言，那法利塞人更好，因为他虽然以为基督不过是人，却不相信罪能借由人得以赦免。由此可见，犹太人比异端者更有见识。犹太人怎么说？\Quote{\BibleQuote{这个赦罪的人是谁？}岂有人敢将此事攫为己有？}异端者又怎么说？\Quote{是我赦免，是我洁净，是我圣化。}不要由我，而由基督回答他：\Quote{人啊，当犹太人以为我不过是人时，我便将罪的赦免赐予信德。（不是我在回答你，而是基督在回答你。）}而你，异端者，你不过是人，却说：\Quote{来吧，妇人，我要使你痊愈。}然而，当犹太人以为我不过是人时，我却说：\BibleQuote{去吧，妇人，你的信德救了你。}\par

9. 他们回答说：\Quote{不知道}，正如宗徒所说：\BibleQuote{不明白自己所说和所主张的}。他们接着说：\Quote{如果人不能赦罪，那么基督所说的\BibleQuote{凡你们在地上所释放的，在天上也要被释放}就是假的。}但你们不知道这句话为何说出，以及按何种意义说出。主将要赐予人圣神，祂愿意让人明白，罪过是借着祂的圣神赐予祂的忠实信徒，而非凭人的功绩。因为你是什么，人啊？不过是一个需要医治的病人。你愿做我的医生吗？与我一同寻求那医生。为使主更清楚地表明这一点，即罪过是由圣神所赦免——祂已将圣神赐予祂的信徒，而非因人的功绩——在祂从死者中复活后，祂在某个地方说：\BibleQuote{你们领受圣神吧；}当祂说了\BibleQuote{你们领受圣神吧}之后，立即接着说：\BibleQuote{你们赦免谁的罪，就给谁赦免；}也就是说，是圣神赦免，而非你们。圣神就是天主。因此是天主赦免，而非你们。但你们之于圣神，又是什么？\BibleQuote{你们不知道你们是天主的宫殿，天主圣神住在你们内吗？}又说：\BibleQuote{你们不知道你们的身体是圣神的宫殿吗？这圣神在你们内，你们是从天主得来的。}所以，天主住在祂的圣殿内，即住在祂神圣的信徒内、在祂的教会内；借着他们，祂赦免罪过，因为他们乃是活的圣殿。\par

10. 但是，那位借着人赦罪的，也可以不借人而赦罪。因为能借他人给予的，同样也能亲自给予。祂曾借若翰的职务赐予某些人。至于若翰本人，祂又是借谁赐予的呢？有充分的理由：天主愿意显明这一点，并为此真理作证，当撒玛黎雅的一些人听了福音、受了洗，而这洗礼是由斐理伯执事——最初选出的七位执事之一——施行时，他们虽然受了洗，却没有领受圣神。消息传到耶路撒冷的宗徒们那里，他们便来到撒玛黎雅，为使那些受洗的人能借覆手领受圣神。果然如此：\BibleQuote{宗徒们就给他们覆手，他们就领受了圣神。}因为当时圣神是那样赐下的，甚至祂明确地显示出祂已被赐予。因为领受祂的人都说起万国的方言，以预表教会在万民中要说万国的方言。于是他们领受了圣神，祂也显然临在他们内。西满看见这事，以为这能力是出于人，便希望自己也能拥有。既然他以为那能力是出于人，便想从人那里购买。\Quote{你们给我多少钱，使我覆手也能赐予圣神？}于是伯多禄诅咒他说：\BibleQuote{你在这种信德上无分无关，因为你以为天主的恩赐可以用金钱买得。你的银子与你一同灭亡吧！}以及他在同一处所说的其他恰当的话。\par

11. 现在，我为何愿意向你们提出这个主题，请留心听，亲爱的诸位。天主首先应显示祂藉人的服务而工作，然后亲身工作，免得人们像西满所想的那样，以为那是人的恩赐，而非天主的恩赐。宗徒们自己早已深知此事。因为有一百二十人聚集在一起，那时没有覆手，圣神就降临在他们身上。当时谁曾给他们覆手呢？然而祂降临了，并首先充满了他们。在西满那件过犯之后，天主做了什么？请看祂藉事实而非言语教导。那位曾给那些人施洗的斐理伯，除非宗徒们聚集并覆手，圣神尚未降临在他们身上，但斐理伯给太监，即甘达刻女王的司库施了洗——他曾在耶路撒冷朝拜，归途中坐在车上读\BibleBook{依撒意亚}{先知书}，却不明白。斐理伯受指引，走近他的车，解释圣经，阐明信德，宣讲基督。太监信了基督，当他们来到一处水边时说：\BibleQuote{看，水！谁阻止我受洗呢？斐理伯对他说：你信耶稣基督吗？他回答说：我信耶稣基督是天主子。他即刻同他下到水里。}当圣洗的奥迹和圣事完成后，为使圣神的恩赐不被认为是来自人，便没有像先前那样等待宗徒到来，而是圣神立即降临。如此，西满的想法被摧毁，免得这种想法有人追随。\par

12. 再者，另一个更奇妙的例子。伯多禄来到百夫长科尔乃略那里，一个未受割损的外邦人，他开始向他及与他在一起的人宣讲耶稣基督。\BibleQuote{当伯多禄还在说话时，}我不是说他还未覆手，而是他甚至尚未给他们施洗，与伯多禄同来的人还怀疑未受割损的人是否应当受洗（因为在信主的犹太人和从外邦人带来信德的人之间，即犹太人与虽未受割损却已受洗的基督徒之间，产生了争端），为消除这个疑问，\BibleQuote{当伯多禄正在说话时，圣神降临了，}充满了科尔乃略，充满了与他在一起的人。藉着如此大事的印证，仿佛有一个响亮的声音对伯多禄说：\Quote{你为什么疑惑水呢？我已经在这里了。}\par

13. 因此，每一个灵魂将藉主的恩宠从诸多邪恶中被解救出来，在教会中如同从污秽的淫乱中被洁净，当以全备的信德相信，走近主的脚前，寻求祂的足迹，用眼泪倾注其上，并以头发擦拭。主的脚是福音的宣讲者。女人的头发是全部多余的财产。让她用头发擦脚，是的，务必擦，让她行哀矜的工作；擦完后，让她亲吻它们，领受平安，好拥有爱德。她曾来到某人那里，由像宗徒保禄这样的人施洗：让她从保禄那里听到：\BibleQuote{你们该效法我，如我效法了基督一样。}但如果她被另一个人施洗，一个\BibleQuote{寻求自己的事，不寻求耶稣基督的事}的人：让她从主那里听到：\BibleQuote{凡他们吩咐你们的，你们要遵守；但不要效法他们的行为。}因此，愿她的信德在于祂，无论她遇到好的福音宣讲者，还是言行不一的人。因为她从主那里以坚定的信德听到：\BibleQuote{女人，你的信德救了妳，平安地回去吧！}\par

\SourceRecord{Translated by R.G. MacMullen.；Nicene and Post-Nicene Fathers, First Series；Vol. 6.；Edited by Philip Schaff.；Buffalo, NY: Christian Literature Publishing Co.,；1888.；<http://www.newadvent.org/fathers/160349.htm>.}

% structure: p0001 source=h1 target=section reason=page-title

\section{论新约 讲道词 第五十篇}

\EditorialNote{C. Ben.}\par

论福音中的话，\BibleRef{路 9:57}等处，论及那三个人：一人说\BibleQuote{我要跟随你，无论你往哪里去}，却未被接纳；另一人不敢自荐，却被激励；第三人想要拖延，而被责备。\par

1. 请你们侧耳聆听主所赐予我，关于福音经课所讲的话。因为我们读到，主耶稣行事不同：一人自荐跟随祂，却未被接纳；另一人不敢自荐，却被激励；第三人拖延，而被责备。因为\BibleQuote{主，我要跟随你，无论你往哪里去}这话，有何如此迅速、如此活跃、如此准备好、如此适宜于这样伟大的善事，如同这\BibleQuote{跟随主，无论祂往哪里去}？你对此感到惊奇，说：\Quote{这是怎么回事？如此准备好的人，竟未得到良善的师傅和主耶稣基督的垂青，尽管祂正在邀请门徒，要将天国赐给他们？}但因祂是这样一位能预见未来的师傅，我们明白，弟兄们，这人若跟随了基督，必定会\Quote{寻求自己的事，而非耶稣基督的事}。因为祂亲自说过：\BibleQuote{凡称呼我\InnerQuote{主啊，主啊}的，不能都进天国。}这样的人正是此人，他并不像医生那样认识自己。因为若他此刻看见自己是伪善者，若他知道自己此时充满双面与诡诈，那么他就不知道他在与谁说话。因为祂正是福音作者所说的：\BibleQuote{祂不需要任何人替祂作证关于人，因为祂自己知道人里面有什么。}那么祂如何回答呢？\BibleQuote{狐狸有洞，天空的飞鸟有窝，但人子却没有枕头的地方。}但祂哪里没有呢？在你的信德中。因为在你心中，狐狸有洞：你充满诡诈；在你心中，天空的飞鸟有窝：你骄傲自大。你既充满诡诈与自傲，就不能跟随我。诡诈之人怎能跟随真诚呢？\par

2. 然后立刻对另一个沉默、一言不发、毫无承诺的人，祂说：\BibleQuote{跟随我！}正如祂在另一个人身上看到许多邪恶，同样祂在这人身上看到许多良善。\BibleQuote{跟随我}，你对一个不想要它的人说。看，这里有一个十分准备好，\BibleQuote{我要跟随你，无论你往哪里去}；然而你对另一个没有这种愿望的人说：\BibleQuote{跟随我}。\Quote{第一个人，}祂说，\Quote{我拒绝，因为我看到他里面有洞，有窝。}\Quote{但为何你催促这另一个人，你挑战他跟随你，而他却找借口？看，你甚至强迫他，而他却不来；你劝勉他，他却不愿跟随。因为他说道：\InnerQuote{我先去埋葬我的父亲。}}他内心的信德在主面前显露出来；但他的孝爱之情使他拖延。但主基督在预备人接受福音时，不容许这种肉体的、暂时的情感作为借口。的确，天主的法律也规定了这些职责，主自己也责备犹太人，因为他们破坏了这条天主的诫命。而宗徒保禄在他的书信中订立并说：\BibleQuote{这是第一条带着应许的诫命。}什么？\BibleQuote{应孝敬你的父亲和你的母亲。}天主确实说过这话。所以这个年轻人愿意服从天主，埋葬父亲；但这里是场合、时间和情况，在此场合、时间和情况必须让路。父亲应当孝敬，但天主必须服从。生养我们的应当爱慕，但创造我们的必须优先。\Quote{我正在召叫你，}祂说，\Quote{到我的福音中；我需要你从事另一项工作：这比你想要做的事情更伟大。\BibleQuote{让死人埋葬他们的死人。}你的父亲死了：有别的死人去埋葬死人。}那些埋葬死人的死人是谁？死人能由死人埋葬吗？如果他们死了，如何停放？如果他们死了，如何抬？如果他们死了，如何哀悼？然而他们确实停放、抬举并哀悼，而他们是死的，因为他们是不信者。\BibleBook{}{雅歌}中所写的教训我们，当教会说：\BibleQuote{请在我内安排爱情。}什么是\BibleQuote{请在我内安排爱情}？建立适当的层次，给各人其所应得的。不要把应该在前的事放在后面。要爱你的父母，但要将天主置于他们之上。看玛加伯的母亲：\Quote{我的儿子们，我不知道你们怎样在我腹中出现。我能怀你们，我能生你们，但我不能塑造你们；所以你们要听祂，将祂置于我之上；不要因我必须留在此地没有你们而困扰。}这样她命令他们，他们听从了她。这位母亲教导她孩子的话，正是主耶稣基督对那祂说\BibleQuote{跟随我}的人所教导的。\par

3. 请看，现在又有一个门徒来了，没有人对他说什么，他却说：\BibleQuote{主，我要跟从祢，但请允许我先去辞别我家里的人。} 我猜想他的意思是：\Quote{让我告诉我的朋友们，免得他们像往常一样寻找我。} 主却说：\BibleQuote{手扶着犁向后看的人，不配进天主的国。} 东方呼唤你，你却向着西方。在这篇训诲中我们学到，主拣选祂所愿意的人。但祂的拣选，正如宗徒所说，是依照祂的恩宠和他们的义德。因为宗徒的话正是这样；他说：\BibleQuote{你注意听厄里亚所说的话：\InnerQuote{主啊，他们杀了你的先知，拆毁了你的祭台，只剩下我一个人，他们还要寻索我的命。}但天主的答复对他说：\InnerQuote{我为自己保留了七千人，是未曾向巴耳屈膝的。}} 你以为只有你是忠信事奉的仆人吗？还有别的人也敬畏我，而且为数不少。因为我在那里有 \BibleQuote{七千人}。然后他又说：\BibleQuote{如今也是这样。} 因为有些犹太人相信了，虽然多数的被弃绝；就像那心里藏着狐狸洞的人。\BibleQuote{如今也是这样，}他说，\BibleQuote{按恩宠的拣选，还有残余的得救。} 这就是说，同一基督现在仍然如此，祂当时也对那厄里亚说：\BibleQuote{我为自己保留了。} 什么叫做\Quote{我为自己保留了}？我拣选了他们，因为我看见他们的心信赖我，而不信赖自己，也不信赖巴耳。他们没有改变，他们仍是我所造的样子。至于你这说话的人，若不是你信赖我，你会在哪里呢？若不是你充满我的恩宠，难道你不也会向巴耳屈膝吗？但你充满我的恩宠，因为你一点也不依靠自己的力量，而是完全依靠我的恩宠。因此，不要因此自夸，以为在你的事奉中没有同仆；还有别的人是我拣选的，正如我拣选了你，就是那些信赖我的人；正如宗徒所说：\BibleQuote{如今也是这样，按恩宠的拣选，还有残余的得救。}\par

4. 基督徒，要谨慎，当心骄傲。因为你虽追随圣人，但总要完全归功于恩宠；因为在你身上能有任何\Quote{余民}，是天主的恩宠所致，而非你自己的功德。因为依撒意亚先知早已着眼于这余民，曾说：\BibleQuote{若不是万军的上主给我们留下生还的苗裔，我们早已像索多玛一样，也早已像哈摩辣一样了。} \BibleQuote{所以，}宗徒说，\BibleQuote{如今也是这样，按恩宠的拣选，还有残余的得救。既是出于恩宠，}他说，\BibleQuote{就不是出于行为了}（就是说，\InnerQuote{不要再仗恃自己的功德了}）；\BibleQuote{否则，恩宠就不成其为恩宠了。} 因为你若建立在自己的行为上，那么赏报就归于你，而非白白赐予的恩宠。但若是恩宠，就是白白赐予的。我问你，罪人啊，\BibleQuote{你信基督吗？}你说：\BibleQuote{我信。} \BibleQuote{你信什么？信你所有的罪能因祂白白地获得赦免？}那么，你就得着你所信的。哦，白白赐予的恩宠！而你，义人啊，你信什么？你信你无法靠天主而持守你的义德吗？那么，你之所以为义，应完全归于祂的慈悲；而你是罪人，则归于你自己的不义。做你自己的控告者，祂必成为你恩慈的拯救者。因为每件罪行、邪恶和罪过都出自我们的疏忽，而一切美德和圣洁都来自天主恩慈的良善。让我们归向主。\par

\SourceRecord{Translated by R.G. MacMullen.；Nicene and Post-Nicene Fathers, First Series；Vol. 6.；Edited by Philip Schaff.；Buffalo, NY: Christian Literature Publishing Co.,；1888.；<http://www.newadvent.org/fathers/160350.htm>.}

% structure: p0001 source=h1 target=section reason=page-title

\section{讲道集第五十一篇：论新约}

\EditorialNote{[CI. Ben.]}\par

\Emph{论福音的话\BibleRef{路 10:2}，\BibleQuote{庄稼确实很多，}等等。}\par

1. 刚才读过的福音训诲提醒我们，要探究主所说的庄稼是什么：\BibleQuote{庄稼确实很多，但工人很少。你们应当求庄稼的主人，派遣工人来收割他的庄稼。}然后，主又拣选了另外七十二人，连同十二门徒（祂也称他们为宗徒），并按照祂的话，似乎都派到那已成熟的庄稼里去。那么，这庄稼是什么呢？因为在那外邦人中，并没有撒下任何种子，这庄稼不可能在他们中间。因此，我们只能明白，这庄稼是在犹太人中。庄稼的主人来到那庄稼中，派收割的人去；但对于外邦人，祂派去的不是收割的人，而是撒种的人。因此，我们应当明白：在犹太人中是收割，在外邦人中是撒种。因为宗徒们是从那庄稼中被选出来的；在那里，庄稼已经成熟，因为先知们已经撒种了。品味天主的耕作，享受祂的恩赐和祂田里的工人，是多么令人愉快的事！因为这耕作中，那说\BibleQuote{我比他们众人更劳苦}的人确实劳动了。但劳动的力量是庄稼之主赐给他的。因此他补充说：\BibleQuote{然而不是我，而是天主的恩宠与我同在。}因为他从事这耕作，他清楚地表明：\BibleQuote{我栽种了，\Person{阿波罗}浇灌了。}但这宗徒，从\Person{扫罗}成为\Person{保禄}，即从傲慢成为最微小的（因为扫罗一名源于\Person{撒乌耳}，而\Person{保禄}是小的；因此，他在解释自己的名字时，说：\BibleQuote{我虽是宗徒中最小的}）。这位\Person{保禄}——小小的、最小的——奉派到外邦人那里，他说他特别被派到外邦人中。他自己这样写，我们读到、相信并宣讲。他在\WorkTitle{致迦拉达人书}中说，被主耶稣召叫后，他来到\Place{耶路撒冷}，\BibleQuote{把福音告诉了宗徒们}，他们向他伸出右手，这是和谐的记号、同意的记号，表明他们从他那里听到的与自己所传的没有差别。后来他说，他和他们之间达成协议：他去外邦人那里，他们去受割礼的人那里；他做撒种的人，他们做收割的人。所以，\Place{雅典}人虽然不知情，却给他起了恰当的名字。因为他们听到他的话，就说：\Quote{这个撒种言论的是谁？}\par

2. 请留意，并和我一起欣赏天主的耕作及其中的两次收割：一次已经过去，一次尚未来到；一次已在犹太人中过去，一次将在外邦人中进行。我们如何证明这一点？唯有通过庄稼之主——天主的圣经。看，在今天的课文中，主说：\BibleQuote{庄稼很多，但工人很少。你们应当求庄稼的主人，派遣工人来收割他的庄稼。}但因为在那庄稼中，将有反对和逼迫的犹太人，主说：\BibleQuote{看，我派遣你们好像羊进入狼群。}让我们从\BibleBook{若望}{福音}中更清楚地说明这收割的事：主疲倦地坐在井旁，确实成就了伟大的奥秘，但时间太短，无法全部讲述。但请听与当前主题相关的内容。因为我们承诺要展示在先知们宣讲的民族中的收割；正因先知们是撒种的人，所以宗徒们才能是收割的人。一位\Place{撒玛黎雅}妇人与主耶稣交谈，当主告诉她应当怎样朝拜天主时，她说：\BibleQuote{我们知道默西亚——就是基督——要来；祂必告诉我们一切。}主对她说：\BibleQuote{和你说话的，就是祂。}信你所听到的；何必追究所见到的呢？\BibleQuote{和你说话的，就是祂。}至于她所说的\BibleQuote{我们知道默西亚要来}，即\Person{摩西}和先知们所宣布的基督，庄稼已经成穗了。当它还生长时，有先知们做撒种的人；现在它成熟了，等待宗徒们做收割的人。她一听这话就信了，丢下水罐，急忙跑去，开始宣传主。那时门徒们去买食物，回来时见主与妇人谈话，感到惊奇。但他们不敢对祂说：\Quote{你和她说什么？}他们心中诧异，却抑制了内心的鲁莽。对这位\Place{撒玛黎雅}妇人来说，基督的名字并不新鲜；她已在等待祂的来临，相信祂会来。她从哪里相信的呢？若不是\Person{摩西}撒了种？但请听更明确的记载。主对门徒们说：\BibleQuote{你们说夏天还远，但举目看看，田地已经发白，可以收割了。}然后祂又说：\BibleQuote{别人劳苦，你们享受他们的劳苦。}\Person{亚巴辣罕}、\Person{依撒格}、\Person{雅各伯}、\Person{摩西}、先知们都劳苦撒种；主来临时，庄稼已成熟。带着福音镰刀的收割者被派去，将麦捆运到主的禾场上，那里\Person{斯德望}将要被碾打。\par

3. 但这里出现了\Person{保禄}，他被派往外邦人那里。他在陈述自己特别领受的恩宠时并未隐瞒此事。因为他在圣经中说，他被派去传福音，到基督的名尚未被传扬的地方。但由于第一次收割已经过去，所有剩下的犹太人已不是收割的对象，让我们考虑我们自身的那次收割。因为它是宗徒和先知们播下的。主亲自播种了它。因为他在宗徒们内，既然基督自己也收割了它。因为他们离了他什么都不能做；他没有他们却仍是完美的。因为他亲自对他们说：\BibleQuote{因为离了我，你们什么也不能做。}那么基督从今以后在外邦人中播种说什么？\BibleQuote{有一个撒种的出去撒种。}\Quote{那里}有收割者\Quote{被派出去}收割，\Quote{这里}一个不知疲倦的撒种者\Quote{出去}撒种。那些\BibleQuote{有的种子落在路旁，有的落在石头地上，有的落在荆棘中}使他害怕了吗？如果他害怕这些难以对付的土地，他就永远不会到达好地。这与我们何干？我们有什么必要现在争论犹太人，谈论糠秕？这唯一关系我们的是，我们不要成为\Quote{路旁}，也不要成为\Quote{石头地}，也不要成为\Quote{荆棘}，而要成为\Quote{好地}。让我们的心预备好，从中可以产生\BibleQuote{三十倍}，或\BibleQuote{六十倍}，或千倍，以及\BibleQuote{一百倍}；有的多一些，有的少一些；但都是麦子。不要让它成为\Quote{路旁}，在那里敌人像鸟一样可以夺走被路人践踏的种子。不要让它成为\Quote{石头地}，那里浅土使它立刻发芽，以致不能忍受太阳。不要让它成为\Quote{荆棘}，即今世的贪欲，无秩序生活的忧虑。因为还有什么比那种生活的忧虑更糟，它不让人到达生命？还有什么比因关心生命而失去生命更悲惨？还有什么比因怕死亡而堕入死亡更不幸？让荆棘被连根拔起，田地预备好，种子播下，让它们长到收割，盼望谷仓，而不是害怕火。\par

4. 因此，我的职分是，主虽不配，却指派我在他的田里作工人，向你们讲这些事，去播种、栽植、浇水，甚至给一些树松土，施用粪筐；我忠实履行这些事，你们忠实领受，主帮助我劳作，帮助你们相信，我们大家一同劳作，但在他内战胜世界。那么，我已说了你们该做的事；现在我想说我们该做的事。但也许你们有些人觉得，我声明想说的话有些多余，心里自言自语：\Quote{巴不得他现在放我们走！他已经说了我们该做的事，至于他该做的事，与我们何干？}我认为，在彼此互爱中，我们属于你们更好。你们如今确实是一家人，我们同一家族的管家，但我们都属于同一位主。我所给的，不是我自己的，而是从他那里，我也从他领受。因为如果我给的是自己的，我将给谎言。\BibleQuote{因为说谎者出于自己说自己的话。}所以你们应当听从属于管家职责的事，或者你们若发现自己如此，就会喜乐，或者即使在这事上你们也会得到教导。因为在这民众中，有多少人将来会成为管家！我也曾像你们现在这样；几年之前，我还在低处与同伴一同领受食物，现在却似乎在高处量给同伴们食物。我现在是主教对平信徒说话；但我知道，我在对他们说话时，也是对将来也会成为主教的许多人说话。\par

5. 让我们看看我们当如何理解主对派去传福音的人所命令的事，并在心里思考这片预备好的收获。他说：\BibleQuote{你们不要带钱囊，不要带行囊，不要带鞋，在路上也不要向人问安。无论进了哪一家，先说：愿这家平安。那里如有和平之子，你们的平安就必停留在他身上；否则，仍归你们。}如果它\Quote{停留了}，别人就失去了它吗？愿这远离圣人们的心意！所以，这不应按肉性的意义来理解；因此，可能钱囊、鞋、行囊也都不是按字面；尤其是那一条，如果我们不加审查地简单接受，似乎命令我们骄傲，即我们\Quote{在路上也不要向人问安}。\par

6. 让我们留意我们的主，我们的真实榜样和援助。让我们证明他是我们的援助；\BibleQuote{因为离了我，你们什么也不能做。}让我们证明他是我们的榜样；\Person{伯多禄}说：\BibleQuote{基督为你们受了苦，给你们留下了榜样，叫你们跟随他的足迹。}主自己在路上带着钱袋，并将这些钱袋托付给了\Person{犹达斯}。诚然，他遭受了盗贼的侵害；但我渴望向主学习，说：\Quote{主啊，你遭受了盗贼的侵害，你从哪里有他能拿走的东西？你告诫我这可怜而虚弱的人，甚至不要带钱囊；你却带着钱袋，有他可以拿走的东西。如果你没有带，他也找不到任何可拿的。}然后他对我说：\Quote{理解你听到的\InnerQuote{不要带钱囊}的意思是什么？钱囊是什么？是封闭的钱财，即隐藏的智慧。\InnerQuote{不要带钱囊}是什么意思？不要只在你们自己内成为智慧。要领受圣神。它应是你内的泉源，而不是钱囊；从那里将分施给他人，而不是封闭在内。}而行囊与钱囊相同。\par

7. 什么是\Quote{鞋子}？我们所用的鞋子，是死兽的皮，是脚的覆盖物。如此，我们受命弃绝死行。这是\Person{梅瑟}在预像中所受的训示，当主对他说话时，说：\BibleQuote{将你脚上的鞋脱下，因为你所站的地方是圣地。} 有什么地比天主的教会更圣呢？因此，让我们站在其中，脱下我们的鞋子，也就是说，弃绝死行。因为关于这些我们行走的鞋子，我主再次向我保证。因为若他本人未曾穿鞋，\Person{若翰}就不会论他说：\BibleQuote{我当不起解他的鞋带。} 那么，要有服从，不要让倨傲的严厉潜入我们中间。有人说：\Quote{我赤脚行走，因此我满全了福音。} 好吧，你能做到，我不能。但让我们俩人都持守我们共同领受的。如何？让我们以爱德燃烧，让我们彼此相爱；如此，我将爱你的刚强，而你将担待我的软弱。\par

8. 然而，你们这些不愿意理解这些话语用意的，并且因你歪曲的解释而被迫甚至诋毁主关于\Quote{钱袋}和\Quote{鞋子}的教导，你们以为如何？那么，当我们途中遇见朋友时，是否应该既不致意于高于我们的人，也不回应低于我们的人的问候？什么，你因被人问候而沉默，就以为满全了福音？但这样你就不像行路的旅人，反而像指路的路标。那么，让我们摒弃这粗俗的解释，正确理解主的话：\BibleQuote{在路上不要问候任何人。} 因为这命令并非无缘无故，他也不会不喜欢我们做他所命令的。那么，\BibleQuote{在路上不要问候任何人}是什么意思？其实可以简单理解为，他命令我们尽快完成他所吩咐的；他的话\BibleQuote{在路上不要问候任何人}就如同他说：\Quote{将其他一切搁置，直到你完成所命令的。} 这是照谈话中常以夸张手法表达的习惯。无需远求；就在同一次讲道中稍后他说：\BibleQuote{葛法翁，你被高举到天上，你将被推下地狱。} \Quote{被高举到天上}是什么意思？那城的城墙碰触云彩，或达到星辰了吗？但\Quote{被高举到天上}是什么意思？你自以为极其幸福、极其强大，你极其骄傲。那么，为了夸张，才向那座并未被高举、并未升到天上的城说：\Quote{你被高举到天上}；同样，为了表达匆忙，才夸张地说：\Quote{如此奔跑，如此执行我所吩咐的，使途中行人丝毫不能拖延你；而是不顾其他一切，奔向摆在你面前的终点。}\par

9. 但这些话还有另一个更深奥的含义，并不难理解，尤其关于我和所有分施者，以及你们这些听道者。问候者，祝愿救恩。因为古人写信时这样写：\Quote{某人向某人致意。} 问候一词源自救恩。那么，\Quote{在路上不要问候任何人}是什么意思？那些\Quote{在路上问候}的人，是\Quote{出于机会}而问候。我看出你们已迅速理解了我，但尽管如此，我还不能结束。因为并非所有的人都理解得那么快。我看到有些人以声音表示理解，我更看到更多人因沉默而要求更多。但既然我们谈论道路，让我们如同在道上行走：快捷的人，等候慢行的人，步调一致。那么我说了什么？\Quote{在路上问候}的人，只是出于机会问候？他并非要走向他所问候的人。他正做一事，另一事拦在他的路上；他寻求一事，却发现某事横在途中要做。那么，\Quote{出于机会问候}是什么意思？\Quote{出于机会}宣布救恩。那么，宣布救恩不是宣讲福音又是什么？因此，若你宣讲，当出于爱德，而非\Quote{出于机会}。于是，有人虽\Quote{寻求自己的事}，但他们并不宣讲别的福音；关于这些人，宗徒叹息说：\BibleQuote{因为众人寻求自己的事，而非耶稣基督的事。} 这些人\Quote{问候}，即宣布救恩，他们宣讲福音；但他们寻求其他事物，因此只是\Quote{出于机会}问候。这是什么意思？若你是这样的人，无论你是谁，你就在做这事；不，并非所有做这事的人都是如此，但可能有些做这事的人是这样的。但若你是这样的人，那不是你做的，而是借着你做的。\par

10. 因为宗徒曾遭如此苦难，但他并没有命令他们如此。这些人做些事，或有些事借着他们做成；他们追求别的东西，却仍宣讲圣言。不要在意讲道者追求什么；你当决心持守他所宣讲的；但他的意图与你无关。从他的口中听取救恩之言，从他的口中持守这救恩。不要判断他的心。若你见他追求其他事物，那与你有何相干？听取那救恩者所说的：\BibleQuote{他们所说的，你们要遵行。}\BibleRef{玛 23:3}那说过\BibleQuote{他们所说的，你们要遵行。}\BibleRef{玛 23:3}的，已使你确知。他们行恶吗？\BibleQuote{不要效法他们的行为。}\BibleRef{玛 23:3}他们行善。他们不\BibleQuote{在路上问候人，}\BibleRef{路 10:4}他们不假意传福音；\BibleQuote{你们要效法我，如同我效法基督一样。}\BibleRef{格前 11:1}好人向你宣讲，你便从葡萄树上摘取葡萄。坏人向你宣讲，你便摘取那挂在篱笆上的葡萄。葡萄串生长在荆棘缠绕的葡萄枝上，却并非从荆棘所长。无论如何，当你见到此类情景且感到饥饿时，你采摘时须谨慎，免得伸手摘葡萄时被荆棘刺伤。我所说的是：你们要以这样的方式听善言，即不要效法那人的恶行。让他\Quote{假意}宣讲，在路上问候人；这对他有害，因为他没有听从基督的诫命：\BibleQuote{不要在路上问候任何人；}\BibleRef{路 10:4}但这不会伤害你，无论你是从路人还是从直接来找你的人那里听到救恩，你都能持守这救恩。请听宗徒，如我所说已使我们明白这事：\BibleQuote{那又怎样呢？无论如何，或是假意，或是真诚，基督总被传扬；为此我欢喜，而且还要欢喜。因为我知道，这事借着你们的祈祷，终必使我得救。}\BibleRef{斐 1:18-19}\par

11. 那么，让这样的基督的宗徒、福音的宣讲者——他们\BibleQuote{不在路上问候人，}\BibleRef{路 10:4}即他们不寻求或做任何其他事，而是以真诚的爱德宣讲福音——让他们进入那房屋，并说\BibleQuote{愿平安归于这家。}\BibleRef{路 10:5}他们不仅用口说话；他们倾注他们所充满的；他们宣讲平安，他们拥有平安。他们不像那些关于他们曾说过\BibleQuote{平安，平安，但并无平安}\BibleRef{耶 6:14}的人。什么是\BibleQuote{平安，平安，但并无平安}\BibleRef{耶 6:14}？他们宣讲平安，却没有平安；他们赞美平安，却不爱平安；他们说，却不做。但你们仍要接受这平安，\BibleQuote{无论是假意，或是真诚，基督总被传扬。}\BibleRef{斐 1:18}那么，谁充满平安，并问候说\BibleQuote{愿平安归于这家，若这家有平安之子，他的平安必降临在他身上；不然，这平安必归于你们。}\BibleRef{路 10:6}因为或许那里没有一个平安的人，但问候者并没有损失什么，\BibleQuote{它必回到你们身上。}\BibleRef{路 10:6}他说，\BibleQuote{回到你们身上。}\BibleRef{路 10:6}它必回到你们身上，虽然它从未离开过你们。因为祂要说的意思是：你们宣告了它，这于你们有益，那没有接受的人却毫无获益；你们没有失去你们的赏报，因为他仍然空虚；这赏报将因你们的善意而给予你们，因你们所施予的爱德而给予你们，那给你们保证的祂——藉着那天使的声音说：\BibleQuote{主爱的人在世享平安。}\BibleRef{路 2:14}——必将它赐给你们。\par

\SourceRecord{Translated by R.G. MacMullen.；Nicene and Post-Nicene Fathers, First Series；Vol. 6.；Edited by Philip Schaff.；Buffalo, NY: Christian Literature Publishing Co.,；1888.；<http://www.newadvent.org/fathers/160351.htm>.}

% structure: p0001 source=h1 target=section reason=page-title

\section{新约讲道集 第五十二讲}

\Emph{[CII. Ben.]}\par

\Emph{关于福音的话，\BibleRef{路 10:16}：\BibleQuote{拒绝你们的，就是拒绝我。}}\par

1. 当时我们的主耶稣基督对祂门徒所说的话，被写了下来，预备给我们听。所以我们听到了祂的话。因为如果祂被看见却没有被听见，对我们有何益处？而如今祂未被看见，却仍被听见，这并无损害。于是祂说：\BibleQuote{轻视你们的，就是轻视我。}如果祂只对宗徒们说：\BibleQuote{轻视你们的，就是轻视我}，那么你们就轻视我们吧。但是，如果祂的话也传到了我们，祂召唤了我们，并安置我们在他们的位置，那么你要小心，不可轻视我们，免得你们对我们所做的不义触及祂。因为如果你们不敬畏我们，要敬畏那说\BibleQuote{轻视你们的，就是轻视我}的那一位。但我们这些不愿意被你们轻视的人，为何对你们说话呢？无非是希望因你们良好的品行而喜乐。愿你们的善工成为我们危险的安慰。活得好，以便死得安好。\par

2. 在我说的话中，\Quote{活得好，以便死得安好}，不要以为那些也许曾活得邪恶、却死在自己床上的人，其葬礼被炫耀，被放入昂贵的棺材，在装饰极为美丽和费心的坟墓中；也不因为你们中有人也许在说：\Quote{我希望这样死去}，就以为我所选择说的话是虚妄的；当我说我希望你们活得好，以便死得安好时，不要这样想。另一方面，也许有某人，他活得很好，但按人的看法却死得不好；也许死于房屋倒塌，死于海难，死于野兽；每个属血肉的人在心中说：\Quote{活得好有什么益处？看这人如此生活，却如此死了。} \Quote{回到你的心中；}如果你是信者，你将在那里找到基督；祂在那里对你说话。因为我大声呼喊，但祂沉默地教导更多。我以声音之词说话；祂以思想之恐惧在内里说话。愿祂将我的话种植在你心中；因为我承担了说\Quote{活得好，以便死得安好}的责任。看，因为信德在你们心中，基督住在那里，由祂来教导我所欲表达的事。\par

3. 记得福音中那个富人和那个穷人；\BibleQuote{富人穿着紫袍和细麻衣，天天奢华宴乐}；穷人\BibleQuote{躺在}富人家门口，饥饿，盼望从富人桌子上掉下来的\BibleQuote{碎屑}，浑身生疮，被狗\BibleQuote{舔}。记得，我说；你们从哪里记得，除非基督在你们心中？告诉我，你们在里面问了祂什么，祂又回答了什么。因为祂接着说：\BibleQuote{后来那穷人死了，被天使带到亚巴郎的怀里。那富人也死了，被埋葬在地狱里。他在痛苦中举目一望，看见拉匝禄在亚巴郎怀里安息。于是他喊叫说：亚巴郎父亲，可怜我吧，请打发拉匝禄用指尖蘸点水来凉凉我的舌头，因为我在这火焰中极其痛苦。}在世界上傲慢，在地狱里却是乞丐！因为那穷人确实得到了他的碎屑；但那人却得不到一滴水。那么，对这两个人，告诉我，谁死得好，谁死得不好？不要问眼睛，回到心中。因为如果你问眼睛，它们会欺骗你。因为那个富人在死时，可能有很多极其华丽、以许多世俗场面伪装的尊荣。可能有多少哭泣的奴仆和使女！何等气派的随从队伍！何等辉煌的葬礼！何等昂贵的墓葬！我想他一定是被香料覆盖。那么，弟兄们，我们该说他死得好，还是死得不好？如果你问眼睛，他死得非常之好；如果你询问你内在的老师，他死得极其不好。\par

4. 如果那些高傲的人，他们只为自己保留财物，不将任何分给穷人，就这样死去；那么那些掠夺他人财物的人又如何死去？因此我有真理由说：\Quote{活得好，以免死得不好}，以免你们像那个富人那样死去。没有比死后的时光更能证明恶死。另一方面，看那穷人；不要用眼睛，因为这样你会错误；让信德看他，让内心看他。把他放在你眼前，躺在地上，\BibleQuote{满身疮痍，狗来舔他的疮}。现在当你以这种外表回忆他时，你立即厌恶他，你扭过脸去，掩住鼻孔：然后用心灵的眼睛看吧。\BibleQuote{他死了，被天使带到亚巴郎的怀里。}那富人的家人可见地为他哀哭；天使却未被看见欢呼。那么亚巴郎如何回答富人？\BibleQuote{孩子，你应记得你生前享过福。}你认为没有什么好的，除了你在今生所拥有的。你已得到了它们；但那些日子过去了；你失去了全部；你被留下在地狱里受折磨。\par

5. 弟兄们，对你们说这些话正是时机。要尊重穷人，无论是躺在地上的，还是行走的；要尊重穷人，行善工。你们一向如此行的，要继续行；你们尚未如此行的，现在就行。让行善者的数目增加，因为信者的数目也在增加。你们尚看不见你们所作之善有多大；因为农夫播种时也看不见收成，但他信赖土地。你们为什么不信赖天主？我们的收成会来到。你们要想，我们现在正在劳苦中忙碌，正在劳苦中工作，但必定会得到，正如经上所记载：\BibleQuote{他们边走边哭，撒下种子；但他们必定会欢呼着前来，带着禾捆。}\par

\SourceRecord{Translated by R.G. MacMullen.；Nicene and Post-Nicene Fathers, First Series；Vol. 6.；Edited by Philip Schaff.；Buffalo, NY: Christian Literature Publishing Co.,；1888.；<http://www.newadvent.org/fathers/160352.htm>.}

% structure: p0001 source=h1 target=section reason=page-title

\section{论新约讲道集 第53篇}

\Emph{[CIII. Ben.]}\par

\Emph{\Quote{论福音之言}}，\BibleRef{路 10:38}\Emph{，\BibleQuote{有一个名叫玛尔大的女人接待他到自己家里}等。}\par

我们主耶稣基督刚才从福音中读到的这些话，让我们明白，在我们于今生忙碌的各种事务中，必须朝向某一件独一无二的事。我们正朝向它，如同仍在旅途中，而非安居之所；尚在途中，而非故乡；尚在渴望，而非享乐。然而，让我们朝向它，不可懈怠、不可间断，以便有朝一日能够抵达。\par

玛尔大和玛利亚是两姐妹，真正的亲属，不只在血统上，也在信仰上；两人都紧紧跟随主，两人都一心一意地服侍那在肉身中临在的主。玛尔大接待他，如同接待素不相识的客人。然而，是婢女接待她的主，是病人接待她的救主，是受造物接待她的造物主。她接待他是为了给肉体食物，而她自己则要在精神上得到喂养。因为主乐意\BibleQuote{取了奴仆的形体}，并且\BibleQuote{取了奴仆的形体}之后，由于他的俯就，而非他的本体，竟愿意由奴仆来喂养。这实在是俯就，竟允许自己被他人喂养。他有一个身体，在其中他确实会饥饿、会口渴；但你们不知道吗？当他在旷野饥饿时，有天使服侍他？因此，他乐意被喂养，乃是向喂养他的人显示恩惠。这有什么稀奇呢？因为他曾向那位寡妇显示过同样的恩惠，关于圣厄里亚，他先前曾藉乌鸦的服侍喂养过他？难道当他打发他到寡妇那里时，他丧失了喂养他的能力吗？绝不。当他打发他到寡妇那里时，他并没有丧失喂养他的能力，但他定意要藉着那位寡妇对他仆人的虔诚服务，来祝福这位虔诚的寡妇。因此，主被当作客人接待，\BibleQuote{他来到自己的地方，自己的人却不接待他；但是，凡接待他的，他赐给他们权能，成为天主的子女}：收纳奴仆，使他们成为弟兄；赎回俘虏，使他们成为同继承者。然而，你们中也许有人会说：\Quote{啊，那些蒙恩在自己家中接待基督的人是有福的！}不要悲伤，不要埋怨，你们生在不再看见主肉身的时代；他没有夺走你们的这种福气。\BibleQuote{因为你们对我这些最小兄弟中的一个所做的，就是对我做的。}\par

限于时间短暂，关于那位乐意在肉身中被喂养、同时在精神上喂养人的主，我就说了这几句话。现在让我们回到我原先提出的关于合一的话题。玛尔大正在安排并准备喂养主，忙于许多服侍。她的妹妹玛利亚却宁愿由主来喂养。她仿佛撇下了忙于许多服侍的姐姐，自己坐在主脚前，静静地听他的话。她那最忠信的耳朵已经听到：\BibleQuote{你们要安静，并要知道我是天主。}玛尔大心烦意乱，玛利亚却正在享受；一个在安排许多事，另一个却注目于唯一的事。两种服侍都是好的；但至于哪一个更好，我们该说什么呢？我们有一位可以询问，让我们一起倾听。哪一个更好，我们刚才在读经时听到了，让我重复时再听一次。玛尔大向她的客人申诉，向审判者提出她虔诚抱怨的请求，说她的姐姐撇下了她，在她忙碌服侍时没有帮助她。玛利亚没有回话，但主在她面前作出了判决。玛利亚宁愿安息，将她的案件交托给审判者，并不打算忙着回答。因为如果她要准备言辞来回答，她就得放松专心听讲。因此，主回答了，他并不缺乏言辞，因为他是圣言。那么他说了什么？\Quote{玛尔大，玛尔大，}请听：\Quote{你为许多事操心忙碌，但只有一件事是必要的；}因为 \OriginalTerm{必要的事}{unum opus est} 的意思并非\Quote{一件工作}，即单独一件工作，而是惟有一件事是必要的、有益的、不可或缺的，玛利亚选择了这一件事。\par

4. 请诸位弟兄深思这\Quote{一件事}，并看看在众多之中，若非\Quote{这一体}，还有什么能令人愉悦。请看，因天主的仁慈，你们的人数何其众多：倘若你们不专注于\Quote{一件事}，谁还能容忍你们呢？在这众多之中，这宁静又从何而来？若有一体，便是子民；若取去一体，便是一群乌合之众。因为一群乌合之众，若非无序的群众，又是什么呢？但请听宗徒所说：\BibleQuote{如今我恳求你们，弟兄们。}他是在对群众说话，却愿意他们全都成为\Quote{一体}。\BibleQuote{如今我恳求你们，弟兄们，你们大家要说话一致，在你们中间不要有分裂，但要同心同德，意念相同。}又在另一处说：\BibleQuote{你们要同心合意，思念同一件事，不存私心，不贪虚荣。}主也为属于祂的人向圣父祈祷：\BibleQuote{使他们合而为一，如同我们原为一体一样。}在宗徒大事录中也记载：\BibleQuote{众信徒都同心合意，没有一人说他的财物归自己所有。}因此，\BibleQuote{请你们与我一同赞美上主，让我们齐声颂扬祂的名。}因为\Quote{一件事}是必要的，那就是天上的合一体，圣父、圣子、圣神在其内合而为一。请看，合一之赞美是如何被推荐给我们的。无疑，我们的天主是圣三。圣父不是圣子，圣子不是圣父，圣神既不是圣父，也不是圣子，而是两者的神；然而这三位不是三位神，也不是三位全能者，而是一位全能的天主，整个圣三就是一位天主；因为\Quote{一件事}是必要的。唯有当我们众多的人拥有同一颗心时，才能达至这\Quote{一件事}。\par

5. 对穷人的服务是好的，尤其是对天主的圣人们所应尽的义务和宗教性的服务。因为这是偿还，而非馈赠，正如宗徒所说：\BibleQuote{我们若在属神的事上给你们播种，若从你们收割属血肉的事，还算大事吗？}这些是好的，我们劝你们去做，甚至藉着主的话来建立你们：\BibleQuote{不可忘记款待}圣徒。有时，那些未曾留意的人，因款待了不认识的人，竟款待了天使。这些事是好的，但\Person{玛利亚}所选择的那件事更好。因为前者因必要而带来诸多烦恼，后者却因爱德而带来甘甜。人在服务时，总希望遇到些什么，但有时却无能为力：缺乏的要去寻找，已有的要准备妥当，心灵便分散了。因为如果\Person{玛尔大}仅凭自己足以应付这些事，她就无需请求妹妹的帮助了。这些事多样而繁杂，因为是属血肉的，是暂时的；虽为善事，却是短暂的。但主对\Person{玛尔大}说了什么呢？\BibleQuote{\Person{玛利亚}选择了更好的一份。}不是你的一份不好，而是她的一份更好。请听，如何更好：\BibleQuote{这一份是不能从她手中夺去的。}总有一天，这些必要职务的重担会从你手中被夺去，而真理的甘甜却是永久的。\BibleQuote{她所选择的这一份，是不能从她手中夺去的。}这不会被夺去，反而会增加。也就是说，在今世会增加，在来世将达至圆满，永不会被夺去。\par

6. 是的，\Person{玛尔大}，在你良好的服务中是有福的，就连你（请容我这样说）也在寻求这劳苦的报酬——安息。如今你忙于许多服务，乐于喂养终有一死的身体，即使是圣人们的身体；但当你到达那国度时，你还会在那里找到任何陌生人，可以接待到家中吗？你还会找到饥饿的人，可以分饼给他吗？或口渴的人，可以递杯给他吗？病人，你可以去探望吗？好争讼的人，你可以使他们和好？死人，你可以埋葬吗？那里将没有这一切，但会有什么呢？\Person{玛利亚}所选择的，将在那里使我们饱足，而不再喂养他人。因此，\Person{玛利亚}在此所选择的，将在那里圆满而完全地实现；她从主的话语这丰盛的筵席上，拾取了些许碎屑。因为你们想知道那里将如何？主自己论及祂的仆人说：\BibleQuote{我实在告诉你们，祂要使他们坐席，自己束上带子前来伺候他们。}什么是\Quote{坐席}呢？岂不是\Quote{安息}吗？什么是\Quote{祂要亲自前来伺候他们}？首先，祂经过，然后伺候。在哪里呢？在那天上的宴席中，祂说：\BibleQuote{我实在告诉你们，将有许多人从东方和西方来，同\Person{亚巴郎}、\Person{依撒格}和\Person{雅各伯}在天国里一起坐席。}在那里，主要喂养我们，但祂先从这里经过。因为（你们应当知道）\Quote{逾越节}按解释就是\Quote{经过}。主来了，行了属天主的事，承受了属人的事。祂还在被吐唾沫吗？还被掌掴吗？还戴茨冠吗？还受鞭打吗？还被钉十字架吗？还被长矛刺透吗？祂已经过。同样，福音告诉我们，当祂与门徒们一起过逾越节时，福音说：\BibleQuote{当耶稣该从这个世界回到父那里去的时候，那时辰到了。}因此祂经过，为要喂养我们；让我们跟随，好使我们被喂养。\par

\SourceRecord{Translated by R.G. MacMullen.；Nicene and Post-Nicene Fathers, First Series；Vol. 6.；Edited by Philip Schaff.；Buffalo, NY: Christian Literature Publishing Co.,；1888.；<http://www.newadvent.org/fathers/160353.htm>.}

% structure: p0001 source=h1 target=section reason=page-title

\section{论新约讲道集第54篇}

\EditorialNote{CIV. Ben.}\par

\Emph{再次，论福音的话，\BibleRef{路 10:38}等，关于玛尔大和玛利亚。}\par

1. 当宣读神圣福音时，我们听见主被一位虔诚的妇女迎入家中，她的名字叫玛尔大。当她忙于侍奉的操劳时，她的妹妹玛利亚却坐在主的脚前，聆听祂的圣言。一位忙碌，另一位安息；一位付出，另一位汲取。然而，玛尔大虽忙于那种侍奉的劳作和辛劳，却向主申诉，抱怨她的妹妹没有帮她工作。但主却为玛利亚答复了玛尔大；祂成了被诉诸为判官的那位的辩护者。\BibleQuote{玛尔大，}祂说，\BibleQuote{你为许多事操心，其实只有一件是必要的。玛利亚选择了更好的一份，是不能从她夺去的。}因为我们既听见了上诉者的申诉，也听见了判官的判决。这判决答复了上诉者，却为另一方辩护。因为玛利亚专注于主话语的甘饴，玛尔大却专注于如何喂养主；玛利亚专注于如何被主喂养。玛尔大为主预备了宴席，而玛利亚正在这宴席中喜悦自己。当玛利亚正以甜蜜的愉悦聆听祂最甘美的圣言，并以最热切的爱情喂养时，她的妹妹向主申诉，我们想，她是多么害怕主会对她说：\Quote{起来，帮你妹妹吧！}因为一种奇妙的甘饴抓住了她；一种心灵的甘饴，无疑比感官的更大。她被免除了，她坐在更大的信心之中。如何免除？让我们思考、查考、尽我们所能深入探究，以便我们也能被喂养。\par

2. 因为，我们难道想象玛尔大的侍奉被责备了吗？她因款待的操心而忙碌，曾将主亲自迎入家中。既然她因如此伟大的客人而喜乐，又怎能被正当地责备呢？如果这是真的，让人们放弃他们对穷人的服务；让他们为自己选择\BibleQuote{更好的一份，是不能从他们夺去的}；让他们完全献身于圣言，让他们渴慕教义的甘饴；让他们为得救的知识而操心；让他们不再关心街上有什么陌生人，有谁需要面包、衣服、或需要被探望、被赎回、被埋葬；让慈悲的工作停止，只热切关注知识。如果这是\BibleQuote{更好的一份}，为什么不是所有人都这样做呢？既然我们有主亲自为我们在这一方面辩护？因为我们在这件事上并不害怕会冒犯祂的公义，因为我们有祂的判断作为支持。\par

3. 然而并非如此；而是如主所说的那样。并非如你所理解，而是如你应当理解的那样。所以请注意：\BibleQuote{你为许多事操心，其实只有一件是必要的。玛利亚选择了更好的一份。}你没有选择坏的一份，但她选择了更好的。怎么更好呢？因为你\BibleQuote{为许多事操心}，她却为\BibleQuote{一件事}。一被优于多。因为多不是来自一，而是多来自一。\par

被造之物有很多，造物主却是一位。天、地、海和其中所有的一切，有多少啊！谁能数算？谁能想象其巨数？谁造了这一切？天主造了这一切。看，\BibleQuote{它们都很好。}祂所造之物非常好，造物主岂不更好！那么，让我们思考我们的\Quote{为许多事操心}。为保养我们的身体，许多服务是必要的。为何如此？因为我们饥饿、口渴。对可怜人施行慈悲是必要的。你给饥饿的人擘饼，因为你遇到了一个饥饿的人；消除饥饿，你给谁擘饼？消除无家可归的流浪，你款待谁？消除赤身露体，你给谁提供衣服？如果没有疾病，你探望谁？没有囚禁，你赎回谁？没有争吵，你与谁和好？没有死亡，你埋葬谁？在将来的世界里，这些灾祸将不再存在，因此这些服务也将不复存在。那么玛尔大在身体方面——我该说是主的缺乏，还是意愿？——侍奉了祂那必朽的血肉。但在这必朽血肉中的祂是谁？\BibleQuote{在起初已有圣言，圣言与天主同在，圣言就是天主：}看玛利亚在聆听什么！\BibleQuote{圣言成了血肉，居住在我们中间：}看玛尔大在侍奉谁！因此\BibleQuote{玛利亚选择了更好的一份，是不能从她夺去的。}因为她选择了那永远常存的；\BibleQuote{是不能从她夺去的。}她希望专注于\Quote{一件事}。她已经明白：\BibleQuote{但我亲近天主对我是好的。}她坐在我们元首的脚前。她坐得越是谦卑，领受得越是丰盛。因为水流向山谷的低洼处，从山丘的高处流下。那么主并没有责备玛尔大的工作，而是区分了她们的服务。\BibleQuote{你为许多事操心，其实只有一件是必要的。}玛利亚已经为自己选择了这一件。多样性的劳苦终将过去，而合一的爱情永存。因此她所选择的\BibleQuote{是不能从她夺去的。}但从你那里，你所选择的（当然这是下文，当然这是不言而喻的）从你那里，你所选择的将被夺去。但为了你的福乐，它将夺去，以便更好的可以赐予。因为劳苦将从你那里被夺去，以便安息可以赐予。你仍在海上，她已在港口。\par

4. 你们看，亲爱的诸位，我想你们已经明白，在这两位女子身上——她们都蒙主悦纳，都是祂所爱的对象，都是门徒——你们看，我说（这是一件重要的事，无论谁明白，都可以由此明白；就连那些不明白的人也应当留心听并知道）：在这两位女子身上，预表了两种生命：今世的生命与来世的生命，劳苦的生命与安息的生命，悲伤的生命与福乐的生命，暂时的生命与永恒的生命。这就是两种生命：你们要更充分地思考它们。这种生命包含什么？我不是指罪恶、不义、邪恶、奢靡或不敬神的生命，而是指劳苦、充满忧愁、被恐惧压制、被诱惑搅扰的生命——我指的是这种无害的生命，就如\Person{玛尔大}所适合的。我说，你们要尽力省察这种生命，并如我所说，比我所说的更充分地思考它。但邪恶的生命远离那房屋，既不在\Person{玛尔大}身上，也不在\Person{玛利亚}身上；若曾有，也在主进入时逃走了。于是，在那接待了主的房屋里，两位女子身上存留了两种生命：都是无害的，都是可称赞的；一种是劳苦的，另一种是安息的；既不邪恶，也不懒惰。两者都是无害的，我说，都是可称赞的：但一种是劳苦，另一种是安息；前者不可邪恶（劳苦的生命必须提防这），后者不可懒惰（安息的生命必须提防这）。于是，在那房屋里有了这两种生命，并有祂自己——生命之泉。在\Person{玛尔大}身上，是现今之事的形象；在\Person{玛利亚}身上，是将来之事的形象。\Person{玛尔大}所做的，我们现今在做；\Person{玛利亚}所做的，我们所盼望的。让我们好好做前者，好能完全获得后者。因为我们现今拥有后者的什么？我们拥有多少？只要我们在此世，我们拥有它多少？因为如今我们在某种程度上从事后者，而你们在放下事务、抛开家务、聚集、站立、聆听时，也是如此。你们这样做，便像\Person{玛利亚}。而且你们做\Person{玛利亚}所做的事，比我——必须分施的人——更容易。然而，若我说什么，那都是基督的；因此它喂养你们，因为是基督的。因为那饼是我们众人共有的，我也与你们一样靠它而活。\BibleQuote{但如今我们活着，如果你们，弟兄们，在主内站稳了。}\BibleRef{得前 3:8}我不愿你们在我们内站稳，而是在主内。\BibleQuote{因为栽种的不算什么，浇灌的也不算什么，只有那使生长的天主才算什么。}\BibleRef{格前 3:7}\par

\SourceRecord{Translated by R.G. MacMullen.；Nicene and Post-Nicene Fathers, First Series；Vol. 6.；Edited by Philip Schaff.；Buffalo, NY: Christian Literature Publishing Co.,；1888.；<http://www.newadvent.org/fathers/160354.htm>.}

% structure: p0001 source=h1 target=section reason=page-title

\section{论新约讲道五十五}

\EditorialNote{CV. Ben.}\par

\Emph{论福音的话，\BibleRef{路 11:5}，\Quote{你们中间谁有朋友，半夜到他那里去，}等等。}\par

1. 我们听到了主，天上的导师和最忠信的劝告者，敦促我们，祂一面敦促我们求，一面在我们求时赐予。我们在福音中听到祂敦促我们要恳切地求，要叩门，甚至像那种冒昧的纠缠。因为祂给我们立了一个榜样，说：\BibleQuote{你们中间谁有朋友，半夜到他那里去，对他说：朋友，借给我三个饼；因为我的朋友从路上来到我那里，我没有甚么给他摆上。那人在里面回答说：不要搅扰我，门已经关上，孩子们也同我在床上了，我不能起来给你。我告诉你们，虽不因他是朋友而起来给他，但因他恬不知耻地直求，就必起来照他所需用的给他。} 他希望多少呢？他只希望三个。主将这个比喻后，又加上劝勉，恳切地鼓励我们祈求、寻找、叩门，直到我们得到所求、所寻、所叩的，用了一个反面的例子：像那\BibleQuote{不敬畏天主，也不尊重人的}审判官，当一个寡妇天天求他时，他因她的纠缠而给了她，并非出于仁慈，而是出于无奈。但我们的主耶稣基督，祂在我们中间是求者，在天主面前是赐予者，若祂不愿赐予，绝不会如此强烈地敦促我们祈求。那么，人的懒惰应当羞愧；祂更愿意赐予，超过我们愿意领受；祂更愿意施恩，超过我们愿意脱离痛苦；无疑，若我们不得救，我们将留在痛苦中。因为祂给我们的劝勉，只是为我们自己的益处。\par

2. 让我们醒寤，相信敦促我们的祂，服从许诺我们的祂，并在赐予我们的祂内喜乐。因为也许有时也有朋友从路上来到我们这里，我们却找不到什么给他摆上；在这种困境中，我们为自己和他得到了所需。因为不可能我们当中有人未曾遇到一个朋友问他某事，而他却无法回答；然后他发现自己没有，当被催迫给予时。一个朋友从路上来到你这里，即从这个世界的生命之路而来，在这路上所有人都如旅客经过，无人作为占有者常住；但每人听到的是：\Quote{你已得享安逸，过去吧，走吧，让给后来者。} 或者，也许从一个邪恶的\Quote{路}，即从邪恶的生活中，你的一位朋友疲惫不堪，未能找到真理（听到并领悟真理才能使他幸福），却在世间的所有情欲和贫乏中耗尽，来到你这里，作为基督徒，说：\Quote{给我解释这个，使我成为基督徒。} 他问的可能正是你因信仰的纯朴而不知道的；于是你没有东西供他充饥，由此你发现了自己的匮乏；当你想教导时，却被迫学习；当你在问他的人面前惭愧，因为自己找不到他所寻求的，你被迫去寻找，好让你配得找到。\par

3. 但你应在哪里寻找？除了主的书卷，还有哪里？他所问的或许记载在书内，但晦涩难懂。或许宗徒在某封书信中已阐明：但阐明的方式是你能读，却不能领会；你无法继续。因为问者催逼你；保禄本人，或伯多禄，或任何先知，都不许你询问。因为这一家现已与他们的主安息，而这生命中的无知是深沉的，也就是说，是深夜，你饥饿的朋友紧逼你。也许简单的信德足能应付你，却不足够应付他。那么，他应被抛弃吗？他应被赶出你的家吗？因此，向主本人，向与祂全家安息的那一位叩门祈祷，祈求，坚持。祂不会像比喻中的那位朋友一样，因纠缠而起来给你。祂愿意赐予；你叩门还未得到，继续叩；祂愿意赐予。祂愿意赐予却延迟，好让你因延迟而更渴望，免得快速赐予而被轻视。\par

4. 但当你得到三个饼，即为了滋养并领悟\OriginalTerm{天主圣三}{Trinity}，你就有了一样东西，既可自己生活，也可喂养别人。现在你无需害怕从路上来的陌生人，反而接待他，使他成为家中公民；你也不必害怕它会用完。那个饼不会用完，反而会结束你的匮乏。那是饼——天主圣父，饼——天主圣子，饼——天主圣神。父是永恒的，圣子与祂永恒，圣神也与祂永恒。父是不变的，圣子是不变的，圣神是不变的。父是创造者，圣子也是，圣神也是。父是牧者、生命赐予者，圣子也是，圣神也是。父是食物和永生的饼，圣子也是，圣神也是。你学习，并教导；你自己生活，并喂养别人。赐予你的天主，赐予你的没有比祂自己更好的。哦贪婪的人，你还在寻求别的什么？若你寻求别的，若天主都不能使你满足，还有什么能足够呢？\par

5. 但你们必须有爱德，有信德，有望德，如此所赐的才成为你们的甘饴。这三样就是信德、望德、爱德。它们也是天主的恩赐。因为信德我们是从祂领受的；祂说：\Quote{天主按各人信德的量分给了各人。} 望德也是从祂领受的，向祂祈祷说：\Quote{祢使我怀有希望。} 爱德也是从祂领受的，经上说：\Quote{天主的爱德藉着赐给我们的圣神，已倾注在我们心中。} 这三样在某种程度上各不相同，但都是天主的恩赐。因为\Quote{现今存在的，有信德、望德、爱德这三样，但其中最大的是爱德。} 在那些饼中，并没有说哪一块饼比别的更大，只是说求了三块饼，也给了三块饼。\par

6. 再看另外三样：\Quote{你们中间有哪个人，儿子向他求饼，反而给他石头呢？或者求鱼，反而给他蛇呢？或者求鸡蛋，反而给他蝎子呢？你们纵然不善，尚且知道把好东西给你们的儿女，何况你们在天之父，岂不更把好东西赐给求祂的人吗？} 那么，让我们再思考这三样，是否就是\Quote{信德、望德、爱德，但其中最大的是爱德}。且把这三样记下：饼、鱼、蛋；其中最大的是饼。因此，在这三样中，我们以\Quote{饼} 很好地理解爱德。为此，祂把石头与饼对立，因为刚硬与爱德相反。我们以\Quote{鱼} 理解信德。有一位圣人说过，我们也乐意说：\Quote{\InnerQuote{好鱼}是虔敬的信德。} 它生活在波涛中，却不为波涛所冲垮或溶解。在今世的试探与风暴中，虔敬的信德生活着；世界狂怒，它却安然无恙。请注意，蛇与信德对立。因为我的信德是新娘，在\BibleBook{}{雅歌}中对她说：\Quote{我的新娘，从黎巴嫩来，从信德的起始向我走来，越过一切。} 因此也称为新娘，因为信德是订婚的开始。新郎应许了某事，并藉此信誓而受约束。主把蛇与鱼对立，即魔鬼与信德对立。为此，宗徒对这位新娘说：\Quote{我已把你们许配给一个丈夫，把你们当作贞洁的童女献给基督。} 又说：\Quote{我怕你们的心意，像蛇用它的狡猾诱惑了厄娃一样，受了败坏，失去那在基督内的纯一；} 也就是说，失去对基督的信德。因为他说：\Quote{但愿基督因着你们的信德住在你们心中。} 因此，不要让魔鬼败坏我们的信德，不要让牠吞噬了鱼。\par

7. 剩下的就是望德，我认为它被比作蛋。因为望德尚未达至拥有；蛋是一种东西，但还不是雏鸡。所以四足动物生幼崽，而鸟类生幼崽的希望。因此，望德劝勉我们：轻视现在，期待将来；\Quote{忘记背后}，让我们与宗徒一起努力向前。因为他这样说：\Quote{我只顾一件事：忘记背后的，努力前面的，向着标竿直跑，要得天主在基督耶稣内从高天召我来得的奖赏。} 因此，没有什么比\Quote{回头} 更敌视望德，即把希望寄托在转瞬即逝的事物上；而应当把希望寄托在尚未赐下、但终将赐下且永不消逝的事物上。但当世界被试炼所淹没，如同索多玛的硫磺火雨，必须畏惧罗特妻子的榜样。因为她\Quote{回头一看}，就在她回头的地方，她停住了。她变成了盐柱，好以她的榜样警醒智慧人。关于这望德，保禄宗徒这样说：\Quote{我们得救是在于望德；看得见的望德不是望德，因为人看得见了，为什么还希望呢？但如果我们希望那看不见的，我们就耐心等待。} 它是蛋，还不是雏鸡。它被蛋壳覆盖着；看不见，因为被覆盖；要耐心等待它；让它感受温暖，从而获得生命。努力，\Quote{努力前面的，忘记背后。因为看得见的是暂时的。} 他说：\Quote{不是顾念所见的，乃是顾念所不见的。因为所见的是暂时的，所不见的是永远的。} 所以，把你们的望德伸展到那看不见的，等候，忍耐。不要回头。要为你们的\Quote{蛋} 畏惧\Quote{蝎子}。请看它如何用身后的尾巴螫人。因此，不要让\Quote{蝎子} 压碎你们的\Quote{蛋}，不要让这个世界用它的毒液（可以说）粉碎你们的望德，因为这毒液是冲着你们来的，因为它在后面。这个世界怎样大声地向你们说话，怎样在你们背后喧嚷，好让你们回头！就是说，好让你们把希望寄托在目前的事物上（其实连眼前也算不上，因为不能称为稳固的事物），并且把心从基督应许却尚未赐下之物转开，而祂是信实的，必将赐下；并且你们满足于在败坏的世界中寻求安息。\par

8. 为此，天主将苦楚掺入世间的幸福中，好让人寻求另一种幸福，其甘甜毫无欺骗；然而，世界竟用这些苦楚试图使你放弃对\Quote{前面的事}的渴求，并引你回转。因为这些苦楚、这些磨难，你抱怨说：\Quote{看，在基督徒的时代，一切都在毁灭。}这是何等怨言！天主并未应许我这些事物不会毁灭；基督并未应许我这些。永恒者应许了永恒的事物：我若相信，必从必死之人变为永恒。这是何等的喧嚷，不洁的世界！何等的怨言！你为何试图引我回转？你这必朽的，想留住我；你若具有任何恒久性，又会做什么？你以全部苦楚向我们施以虚假的滋养，若非如此，还有谁不被你的甘甜所迷惑？至于我，若我有望德，若我坚守我的望德，我的\Quote{蛋}并未被\Quote{蝎子}所伤。\Quote{我时时赞美上主，对祂的赞颂常在我口。}世界繁荣，或世界颠倒，\Quote{我时时赞美上主}，祂创造了世界。是的，实在，我要赞美祂。按肉身而言，我或顺遂，或困厄，\Quote{我时时赞美上主，对祂的赞颂常在我口。}因为若我在顺遂时赞美，在困厄时亵渎，我便受了\Quote{蝎子}之刺，被刺而\Quote{回头观看}——愿这事远离我们。\Quote{上主赐予，上主收回：已照上主所愿成就；愿上主的名受赞美。}\par

9. 那按肉身生养我们的城邑依然存留，感谢天主。愿它获得属灵的重生，并与我们一同进入永恒！若那按肉身生养我们的城邑不存留，那按圣神生养我们的城邑却永远存留。\Quote{上主建造耶路撒冷。}祂岂因睡觉而使祂的工程毁坏，或因不看守而让仇敌进入？\Quote{若不是上主守护城邑，守护者就枉然警醒。}又是哪座\Quote{城邑}？\Quote{保护以色列的那位，既不打盹，也不睡觉。}以色列是什么？岂非亚巴郎的后裔？亚巴郎的后裔是什么？岂非基督？\Quote{和你的后裔}，他说，\Quote{就是基督。}对我们，他又说什么？\Quote{但你们属于基督，因此是亚巴郎的后裔，按照恩许成为继承人。}祂说：\Quote{地上万民都要因你的后裔蒙受祝福。}这圣城、忠信之城、在世上作客旅的城，其根基在天上。忠信者啊，不要败坏你的望德，不要丧失你的爱德，\Quote{束上你们的腰}，点亮并举起你们的灯，\Quote{等候主，看祂何时从婚宴归来。}你为什么惊慌，因为世上的王国正在毁灭？正因如此，一个天上的王国已应许给你，使你不与世上的王国一同灭亡。因为这事已预先宣告，明确宣告它们要毁灭。我们无法否认这是预言。你所等候的主已告诉你：\Quote{民族要起来攻击民族，王国要起来攻击王国。}世上的王国有其变迁；祂必来临，论及祂说：\Quote{祂的王国将永无穷尽。}\par

10. 那些向世上王国应许此事的人，并非受真理指引，而是以谄媚说谎。他们的一位诗人让朱庇特发言，论及罗马人说：\par

\Quote{我授予他们无限的帝国，\newline{}没有年限，永世不绝。}\par

那真理绝不如此回答。你赐予\Quote{没有年限}的这帝国，是在地上，还是在天上？无疑是在地上。即使在天上，\Quote{天地都要过去。}天主亲自创造的事物都要过去，罗慕路斯所建立的，岂不更快过去？或许我们若有意质问维吉尔，并嘲讽地问他为何这样说，他会私下拉我到一旁，对我说：我和你一样清楚这事，但我能做什么？我是向罗马人兜售言辞的人，若不以这种谄媚应许虚假之事，还能怎样？然而，即使在此处我也谨慎了：当我说\InnerQuote{我授予他们无限的帝国}时，我让他们的朱庇特说这话；我并未以自己名义说这谎言，而是将不真实之辞归于朱庇特：正如神是虚假的，诗人也是虚假的。因为你们要知道，我深知真理？在另一处，当我未引入这块被称为朱庇特的石头，而用自己名义说话时，我说：\par

\Quote{罗马国即将倾覆的灾难。}\par

看，我如何论及国度的即将倾覆。我提到了它的倾覆，没有隐瞒。当他以真理说话时，并未隐瞒其倾覆；当以谄媚说话时，却应许它永存。\par

11. 那么，诸位弟兄，我们不要沮丧：所有世上的国度都有一个终结。若这终结就在现在，只有天主知道。因为也许它还未到，而我们出于某种软弱、怜悯或痛苦，希望它不要到来；然而它终究会来，难道不是吗？将你们的希望寄托于天主，渴求永恒之物，期待永恒之物。诸位弟兄，你们是基督徒，我们都是基督徒。基督降生成人并非为使我们安逸度日；让我们忍受而非爱恋眼前的事物；逆境的伤害显而易见，顺境的谄媚却是欺骗。即使在风平浪静时也要畏惧大海。无论如何，不要让我们白白听见：\Quote{请举起我们的心。}为何我们将心放在地上，当我们看到大地正在倾覆？我们不能不劝勉你们，使你们能有话可说，为你们的希望辩护，以对抗那些嘲笑和亵渎基督徒名号的人。不要让任何人借着他的怨言使你们背离对将来的期盼。所有因这些逆境而亵渎我们基督的人，都是\Quote{蝎子}的尾巴。让我们将自己的蛋放在那福音中的母鸡的翅膀下，它向那座虚假而被弃的城市呼喊：\BibleQuote{耶路撒冷，耶路撒冷，我多次愿意聚集你的子女，如同母鸡聚集小鸡，而你却不愿意！}愿它不对我们说：\BibleQuote{我多次愿意，而你却不愿意！}因为那只母鸡就是\OriginalTerm{上智}{Divine Wisdom}；但它取了肉身以使自己适应于小鸡。请看那母鸡，羽毛竖起，翅膀垂下，声音破碎、颤抖、微弱、颓唐地适应于她的小雏。那么，我们的蛋，即我们的希望，让我们放在这只母鸡的翅膀下。\par

12. 你们或许注意到，母鸡如何将蝎子撕碎。哦，但愿福音中的母鸡也能将这些从洞穴中爬出、施放伤人毒刺的亵渎者撕碎并吞下，将他们融入自己的身体，并转化为一个蛋。让他们不要发怒；我们似乎激动了，但我们不以咒骂还咒骂。\BibleQuote{我们被咒骂，就祝福；被毁谤，就劝诫。}但\Quote{不要让他谈论罗马，}人们这样说我：\Quote{哦，但愿他闭嘴不谈罗马；}好像我在侮辱它，而不是为它向主恳求，并劝勉你们这些不配的人。我决不会侮辱它！愿主使这一点远离我的心，远离我良心的痛苦。我们在那里不是有许多弟兄吗？现在没有吗？旅途中的耶路撒冷城不是有很大一部分住在那裡吗？它在那里不是忍受了暂时的苦难吗？但它没有失去永恒的事物。那么，当我谈论罗马时，我能说什么呢？难道他们说我们的基督是罗马的毁灭者，而木石的神祇是它的保卫者，这不是假的吗？加上更贵重的：\Quote{铜神}。加上更为贵重的：\Quote{银神和金神}：\BibleQuote{外邦人的偶像，是金银。}他没有说\Quote{石}，没有说\Quote{木}，没有说\Quote{泥}；而是说他们所珍视的\Quote{银和金}。然而这些金银的偶像\Quote{有眼，却看不见。}金神、木神，就价值而言，是不平等的；但就\Quote{有眼，却看不见}而言，是平等的。看哪，有学识的人将罗马托付给了什么样的守护者——托付给那些\Quote{有眼，却看不见}的。或者，如果他们能保存罗马，为什么他们自己先灭亡了？他们说：\Quote{罗马同时灭亡了。}然而他们灭亡了。\Quote{不，}他们说，\Quote{他们自己没有灭亡，只是他们的雕像灭亡了。}那么，他们连自己的雕像都保存不了，又如何能保存你们的房屋呢？亚历山大里亚一度失去了这样的神祇。君士坦丁堡自成为大城以来——它是由一位基督徒皇帝建造的——早已失去了它的假神；然而它增长了，并且仍在增长，依然存在。只要天主愿意，它就会存在。因为我们这样说，也不是向这座城市承诺永恒。迦太基如今在拥有基督之名的状态下依然存在，然而曾经它的\OriginalTerm{天女神}{Cælestis}被推翻了；因为她不是天上的，而是地上的。\par

13. 他们所说\Quote{一失去诸神，罗马就被攻陷并毁灭}，这并不真实。这完全不真实；他们的偶像先前已被推翻；即便如此，哥特人仍与拉达盖斯一同被征服。请记住，我的弟兄们，请记住；这并非久远之事，就在几年前，回想起来吧。当罗马城中的所有偶像被推翻后，哥特王拉达盖斯率领大军前来，其人数远多于后来的阿拉里克。拉达盖斯是异教徒，他每日向朱庇特祭献。四处传闻拉达盖斯从不停止祭献。于是他们都说：\Quote{看，我们不祭献，他祭献；我们这些不被允许祭献的人，必定被那祭献者征服。} 但天主证明，即使今世的拯救与这些现世国度的保全，也不在于这些祭献；拉达盖斯赖主的帮助奇妙地被击败了。其后另有一些不祭献的哥特人来了；他们虽非基督信仰中的天主教徒，却敌视并反对偶像；他们攻占了罗马；他们征服了那些信赖偶像、仍在寻觅已失偶像并渴望向失落的诸神祭献的人。这些人中也有我们的一些弟兄，他们也一同受苦；但他们已学会说：\BibleQuote{我时时赞美主。} 他们被卷入今世国度的苦难中，却没有失去天国；反而由于经历了磨难的锻炼而得以更好地获得天国。若他们在磨难中没有亵渎，便如纯净的器皿从熔炉中出来，充满了主的祝福。而那些亵渎者，追随并渴慕现世之事，将希望寄托于现世之物，一旦失去这些，无论他们愿意与否，他们还能保有什么？他们能安身何处？外在虚无，内在空虚；空空的箱柜，更空的良心。他们的安息在哪里？他们的救恩在哪里？他们的希望在哪里？那么，让他们来吧，停止亵渎，学习崇拜；让带刺的蝎子被母鸡吞食，让他们转化为那使他们逾越进入祂身体者的身体；让他们在世上受锻炼，在天堂得到冠冕。\par

\SourceRecord{Translated by R.G. MacMullen.；Nicene and Post-Nicene Fathers, First Series；Vol. 6.；Edited by Philip Schaff.；Buffalo, NY: Christian Literature Publishing Co.,；1888.；<http://www.newadvent.org/fathers/160355.htm>.}

% structure: p0001 source=h1 target=section reason=page-title

\section{新约讲道集第五十六篇}

\Emph{[CVI. Ben.]}\par

\Emph{论福音之言，\BibleRef{路 11:39}，\Quote{如今你们法利塞人洗净杯盘的外面}等}\par

1. 你们已听到福音，主耶稣对那些法利塞人所说的话，无疑向自己的门徒传达了一个教训，就是不可认为义德在于身体的洁净。因为法利塞人每天用餐前都要用水沐浴，仿佛每日的沐浴能洁净人心。然后祂显明了他们是怎样的人。祂对他们说话，祂看见他们；因为祂不仅看见他们的外表，也看见他们的内心。为使你们知道这一点，那个法利塞人，基督回答了他，他心中思想，没有出声，但主听见了他。因为他在心里责备主基督，因为祂没有先沐浴就来赴他的宴席。他正想着，主听见了，因此祂回答。祂回答了什么？\BibleQuote{如今你们法利塞人洗净盘子的外面，但你们里面却充满了诡诈和贪婪。}什么！这是赴宴吗！祂如何不饶恕那邀请祂的人？相反，祂借着责备饶恕了他，使他悔改后能在审判中得赦免。祂向我们表明了什么？就是那只举行一次的圣洗圣事，借着信德洁净。如今信德是内在的，不是外在的。因此，在\BibleBook{宗徒}{大事录}中记载并读到了\BibleQuote{借着信德洁净了他们的心。}宗徒\Person{伯多禄}也在他的书信中这样说：\BibleQuote{同样，祂给了你们一个诺厄方舟的比喻，八个人借着水得救了。}然后他补充说：\BibleQuote{同样，这预表的水如今也拯救你们，不是除去肉体的污秽，而是向天主许下纯洁的良心的答复。}法利塞人轻视了\BibleQuote{这纯洁良心的答复}，只清洗\Quote{外在的}，而他们里面却仍充满污秽。\par

2. 此后祂对他们说了什么？\BibleQuote{但是你们施舍，看，一切对你们就洁净了。}看施舍的赞美，去做并证实它。但稍等片刻；这话是对法利塞人说的。这些法利塞人是犹太人，仿佛是犹太人的精英。因为当时那些最有声望和学识的被称为法利塞人。他们尚未受基督的圣洗；他们尚未相信基督，天主的独生子，祂在他们中间行走，却不被他们认识。那么，祂如何对他们说\BibleQuote{你们施舍，看，一切对你们就洁净了}？如果法利塞人听从祂，并施舍，立刻按祂的话\BibleQuote{一切对他们就洁净了}；那么他们还需要相信祂吗？但如果他们不能得洁净，除非相信祂，祂\BibleQuote{借着信德洁净了人心}；那么\BibleQuote{你们施舍，看，一切对你们就洁净了}是什么意思？让我们仔细思考这一点，也许祂自己会解释。\par

3. 祂这样说了以后，无疑他们以为他们确实施舍了。但他们如何施舍？他们献上一切所有，献上所有出产的十分之一，并施舍。很难找到一个基督徒做得如此之多。看看犹太人做了什么。不仅是麦子，还有酒和油；不仅如此，甚至是最微小的东西，如莳萝、芸香、薄荷和茴香，为遵从天主的诫命，他们献上十分之一；也就是说，拿出十分之一，并以此施舍。我想他们于是想起了这一点，认为主基督在空谈，像是对那些不施舍的人说话；而他们知道自己的行为，他们如何献十分之一，并施舍了最微小的出产。他们心里嘲笑祂这样说话，仿佛是对不施舍的人说的。主知道这一点，立刻接着说：\BibleQuote{但是祸哉，你们经师和法利塞人！你们将薄荷、莳萝、芸香和各种草本献上十分之一。}你们要知道，我知道你们的施舍。无疑，这些十分之一就是你们的施舍；是的，连你们最微小的果实也献上十分之一；\BibleQuote{可是你们忽略了法律上更重要的事，即公义和爱德。}注意。你们\BibleQuote{忽略了公义和爱德}，却献上草本十分之一。这不是施舍。祂说：\BibleQuote{这些是你们应当做的，而其他的也不可忽略。}做什么？\BibleQuote{公义和爱德，正义和怜悯}；并且\BibleQuote{不可忽略其他的。}做这些；但要将优先给予其他的。\par

4. 如果这样，祂为什么对他们说：\Quote{你们施舍，那么一切对你们便都洁净了}？什么又是\Quote{施舍}？就是\OriginalTerm{慈悲}{mercy}。什么又是\Quote{慈悲}？如果你明白，就从你自己开始。因为你若对自己残忍，怎能对他人慈悲？\Quote{施舍，那么一切对你们便都洁净了。}要实行真正的施舍。什么是施舍？慈悲。请听圣经；\BibleQuote{怜悯你的灵魂，以悦乐天主。}施舍吧，\BibleQuote{怜悯你的灵魂，以悦乐天主。}你的灵魂在你面前是个乞丐，回到你的良心吧。无论你是谁，生活在邪恶或不信中，回到你的良心；在那里你会发现你的灵魂在乞讨，你发现它贫乏，发现它贫穷，发现它悲伤，甚至也许你发现它并非缺乏，而是因匮乏而哑口无言。因为如果它乞求，它\BibleQuote{饥渴慕义。}现在，当你发现你的灵魂处于这种状态（这一切都在内心，在你心里），首先施舍，给它食物。什么食物？如果法利塞人问了这个问题，主会对他说：\Quote{施舍给你的灵魂。}因为祂确实对他说了这句话；但他不明白，当祂向他们列举他们惯常施舍的事，并以为基督不知道时，祂对他们说：\Quote{我知道你们做这些事，\InnerQuote{你们捐献薄荷、茴香、孜然和芸香；}但我所说的，是另一种施舍；你们轻视了\InnerQuote{判断和爱德}。要在判断和爱德中施舍给你的灵魂。}什么是\Quote{在判断中}？回头看，认识你自己；不喜欢你自己，对你自己的罪下判决。什么是\OriginalTerm{爱德}{charity}？\BibleQuote{你要全心、全灵、全意爱天主你的主；要爱你的近人如同你自己。}这样你就首先在你的良心里，对你的灵魂行了施舍。可是如果你忽略了这种施舍，无论你给出什么，给出多少；保留你的财物，不是十分之一，而是一半；给出九份，只留一份给自己：当你没有对你的灵魂施舍，而自己在内里贫穷时，你什么也没有做。让你的灵魂得到食物，免得它因饥饿而灭亡。给它食物。你会问：什么食物？祂亲自与你说话。如果你愿意听、明白并相信主，祂会亲自对你说：\BibleQuote{我是从天上降下的生活的食粮。}你岂不首先把这食粮给你的灵魂，并向它施舍吗？如果你信，你就应该这样做，以便先喂养你自己的灵魂。相信基督，里面的就洁净了，外面的也必洁净。\Quote{让我们转向主，}等等。\par

\SourceRecord{Translated by R.G. MacMullen.；Nicene and Post-Nicene Fathers, First Series；Vol. 6.；Edited by Philip Schaff.；Buffalo, NY: Christian Literature Publishing Co.,；1888.；<http://www.newadvent.org/fathers/160356.htm>.}

% structure: p0001 source=h1 target=section reason=page-title

\section{新约讲道集第57篇}

\EditorialNote{[CVII. 本笃版]}\par

\Emph{关于福音的话}，\BibleRef{路 12:15}\Emph{，\BibleQuote{祂对他们说：你们要谨慎，防备一切贪念。}}\par

1. 我毫不怀疑，你们敬畏天主的人，必敬畏地听祂的话，并喜乐地遵行；为的是你们现在能盼望，将来能领受祂所应许的。我们刚才听见主耶稣基督、天主子赐给我们一条诫命。那真理本身，既不欺骗人，也不受人欺骗，赐给了我们一条诫命；让我们听，让我们惧怕，让我们谨慎。那么，这条诫命是什么？\Quote{我告诉你们：要防备一切贪念}？什么是\Quote{一切贪念}？什么是\Quote{一切}？祂为什么加上\Quote{一切}？因为祂本来可以说：\Quote{要防备贪念。}祂适合加上\Quote{一切；}并说：\Quote{要防备一切贪念。}\par

2. 祂为什么说这话，其缘由，可说是从这些话语产生的，已在神圣福音中启示给我们。有一个人向祂控告他的弟兄，那弟兄夺走了他的全部家产，没有将他应得的份额还给他。你们看这个申诉者的案情多么好。因为他并非强取别人的东西，而只是寻求自己父母留给他的；他通过向主的审判申诉来要求归还。他有一个不义的弟兄，但面对一个不义的弟兄，他找到了一个公义的审判官。那么，在这样一个好的案件中，他怎能失去这个机会？或者，如果基督不说，谁会对他弟兄说：\Quote{将你弟兄的份额还给他}？那可能被他的更富有和勒索的弟兄用贿赂腐蚀的审判官会说这话吗？他这样孤苦，被剥夺了父亲的财产，当他找到了如此伟大的审判官时，他走上前去，向他申诉，恳求他，用简短的话将他的案子呈在他面前。因为当他向那能看透内心的主说话时，有什么必要详细陈述案情呢？他说：\Quote{师傅，请吩咐我的弟兄与我分家产。}主没有对他说：\Quote{叫你弟兄来。}不，祂既没有派人叫他来，也没有在他面前对那申诉的人说：\Quote{证明你所说的话。}他要求一半的产业，他要求地上的一半产业；主却赐给他天上的全部产业。主所赐的超过他所求的。\par

3. \Quote{对我弟兄说，叫他与我分家产。}公义的案件，简短的案件。但让我们听那同时审判和教导的主。祂说：\Quote{人哪。} \Quote{人啊；} 因为你把这产业看得如此重要，你不过是人吗？祂愿意使他超过人。祂愿意使他成为什么？从祂愿意除去贪念的人，祂愿意使他成为什么？我告诉你们：\Quote{我曾说：你们是神，都是至高者的儿子。} 看，祂愿意使他成为什么——把那没有贪念的人算在\Quote{神}中。\Quote{人哪，谁立我作你们的分家者？} 所以祂的仆人宗徒保禄，当他说：\Quote{我劝你们，弟兄们，你们都说一样的话，不可有分裂，} 他不愿作一个分裂者。后来他这样劝诫那些追随他名字并分裂基督的人：\Quote{你们各人说：\InnerQuote{我是属保禄的，} \InnerQuote{我是属阿颇罗的，} \InnerQuote{我是属刻法的，} \InnerQuote{我是属基督的。} 基督被分裂了吗？保禄曾为你们被钉十字架吗？或者你们受洗是归于保禄的名下吗？} 那么，判断吧，那些希望祂被分裂的人是何等邪恶，祂自己却不愿作一个分裂者。\Quote{谁，} 祂说，\Quote{立我作你们的分家者？}\par

4. 你曾为恩惠祈求；请听劝告。\Quote{我告诉你们：要防备一切贪念。} \Quote{或许，}他会说，\Quote{你称他为贪吝和贪婪，如果他寻求别人的财物；但我告诉你们，连你们自己的东西也不要贪婪地寻求。} 这就是\Quote{一切，要防备一切贪念。} 这是一个沉重的负担！如果这个负担竟加在软弱的人身上，当归向那位加给它的，求祂赐给我们力量。因为这并非小事，我的弟兄们，当我们的主、我们的救赎者、我们的救主，祂为我们而死，用祂自己的血作为赎价赎回了我们，我们的辩护者和审判者，祂说\Quote{要谨慎。}祂很清楚这罪恶有多大；我们不知道，让我们相信祂。\Quote{要谨慎，}祂说。为什么？从哪里？\Quote{一切贪念。} 我只是保存我自己的东西，我没有拿走别人的；\Quote{要防备一切贪念。} 不但是掠夺别人财物的人是贪婪的，而且吝啬地保存自己财物的人也是贪婪的。但如果吝啬地保存自己财物的人受责备，那么掠夺别人财物的人该受何等谴责呢！\Quote{要谨慎，}祂说，\Quote{防备一切贪念：因为人的生命不在于他拥有财产的富裕。} 那人积蓄了大量的财富，他从其中取多少来维持生活？当他取来，并在思想上分出足以生活的部分之后，让他想想其余的部分留给谁；免得当你存着生活所需的时候，你只是在积攒死亡所需。看，基督，看，真理，看，严厉。\Quote{要谨慎，}真理说：\Quote{要谨慎，}严厉说。如果你不爱真理，就要惧怕严厉。\Quote{人的生命不在于他拥有财产的富裕。} 相信祂，祂不欺骗你。另一方面，你说：\Quote{是的，\InnerQuote{人的生命}确实\InnerQuote{在于他拥有财产的富裕。}} 祂不欺骗你；你欺骗自己。\par

5. 于是，借着这次机会，当那个请求者寻求自己的一份，而不愿掠夺他人时，便有了主的那句话，祂说的不是\Quote{小心贪心}，而是加上了\Quote{一切贪心}。不仅如此，祂又举了一个富翁的例子，\Quote{他的田地出产丰盛}。\Quote{有一个}，祂说，\Quote{富翁，他的田地出产丰盛}。什么是\Quote{出产丰盛}？他拥有的田地生出了大量收成。大到何等程度？大到他找不到存放的地方；突然，因着他的丰裕，他变得窘迫——这个老贪心的人。多少年已经过去，那些仓房不也够用了吗？这收成如此之大，以致惯常的地方都不够用了。这可怜的人寻求主意，不是关于如何处理多余出产，而是如何囤积起来；在思想中他找到了一个计策。他自以为聪明，因发现这计策。他有意识地想到了它，聪明地发现了它。他所聪明发现的是什么呢？\Quote{我要拆毁}，他说，\Quote{我的}旧\Quote{仓房，另建更大的，将它们装满；然后我要对我的灵魂说}。你要对灵魂说什么？\Quote{灵魂，你存有许多财物够多年之用，只管享受吧，吃喝快乐吧！}这个聪明的发明者就这样对他的灵魂说。\par

6. \Quote{天主}，祂不屑于与愚昧人说话，\Quote{对他说}。你们中有人也许会说，天主如何与愚昧人说话？啊，我的弟兄们，当福音宣读时，祂在这里与多少愚昧人说话！当它被宣读时，那些听了却不实行的，岂不是愚昧人吗？主又说了什么？因为，我再说，他自以为聪明，因发现他的计策。\Quote{你这愚昧人}，祂说；\Quote{你这愚昧人}，自以为聪明；\Quote{你这愚昧人}，你对你的灵魂说过：\Quote{你存有许多财物够多年之用，今天就要索回你的灵魂了！}你的灵魂，你对它说过：\Quote{你存有许多财物}，今天就被\Quote{索回}，却毫无善工。那么，让她轻视这些财物，使自己成为善的，这样当被\Quote{索回}时，她能在确信的希望中离去。因为，还有什么比一个人希望拥有\Quote{许多财物}，却不愿自身成为善的更乖戾呢？你不配拥有它们，因为你不愿成为你所希望拥有的。你难道希望拥有一栋差的乡间别墅？当然不，而是好的。或是差的妻子？不，而是好的。或是差的风帽？甚至是差的鞋子？为何只有灵魂是差的？祂在此处没有对这个考虑虚妄之事、建造仓房却不关心穷人需要的愚昧人说；祂没有对他说：\Quote{今天你的灵魂将被匆匆带往地狱}；祂没说这样的话，而是\Quote{被索回}。\Quote{我不告诉你你的灵魂将往何处去；但从此地，你为她积攒如此之多东西的地方，它必离去，无论你愿意与否}。看，\Quote{你这愚昧人}，你想填满你的新且更大的仓房，好像你所有的东西没有用场似的。\par

7. 但也许他当时还不是基督徒。那么，让我们听听，弟兄们，对我们这些信徒，福音被宣读，是祂——那说这些话的——被敬拜，祂的印记被我们印在额头并存于心中。因为一个人有基督的印记在何处，是在额头，还是同时在额头和心中，关系重大。你们今天听到了圣先知\Person{厄则克耳}的话，在天主派一人毁灭不敬虔之民前，祂先派一人去标记他们，并对他说：\Quote{你去在那些为我的子民在他们中间所行的罪恶而叹息哀号的人额上画一个记号。}祂没有说\Quote{行在他们外面}，而是\Quote{在他们中间}。但他们\Quote{叹息哀号}，所以他们在\Quote{额上被画记号}：在内里之人的额上，而非外在。因为脸上有额头，良心中也有额头。因此，内里的额头受打击时，外在就变红；或因羞耻而红，或因恐惧而苍白。所以内里之人也有额头。他们在那里被\Quote{画记号}以免被毁灭；因为他们虽没有纠正\Quote{行在他们中间}的罪恶，却为此忧愁，并因这忧愁而分别出来；尽管在天主眼中是分别的，但在人眼中与他们混杂。他们秘密地被\Quote{画记号}，公开地不受伤害。之后，毁灭者被派出，并对他说：\Quote{去，毁灭殆尽，无论老幼男女，都不要饶恕，但不可靠近那些额上有记号的人。}我的弟兄们，你们在这些人中，为在你们中间所行的邪恶叹息哀号，却不实行它们，你们获得了何等大的保障！\par

8. 但为使你们不犯不义，\Quote{要提防各种贪得无厌。} 我要向你们详细说明什么是\Quote{各种贪得无厌。} 在情欲之事上，凡是自己的妻子不能满足的，就是贪得无厌。连拜偶像本身也被称为贪得无厌；因为，在敬拜天主之事上，凡是独一真天主不能满足的，就是贪得无厌。除了贪得无厌的灵魂，还有什么会为自己制造许多神祇？除了贪得无厌的灵魂，还有什么会为自己制造假殉道者？\Quote{要提防各种贪得无厌。} 看，你爱惜自己的财物，并因你不寻求别人的财物而自夸；看看你因不听基督的话——祂说\Quote{要提防各种贪得无厌}——而犯了多大的恶。爱惜你自己的财物，不要夺取别人的财物；你劳苦的果实，公正地属于你；你继承遗产，某人因你获得了他的好感而赠与你；你出海冒险，没有欺诈，没有发假誓，获得了天主愿意你所得的，你贪婪地存着，良心无愧，因为你并非来自邪恶之源，也非追求他人之物。然而，你若不听那说\Quote{要提防各种贪得无厌}的那位，就听听为了自己的财物你将准备犯多大的恶。比如，你偶然被任命为法官。你不会被收买，因为你不寻求别人的财物；没有人给你贿赂说：\Quote{判决不利于我的仇敌。} 这事绝不临到你，一个不寻求他人之物的人，怎么能被说服这样做呢？然而，看看为了自己的财物你将准备犯多恶。或许那希望你判得不公、为他作出有利于他的判决的人，是个有权势的人，有能力向你提出诬告，使你失去所有。你思忖，想到他的权势，想到你所保全的自己的财物，你所爱的：不是你曾经拥有的，而是你自身不幸被束缚于其权下。你这鸟胶，因它你没有自由的美德翅膀；你想到它，就心里说：\Quote{我冒犯了这个人，他在世上很有影响力；他会对我提出恶毒的指控，我会被剥夺法律保护，失去一切。} 这样，你便做出不义的判决，不是在贪求别人的，而是在保全自己的时候。\par

9. 给我一个听从基督的人，给我一个怀着恐惧听见\Quote{要提防各种贪得无厌；}的人，让他不要对我说：\Quote{我是一个穷人，出身低微的平民，普通百姓，我怎能期望成为法官？我不惧怕你放在我眼前的这诱惑的危险。} 然而，看，我要告诉这穷人他应当惧怕什么。某个富有的权势人物呼召你为他作假见证。你现在要做什么？告诉我。你有一点自己的小产业；你劳苦得来，获取并保有它。那人要求你：\Quote{为我作假见证，我就给你若干钱。} 你这位不寻求他人之物的人，说：\Quote{那远离我：我不寻求天主不愿给我的，我不接受；离开我。} \Quote{你不愿接受我给的？我要拿走你已经有的。} 现在看，证明你自己，现在质问你自己。为什么看着我？向内看你自己，看你自己内在，省察你自己内在；坐在你自己面前，传唤你自己到你面前，把你伸展在天主诫命的刑架上，用祂的恐惧折磨你自己，不要宽待你自己；回答你自己。看，如果有人这样威胁你，你会怎么做？\Quote{你若不为我作假见证，我就从你那里拿走你辛苦得来的一切。} 给他那个：\Quote{要提防各种贪得无厌。} \Quote{我的仆人，}祂会对你说，\Quote{是我救赎并释放了你，从仆人收养你为兄弟，立你为我身体的肢体，听我说：他可能拿走你所得的，却不能从我这里拿走你。你保全自己的财物，免得灭亡吗？怎么，我没有对你说过\InnerQuote{要提防各种贪得无厌}吗？}\par

10. 看，你正陷于混乱，颠簸不定；你的心如船一样被风暴摇撼。基督正在沉睡：唤醒沉睡的他，你就不再遭受风暴的侵袭。唤醒他，他乐意在此一无所有，而你却拥有一切，他为你甚至来到十字架上，当他赤身悬挂时，嘲笑他的人清点了他的\Quote{骨骼}，并\Quote{清点了}；并且\Quote{要谨慎，躲避一切贪吝}。贪财并非一切；要\Quote{谨慎贪吝}生命。一种可怕的贪吝，非常可畏的贪吝。有时，一个人会轻视他所拥有的，说：\Quote{我不做假见证；我不做。你对我说，我要夺走你所拥有的。夺走我所拥有的吧；你没有夺走我内在的。因为说过\InnerQuote{主赐予，主取走；愿主的旨意成就}的人并未成为穷人；因此\InnerQuote{愿主的名受赞美}。\InnerQuote{我赤身出于母胎，也要赤身归回大地。}外表赤身，内心穿着华衣。外表穿着这些破布，这些腐朽的破布，内心却穿着衣服。穿什么呢？\InnerQuote{愿你的司铎披上正义。}}但是，如果他对你说，当你轻视你所拥有的之后，他对你说：\Quote{我要杀你}？如果你听从了基督，回答他：\Quote{你要杀我吗？你最好杀死我的身体，也不要我通过说谎的舌头杀死自己的灵魂！你能对我做什么？你将杀死我的身体；我的灵魂将自由地离去，在世界末日再领回这具她轻视过的身体。那么你能对我做什么呢？然而，如果我为你作假见证，我借你的舌头杀死自己；不是在身体上杀死自己；\InnerQuote{因为说谎的口杀害灵魂。}}但或许你不这样说。为什么你不这样说？你想活着；你想活得比天主为你指定的更长？那么你\Quote{要谨慎，躲避一切贪吝}？天主曾愿意你活到这个人来到你面前。他可能会杀死你，使你成为殉道者。那么不要对生命怀有过度的渴望；这样你就不会得到永死。你看，贪吝无处不在，当我们渴求超过所需时，就导致我们犯罪。让我们谨慎躲避一切贪吝，如果我们愿享受永恒智慧。\par

\SourceRecord{Translated by R.G. MacMullen.；Nicene and Post-Nicene Fathers, First Series；Vol. 6.；Edited by Philip Schaff.；Buffalo, NY: Christian Literature Publishing Co.,；1888.；<http://www.newadvent.org/fathers/160357.htm>.}

% structure: p0001 source=h1 target=section reason=page-title

\section{新约讲道第58篇}

\Emph{[CVIII. Ben.]}\par

\Emph{论福音的话}，\BibleRef{路 12:35}\Emph{，\BibleQuote{你们腰里要束上带，灯也要点着；你们自己也要像，}等等。又论}\Emph{\BibleRef{咏 33:12}}\Emph{的话，\BibleQuote{谁愿得生命，}等等。}\par

1. 我们的主耶稣基督曾来到人间，也离开了人间，并将再次来到人间。然而祂来的时候，祂就在这里；祂离开的时候，祂并未离去；祂将要来到那些祂曾对之说过这话的人当中：\BibleQuote{看，我同你们天天在一起，直到今世的终结。}按照祂为我们所取的\BibleQuote{仆人的形象}，祂在特定时间出生，被杀害，又复活了，如今\BibleQuote{祂不再死，死亡再也不能统治祂}；但按照祂的神性——祂与圣父同等——祂早已在这世界上，\BibleQuote{世界是藉着祂造成的，但世界不认识祂。}关于这一点，你们刚才听到了福音，它给了我们何等劝诫，提醒我们保持警惕，希望我们无牵无挂、准备迎接终结；好使这些今世可畏的最终事物过去之后，那没有终结的安息得以到来。有福的是那些将分享这一安息的人。因为那时，那些如今不得安宁的人将得安宁；而那时，那些如今不怕的人将要害怕。为了这等待，为了这希望，我们成了基督徒。我们的希望难道不属于这世界吗？那么，让我们不要爱世界。我们被召唤而脱离对这世界的爱，好使我们盼望并爱上另一个世界。在这世上，我们应当戒绝一切邪欲，也就是说，要\BibleQuote{束上腰}；并且要热忱、以善工发光，也就是说，要\BibleQuote{灯要点着}。因为主亲自在福音的另一处对祂的门徒说过：\BibleQuote{没有人点灯放在斗底下，而是放在灯台上，照亮屋中所有的人。}为了说明祂所说的是什么，祂又接着说：\BibleQuote{照样，你们的光也当在人前照耀，使他们看见你们的善工，光荣你们在天之父。}\par

2. 因此，祂愿意我们\BibleQuote{束上腰，灯要点着}。什么是\BibleQuote{束上腰}？\BibleQuote{远离恶事。}什么是\BibleQuote{点着}？什么是使我们的\BibleQuote{灯点着}？就是\BibleQuote{行善。}祂后来所说的是什么：\BibleQuote{你们也要像等候自己的主人从婚宴回来的人一样：}不正是那圣咏紧接着的话：\BibleQuote{寻求和平，追随和平}吗？这三件事，即\BibleQuote{戒恶行善}，以及对永恒赏报的盼望，记载在宗徒大事录中，其中写着，保禄教导他们关于\BibleQuote{节制和义德}，以及对永生的盼望。节制属于\BibleQuote{你们腰里要束上带}；义德属于\BibleQuote{灯要点着}；对永生的盼望属于等候主。因此，\BibleQuote{远离恶事}，这就是节制，这就是束上腰；\BibleQuote{行善}，这就是义德，这就是\BibleQuote{灯点着}；\BibleQuote{寻求和平，追随和平}，这就是等候将要到来的世界。所以，\BibleQuote{你们要像等候自己的主人从婚宴回来的人一样。}\par

3. 既然我们有了这些诫命和许诺，为什么我们还在世上寻求\Quote{好日子}，而在那里我们却找不到它们？因为我知道你们确实寻求它们，当你们生病，或遭遇到任何在这世上充斥的忧患时。因为当生命接近结束时，老人满口怨言，毫无喜乐。在人类所磨损的一切忧患中，人们只寻求\Quote{好日子}，并希望长寿，而在这里他们无法得到。因为即使一个人的长寿，相较于所有时代的广阔范围，也被压缩在如此短暂的瞬间，仿佛只是整个海洋中的一滴。那么，人的生命，即使被称为长寿的，又是什么呢？他们称之为长寿的，甚至在这世间的历程中也是短暂的；而且如我所说，呻吟甚至直到衰老的晚年也层出不穷。这至多是短暂的，时间不长；然而人们追求它多么急切，以何等的勤勉、何等的辛劳、何等的谨慎、何等的警觉、何等的劳力，人们寻求在这里活得长久，直到老年。然而这长寿本身，除了是奔向终点，又是什么呢？你曾拥有昨天，你也希望拥有明天。但当今天和明天过去之后，你就不再拥有它们。因此你希望破晓，好使那你不愿接近的终点临近你。你和朋友举行某个年度庆典，在那里听到祝福者向你说：\Quote{愿你长命百岁}，你希望他们所说的话得以实现。什么？你希望岁月一年年到来，而这些岁月的尽头却不来吗？你的愿望彼此矛盾；你希望前行，却不愿到达终点。\par

4. 但是，如我所说，人们如此费尽心力，日日以巨大的、持续不断的辛劳，只求稍晚些死去：那么他们岂不更应当努力，以求永不死呢？这一点，却无人思虑。日复一日，人们在此世寻求\Quote{幸福的日子}，却找不到；然而没有谁愿意如此生活，好能到达那能找到幸福之处的所在。因此，同一个圣经劝诫我们说：\BibleQuote{谁愿生命长久，渴望看到幸福的日子？}圣经如此提问，正如它深知人会如何回答；它知道所有人都会\Quote{追求生命和幸福的日子}。它依照人的愿望发问，仿佛所有人心中的回答将是\Quote{我愿意}；它便这样说：\BibleQuote{谁愿生命长久，渴望看到幸福的日子？}正如就在此时此刻我向你们讲话时，你们听到我说：\BibleQuote{谁愿生命长久，渴望看到幸福的日子？}你们都在心中回答：\Quote{我。}因为连我这样与你们讲话的，也\Quote{渴望生命和幸福的日子}；你们寻求的，我也寻求。\par

5. 正如金子为我们众人所必需，我们所有人，包括我和你们，都想要得到金子；倘若在你们田地里的某个地方，在你们权力管辖之处，有一块金子，我见你们在寻找，便对你们说：\Quote{你们找什么？}你们回答我：\Quote{金子。}我便对你们说：\Quote{你们在找金子，我也在找金子；你们找的，我也在找；可是你们找错了地方，不在我们能找到的地方找。那么请听我说，我们该在哪里找；我不是要从你们那里夺走它，而是指给你们地点；}是的，让我们都跟随那位知道我们所寻求之物的所在者。同样，现在你们渴望\Quote{生命和幸福的日子}，我们不能对你们说：\Quote{不要渴望\InnerQuote{生命和幸福的日子}}；而是这样说：\Quote{不要在此世寻求\InnerQuote{生命和幸福的日子}，在此世，\InnerQuote{幸福的日子}不可能存在。}这生命本身难道不就像死亡吗？这些日子匆匆流逝：因为今天抹去了昨天；明天升起，只是为了抹去今天。这些日子本身没有常存；你们为何要与它们一同存留？那么，你们渴望\Quote{生命和幸福的日子}的愿望，我不但不抑制，反而更加强烈地激发。无论如何，要\Quote{寻求}\Quote{生命，寻求幸福的日子}；但要到它们能被找到的地方去寻求。\par

6. 你们是否愿意和我一起聆听那知道\Quote{幸福的日子}和\Quote{生命}何在者的劝告？不要只听我讲，而要与我一同聆听。因为有一位对我们说：\BibleQuote{孩子们，来，听从我。}让我们一同奔跑，站立，竖起耳朵，用心理解那位圣父，他曾说：\BibleQuote{孩子们，来，听从我，我要教导你们敬畏上主。}接着就是他愿意教导我们的内容，以及敬畏上主有何益处。\BibleQuote{谁愿生命长久，渴望看到幸福的日子？}我们众人都回答：\Quote{我们愿意。}那么，让我们听从下文：\BibleQuote{要谨守口舌，不说恶言，嘴唇不吐欺诈。}现在请说：\Quote{我愿意。}刚才我说：\BibleQuote{谁愿生命长久，渴望看到幸福的日子？}我们众人都回答：\Quote{我。}那么来吧，现在让某个人回答\Quote{我}。所以，\BibleQuote{要谨守口舌，不说恶言，嘴唇不吐欺诈。}现在说：\Quote{我。}你们难道想要\Quote{幸福的日子}和\Quote{生命}，却不愿\Quote{谨守口舌，不说恶言，嘴唇不吐欺诈}吗？对奖赏敏感，对工作迟缓！如果不工作，奖赏会给谁呢？我希望在你们家中，你们甚至会奖赏那工作的人！因为对于不工作的人，我确信你们不会付给他报酬。为什么呢？因为你们不欠不工作的人什么。天主也提出了一项奖赏。什么奖赏？\Quote{生命和幸福的日子}，这生命是我们众人所渴望的，这些日子是我们众人努力要达到的。祂会赐予我们应许的奖赏。什么奖赏？\Quote{生命和幸福的日子。}什么是\Quote{幸福的日子}？就是无尽的生命，无忧的安息。\par

7. 祂为我们设置了何等伟大的赏报：在这般伟大的赏报面前，让我们看看祂命令了我们什么。因为受到如此伟大恩许的赏报和对赏报的爱的点燃，让我们立刻准备好我们的力量、我们的腰部、我们的手臂，去执行祂的命令。难道祂命令我们背负重担，或许去挖些什么，或者举起某种器械吗？不，祂并没有吩咐你任何这类劳苦的事，祂只吩咐你\Quote{克制}那个在你所有肢体中你动得如此迅速的肢体。\BibleQuote{使你的口舌远离邪恶。}建造房屋不费劳力，而克制口舌却是劳力吗？\BibleQuote{使你的口舌远离邪恶。}不说谎言，不出口辱骂，不散播诽谤，不作伪证，不说亵渎的话。\BibleQuote{使你的口舌远离邪恶。}你且看自己何等愤怒，若有人毁谤你。正如你对别人愤怒，当他毁谤你时；同样，当你毁谤别人时，也要对自己发怒。\BibleQuote{使你的嘴唇不出欺诈。}你内心有什么，就说出什么。不要让胸膛隐藏一事，而舌头说出另一事。\BibleQuote{远离罪恶，努力行善。}因为我怎能对那至今仍剥去已穿着者衣服的人说\Quote{给赤身者穿衣}呢？因为压迫同乡的人，怎能接纳陌生人呢？所以，按正确次序，首先\Quote{远离邪恶}，然后\Quote{行善}；首先\Quote{束紧你的腰}，然后\Quote{点灯}。当你做完这些，就在确定的希望中等待\BibleQuote{生命和美好日子}。\BibleQuote{寻求和平，追随陪伴。}然后你就能坦然向主说：\Quote{我已做了你所命令的，请赏赐你所恩许的。}\par

\SourceRecord{Translated by R.G. MacMullen.；Nicene and Post-Nicene Fathers, First Series；Vol. 6.；Edited by Philip Schaff.；Buffalo, NY: Christian Literature Publishing Co.,；1888.；<http://www.newadvent.org/fathers/160358.htm>.}

% structure: p0001 source=h1 target=section reason=page-title

\section{新约讲道集 第59篇}

\EditorialNote{[CIX. Ben.]}\par

\Emph{论福音的话：\BibleRef{路 12:56-58}，\BibleQuote{你们知道怎样分辨天地的面色}等；以及论\Quote{当你与你的对头去见官长时，在路上要尽力摆脱他}等话。}\par

1. 我们听到了福音，主在其中斥责那些知道分辨天色，却不知道发现信德的时机、临近的天国的人。这话祂是对犹太人说的，但祂的话也传到了我们这里。如今主耶稣基督亲自开始宣讲祂的福音说：\BibleQuote{你们悔改吧，因为天国临近了。}同样，圣若翰洗者，祂的前驱，也如此开始：\BibleQuote{你们悔改吧，因为天国临近了。}而现在主斥责那些在\Quote{天国临近}时仍不愿悔改的人。\Quote{天国}，正如祂自己所说，\BibleQuote{不是以观察而来的。}祂又说：\BibleQuote{天国就在你们中间。}那么，愿每一个人明智地接受师傅的训诫，以免错失救主仁慈的时机，这仁慈如今正被分施，只要人类还被宽容。因为人被宽容正是为了使他悔改，免得他被定罪。只有天主知道世界的终结何时到来，但如今是信德的时候。世界的终结是否会在这里找到我们任何人，我不知道，也许不会找到我们。我们每个人的时间都非常接近，因为我们是必死的。我们在偶然中行走。如果我们是用玻璃做的，我们害怕的意外会比现在少。有什么比玻璃器皿更脆弱呢？然而它被保存，并持续数世纪。因为虽然人们担心玻璃器皿会摔倒，但它不怕发烧或衰老。我们则更脆弱、更软弱，因为我们在人的事务中不断发生的所有意外，无疑因我们的软弱而每日恐惧；即便这些意外不来，时间也在流逝；人避开了这一打击，就能避开他的终结吗？他避开了从外而来的意外，但内在生出的东西能驱散吗？再者，内脏会长虫，有时某种疾病突然袭来；最后，即使人长久被宽容，当老年到来时，也无法推迟。\par

2. 因此，让我们听从主，让我们在自己内实践祂所命令的。让我们看看那对头是谁，祂警告我们说：\Quote{\BibleQuote{当你与你的对头去见官长时，在路上要尽力摆脱他，免得对头把你交给官长，官长把你交给狱警，你被投入监狱，从那里你不得出来，直到你还清最后一分钱。}}这对头是谁？如果是魔鬼，我们已经被救赎脱离了他。我们为他付出了何等代价，才从他手中被赎出！在这方面，宗徒谈到我们的救赎说：\BibleQuote{祂拯救了我们脱离黑暗的权势，将我们迁到祂爱子的国里。}我们已被救赎，我们弃绝了魔鬼，我们怎能\Quote{尽力摆脱他}，免得他使我们这些罪人再次成为他的俘虏呢？但这不是主警告我们的\Quote{对头}。因为在另一处，另一位圣史这样表达，如果我们将两处结合起来，并比较两位圣史的说法，我们很快就会明白这个对头是谁。请看，路加在这里说了什么？\Quote{\BibleQuote{当你与你的对头去见官长时，在路上要尽力摆脱他。}}但另一位圣史却这样表述同一件事：\Quote{趁你和你的对头还在路上时，赶快与他和解。}其余的都相似：\Quote{免得对头把你交给审判官，审判官交给狱警，你被投入监狱。}两位圣史都这样解释了。一位说：\Quote{在路上要尽力摆脱他，}另一位说：\Quote{与他和解。}因为你不能\Quote{摆脱他}，除非你\Quote{与他和解}。你想\Quote{摆脱他}吗？\Quote{与他和解。}但这是什么？难道基督徒应当与魔鬼\Quote{和解}吗？\par

3. 那么，让我们找出这个\Quote{对头}，我们应当与他\Quote{和解，免得他把我们交给判官，判官交给差役}；让我们找出他，\Quote{并与他和解}。如果你犯罪，天主的圣言就是你的对头。例如，你或许喜好酗酒；它向你说：\Quote{不可这样做。}你好频繁于戏院和这类无聊之事；它向你说：\Quote{不可这样做。}你好行奸淫；天主的圣言向你说：\Quote{不可这样做。}无论你在什么罪中想要按自己的意愿行事，它都向你说：\Quote{不可这样做。}它是你意愿的对头，直到它成为你救恩的创始者。哦，这\Quote{对头}多么美好，多么有益！它不寻求我们的意愿，而是寻求我们的益处。它是我们的\Quote{对头}，只要我们是我们自己的对头。只要你是你自己的仇敌，你就有天主的圣言作你的仇敌；成为你自己的朋友，你就与它\Quote{和解}了。\Quote{不可杀人；}听从，你就与它\Quote{和解}了。\Quote{不可偷盗；}听从，你就与它\Quote{和解}了。\Quote{不可奸淫；}听从，你就与它\Quote{和解}了。\Quote{不可作假见证；}听从，你就与它\Quote{和解}了。\Quote{不可贪恋你近人的妻子；}听从，你就与它和解了。\Quote{不可贪恋你近人的财物；}听从，你就与它\Quote{和解}了。在这一切事上，你已与这个\Quote{你的对头}和解了，你对自己损失了什么？你不仅一无所失，反而找回了先前迷失的自己。\Quote{道路}就是今生；如果我们\Quote{与对头和解}，如果我们与他达成协议；当\Quote{道路}终结时，我们就不必惧怕判官、差役和监狱。\par

4. \Quote{道路}何时终结？并非在同一时刻终结于众人。各人有他自己的时刻，结束他的\Quote{道路}。此生被称为\Quote{道路}；当你结束此生时，你就结束了\Quote{道路}。我们正在前行，生活本身就在前进。除非你们或许以为时间在前进，而我们静止不动！那是不可能的。时间在前进，我们也在前进；年岁并非向我们走来，而是离我们而去。人们大错特错，当他们说：\Quote{这孩子尚少智慧，但年岁将会来到他身上，他将会明智。}思想你所说的话：\Quote{将会来到他身上}，你曾说过；\Quote{我将证明它们离去}，而你说\Quote{它们来到}。请听我多么容易证明它。让我们假设我们从他的出生就已知道他的年数；例如（愿他平安），他必须活八十年，他将到达老年。写下八十年。他已活了一年；总数中还有多少？你写下了多少？八十年！减去一年。他已活了十年；还有七十年。他已活了二十年；还有六十年。然而，定然有人会说，年岁确实来了；这是什么意思？我们的年来到是为了离去；它们到来，我说，是为了离去。因为它们的到来，不是要留在我们身边，而是当它们经过我们时，磨蚀我们，使我们越来越不强壮。这就是我们所进入的\Quote{道路}。那么，我们与那个\Quote{对头}，即天主的圣言，有什么关系呢？\Quote{与他和解}。因为你不知道\Quote{道路}何时可能终结。当\Quote{道路}终结时，留下的有\Quote{判官}、\Quote{差役}和\Quote{监狱}。但如果你对你的\Quote{对头}保持善意，并\Quote{与他和解}；你将发现，取代\Quote{判官}的是一位父亲，取代残酷的\Quote{差役}的是一位天使，带你去到亚巴郎的怀抱，取代\Quote{监狱}的是乐园。你在\Quote{道路上}何等迅速地改变了一切，因为你已\Quote{与你的对头和解}！\par

\SourceRecord{Translated by R.G. MacMullen.；Nicene and Post-Nicene Fathers, First Series；Vol. 6.；Edited by Philip Schaff.；Buffalo, NY: Christian Literature Publishing Co.,；1888.；<http://www.newadvent.org/fathers/160359.htm>.}

% structure: p0001 source=h1 target=section reason=page-title

\section{论新约讲道六十}

[CX. Ben.]\par

\Emph{关于福音的言语}，\BibleRef{路 13:6} \Emph{，那里我们被告知关于无花果树，三年不结果子；以及关于患病十八年的妇人；以及关于}\Emph{\BibleRef{咏 9:19}} \Emph{，\BibleQuote{上主，起来，不要让世人得胜；愿万民在你面前受审判。}}\par

1. 论及\Quote{无花果树}受了三年试验而不结果子，以及\Quote{患病十八年的妇人}，请听主赐我所言。无花果树象征着人类。三年代表三个时期：一个是法律之前，一个是法律之下，一个是恩宠之下。将\Quote{无花果树}理解为人类并无不妥。因为第一个人犯罪时，他用无花果树叶遮盖了自己的裸体；遮盖了那使我们出生的肢体。因为他犯罪前本应荣耀之处，犯罪后却成了羞耻。因此之前，\Quote{他们赤身露体，并不觉得羞耻。}因为他们没有理由脸红，因为没有罪在先；他们也不能为造物主的作品而脸红，因为他们还没有在造物主的善工中掺杂自己任何恶行。因为他们还没有吃那棵不可吃的知善恶树的果子。在他们吃了并犯罪之后，人类便从他们而生；即人从人生，负债者从负债者生，有死之人从有死之人生，罪人从罪人生。因此，在这\Quote{树}上，他指那些在整个时代范围内不结果子的人；为此，斧子正悬在不结果子的树上。园丁为它求情，刑罚被推迟，以便给予帮助。那求情的园丁，就是在教会内为教会外的人祈祷的每一位圣人。他祈祷什么？\Quote{主，容它这一年也留着；}即，在这恩宠时期，宽恕罪人，宽恕不信者，宽恕不结果者，宽恕不结果子者。\Quote{我要给它松土，加上粪肥；如果结果子，便好；否则，你就要来砍掉它。}\Quote{你就要来：}何时？你将在审判中来，当你来审判生者死者时。此时他们被宽恕。但\Quote{松土}是什么？什么是\Quote{给它松土}，岂不就是教导谦卑与悔改？因为沟是低地。粪肥要理解其有利效果。它肮脏，却结出果实。园丁的污秽就是罪人的忧苦。那些悔改的人，在污秽的袍子中悔改；也就是说，如果他们正确理解，并真诚悔改。对这棵树说，\BibleQuote{你们悔改吧，因为天国临近了。}\par

2. 那\Quote{患病十八年的妇人}是什么？六天中天主完成了祂的工程。三乘以六是十八。那么，\Quote{树}中\Quote{三年}所象征的，在这妇人中就是\Quote{十八年}。她弯腰不能抬头，因为她徒然地听了\Quote{请举心向上}。但主使她挺直。那么，就有希望，对于人，直到审判之日来临。人很自大。然而，人算什么？一个义人是伟大的。但义人只靠天主的恩宠为义。\BibleQuote{人算什么，竟使你怀念他？}你要看人是什么？\BibleQuote{众人都是说谎者。}我们曾咏唱：\BibleQuote{上主，起来，不要让世人得胜。}什么是\Quote{不要让世人得胜}？宗徒们不是人吗？殉道者不是人吗？主耶稣自己，在不停止是天主的同时，也屈尊成为人。那么，\Quote{上主，起来，不要让世人得胜}是什么意思？如果\Quote{众人都是说谎者；起来，}真理，\Quote{不要让}谎言\Quote{得胜。}因此，人若想成为任何善，决不可出于自己的任何东西。因为他若想凭自己成为什么，就会是\Quote{说谎者。}他若想成为真实的，就必须出于那从天主而来的，而非出于自己的任何东西。\par

3. 因此，\BibleQuote{主，起来，让人不要得胜。} 在洪水之前，谎言如此盛行，以致洪水后只留下了八个人。这些人又使大地充满说谎的人，从他们中选出了天主的子民。行了许多奇迹，赐予了神圣的恩惠。他们被带领直抵许诺之地，从埃及的奴役中被解救出来：在他们中兴起先知，他们领受了圣殿，领受了司祭职，领受了傅油，领受了法律。然而，关于这同一子民，后来有话说：\BibleQuote{外邦的子女向我说了谎。} 最终，那位由先知们预许的被派遣来了。\BibleQuote{让人不要得胜，} 甚至更甚，因为天主成了人。但即使是他，虽然行了神圣的作为，却被轻视，尽管他显示了如此多的仁慈之举，他被捕，被鞭打，被悬挂。至此\Quote{人得了胜，} 以至于逮捕了天主子，鞭打了天主子，给天主子戴上荆棘冠冕，将天主子挂在树上。至此\Quote{人得了胜：} 但到什么地步呢？不过是直到他从树上被取下，被安放在坟墓里？如果他留在那里，人就真的\Quote{得胜}了。但这预言正是对主耶稣本人说的：\BibleQuote{主，起来，让人不要得胜。} 主，你屈尊降来成为血肉，圣言成了血肉。圣言在我们之上，血肉在我们之中，圣言血肉介于天主与人之间：你选择了一位童贞女，按血肉而诞生，当你将被孕育时，你找到了一位童贞女；当你诞生时，你仍保留了一位童贞女。但你未被承认；你被看见，却隐藏着。软弱被看见，大能被隐藏。这一切都是为了让你流下那血，那是我们的代价。你行了如此伟大的奇迹，治愈了病弱者的软弱，显示了如此多的仁慈，却以恶报善。他们嘲弄你，你挂在树上；不敬虔的人在你面前摇头，说：\BibleQuote{你若是天主子，就从十字架上下来吧。} 难道你那时失去了你的大能，或者你是在显示你的忍耐？然而他们嘲弄你，他们讥讽你，当你被杀害时，他们好像得胜者离开了。看，你被安放在坟墓里：\BibleQuote{主，起来，让人不要得胜。} \Quote{不要让}不敬虔的敌人\Quote{得胜，不要让}盲目的犹太人\Quote{得胜。}因为当你被钉十字架时，犹太人自以为是地\Quote{得了胜}。\BibleQuote{主，起来，让人不要得胜。} 这事已成就，是的，已成就。如今还剩下什么？岂不是\BibleQuote{万民要在你面前受审判}？因为他已复活，如你们所知，并升了天；从那里他将要来审判生者死者。\par

4. 啊！不结果实的树，不要讥讽，因为你尚且被宽恕；斧头被延迟，不要自以为安全；他必来，你将被砍倒。相信他必来。你们如今所见的一切，曾经并不存在。曾经，基督徒民族并不遍布全世界。这写在预言中，而非见于地上；如今既被读到，也被看到。同样，教会本身也得以完成。并非对她说：\BibleQuote{女儿，请看而听；} 而是：\BibleQuote{请听而看。} 听预言，看应验。因此，我亲爱的弟兄们，基督曾未从童贞女诞生，但他的诞生被预许，于是他诞生了；他曾未行奇迹，它们被预许，于是他行了；他曾未受难，它被预许，于是成就了；他曾未复活，它被预告，于是应验了；他的名未曾遍及世界，它被预告，于是应验了；偶像未曾被摧毁和打破，它被预告，于是应验了；异端者未曾攻击教会，它被预告，于是应验了。同样，审判之日尚未到来，但既然它被预告，必要应验。难道那位在如此多事上显示自己为真实者，会在审判之日上说谎吗？他给了我们他许诺的保证。因为天主使自己成为债务人，并非因欠债，即非因借贷，而是因许诺。所以我们不能对他说：\Quote{归还你所领受的。} 因为\BibleQuote{谁先给了他，而使他偿还呢？} 我们不能对他说：\Quote{归还你所领受的；} 但我们毫无顾虑地说：\Quote{归还你所许诺的。}\par

5. 因此，我们敢天天说：\BibleQuote{愿你的国来临；} 使他的国来临时，我们也能与他一同为王。这已用这些话许诺了我们：\BibleQuote{那时，我要对他们说：你们蒙我父降福的，来吧！承受自创世以来为你们预备的国度。} 但确然，只有我们做了那处随后所说的事才行：\BibleQuote{因为我饿了，你们给了我吃的，} 等等。他向我们的祖先作了这些许诺；但他给了我们一份保证书，也让我们来读。如果那位屈尊赐予我们这保证书的人要与我们算账，并说：\Quote{读我的债务，即我许诺的债务，并计算我已偿付的，也计算我仍欠的；看，我已偿付了多少；我所欠的只是很少；你们会为了那少许的剩余，而认为我是不可信的吗？} 面对这极其明显的真理，我们能回答什么呢？所以，让那不结果实的人悔改，结出\BibleQuote{与悔改相称的果实}。那弯腰弓背、只注目于地、以地上的幸福为乐、认为唯有这是幸福生活、并相信没有另一种生活的人；无论他是什么人，这样弯腰弓背的，让他挺直起来；若不能靠自己，就呼求天主。因为那妇人岂是靠她自己挺直的吗？如果他没有伸出他的手，她就有祸了。\par

\SourceRecord{Translated by R.G. MacMullen.；Nicene and Post-Nicene Fathers, First Series；Vol. 6.；Edited by Philip Schaff.；Buffalo, NY: Christian Literature Publishing Co.,；1888.；<http://www.newadvent.org/fathers/160360.htm>.}

% structure: p0001 source=h1 target=section reason=page-title

\section{新约讲道集 第61篇}

\Emph{[CXI. Ben.]}\par

\Emph{论福音之言，其中记载：\BibleRef{路 13:21}，天主国被比喻为\BibleQuote{酵母，妇人取来藏在三斗面中}；同章又记载：\BibleQuote{主，得救的人真的不多吗？}}\par

1. 主所说的\Quote{三斗面}是指人类。回想洪水时代：仅存三人，人类从他们再次繁衍。诺厄有三个儿子，人类由他们重建。那位\Quote{藏酵母的圣妇}是智慧。看，普世教会都在呼喊：\BibleQuote{我知道上主何其伟大。}然而无疑只有少数人得救。你们还记得前不久福音中提出一个问题：\Quote{主，得救的人真的不多吗？}主如何回答？祂没有说：\Quote{得救的人不多，而是许多。}祂没这样说。当祂听到\Quote{得救的人真的不多吗？}时，祂说：\BibleQuote{你们竭力由窄门进入吧。}所以当你们听到\Quote{得救的人真的不多吗？}主肯定了这问题。通过\Quote{窄门}只有\Quote{少数}能\Quote{进入}。在另一处祂自己说：\BibleQuote{窄而狭的路通向生命，找到它的人很少；宽而大的路通向丧亡，走那条路的人很多。}我们为何因人多而欢喜？请听我说，你们这些\Quote{少数}。我知道听我的人很多，但真正听从的却是\Quote{少数}。我看见打谷场，我寻找麦粒。当打谷场被碾打时，几乎看不见麦粒；但将来必会簸扬。与将丧亡的多数相比，得救的只是少数。因为这群\Quote{少数}本身将形成一大群。当簸扬者手持簸箕来临时，\BibleQuote{祂要清洁祂的禾场，将麦粒收入粮仓，用不灭的火焚烧糠秕。}愿糠秕不要嘲笑麦粒；祂说真理，不欺骗任何人。你们在众人中要成为众多，即使与某些众多相比你们是少数。从这个禾场将产生如此庞大的群体，足以充满天上的粮仓。因为主基督不会自相矛盾，祂曾说：\BibleQuote{许多人由窄门进入，许多人由宽门走向灭亡}；祂在另一处又说：\BibleQuote{许多人要从东方和西方来。}所以\Quote{少数}就是\Quote{许多}；既是\Quote{少数}又是\Quote{许多}。难道\Quote{少数}是一种，\Quote{许多}是另一种？不。而是\Quote{少数}本身即\Quote{许多}；与丧亡者相比是\Quote{少数}，与天使的共融相比是\Quote{许多}。请听，亲爱的诸位。\WorkTitle{默示录}记载：\BibleQuote{此后我看见各民族、各邦国、各语言的一大队人，没有人能数清，他们身穿白衣手持棕榈枝。}这就是圣徒的群体。当禾场被簸扬后，从不虔敬、邪恶、伪基督徒的群众中分离出来——那些\Quote{拥挤}却不\Quote{触摸}的人（福音中一位妇人\Quote{触摸}基督，而众人\Quote{拥挤}祂）——被投入永火；当所有该受罚的人被分离后，这被净化的群众将站在右边，不再害怕恶人混杂，也不担忧善人失落，即将与基督为王，他们将何等自信地宣告：\BibleQuote{我知道上主何其伟大}！\par

2. 所以，我的弟兄们（我是对麦粒说话），如果你们承认我所说的，被预定得永生的人，要以行为说话，而非以言语。我不得不对你们说些我不该说的话。我本应在你们身上找到可赞美之处，而非寻找劝诫的理由。然而我只说几句，不会停留。要承认好客之德，有人因此达到天主面前。你们接待一位陌生人，你们自己也是旅途中的同伴；因为众人都是旅客。一个基督徒，即使在自己的家中和本国，也承认自己是旅客。因为我们的家乡在天上，在那里我们不再是旅客。因为每个人在世上，即使在自己家中，也是旅客。如果不是旅客，他就不会从此处离开。既然必须离开，他就是旅客。不要自欺，他是旅客；无论愿意与否。他把房子留给儿女，一个旅客留给其他旅客。为什么？如果你在客栈，当另一个人来时，你不也会离开吗？即使在自己家中也是如此。你的父亲把地方留给你，你将来也会留给你的儿女。你在此地不是永久居住，也不会永久留给住在此地的人。既然我们都要离去，让我们做那不能逝去的事，好让我们离去之后，到达那永不再离去的地方时，能在那里找到我们的善工。基督是保管者，你为何害怕因施舍穷人而有所损失？\Quote{让我们归向上主}等等。\par

\Emph{讲道后。}\par

我提醒你们，亲爱的诸位，你们已经知道的事。明天是可敬的奥勒留主教祝圣周年纪念日；他借着我的卑微职务请求并劝告你们，亲爱的弟兄们，请你们虔诚地聚集在福斯多圣殿。感谢天主。\par

\SourceRecord{Translated by R.G. MacMullen.；Nicene and Post-Nicene Fathers, First Series；Vol. 6.；Edited by Philip Schaff.；Buffalo, NY: Christian Literature Publishing Co.,；1888.；<http://www.newadvent.org/fathers/160361.htm>.}

% structure: p0001 source=h1 target=section reason=page-title

\section{新约讲道集第六十二篇}

\Emph{[XCII. 本笃版]}\par

\Emph{论福音之言}，\BibleRef{路 14:16}\Emph{，\BibleQuote{某人设了盛大的晚宴}等。}\par

\EditorialNote{在雷斯蒂图塔圣殿讲道。}\par

1. 神圣的读经已陈设在我们的面前，我们应当侧耳恭听，并且靠着主的帮助，我将作一些论述。在宗徒读经中，为外邦人的信德感谢主，当然，因为这是祂的工程。在圣咏中我们曾诵念：\BibleQuote{万军的天主呀，求祢转变我们，显示祢的圣容，我们就得救。}在福音中我们被召赴宴；是的，不如说别人被召，我们未被召，而是被引导；不仅被引导，甚至被强迫。因为我们听到：\BibleQuote{某人设了盛大的晚宴。}这人是\BibleQuote{天主与人之间的中保，那人是基督耶稣}。祂打发人去，使被邀请的人来，因为时候已经来到，他们应当来。谁是被邀请的人？岂不是那些先前被派遣的先知所召唤的人？何时？古时，自从先知被派遣以来，他们就邀请人赴基督的晚宴。他们被派遣到以色列百姓那里。他们屡次被派遣，屡次召唤人，叫他们在晚宴的时刻来。但他们接纳了邀请他们的人，却拒绝了晚宴。\Quote{他们接纳了邀请他们的人，却拒绝了晚宴}是什么意思？他们读了先知，却杀害了基督。但当他们杀害祂时，虽然他们不知道，却为我们预备了晚宴。当晚宴已经预备好，当基督已被献祭，当主的晚宴——就是信徒所知道的——在基督复活后设立，并由祂的手和嘴建立时，宗徒被派遣到他们那里，就是先前先知被派遣到的人那里。\BibleQuote{你们来赴晚宴吧。}\par

2. 他们不愿来，便推辞。他们如何推辞？有三种推辞：\BibleQuote{一个说：我买了一块田地，必须去看看；请准我辞。另一个说：我买了五对牛，要去试一试；请准我辞。第三个说：我娶了妻子，请准我辞；我不能来。}我们猜想这些岂不是阻挡所有人、使他们不肯赴宴的推辞？让我们审视、讨论、找出它们；但只是为叫我们警惕。在买田地的事上，标记出统治的灵；因此骄傲受到谴责。因为人们喜爱有田地，去持有、占有它，有手下的人，有统治权。这是恶习，第一个恶习。因为第一个人想统治，就是不愿有人统治他。统治是什么？岂不就是以自己的力量为乐？有一个更大的力量，让我们服从祂，好能得到安全。\BibleQuote{我买了一块田地，请准我辞。}发现了骄傲，他便不肯来。\par

3. \BibleQuote{另一个说：我买了五对牛。}说\Quote{我买了牛}不就行了吗？无疑有些东西，因其隐晦而促使我们寻求并理解；正因其封闭，祂劝勉我们叩门。五对牛就是这身体的感官。这身体有五种感官，这是众所周知的；那些或许不曾留意的人，一经提醒，定然会明白。于是这身体有五种感官。视觉在眼睛，听觉在耳朵，嗅觉在鼻子，味觉在口，触觉在全身。我们通过视觉感知白和黑，以及各种颜色、明暗；通过听觉感知刺耳和悦耳的声音；通过嗅觉感知香臭；通过味觉感知甜苦；通过触觉感知软硬、光滑粗糙、冷热、轻重。它们是五种，且是成对的。它们成对这一点，在最先三种感官中最易见：有两只眼睛、两只耳朵、两个鼻孔——这是三对。在口中，即味觉中，某种成对性可见，因为除非触到舌头和上颚，没有东西影响味觉。与触觉有关的肉体的快乐，其成对性较不明显。因为既有外在的触觉，也有内在的触觉。因此它也是成对的。为何称为\Quote{牛}？因为这身体的感官寻求的是地上的事物。牛翻动土地。所以有些远离信德的人，沉溺于世俗之事，忙于肉体之事；他们除藉身体的五种感官所达到的以外，什么都不信。他们将整个意志的规则建立在那五种感官上。\Quote{我什么都不信，}一个说，\Quote{除非我看见。看，这是我所知并确信的。某物是白或黑，或圆或方，或某种颜色；我知道，感觉到，把握住了；自然本身教导我。我不必被迫相信你不能向我显示的东西。或者那是一种声音：我感知那是声音；它唱得好，唱得坏，甜美或刺耳。我知道，我知道，它到我这里来了。有好的或坏的气味：我知道，我闻到。这是甜，这是苦；这是咸，这是淡。我不知道你还要告诉我什么。通过触觉我知道什么是硬，什么是软；什么是光滑，什么是粗糙；什么是热，什么是冷。你还要向我显示什么？}\par

4. 因这阻隔，我们的宗徒\Person{多默}受了牵制；他对主基督、即基督的复活，竟连自己的眼睛也不肯相信。\BibleQuote{除非，}他说，\BibleQuote{我把我的指头放在钉痕和伤处，除非我把我的手探入祂的肋膀，我决不信。}主本来可以毫无伤痕地复活，却保留了疤痕，让那怀疑的宗徒触摸，好治愈他内心的创伤。然而，祂有意召叫其他人赴祂的宴席，针对\Quote{五对牛}的托辞，祂说：\BibleQuote{那些没有看见而相信的，才是有福的。}我们，我的弟兄们，已被召叫赴这宴席，没有被这\Quote{五对}所阻。因为我们没有在今世渴望看见主的身体的容颜，也没有渴望听见从那身体口中发出的声音；我们没有在祂身上寻求任何转瞬即逝的香气。有一个\BibleQuote{妇人用极贵重的香液傅抹了祂，}那\BibleQuote{屋子充满了香气；}但我们不在那里；看，我们没有闻到，却相信了。祂亲手祝圣了晚餐给门徒；但我们没有参加那次筵席，却每天凭信德吃这同一晚餐。不要以为奇怪：在祂亲手赐予的那晚餐中，有一个没有信德的人在场；后来显现的信德远远弥补了那时的无信。\Person{保禄}不在那里，他却相信；\Person{犹达斯}在那里，他背叛了。如今在这同一晚餐中，虽然后来他们没有看见那张桌子，也没有用眼睛看见、用口品尝主拿在手里的饼，但因为这与现在预备的是同一的，如今在这同一晚餐中，有多少人\BibleQuote{吃喝自己的判决}呢？\par

5. 然而，是什么机会让主说起这晚餐呢？与祂同席的人中（因为祂应邀赴宴），有一个说：\BibleQuote{在{\Person{天主}}的国里吃饭的人是有福的。}他仿佛叹息遥远的事，而那位\Quote{食粮}本身正坐在他面前。\Quote{\Person{天主}}国的食粮是谁？不就是那说\BibleQuote{我是从天上降下的生活的食粮}的那一位吗？不要预备你的口，而要预备你的心。就在这机会上，祂设了这晚餐的比喻。看，我们信基督，我们以信德接受祂。接受祂时，我们知道该想什么。我们领受得少，却在心里得到滋养。所以，喂养我们的不是所见，而是所信。因此，我们也没有寻求那种外在的感觉；也没有说：\Quote{让那些在祂复活后亲眼看见、亲手摸过主本身的人相信吧，如果所说的是真的；我们摸不到祂，为什么要信？}如果我们这样想，就会被那\Quote{五对牛}阻隔，不能赴宴。为使你们知道，弟兄们，祂所指的并非那五感官的快感，即舒适和享乐，而是一种好奇心，祂没有说：\Quote{我买了五对牛，要去喂养它们}；却说：\Quote{我去试它们。}那想用\Quote{牛对}去\Quote{试}的人，不愿存疑，正如圣\Person{多默}用这些\Quote{对}不愿存疑。\BibleQuote{让我看，让我摸，让我伸指头进去。}\Quote{看，}主说，\BibleQuote{伸你的指头沿着我的肋旁，不要作无信的人。为了你的缘故，我被杀；在你想要触摸的地方，我流了我的血，为救赎你；你还怀疑我，除非你摸我？看，这一点我也允许；看，这一点我也向你显示；摸吧，相信吧；找到我伤处的地方，治愈你疑虑的创伤。}\par

6. \Quote{第三个说：我娶了妻子。}这是肉身的享乐，阻碍了许多人；我但愿这只是在外面，而不是在里面！有的人说：\Quote{人若不享有肉身的快乐，便无幸福。}这些人正是宗徒所斥责的，\Person{保禄}说：\BibleQuote{我们吃喝吧，因为明天我们就要死了。}谁从那边复活到今生来？谁曾告诉我们那边的事？我们带走的，只是今世使我们快乐的东西。这样说的人，\Quote{娶了妻子}，依附肉身，以肉身的快乐为乐，从晚餐中推辞；他要小心，免得死于内在的饥饿。要听圣\Person{若望}宗徒兼福音作者的话：\BibleQuote{不要爱世界，也不要爱世界上的事。}哦，你们这些来赴主晚餐的人，\BibleQuote{不要爱世界，也不要爱世界上的事。}他没有说：\Quote{不要有；}而是说：\Quote{不要爱。}你曾有过、拥有过、爱过。对世俗事物的爱，是灵魂翅膀的粘鸟胶。看，你曾贪求，你已粘住。\Quote{谁给你翅膀如鸽子？}你几时才能飞向那当去之地？因你乖谬地愿在此安息，反而有害地粘住了。\BibleQuote{不要爱世界，}这是天主的号角。这号角的声音不断地向大地四周、向全世界宣告：\BibleQuote{不要爱世界，也不要爱世界上的事。凡爱世界的，父的爱就不在他内。因为凡在世界上的，就是肉身的贪欲、眼目的贪欲和生活的奢傲。}祂从福音结束的地方开始。祂从福音终止的地方开始。\Quote{肉身的贪欲：我娶了妻子。眼目的贪欲：我买了五对牛。生活的奢傲：我买了一块田。}\par

7. 现在，这些感官仅通过提及眼睛来表达，以部分代表整体，因为在五种感官中，眼睛居首位。因此，虽然视觉专属眼睛，但习惯上我们通过所有五种感官使用\Quote{看}这个词。如何？首先，就眼睛本身而言，我们说：\Quote{看它有多白，看并看它有多白：}这与眼睛有关。听并看它有多悦耳！反过来，我们能说：\Quote{听并看它有多白}吗？这个表达\Quote{看}贯穿所有感官；而其他感官的独特表达却不能反过来贯穿它。\Quote{注意并看它有多悦耳；闻并看它有多宜人；尝并看它有多甜；触并看它有多柔软。}然而，既然它们都是感官，我们倒不如这样说：\Quote{听并感知它有多悦耳；闻并感知它有多宜人；尝并感知它有多甜；触并感知它有多热；抚摸并感知它有多光滑；抚摸并感知它有多柔软。}但我们没说这些。因为主自己复活后向门徒显现时，虽然他们看见祂，但仍怀疑信德，以为看见的是神灵，祂说：\Quote{你们为什么怀疑，为什么心里起疑虑？看我的手和我的脚。}只说\Quote{看}还不够，祂说：\Quote{触摸，抚摸，并看。}\Quote{看并看，抚摸并看；仅用眼睛看，并通过所有感官看。}因为祂寻找内在的信德感官，便将祂自己献给身体的外在感官。通过这些外在感官，我们在主内一无所获；我们用耳朵听见，用心相信；而这听见不是来自祂的口，而是来自祂宣讲者的口，来自那些已经在晚宴上的他们，他们倾倒出他们在那里饱饮的，邀请了我们。\par

8. 那么，让我们抛弃虚妄而邪恶的借口，来到这能使我们内在饱足的晚宴。不要让骄傲的膨胀阻止我们，不要让它抬举我们，也不要让非法的好奇惊吓我们，使我们远离天主；不要让肉体的欢愉阻碍心灵的欢愉。让我们来，并得饱足。谁来了呢？岂不是乞丐、\Quote{残缺的}、\Quote{瘸腿的}、\Quote{瞎眼的}吗？但富人和健全的人没有来，他们自以为行走端正，目光敏锐；他们对自己大有信心，因此越是骄傲，就越处于绝望的境地。让乞丐来吧，因为祂邀请他们，\Quote{祂本是富有的，为了我们却成了贫穷的，好使我们这些乞丐因祂的贫穷而成为富足的。}让残缺的人来吧，\Quote{因为健康的人不需要医生，而是有病的人。}让瘸腿的人来吧，他们可以对祂说：\Quote{请在我的道路上安排我的脚步。}让瞎眼的人来吧，他们可以说：\Quote{照亮我的眼睛，免得我沉睡至死。}这些人来时，那些最初被邀请的人因自己的借口而被摒弃了；他们来时，从城中的大街小巷进来。那被派去的仆人回报说：\Quote{主，已照祢所吩咐的做了，可是还有空位。}祂说：\Quote{出去，到大道和篱笆那里，勉强你所遇见的人进来。}你所遇见的人，不要等他们自愿来，勉强他们进来。我准备了盛大的晚宴，宽大的房屋，我不能让其中任何地方空着。外邦人从大街小巷来了；让异端者从篱笆那里来，在这里他们将找到平安。因为那些制作篱笆的人，目的是制造分裂。让他们从篱笆中被拉出来，从荆棘中被拔出。他们卡在篱笆里，不愿被勉强。他们说，我们自愿进来。这不是主的命令，\Quote{勉强他们，}祂说，\Quote{进来。}让勉强在外面找到，意志在里面兴起。\par

\SourceRecord{Translated by R.G. MacMullen.；Nicene and Post-Nicene Fathers, First Series；Vol. 6.；Edited by Philip Schaff.；Buffalo, NY: Christian Literature Publishing Co.,；1888.；<http://www.newadvent.org/fathers/160362.htm>.}

% structure: p0001 source=h1 target=section reason=page-title

\section{新约讲道集第六十三篇}

[CXIII. Ben.]\par

\Emph{论福音中的话，\BibleRef{路 16:9}，\BibleQuote{你们要用不义的钱财结交朋友}等等。}\par

我们的本分是向他人施予我们自己所接受的劝诫。最近的福音教训劝诫我们要用不义的钱财结交朋友，好叫他们也接那这样行的人进入永久的帐幕。但谁将有永久的帐幕呢？岂非天主的圣人？谁将被他们接入永久的帐幕呢？岂非那些供应他们需要、并欣然施舍他们必需的人？因此，我们要记住，在最后的审判中，主将对站在祂右边的人说：\BibleQuote{我饿了，你们给了我吃的；}以及其余你们所知道的。当他们问何时向他们行了这些善事时，祂回答说：\BibleQuote{你们对我这些最小兄弟中的一个所做的，就是对我做的。}这些最小的人就是那些接人进入永久帐幕的人。祂对右边的人这样说，因为他们行了这些事；对左边的人则说相反的话，因为他们不肯行。但那行了这些事的右边的人，他们得了什么，或者，他们将要得什么？祂说：\BibleQuote{来吧，我父所祝福的，承受自创世以来为你们预备的国度吧。因为我饿了，你们给了我吃的。你们对我这些最小兄弟中的一个所做的，就是对我做的。}那么，基督的这些最小的人是谁呢？他们是那些舍弃一切所有、跟随祂的人，将他们所有的分施给穷人；好使他们毫无牵挂、不受任何世俗束缚地事奉天主，卸下世界的重担，展翅高飞。这些就是最小的。为何最小？因为谦卑，因为不张狂，不骄傲。但若称量这些最小的，你会发现他们分量极重。\par

但祂说他们是\Quote{不义钱财的朋友}是什么意思？什么是\Quote{不义钱财}？首先，什么是\OriginalTerm{玛门}{mammon}？因为它不是拉丁词。它是个希伯来词，与布匿语同源。这两种语言因意义相近而彼此关联。布匿语称为玛门，拉丁语称为\OriginalTerm{利益}{lucre}。希伯来语称为玛门，拉丁语称为\OriginalTerm{财富}{riches}。因此，我们用拉丁语表达全部，我们的主耶稣基督这样说：\BibleQuote{你们要用不义的钱财结交朋友。}有些人错误理解这点，掠夺他人财物，施舍一部分给穷人，以为这样就履行了诫命。他们主张：掠夺他人财物是不义钱财；将其中一部分，尤其是施舍给圣徒穷人，这就是用不义钱财结交朋友。这种理解必须纠正，是的，必须从你的心版上彻底抹去。我不希望你们这样理解。从你们正义的劳动中施舍；从你们合法拥有的财物中施舍。因为你们不能贿赂基督，你们的审判者，使祂不听你们与你们所抢夺的穷人一同的控告。因为如果你掠夺一个弱小者，你更强壮更有势力，他与你同到审判者面前，任何你所喜欢的世上审判者，他有权审判，并要与你对质；如果你从那个穷人的掠夺物中给审判者一些东西，让他判你胜诉，那个法官会令你满意吗？诚然，他判了你胜诉，然而正义的力量如此之大，连你也会对他不满。因此，不要这样想象天主。不要在你的心殿中设立这样的偶像。你的天主不是那样，而你自己也不应当那样。如果你自己不会这样审判，而是公正审判；同样，你的天主比你更美善：祂不比你低劣，祂更公正，祂是正义的源泉。凡你行的善事，都是从祂得来的；凡你倾注的善言，都是从祂汲取的。你夸赞器皿，因它从祂有所得，却责备源泉？不要用高利贷和利息来施舍。我对信友说话，对受我们分施基督圣体的人说话。你们当畏惧并改正自己：免得我以后要说，你们这样做，你们也这样做。但我认为，如果我这样说，你们不该向我发怒，而应向自己发怒，好能改正。因为这就是圣咏中那句话的意思：\BibleQuote{你们当发怒，却不要犯罪。}我要你们发怒，但只为了让你们不犯罪。现在，为了不犯罪，你们应该向谁发怒，除了向自己？什么是悔罪的人，不就是对自己发怒的人吗？为了获得宽恕，他对自己施加惩罚；因此他有权向天主说：\BibleQuote{求祢转面不顾我的罪，因为我承认我的罪。}如果你承认，祂就会赦免。你们这些做错事的人，不要再做了：这是不合法的。\par

3. 但如果你已经这样做了，并且拥有这样的金钱，从而装满了你的钱柜，并以此积聚珍宝：你所拥有的来自邪恶，那么现在不要再增添邪恶，要藉着\OriginalTerm{不义的钱财}{mammon of iniquity}为自己结交朋友。\Person{匝凯}所拥有的来自好源头吗？读一读便知。他是税吏长，也就是说，是公共税金缴付给他的人：他因此获得财富。他压迫了许多人，从许多人那里索取，从而积聚了大量财富。基督进入了\Person{匝凯}的家，救恩降临到他的家；因为主亲自说：\BibleQuote{今天救恩临到了这一家}。现在请注意这救恩的方式。首先他切望看见主，因为他身材矮小；但当人群阻挡他时，他爬上了一棵野桑树，看见主经过。但\Person{耶稣}看见了他，说：\BibleQuote{\Person{匝凯}，快下来，今天我必须住在你家里}。你挂在那里，但我不会让你悬着。也就是说，我不会推迟。你想看我经过，今天你将发现我住在你家里。于是主进入了\Person{匝凯}的家，他满心欢喜地说：\BibleQuote{我把我财产的一半分给穷人}。看，他跑得多快，他跑着去用不义的钱财结交朋友。而且唯恐在其他方面被定罪，他说：\BibleQuote{如果我从任何人那里拿了什么，我要四倍偿还}。他给自己判了罪，以免遭受定罪。所以，你们这些从邪恶来源拥有任何东西的人，要用它行善。你们这些没有的人，不要渴望以邪恶手段获得。你自己要成为善人，用邪恶所得行善：并且当你开始用这邪恶行善时，不要自己仍留在邪恶中。你的钱正在转化为善，而你自己却继续作恶吗？\par

4. 确实还有另一种理解方式；我也不隐瞒它。\OriginalTerm{不义的钱财}{mammon of iniquity}就是这世界的一切财富，无论其来源如何。因为无论它们如何积聚，都是不义的钱财，也就是说，是罪恶的财富。\BibleQuote{它们是罪恶的财富}是什么意思？那是罪恶称之为财富的钱财。因为如果我们寻求真正的财富，它们与这些不同。\Person{约伯}在这些财富上富足，尽管他一无所有，但当他向他的天主怀着饱满的心，并像最珍贵的宝石一样向他的天主倾泻赞美时，他已失去了一切。如果他什么都没有，他从什么宝库里这样做呢？这些才是真正的财富。但另一种被罪恶称为财富。你拥有这些财富。我不责备它：你继承了遗产，你父亲富有，留给了你。或者你诚实获得了它们：你家里充满了正义劳动之果；我不责备它。然而即便如此，不要称它们为财富。因为如果你称它们为财富，你会爱它们；如果你爱它们，你会与它们一同灭亡。舍弃，以免被舍弃；施与，以便获得；播种，以便收割。不要称这些为财富，因为它们不是\Quote{真正的}。它们充满贫乏，且容易遭遇意外。这些是哪种财富，为了它们你害怕强盗，为了它们你害怕自己的仆人，怕他杀了你，拿走它们逃跑？如果它们是真正的财富，它们会给你安全。\par

5. 所以，那些才是真正的财富，当我们拥有它们时，不会失去。并且唯恐你因它们而害怕贼，它们会在一个无人能拿走的地方。听听你的主：\BibleQuote{你们要为自己积蓄财宝在天上，那里没有贼接近}。那么当你将它们从这里移去时，它们才会成为财富。只要它们在地上，就不是财富。但世界称它们为财富，罪恶也这样称呼。因此，天主称它们为\OriginalTerm{不义的钱财}{mammon of iniquity}，因为罪恶称它们为财富。听听圣咏：\BibleQuote{上主，求你救我脱离外方人的手，他们的口说出虚妄的话，他们的右手是罪恶的右手。他们的儿子如同新栽的树苗，从幼年就扎根牢固。他们的女儿装饰打扮，如同圣殿般华丽。他们的仓库满溢，流溢到各处。他们的牛犊肥壮，羊群繁殖，不断增多。没有城墙的缺口，没有出征，街上没有哭喊声。} 看，圣咏作者描述了怎样的幸福：但请听他对那些他描述为罪恶之子的人的情况。\BibleQuote{他们的口说出虚妄的话，他们的右手是罪恶的右手。} 他就是这样描述他们，并说他们的幸福只在地上。他还加了什么？\BibleQuote{拥有这些的人民是有福的。} 但谁这样称他们？\BibleQuote{外方人}，与种族无关的异乡人，不属于亚巴郎的后裔：他们\BibleQuote{称拥有这些的人民是有福的}。谁这样称他们？\BibleQuote{那些口出虚妄的人。} 所以，称拥有这些的人为有福是虚妄的事。然而，他们却这样被称，\BibleQuote{那些口出虚妄的人}。福音中不义的钱财被他们称为财富。\par

6. 但你说什么？既然这些\Quote{外人}——他们\Quote{口说虚妄的话}——\Quote{称那拥有这些的百姓为有福}，你说什么？这些都是虚假的财富，请给我看真正的财富。你指摘这些，请给我看你所赞美的。你希望我轻视这些，请给我看你所优先选择的。让圣咏作者亲自发言。因为那说\Quote{他们称那拥有这些的百姓为有福}的那位，给我们这样的回答，好像我们对他说了，也就是说，对圣咏作者本人说：\Quote{看，你从我们这里拿走了这些，什么也没有给我们；看，这些，看，这些我们轻视；我们靠什么生活，我们靠什么成为有福？因为他们说了，他们要亲自为自己回答。因为他们\InnerQuote{称}拥有财富的人\InnerQuote{有福}。但你说什么？}好像他这样被问了，他回答并说：他们称有钱人为有福；但我却说：\BibleQuote{有福的百姓，他们的天主是主。}那么，你已经听到了真正的财富，用不义的钱财结交朋友，你将成为\BibleQuote{有福的百姓，他们的天主是主。}有时我们走在路上，看到非常宜人且多产的庄园，我们就说：\Quote{那是谁的庄园？}我们被告知：\Quote{某人的}；我们便说：\Quote{有福的人！}我们\Quote{说虚妄的话。}他的房屋有福，他的庄园有福，他的羊群有福，他的仆人，他的家眷有福。除去虚妄，如果你愿意听真理。\BibleQuote{他的天主是主的那个人有福。}因为不是拥有那庄园的人有福，而是拥有那位\Quote{天主}的人有福。但为了最清楚地说明财富的幸福，你说你的庄园使你幸福。为什么？因为你靠它生活。因为你极力赞扬你的庄园时，你这样说：\Quote{它给我提供食物，我靠它生活。}请考虑你真正靠什么生活。你靠祂生活的那一位，就是你对祂说：\BibleQuote{在祢那里有生命的泉源}的那位。\BibleQuote{那百姓有福，天主是主。}哦主我的天主，哦主我们的天主，请祢使我们幸福，好让我们来到祢面前。我们不愿从黄金、白银、土地，从这些尘世、最虚妄、易逝的现世财富中得到幸福。愿\Quote{我们的口不说虚妄的话。}请祢使我们幸福，因为我们永远不会失去祢。当我们一旦得到祢，我们既不会失去祢，也不会失去自己。请祢使我们幸福，因为\BibleQuote{有福的百姓，他们的天主是主。}如果我们说祂是我们的产业，主也不会生气。因为我们读到：\BibleQuote{主是我继承的份。}何等伟大的事，弟兄们，我们既是祂的产业，祂也是我们的产业，因为我们既侍奉祂，祂也栽培我们。祂栽培我们并不有损祂的尊荣。因为如果我们侍奉祂如我们的天主，祂也栽培我们如祂的田地。并且（为了让你知道祂栽培我们），请听祂所派到我们这里的那一位说：\BibleQuote{我，}祂说，\BibleQuote{是葡萄树，你们是枝条，我父是农人。}所以祂确实栽培我们。如果我们结出果实，祂为我们预备祂的仓房。但如果在如此伟大之手的照料下我们仍不结果实，并结出荆棘而非好果实，我不愿说出随后的事。让我们以喜乐的主题结束。让我们转向主，等等。\par

\SourceRecord{Translated by R.G. MacMullen.；Nicene and Post-Nicene Fathers, First Series；Vol. 6.；Edited by Philip Schaff.；Buffalo, NY: Christian Literature Publishing Co.,；1888.；<http://www.newadvent.org/fathers/160363.htm>.}

% structure: p0001 source=h1 target=section reason=page-title

\section{新约讲道集第64篇}

\Emph{[114. 本笃会版]}\par

\Emph{论福音之言}，\BibleRef{路 17:3}\Emph{，\BibleQuote{如果你的兄弟犯罪，规劝他}等等，论及罪的赦免。}\par

\EditorialNote{曾在圣西彼廉祭台前，在博尼法爵伯爵面前宣讲。}\par

1. 我们刚才听到的福音诵读了关于罪赦的劝诫。现在我的讲道也当就这题目劝诫你们。因为我们是言语的仆役，不是我们自己的话，而是我们的天主和主的话；事奉祂没有不得光荣的，轻慢祂没有不受惩罚的。那么，我们的主天主，祂与圣父同在造了我们，又为了我们而成为受造者，重新造了我们；我们的主天主耶稣基督亲自对我们说了我们刚才在福音中所听到的话。\BibleQuote{如果，}祂说，\BibleQuote{你的兄弟得罪你，你要规劝他；如果他后悔，就宽恕他；如果他一天七次得罪你，又回来向你说：我后悔了，你要宽恕他。}祂不愿意人们把\InnerQuote{一天七次}理解为别的，只理解为\InnerQuote{无论多少次}，以免他得罪八次，你就不愿意宽恕。那么什么是\InnerQuote{七次}？就是常常，每次他犯罪并后悔。为此，\BibleQuote{一日七次我要赞美你}与另一篇圣咏中的\BibleQuote{祂的赞美常在我口中}是同样的意思。而七次用来代表常常有着最充分的理由：因为整个时间的过程都在七天循环的来回中旋转。\par

2. 所以，无论你是谁，你若思念基督，愿意领受祂所应许的，就不要迟延去做祂所命令的。那么祂应许了什么？\BibleQuote{永生。}祂命令了什么？就是宽恕你的兄弟。祂仿佛对你说：\Quote{人啊，宽恕一个人，好让我——我是天主——能来到你那里。}但暂且不提那些更高超的神圣应许，在其中我们的创造主应许使我们与祂的天使同等，好叫我们与祂、在祂内、藉着祂永远生活；现在不说这些，你不是希望从天主那里领受这个——你所被命令赐给你兄弟的——东西吗？正是这个，我说，你所被命令赐给你兄弟的东西，你难道不希望从你的主那里领受吗？告诉我，如果你不希望，就不要给。这岂不是说，如果有人向你请求宽恕，你就该宽恕他，如果你也希望被宽恕吗？但如果你没有什么需要被宽恕的，我敢说，你就不要宽恕。虽然我甚至不该说这话。即使你没有什么需要被宽恕的，也要宽恕。\par

3. 你正要对我说：\Quote{但我不是天主，我是人，是个罪人。}天主感谢你承认你有罪。那么宽恕吧，好使你得到宽恕。然而主天主亲自劝勉我们效法祂。首先天主自己——基督——劝勉我们，关于祂，宗徒伯多禄曾说：\BibleQuote{基督为我们受了苦，给你们留下榜样，叫你们跟随祂的足迹；祂没有犯过罪，口中也找不到诡诈。}祂确实没有罪，却为我们的罪死了，并倾流祂的血以赦罪。祂为我们的缘故承担了本不属于祂的，好把我们从所应得的当中拯救出来。死亡不属于祂，生命也不属于我们。为什么？因为我们是有罪的。死亡不属于祂，生命也不属于我们；祂接受了不属于祂的，祂赏赐了我们不属于我们的。但既然我们正在谈论罪赦，恐怕你们认为效法基督是太高的事，请听宗徒说：\BibleQuote{彼此宽恕，如同天主在基督内宽恕了你们一样。}\BibleRef{弗 4:32}\BibleQuote{所以你们要效法天主。}\BibleRef{弗 5:1}\BibleQuote{你们要效法天主，如同蒙爱的儿女一样。}这是宗徒的话，不是我的话。效法天主难道是一件傲慢的事吗？请听宗徒：\BibleQuote{你们要效法天主，如同蒙爱的儿女一样。}你被称为儿女：如果你拒绝效法祂，为什么寻求祂的产业？\par

4. 即使你没有任何罪想得到赦免，我也要说这话。但实际上，无论你是谁，你是人；即使你是义人，你是人；无论你是平信徒、隐修士、文书、主教，还是宗徒，你是人。请听宗徒的声音：\BibleQuote{如果我们说我们没有罪，便是欺骗自己。}那位著名的若望，一位福音作者，主基督在众人中特别爱他，他曾靠在主怀里，他说：\BibleQuote{如果我们说。}他没有说：\BibleQuote{如果你们说你们没有罪，}而是\BibleQuote{如果我们说我们没有罪，便是欺骗自己，真理也不在我们内。}他把自己也列在罪过中，好让自己也被列入宽恕之中。\BibleQuote{如果我们说。}请想是谁在说：\BibleQuote{如果我们说我们没有罪，便是欺骗自己，真理也不在我们内。但如果我们承认自己的罪，祂是忠信和公义的，必要赦免我们的罪，洗净我们的一切不义。}\BibleRef{若一 1:8-9}祂如何洗净？藉着宽恕，不是好像祂没有找到可惩罚的，而是找到可宽恕的。所以，弟兄们，如果我们有罪，让我们宽恕那些向我们请求的人。不要在心里对别人怀恨。因为怀恨比任何事更能败坏我们这颗心。\par

5. 那么，我愿你们宽恕，因为我看见你们求宽恕。有人求你们，就宽恕；你们求人，自己也要去求；有人求你们，就宽恕；你们将要求人宽恕；因为，看，祈祷的时刻将至：我使你们紧握你们将要说出的话语。你们将说：\BibleQuote{我们在天的父。}因为你们若不念\BibleQuote{我们的父}，就不在子女之列。所以你们将说：\BibleQuote{我们在天的父。}接着说：\BibleQuote{愿你的名被尊为圣。}再说：\BibleQuote{愿你的国来临。}继续：\BibleQuote{愿你的旨意奉行在人间，如同在天上。}看你们接着添加什么：\BibleQuote{求你今天赏给我们日用的食粮。}你们的财富在哪里？所以你们是乞丐。然而与此同时（这正是我所讲的），说出紧接\BibleQuote{求你今天赏给我们日用的食粮}之后的话。说出这一句后面的话：\BibleQuote{宽免我们的罪债。}现在你们触及我的话语：\BibleQuote{宽免我们的罪债。}有何权利？何种盟约？何种条件？何种明文规定？\BibleQuote{如同我们也宽免得罪我们的人。}你们不宽恕，这还算小事；不仅如此，你们还向天主撒谎。条件已定，律法已立。\Quote{如同我宽恕，你们也要宽恕。}因此，除非你们宽恕，否则祂也不宽恕。\Quote{如同我宽恕，你们也要宽恕。}你们希望在祈求时得到宽恕，就宽恕那求你们的人。精通天上法律的祂口授了这些祷词：祂不欺骗你们；按祂天乡之声的要旨祈求：说\BibleQuote{宽免我们，如同我们也宽免}，并实行你们所说的。在祈祷中说谎的人，失去他所求的恩惠：在祈祷中说谎的人，既输了官司，又遭刑罚。若有人在皇帝面前说谎，他一来到就被揭穿；但你在祈祷中说谎，你的祈祷本身就将你定罪。因为天主无需见证人来定你们的罪。口授祷词给你们的，就是你们的辩护者：若你们说谎，祂就是反对你们的见证人；若你们不改过，祂将作你们的审判者。所以，既要说出，也要做到。因为若不说，你们就不能得到所求，违反律法；但若说了而不做，更是犯上说谎的罪。除非我们实行所说的，否则无法逃避那句经文。我们能将这经文从祈祷中抹去吗？你们愿意保留\BibleQuote{宽免我们的罪债}这一句，而抹去后面\BibleQuote{如同我们也宽免得罪我们的人}吗？你们不可抹去它，免得自己先被抹去。所以在这祈祷中，你们说\Quote{赐给}，也说\Quote{宽赦}：好叫你们得到所没有的，并得宽赦所犯的过错。所以，你们希望获得，就给予；你们希望被宽恕，就宽恕。这是个简短总结。在另一处，听基督亲自说：\BibleQuote{你们宽恕，就必蒙宽恕。}你们要宽恕什么？别人得罪你的事。你们要蒙宽恕什么？你们自己所犯的罪。\Quote{宽恕。} \BibleQuote{你们给，就必给你们所渴望的，}永生。支持穷人的现世生命，维持穷人当下的生活，为了这如此微小和地上的种子，你们将收获永生。阿们。\par

\SourceRecord{Translated by R.G. MacMullen.；Nicene and Post-Nicene Fathers, First Series；Vol. 6.；Edited by Philip Schaff.；Buffalo, NY: Christian Literature Publishing Co.,；1888.；<http://www.newadvent.org/fathers/160364.htm>.}

% structure: p0001 source=h1 target=section reason=page-title

\section{新约讲道集 第65篇}

\Emph{[CXV. Ben.]}\par

\Emph{论福音之言}，\BibleRef{路 18:1} \Emph{，\BibleQuote{人应当时常祈祷，不要灰心}等，以及论及上圣殿祈祷的两个人，和那些被带到基督面前的小孩子。}\par

1. 圣福音的教训建立我们于祈祷和相信的责任，以及不自信而信赖主。还有什么比那比喻——不义的判官——给予我们更大的祈祷鼓励？因为一个不义的判官，不敬畏天主，也不尊重人，却因寡妇的纠缠而屈服，听允了她，并非出于仁慈。如果连那个讨厌被求的人尚且听了她的祈祷，那么那个鼓励我们祈求的祂，岂不更要垂听吗？因此，主用这个反面比较教导说：\BibleQuote{人应当时常祈祷，不要灰心。}祂又接着说：\BibleQuote{但是人子来临时，你认为祂在地上能找到信德吗？}如果信德失败，祈祷便消亡。因为谁会为他所不信的事祈祷呢？所以蒙福的宗徒在劝勉祈祷时也说：\BibleQuote{凡呼号主名的人，必将得救。}为显明信德是祈祷的泉源，他接着说：\BibleQuote{但怎能呼号他们所不信的祂呢？}因此，为使我们可以祈祷，让我们相信；并且为使这祈祷所用的信德不致失落，让我们祈祷。信德涌出祈祷，而祈祷的涌出获得信德的坚固。我说，信德涌出祈祷，而祈祷的涌出甚至使信德本身得到坚固。为免信德在诱惑中失败，因此主说：\BibleQuote{你们要警醒祈祷，免得陷于诱惑。}\BibleQuote{警醒，}祂说，\BibleQuote{并祈祷，免得陷于诱惑。}所谓\BibleQuote{陷于诱惑}，不就是离弃信德吗？因为诱惑前进多远，信德就退后多远；诱惑退后多远，信德就前进多远。为使你们更清楚地知道，亲爱的，主所说\BibleQuote{你们要警醒祈祷，免得陷于诱惑}，是指信德免得它失败和消亡；祂在福音的同一处说：\BibleQuote{今夜撒殚想要筛你们像筛麦子一样，但我已为你祈祷，伯多禄，使你的信德不至失落。}那保护者祈祷，那处于危险中的岂不更应该祈祷吗？因为主的话说：\BibleQuote{当人子来临时，你认为祂在地上能找到信德吗？}祂指的是那完美的信德。因为在地上几乎找不到。看！这个天主的教会座无虚席：若没有信德，谁会到这里来？但若有圆满的信德，谁不能移山呢？看看那些宗徒：他们若没有大信德，就不会舍弃所有，践踏此世的希望，跟随主；但若他们有圆满的信德，就不会向主说：\BibleQuote{请增加我们的信德。}再看那个人，他承认自己的状况（看，有信德，但不圆满），当他带儿子来求主治愈恶灵，并被问是否相信时，他回答说：\BibleQuote{主，我信，但请帮助我的不信。}\BibleQuote{主，}他说，\BibleQuote{我信，}我相信；因此有信德。但\BibleQuote{请帮助我的不信}，因此信德并不圆满。\par

2. 但因信德不属于骄傲的人，而属于谦卑的人，\BibleQuote{祂向那些自充为义人而轻视他人的人讲了这比喻：有两个人上圣殿去祈祷，一个法利塞人，另一个税吏。法利塞人说：天主，我感谢你，因为我不像其他的人。} 他至少可以说\Quote{像许多人}。\Quote{不像其他的人}是什么意思？岂不就是除了他自己以外的所有人？他说\Quote{我}，\Quote{是正义的，其余的都是罪人。}\BibleQuote{我不像其他的人，勒索、不义、奸淫。} 并且，看，从你的邻人，那个税吏身上，你找到了更骄傲的借口。他说\Quote{像}，\Quote{这个税吏。} 他说\Quote{我}，\Quote{独自一人，他属于其余的人。} 他说\Quote{我不像他那样}，\Quote{因我的义行，我没有任何不义。}\BibleQuote{我每周禁食两次，凡我所得的，都捐献十分之一。} 在他所有的话中，找一找他向天主祈求了什么，你必一无所获。他上去祈祷，却无意向天主祈祷，只愿赞美自己。不，这还算小事：他不仅没有向天主祈祷，反而赞美了自己。不仅如此，他还嘲笑那真正祈祷的人。\BibleQuote{但税吏远远地站着}；然而他实际靠近天主。心的自觉使他远离，虔诚使他亲近。\BibleQuote{但税吏远远地站着}；然而主视他为近。\BibleQuote{因为上主虽崇高，却垂顾卑微的人。} 但\Quote{那些高傲的人}，如这个法利塞人，\BibleQuote{祂只从远处认识}。的确，\Quote{高傲的人}\BibleQuote{天主只从远处认识}，但祂并不宽赦他们。请再听税吏的谦卑。他远远地站着还只是小事；\BibleQuote{他连举目望天都不敢}。他不观看，为使自己被观看。他不敢向上看，良心压迫他；但望德提升他。再听：\BibleQuote{他捶着自己的胸膛}。他惩罚自己；因此主因他的承认而宽恕了他。\BibleQuote{他捶着胸膛说：天主，可怜我这个罪人吧！} 看那祈祷的是谁？既然他承认自己的罪，你为何惊奇天主会宽恕？如此你已听了法利塞人和税吏的案例；现在听判决；你听了骄傲的控告者，你听了谦卑的罪犯；现在听审判者。\BibleQuote{我实在告诉你们}。真理在说，天主在说，审判者在说。\BibleQuote{我实在告诉你们：那个税吏从圣殿下去，成了正义的，而不是那个法利塞人。} 请教我们，主，原因是什么。看！我看到税吏从圣殿下去，成了正义的，而不是那个法利塞人。我询问为什么？你问为什么？听为什么。\BibleQuote{因为凡高举自己的，必被贬抑；凡贬抑自己的，必被高举。} 你已听了判决，要提防它的恶因。换言之，你已听了判决，要提防骄傲。\par

3. 现在让那些不敬的妄谈者，无论他们是谁，那些倚靠自己能力的人，让他们听见并看见这些事：让他们听见那些说\Quote{天主造了我成人，我使自己成为义人}的人。哦，你比法利塞人更坏、更可憎！福音中的法利塞人确实称自己为义人，但他仍为此感谢天主。他称自己为义人，但仍感谢天主。\BibleQuote{天主，我感谢你，因为我不像其他的人。}\BibleQuote{天主，我感谢你。}他感谢天主，因为他不像其他的人；然而他仍被责备为骄傲自大；不是因他感谢天主，而是因他似乎不再需要加添什么。\Quote{我感谢你，因为我不像其他的人，不义。}那么你就是义人；那么你就不求什么；那么你已满盈；那么人生在世不是考验；那么你已满盈；那么你已丰盛；那么你没有理由说\BibleQuote{宽免我们的罪债！}那么，如果那骄傲地感谢的人尚且被责备，那冒犯恩宠的人又当如何呢？\par

4. 并且，看，在案件陈述和判决宣告之后，小孩子也来了，其实，是被抱着带来接受触摸。接受谁的触摸，岂不是那位医生？当然，人们会说，他们是健康的。婴孩们被带来接受触摸，是给谁？给谁？给救主。若是给救主，他们就被带来得救。给谁，岂不是给那\BibleQuote{来寻找并拯救迷失的人}？他们如何迷失的？就他们个人而言，我看见他们无罪，我在寻找他们的罪过。它从何而来？我听从宗徒的话：\BibleQuote{藉着一个人，罪恶进入了世界。藉着一个人，}他说，\BibleQuote{罪恶进入了世界，死亡藉着罪恶也进入了世界；这样死亡就殃及了众人，因为众人都犯了罪。}让小孩子来吧，让他们来：愿主被听从。\BibleQuote{让小孩子到我这里来！}让小的来，让病人到医生那里，迷失者到他们的救赎主那里：让他们来，不要有人阻止他们。在枝条上，他们还没有犯任何恶，但他们在根上已经败坏了。愿主祝福小的与大的。让医生触摸大小。小孩子的案件我们托付给他们的长者。你们为那些哑巴说话，为那些哭泣的祈祷。如果你们不是无益地作他们的长者，就作他们的保护者：为那些尚不能管理自己案件的人辩护。损失是共同的，愿找回也是共同的：我们所有人都一同迷失了，一同在基督里被找到。功绩不均，但恩宠是共同的。他们没有别的恶，只有从源头汲取的；他们没有别的恶，只有从最初根源所继承的。不要让那些在他们所继承的之上又加添更多恶的人阻止他们得救。年长者在年龄上是长者，在邪恶上也是长者。但天主的恩宠抹去你所继承的，也抹去你所加添的。因为\BibleQuote{罪恶在那里越多，恩宠在那里也越格外丰富。}\par

\SourceRecord{Translated by R.G. MacMullen.；Nicene and Post-Nicene Fathers, First Series；Vol. 6.；Edited by Philip Schaff.；Buffalo, NY: Christian Literature Publishing Co.,；1888.；<http://www.newadvent.org/fathers/160365.htm>.}

% structure: p0001 source=h1 target=section reason=page-title

\section{\Emph{论新约的讲道 66}}

\Emph{[CXVI. Ben.]}\par

\Quote{论福音之言：\BibleBook{路加}{福音} 24:36}，\Emph{\Quote{祂亲自站在他们中间，对他们说：愿你们平安。}}等。\BibleRef{路 24:36}\par

1. 主在祂复活后向祂的门徒显现，正如你们所听见的，并问候他们说：\Quote{愿你们平安。}这实在是平安，是救恩的问候；因为\Quote{问候}一词本身即源自\Quote{救恩}。还有什么比救恩本身问候人更好的呢？因为基督就是我们的救恩。祂是我们的救恩，祂为我们受伤，被钉在木头上，从木头上被取下，安放在坟墓中。祂从坟墓中复活了，伤口愈合了，伤痕却保留着。祂认为这对祂的门徒是有益的，祂的伤痕得以保留，以医治他们心中的创伤。什么创伤？不信的创伤。因为祂显现于他们眼前，显示真实的肉身，而他们却以为看见了神魂。这不是轻微的创伤，这是心灵的创伤。是的，那些停留在这创伤中的人形成了恶毒的异端。但我们难道以为门徒们没有受伤，因为他们这么快就被治愈了？亲爱的，只要想一想，如果他们继续存留在这创伤中，认为那被埋葬的身体不能复活，而是一个神魂以身体的形象欺骗了人的眼睛；如果他们继续持守这样的信念，甚至是不信，那么我们就不该哀悼他们的伤口，而该哀悼他们的死亡了。\par

2. 但主耶稣说了什么？\Quote{你们为什么烦乱？为什么心里起疑念？}如果疑念升入你们的心，那疑念来自地上。但一个人好的不是疑念升入他的心，而是他的心本身升到高处，正如宗徒愿意信者将心放在那里，他对他们说过：\Quote{你们既然与基督一同复活了，就该追求天上的事，那里有基督坐在天主的右边。你们要思念上面的事，不要思念地上的事。因为你们已经死了，你们的生命与基督一同藏在天主内。当基督、你们的生命显现时，你们也要与祂一同在光荣中显现。}在什么光荣中？复活的光荣。在什么光荣中？听宗徒论及这身体说：\Quote{播种的是可羞辱的，复活起来的是光荣的。}门徒们不愿将这光荣归于他们的师傅、他们的基督、他们的主；他们不信祂的身体能从坟墓中复活；他们看见祂的肉身，却以为祂是神魂，他们连自己的眼睛也不信。然而，我们相信那些宣讲祂却不显示祂的人。看，他们不信那亲自显现于他们的基督。恶毒的创伤！让这些伤痕的药方出来吧。\Quote{你们为什么烦乱？为什么心里起疑念？看我的手和我的脚，}就是我曾被钉子钉穿的地方。\Quote{你们摸一摸，看一看。}但你们看见了，却仍然看不见。\Quote{摸一摸，看一看。}什么？\Quote{因为神魂没有骨肉，你们看我是有的。}祂说了这话，就如经上所记，\Quote{祂就把手和脚给他们看。}\par

3. \Quote{他们正在犹豫不决、欢喜惊奇的时候。}这时已经有喜乐，但犹豫依然存在。因为一件不可思议的事发生了，却又确实发生了。难道今天主的身體从坟墓中复活仍是一件不可思议的事？整个洁净了的世界已经相信了；谁不相信，就仍留在他的不洁中。然而在那时，这是不可思议的；不单是眼睛，连手也受了说服，好让信德借着身体的感官进入他们的心，并且那信德如此进入他们的心，然后传遍世界，给那些既未看见也未触摸，却毫不怀疑就相信的人。祂说：\Quote{你们这里有吃的没有？}这良善的建筑师为了建造信德的殿宇，还做了多少事？祂并不饥饿，却求食。祂吃，是因祂的权能，并非出于需要。所以，让门徒们承认祂身体的真实性吧，这个世界因他们的宣讲而承认了。\par

4. 若偶然有异端者，心中仍坚持说基督显现给人看，但基督的并不是真正的血肉；让他们现在抛开这错误吧，让福音说服他们。我们只是责备他们怀有这种想法；如果他们坚持不改，祂将定他们的罪。你是谁，竟不信那放入坟墓的身体能复活？如果你是摩尼教徒，你连祂被钉十字架也不信，因为你连祂降生也不信，你声称祂所显示的一切都是虚假的。祂显示了虚假，而你说的是真理？你口里不说谎，难道祂的身体说谎了吗？看，你认为祂显现于人的眼睛的并非祂真实所是，祂是神魂，不是血肉。听祂说话：祂爱你，愿祂不定你的罪。听祂说话：看，祂对你说，你这不幸的人，祂对你说：\Quote{你们为什么烦乱？为什么心里起疑念？}\Quote{看，}祂说，\Quote{我的手和我的脚。摸一摸，看一看，因为神魂没有骨肉，你们看我是有的。}这是真理说的，难道祂会欺骗吗？那么，那是身体，是血肉；那被埋葬的，显现了。让怀疑消失，赞美随之而来吧。\par

5. 那时，祂向门徒们显示了自己。\Quote{祂自己}是什么？祂的教会之头。教会被祂预见，如同在你们中将要遍布全世界，但门徒们尚未看到。祂显示了头，祂预许了身体。因为祂接着说了什么？\BibleQuote{这些话是我还同你们在一起的时候对你们说过的。}这是什么，\Quote{我还在同你们一起的时候}？难道那时祂不是与他们在一起，当祂对他们说话时？\Quote{当我还在与你们一起的时候}是什么意思？与你们一起如同可死的，现在我已不是。我与你们一起，那时我还需要死。\Quote{与你们一起}是什么意思？与你们这些将要死的人一起，我自己也要死。如今我不再与你们一起：因为我与那些将要死的人一起，我自己永远不再死。这就是我对你们所说的话。什么？\BibleQuote{凡在法律、先知和圣咏上指向我所写的一切，都必须应验。}我曾告诉你们一切都必须应验。\BibleQuote{于是祂开启了他们的明悟。}来吧，主啊，请用祢的钥匙开启，好让我们能明白。看，祢述说一切，却不被相信。祢被认为是一个幽灵，被触摸，被粗鲁地对待，然而那些摸祢的人却犹豫。祢从圣经中告诫他们，但他们仍不明白祢。他们的心关闭了，请打开，并进入。祂这样做了。\BibleQuote{于是祂开启了他们的明悟。}开启吧，主啊，是的，开启那对基督怀疑之人的心。开启\Quote{他}的明悟，使他相信基督不是一个幻象。\BibleQuote{于是祂开启了他们的明悟，叫他们理解圣经。}\par

6. 并且\BibleQuote{祂对他们说。}什么？\BibleQuote{必须这样。经上这样记载，必须这样。}什么？\BibleQuote{基督必须受苦，并第三天从死者中复活。}这事他们看见了，他们看见祂受苦，他们看见祂被悬挂，他们看见祂复活后与他们一起活着。他们那时没有看见什么？身体，即教会。他们看见了祂，却没有看见她。他们看见了新郎，新娘还隐藏着。让他也应许她。\BibleQuote{经上这样记载，基督必须受苦，并第三天从死者中复活。}这是新郎，新娘怎么样？\BibleQuote{并且悔改与罪过的赦免，必须因祂的名从耶路撒冷开始，宣讲给万民。}这事门徒们还没有看见：他们还没有看见遍布万国的教会，从耶路撒冷开始。他们看见了头，并相信头关乎身体的事。藉着他们所看见的，他们相信了他们没有看见的。我们也像他们：我们看见了一些他们没有看见的，也看不见一些他们看见了的东西。我们看见什么，是他们没有看见的？遍布万国的教会。我们看不见什么，是他们看见的？基督在肉身中显现。正如他们看见了祂，并相信了关于身体的事，同样我们看见了身体；让我们相信关于头的事。让我们各自所看见的帮助我们自己。看见基督帮助他们相信未来的教会；看见教会帮助我们相信基督已复活。他们的信德得以成全，我们的信德也得以成全。他们的信德从看见头得以成全，我们的信德因看见身体得以成全。基督被\Quote{完整地}启示给他们，也同样启示给我们：但他们未曾\Quote{完整地}看见祂，我们也未曾\Quote{完整地}看见祂。在他们，头被看见了，身体被相信了；在我们，身体被看见了，头被相信了。然而基督对任何人都没有欠缺：在一切中祂是完整的，尽管直到今天祂的身体仍然不完全。宗徒们相信了；藉着他们，许多耶路撒冷的居民相信了；犹太相信了。撒玛黎雅相信了。让肢体被添加上，建筑物被添加上到根基上。\BibleQuote{因为没有人能奠立别的根基，}宗徒说，\BibleQuote{除了那已奠立的，即耶稣基督。}让犹太人大发狂怒，满怀嫉妒：斯德望被石头砸死，扫罗看守那些用石头砸他之人的衣服，扫罗，有一天成为宗徒保禄。让斯德望被杀，耶路撒冷教会混乱地分散：从其中燃起燃烧的火把，散开并扩展其火焰。因为在耶路撒冷教会中，当他们都一心一意、一心向主时，就好像燃烧的火把被圣神点燃。当斯德望被石头砸死时，那堆柴遭了迫害：火把四散，世界被点燃了。\par

7. 然后，扫禄一心想着他狂怒的计划，从司祭长那里领取了信函，开始他残酷暴怒的旅程，口吐杀机，渴求鲜血，要把凡能找到的人，从四面八方拖来捆绑，押送去受刑，以他们的血来满足自己。但天主在哪里？基督在哪里？那给斯德望戴上荣冠的主在哪里？除了在天上，还能在哪里？愿祂现在注视着扫禄，嘲笑他的暴怒，并从天上呼喊：\Quote{\InnerQuote{扫禄，扫禄，你为什么迫害我？}我在天上，你在地上，你却迫害我。你碰不到身体，却践踏我的肢体。然而你在做什么？你得到了什么？\InnerQuote{你用脚踢刺棒是难的。}放下你的怒气吧，恢复健全。放弃恶谋，寻求善助。}那声音把他击倒在地。谁被击倒了？迫害者。看，仅此一句话他就被战胜了。你往哪里去？你的狂怒把你带到哪里去了？那些你正搜寻的人，如今你跟随他们；你所迫害的人，如今你为他们受苦。他从地上站起来，成为传道者，而那迫害者被击倒在地。他听见了主的声音。他眼瞎了，但只是身体上的瞎，好叫他内心得以光照。他被带到阿纳尼雅那里，接受多方面的要理讲授，受了洗，然后作为一位宗徒出来。那么，讲吧，宣讲吧，宣讲基督，传播祂的教义，你这羊群的好牧者，不久前还是一头狼。看他，注视他，这个曾经狂怒的人。\BibleQuote{至于我，我绝不夸耀，除非夸耀我们的主耶稣基督的十字架，借着它，世界对我来说已被钉在十字架上，我对于世界也被钉在十字架上。}传播福音：用你的口散布你在心中所怀的。让万民听见，让万民相信；让万民增多，让主那披着紫袍的新娘从殉道者的血中萌发。从她那里已经出来了多少人，多少肢体已紧附于元首，并仍然紧附着祂，并且相信！他们受了洗，他人也将受洗，之后还会有别人来。那么，我说，在世界终结时，石头将被连接到根基上，活石，圣石，好让整个建筑由那教会，是的，正是这教会，在房屋建造的同时唱着新歌。因为圣咏本身这样说：\BibleQuote{当房屋在被掳之后建造时；}它说什么？\BibleQuote{你们要向主唱新歌，普世大地，你们要向主歌唱。}这房屋何等伟大！但它在何时唱新歌呢？当它被建造时。它何时被奉献？在世界终结时。它的根基已经奉献了，因为祂升了天，不再死亡。当我们也要复活不再死亡时，那时我们将被奉献。\par

\SourceRecord{Translated by R.G. MacMullen.；Nicene and Post-Nicene Fathers, First Series；Vol. 6.；Edited by Philip Schaff.；Buffalo, NY: Christian Literature Publishing Co.,；1888.；<http://www.newadvent.org/fathers/160366.htm>.}

% structure: p0001 source=h1 target=section reason=page-title

\section{新约讲道集第67篇}

\Emph{[一一七．本笃版]}\par

\Emph{论福音之言}，\BibleRef{若 1:1}，\Emph{，\BibleQuote{在起初已有圣言，圣言与天主同在，圣言就是天主}，等等。驳斥亚略派。}\par

1. 所读的福音段落，最亲爱的弟兄们，期待心灵的纯洁之眼。因为从若望福音我们已明白我们的主耶稣基督，按祂的天主性是为了创造万有，按祂的人性是为了拯救堕落的受造物。如今在同一福音中，我们找到若望是怎样的一个人，他有多么伟大，以便从这分施者的尊严可以明白那能被如此之人所宣告的圣言有多么珍贵；其实，更应说那超越万有的圣言是无法估价的。因为任何可购买之物要么等于其价格，要么低于它，要么超过它。当一人以价值相当的价格购得一物时，价格等于所购之物；当以较低价格时，则低于它；当以较高价格时，则超过它。但对于天主的圣言，没有一样东西能与之相等，也没有任何东西可以低于它或高于它作为交换。因为万物都可以低于天主的圣言，因为\BibleQuote{万物是藉着祂造成的}；然而它们并非如此低于，仿佛它们是圣言的价格，使人可以给出某物以领受那圣言。但如果我们可以这样说，并且如有任何原则或说话习惯容许这种表达，那么获取圣言的价格，就是那获取者本人，他将自己为自己给了这圣言。因此，当我们购买任何东西时，我们寻找某物来给出，以便用我们给出的价格获得我们希望购买的东西。而我们给出的东西是在我们之外的；如果它本在我们之内，我们给出的东西就变得在我们之外，以便我们所获取的可以在我们之内。无论购买者发现怎样的价格，那价格必须如此：他给出他所拥有的，并接受他所没有的；然而，如此，那付出价格的人自己仍然存在，而他所支付价格得到的东西是加给他的。但谁若想获取这圣言，谁若想拥有它，不要在他自己之外寻找任何东西来给出，让他给出自己。而当他这样做时，他并不失去自己，如同他购买东西时失去价格一样。\par

2. 天主的圣言于是呈现在所有人面前；谁若能够，就获取它，而有虔诚意愿的人就能够。因为在那圣言中有平安；而\BibleQuote{平安在地上归于善人}。因此，谁若愿意获取它，就让他给出自己。这就像是圣言的价格，如果可以这样说的话，当给出者并不失去自己，而且获得了那圣言，他为了这圣言给出自己，并且在他将自己给与的圣言中也获得了自己。而他给圣言什么呢？不是任何属于其他而非属于那圣言的东西，因为他把自己给与它；而是由同一圣言所造成的，就归还给祂以便被重新造成；\BibleQuote{万物是藉着祂造成的}。如果万物，那么当然人也包括在内。如果天、地、海和其中的万物，如果整个受造物；那么当然更明显的是，那按天主的肖像由圣言造成的人。\par

3. 弟兄们，我现在不是在讨论如何理解\BibleQuote{在起初已有圣言，圣言与天主同在，圣言就是天主}这些话。它们可以以一种不可言喻的方式被理解；无法用人的言语使人理解。我正在讲论天主的圣言，并告诉你们它为何不被理解。我现在不是为使它被理解而讲，而是告诉你们是什么阻碍了它被理解。因为祂是一种\OriginalTerm{形式}{Form}，一种未受造的形式，却是所有受造之物的形式；一种不变的形式，没有欠缺，没有朽坏，没有时间，没有空间，超越万有，存在于万有之中，如同一种基础，万物在其中，又如同一种冠石，万物在其下。如果你说万物都在祂内，你并没有说谎。因为这圣言被称为天主的\OriginalTerm{智慧}{Wisdom}；我们也有经上写着：\BibleQuote{祢以智慧造成了万物}。看，那么在祂内一切都是；然而，既然祂是天主，那么万物都在祂之下。我正在表明所读的是如何不可理解；然而它被读出来，不是为了被人理解，而是为了让人悲哀他不能理解，并找出是什么阻碍他理解，并除去这些障碍，使自己从较坏变为较好，渴望领悟那不变的圣言。因为圣言并不因认识祂的人的增加而前进或增长；而是完整的，如果你留下；完整的，如果你离开；完整的，当你回来时；祂在自身内存在，并更新一切。那么祂是万有的形式，一种未受造的形式，没有时间，如我所说，也没有空间。因为凡在空间内的，都被限定。每一种形式都由界限所限定；它有达到何处的极限。再者，凡在地方内并有一定体积和空间伸展的，其各部分小于整体。愿天主使你们能理解。\par

4. 现在，从每日出现在我们眼前的物体，我们看见的、触摸的、生活在其中的物体，我们能够判断，每个物体在空间中都有一形式。现在，占据一定空间的每样东西，其各部分小于整体。例如，手臂是人体的一个部分；当然手臂小于整个身体。如果手臂较小，它占据的空间也较小。同样，头部，作为身体的一部分，被包含在较小的空间中，小于它作为头部的整个身体。所以，所有在空间中的东西，其各部分分别小于整体。对于那圣言，我们不要持有这样的观念，这样的思想。我们不要根据肉体的暗示来形成对属灵事物的概念。那圣言，那天主，在部分中并不小于在整体中。\par

5. 但你们无法想象任何这样的事。这种无知比傲慢的知识更为虔诚。因为我们谈论的是天主。经上说：\BibleQuote{圣言就是天主。}我们谈论的是天主；你若不明白，又何足为奇？因为若你明白，祂就不是天主。愿有对无知虔诚的承认，而非轻率的自称有知。凭心智在任何程度上触及天主，乃莫大的真福；但要完全领悟祂，却全然不可能。天主是心智的对象，祂是应被理解的；身体是眼睛的对象，它是被看见的。但你认为你能用眼睛领悟一个身体吗？你根本不能。因为无论你看什么，你都没有看见整体。如果你看到一个人的脸，你在看到脸的时候并没有看到他的后背；当你看到后背时，你当时并没有看到脸。所以你并没有以领悟的方式来看；但当你看到另一个你以前没有看到的部分时，除非记忆帮助你记住你曾看到过你正在离开的那个部分，否则你永远不能说你在表面上理解了任何东西。你处理你所看到的，把它翻来覆去，或者你自己绕着它走去看整体。因此，凭一眼你无法看到整体。只要你转动它去看，你就只是在看部分；并通过把你已经看到其他部分拼凑起来，你幻想你看到了整体。但这不应被理解为眼睛的视觉，而是记忆的活动。那么，弟兄们，关于那圣言还能说什么呢？看，对于眼前的身体，我们无法凭一瞥就领悟它们；那么什么样的心灵之眼能领悟天主呢？只要眼睛纯洁，触及祂就够了。但如果触及，那是通过一种无形和属神的接触，却并不领悟；而且，那只靠纯洁。人藉着用心接触那永远有福者而成为有福的；那就是这永恒的真福本身，那永恒的生命，人由此得生；那完美的智慧，人由此成为智慧；那永恒的光明，人由此得蒙光照。且看，通过这一接触，你成了你以前所不是的，你并没有使你接触的那位变成祂以前所不是的。我再重复一遍，认识天主的人并没有给天主带来任何增加，而是认识祂的人，从天主的认识中获得益处。亲爱的弟兄们，让我们不要以为我们赠予天主任何好处，因为我说过我们好像给了祂一个代价。因为我们没有给祂任何可以增加祂的东西；当你堕落时，祂是圆满的；当你回归时，祂仍保持圆满，随时愿意显现自己，以赐福那些转向祂的人，并以盲目惩罚那些远离祂的人。因为通过这种盲目，作为惩罚的开始，祂首先向那远离祂的灵魂施行报复。因为谁远离真光，即天主，就立刻变瞎。他尚未察觉自己的惩罚，但他已经接受了。\par

6. 因此，亲爱的弟兄们，让我们明白，天主的圣言是无形体、不可侵犯、不变、无时间性的诞生，却出自天主。我们能否以为，我们能以某种方式说服某些不信者，说那不是真的，与按照天主教信仰我们所说的真理一致，而这正与亚略派相反，天主的教会常被他们所考验，因为属血肉的人更容易接受他们习以为常看到的事物？因为有些人竟敢说：\Quote{父大于子，并在你们内先于祂；}也就是说，父大于子，子小于父，并在你们内被父在先。他们这样辩论：\Quote{如果祂被生，当然父先于祂的儿子被生给祂。}请听；愿祂与我同在，因你们的祈祷帮助我，并以虔诚的心渴望接受祂所赐予、所提示给我的；愿祂与我同在，使我能在某种程度上解释我所开始的。然而，弟兄们，在我开始之前我告诉你们，若我不能解释，不要以为这是论证的失败，而是人的失败。因此我劝勉并恳求你们祈祷；愿天主的仁慈与我同在，并使这问题被我解释得合宜于你们听，也合宜于我说。他们就这样说：\Quote{如果祂是天主子，祂就是被生的。}这是我们承认的。因为祂若不是被生的，就不是儿子。显然，信仰承认这一点，天主教会赞同它，这是真理。他们接着说：\Quote{如果子是父所生，父就存在于子被生给祂之前。}这一点信仰拒绝，天主教徒的耳朵拒绝它，它是被诅咒的；谁怀有这种想法，就是外人，他不属于圣人的团体和共融。于是他说：\Quote{请给我一个解释，子如何能为父所生，却又与生祂的那一位同龄？}\par

7. 弟兄们，当我们向属血气的人传授属神之事的道理时，即使我们自己并非属血气的人，当我们向属血气的人——那些习惯于尘世生育、看见受造物的秩序、在那里继承和分离使生育者与被生者在年龄上分开的人——暗示这些属神的真理时，我们能做什么呢？因为儿子是在父亲之后出生，为了继承父亲，而父亲在时间中当然必定会死亡。我们在人身上发现这一点，在其他动物中也发现这一点：父母在先，子女在时间上在后。透过这种观察的习惯，他们想要将属血气的事转移到属神的事上，并且由于专注于属血气的事，更容易陷入错误。因为并非听者的理性跟随那些宣讲这些事的人，而是习惯甚至纠缠他们自己，以至于他们宣讲这些事。我们能做什么呢？我们应保持沉默吗？但愿我们能如此！因为也许通过沉默，某种配得上这不可言说之事的思想可以被想到。因为凡不能被言说的，就是不可言说的。天主是不可言说的。因为如果宗徒\Person{保禄}说，他\BibleQuote{被提到第三层天，并听见隐秘的话}，那么那显示这些事——这些事甚至不能由那被显示的人所说出——的主，更是何等不可言说呢？所以，弟兄们，如果我们能保持沉默并说：\Quote{信德包含这些；我们如此相信；你尚且不能接受，你还是个婴孩；你必须耐心忍受，直到你的翅膀长成，免得你在没有翅膀时想要飞翔，那将不是自由的翱翔，而是鲁莽的堕落。}那更好。对此他们反诘说：\Quote{哦，如果他有什么要说的，他会对我说。这只是错者的借口。那不愿回答的人被真理征服了。}听到这话的人，如果他不回答，虽然他自己未被征服，却在动摇的弟兄们中被征服了。因为软弱的弟兄们听到后，会认为真的没什么可说的；也许他们正确地认为没什么可说的，但却不是没什么可感觉的。因为一个人无法表达他不能感觉到的东西；但他可能感觉到一些他无法表达的东西。\par

8. 然而，保留那至上威能的不可言说性，免得当我们提出某些比拟反对他们时，有人会以为我们藉这些比拟已经达到了那不能被婴孩表达或理解的事（即使较有长进的人也只能部分地、在谜语中、仅\BibleQuote{透过镜子}；但尚未\BibleQuote{面对面}），让我们也提出某些比拟来反对他们，借此他们可以被驳斥，而不是\Quote{它}被理解。因为当我们说，很可能发生、可以被理解：祂既可以被生，却又与那生祂者同等永存，为了驳斥这点并证明它似乎是假的，他们就提出比拟反对我们。从何处？从受造物而来，他们对我们说：\Quote{当然，每个人在他生儿子之前，他在年龄上大于他的儿子；同样，马在生马驹之前，羊及其他动物也是如此。}他们就这样从受造物提出比拟。\par

9. 什么！难道我们也必须劳苦，以便找到我们所确立之事的相似物吗？如果我找不到任何相似物，难道我不能正当地说：\Quote{也许造物主的出生在受造物中没有自己的相似物？因为祂超越在此处的事物，一如祂在彼处，同样地，祂超越在此处出生的事物，一如祂在彼处出生。此处的一切事物都来自天主；但有什么能与天主相比呢？同样，此处出生的一切事物都是由祂的作为而出生。所以也许找不到祂出生的相似物，如同找不到祂的本质、不变性、天主性、威能的相似物一样。因为在此处能找到什么像这些呢？如果那么巧，连祂出生的相似物也找不到，那么我是否就因此被压倒了，因为当我渴望在受造物中找到与造物主相似之物时，却没有找到呢？}\par

10. 而在事实上，诸位弟兄，我不太可能找到任何暂时的比拟能与永远相比。但至于你们所发现的，它们是什么呢？你们发现了什么？那就是一个父亲在时间上大于他的儿子；因此你们就想要天主子在时间上小于永远圣父，因为你们发现一个儿子小于一个在时间中出生的父亲。请给我在这里找到一个永远的父亲，你们就找到了一个比拟。你们在时间中找到一个儿子小于父亲，一个在时间中的儿子小于一个在时间中的父亲。你们可曾给我找到一个在时间中的儿子比永远的圣父更年轻？既然在永远中有稳定，在时间中有变化；在永远中万物静止，在时间中一事前来，另一事接替；你们能在时间的变化中找到一个年岁较小的儿子继承他的父亲，因为他自己也继承了他的父亲，而不是一个在时间中出生的儿子继承一位永远的圣父。那么，诸位弟兄，我们如何在受造物中找到一样\OriginalTerm{共永恒}{coeternal}的东西，既然在受造物中我们找不到任何永恒的东西？你们在受造物中找到一个永远的圣父，我就要找到一个共永恒的圣子。但如果你们找不到一个永远的圣父，并且一个在你们中超越另一个；那么，为了比拟，我找到一样\OriginalTerm{共年代}{coeval}的东西就已足够。因为共永恒是一回事，共年代是另一回事。我们每天都称那些有着相同时间尺度的人为共年代；一个在时间上并未先于另一个，然而我们称为共年代的两者都曾开始\Quote{存在}。现在，如果我能发现一样东西，它与生出它的那东西是共年代的；如果能发现两样共年代的东西，即那生出者和那被生者；那么我们在这种情况下发现共年代的东西，就让我们在另一种情况下理解共永恒的东西。如果我在这里发现，被生者从生出者开始存在时就开始存在，那么我们至少可以理解，天主子并没有开始存在，因为那生祂的从无始就已存在。请看，诸位弟兄，我们也许能在受造物中发现一样东西，它由别的东西所生，却与生出它的那东西同时开始存在。在后一种情况中，一个在另一个开始时开始存在；在前一种情况中，一个没有开始存在，因为另一个从无始就已存在。因此，第一种是共年代，第二种是共永恒。\par

11. 我想你们的圣德已经明白我所说的：暂时之物不能与永远之物相比；但藉着一些轻微渺小的比拟，共年代之物可与共永恒之物相提并论。因此，让我们找出两样共年代的东西；并让我们从圣经中获取关于这些比拟的提示。我们在圣经中读到关于智慧的话：\Quote{她是不灭之光的辉映。}又读到：\Quote{天主德能的无玷明镜。}智慧本身被称为\Quote{不灭之光的辉映，}被称为\Quote{圣父的肖像；}让我们从中取一个比拟，好找出两样共年代的东西，从而理解共永恒之物。哦，你们\OriginalTerm{亚略派}{Arian}，如果我发现那生出者并不在时间上先于它所生者，被生者并不在时间上小于它所由生者；那么你们就得同意我，这些共永恒之物可以在造物主中找到，既然共年代之物可以在受造物中找到。我想这事确实已经出现在一些弟兄心中。因为我一说\Quote{她是不灭之光的辉映}，就有人预先想到了我。因为火发出光，光从火发出。如果我们问哪一个来自哪一个，我们每天点燃蜡烛时，就被提醒某些看不见、说不出的东西，好让我们理解的蜡烛仿佛在这世界的黑夜中被点燃。请看那点蜡烛的人。蜡烛未被点燃时，还没有火，也没有从火发出的光。但我要问：\Quote{光是从火而来，还是火从光而来？}每个灵魂都回答我（因为天主喜悦在每个灵魂中播下理解和智慧的开端）；每个灵魂都回答我，无人怀疑，那光来自火，而非火来自光。那么，让我们把火看作那光的父亲；因为我先前说过，我们是在寻找共年代之物，而非共永恒之物。如果我想要点蜡烛，那里还没有火，也还没有那光；但我一点燃它，火与光就一同出现。那么，请给我一个没有光的火，我就相信你：圣父曾经没有圣子。\par

12. 请注意；此事已按我所能作的最大程度解释过了，主旨由主助长了你们祈祷的热忱和心灵的准备，你们已尽所能领受了。然而这些事是无可言喻的。不要以为任何配得上此主题的话已被说出，只因有形的与同永恒的相比，暂时的与恒常的相比，可朽的与不朽的相比。但既然圣子也被称为\OriginalTerm{圣父的肖像}{Image of the Father}，让我们也从此得出某种相似，尽管是在非常不同的事物中，正如我先前所说。一个人照镜子时，他的影像从镜子中反射出来。但这并不能帮助我们阐明我们试图以某种方式解释的事。因为有人对我说：\Quote{当然，照镜子的人早已存在，并且在此之前已出生。影像只是在他看自己时出现。因为照镜子的人在来到镜子前就已存在。}那么，我们将找到什么，才能像我们从火与\OriginalTerm{光辉}{brightness}中那样引出一种相似呢？让我们从一件微小的事物中找吧。你毫不困难地知道水如何常常反射出物体的影像。我说的是，当有人经过或站在水边时，他看见自己的影像。那么，假设在水边生出一样东西，例如一株灌木或一棵草，它难道不是与它的影像一同出生的吗？它一开始存在，它的影像就与它一同开始存在；它的出生并不先于它自己的影像；我不能向你证明，有东西在水边出生，而它的影像后来才出现，起初它出现时没有影像；而是它与它的影像一同出生；然而影像来自它，而不是它来自影像。那么它与其影像一同出生，灌木和它的影像同时开始存在。你难道不承认影像是从那灌木所生，而非灌木从影像所生吗？那么你就承认影像来自那灌木。因此，那生育的和被生育的同时开始\Quote{存在}。所以它们是同岁。如果灌木一直存在，那么来自灌木的影像也会一直存在。现在，那从别的东西得到存在的东西，当然是从它而生。那么，生育者可能永远存在，并永远与那从它而生者同在。因为在这里我们曾困惑和苦恼，如何理解\OriginalTerm{永恒诞生}{Eternal Nativity}。因此，天主子如此被称呼，其原则是：圣父也存在，祂从祂获得存在；而不是圣父在时间上先行，圣子在后面。圣父永远是，圣子永远来自圣父。并且，因为凡从另一事物而\Quote{存在}的，都是生的，因此圣子永远是生而有的。圣父永远是，那来自祂的肖像永远存在；正如那灌木的影像从灌木而生，如果灌木一直存在，那么影像也会一直从灌木而生。你无法找到与永恒生育者同永恒的被生者，但你找到了与时间中生育它们者同岁的被生者。我理解圣子与那生育祂的永恒者\OriginalTerm{同永恒}{coeternal}。因为在时间事物中称为同岁的，在永恒事物中就称为同永恒。\par

13. 兄弟们，这里有你们当思考的，以防备亵渎。因为常有人说：\Quote{看，你提出了一些相似；但从火发出的光辉，不如火本身明亮，灌木的影像也比那作为其影像的灌木本身具有更少的实在存在。这些例子有相似性，但没有完全的平等：因此它们似乎不是同一实体。}那么，若有人说：\Quote{那么圣父对于圣子，就如同光辉对于火，影像对于灌木？}我们该说什么呢？看，我已理解圣父是永恒的，圣子是与祂同永恒的；然而我们说祂如同从火发出且不如火明亮的光辉，或如同从灌木反射且不如灌木实在的影像吗？不，而是有完全的平等。\Quote{我不信，}他会说，\Quote{因为你没有发现一个相似。}那么，请相信宗徒，因为他能看见我所说的话。因为他说：\Quote{祂不以自己与天主同等为\OriginalTerm{僭夺}{robbery}。}\OriginalTerm{平等}{equality}是各方面完全相似。他说了什么？\Quote{不是僭夺。}为什么？因为那是属于别人的才是僭夺。\par

14. 然而，从这两个比较、这两种类比中，我们或许能在受造物内找到一种相似，藉以理解圣子如何与圣父同永，并且在任何方面都不比祂小。但我们无法单独从一种类比中找到这点：让我们将两种类比结合起来。哪两种？一种，是他们自己所举的相似例子；另一种，是我们所举的。因为他们所举的相似例子，来自那些在时间中出生、并且被生者所先在（在时间上）的事物，如人从人生。先出生的在时间上更大；但仍是人与人的相同实体。因为人生人，马生马，羊生羊。这些是按相同实体生育，却不在相同时间。它们在时间上不同，但在本性上并非不同。那么，我们在此生育中称赞什么？当然是本性的平等。但缺少了什么？时间的平等。让我们保留这里被称赞的一点，即本性的平等。但在我们提出的另一种类比中，从火的光辉和灌木的影像，你们找不到本性的平等，却找到了时间的平等。我们在这里称赞什么？时间的平等。缺少了什么？本性的平等。把你们所称赞的结合起来。因为在受造物中，缺少了你们所称赞的某些东西；在造物主中，没有什么可以缺少：因为你们在受造物中找到的，皆出自造物主之手。那么，在同时存在的事物中，岂不该将你们在这里所称赞的归于天主？但所缺少的，不可归于那至上威严，在祂内毫无缺陷。看，我向你们呈上同时存在的生育者与被生者：在这些中，你们称赞时间的平等，却指责本性的不平等。你们所指责的，不要归于天主；你们所称赞的，归于祂；如此，从这类相似中，你们将同时性（而不仅是同时同）归于祂，使圣子与生祂的圣父同永。但从另一类相似，它本身也是天主的受造物，应当赞美造物主，你们在其中称赞什么？本性的平等。你们先前因第一区分而将同永性归于祂；因这最后一点，将平等归于祂；于是同一实体的生育便完满了。因为，我的弟兄们，还有什么比这更疯狂：我在受造物中称赞某物，而那物在造物主内却不存在？在人身上我称赞本性的平等，难道在造人的天主那里，我不该相信吗？从人生的是人；从天主生的，难道不就是那生祂的天主本身？我与天主未曾创造的受造物毫无关系。那么，让造物主的一切受造物都赞美祂。我在一类中发现同时性，在另一类中得知同永性。在第一类中我找到本性的平等，在另一类中我理解实体的平等。因此，在这（天主）内，是\Quote{全然}，即那在别处分散于各个部分和各样事物中的。那么，在此是完全的全体，并且不仅是受造物中的；我在此找到它完全，但如同在造物主内那样，以更卓越的方式，因为一个是可见的，另一个是不可见的；一个是暂时的，另一个是永恒的；一个是可变的，另一个是不变的；一个是可朽的，另一个是不可朽的。最后，就人本身而言，我们所找到的人与人，是两个男人；而在这里，圣父与圣子是唯一的天主。\par

15. 我向我们的主天主献上说不尽的感谢，因祂俯允了你们的祈祷，使我的软弱从这极其复杂和困难之处得解脱。然而，首先请记住：造物主远远超越我们所能从受造物中搜集的一切，无论是通过感官还是心智的思考。但你想用心智达到祂吗？洁净你的心智，洁净你的心。使那能看见祂（无论祂是什么）的眼睛得以清净。因为\BibleQuote{心里洁净的人是有福的，因为他们要看见天主。}但心尚未洁净时，祂能更慈悲地提供和赐下的，岂不是那成为我们这样的圣言——我们曾对祂说了如此多而大的事情，却仍未说得相称——那\BibleQuote{万物是藉着祂而造成的}圣言成为我们所是的，使我们得以达到我们所不是的？因为我们不是天主，但藉着心智或内心的眼睛我们能看见天主。我们的眼睛被罪恶蒙蔽、被昏盲、因软弱而无力，渴望看见；但我们仍在希望中，尚未得着。我们是天主的子女。那说\BibleQuote{在起初已有圣言，圣言与天主同在，圣言就是天主}的若望这样说，他曾靠在主的胸膛上，从祂的心怀中汲取这些奥秘；他说：\BibleQuote{亲爱的，我们已是天主的子女，将来如何还未显明；我们知道，当祂显现时，我们要像祂，因为我们要看见祂实在怎样。}这就是所应许给我们的。\par

16. 但为了使我们能够达到——若我们还不能看见天主圣言，就让我们听取成了血肉的圣言吧；既然我们是属血肉的，就让我们听取降生成人的圣言吧。因为祂为此而来，为此承担了我们的软弱，好使你们能够接受一位承担你们软弱的天主的刚强言语。祂确实被称为\Quote{乳}。因为祂将乳给予婴孩，好将智慧之肉给予较成熟的人。那么，现在你们要耐心吮乳，好使你们能按心中最热切的愿望得喂养。因为那喂养婴孩的乳，是如何制成的？它岂不原先是在桌上的固体食物吗？但婴孩不够强壮去吃桌上的食物；母亲做了什么？她将食物转化为自己肉体的实质，并从中制成乳。为我们制成我们能够接受的。同样，圣言成了血肉，好使我们这些小人物——在食物方面确实是婴孩——能靠乳滋养。但有一个区别：当母亲将食物变成肉再变成乳时，食物变成了乳；而圣言却保持自身不变地取了血肉，好使两者仿佛交织在一起。祂之所是，祂没有腐化或改变，以便祂以你们的方式向你们说话，而非转变并变成人。因为祂保持不变、不改变、完全不可侵犯，就你们而言，祂成了你们所是的；就祂本身而言，祂在圣父面前仍是祂所是的。\par

17. 因为祂自己向软弱者说什么，为使他们恢复视力，得以在某种程度上达到那创造万物的圣言？\Quote{凡劳苦和负重担的，你们都到我这里来，我要使你们安息。你们背起我的轭，向我学习，因为我是良善心谦的。}师傅，天主子，天主的智慧，万物藉祂而造，祂在宣扬什么？祂呼唤人类说：\Quote{凡劳苦的，你们都到我这里来，向我学习。}你们或许以为天主的智慧会说：\Quote{学习我如何创造了诸天和星辰；万物在受造之前如何已在我内被数算；头发如何因不变的原理而被数算。}你们以为智慧会说这些和类似的事吗？不。而是先说了这句话：\Quote{我是良善心谦的。}弟兄们，看，这里你们可以领会的事，固然是小事。我们正走向大事，让我们接受小事，我们就将成为伟大的。你们要领会天主的崇高吗？先领会天主的谦卑。为了你们自己的缘故，屈尊就卑，因为天主也为了你们的缘故屈尊就卑了；祂不是为了自己的缘故。那么，领会基督的谦卑，学习谦卑，厌恶骄傲。承认你的软弱，耐心躺在医生面前；当你领会了祂的谦卑时，你就与祂一同上升；并非祂自己上升，因祂是圣言；而是你上升，好使祂越来越被你领会。起初你理解得犹豫不定，后来你将更确实、更清晰地理解。祂并没有增加，而是你进步了，祂仿佛与你一同上升。弟兄们，情况就是这样。相信天主的诫命，并付诸实践，祂将赐给你理解的力量。不要把最后的放在最先，也不要把知识置于天主的诫命之上，以免你们只是更低，却不更加根深蒂固。想想一棵树：它先向下扎根，然后向上生长；将根扎在低处，以便树梢伸向天空。它若不先从卑微开始，能努力生长吗？而你们若没有爱德，却想领会这些超越的事，如无根之树射向天空？这将是毁灭，而非成长。那么，当\Quote{基督}因信德住在你们心中时，你们要在爱德上根深蒂固，好使你们充满天主的一切圆满。\par

\SourceRecord{Translated by R.G. MacMullen.；Nicene and Post-Nicene Fathers, First Series；Vol. 6.；Edited by Philip Schaff.；Buffalo, NY: Christian Literature Publishing Co.,；1888.；<http://www.newadvent.org/fathers/160367.htm>.}

% structure: p0001 source=h1 target=section reason=page-title

\section{新约讲道集 第六十八篇}

\Emph{[CXVIII. Ben.]}\par

\Emph{论同一福音之言，\BibleBook{若望}{福音} 第一章，\Quote{在起初已有圣言}等。}\par

1. 你们众人寻求人的众多言语，却要认识天主的唯一圣言：\BibleQuote{在起初已有圣言。} 然而，\BibleQuote{在起初天主创造了天地。} 但是，\Quote{圣言是自有的，} 因为我们听到 \BibleQuote{在起初天主创造了。} 我们当在祂内承认创造者；因为创造者是那造物的，受造物是祂所造的。因为没有任何受造物是 \Quote{自有} 的，如同天主圣言 \Quote{自有} 一样，万物都是藉祂永远造成的。当我们听到 \Quote{圣言是自有的} 时，祂与谁同在？我们理解圣父并没有制造或创造这同一位圣言，而是生了祂。因为 \BibleQuote{在起初天主创造了天地。} 祂藉什么创造了它们？\BibleQuote{圣言与天主同在；} 但这是何种圣言？它是否发出了声音而后消逝？它是否只是一种思想、一种心灵的动态？不是。它是否由记忆所提示而出于口？不是。那么是何种圣言？你为什么从我这里寻求许多言语？\BibleQuote{圣言就是天主。} 当我们听到 \BibleQuote{圣言就是天主} 时，我们并不制造第二个天主；而是明认圣子。因为圣言是天主之子。看，圣子，除了是天主还能是什么？因为 \BibleQuote{圣言就是天主。} 圣父是什么？自然是天主。如果圣父是天主，圣子是天主，我们是否制造了两个天主？绝对不可。圣父是天主，圣子是天主；但圣父与圣子是一个天主。因为天主的独生子不是受造的，而是受生的。\BibleQuote{在起初天主创造了天地；} 但圣言出自圣父。那么圣言是由圣父创造的吗？不是。\BibleQuote{万物是藉着祂造成的。} 如果万物是藉着祂造成的，那么祂自己也由祂自己造成吗？不要想象那你们所听到的造成万物的那一位，祂自己竟是在万物中受造的。因为如果祂自己受造，那么万物就不是由祂造成，而是祂自己也在其余受造物中受造了。你说：\Quote{祂是受造的；} 什么，由祂自己？谁能造自己？如果祂是受造的，万物又怎能由祂造成？看，连祂自己也受造，如你所说，不是我说的，因为祂是受生的，我并不否认。如果你说祂是受造的，我问：藉什么，藉谁？由祂自己？那么祂 \Quote{是}，在祂受造之前，为要造自己。但是如果万物是藉着祂造成的，那么当知道祂自己并非受造。如果你们不能领悟，要相信，好使你们能领悟。信德在先；领悟在后；因为先知说：\BibleQuote{除非你们相信，你们必不能理解。} \Quote{圣言是自有的。} 但你们说：\Quote{有过一段时间，圣言不存在。} 你们说谎；你们在任何地方都读不到这个。但我为你们读到：\BibleQuote{在起初已有圣言。} 你们在起初之前寻找什么？如果你们能在起初之前找到任何东西，那便是起初。凡在起初之前寻找任何东西的人，是疯狂的。那么他说什么东西在起初之前？\BibleQuote{在起初已有圣言。}\par

2. 但你们要说：\Quote{圣父\Quote{是}，并且是在圣言之前。} 你们在寻求什么？\BibleQuote{在起初已有圣言。} 你们所找到的，要明白；不要寻求你们不能找到的。在起初之前没有一物。\BibleQuote{在起初已有圣言。} 圣子是圣父的光辉。论到圣父的智慧，即圣子，经上说：\BibleQuote{因为祂是永恒之光的辉映。} 你们寻求一个没有圣父的圣子吗？给我一个没有光辉的光。如果有一段时间圣子不存在，圣父便是昏暗的光。因为如果光明没有光辉，它怎能不是昏暗的光？因此，圣父永远存在，圣子永远存在。如果圣父永远存在，圣子也永远存在。你们问我，圣子是否受生？我回答：\Quote{受生。} 因为如果不是受生，祂就不是圣子。所以当我说圣子永远存在时，我实际上是说祂永远是受生的。谁能明了 \Quote{永远受生}？给我一个永恒的火，我就给你一个永恒的光辉。我们赞美天主，祂赐给了我们圣经。不要在这光的辉映中成为瞎眼的。光辉由光而生，然而光辉与生它的光同永恒。光永远存在，它的光辉也永远存在。光生了它自己的光辉；但它曾否没有它的光辉？让天主被允许生一个永恒的圣子。我恳求你们听我们在说谁；要听，要注意，要相信，要明白。我们是在谈论天主。我们宣认并相信圣子与圣父同永恒。但你们要说：\Quote{当人生子时，生者是年长的，所生的是年幼的。} 这话不错；在人方面，生者是年长的，所生的是年幼的，并且他渐渐长大达到父亲的力气。但为什么呢？岂不是因为一个在成长，另一个在衰老？让父亲暂时静止，在儿子成长时，他将会追赶父亲，你将看到他们相等。但是看，我给你一个理解之方。火产生同代的光辉。在人中，你只找到幼子、老父；你找不到他们同代：但正如我所说的，我指给你看光辉与生它的火同代。因为火产生光辉，却从没有不带光辉。既然你看到光辉与火同代，就让天主生一个同永恒的圣子吧。谁明白，就让他喜乐；谁不明白，就让他相信。因为先知的话不会落空：\BibleQuote{除非你们相信，你们必不能理解。}\par

\SourceRecord{Translated by R.G. MacMullen.；Nicene and Post-Nicene Fathers, First Series；Vol. 6.；Edited by Philip Schaff.；Buffalo, NY: Christian Literature Publishing Co.,；1888.；<http://www.newadvent.org/fathers/160368.htm>.}

% structure: p0001 source=h1 target=section reason=page-title

\section{论新约 讲道词第六十九篇}

\Emph{[CXIX. Ben.]}\par

\Emph{论同一句话，\BibleBook{若望}{福音}1。\BibleQuote{在起初已有圣言，}等等。}\par

1. 我们的主耶稣基督为寻找失丧的人而成了人，我们的宣讲从未隐瞒，你们的信德也一直持守；此外，这位为了我们而成了人的主，始终与圣父同是天主，并将永远如此，不，他始终是\Quote{是}；因为那里没有时间的先后，就没有\Quote{曾存在}和\Quote{将存在}。因为所谓\Quote{曾存在}的那物，现已不再存在；所谓\Quote{将存在}的那物，尚未存在；但他始终是\Quote{是}，因为他确实是\Quote{是}，即是不变者。因为福音选段刚刚教导我们一个崇高而神圣的奥迹。因为圣若望倾注了这福音的开端，是因为他从主的胸怀领受了它。你们记得，最近刚刚向你们读过，圣若望宗徒如何侧卧在主怀中。为清楚解释这一点，他说：\BibleQuote{在主胸前}；好让我们明白他所谓\BibleQuote{在主怀中}的意思。我们认为，那侧卧在主胸前的人，喝下了什么呢？不，让我们不去思想，而是去喝；因为我们刚刚也听到了我们可以喝下的东西。\par

2. \BibleQuote{在起初已有圣言，圣言与天主同在，圣言就是天主。}啊，光荣的宣讲！啊，主胸怀丰盛宴席的成果！\BibleQuote{在起初已有圣言。}你为什么寻求在此之前有什么？\BibleQuote{在起初已有圣言。}如果圣言是被造的（因为借以造万物的那一位确实不是被造的）——如果圣言是被造的，圣经就会说：\BibleQuote{在起初天主造了圣言，}正如\BibleBook{}{创世记}所说：\BibleQuote{在起初天主创造了天地。}所以天主并没有在起初造圣言，因为\BibleQuote{在起初已有圣言。}这起初就存在的圣言，在哪里呢？往下读：\BibleQuote{圣言与天主同在。}但我们因天天听见人的言语，便轻看了这\Quote{圣言}的名号。在此不要轻看\Quote{圣言}的名号：\BibleQuote{圣言就是天主。这位圣言在起初与天主同在。万物是借着他造的；凡受造的，没有一样不是借着他造的。}\par

3. 敞开你们的心怀，弥补我言辞的贫乏。凡我能表达的，请侧耳听；凡我不能表达的，请默想。谁能领悟那常存的圣言？所有的话语都发出声音，然后消逝。谁能领悟那常存的圣言，除非那住在他内的人？你想要领悟那常存的圣言吗？不要追随肉身的潮流。因为这肉身确实是一个潮流，因为它没有常存之物。从某种隐秘的自然泉源，人们出生、生活、死亡；他们从何而来，往何而去，我们不知道。它是隐藏的水，直到它从源头涌出；它流去，在过程中可见；然后又隐没在海中。让我们鄙视这条溪流——流去、奔流、消失——让我们鄙视它。\BibleQuote{凡有血肉的都似草，而他的一切荣美都似草上的花；草必枯萎，花必凋谢。}你要持久吗？\BibleQuote{但主的圣言却永远常存。}\par

4. 但为救助我们，\BibleQuote{圣言成了血肉，居住在我们中间。}\BibleQuote{圣言成了血肉}是什么意思？金子变成了草。它变成草是为了被焚烧；草被焚烧，但金子仍存留；在草中它不会消亡，反而改变了草。它怎样改变了草？它使之复活，使之获得生命，将它提升到天上，放在圣父的右边。但为使我们可以说\BibleQuote{圣言成了血肉，居住在我们中间，}让我们回想一下之前的话。\BibleQuote{他来到自己的地方，自己的人却没有接受他。但是凡接受他的，他赐给他们权能，成为天主的子女。}\BibleQuote{成为，}因为他们\BibleQuote{不是}；但他自己在起初就\BibleQuote{是。}\BibleQuote{他赐给他们权能，成为天主的子女，就是信他名的人；他们不是由血，也不是由肉欲，也不是由男欲，而是由天主生的。}看，他们诞生了，不管在肉身的哪个年龄；你们看见婴孩；看见了，就喜乐吧。看，他们诞生了；但他们是由天主生的。他们母亲的子宫就是圣洗之水。\par

5. 愿没有人怀着心灵的贫乏抱持这种看法，在心中反复思考这种最可怜的念头，对自己说：\Quote{怎么\BibleQuote{在起初已有圣言，圣言与天主同在，圣言就是天主：万物是藉着他造的}；而看，\BibleQuote{圣言成了血肉，居住在我们中间}？}听一听为什么这样做。\BibleQuote{凡接受他的}我们知道\BibleQuote{信他的人，他赐给他们权能，成为天主的子女。}那么，不要让那些他赐给权能成为天主的子女的人，认为成为天主的子女是不可能的。\BibleQuote{圣言成了血肉，居住在我们中间。}不要以为你们成为天主的子女是一件太过伟大的事；因为为了你们，那原是天主之子者成了人子。如果他成为较小的，而他原是较大的，难道他不能使我们这原先较小的，变成更为伟大吗？他下降到我们这里，难道我们不应上升到到他那里吗？他为了我们接受了我们的死亡，难道他不会赐给我们他的生命吗？他为你忍受了你邪恶的事物，难道他不会赐给你他的美善事物吗？\par

6. \Quote{但是，}有人会说，\Quote{天主的圣言——万物借祂而治理，万物借祂而受造，过去如此，现在也如此——怎么会将自己收缩到一位童贞女的胎中？怎么会舍弃世界，离开天使，被关闭在一个女人的子宫里呢？}你不懂如何领会属神的事。天主的圣言（人哪，我是在对你说话，我是在对你讲论天主的圣言的无所不能）确实能做一切，因为天主的圣言是无所不能的：同时留在圣父那里，又来到我们这里；同时以肉身来到我们当中，又隐藏于祂内。因为即使祂没有从肉身而生，祂也不会因此有所减少。祂\Quote{是}存在于自己的肉身之前；祂创造了祂自己的母亲。祂选择了她，使她在祂内受孕；祂创造了祂，使祂由她受造。你为什么惊奇？我是在对你谈论天主：\BibleQuote{圣言就是天主。}\par

7. 我正在论述圣言，也许人的言语可以提供某种类似的东西；虽然极不相等，相差甚远，完全无法相比，却仍能以相似的方式向你传递一点提示。看，我对你所说的言语，先前就在我心里；它出来到你那里，却并未离开我；那在你内开始的，原本不在你内；它出来到你那里时，仍与我同在。正如我的言语被带到你的感官，却没有离开我的心；同样，那圣言来到我们的感官，却没有离开祂的圣父。我的言语与我同在，并出来成为声音；天主的圣言与圣父同在，并出来成为肉身。但我能用我的声音做到祂用祂的肉身所能做的吗？因为我不能支配我的声音在它飞逝之后；祂不仅支配祂的肉身，使它出生、生活、行动，甚至当它死后，祂还复活了它，并将那仿佛祂来我们这里时所用的运载工具高举到圣父那里。你可以称基督的肉身为一件衣服，你可以称它为运载工具，或许正如祂亲自屈尊教导我们的，你可以称它为祂的坐骑；因为在这坐骑上，祂扶起了那个被强盗打伤的人；最后，正如祂自己更明确地说的，你可以称它为圣殿；这圣殿不再经历死亡，它的位置在圣父的右边；在这圣殿中，祂将要来审判活人死人。祂以训诲所教导的，以榜样彰显了。祂在自己肉身中所显示的，你也应当在你的肉身中盼望。这就是信德；紧紧抓住你目前还看不见的。有必要，使你借着相信而坚定于你所看不见的；免得你看见时，反觉羞愧。\par

\SourceRecord{Translated by R.G. MacMullen.；Nicene and Post-Nicene Fathers, First Series；Vol. 6.；Edited by Philip Schaff.；Buffalo, NY: Christian Literature Publishing Co.,；1888.；<http://www.newadvent.org/fathers/160369.htm>.}

% structure: p0001 source=h1 target=section reason=page-title

\section{新约讲道集 第七十篇}

\Emph{【CXX. Ben.】}\par

\Emph{论\BibleBook{若望}{福音}第1章同样的词句}，\Quote{在起初已有圣言}\Emph{，等等。}\par

1. \BibleBook{若望}{福音}的开端：\BibleQuote{在起初已有圣言}。他这样开始，他看见了这事，并超越整个受造界，山岭、空气、诸天、星辰、上座者、宰制者、率领者、掌权者、所有天使和总领天使，超越一切；他看见了在起初的圣言，并畅饮了它。他看见在一切受造物之上，他从主的胸前畅饮。因为这位圣\Person{若望}宗徒正是\Person{耶稣}特别爱的那一位，以至于晚餐时他曾靠在祂的胸膛上。内中隐藏着这奥秘，为的是从那胸中汲取，然后在福音中涌流出来。听见又明白的人是有福的。次一等有福的，是那些虽不明白却相信的人。因为看见这天主的圣言是多么伟大的事，谁能用人言解释呢？\par

2. 你们要举心向上，我的弟兄们，尽你们所能地举起它们；凡是根据任何有形物体所想到的，都要摒弃。如果你们想到天主的圣言时，以这太阳的光为联想，那么无论你们怎样扩大、延伸，在思想中不为那光设限，它仍与天主的圣言无可比拟。要认为圣言处处都是完整的。你们要明白我所说的；由于时间紧迫，我尽量为你们简短而言。看，这来自天上的光，称为太阳，它出来时，照亮大地，展开白昼，显明形貌，分别颜色。这是极大的祝福，是天主赐给所有凡人的伟大恩赐；愿祂的化工赞美祂。太阳如此美丽，那么太阳的创造者岂不更加美丽？然而请注意，弟兄们；看，它向整个地球倾泻光芒；穿透敞开之处，关闭之处阻隔它；它透过窗户送进光，但能透墙而入吗？对于天主的圣言，一切都是敞开的，没有隐藏的。注意另一个区别，受造物，特别是物质受造物，离创造者有多远。当太阳在东方时，它不在西方。它的光从这巨大的天体发出，确实能到达西方；但它本身不在那里。当它开始落下时，它就会在那里。它升起时在东方；落下时在西方。因着这些运行，它赋予了那些区域名称。因为它在东方升起，故被称为\Quote{旭日}；它在西方落下，故被称为\Quote{落日}。夜间它无处可见。天主的圣言也是如此吗？它在东方时，不在西方吗？或者它在西方时，不在东方吗？它是否曾离开大地，下到或转到地后？它处处都是完整的。谁能用言语解释这个？谁能看见它？我用什么证据向你们证实我所说的话呢？我是作为一个人讲话，是向人讲话；我是作为软弱者讲话，是向更软弱的人讲话。然而，我的弟兄们，我敢说，我确实在某种程度上看见了我对你们所说的，虽然\Quote{通过镜子}，或\Quote{模糊不清}，我也在某种程度上在心中领悟到关于这事的一个话。但它想要出去到你们那里，却找不到合适的载体。话的载体是声音。我在我自己内所讲的，我想讲给你们，但言语不足。因为我想谈论天主的圣言。这是何等伟大的话，是怎样的话？\BibleQuote{万物是藉着祂而造成的。}看看这些工程，敬畏那位工作者。\BibleQuote{万物是藉着祂而造成的。}\par

3. 回来吧，人的软弱，我说，回来吧。让我们若能理解这些人间的事吧。我们是人，我说话的是一个人，向人说话，发出我的声音。我把声音送入人的耳朵，借着声音也通过耳朵将理解存入心里。那么，让我们就这一点尽量说说，尽量理解它。但如果我们连这个都不能理解，对于那\Quote{另一位}我们算什么？看，你们在听我说话；我正讲一个话。如果有人从我们中间出去，在外面被问这里在做什么，他回答：\Quote{主教在说话。}我正讲一个关于圣言的话。但这一个话，是关于怎样的圣言呢？一个会死的话，关于不朽的圣言；一个可变的话，关于不变的圣言；一个短暂的话，关于永恒的圣言。然而，请考虑一下我的话。因为我已告诉你们，天主的圣言处处都是完整的。看，我正向你们说话；我所说的到达了所有人。现在，我所说的话能传达给你们众人，你们是否分了我的话？如果我是喂你们，要充满的不是你们的理智，而是你们的身体，把饼放在你们面前使你们饱饫；你们岂不把我的饼分给你们？我的饼能来到你们每个人吗？如果它只来到一个人那里，其余的人就没有了。但现在看，我说话，你们全都领受。不止全都领受，而且全都完整地领受。它完整地来到所有人，给每个人是完整的。啊，我这言语的奇迹！那么，天主的圣言是什么呢？再听：我讲了话；我所讲的，已经出去到你们那里，并没有离开我。它到达你们，却没有与我分离。我讲之前，我拥有它，而你们没有；我讲了，你们开始拥有，而我什么也没有失去。啊，我这言语的奇迹！那么，天主的圣言是什么？从小事来推测大事。想想地上的事，赞美天上的事。我是一个受造物，你们是受造物；我的言语在我心里，在我口中，在我声音里，在你们耳朵里，在你们心里，行出如此伟大的奇迹。那么，创造者是谁？主啊，求祢垂听我们。祢既然造了我们，求祢塑造我们。祢既然造了我们为受光照的人，求祢使我们成为善。这些身着白衣的受光照者，通过我听到祢的圣言。因为他们因祢的恩宠受光照，站在祢面前。\BibleQuote{这是主所安排的一天。}只求他们劳苦，为此祈祷，好使这些日子过去之后，那已成为天主奇妙与祝福之光的人，不致变成黑暗。\par

\SourceRecord{Translated by R.G. MacMullen.；Nicene and Post-Nicene Fathers, First Series；Vol. 6.；Edited by Philip Schaff.；Buffalo, NY: Christian Literature Publishing Co.,；1888.；<http://www.newadvent.org/fathers/160370.htm>.}

% structure: p0001 source=h1 target=section reason=page-title

\section{新约讲道71}

\Emph{[CXXI. Ben.]}\par

\Emph{论福音的话}，\BibleRef{若 1:10}，\Emph{、\Quote{世界是藉着他造的}等。}\par

1. 主\Quote{造了世界，世界却不认识他}。哪一个世界是藉着祂造的，哪一个世界不认识祂？因为藉着祂所造的世界，并不是不认识祂的那个世界。藉着祂造的世界是什么？天和地。天如何不认识祂？当祂受难时，太阳昏暗了。地如何不认识祂？当祂悬在十字架上时，地震动了。但是\Quote{世界不认识他}，这世界的首领就是他，关于这首领曾说：\Quote{看，这世界的首领来了，他在我身上一无所获。}恶人被称为世界；不信的人被称为世界。他们因所爱之物而得名。藉着对天主的爱，我们成为神；同样，藉着对世界的爱，我们被称为世界。但\Quote{天主在基督内使世界与自己和好}。因此\Quote{世界不认识他}。什么？\Quote{所有的人？}\par

2. \Quote{他来到自己的领域，自己的人却没有接受他。}万物都是祂的，但那些人被称为\Quote{自己的人}，从他们中祂的母亲出生，祂在他们中取了肉身，祂曾派遣前驱预报祂的来临，祂曾赐给他们法律，祂曾从埃及奴役中拯救他们，祂曾按肉身拣选了他们的祖先亚巴郎。因为祂曾说真理：\Quote{在亚巴郎以前，我就有。}祂并没有说：\Quote{在亚巴郎以前，}或\Quote{在亚巴郎被造以前，我被造。}因为\Quote{在起初已有圣言}，并非\Quote{被造}。所以\Quote{他来到自己的领域}，祂来到犹太人中间。\Quote{自己的人却没有接受他。}\par

3. \Quote{但是，凡接受他的，}因为宗徒们当然在那里，他们\Quote{接受了他}。有那些在祂驴驹前举着树枝的人。他们前行后随，铺开衣服，大声呼喊：\Quote{贺三纳于达味之子！奉主名而来的当受赞美！}法利塞人对祂说：\Quote{约束这些孩子，不要让他们这样向你呼喊。}祂却说：\Quote{这些人若不作声，石头就要喊叫了。}祂说这些话时看见了我们；\Quote{这些人若不作声，石头就要喊叫了。}谁是石头？岂不是那些敬拜石头的人？如果犹太人的孩子不作声，年长的和年幼的外邦人就要呼喊。谁是石头？岂不是那位来\Quote{为光作证}的若翰所说的人？因为当他看见这些自夸出于亚巴郎的犹太人时，他对他们说：\Quote{毒蛇的种类。}他们自称为亚巴郎的子孙；他却对他们说：\Quote{毒蛇的种类。}他冤枉了亚巴郎吗？决不是！他按他们的品性给他们起了名字。因为如果他们真是亚巴郎的子孙，就会效法亚巴郎；正如主对那些对祂说\Quote{我们是自由的，从未做过任何人的奴隶；我们有亚巴郎为父}的人所说的话。祂说：\Quote{如果你们是亚巴郎的子女，你们必行亚巴郎的事。你们却想杀我，因为我给你们讲真理；亚巴郎并没有做过这样的事。}你们是他的后代，却是堕落的后代。那么，若翰说了什么？\Quote{毒蛇的种类，谁指教你们逃避那即将来临的愤怒？}因为他们来受若翰的悔改洗礼。\Quote{谁指教你们逃避那即将来临的愤怒？你们要结出与悔改相称的果实来。心里不要想：我们有亚巴郎为父。我告诉你们：天主能从这些石头中给亚巴郎兴起子孙来。}因为天主能从这些石头中——他在神视中所见的——兴起子孙来；他对他们说话，他预见了我们：\Quote{天主能从这些石头中给亚巴郎兴起子孙来。}哪些石头？\Quote{这些人若不作声，石头就要喊叫了。}你们刚才听见了，并喊叫了。这话应验了：\Quote{石头要喊叫。}因为我们来自外邦人，我们的祖先曾崇拜石头。因此我们也被称为狗。请回想那向主呼喊的妇人所听到的话，因为她是一个客纳罕妇人，崇拜偶像，魔鬼的仆人。耶稣对她说：\Quote{拿儿女的饼扔给小狗吃是不对的。}你们难道没有注意到狗如何舔油腻的石头？所有崇拜偶像的人都是如此。但是恩宠已临到你们。\Quote{但是，凡接受他的，他给他们权能，成为天主的子女。}看，这里有些刚刚出生的人：祂给了他们\Quote{成为天主子女的权能}。祂给了谁？\Quote{给那些信他名字的人。}\par

4. 他们如何成为天主的子女？\Quote{他们不是由血，也不是由肉欲，也不是由男欲，而是由天主生的。}他们获得成为天主子女的权能，由天主而生。请注意：他们由天主而生，\Quote{不是由血}，像他们的第一次出生那样，那悲惨的出生，出自悲惨。但是由天主而生的人，他们曾是什么？他们最初是如何出生的？由血；由男女的血共同结合，由男女肉身的结合，就这样出生。现在呢？他们由天主而生。第一次出生来自男女；第二次出生来自天主和教会。\par

5. 看，他们由天主而生；这如何成就呢？他们本是由人而生，却要由天主而生？这如何成就，如何成就？\BibleQuote{圣言成了肉躯，为能居住在我们中间。}奇妙的交换：祂成了肉躯，他们成了灵。这是什么？何等屈尊就卑，我的弟兄们！请将你们的心志提升，去盼望并领会更美的事。不要顺从世俗的欲望。\BibleQuote{你们是重价买来的；}为了你们的缘故，圣言成了肉躯；为了你们的缘故，那位本是天主之子的，成了人子：为使你们这些本是人的儿子的，得以成为天主的儿子。祂是什么，祂成了什么？你们是什么，你们成了什么？祂本是天主之子。祂成了什么？人子。你们本是人的儿子。你们成了什么？天主的儿子。祂与我们分享我们的不幸，为要把祂的美好赐给我们。但即使在祂成为人子这件事上，祂也与我们大不相同。我们成为人的儿子是因肉身的欲望；祂成为人子是因童贞女的信德。任何其他人的母亲都是通过肉体的结合而怀孕；每个人都是由人的父母，即父亲和母亲所生。但基督是由圣神和童贞玛利亚所生。祂来到我们中间，却从未远离祂自己；是的，祂作为天主从未离开祂自己；但祂将祂所是的加给了我们的本性。因为祂来到了祂所不是的，却没有失去祂所是的。祂成了人子，但并未停止是天主子。因此祂是中保，在中间。什么是\Quote{在中间}？既不在上，也不在下。怎能既不在上也不在下？不在上，因为祂是肉躯；不在下，因为祂不是罪人。然而，就祂是天主而言，祂总是在上。因为祂来到我们这里，并没有离开圣父。祂从我们这里去了，却没有离开我们；祂将再次来到我们这里，也不会离开祂。\par

\SourceRecord{Translated by R.G. MacMullen.；Nicene and Post-Nicene Fathers, First Series；Vol. 6.；Edited by Philip Schaff.；Buffalo, NY: Christian Literature Publishing Co.,；1888.；<http://www.newadvent.org/fathers/160371.htm>.}

% structure: p0001 source=h1 target=section reason=page-title

\section{新约讲道集 第72篇}

\Emph{[CXXII. Ben.]}\par

\Emph{论福音的话}，\BibleRef{若 1:48} \Emph{，\BibleQuote{当你还在无花果树下时，我就看见了你，}等等。}\par

1. 我们的主耶稣基督对\Person{纳塔乃耳}所说的话，若我们正确理解，便不仅关乎他一人。因为我们的主耶稣看见整个人类种族都在无花果树下。因为在经文中，无花果树在此象征罪。但并非总是如此，只是在此处，根据象征的适宜性，如你们所知，原祖犯罪后，曾用无花果树叶子遮盖自己。因为他们犯罪后，便用这些叶子遮盖自己的裸体，为自己的罪而羞耻；而天主为他们造作肢体，他们却为自己制造了羞耻的缘由。其实他们不必为天主的工程而羞耻；是罪的起因先于羞耻。如果罪愆未曾先行，裸体决不会令人脸红。因为\BibleQuote{他们赤身露体，并不羞耻。}因为他们尚未犯下任何可羞耻的事。但我为何说这一切？为使我们明白无花果树象征罪。那么，\BibleQuote{当你还在无花果树下时，我就看见了你}是什么意思？就是：当你还在罪下时，我就看见了你。\Person{纳塔乃耳}回顾所发生的事，想起自己曾在一棵无花果树下，而基督并不在那里。祂不在那里，即按肉身临在不在此；但按圣神内的认识，祂何处不在？因为他知道自己独自在无花果树下，而主基督并不在那里；当主对他说：\BibleQuote{当你还在无花果树下时，我就看见了你}，他便认出了祂的天主性，并喊道：\BibleQuote{你乃以色列的君王。}\par

2. 主说：\BibleQuote{因为我对你说了：你还在无花果树下时，我就看见了你，你就惊奇吗？你将要看见比这些更大的事。}这些更大的事是什么？祂又说：\BibleQuote{你将看见天敞开，天主的众天使在人子身上上去下来。}让我们回忆圣书中所记载的古老故事。我说的是\WorkTitle{创世纪}。当\Person{雅各伯}在某处睡觉时，他把一块石头放在头下；在梦中他看见一个梯子从地直立到天，主靠在上面；众天使在梯子上上去下来。\Person{雅各伯}看见了这一切。一个人的梦不会被记录下来，若非其中预表了某种伟大的奥迹，若非在那异象中应理解某种伟大的预言。因此，\Person{雅各伯}本人因为明白他所见的事，就在那里立了一块石头，并傅油在上面。现在你们认出这傅油；也认出那位受傅者。因为祂就是\BibleQuote{建筑者弃而不用的石头，却成了屋角的基石。}祂就是那块石头，祂自己曾说：\BibleQuote{凡跌在这石头上的，必被粉碎；但这石头落在谁身上，必要压碎他。}当祂躺在地上时，人跌在祂上面；但那石头从高处落下，击打人，便是当祂从天降来审判生者死者之时。犹太人有祸了，因为当基督谦卑降临时，他们却跌在祂身上。他们说：\BibleQuote{这人不是从天主来的，因为他不守安息日。} \BibleQuote{如果他是天主子，就让他从十字架上下来。}狂人啊，石头躺在地上，你便讥笑它。但你既讥笑它，你便瞎眼；你既瞎眼，你便跌在上面；你既跌在上面，你便被它摇动；它如今躺在地上，你既被它摇动，将来它从高处落下，你必被它压碎。于是\Person{雅各伯}傅油于石头。他以此立为偶像吗？他在其中表明了意义，但并未朝拜它。现在请你们侧耳，听\Person{纳塔乃耳}的事，主耶稣籍祂的机缘，乐意为我们解释\Person{雅各伯}的异象。\par

3. 你们这些在\Person{基督}的学校里受过良好教导的人，知道这\Person{雅各伯}也就是\Person{以色列}。这是两个名字，却是同一个人。他最初的名字\Person{雅各伯}，意为\Quote{取代者}，是在出生时得的。因为那对孪生子出生时，他的兄弟\Person{厄撒乌}先出生，而弟弟的手正抓着哥哥的脚后跟。他抓住了在他之前出生的兄弟的脚后跟，而自己随后出来。由于这件事，因他抓住了他兄弟的脚后跟，他便被称为\Person{雅各伯}，即\Quote{取代者}。后来，当他从\Place{美索不达米亚}返回时，有一位天使在路上与他角力。天使的力量与人的力量如何相比？因此这是一个奥秘、一件圣事、一个预言、一个预象；让我们来理解它。也要思考角斗的方式。当他角斗时，\Person{雅各伯}战胜了天使。这里有深意。当人战胜了天使后，他紧紧抓住祂；是的，这战胜者紧紧抓住了被他征服的那一位。并对祂说：\Quote{除非你祝福我，否则我不放你走。}当征服者被被征服者祝福时，基督便被预象了。因此，那位被认为是\Person{主耶稣}的天使对\Person{雅各伯}说：\Quote{你不可再叫雅各伯，要叫以色列，}这名字的意思是\Quote{看见天主}。此后，祂触摸了他的大腿筋，即大腿的宽处，它就枯干了；\Person{雅各伯}就瘸了。那被征服的就是这样。这被征服者拥有如此大能，一摸大腿就使人瘸了。祂当时是甘心被征服的。因为祂\BibleQuote{有能力放下祂的力量，也有能力再取回来}。祂不因被征服而发怒，因为祂也不因被钉在十字架上而发怒。祂甚至祝福了他，说：\Quote{你不可再叫雅各伯，要叫以色列。}于是\Quote{取代者}被变成了\Quote{看见天主的人}。祂如前所述触摸了他的大腿，使他瘸了。要在\Person{雅各伯}身上观察犹太人民，那成千上万跟随并走在主驴驹之前的人，他们与众宗徒一起朝拜主，并高呼：\BibleQuote{贺三纳于达味之子！奉主名而来的当受赞美！}看，\Person{雅各伯}受了祝福。他直到如今在这些今日的犹太人中仍然瘸着。因为大腿的宽处象征人数的增多。关于这些人，圣咏预言外邦人将相信时说道：\BibleQuote{我不认识的民族，竟服事了我；他们一听见就服从了我。}我不在那里，却被听见；我在这里，却被杀害。\BibleQuote{我不认识的民族，竟服事了我；他们一听见就服从了我。}因此，\BibleQuote{信德出于听闻，听闻则凭基督的话。}接着又说：\BibleQuote{陌生的儿子们对我说了谎话；}这是指犹太人。\BibleQuote{陌生的儿子们对我说了谎话；陌生的儿子们衰败萎缩，行走蹒跚。}我已经向你们指出了\Person{雅各伯}，\Person{雅各伯}受了祝福，\Person{雅各伯}又瘸了。\par

4. 但由此引出的一个问题不可放过，它或许会困扰你们中的一些人：为何当这\Person{雅各伯}的祖父\Person{亚巴辣罕}的名字被改变时（他最初也叫\Person{亚巴郎}，天主改变了他的名字，说：\BibleQuote{你不可叫亚巴郎，要叫亚巴辣罕}），从那以后他就不再被称为\Person{亚巴郎}了。你们去查考圣经，就会看到，在得到新名字之前，他只被称为\Person{亚巴郎}；得到之后，他只被称为\Person{亚巴辣罕}。但这\Person{雅各伯}得到另一个名字时，也听到了同样的言词：\BibleQuote{你不可叫雅各伯，要叫以色列。}你们查考圣经，就会看到他总是有两个名字，既叫\Person{雅各伯}也叫\Person{以色列}。\Person{亚巴郎}得到新名后只被称为\Person{亚巴辣罕}；\Person{雅各伯}得到新名后却既叫\Person{雅各伯}又叫\Person{以色列}。\Person{亚巴辣罕}的名字要在今世发展；因为在此他被立为万民之父，故得此名。但\Person{以色列}的名字关联另一个世界，在那里我们将看见天主。因此天主的人民，即现今的基督徒子民，既是\Person{雅各伯}也是\Person{以色列}：实际上是\Person{雅各伯}，在希望中是以\Person{以色列}。因为年幼的民族被称为其年长民族的\Quote{取代者}。什么？我们取代了犹太人？不，而是说我们是他们的取代者，因为为了我们的缘故他们被取代了。如果他们没有被蒙蔽，\Person{基督}就不会被钉在十字架上；祂的宝血就不会流出；如果那血没有流出，世界就不会被救赎。因为他们的蒙蔽便于我们得益，所以年长的兄弟被年幼的取代了，年幼的被称为\Quote{取代者}。但这要到何时为止呢？\par

5. 时候将到，世界的终穷将到，全以色列都要相信；不是现今的这些人，而是那时他们的子孙。因为现今这些行走在自己道路中的，要往自己的地方去，将进入永罚。但当他们成为同一子民时，我们所歌唱的就要实现：\Quote{当祢的光荣显现时，我必心满意足。}那时赐予我们的恩许，即我们\Quote{面对面}看见，将要来到。\Quote{如今我们藉着镜子观看，模糊不清，}且\Quote{仅部分；}但当两子民如今净化、如今复活、如今加冕、如今改变成不朽的形态和永恒的不腐坏时，要面对面看见天主，\Person{雅各伯}不再存在，只有\Person{以色列}；那时主要在这一位圣\Person{纳塔乃耳}身上看见他，并说：\Quote{看，这确是一个以色列人，在他内毫无诡诈。}当你们听见\Quote{看，一个以色列人；}愿以色列进入你们心中；当以色列进入你们心中时，愿他的梦进入你们心中，就是他看见一个梯子从地直达天，主站在上面，天主的众天使上升下降的那个梦。这个梦是\Person{雅各伯}看见的。但此后他被称为以色列；即稍后当他从\Place{美索不达米亚}归来，在旅程中。如果\Person{雅各伯}看见了梯子，他又被称为以色列；而这\Person{纳塔乃耳}是一个\Quote{在他内毫无诡诈的以色列人}；因此当他因主对他说\Quote{我看见你在无花果树下}而惊奇时，祂对他说：\Quote{你将看见比这更大的事。}于是祂向他宣告了\Person{雅各伯}的梦。祂向谁宣告？向祂称为\Quote{一个毫无诡诈的以色列人}的那一位。好像祂说：我用以称呼你的他的梦，要显示在你身上；不要急于惊奇，\Quote{你将看见比这更大的事。你将看见天开，天主的众天使在人子身上上升下降。}看\Person{雅各伯}所看见的；看\Person{雅各伯}为何用油傅抹石头；看\Person{雅各伯}为何先知性地预示和预表了受傅者。因为那行动是一个预言。\par

6. 现在我知道你们在等待什么；我明白你们想听我说什么。我也要简短地宣告，按主的加能；\Quote{在人子身上上升下降。}如何——如果他们下降到祂那里，祂在这里；如果他们上升到祂那里，祂在上。但如果他们上升到祂那里，又下降到祂那里，祂同时在上和在这里。绝不可能，除非祂既在他们上升之处，又在这里他们下降之处——我们如何证明祂既在那里，又在这里？让\Person{保禄}——他先是\Person{撒乌耳}——回答我们。他亲身体验，他先是迫害者，后来成为宣讲者；先是\Person{雅各伯}，后是\Person{以色列}；他自己也是\Quote{出自以色列支派，本雅明支派。}在他身上，让我们看见\Person{基督}在上，\Person{基督}在下。首先，主从天上的声音本身就显示了这一点；\Quote{撒乌耳，撒乌耳，你为什么迫害我？}什么！\Person{保禄}升了天吗？\Person{保禄}甚至向天投石了吗？他迫害基督徒，捆绑他们，拖去处死，在各地搜索隐藏的人，一旦找到丝毫不宽容。主\Person{基督}对他说：\Quote{撒乌耳，撒乌耳，你为什么迫害我？}祂从何处喊？从天。所以祂在上。\Quote{你为什么迫害我？}所以祂在下。亲爱的，我已尽我所能在简短中向你们解释了一切。我按职责服事了你们，现在你们的职责是纪念穷人。让我们转向主，等等。\par

\SourceRecord{Translated by R.G. MacMullen.；Nicene and Post-Nicene Fathers, First Series；Vol. 6.；Edited by Philip Schaff.；Buffalo, NY: Christian Literature Publishing Co.,；1888.；<http://www.newadvent.org/fathers/160372.htm>.}

% structure: p0001 source=h1 target=section reason=page-title

\section{讲道第七十三篇 论新约}

\Emph{[CXXIII. Ben.]}\par

\Emph{论福音的话}，\BibleRef{若 2:2} \Emph{，\BibleQuote{耶稣和祂的门徒也被请去赴婚宴。}}\par

1. 你们知道，弟兄们，因为你们相信基督而学到了，并且我们也通过我们的职务不断地向你们强调，基督的谦逊是对人类肿胀的骄傲的良药。因为人若没有因骄傲而膨胀，就不会灭亡。因为\BibleQuote{骄傲，}如圣经所说，\BibleQuote{是万恶的开端。}对抗罪恶的开端，需要正义的开端。若骄傲是万恶的开端，那么骄傲的肿胀如何能得到医治，除非天主屈尊谦卑自己？让人羞于骄傲，因为天主已谦卑自己。因为当人被告知要谦卑时，他藐视它；而当人受伤害时，是骄傲使他们想要报复。因为他们藐视谦卑，他们想要报复；仿佛别人的惩罚对任何人有什么益处。一个受伤受屈的人想要报仇；他从别人的惩罚中寻求自己的医治，却获得了巨大的痛苦。因此主基督屈尊在一切事上谦卑自己，向我们展示道路；只要我们愿意行走其上。\par

2. 在祂的其他作为中，看，童贞女的儿子来赴婚宴；祂与圣父一同设立了婚姻。正如第一个女人，因她有了罪，是从男人而出，没有女人；同样，那除去了罪的人子，是从女人而出，没有男人。通过前者我们堕落，通过后者我们上升。祂在这婚宴上做了什么？祂将水变成了酒。还有什么比这更大的权能标志？祂有行这样事的能力，却屈尊成为需要者。祂能将水变酒，也能将石头变成面包。权能是相同的；但那时魔鬼试探祂，因此基督没有做。因为你们知道，当主基督受试探时，魔鬼向祂提出这个。因为祂饿了，这也是祂屈尊承受的，因为这也是祂降卑的一部分。那面包饿了，正如道路疲乏，正如救恩受伤，正如生命死亡。当祂饿了，正如你们所知，试探者对祂说：\BibleQuote{你若是天主子，就命这些石头变成面包吧。}祂回答试探者，教导你如何回答试探者。因为元帅之所以战斗，是为了让士兵学习。祂如何回答？\BibleQuote{人生活不只靠面包，也靠天主口中的一切话。}祂没有用石头做面包，当然祂很容易就能做，正如祂将水变成酒一样。因为用石头做面包是同样的权能；但祂没有做，是为了藐视试探者的意愿。因为战胜试探者没有别的方法，只有藐视。当祂战胜了魔鬼的试探，\BibleQuote{天使前来伺候祂。}那么，祂有如此伟大的权能，为什么没有做这个，却做了那个？请读，是的，请回想你们刚才听到的，当祂做这事时，也就是，当祂将水变成酒时，福音作者加上了什么？\BibleQuote{祂的门徒就信从了祂。}在另一个场合，魔鬼会相信祂吗？\par

3. 因此，祂能行如此伟大的事，却饿了，渴了，疲乏了，睡觉了，被捉拿，被鞭打，被钉死，被杀害。这就是道路；通过谦逊行走，好能到达永恒。基督——天主是我们前往的故乡；基督——人是我们所行的道路。我们走向祂，借着祂行走；为何害怕我们走迷？祂没有离开圣父，却来到我们中间。祂吮吸母乳，却容纳世界。祂躺在马槽里，却喂养天使。天主和人，同一天主也是人，同一个人也是天主。但祂作为人时不是天主，天主是圣言；人是圣言成了血肉；同时继续作为天主，并取了人的血肉；增添了祂所不是的，没有失去祂所是的。因此从此以后，在祂的降卑中受苦、死亡、埋葬之后，祂已复活，升了天，在那里，坐在圣父的右边；而在这里，祂在祂的穷人身上仍有所缺。昨天我也借祂对\Person{纳塔乃耳}说的话向你们爱德阐述了这点：\BibleQuote{你将要看见比这更大的事。因为我实实在在告诉你们：你们要看见天开了，天主的天使在人子身上，上去下来。}我们探究了这是什么意思，并讲了一些时间；今天要重述吗？让那些在场的人记住；但我将简要概述。\par

4. 祂不会说\BibleQuote{在人子身上上去，}除非祂是上边；祂不会说\BibleQuote{在人子身上下来，}除非祂也是下边。祂同时是上边和下边；在上边是祂自己，在下边是祂的肢体；在上边与圣父同在，在下边在我们内。因此也有那声音对\Person{扫禄}说：\BibleQuote{扫禄，扫禄，你为什么迫害我？}祂不会说\BibleQuote{扫禄，扫禄，}除非祂是上边。但\Person{扫禄}并没有迫害上边的祂。那么，那在上边的祂不会说\BibleQuote{你为什么迫害我？}除非祂也在下边。要敬畏上边的基督；要认出下边的基督。有上边的基督广施恩惠，也要认出祂在这里有需要。在这里祂是贫穷的，在那里祂是富有的。基督在这里贫穷，祂自己对我们说：\BibleQuote{我饿了，我渴了，我赤身露体，我作客，我在监里。}祂对一些人说：\BibleQuote{你们服事了我，}对另一些人说：\BibleQuote{你们没有服事我。}看，我们证明了基督贫穷；基督富足，谁不知道？甚至在这里，将水变成酒是这些财富的特性。如果拥有酒的人是富有的，那么制造酒的人何等富有？所以基督既富有又贫穷；作为天主，富有；作为人，贫穷。是的，如今作为真人，祂已升了天，坐在圣父的右边；但在这里祂仍然贫穷、饥饿、渴了、赤身露体。\par

5. 你是什麼？富有的，還是貧窮的？許多人對我說：\Quote{我是貧窮的。}他們說的是真話。我看出有些貧窮的人擁有一些東西，有些則一無所有。但有些人有許多金銀。但願他們承認自己是貧窮的！他們若承認身邊的窮人，就會承認自己是貧窮的。因為怎麼說呢？無論你擁有多少，你這個富有的人，不論你是誰，你都是天主的乞丐。祈禱的時刻來臨，我就在那裡考驗你。你提出你的祈求。你怎麼不是貧窮的呢？你提出祈求。我再說，你為麵包祈求。你豈不是要說：\BibleQuote{求你今天賞給我們日用的食糧}？你這個為日用食糧祈求的人，是貧窮的，還是富有的？然而基督對你說：\Quote{把你從我領受的，給我一些。因為你來到這裡時，帶來了什麼？我所創造的一切，你在這裡都找到了；你什麼也沒有帶來，什麼也不能帶走。你為什麼不把你從我得到的給我呢？因為你飽足了，而窮人空虛。看看你的最初起源；你們倆生來都是赤身露體的。那麼，你也是生來赤身露體的。你在這裡找到了大量的財富；你帶了什麼來嗎？我要求我自己的；給吧，我會償還的。你發現我是慷慨的施予者，立刻使我成為你的債務人。\InnerQuote{你發現我是慷慨的施予者，立刻使我成為你的債務人}還不夠；讓我視你為放貸取利。你給我少許，我將償還更多。你給我屬地的事物，我將償還屬天的事物。你給我暫時的事物，我將歸還永恆的事物。當我把你歸還給我時，我將把你歸還給你自己。}\par

\SourceRecord{Translated by R.G. MacMullen.；Nicene and Post-Nicene Fathers, First Series；Vol. 6.；Edited by Philip Schaff.；Buffalo, NY: Christian Literature Publishing Co.,；1888.；<http://www.newadvent.org/fathers/160373.htm>.}

% structure: p0001 source=h1 target=section reason=page-title

\section{\WorkTitle{新约讲道第七十四篇}}

\Emph{[第一二四篇·本笃版]}\par

\Emph{论福音的话}，\BibleRef{若 5:2}，\Emph{\BibleQuote{在耶路撒冷羊门旁有一个池子}，等等。}\par

1. 福音的训诲刚刚在我们耳中响起，使我们渴望知道所读的经文有何意义。我想，这是人们对我的期望；我承诺，靠主的助佑，尽我所能解释。因为无疑，那些奇迹不是没有意义的，它们为我们预示了关乎永恒救恩的某事。因为这个人的身体得回的健康，能持续多久呢？\Quote{因为你们的生命是什么？}圣经说，\Quote{它不过是一片云雾，出现片刻，就消散了。}因此，在健康被暂时恢复给这人时，某种持久性被恢复给了云雾。所以这并不值得太重视：\Quote{人的健康是虚幻的。}弟兄们，请回忆那先知和福音的见证，因为它在福音书中读过：\BibleQuote{凡有血肉的都像草，一切荣华都像草上的花；草必枯干，花必凋谢，}好像我们应该问：\Quote{草有什么希望？草上的花有什么稳定？}它说：\BibleQuote{但主的圣言永远常存。}主的圣言甚至把荣耀赐给了草，而且是不会消逝的荣耀；因为祂甚至把不朽赐给了血肉。\par

2. 但是首先这今生的苦难消逝，从其中祂给了我们帮助，我们对祂说过：\Quote{求你从苦难中救助我们。}而这一切生命确实是对理解的磨难。因为有两个折磨灵魂者，不是同时折磨，而是轮流折磨。这两个折磨者的名字是：恐惧和悲伤。当你顺利时，你恐惧；当你倒霉时，你悲伤。这世界的顺境，谁不被它欺骗？它的逆境，谁不被它击垮？在这草中，在草的日子里，必须持守更稳妥的道路，就是天主的\OriginalTerm{圣言}{Word}。因为当它被说过：\BibleQuote{凡有血肉的都像草，一切荣华都像草上的花；草必枯干，花必凋谢；}好像我们应该问：\Quote{草有什么希望？草上的花有什么稳定？}它说：\BibleQuote{但主的圣言永远常存。}你将会说：\Quote{那圣言对我来说怎样？}\BibleQuote{圣言成了\OriginalTerm{血肉}{Flesh}，居住在我们中间。}因为主的圣言对你说：\Quote{不要拒绝我的许诺，因为我没有拒绝你的草。}那么，主的圣言所赐给我们的——使我们能持守祂，不随草上的花消逝——我所说的，祂所赐给我们的，就是圣言成了血肉，取了血肉，并非变为血肉，而是保持原样并取了我们；保持祂所是的，取了祂所不是的。——我所说的，祂所赐给我们的，那个池子也象征了这一点。\par

3. 我说得简短。那水是犹太人民，五个廊子是法律。因为梅瑟写了五本书。因此，水被五个廊子围住，正如那个民族被法律束缚。水的搅动是主在那民族中的\OriginalTerm{受难}{Passion}。那下去的人就痊愈了，而且只有一个；因为这象征合一。凡对基督的受难反感的人都是骄傲的；他们不肯下去，就不痊愈。他们说：\Quote{我能相信天主\OriginalTerm{降生成人}{Incarnate}，天主由女人所生，天主被钉十字架、被鞭打、死亡、受伤、埋葬吗？}我绝不能相信天主是这样的，这有损于祂的尊严。让心说话，不要颈项。对骄傲的人，主的谦卑似乎有损于祂的尊严；因此救恩远离这样的人。不要自高；你若想痊愈，就下去。如果只谈论基督在可变化之肉体内，虔诚可能会惊慌。但现在真理向你表明，基督在其本性上，作为\OriginalTerm{圣言}{Word}，是不变的。因为\BibleQuote{在起初已有圣言，圣言与天主同在；}因为\BibleQuote{圣言就是天主。}所以你的天主是不变的。噢，真虔诚！你的天主常存，不要害怕。祂不会消灭，借着祂，你也不会消灭。祂常存，祂由女人所生，但只是在\OriginalTerm{血肉}{Flesh}中。祂在祂被造成之前就存在，祂造了祂要在其中被造的她自己。祂是婴儿，但只是在血肉中。祂吃奶，祂成长，祂摄取营养，祂经历了人生的各个阶段，祂到了成年，但只是在血肉中。祂疲倦了，祂睡了，但只是在血肉中。祂又饥又渴，但只是在血肉中。祂被捕、被绑、被鞭打、被辱骂、终于被钉十字架、被杀，但只是在血肉中。你为什么惊慌？\BibleQuote{主的圣言永远常存。}谁拒绝天主的这一谦卑，就不愿因骄傲的致命肿胀而得医治。\par

4. 因此，主耶稣基督藉祂的肉躯，将希望赐予我们的肉躯。因为祂取了我们在世上所熟知、此处充盈的那一切——即出生与死亡。出生与死亡在此处充盈；而复活并永远生活，却不在此处。祂在此处找到了可怜的地上货物，却带来了陌生而属天的货物。你若因死亡而惊慌，就当热爱复活。祂已从患难中赐你援助；否则你的健康便永远徒然。为此，让我们承认并热爱这世上陌生的救恩，即永生的健康，并在此世如旅客般生活。让我们思想自己不过是在流逝，如此我们便会少犯罪。让我们更感谢主天主，因祂乐意使此生的最后一日既近在眼前又不确定。从最早婴孩至衰老暮年，不过是短暂的一瞬。倘若亚当今日死去，他活了那么久又有什么益处呢？在其中有终结的事物中，何谓\Quote{长久}？无人记起昨日；今日被明日催迫，好使之流逝。在这短暂的一瞬中，让我们善度此生，好使我们能前往那永不再逝去之处。此刻，即使我们正在谈论，也的确在流逝。我们的言语流转，时光飞逝；我们的年岁、我们的行为、我们的尊荣、我们的苦难、我们在这世上的幸福，无不如此。一切皆逝，但让我们不要惊慌：\Quote{天主的圣言永远常存}。让我们转向主，等等。\par

\SourceRecord{Translated by R.G. MacMullen.；Nicene and Post-Nicene Fathers, First Series；Vol. 6.；Edited by Philip Schaff.；Buffalo, NY: Christian Literature Publishing Co.,；1888.；<http://www.newadvent.org/fathers/160374.htm>.}

% structure: p0001 source=h1 target=section reason=page-title

\section{新约讲道第七十五篇}

\EditorialNote{CXXV. Ben.}\par

\Emph{再次论述} \BibleRef{若 5:2} \Emph{，关于那五个廊子，那里躺着大批患病的群众，以及史罗亚池。}\par

1. 现在再次叙述既非你们的耳朵也非你们的心所陌生的事题：但它们却唤起听者的情感，并在某种程度上借重复使我们更新；且听已经熟知的事并不厌烦，因为主的话语常是甘甜的。对圣经的诠释就像圣经本身：虽然它们已熟知，但仍被诵读以加深对它们的记忆。因此，对它们的诠释，虽已熟知，仍须重复，使那些遗忘的人得以记起，或那些未曾听见的人得以听见；并使那些确实记得惯常听到的内容的人，借重复而得以不致忘记。因为我记得，亲爱的诸位，我已经就这篇福音章节向你们讲过。然而向你们重复同样的解释并不使我厌烦，正如向你们重复同样的课文也不厌烦一样。宗徒保禄在某封书信中说：\BibleQuote{向你们写同样的事，为我并不厌烦，但为你们却是必要的。} 同样，我自己向你们说同样的事，为我并不厌烦，但为你们却是稳妥的。\par

2. 那五个患病的人躺卧的廊子象征法律，这法律最初是由天主的仆人梅瑟赐给犹太人和以色列人民的。因为这位法律的代理人梅瑟写了五本书。因此，根据他所写书卷的数目，五个廊子预表了法律。但因法律赐下不是为医治病人，而是为了发现并显露他们；因为宗徒这样说：\BibleQuote{因为假若曾颁给一种能赐与生命的法律，那么义德确实是出于法律了。但圣经却把众人都禁锢在罪恶之中，为使恩许藉耶稣基督的信德赐与那些相信的人。}因此，那些人躺在那些廊子里，但未被治愈。因为他怎么说？\BibleQuote{假若曾颁给一种能赐与生命的法律。}因此，那些预表法律的廊子不能治愈病人。有人会对我说：\Quote{为何那时赐与了呢？}宗徒保禄自己解释说：\BibleQuote{圣经}，他说，\BibleQuote{却把众人都禁锢在罪恶之中，为使恩许藉耶稣基督的信德赐与那些相信的人。}因为这些患病的人自以为健康。他们接受了法律，却不能遵守；他们得知了自己的病症，就恳求医生的帮助；他们希望被治愈，因为他们意识到自己处于困境，若非他们未能遵守所赐的法律，他们本不会知道。因为人自以为无罪，并因这虚假无罪的骄傲而更加疯狂。因此，为了驯服这骄傲并揭露它，法律被赐下；不是为了解救病人，而是为了说服骄傲的人。亲爱的诸位，请留意：法律被赐下正是为此目的：为了发现疾病，而非除去疾病。因此，那些病人，如果没有那五个廊子，他们本可在自己家中更为隐蔽地生病，却被置于那些廊子中，暴露在众人眼前，但并未被廊子治愈。所以法律有助于发现罪，因为人因违犯法律而罪过更重，便在驯服骄傲后，恳求那怜悯者的帮助。请听宗徒的话：\BibleQuote{法律本是外来的，是为叫过犯增多；但是罪恶在哪里增多，恩宠在那里也格外丰富。}这是什么意思？\BibleQuote{法律本是外来的，是为叫过犯增多}？正如他在另一处所说：\BibleQuote{因为那里没有法律，那里就没有违犯。}人在法律之前可被称为罪人，但不能称为违犯者。但他犯罪后，接受了法律，他不仅成了罪人，也成了违犯者。因此，既然罪上加上了违犯，所以\BibleQuote{罪恶增多}。当罪恶增多时，人类的骄傲终于学会屈服，并向天主认罪，说：\BibleQuote{我是软弱的。}这句只有谦虚的灵魂才能说的圣咏的话：\BibleQuote{我说，上主，求你怜悯我，医治我的灵魂，因为我得罪了你。}让那软弱的灵魂这样说吧，它至少因违犯而被说服，并非被治愈，而是被法律所显露。请再听保禄本人向你们显示，法律是善的，但只有基督的恩宠能使人脱离罪恶。因为法律能禁止和命令；它施加药方，但那使人不能遵守法律的，它无法治愈，只有恩宠才能做到。因为宗徒说：\BibleQuote{因为我按内心的人喜爱天主的法律。}也就是说，我现在看到法律所责备的是邪恶，法律所命令的是良善。\BibleQuote{因为我按内心的人喜爱天主的法律。但我发现在我肢体内另有一种法律，抗拒我理智的法律，并把我掳去成为在我肢体内罪恶法律的俘虏。}这来自罪的惩罚、死亡的流传、对亚当的定案，\BibleQuote{抗拒理智的法律，并把它掳去成为在肢体中罪恶法律的俘虏。}他被说服了；他接受了法律，为使自己被说服：看他被说服后得了什么益处。请听以下的话：\BibleQuote{我这个人真不幸！谁能救我脱离这死亡的身体？藉我们的主耶稣基督，天主的恩宠。}\par

3. 请留意。那五个门廊象征法律，承载病人，但不医治他们；揭露疾病，却不治愈。但谁治愈病人呢？就是那下到池子里的那位。病人何时下到池子里呢？当天使通过水的波动给出征兆时。因为那池子就是这样被祝圣的，当天使下来搅动水时。人们看见水，从被搅动的水的震动中，他们明白天使的临在。若有人那时下去，便得痊愈。那么，为什么那病人没有痊愈呢？让我们思考他自己的话：\Quote{我没有人}，他说，\Quote{在水动的时候，把我放入池子；但我正去的时候，别人先我而下。}难道你不能在那人之后下去吗？这里向我们显示，每次水动时只有一人得痊愈。谁先下去，谁就独得痊愈；但随后下去的人，在那次水动时不得痊愈，而要等待下一次搅动。那么，这奥秘是什么意思呢？因为它不是没有意义的。亲爱的，请留意。在\BibleBook{}{默示录}中，水被用作民众的象征。因为当\Person{若望}在\BibleBook{}{默示录}中看见许多水时，他问这是什么意思，被告知那是民众。因此，池子里的水象征犹太民众。正如那个民族被\Person{梅瑟}五书拘束在法律中，那水也被五个门廊围住。水何时被搅动？当犹太民众被搅动时。犹太民众何时被搅动，岂不是当主耶稣基督来临时？主的受难就是水的搅动。因为主受苦时，犹太人就被搅动了。看，刚才读到的经文正与这搅动有关：\BibleQuote{犹太人想要杀祂，不但因为祂在安息日做这些事，而且因为祂称自己为天主子，使自己与天主平等。}因为基督以一种方式称自己为子，而另一种方式是对人说的：\BibleQuote{我曾说：你们是神，都是至高者的儿子。}因为若祂像任何人那样被称为天主子（因天主的恩宠，人被称为天主子），犹太人就不会发怒。但他们理解祂以另一种方式自称天主子，即如\BibleQuote{在起初已有圣言，圣言与天主同在，圣言就是天主；}又如宗徒所说：\BibleQuote{祂虽具有天主的形体，却没有以自己与天主同等为僭越；}他们看见一个人，就发怒，因为祂使自己与天主平等。但祂深知自己平等，然而在他们看不见的方面。因为他们所看见的，他们想钉死；藉他们所看不见的，他们被审判。犹太人看见了什么？也是宗徒们所看见的，当\Person{斐理伯}说：\BibleQuote{求父显示给我们，我们就心满意足了。}但犹太人没有看见什么？连宗徒们也没有看见的，当主回答说：\BibleQuote{这么长久与你们在一起，你们还不认识我吗？看见我的，就是看见了父。}因为犹太人不能在祂身上看见这一点，他们就视祂为一个骄傲不敬的人，使自己与天主平等。这就是一次搅动，水被搅动了，天使已经来了。因为主也被称为\OriginalTerm{大议会的天使}{Angelus magni consilii}，因为祂是父旨意的使者。因为希腊语的\Quote{天使}在拉丁语中是\OriginalTerm{使者}{messenger}。所以你们有主说祂向我们宣报天国。祂已经来了，\Quote{大议会的天使}，但祂是众天使的主。称为\Quote{天使}，因为祂取了肉身；\Quote{众天使之主}，因为\BibleQuote{万物是藉着祂造的，凡受造的，没有一样不是藉着祂造的。}因为如果万物都是藉着祂造的，天使也是。因此祂自己不是受造的，因为万物是藉着祂造的。凡受造的，没有圣言的运作就不能被造。但成为基督母亲的肉身，若非由以后从她所生的圣言所创造，就不能诞生。\par

4. 那时犹太人便不安了。这是什么意思？\BibleQuote{祂为什么在安息日做这些事？}尤其是当主说：\BibleQuote{我父至今工作，我也工作。}他们由于肉性地理解天主在第七日歇了祂一切的工作，便\Quote{不安起来。}因为这记载于\BibleBook{}{创世纪}，且是极为优美的记载，也有极好的理由。但他们愚昧地以为天主在第七日仿佛是从疲劳中安息了，且因此祝福了它，因为在那日祂从困倦中恢复了精神，却不明白那位以圣言创造万物的天主，是不会疲倦的。让他们去读，并告诉我，那位说\BibleQuote{让它被造，它就被造了}的天主，怎么可能疲倦呢？今天如果有人能像天主那样做，他怎么会疲倦呢？祂说：\BibleQuote{要有光，就有了光。}又说：\BibleQuote{要有穹苍，就有了穹苍。}如果祂说了，却没有成就，那么祂才是疲倦了。在另一处简短地说：\BibleQuote{祂说话，它们就造成；祂命令，它们就受造。}那么，这样工作的祂，怎会劳苦呢？但若祂不劳苦，祂又怎会安息呢？然而，在记载天主歇了祂一切工作的那个安息日里，所象征的是天主安息中的我们的安息；因为当六个时期过去之后，这世界的安息日就要来到。这世界的六天正在消逝。第一天已经过去，从亚当到诺厄；第二天从洪水到亚巴郎；第三天从亚巴郎到达味；第四天从达味到被掳往巴比伦；第五天从被掳往巴比伦到我们的主耶稣基督的来临。现在第六天正在消逝。我们正处于第六时期，第六天。那么，让我们按天主的肖像被改造，因为在第六天，人是按天主的肖像受造的。当初的塑造，让改造在我们内完成；当初的创造，让再造在我们内完成。在这我们现在所处的日子之后，在这个时期之后，应许给圣人们并在那些日子预表过的安息将要来到。因为实在的，在祂创造了世界上的一切之后，祂再也没有在受造界中创造新的东西。受造物本身将要被转化和改变。因为自从受造物成形以来，再没有增加什么。然而，如果创造者不治理世界，被造物就会崩毁：祂不能不管理祂所造的。因此，既然没有给受造物增加什么，祂就被称为歇了祂一切的工作；但祂并没有停止管理祂所造的，所以主正确地宣称：\BibleQuote{我父至今工作。}请注意，亲爱的。祂完成了，被称为安息了；因为祂完成了祂的工作，并且没有再增加。祂治理祂所造的，因此祂并没有停止工作。但祂以创造时同样的容易来治理。因为，弟兄们，不要以为祂在创造时不劳苦，而在治理时劳苦：好比在船上，建造船的人劳苦，管理船的人也劳苦；因为他们是人。因为祂以\BibleQuote{祂说话，它们就造成}的同样容易，以同样的容易和判断，藉着圣言治理万有。\par

5. 我们不要因为人世间的事看似混乱，就以为人世间的事没有治理。因为众人各自在适当的位置上被安排；但在每人看来似乎毫无秩序。你只看你愿意成为什么；因为你愿意成为什么，主人知道把你安置在哪里。看看一位画家。在他面前摆着各种颜色，他知道每种颜色该放在哪里。毫无疑问，罪人选择了做黑色；那么，艺术家难道不知道把他放在哪里吗？画家用黑色完成了多少部分？用它做出了多少装饰？他用它画头发、胡须、眉毛；他只将脸用白色来画。那么，你只看你愿意成为什么；不要操心那位不会错误的主会把你安排在哪里，祂知道把你放在哪里。因为我们也看到世上普遍的法律也是如此。比如，有人选择了做窃贼：法官的法律知道他违反了法律；法官的法律知道把他安置在哪里；并极恰当地处治他。他确实生活邪恶；但法律并没有邪恶地处置他。从窃贼他将被判罚到矿场；从这种劳役中建造出多么大的工程！那被定罪者的惩罚成了城市的装饰。同样，天主也知道把你安置在哪里。如果你存心要混乱，不要以为你扰乱了天主的计划。那知道如何创造你的，难道不知道如何安置你吗？你努力争取这个，即被安置在好的地方，这对你是有益的。宗徒论及犹达斯说了什么？\BibleQuote{他去了他自己的地方。}这当然是天主管照的运作；因为他以邪恶的意志选择成为邪恶，但天主并没有藉着安排邪恶而制造邪恶。但那个邪恶的人自己选择了成为罪人，他做了他所愿意的，并承受了他所不愿意的。在他做了他所愿意的事上，他的罪显明了；在他承受了他所不愿意的事上，天主的秩序受到了赞美。\par

6. 我为何说这一切呢？弟兄们，是要你们明白主耶稣基督所说最卓越的话：\BibleQuote{我父工作至今。} 在祂不放弃祂所造的受造物这件事上。祂又说：\BibleQuote{祂怎样工作，我也怎样工作。} 在此祂立刻表明了自己与天主相等。祂说：\BibleQuote{我父}，祂说，\BibleQuote{工作至今，我也工作。} 他们属肉体的认知关于安息日被扰乱了。因为他们认为主疲乏了，所以休息，不再工作。他们听到\BibleQuote{我父工作至今，}就惊惶。\BibleQuote{我也工作：}祂使自己与天主平等，他们便惊惶。但不要惊慌。水被搅动，现在病人将要被医治。这是什么意思？所以他们惊惶，好让主受难。主确实受难，宝贵的血被倾流，罪人被救赎，恩宠赐给罪人，赐给那说\BibleQuote{我真不幸！谁能救我脱离这该死的肉身呢？天主的恩宠，藉着我们的主耶稣基督。}的人。但他如何被医治？如果他下去。因为那水池是这样建造的，人应当下去，而不是上去。因为可能有另一种水池，建造得使人必须上去。但为何这样建造，使人必须下去？因为主的受难寻找谦卑的人。让谦卑的人下去，不要骄傲，如果他希望被医治。但为何只有一个？因为教会全世界只有一个，合一得以保存。那么当一个人痊愈时，合一被象征。由一理解合一。因此不要离开合一，如果你不想在这救人的医治中无份。\par

7. 那么，那人患病三十八年是什么意思？弟兄们，我知道我已经讲过这个；但那些读书的人尚且忘记，何况那些少听的人呢？因此，亲爱的诸位，请稍加留意。在数字四十中，象征了公义的圆满。公义的圆满，在于我们在此世生活在劳苦、辛劳、自制、禁食、醒寤、磨难中；这就是公义的操练，承受现世的时光，并仿佛从这世界禁食；不是从身体的饮食禁食，那是我们很少做的；而是从对世界的爱禁食，那是我们应当常常做的。那么，凡戒绝这世界的人就成全了法律。因为他不能爱永恒的事物，除非他不再爱属时间的事物。请思想一个人的爱：把它想象为灵魂的手。如果它握着什么东西，就不能再握别的东西。但为了能握住给它的东西，它必须放开它已握着的东西。我说这话，你们看我说得多么明白：\Quote{凡爱世界的，不能爱天主；他的手被占用了。} 天主对他说：\Quote{握住我所给的。} 他不肯放开他正握着的东西，就不能接受所给的。我是否说过人不应拥有任何东西？如果他能，如果成全要求他如此，就让他不拥有。如果因某种必要而不能，就让他拥有，但不要被拥有；让他支配，不要被支配；让他做财产的主人，而非奴隶；正如宗徒所说：\BibleQuote{不过，弟兄们，时限是短促的；今后有妻子的，要像没有一样；购买的，要像一无所得的；喜庆的，要像不喜庆的；哭泣的，要像不哭泣的；享用这世界的，要像不享用的，因为这世界的局面正在逝去。我愿你们无所挂虑。} 那\Quote{不要爱你在世上所拥有的}是什么意思？不要让它抓紧你的手，那手必须抓住天主。不要让你的爱被占用，藉着它你才能走向天主，并紧紧依附于造你的天主。\par

8. 你们会对我说，并回答我说：\Quote{的确，天主知道，我清白地拥有我所拥有的。}试探在考验你。你们的财产受到扰动，你就亵渎。不久前我们就曾遭遇这样的情形。你们的财产受到扰动，你就不再是原先的样子，你显明今天你口中的是一回事，昨天你口中的是另一回事。我宁愿你们即使激烈地捍卫自己的东西，也不愿你们大胆地夺取别人的东西；更糟的是，为了逃避责备，你们坚持别人的东西是你们自己的。但我何必再多说？这就是我的劝告，这就是我所说的，诸位弟兄，我如同兄弟一样劝告：天主命令，我劝告，因为我也受到劝告。祂使我警醒，不容我沉默。祂向我索回祂所给予的。因为祂赐予它是为了分施，而非保存。如果我将它保存并隐藏起来，祂会对我说：\BibleQuote{你这又恶又懒的仆人，为什么没有把我的银子交给兑换银钱的人，好使我回来时，可以连本带利收回？}我若未失去所领受的任何东西，对我有何益处？这对我的主还不够，祂是贪求的；但天主的贪求却是我们的救恩。祂是贪求的，祂寻找自己的银子，祂聚集自己的肖像。\BibleQuote{你应当给，}祂说，\BibleQuote{把银子交给兑换银钱的人，等我回来时，可以连本带利收回。}如果因偶然的遗忘，使我未能劝勉你们，那么我们正在经受的试探和患难，至少会成为对你们的劝勉。你们至少听到了天主的话。愿主和祂的荣耀受赞美。因为你们聚集在这里，渴望聆听天主仆人的话语。不要留意我们的肉身，因为天主的话是通过肉身传给你们；因为饥饿的人不介意器皿的卑贱，而看重食物的珍贵。天主在考验你们。你们聚集在一起，赞美天主的话；试探将考验你们如何聆听：你们将通过生活中的忙碌事务显露出你们的真面目。因为今天大声辱骂的人，昨天还是一个乐于聆听的人。因此我预先警告，因此我告诉你们，因此我不隐瞒，我的诸位弟兄，审问的时刻将要来到。因为主察问义人和恶人。你们知道你们曾歌唱过，我们也曾一同歌唱过：\BibleQuote{主察问义人和恶人。}接下来是什么？\BibleQuote{喜爱不义的人，憎恨自己的灵魂。}在另一处：\BibleQuote{对不虔敬者的意念，必会加以审问。}天主不在我审问你们的地方审问你们。我审问你们的舌头，天主审问你们的意念。因为祂知道你们如何聆听，也知道如何索取，祂命令我给予。祂愿意我作分施者，而将索取的权利保留给自己。劝勉、教导、斥责是我们的本分；但拯救、加冕，或定罪、投入地狱，不是我们的本分；\BibleQuote{判官把他交给狱警，狱警把他投在监狱里。我实在告诉你们，非到你还清最后一文钱，你决不能从那里出来。}\par

9. 让我们回到我们的主题。义德的成全由四十这个数目所显示。满全四十这个数目是什么意思？就是克制自己脱离对这个世界的爱。克制自己脱离暂时之物，以免我们爱它们而自取灭亡，就如同与此世界禁食。因此，主禁食了四十天，\Person{梅瑟}和\Person{厄里亚}也是。那位赐予祂仆人们禁食四十天能力的那一位，难道不能禁食八十天或一百天吗？那么，祂为什么不愿意禁食超过祂给予祂仆人们所行的呢？岂不是因为在这四十这个数目中，隐藏着禁食的奥迹，即克制自己脱离这个世界？这意味着什么？宗徒所说的：\BibleQuote{世界已钉在我身上，我也钉在世界上。}那么，他就满全了四十这个数目。主又显示了什么呢？既然\Person{梅瑟}这样做了，\Person{厄里亚}这样做了，基督这样做了，说明法律、先知和\BibleBook{}{福音}都这样教导；免得你们认为在法律中有一样，在先知中有另一样，在\BibleBook{}{福音}中又有另一样。全部圣经只教导你们一件事，就是克制自己脱离对世界的爱，好使你们的爱能飞奔向天主。作为法律教导此事的预象，\Person{梅瑟}禁食了四十天。作为先知教导此事的预象，\Person{厄里亚}禁食了四十天。作为\BibleBook{}{福音}教导此事的预象，主禁食了四十天。因此，在那山上，这三者也出现了，主在中间，\Person{梅瑟}和\Person{厄里亚}在两边。为什么？因为\BibleBook{}{福音}本身从法律和先知获得证言。但是，为什么在四十这个数目中是义德的成全呢？在\BibleBook{}{圣咏集}中记载说：\BibleQuote{天主，我要向祢唱新歌，用十弦琴向祢歌唱赞美。}这寓意法律的十诫，主来不是为废除，而是为满全。法律本身遍及全世界，显然有四个方向：东、西、南、北，正如圣经所说。因此，那载着各种象征动物的器皿——就是向\Person{伯多禄}展示的，当他对他说：\BibleQuote{宰了吃！}为显示外邦人应当相信并进入教会的身体，正如我们所吃的东西进入我们的身体——那由四角从天降下的器皿（这四角就是世界的四方），显示整个世界应当相信。因此，在四十这个数目中是对世界的克制。这就是法律的满全：法律的满全是爱德。因此，我们在逾越节前禁食四十天。因为逾越节前的这段时间是我们劳苦生活的标记，在其中，我们通过劳苦、忧虑和节制来满全法律。但之后我们庆祝逾越节，即主复活的日子，象征我们自己的复活。因此庆祝五十天，因为\OriginalTerm{银币}{denarius}的赏报加到四十上，就成了五十。为什么赏报是一枚\OriginalTerm{银币}{denarius}？难道你们没有读过，那些被雇进葡萄园的人，无论是在第一时辰、第六时辰还是第十一时辰，都只能得到一枚银币吗？当我们的义德加上它的赏报时，我们就将处于五十的数目中。是的，那时我们将没有任何其他职务，只有赞美天主。因此，在整个那些日子里，我们说：\BibleQuote{阿肋路亚。}因为\OriginalTerm{阿肋路亚}{Hallelujah}就是赞美天主。在这必朽坏的脆弱境况中，在这四十的数目中，如同在复活之前，让我们在祈祷中叹息，以便那时我们能歌唱赞美。现在是渴望之时，那时将是拥抱与享受之时。让我们不要在这四十的时期中灰心，以便我们在五十的时期中欢乐。\par

10. 那么，谁是满全法律的呢？不就是那拥有爱德的人吗？请问宗徒：\BibleQuote{爱德是法律的满全。因为全部法律都在一句话内得以满全，就是\InnerQuote{你应爱你的近人如你自己。}}但爱德的诫命是双重的：\BibleQuote{你应全心、全灵、全意爱上主你的天主。这是最大的诫命。第二条与此相似：你应爱你的近人如你自己。}这是主在\BibleBook{}{福音}中的话：\BibleQuote{这两条诫命是全部法律和先知的总纲。}没有这双重的爱，法律就不能满全。只要法律没有满全，就有软弱。因此，那病了三十八年的人，是缺少了两样。\Quote{缺少两样}是什么意思？他没有满全这两条诫命。如果其余的都满全了，而这两条没有满全，那有什么益处呢？你有三十八吗？如果你没有那两条，其余的对你就毫无益处。你缺少两样，没有这两样，其余的都无效，如果你没有那两条导向救恩的诫命。\BibleQuote{我若能说人间的语言和天使的语言，却没有爱德，我就成了鸣的锣、响的钹。我若有先知之恩，又明白一切奥秘和各种知识；我若有全备的信德，甚至能移山，却没有爱德，我什么也不算。我若把我所有的财产全施舍了，又舍身被火焚烧，却没有爱德，为我毫无益处。}这些都是宗徒的话。因此，他所提到的这一切，就如同三十八年；但因为爱德不在那里，就有软弱。那么，谁能使那软弱痊愈呢？不就是那位来赐予爱德的吗？\BibleQuote{我给你们一条新诫命：你们要彼此相爱。}因为祂来赐予爱德，而爱德满全法律，所以祂有充分的理由说：\BibleQuote{我来不是为废除法律，而是为满全。}祂治愈了那个病人，并叫他拿起床，回家去。祂也这样对他治愈的瘫子说。拿起我们的床是什么意思？我们肉体的快乐。我们在软弱中躺卧之处，就如我们的床。但那些被治愈的人驾驭并拿起它，并不被这肉体所驾驭。所以，你这痊愈的人，要驾驭你肉体的脆弱，以便在这与此世界禁食四十天的标记中，你满全四十这个数目，因为那位使病人痊愈的，\BibleQuote{祂来不是为废除法律，而是为满全。}\par

11. 你们听了这话，当把心转向天主。不要自欺。当你们在世上一帆风顺时，问问自己：你们是爱世界，还是不爱它？学会在你们被放下之前放下它。什么是放下它？不是诚心地爱它。当你们还有些东西迟早要失去，或在生前或在死时放下，它不会永远与你们同在；我说，当它还与你们同在时，松开你们的爱；准备好服从天主的旨意，依附于天主。紧紧抓住祂，祂是不能违背你们的意愿而失去的；这样，如果你们偶然失去这些短暂之物，你们可以说：\Quote{主赐予，主收回；如主所喜悦的，就如此成就；愿主的名受赞美。}但若偶然，且天主如此愿意，你们所拥有的东西甚至与你们同在直到最后：因你们脱离此世，你们领受\OriginalTerm{银币}{denarius}，五十，圆满的幸福将在你们内成就，那时你们将歌唱\Quote{阿肋路亚}。愿我如今向你们记忆中所陈明的这些事，有助于推翻你们对世俗的爱。它的友谊是邪恶的，骗人的，它使人成为天主的仇敌。很快，在一次诱惑中，人就得罪天主，成为祂的仇敌。不，那时他不是成为祂的仇敌；而是被发现曾是祂的仇敌。因为当他爱慕和赞美祂时，他是仇敌；但他自己不知道，别人也不知道。诱惑来了，脉搏被触动，热病被发现。所以，弟兄们，对世俗的爱和世俗的友谊使人成为天主的仇敌。它不履行它所承诺的善，它是说谎者，欺骗人。因此，人们从不停止在这世界上希望，又有谁达到了他所希望的一切呢？然而无论他达到什么，他所达到的立刻被他轻视。其他事开始被渴望，其他虚幻的事被盼望；当它们来到时，无论来到你面前的是什么，都被轻视。那么，你们当紧紧依附于天主，因为祂永不会被轻视，因为没有比祂更美的。正因为这些事被轻视，因为它们不能持久，因为它们不是祂是什么。因为虚无，啊，灵魂，不能满足你，除非那创造你的祂。无论你抓住别的什么，都是悲惨的；因为惟独祂能充足地满足你，祂按自己的肖像造了你。因此清楚地说：\Quote{主，将圣父显示给我们，我们就满足了。}惟有那里有安全；而在有安全的地方，将在某种程度上有一种不饱足之饱足。因为你既不会饱足到想要离开，也不会缺乏什么，好像你会遭受匮乏。\par

\SourceRecord{Translated by R.G. MacMullen.；Nicene and Post-Nicene Fathers, First Series；Vol. 6.；Edited by Philip Schaff.；Buffalo, NY: Christian Literature Publishing Co.,；1888.；<http://www.newadvent.org/fathers/160375.htm>.}

% structure: p0001 source=h1 target=section reason=page-title

\section{新约讲道集 第七十六篇}

\Emph{〔一二六．本笃〕}\par

\Emph{论福音的话}，\BibleRef{若 5:19}，\Emph{，\Quote{圣子不能由自己做什么，除非看见圣父做什么。}}\par

1. 天主国的奥迹与秘密首先寻求相信的人，好使他们获得理解。因为信德是理解的阶梯，而理解是信德的成果。先知明确地对那些过早、不合秩序地寻求理解而忽略信德的人说：\BibleQuote{除非你们相信，你们不会理解。}信德本身在圣经、先知书、福音和宗徒教训中也有自己的光。因为今世读给我们的一切，都是在黑暗处的光，使我们得以滋养直到那日子。宗徒伯多禄说：\BibleQuote{我们更有先知确切的预言，你们应留心它，如同在黑暗处的光，直到黎明破晓，晨星在你们心中升起。}\par

2. 你们看，弟兄们，那些像未成熟的概念一样，在出生前寻求过早分娩的人是多么无序而急躁。他们对我说：\Quote{你为什么叫我相信我没看见的？让我看见什么，好使我相信。你叫我相信，而我还看不见；我愿看见，并因看见而相信，不是因听见。}让先知说话：\BibleQuote{除非你们相信，你们不会理解。}你们想上升，却忘记了阶梯。显然，这是颠倒秩序。人哪，如果我能让你已经看见你所见的，我就不用劝你相信了。\par

3. 信德，正如别处所定义的，是\BibleQuote{所希望之事的坚实支撑，未见之事的明证。}如果它们未见，如何被证明存在？你所见的这些事物，难道不来自那未见者吗？诚然，你看见一些，好使你相信一些；从你所见的，你能相信你尚不能见的。不要忘恩于那使你看见的那一位，好使你能够相信你还不能见的。天主给了你身体的眼睛和心中的理性；唤醒心中的理性，唤醒你内在眼睛的内在居民，让他使用他的窗户，察看天主的受造物。因为有一个内在者通过眼睛来看。当你内里的思想转向其他事物，内在居民转离时，眼前的事物你便看不见。因为窗户空开，当透过它们观看者不在时，便毫无用处。所以，不是眼睛在看，而是某人通过眼睛在看；唤醒他，唤醒他。因为天主并没有拒绝你；祂造了你为理性的动物，使你高于牲畜，按祂的肖像造了你。你难道要像牲畜一样使用它们，只看见为肚子添加什么，而不为灵魂添加吗？我说，唤醒理性的眼睛，用人应有的方式使用你的眼睛；观察天地、天上的装饰、地上的果实、飞鸟的飞翔、鱼类的游动、种子的力量、季节的秩序；观察作品，寻求作者；查看你所见的，寻求你所不见的。因这些你所见的，相信那不见者。免得你以为是我用我自己的话劝你；听宗徒说：\BibleQuote{自创世以来，天主那看不见的事，藉着所造的之物，都可被清楚看见。}\par

4. 你轻视这些事，不以人的方式看待，却像无理性的牲畜。先知向你呼喊，却是徒然：\BibleQuote{不要像马和骡子，它们没有领悟。}我说，你看见了这些，却轻视。天主每日的奇迹被轻视，不是因为其容易，而是因为其不断重复。因为有什么比人的出生更难理解？一个人存在，却因死亡而离去进入黑暗；一个不存在的人，却因出生而来到光明。何等奇妙，何等难以理解？但对天主来说却容易实现。对这些事惊叹，醒过来吧！你为祂非同寻常的作为而惊叹，难道那些比你所习惯看见的更大吗？人们惊叹我们的主天主耶稣基督用五个饼使数千人吃饱，却不惊叹通过几粒种子，整个大地被庄稼充满。水变成酒，人们看见并惊异；雨水落在葡萄根上，岂不也同样发生？祂行了一件，也行另一件；一件使你得到喂养，另一件使你惊叹。但两者都奇妙，因为两者都是天主的工程。人看见不寻常的事就惊叹；那惊叹的人本身是谁？他在哪里？他从何处来？怎么形成？他身体的样子如何？肢体的区分如何？那美丽的形式从何而来？从什么开始？何等卑微的开始？而他却惊叹别的事物，他自己这个惊叹者本身就是一个大奇迹。所以，你所见的这些事物，除了从那不见者而来，还能从哪里来？但如我开始所说，因为你轻视了这些，祂亲自来做不寻常的事，好使你在这些寻常之事中也承认你的造物主。祂来了，向祂曾说：\BibleQuote{更新奇事。}向祂曾说：\BibleQuote{显明你奇妙的仁慈。}因为祂一直在施予；祂施予，却无人惊叹。因此祂作为婴孩来到婴孩中间，作为医生来到病人中间；祂能按祂的意愿来，按祂的意愿回去，做祂愿意做的一切，按祂的意愿审判。这祂的意愿就是真正的正义；是的，祂所意愿的就是真正的正义。因为祂所意愿的没有不义，祂所不意愿的也不可能是正义。祂来了，使死人复活，人们惊叹祂使已在光明中的人重见光明，而祂每日使不存在的人来到光明。\par

5. 祂做了这些事，却仍被许多人轻视，他们不太考虑祂做了多么伟大的事，而是看祂多么卑微；好像他们心里说：\Quote{这些是神圣的事，但祂是一个人。}那么你看见两件事：神圣的作为和一个人。如果神圣的作为只能由天主做成，就要当心在这人内天主是否隐藏了。我说，留心你所看见的；相信你所没有看见的。祂没有抛弃你，祂召叫你相信；尽管祂命令你相信你所不能看见的，但祂也没有让你什么也看不见，以便你能相信你所看不见的。受造物本身不就是一个小小的记号，一个小小创造者的迹象吗？祂也来了，祂行了奇迹。你不能看见天主，但你能看见一个人，所以天主成了人，好使你在同一者身上，既有可看见的，也有可相信的。\BibleQuote{在起初已有圣言，圣言与天主同在，圣言就是天主。}这样你听见了，却仍没有看见。看，祂来了，看，祂生出来了，看，祂从女人中出来，祂造了男人和女人。祂造了男人和女人，却不是由男人和女人所造。因为你们或许会因祂的出生而轻视祂，但祂的出生方式你们不能轻视；因为祂在出生之前就已经存在。我说，看，祂取了身体，祂穿上了肉身，祂从母胎中出来。现在你看见了吗？我说，现在你看见了吗？我问肉身，但我指向那肉身；有些东西你看见，有些东西你没有看见。看，就在这出生中，同时有两件事，一件你可以看见，另一件你不能看见；但这样，借着你所看见的，你可以相信你所没有看见的。你已经开始轻视，因为你看见那出生的人；相信你所没有看见的，即祂是由童贞女所生的。\Quote{多么微不足道的一个人，}有人说，\Quote{这个出生的人！}但那由童贞女所生的何其伟大！那由童贞女所生的给你带来一个暂时的奇迹；祂不是由父亲生的，我是指任何男人，祂的父亲，然而祂是由肉身生的。但不要觉得不可能：祂由母亲独自所生，祂在父亲和母亲之先造了人。\par

6. 祂给你们带来一个暂时的奇迹，好使你们寻求并赞美那永恒者。因为祂\BibleQuote{如新郎从洞房出来}，即从童贞女的胎中，在那里圣言与肉身举行了神圣的婚配：我说，祂带来了暂时的奇迹；但祂自己是永恒的，祂与圣父同为永恒，祂就是那\BibleQuote{在起初已有圣言，圣言与天主同在，圣言就是天主}。祂为你们做了这些，好使你们得以痊愈，能够看见你们所没有看见的。你们在基督里轻视的，还不是痊愈者的默观，而是病人的药物。不要急于看见完全的健康。天使看见，天使喜乐，天使以此维生；他们所喂养的不会短缺，他们的食物不会减少。在荣耀的宝座上，在诸天之上，在超越诸天的部分，圣言被天使看见，成为他们的喜乐；成为他们的食物，且常存不变。但为了使人能吃到天使的食粮，天使之主成了人。这就是我们的救恩，软弱者的药物，健康者的食粮。\par

7. 祂向人说话，说了你们刚才所听见的：\BibleQuote{圣子不能凭自己做甚么，而只看见圣父做甚么，祂才能做。}现在，我们以为有谁能理解这个呢？我们以为有谁身上肉眼的眼药已经生效，能多少辨别神性的光辉呢？祂说了，让我们也说；因为祂是圣言，我们因圣言而存在。我们为何谈论圣言呢？无论我们怎样做，都是关于圣言。因为我们是由圣言按照圣言的肖像造成的。那么，就我们所能的，就我们所能分享那不可言喻之事的程度，让我们也来说话，并且不被反驳。因为我们的信德已经先行，所以我们能说：\BibleQuote{我信了，所以我发言。}那么我讲我所相信的；无论我是否看见，或我如何看见；祂倒看见了；你们不能看见。但当我发言之后，那看见我所讲之事的人，是否相信我也看见了我所讲的，或他不相信，那对我有什么关系呢？只让他真实看见，并任由他对我随意相信吧。\par

8. \BibleQuote{圣子由自己什么也不能做，只有看见圣父做什么，才做什么。}此处兴起亚略派的错误；但它兴起是为了跌倒，因为它没有被谦卑下来，好能兴起。是什么使你偏离了正道？你会说圣子小于圣父。因为你听到了\BibleQuote{圣子由自己什么也不能做，只有看见圣父做什么，才做什么。}由此你便称圣子为较小的；这一点我知道，我知道正是这一点使你偏离了正道；要相信祂并不小，你尚不能看见，要相信，这就是我刚才所说的。\Quote{但是，}你会说，\Quote{我怎能相信与祂自己的话相反的事呢？}祂自己说：\BibleQuote{圣子由自己什么也不能做，只有看见圣父做什么，才做什么。}也要注意随后的话：\BibleQuote{因为圣父所作的事，圣子也照样作；}祂没有说\Quote{相似的事，}亲爱的，请思考一会儿，免得你们自相混淆。需要一颗宁静的心，一份虔敬与笃定的信德，一种虔诚的认真关注；要关注，不是关注我这可怜的器皿，而是关注那将饼放入器皿的那一位。那么请思考一会儿。因为在我以上为劝勉你们信德所说的一切，即那以信德浸染的心灵能领悟的一切，所说的一切都带着悦耳、欣喜而轻松的声音，振奋了你们的心灵，你们跟随了它，你们理解了我所说的。但我现在要说的，我希望有些人会理解；但我担心并非所有人都能理解。既然天主借着福音的读经给我们提出了一个要讲的题目，我们不能回避主所提出的题目；我担心那些不理解的人——或许会是多数——会认为我对他们白说了；然而为了那些将会理解的人，我并非白说。愿理解的人欢喜，愿不理解的人耐心忍受；他不理解的，让他忍受，且为能理解，让他忍受延迟。\par

9. 那么祂并没有说：\Quote{圣父所作的事，圣子也作相似的事：}仿佛圣父做些事，圣子另做些别的。因为当祂在上文说\BibleQuote{圣子由自己什么也不做，只有看见圣父做什么，才做什么}时，似乎正是此意。请注意；祂在那里也没有说：\Quote{只有听见圣父命令什么，才做什么；}而是\Quote{看见圣父做什么，才做什么。}这仍是血肉的观点。然而，为能理解那些更高的事，让我们不要拒绝这些较低和卑微的事。首先，让我们以这样的方式将某物放在眼前；让我们假设有两个工人，父亲和儿子。父亲做了一只箱子，儿子除非看见父亲做，否则不能做；他注视着父亲做的箱子，也做了一只相似的箱子，但不是同一只。我暂时搁置后面的话，现在我问亚略派：\Quote{你是在这种假设的意义上理解的吗？圣父做了某件事，圣子看见祂做了，也做了一件相似的事？因为使你困惑的那些话似乎有此含义吗？}如今祂并没有说：\Quote{圣子由自己什么也不能做，只有听见圣父命令什么，才做什么。}而是：\BibleQuote{圣子由自己什么也不能做，只有看见圣父做什么，才做什么。}看，若你这样理解：圣父做了某事，圣子注意观看，好看见祂自己也要做的事；而且那是与圣父所做的不同的另一件事。这圣父所做的，是藉谁做的呢？若不是藉圣子，若不是藉圣言，你便犯了亵渎福音的罪。\BibleQuote{因为万物是藉祂造的。}那么，圣父所做的，是藉圣言做的；若藉圣言做的，便是藉圣子做的。那么那位另一个人是谁，他注意观看，好做他看见圣父所做的另一件事？你不常说圣父有两个儿子：有一个独生子——唯一的。但因祂的仁慈，就祂的天主性而言是唯一的，就嗣业而言并非唯一。圣父使许多人与祂的独生子同为嗣业继承者；并未像祂那样由自己的本体生他们，而是由祂自己的家藉着祂收纳他们为义子。因为如圣经所证，\BibleQuote{我们蒙召，是为获得义子的名分。}\par

10. 那么你说什么？正是那位独生子亲自说话；那位独生子在福音中说话；圣言亲自将言语赐给了我们，我们亲耳听见祂自己说：\Quote{子不能由自己作什么，除非祂看见父作什么。}那么，父作事，为叫子看见该作什么；然而，父若不藉着子，什么也不作。的确你感到困惑了，异端者，的确你感到困惑了；但你的困惑如同服用了藜芦，使你得以痊愈。即使现在你也找不到你自己，我想，你甚至自己谴责自己的判断和属肉体的观点。放下肉体的眼睛，举起你心中所有的眼睛，注视神圣的事。你听见的确实是人言，而且是通过一个人、通过福音作者、通过福音，你作为人听见人言；但你听见的是关于天主圣言的事，为叫你听见人性的，从而认识神圣的。师傅制造了麻烦，为要教导；播下了难题，为要激发认真的关注。\Quote{子不能由自己作什么，除非祂看见父作什么。}祂本当接着说：\Quote{因为父所作的事，子也照样作。}但祂没有这样说，而是说：\Quote{凡父所作的，子也照样作。}父不作一些事，子另作一些事；因为父所行的一切，都是藉着子而行的。子复活了拉匝禄，难道父没有复活他吗？子使瞎子看见，难道父没有使他看见吗？父藉着子在圣神内行动。这是天主圣三；但天主圣三的行动是同一个，尊威是同一个，永恒是同一个，同样永恒是同一个，作为也是相同的。父不造一些人，子另造一些人，圣神又另造一些人；父、子和圣神造同一个人；父、子和圣神，独一天主，造了他。\par

11. 你观察到位的多位格，但承认天主性的统一性。因为由于位的多位格，才说：\Quote{让我们按我们的肖像和模样造人。}祂没有说：\Quote{我要造人，当我造他时你留心，好使你也能够造另一个。}祂说：\Quote{让我们造}；我听见多位格；\Quote{按我们的肖像}；我又听见多位格。那么天主性的单一性在哪里？读下文：\Quote{天主造了人。}经上说了\Quote{让我们造人}，却没有说\Quote{众神造了人。}统一性在它说\Quote{天主造了人}时被认出来了。\par

12. 那么那个属肉体的观点在哪里？愿它被羞愧、隐藏、消灭；愿天主的圣言向我们说话。即使现在，作为虔诚的人，作为已经相信的人，作为已经浸润于信德并已获得一些领悟的人，让我们转向圣言本身，转向光明之源，让我们一起说：\Quote{主啊，父永远行同样的事如同祢；因为凡父所作的，祢藉着祂而作。我们听见祢是起初就有的圣言；我们没有看见，却相信了。在那里我们也听见了随后的话，就是\InnerQuote{万物是藉着祢而造的。}那么凡父所作的，祂都藉着祢而作。因此祢所行的同样的事如同父。那么为什么祢愿意说：\InnerQuote{子不能由自己作什么}？因为我在祢身上看到某种与父的平等，因为我听见：\InnerQuote{凡父所作的，子也照样作。}我由此认出了平等，我领悟并尽可能理解：\InnerQuote{我与父原是一体。}那么祢不能作什么，除非祢看见父作什么，这是什么意思？这有什么含义？}\par

13. 或许祂会对我说，是的，对我们所有人说：\Quote{至于我所说的：\InnerQuote{子不能作什么，除非祂看见父作什么；}我的\InnerQuote{看见}你如何理解？我的\InnerQuote{看见}是什么？暂时放下祂为你们而取的奴仆的形象。因为在那奴仆的形象中，我们的主在肉身中有眼睛和耳朵，那个人形有同样的身体形状，如同我们所具有的，同样的肢体轮廓。那肉躯来自亚当；但祂不像亚当。所以主无论是行走在地上或海上，如祂所喜悦，如祂所愿，为凡祂愿的事，祂都能；祂看祂愿看的；祂定睛看，祂看见了；祂转开眼，没有看见；跟在祂后面的人，在祂前面的人，凡能被看见的；用祂肉身的眼睛，祂只看见面前的事。但从祂的天主性中，没有事物是隐藏的。暂时放下、放下，我再说，奴仆的形象，注视天主的形象，那是祂在世界被造之前所拥有的；在那形象中祂与父平等；由此接受并理解祂对你们所说的：\InnerQuote{祂虽具有天主的形象，却不以自己与天主同等为抢夺的。}在那里看祂，如果你能够，好使你能够看见祂的\InnerQuote{看见}是什么。}\Quote{在起初已有圣言。}圣言如何看见？圣言有眼睛吗，或者我们的眼睛是在祂里面，不是肉身的眼睛，而是虔诚心灵的眼睛？因为\Quote{心里洁净的人是有福的，因为他们要看见天主。}\par

14. 基督，你看见祂是人又是天主；祂向你显示为人，天主祂为你保留。现在请看祂如何为你保留天主，祂却向你显示为人。\BibleQuote{谁爱我，就遵守我的诫命}，祂说，\BibleQuote{谁爱我，必蒙我父爱他，我也要爱他。} 好像有人问：\Quote{祢将给祢所爱的人什么？} \BibleQuote{我也要向他显示我自己，}祂说，\BibleQuote{向他。} 这是什么意思，弟兄们？那位他们早已看见的，许诺祂要向他们显示自己。向谁？向那些已经见过祂的人，还是向那些未曾见过祂的人？祂曾对某一位宗徒说话，那位宗徒请求看见父，以为这样便满足，就说：\BibleQuote{请将父显示给我们，我们就满意了}——那时祂站在这个仆人的眼前，取了\OriginalTerm{仆人的样子}{the form of a servant}，为祂将来受荣显时，为祂的双眼保留\OriginalTerm{天主的形式}{the Form of God}，就对他说：\BibleQuote{我和你们在一起这么长久，你们还不认识我吗？谁看见我，就是看见了父。} 你请求看见父；看见我，你看见我，却没有看见我。你看见我为了你所取的东西，你并没有看见我为你所保留的东西。请听我的诫命，洁净你的眼睛。\BibleQuote{因为谁爱我，就遵守我的诫命，我也要爱他。} 对于遵守我诫命、并藉着我的诫命成为健全的人，我将向他显示我自己。\par

15. 那么，弟兄们，如果我们不能看见什么是圣言的\OriginalTerm{看见}{Seeing}，我们将往何处去？我们是否过于急切地要求什么样的\OriginalTerm{看见}{Seeing}？为什么我们渴望有人向我们显示我们目前无法看见的事物？因此，这些事是我们渴望看见的，而不是我们已经能够理解的。因为如果你看见圣言的\OriginalTerm{看见}{Seeing}，也许在你看见圣言的\OriginalTerm{看见}{Seeing}时，你将会看见圣言本身；这样圣言就不至于是一回事，而圣言的\OriginalTerm{看见}{Seeing}是另一回事，免得其中有任何结合、连接、双重、组合。因为它是单纯的，一种不可言喻的单纯。不像人一样，人和人的看见是两回事。因为有时人的看见会熄灭，而人却依然存在。这就是我所说的话，我要说一些并非所有人都能理解的事；愿主甚至赐予一些人能够明白。我的弟兄们，祂劝勉我们的目的，就是让我们看见，圣言的\OriginalTerm{看见}{Seeing}超出我们的能力；因为我们的能力很小，让它们被滋养、被成全。藉着什么？藉着诫命。什么诫命？\BibleQuote{谁爱我，就遵守我的诫命。} 什么诫命？因为我们早已渴望增长、坚强、成全，好使我们能看见圣言的\OriginalTerm{看见}{Seeing}。主，请告诉我们，现在是什么诫命？\BibleQuote{我给你们一条新诫命，你们要彼此相爱。} 那么，弟兄们，让我们从这泉源的丰盈中汲取爱德，让我们接受它；被它所滋养。接受那能使你能够接受的。让爱德生你，让爱德滋养你；爱德带你到成全，爱德坚强你；好使你能看见圣言的这\OriginalTerm{看见}{Seeing}，即圣言和祂的\OriginalTerm{看见}{Seeing}不是两回事，而是圣言的\OriginalTerm{看见}{Seeing}就是圣言本身；这样你或许很快就能明白所说的话：\BibleQuote{子不能凭自己做什么，除非祂看见父做什么}，就好像祂说：\Quote{子原本不会存在，如果祂不是由父所生。}\par

\SourceRecord{Translated by R.G. MacMullen.；Nicene and Post-Nicene Fathers, First Series；Vol. 6.；Edited by Philip Schaff.；Buffalo, NY: Christian Literature Publishing Co.,；1888.；<http://www.newadvent.org/fathers/160376.htm>.}

% structure: p0001 source=h1 target=section reason=page-title

\section{论新约讲道集第77篇}

\Emph{[CXXVII. Ben.]}\par

\Emph{论\BibleBook{若望}{福音}之言，\BibleRef{若5:25}，\BibleQuote{我实实在在告诉你们：时候要到，且现在就是，死人将要听见天主子的声音；听见的人将要生活。}等等；以及宗徒之言，\BibleQuote{眼睛未曾看见……}等等，\BibleRef{格前2:9}}\par

1. 诸位弟兄，我们的希望不属于现世，也不属于这个世界，也不在于那使人忘掉天主的幸福。这是我们应该首先知道的，并在基督徒的心中持守：我们成为基督徒不是为现世的好事，而是为别的东西，天主立刻应许，但人尚未理解。关于此善，经上说：\BibleQuote{眼所未见，耳所未闻，人心也未曾想到的，就是天主为爱祂的人所预备的。}因为此善如此伟大、如此卓越、如此不可言说，未进入人的理解，所以需要天主的应许。因为所应许的，人心盲目的现在不领悟；现在也不能向他显示应许将来要成为什么。因为一个婴孩，若能理解说话者的言语，他自己不能说话、行走、做任何事，而是如我们所见那样软弱，不能站立，需要他人帮助，他若只能理解对他说话的人并告诉他说：\Quote{你看，像我这样行走、工作、说话，几年后你也会像我一样。}当他考虑自己和对方时，虽然他看到了所应许的；但考虑到自己的软弱，他不会相信，然而他会看到所应许的。但我们像婴孩一样躺在这肉体和软弱中，所应许的既伟大又看不见；因此信德被激发，使我们相信我们看不见的，好使我们达到相信我们所看见的。凡嘲笑此信德，认为不该相信他自己看不见的，当他不信的事情到来时，便羞愧难当；慌乱之后被分离，分离之后被定罪。但谁若相信，将被放在右边，并要怀着极大的信心和喜乐站立在听到\BibleQuote{我父所祝福的，你们来吧！承受自创世以来为你们预备的国度。}的人中间。但主说完这些话后，这样结束：\BibleQuote{这些人要去进入永火，而义人却要进入永生。}这就是所应许给我们的永生。\par

2. 因为人类爱在这地上生活，所以生命被应许给他们；因为他们极其害怕死亡，所以永生被应许给他们。你爱什么？活着。你将要拥有这个。你怕什么？死亡。你将不会遭受它。这对人性的软弱来说似乎足够了，以至于应该说：\Quote{你将拥有永生。}人的心意可以领会这一点，藉其目前的状况，它可以在某种程度上领会将要是什么。但由于目前状况的不完美，它能领会到什么程度呢？因为他活着，不愿死去；他爱永生，希望永远活着，永不死亡。但那些将在惩罚中受折磨的人，甚至希望死去，却不能够。那么，活得很长，或永远活着，并非什么大事；但幸福地活着才是一件大事。让我们爱永生，并由此可知我们应当如何努力为永生劳苦，当我们看见那些爱现世生命的人——这生命只持续一时，而且必会结束——如此为它劳苦，以致当死亡的恐惧来临时，他们会做他们所能做的一切，不是要消除死亡，而是要推迟死亡。当死亡威胁时，一个人如何劳苦呢？藉着逃跑、躲藏、舍弃一切、赎回自己、劳苦、忍受折磨和不安、请医生，以及一个人能做的一切？看，当耗尽了他的一切劳苦和资源，他只能设法多活一点；永远活着，他不能。如果人们用如此大的劳苦、如此大的努力、如此大的代价、如此的热诚、如此的警醒、如此的小心，只为了能多活片刻；那么他们应当如何努力，才能永远活着呢？如果那些用尽一切办法努力推迟死亡、多活几天，以免失去几天的人被称为明智；那么那些如此生活以致失去永恒日子的人是多么愚蠢！\par

3. 那么，唯有这一点可以许给我们：天主的恩赐，按我们目前所拥有的程度，可能对我们何等甘饴；因为我们拥有它，我们生活，我们健康，这都是祂的恩赐。既然应许了永生，让我们在眼前呈现这样一种生命，它除去我们在此所遭受的一切不愉快。因为对我们来说，更容易找出那里没有什么，而不是那里有什么。看，我们在此生活；我们也将在那里生活。在此我们健康，当我们不生病时，身体没有痛苦；在那里我们也必健康。在此生顺利时，我们不受鞭笞；在那里我们也将不受任何鞭笞。假设这里有一个活在世上的人，健康，不受鞭笞；如果有人赐给他永远如此，且这美好景况永不终止，他将何等喜乐？何等欢欣？他将如何在无痛苦、无折磨、无生命尽头的喜乐中不能自制？如果天主只应许了我们刚才提到的、我刚才用我能用的言语描述和陈述的这些，那么如果这东西要出售，该用什么价格购买？该付出多大的金额？即使你拥有全世界，你所有的一切足够吗？然而它是要出售的；如果你愿意，就买吧。不要因为如此伟大的事物而过于不安，因为价格很大。它的价格不过是你所拥有的。现在，要获得任何伟大宝贵的东西，你会准备金、银、钱，或任何牲畜的增加，或你的产业所产出的果实，来购买我不知道的某种伟大卓越之物，以便在此世幸福生活。如果你愿意，也购买这个吧。不要看你有什么，而是看你是什么。这东西的价格就是你自己。它的价格就是你自己本身。付出你的自我，你便会拥有它。你为什么烦恼？为什么不安？怎么？你要寻找自己，还是购买自己？看，付出你自身，如你之所是，如此献给那东西，你便会拥有它。但你也许会说：\Quote{我是邪恶的，也许它不会接纳我。} 通过把自己献给它，你将会成为善的。把自己献给这信仰和应许，就是成为善。当你成为善时，你将成为这东西的价格；并且你将拥有，不仅我刚才提到的健康、安全、生命和无尽的生命；你不仅拥有这些，我还要除去其他东西。那里将没有疲倦和睡眠；那里将没有饥饿和饥渴；那里将没有成长和衰老；因为那里也不会有生育，数目保持完整。那里的数目是完整的；也不需要增加，因为那里没有减少的可能。看，我除去了多少东西，却还没有说那里有什么。看，那里已经有生命和安全；没有鞭笞，没有饥饿，没有饥渴，没有衰竭，没有这些；然而我还没有说：\BibleQuote{眼未见过，耳未听过，人心也未尝想到的。} 因为如果我已说了，那么所写的就是假的：\BibleQuote{眼未见过，耳未听过，人心也未尝想到的。} 因为那没有上升到人心中的东西，怎能上升到我的心中，使我说 \BibleQuote{那没有上升到人心中的东西}？它是被相信，而不是被看见；不仅未被看见，甚至未被表达。那么，如果未被表达，怎能被相信？谁相信他所没有听到的？但如果他听到以便相信，它就是被表达；如果被表达，它就被想到；如果被想到并被表达，它就进入人的耳中。而因为它若未被想到就不会被表达，它也已经上升到了人的心中。看，仅仅提出如此伟大的事物，已经使我们不安，以致我们无法用言语清楚地表达。那么谁能解释事物本身呢？\par

4. 让我们注意\WorkTitle{福音}；刚才主在说话，让我们做祂所说的事。\Quote{那信我的，} 祂说，\Quote{出死入生，不受审判。我实实在在告诉你们：时候要到，且现在就是，死人要听见天主子的声音，凡听见的必会生存。就如父在自己内有生命，祂也赐给子在自己内有生命。} 祂借着生祂而赐给了；祂在生祂时赐给了。因为子来自父，不是父来自子；父是子的父，子是父的子。我说子由父所生，不是父由子所生；子永远是，因此永远受生。谁能理解这个 \Quote{\OriginalTerm{永远受生}{always begotten}}？因为当任何人听到\InnerQuote{受生}时，他便会想：\Quote{因此有一个时候，那受生者不存在。} 那么我们要说什么？并非如此；在子之前没有时间，因为 \BibleQuote{万有是藉着祂造成的。} 如果万有是藉着祂造成的，时间也是藉着祂造成的；时间怎能先于子，因为时间是由子造成的？那么除去一切时间，子常与父同在。如果子常与父同在，而子又是受生的，祂永远是受生的；如果永远受生，那受生者常与生祂者同在。\par

5. 你会说：\Quote{这件事我从未见过，一位父亲生育，且始终与他所生之子同在；但生育者先来了，而他所生者随后在时间中到来。}你说得好：\Quote{我从未见过此事；}因为这属于\BibleQuote{眼所未见的事}。你问这事如何表达？它无法表达：\BibleQuote{耳朵未曾听见，也未曾上到人心里。}当信而朝拜，当我们相信时，我们朝拜；当我们朝拜时，我们成长；当我们成长时，我们领悟。因为当我们仍在这肉身上时，只要我们远离主，就那些看见这些事的圣天使而言，我们仍是婴孩，以信德为乳养，日后将以眼见为食。因为宗徒如此说：\BibleQuote{我们身在肉身时，是远离主的。因为我们凭信德行走，并非凭眼见。}我们终有一天会达到眼见，这是\Person{若望}在他书信中所应许的：\BibleQuote{亲爱的诸位，我们是天主的子女，将来如何，还未显明。}我们现在因恩宠、因信德、因圣事、因基督的血、因救主的救赎，已是天主的子女；\BibleQuote{我们是天主的子女，将来如何，还未显明。我们知道当祂显现时，我们将相似祂，因为我们将看见祂实在怎样。}\par

6. 看，我们被滋养是为了领悟什么？看，我们被滋养是为了拥抱和喂养什么？然而，所喂养之物不会减少，而喂养者得以支撑。因为现在食物通过食用支持我们；但所吃的食物会减少；但当我们开始以正义为食，以智慧为食，以那不死的食物为食时，我们立刻得到支撑，而那食物并不减少。因为如果眼睛知道如何以光为食，却不减少光；光不会因为被更多人看见而变少；它滋养更多眼睛，却仍与先前一样大：它们被滋养，而光不减少；如果天主将这赐给祂为肉身眼睛所造的光，祂自己——心眼的真光——又当如何？若有人向你赞美天主，就预备心吧。\par

7. 看哪，你的主对你说：祂说：\BibleQuote{时候将到}，\BibleQuote{且就是现在}。\BibleQuote{时候将到}，是的，那个时刻，\BibleQuote{就是现在，当}——什么时候？\BibleQuote{当死者听到天主子的声音，那些听到的人将活。}那么那些不听到的人，便不得活。什么是\Quote{那些听到的人}？那些服从的人。什么是\Quote{那些听到的人}？那些相信而服从的人，他们将活。因此，在他们相信并服从之前，他们是死的；他们行走，却是死的。他们行走虽是死的，这对他们有什么益处呢？然而，如果其中有人要经历肉身的死亡，他们会跑去，预备坟墓，包裹他，抬出去，埋葬他——死者埋葬死者；关于这些人曾说：\BibleQuote{让死人埋葬他们的死人。}这样的死人，因天主的话而被复活，以致活于信德。那些死于不信的人，被话所唤醒。主论及这个时刻说：\BibleQuote{时候将到，且就是现在。}因为祂以自己的话，唤醒了那些死于不信的人；宗徒论及这些人说：\BibleQuote{你这睡着的人，醒来吧，从死者中起来，基督必赐给你光明。}这就是心灵的复活，这就是内里之人的复活，这就是灵魂的复活。\par

8. 但这并非唯一的复活，还有身体的复活。凡在灵魂上复活的，在身体上也将复活，归于他的福祉。因为并非所有人在灵魂上复活；但所有人在身体上都将复活。我说，并非所有人在灵魂上复活，而是那些相信并服从的人；因为\BibleQuote{那些听到的人将活}。但正如宗徒所说：\BibleQuote{并非所有人都有信德。}如果并非所有人都有信德，那么也并非所有人都在灵魂上复活。当你们身体复活的时刻到来时，所有人都将复活；无论好人坏人，所有人都将复活。但谁若先已在灵魂上复活，身体复活将归于他的福祉；谁若未先在灵魂上复活，身体复活将归于他的咒诅。凡在灵魂上复活的，身体复活将归于生命；凡不在灵魂上复活的，身体复活将归于刑罚。既然主已将这灵魂的复活印在我们心中，我们众人皆应奔向它，并努力在其中生活，且持守至终，那么还有必要请祂也把身体的复活印在我们心中，这身体的复活将在世界末了发生。现在，请听祂是如何印下这一点的。\par

9. 当他说了：\BibleQuote{我实实在在地告诉你们：时候将要来到，且现在就是，死者，}即不信者，\BibleQuote{将要听见天主子的声音，}即福音，\BibleQuote{听见的人，}即服从的人，\BibleQuote{将要生活，}即成义，不再是不信者；当我说，他说了这话时，因为他看出我们也需要受教关于肉身的复活，并且不应当就此留下，他就接着说：\BibleQuote{因为圣父怎样在自己内有生命，也赐给圣子在自己内有生命。}这指的是灵魂的复活，灵魂的被赋予生命。然后他加上：\BibleQuote{并且赐给他行审判的权柄，因为他是人子。}这天主子，就是人子。因为如果天主子一直是天主子，而没有成为人子，他就不会拯救人子。那造人的，他自己成为他所造的，好使他造的不会灭亡。但他成为人子，仍继续是天主子。他成为人，是取了非他所是的，并非失去他所是的；仍是天主，成为人。他取了你，并未被你消融。因此他这样来到我们这里：天主子与人子，造物者与被造者，创造者与受造者；他母亲的创造者，由他母亲所生；他这样来到我们这里。作为天主子，他说：\BibleQuote{时候将要来到，且现在就是，死者将要听见天主子的声音。}他没有说：\BibleQuote{人子的声音，}因为他是在强调真理，即他与圣父相等。\BibleQuote{听见的人将要生活。因为圣父怎样在自己内有生命，也赐给圣子在自己内有生命；}不是通过分沾，而是在我们的天主内。但圣父本身有生命，他生了这样的圣子，使圣子本身有生命；圣子不是成为生命的分享者，而是他自己就是生命，我们应分享这生命；也就是说，他本身有生命，他自己就是生命。但他成为人子，是从我们取来的。天主子本身是他自己的；他成为人子是从我们取来的。天主子属于他自己，人子属于我们。那较小的，他从我们取来；那较大的，他给了我们。所以他死，是因为他是人子，而非因为他是天主子。然而天主子死了，但他是在肉身方面死了，并非就\BibleQuote{圣言成了血肉，居住在我们中间}而言。所以他死，是死于属于我们的；我们活，是活于属于他的。他不能死于他自己，我们也不能活于我们自己。所以，主耶稣作为天主，作为独生子，与生他的圣父同等，把这印在我们心中：我们若听，就必生活。\par

10. 但是，他说：\BibleQuote{并且赐给他行审判的权柄，因为他是人子。}因此，那形象要来审判。人子的形象要来审判；所以他说：\BibleQuote{并且赐给他行审判的权柄，因为他是人子。}这里的法官将是人子；那曾被审判的形象将要审判。请听且明白：先知早已说过：\BibleQuote{他们要仰望他们所刺透的。}他们将要看见那被长枪刺穿的形象。他将坐在审判席上，他曾站在审判席位。他将定真正罪人的罪，他曾被诬为罪人。他亲自来临，那形象必来临。这在福音中也可找到：当他在门徒眼前升天时，他们站着观看，天使的声音说：\BibleQuote{加里肋亚人，你们为什么站着？}等等。\BibleQuote{这位耶稣怎样升天，也要怎样降临。}什么是\BibleQuote{怎样降临}？就是以这同样的形象降临。因为\BibleQuote{赐给他行审判的权柄，因为他是人子。}现在看，这是多么合宜且正当，好使被审判者能看见审判者。因为被审判的有善有恶。\BibleQuote{但心地纯洁的人有福了，因为他们要看见天主。}因此，在审判中，仆人的形象要向善恶显露，天主的形象只保留给善人。\par

11. 善人们将要领受的是什么？看，我现在说出我先前没有说出的；然而我说出时，却又没有说出。因为我说了，在那里我们将会健康、平安、生活、没有鞭打、没有饥渴、没有衰竭、没有眼目的丧失。这一切我都说了；但我们将要得到的更大的是什么，我却没有说。我们要看见天主。这将是如此伟大，甚至如此伟大，以至于相比之下，其他一切都算不得什么。我说了我们将要生活，将要平安无恙，将不受饥渴，将不陷于疲惫，睡眠不会压迫我们。这一切，与那使我们看见天主的幸福相比，又算得了什么？因为天主现在不能如他本来那样显示，然而我们将要看见他；因此，\BibleQuote{眼所未见，耳所未闻，}这正是善人所要看见的，虔诚的人所要看见的，仁慈的人所要看见的，忠信的人所要看见的，凡在肉身复活中分得好命运的人将要看见的，因为他们曾在心灵的复活中有良好的服从。\par

12. 那么，恶人也能看见天主吗？关于此人，依撒意亚说：\BibleQuote{但愿不敬神的人被带走，使他看不见天主的荣耀。}不敬神的人和敬神的人都将看见那形态；而当宣告了\BibleQuote{但愿不敬神的人被带走，使他看不见天主的荣耀}这句话时，剩下的就是：对于敬神和善人，要成就主自己曾应许的事——当祂在此世居于血肉时，不仅被善人看见，也被恶人看见。祂在善人和恶人中讲论，并被人人看见：作为天主，祂是隐藏的；作为人，祂是显现的；作为天主，祂治理人；作为人，祂出现在人中间。我告诉你们，祂在他们中间说：\BibleQuote{谁爱我，就遵守我的诫命；那爱我的，必蒙我父爱他，我也要爱他。}就好像有人问祂：祢将给他什么？祂说：\BibleQuote{我要向他显明我自己。}祂何时说这话？当祂被人看见时。祂何时说这话？当祂甚至被那些不爱祂的人看见时。那么，祂要如何向那些爱祂的人显明自己呢？除非以一种那些爱祂的人当时未曾看见的形态。因此，既然天主的形态被保留，人的形态被显现；通过人的形态向人说话，显明可见，祂向众人——善人和恶人——显现了自己，却为爱祂的人保留了自己。\par

13. 祂何时要向爱祂的人显现自己？在肉身复活之后，那时\BibleQuote{不敬神的人将被带走，使他看不见天主的荣耀。}因为那时\BibleQuote{当祂显现时，我们将相似祂，因为我们将看见祂实在怎样。}这就是永生。因为我们先前所说的一切，与那永生相比都算不得什么。我们活着，算什么呢？我们健康，算什么呢？我们将看见天主，这才是大事。这就是永生；祂亲自说过：\BibleQuote{永生就是：认识祢，唯一的真天主，和祢所派遣的耶稣基督。}这就是永生：认识、看见、领悟、了解他们所相信的，感知他们尚不能领悟的。那时，心灵将看见\BibleQuote{眼所未见，耳所未闻，人心所未想到的}；在末了，将对他们说：\BibleQuote{我父所祝福的，你们来吧！承受自创世以来为你们预备了的国度。}那些恶人于是要进入永火。但义人呢？进入永生？什么是永生？\BibleQuote{永生就是：认识祢，唯一的真天主，和祢所派遣的耶稣基督。}\par

14. 因此，当谈到将来肉身的复活时，主没有让我们停留于此，祂说：\BibleQuote{祂并且赐给祂权柄施行审判，因为祂是人子。你们不要惊奇这事，因为时辰要来。}祂在此处没有加上\BibleQuote{现在就是}，因为这时辰将在将来，这时辰将在世界末了，这将是最后的时辰，将在最后的号角吹响时。\BibleQuote{你们不要惊奇这事，}因为我曾说：\BibleQuote{祂并且赐给祂权柄施行审判，因为祂是人子。不要惊奇。}我说这话的原因在于：祂作为人，理当被人审判。祂要审判哪些人呢？那些祂发现活着的人吗？不仅是他们，还有谁？\BibleQuote{时辰要来，那时，凡在坟墓里的。}祂如何表达那些在肉身中死了的人？\BibleQuote{凡在坟墓里的，}他们的尸首被埋葬，骨灰被覆盖，骨头四散，血肉不再为血肉，但为天主仍是完整的。\BibleQuote{时辰要来，那时，凡在坟墓里的，都要听见祂的声音，而出来。}无论善人恶人，都要听见那声音，而出来。坟墓的所有束缚都要被扯断；所有失去的，更确切地说，被认为失去的，都要被恢复。因为如果天主创造了那个本不存在的人，祂岂不能重塑那个曾经存在的人吗？\par

15. 我认为，当人们说\Quote{天主将再次使死者复活}时，并未说什么不可信的事，因为这是论及天主，而非论及人。这将要成就的事是伟大的，是的，是件不可信的事。但不要让它显得不可信，因为请看，是谁要成就它。那创造你的，经上说，将复活你。你曾经不存在，而现在你存在；一旦被造，难道你将不再存在吗？愿天主阻止你这样想！天主在创造那不存在的时，行了更奇妙的事；然而祂确实创造了那不存在的；难道借着那些祂使其成为自己所不是的人，竟不能相信祂能重造那曾经存在的吗？这就是我们这些原先不存在、后被造成的人，对天主的回报吗？这就是我们对祂的回报吗：我们不愿相信祂能复活祂所造的？这就是祂的受造物对祂的回报吗？\Quote{因此，}天主对你说，\Quote{我造了你，人啊，在你尚未存在之前，是为了使你不信我吗？你将成为你所是的，你曾能够成为你所不是的？}但你会说，\Quote{看，我在坟墓中所见的，是尘土、灰烬、骸骨；这些能重获生命，皮肤、组织、肉体，并复活吗？什么？我在坟墓中所见的这些灰烬、这些骸骨？}好吧。至少你在坟墓中看见了灰烬和骸骨；在你母亲的子宫里本是一无所有。你现在所见的，至少有灰烬和骸骨；在你存在之前，既无灰烬，也无骸骨；然而你却被造了，当时你完全不存在；难道你不相信这些骸骨（无论它们处于何种状态、何种形态，它们确实存在）将重新获得它们曾有的形态，而你也曾接受了你不曾拥有的吗？相信吧；因为如果你相信这一点，那么你的灵魂就会被提升。你的灵魂将被提升，\BibleQuote{现在}；\BibleQuote{时刻要来，且现在就是}；然后你的肉体将为了你的祝福而复活，\BibleQuote{那时，凡在坟墓里的，都要听见祂的声音，而出来}。因为你不应立刻欢喜，因为你听到\Quote{而出来}；要听下文，\BibleQuote{行过善的，复活进入生命；作过恶的，复活进入审判}。转向主，等等。\par

\SourceRecord{Translated by R.G. MacMullen.；Nicene and Post-Nicene Fathers, First Series；Vol. 6.；Edited by Philip Schaff.；Buffalo, NY: Christian Literature Publishing Co.,；1888.；<http://www.newadvent.org/fathers/160377.htm>.}

% structure: p0001 source=h1 target=section reason=page-title

\section{论新约的讲道词 78}

\Emph{[CXXVIII. Ben.]}\par

\Emph{论福音的话}，\BibleRef{若 5:31} \Emph{，\BibleQuote{我若为自己作证}等等；以及论宗徒的话}，\BibleRef{迦 5:16} \Emph{，\BibleQuote{你们应当顺从圣神行事，决不满足肉身的贪欲，因为肉身贪欲，}等等。}\par

1. 我们听到了神圣福音的话；主耶稣所说的这句话：\BibleQuote{我若为自己作证，我的见证便不真}，可能会使一些人不解。那么，真理的见证怎么不真呢？难道不是祂亲自说过：\BibleQuote{我是道路、真理、生命}吗？如果我们不可相信真理，那么该相信谁呢？因为人若不愿相信真理，他必定只想相信虚假。所以，这话是从他们的角度说的，你应当这样理解，并从这些话中得出此意：\BibleQuote{我若为自己作证，我的见证便不真}，也就是说，按照你们的想法。因为祂深知祂自己对自己的见证是真的；但为了软弱、难信、无悟性的人，太阳寻求了灯。因为他们视力的软弱不能承受太阳的耀眼强光。\par

2. 因此，若翰被找来为真理作证；你们听见了祂（基督）所说的话: \Quote{你们曾到若翰那里去，他是燃烧而发亮的明灯，你们情愿暂时在他的光中欢欣。} 这盏灯是为使他们羞愧而准备的，因为许久以前在圣咏中早已如此说过：\Quote{我为我的受傅者预备了一盏明灯。} 什么！为太阳而备的明灯！\Quote{我将使祂的仇敌蒙受羞辱：但在祂身上，我的圣化必要兴盛。} 因此，他们曾在某个地方藉着这若翰自己而蒙羞，当时犹太人对主说：\Quote{你凭什麼权柄做这些事？请告诉我们。} 祂回答他们说：\Quote{你们也告诉我，若翰的圣洗是从天上来的，还是从人来的？} 他们听后，便默不作声。因为他们心里立刻想：\Quote{我们若说\InnerQuote{从人来的}，百姓会用石头砸我们，因为他们认为若翰是先知。我们若说\InnerQuote{从天上来}，祂会对我们说：\InnerQuote{那么你们为什么没有相信他？}} 因为若翰为基督作证。所以，他们被自己的问题困住了心，陷入自己的网罗，便回答说：\Quote{我们不知道。} 黑暗的声音还能是什么呢？人不知道时，说\Quote{我不知道}是对的。但若知道却说\Quote{我不知道}，那便是反对自己的见证。他们深知若翰的高贵，也知道他的圣洗是从天上来的，但他们不愿同意若翰所见证的那一位。但当他们说\Quote{我们不知道}时，耶稣回答他们：\Quote{我也不告诉你们我凭什麼权柄做这些事。} 他们便蒙羞了；于是应验了：\Quote{我为我的受傅者预备了一盏明灯，我将使祂的仇敌蒙受羞辱。}\par

3. 难道殉道者不是基督的见证人吗？他们不为真理作证吗？但若我们更仔细地思考，当那些殉道者作证时，是祂亲自为自己作证。因为祂居住在殉道者内，使他们能为真理作证。请听一位殉道者，就是宗徒保禄：\Quote{你们愿意接受基督在我内说话的证明吗？} 当若翰作证时，住在若翰内的基督亲自为自己作证。让伯多禄作证，让保禄作证，让其余的宗徒作证，让斯德望作证，是住在他们众人内的祂亲自为自己作证。因为祂没有他们，仍是天主；他们没有祂，算什么？\par

4. 论到祂，经上说：\Quote{祂上升至高处，掳掠了俘虏，赐恩惠予人。} 什么是\Quote{掳掠了俘虏}？祂战胜了死亡。什么是\Quote{掳掠了俘虏}？魔鬼是死亡的主谋，而魔鬼自己藉着基督的死亡被掳掠。\Quote{祂上升至高处。} 我们知道有什么比天更高？在门徒们眼前，祂明明地升了天。我们知道，我们相信，我们宣认。\Quote{祂赐恩惠予人。} 什么恩惠？圣神。赐予如此恩惠的那一位，祂自己是什么呢？因为天主的仁慈何其伟大；祂赐予与人相等的恩惠；因为祂的恩惠是圣神，整个天主圣三，圣父、圣子、圣神，是一体天主。圣神给我们带来了什么？请听宗徒的话：\Quote{天主的爱，}他说，\Quote{已倾注在我们心中。} 你这乞丐，天主的爱是怎样倾注在你心中的？天主的爱如何或在哪里倾注在人心中？他说：\Quote{我们把这宝藏存在瓦器中。} 为什么在瓦器中？\Quote{为使那卓著的能力归于天主。} 最后，当他说了\Quote{天主的爱已倾注在我们心中，} 免得有人以为他靠自己拥有这天主的爱，他立刻补充说：\Quote{藉着赐予我们的圣神。} 因此，为使你爱天主，让天主居住在你内，并在你内爱祂自己，也就是说，让祂感动你、点燃你、光照你、唤醒你，去爱祂。\par

5. 因为我们这身体中有一场战斗；只要我们活着，我们就在争战；只要我们还在争战，我们就身处险境；但是，\BibleQuote{在这一切事上，我们靠着那爱我们的主，大获全胜。} 我们的这场争战，你们刚才在诵读宗徒书信时已经听到了。\Quote{全部法律，}他说，\Quote{都包含在一句话内：\BibleQuote{你应当爱你的近人如你自己。}} 这爱来自圣神。\BibleQuote{你应当爱你的近人如你自己。} 首先，看看你是否知道如何爱自己；然后我才会将你应当如爱自己一样去爱的近人托付给你。但若你还不知道如何爱自己，我恐怕你会欺骗你的近人如同欺骗你自己。因为若你喜爱不义，你就不爱自己。圣咏作证说：\BibleQuote{但喜爱不义的，恨恶自己的灵魂。} 既然你恨恶自己的灵魂，爱肉身又有什么益处呢？若你恨恶自己的灵魂却爱肉身，你的肉身固然要复活，但只是为了你的灵魂受折磨。因此，灵魂必须先被爱，它应当服从于天主，如此这服务才能维持其应有的秩序：灵魂归于天主，肉身归于灵魂。你愿意你的肉身服事你的灵魂吗？那就让你的灵魂服事天主吧。你应当被治理，才能治理别人。因为这场战斗如此危险，若你的主宰舍弃你，毁灭必然来临。\par

6. 什么争战？\BibleQuote{但如果你们彼此咬伤、吞噬，你们要小心，免得彼此消灭。但我对你们说：你们应随从圣神。} 我是在引用宗徒的话，这些话刚才从他的书信中诵读过。\BibleQuote{但我对你们说：你们应随从圣神，就绝不会去满足肉身的贪欲。} \BibleQuote{但我对你们说：你们应随从圣神，至于肉身的贪欲，}他没有说\Quote{你们将不会有，}也没有说\Quote{你们将不去做，}而是说\BibleQuote{你们将不会去满足。} 现在这是怎么回事，我要靠主的帮助，尽我所能告诉你们；请留心听，好使你们能明白，如果你们是在随从圣神。\BibleQuote{但我对你们说：你们应随从圣神，就绝不会去满足肉身的贪欲。} 让他继续说下去；也许这晦涩之处可以通过他后面的话更容易理解。因为我说过，宗徒不愿说\Quote{你们将不会有肉身的贪欲，}也不愿说\Quote{你们将不去做肉身的贪欲，}而是说\BibleQuote{你们将不会去满足肉身的贪欲，}这并非没有意义。他已将这争战摆在我们面前。如果我们事奉天主，我们就投身在这战斗中。那么接下来呢？\BibleQuote{因为肉身贪欲与圣神作对，圣神也与肉身作对；二者彼此敌对，以致你们不能做你们所愿意的事。} 这段话，若不被理解，听来极为危险。因此我焦虑不安，唯恐人因错误的解释而丧亡，遂靠主的帮助，承担起向你们的爱心解释这些话的责任。我们有足够的时间，我们一早就开始，午饭时间也不紧迫；在安息日这一天，那些饥渴天主圣言的人尤其惯于聚集。请听并留意，我将尽我所能谨慎地讲述。\par

7. 那么，我所说的\Quote{若不被理解，听来极为危险}是什么意思呢？许多人被肉性而该死的贪欲所胜，犯下各种罪恶和不洁，沉溺于如此可憎的污秽中，甚至提及都是可耻的；他们对自己说宗徒的这些话。请看宗徒所说的：\BibleQuote{以致我们不能再做你们所愿意的事。} 我不愿做这些事，我是被迫的，我被驱策，我被击败，\BibleQuote{我做我所不愿做的事，}正如宗徒所说。\BibleQuote{肉身贪欲与圣神作对，圣神也与肉身作对，以致你们不能做你们所愿意的事。} 你看到这段话若不被理解，听来多么危险。你看到这多么关乎牧者的职责：开启封闭的泉源，向干渴的羊群供应纯净无害的水。\par

8. 那么，战斗时不要甘愿被击败。请看他所描述的是怎样的战争、怎样的战斗、怎样的冲突，是在内里，在你自身之内。\BibleQuote{肉身贪欲与圣神作对。} 如果圣神不与肉身作对，那就去行奸淫。但如果圣神与肉身作对，我便看到一场斗争，但未见胜利，这是一场角力。\BibleQuote{肉身贪欲与圣神作对。} 奸淫有其快感。我承认它有快感。但是，\BibleQuote{圣神与肉身作对：}贞洁也有其快感。因此，让圣神战胜肉身；无论如何不要让肉身战胜圣神。奸淫寻求黑暗，贞洁渴望光明。你希望在人前怎样，就怎样生活；你希望在人前怎样，即使在人们视线之外也要那样生活；因为那位造你的，甚至在黑暗中也能看见你。为什么贞洁受到众人的公开赞美？为什么连奸淫的人也不赞美奸淫？\BibleQuote{那履行真理的，却来就光明。} 但奸淫有其快感。要反驳，要抵抗，要反对。因为你并非没有可用来战斗的武器。你的天主在你内，善神已赐给了你。尽管如此，我们的肉身仍被允许通过邪恶的暗示和真实的快乐与圣神作对。要确保宗徒所说的：\BibleQuote{不要让罪恶在你们必死的身体上为王。} 他没有说\Quote{不要让罪恶存在那里。} 它已经在那里了。这被称为罪，因为它藉着罪的代价临到我们。因为在乐园里，肉身不与圣神作对，那里也没有这场斗争，只有和平；但在悖逆之后，在人不愿服事天主而被交付给自己之后——然而并非被交付给自己以致他能拥有自己，而是被那欺骗他的所拥有——肉身开始与圣神作对。如今在善人身上，它与圣神作对；因为在恶人身上，它没有什么可与之作对的。因为哪里有圣神，它就在哪里与圣神作对。\par

9. 因为当祂说：\Quote{肉性欲相反圣神，圣神也相反肉性；}你不要以为这归功于人的神太多了。在你内与你作战，反对那个在你内反对你的，是天主的神。因为你本不站在天主面前；你跌倒了，破碎了；如同器皿从人手中掉到地上，你破碎了。因为你破碎了，因此你与自己为敌；因此你与你自己相反。愿你内没有与你相反之物，你就能在完整中站立。为使你知道这职责属于圣神，宗徒在另一处说：\Quote{如果你们随从肉性生活，必要死亡；但如果你们依赖圣神，治死身体的恶行，必要生活。}从这些话人立刻自高自大，好像凭自己的神就能治死身体的恶行。\Quote{如果你们随从肉性生活，必要死亡；但如果依赖圣神，治死身体的恶行，必要生活。}宗徒，请为我们解释，依赖什么神？因为人也有属于其本性的神，借此他是人。因为人由身体和神组成。关于这人的神，经上说：\Quote{除了人内里的神，谁知道人的事？}那么我看见人自己有属于其本性的神，我听到你说：\Quote{但如果依赖圣神，治死身体的恶行，必要生活。}我问，依赖什么神：我自己的，还是天主的？因为我听到你的话，仍为这含糊不解。因为当使用\Quote{神}这个词时，有时用于人的神，也用于牲畜的神，如经上记载：\Quote{凡在旱地上有生气的神都死了。}所以神这个词用于牲畜，也用于人。有时连风也称为神，如在圣咏中：\Quote{火、冰雹、雪、雾，执行祂命令的风暴之神。}既然\Quote{神}这个词有多种用法，宗徒啊，你说恶行要依赖什么神治死：我自己的，还是天主的神？听下文，就明白了。困难被下文所解除。因为当祂说了：\Quote{但如果依赖圣神，治死身体的恶行，必要生活；}立刻加上：\Quote{因为凡被天主的圣神引导的，都是天主的子女。}你行动，若你被引导；并且行动得好，若你被至善者引导。所以当祂对你说：\Quote{如果依赖圣神，治死身体的恶行，必要生活；}而你疑惑祂说的是什么神时，在后来的话中要认识那位师傅，承认那位救赎者。因为那位救赎者赐给了你圣神，使你借此能治死身体的恶行。\Quote{因为凡被天主的圣神引导的，都是天主的子女。}他们若不从天主的圣神引导，就不是天主的子女。但若被天主的圣神引导，他们就争战；因为他们有一位大能的助佑者。因为天主看我们的战斗不像人看角斗士。人们可以偏爱角斗士，但当角斗士遇险时，他们不能帮助他。\par

10. 所以在这里也一样：\Quote{肉性欲相反圣神，圣神也相反肉性。}\Quote{以致你们不能行你们所愿意的}是什么意思？因为对于误解的人，这里有危险。现在我当尽我的职责来解释它，尽管我多么不称职。\Quote{以致你们不能行你们所愿意的。}请听，圣人们，无论你们是谁，正在争战的人。我对在战斗的人说话。争战的人了解我；不争战的人不了解我。是的，争战的人，我不说他了解我，而是他先于我了。贞洁者的愿望是什么？是没有任何相反贞洁的欲情在他肢体中兴起。他愿享平安，但尚未得到。因为当我们到达那样的境界，不再有相反的欲情兴起时，就没有敌人需要与之争斗了；在那里胜利不是期待的事，因为已经战胜了已被制伏的敌人。听听宗徒自己的话怎样论及这胜利：\Quote{这必朽坏的必须穿上不朽坏的，这必死的必须穿上不死的。当这必朽坏的穿上不朽坏的，这必死的穿上不死的时，那时就应验了经上的话：\InnerQuote{死亡被胜利吞灭了。}}听听他们胜利时的欢呼：\Quote{死亡，你的胜利在哪里？死亡，你的刺在哪里？}你打击了，你伤害了，你推倒了；但为我受了伤的那位造了我。死亡啊，死亡，造我的那位为我受了伤，并藉祂的死亡战胜了你。然后他们在胜利中说：\Quote{死亡，你的胜利在哪里？死亡，你的刺在哪里？}\par

11. 但现在，当\BibleQuote{肉身切望相反圣神，圣神也切望相反肉身}时，这就是死亡的争斗；我们不能做我们所愿意的。为什么？因为我们愿意没有私欲偏情，但我们无法阻止。无论我们愿意与否，我们都有它们；无论我们愿意与否，它们诱惑、勾引、刺痛、扰乱我们，它们会起来。它们被压制，尚未被熄灭。肉身切望相反圣神，圣神也切望相反肉身，要到几时呢？即使人死了，也会这样吗？断乎不可！你脱去肉身，又怎能带着肉身的私欲偏情呢？不然，如果你战斗得好，你将被接引到安息中。从这安息，你被领受冠冕，而非被定罪；以便此后你被带到天国。我的弟兄们，只要我们生活在此世，情形就是这样；即使我们这些在这场战争中已年老的人也是如此，确实我们有较弱的敌人，但我们仍有它们。我们的敌人在某种程度上甚至因年龄而疲惫；然而，尽管它们疲惫，它们仍不停用它们所能的刺激来骚扰老年的安宁。年轻人的战斗更激烈；我们很清楚，我们经历过：因此\BibleQuote{肉身就切望相反圣神，圣神也切望相反肉身，使你们不能行你们所愿意的。}因为你们愿意什么，哦圣洁的人，善良的战士，基督的勇敢士兵？你们愿意什么？就是完全没有邪恶的私欲偏情。但你们无法避免。维持战争，盼望胜利。因为此刻你们必须战斗。\BibleQuote{肉身切望相反圣神，圣神也切望相反肉身，使你们不能行你们所愿意的；}也就是说，使你们完全没有肉身的私欲偏情。\par

12. 但你们要做你们所能做的；正如宗徒自己在另一处所说的，我已开始重复的：\BibleQuote{不要让罪在你们必死的身体上为王，使你们顺从它的情欲。}看，我所不愿的；邪恶的私欲偏情起来了；但不要顺从它们。武装自己，拿起战争的武器。天主的诫命是你们的武器。如果你们按应当的听我讲，你们甚至因我所说的而被武装了。\BibleQuote{不要让罪在你们必死的身体上为王。}因为只要你们带着一个必死的身体，罪就与你们战斗；但不要让它为王。什么是\BibleQuote{不要让它为王}？就是\BibleQuote{使你们顺从它的情欲}。如果你们开始顺从，它就为王了。顺从是什么，不就是\BibleQuote{将你们的肢体献给罪作不义的武器}吗？没有什么比这位老师更卓越的了。你们要我还有什么可解释的呢？做你们所听到的。不要将你们的肢体献给罪作不义的武器。天主藉着祂的圣神赐给了你们能力来约束你们的肢体。私欲偏情起来，约束你们的肢体；它起来后能做什么呢？约束你们的肢体；不要将你们的肢体献给罪作不义的武器；不要武装你们的敌人反对你们自己。约束你们的脚，使它们不行于不法之事。私欲偏情起来了，约束你们的肢体；约束你们的手远离一切邪恶；约束眼睛，使它们不游荡失迷；约束耳朵，使它们不喜悦地听情欲的话语；约束整个身体，约束两肋，约束其最高和最低的部分。私欲偏情能做什么？它知道如何起来。它不知道如何战胜。通过不断地无效起来，它甚至学会了不起来。\par

13. 让我们回到我最初从宗徒引用的被认为晦涩的话，现在我们将看到它们是清楚的。因为我引用了这句：宗徒并没有说\BibleQuote{随从圣神行事，就不会有肉身的私欲偏情}；因为我们必然有它们。那么他为什么不说\BibleQuote{你们不要满足肉身的私欲偏情}？因为我们满足它们；因为我们有私欲偏情。有私欲偏情本身就是在做。但宗徒说：\BibleQuote{现在已不是我在做，而是住在我内的罪。}那么你们要提防什么呢？无疑是这样：你们不要满足它们。一个该受罚的私欲偏情起来了，它起来，作了暗示；不要听它。它燃烧，而不平息，你们愿意它不燃烧。那么\BibleQuote{使你们不能行你们所愿意的}在哪里？不要给它你们的肢体。让它无效地燃烧，它会耗尽自己。所以在你们内，这些私欲偏情被满足。必须承认，它们被满足了。因此他说：\BibleQuote{你们不要满足。}那么就让它们不被满足。你们决定去做，你们就满足了。因为如果你们决定犯奸淫，而没有犯，因为没有找到地方，因为没有机会，或许因为你们似乎为之不安的女子是贞洁的；看，现在她是贞洁的，而你们是奸夫。为什么？因为你们满足了私欲偏情。什么是\BibleQuote{满足}？就是在思想中决定了犯奸淫。如果现在，但愿不如此，你们的肢体也行动了，你们就一头栽进死亡中了。\par

14. 基督复活了会堂长那死在屋里的女儿。她还在屋里，尚未被抬出。那在心里决意作恶的人也是这样；他是死的，却躺在里面。但如果他到了付诸行动的地步，他就已经被人从屋里抬出来了。但主也复活了那寡妇的儿子，是一位青年，在他被抬出城外死亡之时。所以我敢说，你在心里决意，如果你从自己的行为中回转，你未付诸行动之前就会痊愈。因为如果你心里悔改，你曾决意作某件恶劣、邪恶、可憎、该受诅咒的事；就在你当时躺着死去的地方，在里面，你就在里面复活了。但如果你已经完成了，现在你就被抬出来了；但你有一位会对你说：\BibleQuote{青年，我命你起来！}纵使你已行了那事，要悔改，立刻回头，不要走向坟墓。但就在这里，我还看到第三个死人，他甚至已被带到坟墓。他身上已有习惯的重量，泥土堆重重地压着他。因为他惯于不洁的行为，被无度的习惯重重地压着。在这里，基督也呼喊：\BibleQuote{拉匝禄，出来！}因为一个习惯极恶的人\BibleQuote{如今已经臭了。}基督在那情况下呼喊是理所应当的；不仅呼喊，而且大声呼喊。因为基督的呼喊，即使是这样的人，虽是死的，虽是埋葬的，虽是发臭的，但他们仍要复活，他们要复活。因为在这位复活者面前，没有一个躺在死亡中的人是我们应当绝望的。让我们转向主，等等。\par

\SourceRecord{Translated by R.G. MacMullen.；Nicene and Post-Nicene Fathers, First Series；Vol. 6.；Edited by Philip Schaff.；Buffalo, NY: Christian Literature Publishing Co.,；1888.；<http://www.newadvent.org/fathers/160378.htm>.}

% structure: p0001 source=h1 target=section reason=page-title

\section{\WorkTitle{论新约讲道集 79}}

\EditorialNote{CXXIX. Ben.}\par

\Emph{论福音的话}，\BibleRef{若 5:39}\Emph{，}\Quote{你们查考圣经，以为在其中得永生；正是这些为我作证。}\Emph{等等。驳斥\OriginalTerm{多纳徒派}{Donatists}。}\par

1. 亲爱的诸位，请留意刚才在我们耳中响起的福音读经；我要按天主赐予我的能力说几句话。主耶稣向犹太人说话，对他们说：\Quote{你们查考圣经，以为在其中得永生；正是这些为我作证。}过了一会儿，祂又说：\Quote{我因我父的名而来，你们却不接待我；若有别人因自己的名而来，你们反要接待他。}又过了一会儿：\Quote{你们既然彼此寻求光荣，而不寻求惟独从天主来的光荣，怎能相信我呢？}最后祂说：\Quote{我不在父前控告你们；有一位控告你们的，就是你们所信赖的梅瑟。如果你们相信梅瑟，必会相信我，因为他曾论及我写了书。既然你们不信他的话，怎能信我的话呢？}这些话是出自神圣的默感，借着读经者的口，但却是借着救主的服务摆在我们的面前。请听我讲几句话，不要以数量来估计，而要按分量来权衡。\par

2. 所有这些事就犹太人而言是容易理解的。但我们须当心，免得过于注意它们，而把我们的视线从自己身上移开。因为主是向祂的门徒讲话；无疑地，祂对他们所说的话，也是对我们这些他们的后代说的。祂所说的\Quote{看，我同你们天天在一起，直到今世的终结}，不仅适用于他们，也适用于以后及直到世末的所有基督徒。因此，祂向他们说：\Quote{你们要谨慎，提防法利塞人的酵母。}那时他们以为主说这话是因为他们没有带饼；他们不明白\Quote{提防法利塞人的酵母}就是\Quote{提防法利塞人的教训。}法利塞人的教训是什么？就是你们刚才听到的：\Quote{彼此寻求光荣，互相寻求光荣，而不寻求惟独从天主来的光荣。}关于这些人，宗徒保禄这样说：\Quote{我为他们作证，他们对天主有热心，但不是按着真知灼见。}他说：\Quote{他们对天主有热心；}我知道，我确信这点；我从前也在他们中间，像他们一样。\Quote{他们对天主有热心，但不是按着真知灼见。}宗徒啊，\Quote{不是按着真知灼见}是什么意思？请向我们解释你所提出的真知灼见是什么，你所遗憾他们没有而希望我们具有的真知灼见？他接着补充并阐释了他所提出的内容。\Quote{他们对天主有热心，但不是按着真知灼见}是什么意思？\Quote{因为他们不认识\OriginalTerm{天主的正义}{righteousness of God}，而想建立自己的正义，便不服从天主的正义。}因此，不认识天主的正义，而想建立自己的正义，这就是\Quote{彼此寻求光荣，而不寻求惟独从天主来的光荣。}这就是法利塞人的酵母。主嘱咐我们要提防它。如果主是仆人，而祂要我们提防，那么让我们提防；免得我们听到：\Quote{你们为什么称呼我：主啊！主啊！却不照我所说的去做？}\par

3. 那么，让我们暂时离开当时主向之讲话的犹太人。他们外在于我们，不愿听我们；他们仇恨福音本身，他们曾买通假见证控告主，要定活着的祂的罪；祂死后，他们又花钱买了别的见证。当我们对他们说：\Quote{请相信耶稣，}他们回答我们说：\Quote{难道我们要信一个死人吗？}但我们加上：\Quote{但祂复活了；}他们回答：\Quote{根本没有；}是祂的门徒从坟墓里偷走了祂。犹太的买主喜爱谎言，鄙视主的真理、救赎主。犹太人哪，你所说的，你的父母用钱买来的；他们所买的仍存留于你。你该留意那位赎买了你的，而不是那位为你买来谎言的人。\par

4. 但正如我所说，让我们撇下这些人，而更留意这些与我们有关的弟兄们。因为基督是身体的头。头在天上，身体在地上；头是主，身体是祂的教会。但你们记得经上说：\Quote{二人成为一体。}宗徒说：\Quote{这是伟大的奥秘，但我是指基督和教会说的。}既然他们二人成为一体，那么他们就是同一声音。我们的头主基督向犹太人说了我们听福音时听到的这些话——头向祂的仇敌说话；也让身体，即教会，向它的仇敌说话。你们知道它该向谁说话。它要说什么？我所说的声音是同一的，不是出于我自己；因为身体是同一的，声音是同一的。让我们向他们这样说；我是以教会的声音说话。哦，弟兄们，失散的子女，迷途的羊，砍下的枝条，你们为什么诽谤我？为什么不承认我？\Quote{你们查考圣经，以为在其中得永生；正是这些为我作证；}我们的头对犹太人这样说，身体也对你们这样说：\Quote{你们要找我，必找不着。}为什么？因为你们不\Quote{查考那为我作证的圣经。}\par

5. 关于头的证言；\BibleQuote{曾向亚巴郎和他的后裔颁赐恩许。并没有说\Quote{后裔们}，指许多人，而是说\Quote{你的后裔}，指一位，就是基督。} 关于身体向亚巴郎的证言，是宗徒提出的。\BibleQuote{曾向亚巴郎颁赐恩许。主说：我指着自己起誓，因为你听从了我的话，没有为了我怜惜你的独生子，我必赐福与你，使你的后裔繁多，如同天上的星辰，海边的沙粒，地上万民都要因你的后裔蒙福。} 这里你们有一个关于头的证言，和一个关于身体的证言。请听另一个简短的，几乎一句话就包含了头与身体的证言。圣咏谈到基督的复活；\BibleQuote{天主，请祢超越诸天。} 随即关于身体：\BibleQuote{愿祢的光荣普照大地。} 请听关于头的证言；\BibleQuote{他们穿透了我的手足，数遍了我所有的骨骸；他们且注视着我；他们分了我的衣服，为我的长衣拈阄。} 请听随即关于身体的证言，稍后几句：\BibleQuote{大地四极将回心转意，归向上主；万邦诸族将在祂面前叩拜；因为王国是主的，祂要统治万邦。} 请听关于头的证言；并且\BibleQuote{祂如同新郎走出洞房。} 在同一篇圣咏中，请听关于身体的证言：\BibleQuote{它们的声音传遍大地，它们的言语直到地极。}\par

6. 这些经文是针对犹太人和我们自己的弟兄们的。为什么这样？因为旧约圣经这些经文，犹太人和我们这些弟兄们都接受。但是基督自己，是那些人所不接受的，我们看看这些弟兄是否接受。让祂自己发言，既为祂自己——祂是头——说话，也为祂的身体——就是教会——说话；因为在我们身上，头也为身体说话。请听关于头的：祂从死者中复活了，发现门徒们犹豫、怀疑、因喜乐而不信；祂\BibleQuote{开启了他们的明悟，叫他们理解圣经，并对他们说：经上记载，基督必须受苦，第三天从死者中复活。} 这是关于头的；让祂也为身体说话：\BibleQuote{并且因祂的名向万邦宣讲悔改与赦罪，从耶路撒冷开始。} 那么让教会向她的敌人发言，让她发言。她确实清晰地发言，她并不沉默；只愿他们侧耳倾听。弟兄们，你们已听到了证言，如今请承认我。\BibleQuote{你们查考圣经，因你们认为在其中可以获得永生；正是这些为我作证。} 我所说的不是出于我自己，而是出于我的主；然而你们仍然转身，仍然背弃。\BibleQuote{你们彼此求光荣，而不求唯独从天主来的光荣，怎么能相信我？} \BibleQuote{因为他们不认识天主的正义，却对天主怀有热情，但不按照知识。因为他们不认识天主的正义，想建立自己的正义，就不服从天主的正义。} 除了说\Quote{是我圣化，是我成义；我所给的便是圣的}，还有什么是不认识天主的正义而想建立自己的正义呢？把属于天主的归于天主；人啊，认出属于人的。你不认识天主的正义，却想建立你自己的。你想成义我；你与我一同成义就够了。\par

7. 关于假基督，所有人都明白主所说的：\BibleQuote{我因我父的名而来，你们却不接待我；如果另有人因自己的名而来，你们反要接待他。} 但是我们也听听若望的话：\BibleQuote{你们听说过假基督要来，如今已出现了许多假基督。} 假基督中我们害怕的，无非是他尊崇自己的名，藐视主的名。那说\Quote{是我成义}的人又做什么呢？我们回答他：\Quote{我来到基督前，不是用脚，而是用心；我在哪里听到福音，就在哪里相信，就在哪里受洗；因为我信基督，我信天主。} 但他却说：\Quote{你不洁净。} \Quote{为什么？} \Quote{因为那时我不在那里。} \Quote{告诉我为什么我不洁净？我是一个在耶路撒冷受洗的人，例如在厄弗所人中间受洗，而你们所读的书信就是写给他们的，你们却藐视他们的平安。看，宗徒写信给厄弗所人；一个教会建立起来，存留至今；是的，存留着更丰硕的果实，存留着更多的数目，持守从宗徒领受的：\InnerQuote{若有人给你们宣讲福音，与你们所接受的不同，当受诅咒。} 如今怎样？你对我说什么？我不洁净吗？我在那里受洗，不洁净吗？} \Quote{不，你也不洁净。} \Quote{为什么？} \Quote{因为我不在那里。} \Quote{但那无处不在的曾在那里。那无处不在的，我信了祂的名，祂曾在那里。你从我不知道的地方来，其实不是来，而是希望我到你这固定不动的地方去，你对我说：\InnerQuote{你受洗不合法，因为我不在那里。} 试想谁在那里。若翰曾被告知：\InnerQuote{你看见圣神降下如同鸽子，并停留在谁身上，谁就是那以圣神施洗的。} 你寻求祂为你；但因为你忌妒我受了祂的洗，你反而失去了祂。}\par

8. 那麽，我的弟兄们，请理解我们和他们的言语，并看看你们会选择哪一个。我们这样说：\Quote{我们是圣的，天主知道；我们是不义的，这一点祂更清楚。不要将你们的希望寄托于我们，无论我们是怎样的人。如果我们好，就照着所写的去做：\InnerQuote{你们当效法我，如同我效法基督一样。}但如果我们坏，你们也不会因此被遗弃，也不会因此没有忠告：要听从祂说：\InnerQuote{他们所说的，你们要遵行；但他们所行的，你们不要效法。}}而他们却相反地说：\Quote{如果我们不好，你们就灭亡了。}看，这就是\Quote{另有一个要凭自己的名字而来}。那麽，我的生命要依赖你们，我的救恩要系于你们吗？难道我忘记了根基？基督不是磐石吗？那建在磐石上的，既不被风也不被洪水冲倒，岂不是真的吗？那麽，如果你们愿意，就同我一起到磐石上来，不要希望作我的磐石。\par

9. 让教会也说这些话：\Quote{如果你们相信梅瑟，也会相信我，因为他写的是关于我的事。}因为我是祂的身体，祂所写的正是关于我。梅瑟也写了关于教会的事。因为我引用了梅瑟的话：\Quote{在你后裔中，地上万民都要蒙受祝福。}梅瑟在第一卷书中写了这话。如果你们相信梅瑟，也会相信基督。因为你们轻视梅瑟的话，也必然轻视基督的话。祂说：\Quote{他们那里有梅瑟和先知，让他们听从他们吧。}\Quote{不，亚巴郎父亲，但如果有人从死者中到他们那里去，}他们就会听他。\Quote{祂说：如果他们不听从梅瑟和先知，即使有人从死者中复活，他们也不会相信。}这话是对犹太人说的，难道就不是对异端者说的吗？那位从死者中复活的说过：\Quote{基督必须受难，并从死者中复活，第三日复活。}我相信这事。他说：我相信。你相信吗？那你为什麽不相信接下来所说的呢？你相信\Quote{基督必须受难，并从死者中复活，第三日复活}——这是关于元首说的；也请相信接下来关于教会所说的：\Quote{要因祂的名，向万邦宣讲悔改和赦罪。}你为什麽相信关于元首的，却不相信关于身体的？教会对你做了什麽，以至你要，可以说是，斩首？你想除去教会的头，却相信元首，留下身体如同无生命的躯干。你向元首献媚，像个忠心的仆人，却毫无用处。那要斩首的人，会尽力同时杀死头和身体。他们耻于否认基督，却不耻于否认基督的话。基督，我们和你们都没有亲眼见过。犹太人见过并杀死了祂。我们没有见过祂，却相信；祂的话与我们同在。你们把自己与犹太人相比：他们轻视挂在树上的祂，你们轻视坐在天上的祂；因他们的怂恿，基督的牌子被立起；因你们的自高自大，基督的圣洗被抹去。但是，弟兄们，除此之外还有什麽呢？我们甚至要为骄傲的人祈祷，甚至要为自高自大的人祈祷，就是那些如此抬高自己的人。让我们替他们向天主说：\Quote{愿他们知道：主是祢的名字，}而\Quote{不是}人，\Quote{唯独祢是至高无上的，超越全地。}让我们转向主，等等。\par

\SourceRecord{Translated by R.G. MacMullen.；Nicene and Post-Nicene Fathers, First Series；Vol. 6.；Edited by Philip Schaff.；Buffalo, NY: Christian Literature Publishing Co.,；1888.；<http://www.newadvent.org/fathers/160379.htm>.}

% structure: p0001 source=h1 target=section reason=page-title

\section{新约讲道第八十篇}

\Emph{[CXXX. Ben.]}\par

\Emph{论福音经文}，\BibleRef{若 6:9}\Emph{，其中叙述了五个饼和两条鱼的奇迹。}\par

1. 这是一个伟大的奇迹，亲爱的诸位：用五个饼和两条鱼使五千人吃饱，剩下的碎块装满了十二篮。这是一个伟大的奇迹；但我们若注意行这奇迹的那一位，就不会对这所行的事感到惊奇了。祂使五个饼在分饼的人手中增多，正如祂使土里生长的种子增多一样：播下几粒种子，便装满整个谷仓。但是，因为祂每年都这样做，便没有人惊奇。使人不惊奇的，不是所行之事的微小，而是它的恒常。但当主行这些事时，祂对理解的人说话，不仅用言语，也用奇迹本身。五个饼代表了梅瑟法律的五卷书。旧法律与福音的麦子相比，只是大麦。在这些书中包含着关于基督的伟大奥秘。因此祂自己说：\BibleQuote{如果你们信了梅瑟，也必信我，因为他写的是关于我的。}但是，如同大麦的核仁藏在谷壳下，同样，在法律的奥秘的面纱下，基督也被隐藏着。当这些法律的奥秘被展开和揭示时，那些饼在被擘开时也就增多了。而我所向你们解释的这一切，就是为你们擘开了饼。那五千人代表了受五卷法律书治理的子民。十二篮是十二位宗徒，他们自己也被法律的碎块所充满。那两条鱼要么是爱天主和爱近人的两条诫命，要么是受割损和未受割损的两个民族，或是君王和司祭这两种神圣身份。当这些事被解释时，就被擘开；当被理解时，就被吃下。\par

2. 让我们转向行这些事的那一位。祂自己就是\BibleQuote{自天而降的食粮}；这食粮使衰败者焕然一新，而自身不变坏；这食粮可以被品尝，不会被消耗。这食粮也是玛纳所预表的。因此经上说：\BibleQuote{祂赐给他们天上的食粮，人吃了天使的食粮。}谁是那天上的食粮呢？不就是基督吗？但是，为了使人类能吃天使的食粮，天使的主成了人。因为如果祂没有成为人，我们就不会拥有祂的肉体；如果我们没有祂的肉体，我们就不会吃祭台上的食粮。让我们赶快奔向产业，既然我们已藉此接受了产业的伟大凭据。我的弟兄们，让我们渴望基督的生命，既然我们把握着基督的死亡作为凭据。祂既承受了我们的恶事，岂不将祂的美物赐给我们呢？在这我们的地上，在这邪恶的世界里，什么占多数呢？不过是出生、劳苦和死亡。仔细考察人的处境，如果我撒谎，就定我的罪：请思考所有人类，难道他们在这世上还有别的目的吗？除了出生、劳苦和死亡？这就是我们国家的货物：这些在这里充满。那位商人就是为此货物而降临的。因为每个商人都是给予和接受：给予他所拥有的，接受他所没有的；当他购买某物时，他付出金钱，并收下所买之物。同样，基督在祂的这交易中也给予了和接受了。但祂接受了什么？就是这里充满的事物：出生、劳苦和死亡。祂给予了什么？重生、复活和永远为王。啊，善良的商人，买我们吧。我为什么说买我们呢？我们应当感谢祢，因为祢已经买了我们。祢把我们的赎价分施给我们，我们饮祢的血；这样祢就把我们的赎价分施给我们。我们阅读福音，我们的契据。我们是祢的仆人，我们是祢的受造物：祢造了我们，祢救赎了我们。任何人都可以买他的仆人，却不能创造他；但主却创造了祂的仆人并救赎了他们：创造他们，使他们存在；救赎他们，使他们不再成为俘虏。因为我们落入了今世的首领之手，他诱惑了亚当，使他成为自己的仆人，并开始占有我们作他的奴隶。但是救赎者来了，那诱惑者被击败了。我们的救赎者对那抓住我们为俘虏的做了什么？祂伸出祂的十字架作为陷阱，作为赎价；在十字架上祂放置了祂的血作为诱饵。那诱惑者确实有能力倾流祂的血，却不能饮祂的血。既然他倾流了并非债务人的血，他就被命令交出债务人；他倾流了无罪者的血，就被命令从有罪者身上退出。祂倾流祂的血正是为了这个目的，为要洗去我们的罪过。那曾经抓住我们的东西，就被救赎者的血抹去了。因为那诱惑者只能靠我们自己的罪过捆绑我们。它们是俘虏的锁链。祂来了，用祂苦难的锁链捆绑了那强者；祂进入他的家，即那些他居住者的心中，夺走了他的器皿。我们就是他的器皿。他曾用他自己的苦味充满我们。这苦味他也曾用苦胆向我们的救赎者献上。他既充满我们作他的器皿；但是我们主抢夺他的器皿，使之成为自己的，倒出苦味，代以甘饴。\par

3. 那么，让我们爱祂，因为祂是甘饴的。\BibleQuote{请尝并看看主是何等甘饴。} 祂当受敬畏，但更当受爱慕。祂是人与天主；唯一的基督是人与天主；正如一人由灵魂与身体组成：但天主与人并非两个位格。在基督内确实有两个性体：天主与人；但只有一个位格，为使天主圣三得以保持，而不因人性的加入而引入一个\OriginalTerm{四位}{quaternity}。那么，天主怎么会不怜悯我们呢？为了我们，天主竟成了人。祂已做的何其伟大，祂所做的比祂所应许的更奇妙；并且凭借祂已做的，我们应当相信祂所应许的。因为祂所做的，我们几乎不会相信，除非我们也看见了。我们在哪里看见呢？在信众之中，在已归向祂的群众中。因为向亚巴郎所应许的已经应验了；从我们看见的这些事，我们相信我们所看不见的。亚巴郎是一个独一的人，曾对他说：\BibleQuote{地上万民都要因你的后裔蒙福。} 如果他只看自己，他何时会相信呢？他是一个独一的人，且已年老；他有一个不生育的妻子，而且她年事已高，即使不是不生育，也不能怀孕。没有任何可寄托希望的东西。但他仰望那赐下应许者，相信了他所看不见的。看，他所相信的，我们看见了。因此，从我们看见的这些事，我们应当相信我们所看不见的。他生了依撒格，我们没有看见；依撒格生了雅各伯，我们也没有看见；雅各伯生了十二个儿子，我们也没有看见；他的十二个儿子生了以色列民族；这个伟大的民族我们看见了。我现在开始提到那些我们确实看见的事。从以色列民族中生出了童贞玛利亚，她生了基督；并且，看，在基督内万民都蒙福了。还有什么比这更真实？更确定？更明显？与我们一起，在未来的世界之后，你们这些从万民中聚集而来的人。在这个世界上，天主已实现了祂关于亚巴郎后裔的应许。祂使我们成为亚巴郎的后裔，怎会不赐给我们祂永恒的应许呢？因为宗徒说：\BibleQuote{如果你们属于基督（这是宗徒的话），那么你们就是亚巴郎的后裔。}\par

4. 我们已经开始成为某种伟大之物；愿无人轻视自己：我们曾经什么都不是，但现在我们成了什么。我们曾对主说：\BibleQuote{求祢记得我们不过是尘土；} 但祂从尘土造了人，赐生命给尘土，并已在我们主基督内将这同样的尘土提升到天国。因为祂从这尘土取了肉身，从这尘土取了泥土，并将泥土升到天上，祂是造天地的主。如果这两件新事，尚未成就，摆在我们面前，并有人问我们：\Quote{哪一件更奇妙？——祂是天主却成了人，或他是人却成了天主的人？哪一件更奇妙？哪一件更难？} 基督应许了我们什么？是我们至今尚未看见的，即我们成为祂的人，与祂一同为王，永不再死。这可以说是难以相信的，即一个生而为人的人能达到那永不再死的生命。这是我们要以洁净的心相信的，我指的是洁净了世界尘土的心；免得这尘土遮蔽我们信德的眼目。这就是我们被要求相信的：在我们死后，我们将与我们的尸体一起进入生命，永不再死。这很奇妙；但基督所做的更加奇妙。因为哪一件更难以置信：人永远活着，还是天主曾死过？人从天主那里接受生命是更可信的；天主从人那里接受死亡，我想是更难以置信的。然而这已经实现了；那么让我们相信那将要来的吧。如果那更难以置信的事已经实现，祂岂不会把那更可信的事赐给我们吗？因为天主有能力使凡人成为天使，祂从卑贱污秽的产物造了人。我们将成为什么？天使。我们曾经是什么？我不愿回想；我被迫思考，却羞于说出。我们曾经是什么？天主从哪里造了人？我们存在之前是什么？我们什么都不是。当我们在母腹中时，我们是什么？你们记得就够了。把你们的心思从你们被造的原料移开，想想你们是什么。你们活着；但草木也活着。你们有感觉，牲畜也有感觉。你们是人，你们超越了牲畜，你们高于牲畜；因为你们明白祂为你们行了何等大事。你们有生命，有感觉，有理智，你们是人。但这个恩惠有什么可比呢？你们是基督徒。因为如果我们没有领受这个，我们是人又有什么益处呢！所以我们是基督徒，我们属于基督。尽管全世界都发怒，它不能摧毁我们，因为我们属于基督。尽管全世界都抚慰，它不能引诱我们，因为我们属于基督。\par

5. 弟兄们，我们找到了一位大保护者。你们知道，人很依赖他们的保护者。一个有权势者的依附者会对任何威胁他的人回答说：\Quote{只要我主人的头安全，你对我便无可奈何。} 我们更能勇敢而确定地说：\Quote{只要我们的头安全，你对我们便无可奈何。} 因为我们的保护者就是我们的头。凡依赖任何人作为保护者的，都是他的依附者；我们是保护者的肢体。愿祂将我们怀在祂内，愿无人将我们从祂身边夺走。因为我们在今世将忍受的任何劳苦，一切都是暂时的，转瞬即逝。那不会消逝的美善事物将要到来；我们经过劳苦到达它们。但当我们到达后，无人能将我们从它们那里夺走。耶路撒冷的城门是关闭的；它们还上了门闩，好对那座城说：\BibleQuote{耶路撒冷，请你赞颂上主！熙雍，请你赞颂你的天主！因为祂巩固了你的门闩，祝福了在你内的子女；并使你的四境平安。} 当城门关闭，门闩插上时，没有朋友出去，没有仇敌进来。如果我们在此世没有放弃真理，在那里我们将享有真实而确定的平安。\par

\SourceRecord{Translated by R.G. MacMullen.；Nicene and Post-Nicene Fathers, First Series；Vol. 6.；Edited by Philip Schaff.；Buffalo, NY: Christian Literature Publishing Co.,；1888.；<http://www.newadvent.org/fathers/160380.htm>.}

% structure: p0001 source=h1 target=section reason=page-title

\section{论新约（讲道第81篇）}

[CXXXI. 本笃版]\par

\Emph{论福音之言}，\BibleRef{若 6:53}\Emph{，\BibleQuote{你们若不吃我的肉}等等，以及论宗徒之言。与\BibleBook{}{圣咏}。驳斥\OriginalTerm{白拉奇派}{Pelagians}。}\par

\EditorialNote{在殉道者圣西彼廉的祭台，于九月二十三日主日宣讲。}\par

1. 我们已听见了真正的导师、神圣的救赎者、人类的救主，向我们推荐我们的赎价，即他的血。因为他向我们谈及了他的身体和血；他称他的身体为肉，他的血为饮料。信徒们认出这属于信徒的圣事。但听者，除了听，还能做什么呢？因此，当推荐这样的肉和这样的饮料时，他说：\BibleQuote{你们若不吃我的肉，不喝我的血，在你们内便没有生命；}（这有关生命的话，除了生命本身，谁还能说呢？但那认为生命是虚假的人，将拥有死亡，而非生命。）门徒们听了，便起了反感，并非全部，但很多人，心中说：\BibleQuote{这话生硬，谁能听得下去呢？}但主在自己内知道这事，并听到他们心中的怨言，便回答他们，虽未发声，却思考着，好使他们明白他们被听见了，而不再存有这种念头。那么他回答了什么？\BibleQuote{这使你们跌倒吗？}\BibleQuote{那么，如果你们看见人子升到他先前所在的地方，又怎样呢？}这是什么意思？\BibleQuote{这使你们跌倒吗？}\BibleQuote{你们以为我要将我现所显示的这身体分割，并将我的肢体切开，分给你们吗？\InnerQuote{那么，如果你们看见人子升到他先前所在的地方，又怎样呢？}}确实，他能完整地升天，就不能被消耗。因此，他既将自己的身体和血赐给我们作为健康的滋养，又简短地解决了关于他自身\OriginalTerm{完整性}{Entireness}的如此重大疑问。那么，让那吃的人继续吃，喝的人继续喝；让他们饥渴；吃生命，喝生命。这吃，就是被更新；但你被更新，而那更新你的事物却不会减少。这喝，岂不就是生活？吃生命，喝生命；你将拥有生命，而生命是完整的。但这事将如此，即基督的身体和血将成为每人的生命；如果那在圣事中可见地领受的，在真理本身中属神地吃、属神地喝。因为我们亲耳听见了主自己说：\BibleQuote{使生活的是神，肉一无所用；我给你们所讲论的话，就是神，就是生命。但你们中间有些人，}他说，\BibleQuote{不信。}这些人就是曾说\BibleQuote{这话生硬，谁能听得下去呢？}的人。这话生硬，但只对心硬者而言；也就是说，它难以置信，但只对不信者而言。\par

2. 但是，为教训我们，这信德本身是恩赐，而非功绩，他说：\BibleQuote{正如我对你们说过的，除非蒙我父恩赐，没有人能到我这里来。}至于主在何处说这话，如果我们回想福音的前文，我们会发现他曾说：\BibleQuote{除非派遣我的父吸引他，没有人能到我这里来。}他不是带领，而是吸引。这种暴力是施加于心灵，而非身体。那么你为什么惊讶？相信，你就来了；爱，你就被吸引。不要以为这里有任何粗暴艰难的暴力；它是温柔的，是甘饴的；正是这甘饴吸引你。当一只饥饿的羊看到鲜草时，它难道不是被吸引吗？然而我认为它并非被身体驱赶，而是被欲望牢牢束缚。你也应以这样方式来到基督前；不要想象漫长的旅程；你相信的地方，你就到了那里。因为对于那位无处不在者，我们是藉着爱而来，而非藉着航行。但即便如此，在这航行中，各种诱惑的波涛和风暴层出不穷；相信被钉者；好使你的信德能登上那\OriginalTerm{木头}{Wood}。你不会沉没，反而会被木头托起。这样，甚至这样，在世俗波涛中航行的人正是那曾说：\BibleQuote{但我只以我们的主耶稣基督的十字架来夸耀}的人。\par

3. 但奇妙的是，当宣讲被钉的基督时，两人听见，一人轻视，另一人登上。让那轻视的人归咎于自己；让那登上的人不要自夸。因为他从真正的导师那里听见：\BibleQuote{除非蒙我父恩赐，没有人能到我这里来。}让他因已蒙恩赐而喜乐；让他以谦卑而非骄傲之心感谢那施恩者；以免他因谦卑所得，却因骄傲而失去。因为即使是那些已经在义路上行走的人，如果他们将此归功于自己或自己的力量，就会从此路中灭亡。因此，圣经教训我们谦卑，藉着宗徒说：\BibleQuote{你们要怀着恐惧战栗，成就你们的得救。}并且，以免他们因此将任何事归功于自己，因为他说了\BibleQuote{工作，}他立刻补充说：\BibleQuote{因为是天主在你们内工作，使你们愿意，并使你们力行，为成就他的善意。}\BibleQuote{是天主在你们内工作；}因此\BibleQuote{怀着恐惧战栗，}成为山谷，接受雨水。低地充满，高地干涸。恩宠是雨水。那么，如果\BibleQuote{天主拒绝骄傲人，却赏赐恩宠给谦卑人}，你有什么好惊讶的呢？因此，\BibleQuote{怀着恐惧战栗，}就是谦卑。\BibleQuote{不可心高气傲，而要恐惧。}恐惧，好使你充满；不要心高气傲，以免你干涸。\par

4. 但你会说：\Quote{我已在行这条路；一度我需要学习，需要藉法律的教导知道该做什么；如今我有自由意志的选择；谁能使我从这路上偏离呢？}你若仔细阅读，就会发现曾有一个人开始自高，因他拥有某一种丰裕，但这丰裕毕竟是领受的；然而仁慈的主为教导他谦逊，收回了所赐予的，他顿时陷入贫穷，在回忆中承认天主的仁慈，说道：\Quote{\BibleQuote{我在丰裕时说：我永不动摇。}} \Quote{\BibleQuote{我在丰裕时说。}} 但我说的，我这人说的；\Quote{\BibleQuote{所有人都是说谎者，我说。}} 因此，\Quote{\BibleQuote{我在丰裕时说；}}丰裕如此之大，竟使我敢说：\Quote{\BibleQuote{我永不动摇。}} 然后呢？\Quote{\BibleQuote{上主，祢以恩惠加增了我的尊荣。}} 但\Quote{\BibleQuote{祢转面不顾我，我就惊惶。}} \Quote{祢指示了我，}他说，\Quote{我之所以丰裕，是出于祢。祢指示了我当从哪里寻求，当归功于谁，当感谢谁，当在干渴时投奔谁，由谁得以充满，并与谁保守那使我得以充满的。\InnerQuote{因为我的力量，我要为祢保守；}我因祢的慷慨而充满，因祢的保守而不致失去。\InnerQuote{我的力量，我要为祢保守。}祢为指示我这事，\BibleQuote{祢转面不顾我，我就惊惶。}\InnerQuote{惊惶，}是因为干涸；干涸，是因为高傲。那么，你这干涸枯槁的人，为能重新充满，就当说：\BibleQuote{我的灵魂在祢面前如同无水之地。}说：\BibleQuote{我的灵魂在祢面前如同无水之地。}因为是你说的，不是主说的：\BibleQuote{我永不动摇。}你说了，仗恃自己的力量；但那并非出于你自己，你却以为是。}\par

5. 那么主说什么？\Quote{\BibleQuote{你们当怀着敬畏事奉上主，战战兢兢向他欢呼。}}同样，宗徒也说：\Quote{\BibleQuote{你们要怀着恐惧战兢，成就你们救恩的工夫，因为是天主在你们内工作。}}因此，要战战兢兢地欢乐：\Quote{\BibleQuote{免得上主发怒。}}我看你们用呼喊预先说出了我的话。因为你们知道我要说什么，你们用呼喊预先说出了。这从何而来？岂不是那因信而投奔的主教导了你们吗？那么，祂这样说；你们听已经知道的；我不是教导，而是讲道提醒你们；实在，我不是教导，因为你们已经知道，也不是提醒，因为你们记得，但让我们一起说出你们与我们一同保留的。\Quote{\BibleQuote{要接受训诲，且欢乐，}}但要\Quote{战战兢兢，}好谦卑地永远持守所领受的。\Quote{\BibleQuote{免得上主发怒；}}当然是对骄傲的人发怒，他们归功于自己，不感谢赐予他们的那一位。\Quote{\BibleQuote{免得上主发怒，你们便从义路上灭亡。}}他说的是\Quote{免得上主发怒，你们就不进入义路}吗？他说的是\Quote{免得主发怒，祂就不带你们进入义路}吗？或\Quote{不让你们进入义路}？你们已在其中行走，不要骄傲，免得甚至从其中灭亡。\Quote{\BibleQuote{你们便从义路上灭亡，}他说。}\Quote{\BibleQuote{当祂的怒火转瞬爆发时}}攻击你们。就在不远处。你们一骄傲，立刻失去所领受的。好像受这一切惊吓的人要说：\Quote{那么我该怎么办？}接着是：\Quote{\BibleQuote{凡投靠祂的，都是有福的：}}不是投靠自己，而是投靠祂。\Quote{\BibleQuote{我们得救是因恩宠，不是出于我们自己，而是天主的恩赐。}}\par

6. 或许你会说：\Quote{他是什么意思？为什么总说这个？一次二次三次，他总这么说；几乎每次说话，都不离此。}但愿我不是白说！因为有人对恩宠忘恩负义，过分归功于那贫乏软弱的本性。诚然，人受造时领受了巨大的自由意志能力；但由罪而失去了。他陷入了死亡，变得软弱，被强盗留在路上半死不活；那撒玛黎雅人——意即护卫者——经过时，把他扶上自己的牲口；他仍在被送往客栈的途中。为什么被扶起？他仍在治疗过程中。\Quote{但是，}他会说，\Quote{我在圣洗中领受了所有罪的赦免，这为我已经够了。}因为不义被涂去了，难道软弱也终结了吗？\Quote{我领受了，}他说，\Quote{所有罪的赦免。}这完全是真确的。所有罪都在圣洗圣事中被涂去，完全彻底，言语、行为、思想，一切都涂去了。但这正是那路上所倾倒的\BibleQuote{油与酒}。亲爱的弟兄们，你们记得那个被强盗打伤、半死不活地躺在路上的人，他如何因伤口敷上油与酒而得到力量。他的错误确实已被赦免，但软弱仍在客栈中接受治疗。客栈，如果你们认出，就是教会。在现世，它是客栈，因为我们在生命中经过；将来它将成为家，我们永不再离开，那时我们将完全康复，进入天国。同时，让我们欣然接受客栈中的治疗，既然仍旧软弱，就不可夸耀健康；免得因骄傲，除了永不受治之外，一无所获。\par

7. \BibleQuote{我的灵魂，请赞颂上主。}\BibleRef{咏103:1}你要说，是的，要对你的灵魂说：你还在这生命中，仍然随身带着脆弱的肉身，仍然\BibleQuote{可朽坏的肉身压迫着灵魂}\BibleRef{智9:15}；在获得完全的赦免之后，你仍接受了祈祷的救治；因为当你软弱仍在被医治时，你说：\BibleQuote{宽免我们的罪债}\BibleRef{玛6:12}。所以要对你的灵魂说，你这卑微的山谷，而非高耸的山丘；对你的灵魂说：\BibleQuote{我的灵魂，请赞颂上主，不要忘记祂的一切恩惠}\BibleRef{咏103:2}。是什么恩惠？述说它们，一一列举，感恩谢恩。是什么恩惠？\BibleQuote{祂赦免你的一切罪愆}\BibleRef{咏103:3}。这是在圣洗中发生的。现在发生什么？\BibleQuote{祂医治你的一切软弱}\BibleRef{咏103:3}。这现在发生，我承认。但当我仍在这里时，\BibleQuote{可朽坏的肉身压迫着灵魂}\BibleRef{智9:15}。所以也要说出那随后的话：\BibleQuote{祂救赎你的生命脱离败坏}\BibleRef{咏103:4}。在脱离败坏之后，还有什么呢？\BibleQuote{当这可朽坏的穿上了不朽，这必死的穿上了不死，}\BibleRef{格前15:54}那时就要应验那记载的话：\BibleQuote{死亡在胜利中被吞灭。}\BibleRef{依25:8}\BibleQuote{死亡啊，你的争战在哪里？}\BibleRef{格前15:55}那里正当地说：\BibleQuote{死亡啊，你的毒刺在哪里？}\BibleRef{格前15:55}你寻找它的所在，却找不到。\BibleQuote{死亡的毒刺}\BibleRef{格前15:56}是什么？\BibleQuote{死亡啊，你的毒刺在哪里？}\BibleRef{格前15:55}罪在哪里？你寻找，它无处可寻。因为\BibleQuote{死亡的毒刺就是罪}\BibleRef{格前15:56}。这是宗徒的话，不是我的。那时将要说：\BibleQuote{死亡啊，你的毒刺在哪里？}\BibleRef{格前15:55}罪将无处可寻，既不使你惊异，也不攻击你，也不灼烧你的良心。那时将不再说：\BibleQuote{宽免我们的罪债}\BibleRef{玛6:12}。但将说什么？\BibleQuote{上主我们的天主，求祢赐给我们平安；因为祢将一切赐给了我们。}\BibleRef{依26:12}\par

8. 最后，在脱离一切败坏之后，除了正义的冠冕，还有什么呢？这至少存留着，但即使在它之中或之下，也不要让头肿胀，好能接受冠冕。请听，请注意圣咏，那冠冕将如何不容许肿胀的头。在他说了\BibleQuote{祂救赎你的生命脱离败坏}\BibleRef{咏103:4}之后，他说：\BibleQuote{祂给你加冕}\BibleRef{咏103:4}。在这里你立即就要说：\Quote{\InnerQuote{给你加冕}是对我功绩的承认，我自己的卓越成就了它；这是债务的偿还，而非恩赐。}还是倾听圣咏的话吧。因为说这话的又是你；而\BibleQuote{所有人都说谎}\BibleRef{咏116:11}。请听天主所说：\BibleQuote{以慈悲和怜悯给你加冕}\BibleRef{咏103:4}。祂出于慈悲给你加冕，出于怜悯给你加冕。因为你本不值得祂召叫你，召叫后又使你成义，成义后又光荣你。\BibleQuote{剩下的因恩宠的拣选而得救。但如果是由于恩宠，就不再是由于行为；否则恩宠就不再是恩宠了。}\BibleRef{罗11:5-6} \BibleQuote{对那工作的，报酬不按恩宠计算，而是按债务计算。}\BibleRef{罗4:4}宗徒说：\BibleQuote{不按恩宠，而是按债务}\BibleRef{罗4:4}。但\BibleQuote{祂以怜悯和慈悲给你加冕}\BibleRef{咏103:4}；如果你的功绩先于恩宠，天主对你说：\Quote{好好检验你的功绩，你将看到它们是祂的恩赐。}\par

9. 这就是天主的正义。正如它被称为\Quote{主的救恩}\BibleRef{咏3:9}，并非主借此得救，而是祂赐给祂所拯救之人的；同样，借着我们的主耶稣基督而来的天主的恩宠，也被称为天主的正义，并非主借以成为正义的，而是祂借以使那不虔敬的人成义。然而有些人，如同古代的犹太人一样，既愿意被称为基督徒，却又不认识天主的正义，渴望建立自己的正义，甚至在我们这个时代，在公开恩宠的时代，在从前隐藏的恩宠完全启示的时代；在恩宠如今显现于打谷场上的时代，而这恩宠曾隐藏在羊毛里。我看到少数人理解了我，多数人没有理解，我绝不以沉默欺骗他们。古时的义人\Person{基德红}向主求一个记号，说：\BibleQuote{主啊，我求祢，让我放在打谷场上的这块羊毛沾满露水，而打谷场干燥。}\BibleRef{民6:36-40}果然如此；羊毛沾满露水，整个打谷场都是干的。早晨他在一个盆里拧干羊毛；因为恩宠赐给了谦卑的人；而在一个盆里，你们知道主对祂的门徒做了什么。他又求另一个记号：\BibleQuote{主啊，我愿意羊毛干燥，打谷场沾满露水。}\BibleRef{民6:36-40}果然如此。回想旧约时代，恩宠隐藏在云中，如同雨在羊毛里。现在请看新约时代，仔细观察犹太民族，你会看到它就像一块干羊毛；而整个世界，如同那打谷场，充满了恩宠，不是隐藏的，而是显现的。因此我们不得不极为哀叹我们的弟兄们，他们不是与隐藏的恩宠争战，而是与公开显现的恩宠争战。对于犹太人，尚有可原谅之处。我们该如何说基督徒呢？为什么你们是基督恩宠的敌人？为什么依靠自己？为什么忘恩负义？因为基督为什么来了？难道本性不早已在此吗？难道本性不在此吗？而你们不过是以过度的赞美欺骗它。难道法律不在此吗？但宗徒说：\BibleQuote{如果正义由法律而来，那么基督就白白死了。}\BibleRef{迦2:21}宗徒对法律所说的，我们也对这些人说关于本性的话：\Quote{如果正义由本性而来，那么基督就白白死了。}\par

10. 那么，从前关于犹太人所说的话，如今我们在这等人身上同样看到了。\BibleQuote{他们对天主有热情：我听见他们见证说，他们对天主有热情，但并非按照真知。} 什么是\BibleQuote{并非按照真知}？ \BibleQuote{因为他们不认识天主的义，并图谋建立自己的义，而没有服从天主的义。} 我的弟兄们，你们要同我一起分担忧伤。你们若见到这样的人，不要隐藏他们；愿你们没有这种误入歧途的慈悲；无论如何，你们见到这样的人，不要隐藏他们。要说服那些反对者，把那些抗拒的人带到我们这里来。因为关于这个问题，已有两个会议送交宗座，并且也有诏书从那里发出。这问题已得出结论；但愿他们的错误也有朝一日得以解决！因此我们劝告他们留心，我们教导他们受教，我们祈祷他们改过。让我们转向主等等。\par

\SourceRecord{Translated by R.G. MacMullen.；Nicene and Post-Nicene Fathers, First Series；Vol. 6.；Edited by Philip Schaff.；Buffalo, NY: Christian Literature Publishing Co.,；1888.；<http://www.newadvent.org/fathers/160381.htm>.}

% structure: p0001 source=h1 target=section reason=page-title

\section{论新约讲道集 第八十二篇}

\Emph{[CXXXII. Ben.]}\par

\Emph{关于福音的话}，\BibleRef{若 6:55} \Emph{，\Quote{我的肉确实是食物，我的血确实是饮料。谁吃我的肉，}等等}\par

1. 正如我们听到福音宣读时，主耶稣基督以永生的应许劝勉我们吃祂的肉、喝祂的血。你们这些听到了这些话的人，还没有全部理解它们。因为在你们当中，那些已经受过洗的信友确实知道祂的意思。但你们当中那些被称为慕道者或听道者的人，在宣读时可以做听众，但也能够做理解者吗？因此，我们的讲道是针对两者的。让那些已经吃主的肉、喝主的血的人，思想他们所吃所喝的是什么，免得如宗徒所说：\Quote{他们吃喝自己的审判。}但那些尚未吃喝的人，当他们被邀请到如此盛筵时，要赶紧前来。这些日子里，教师们喂养你们。基督每日喂养你们，祂的筵席常摆在你们面前。这是什么原因呢？哦，听道者啊，你们看见筵席，却不来赴宴？也许就在刚才福音被宣读时，你们心里说：\Quote{我们在想祂所说的话\InnerQuote{我的肉确实是食物，我的血确实是饮料}。主的肉如何吃，主的血如何喝？我们在想祂说了什么。}谁向你们关闭了它，使你们不知道这事？它上面蒙着一层帕子；但如果你们愿意，帕子就会被揭开。来宣认信仰，你们就解决了难题。因为主耶稣所说的，信友早已知道。但你们被称为慕道者，被称为听道者，却是聋子。因为你们身体的耳朵是开的，听见了所说的话；但你们心里的耳朵仍是闭的，因为你们不明白所说的话。我恳求，我不争论。看，复活节近了，报名受洗吧。如果节日不能唤醒你们，就让好奇心驱动你们：好使你们知道\InnerQuote{谁吃我的肉、喝我的血，就住在我内，我也在他内}的含义。好使你们与我一同知道：\Quote{敲，就给你们开；}正如我对你们说：\Quote{敲，就给你们开，}我也同样敲，请给我开。当我大声向耳朵说话时，我是在敲胸膛。\par

2. 但是，我的弟兄们，如果连慕道者也需要被劝勉不要迟延归向如此伟大的重生恩宠，那么我们对坚固信友应当何等谨慎，使他们的接近有益，使他们不把这样的盛筵吃喝到自己的审判中。现在，为了不让他们吃喝到审判，让他们好好生活。做劝诫者，不是用言语，而是用你们的行为；使那些尚未受洗的人，能这样赶紧跟随你们，而不会因效法你们而灭亡。你们这些已婚的人，要对你们的妻子保持婚床的忠贞。付出你们所要求的。作为丈夫，你要求妻子贞洁；给她榜样，而不是言语。你是头，看看你往哪里去。因为你应当去她跟随没有危险的地方；是的，你应当行在你要她跟随的路上。你要求从较弱的性别那里得到力量；你们两人都有肉身的私欲；让更强的人先得胜。然而，可悲的是，许多男人被女人征服。女人保持贞洁，男人却不愿保持；而他们不愿保持时，却想显得像男人：仿佛只有在性别上更强，以便敌人更容易征服他。这是一场斗争、一场战争、一场战斗。男人比女人更强，\Quote{男人是女人的头。}女人战斗并得胜；你却向敌人屈服？身体站立，头却低垂？但你们这些尚未有妻子，却又已接近主的祭台，吃基督的肉、喝祂的血的人，如果你们将要结婚，要为你们的妻子保守自己。正如你们愿意她们来到你们面前那样，她们也应当发现你们如此。哪个年轻人不愿娶一位贞洁的妻子？如果他娶了一位童贞女，谁不愿她是无玷污的？你寻求无玷污的，你自己也要无玷污。你寻求纯洁的，你自己不要不洁。因为并不是她能而你不能。如果是不可能的，那么她也不能如此。但是，既然她能，这就教你那是可能的。而她能有这能力，是天主管理她。但如果你做了，你将得到更大的光荣。为什么更大的光荣？父母的警惕约束她，较弱性别的羞耻心是她的缰绳；最后，她害怕你不害怕的法律。因此，如果你做了，你将得到更大的光荣；因为如果你做了，你敬畏天主。她有许多除了天主以外要怕的事，而你就怕天主一人。但你所怕的那位大于一切。祂在公开场合当怕，祂在暗中当怕。你出去，祂看见你；你进来，祂看见你；灯点亮，祂看见你；灯熄灭，祂看见你；你进入内室，祂看见你；在你内心的隐密处，祂看见你。你要怕祂，祂留心看你；甚至因这怕而成为贞洁。或者如果你要犯罪，找一个祂看不见你的地方，然后做你想做的事。\par

3. 但你们这些已经发愿的人，要更严格地克苦自己的身体，不要让自己放纵欲望，甚至对那些被允许的事；这样你们不仅远离非法的结合，而且轻视合法的眼神。记住，无论你们是男性还是女性，你们在地上过着天使的生活：\BibleQuote{因为天使不娶不嫁}。当我们复活后，也将如此。你们有多么更美好，在死亡之前就开始成为人在复活后将是的！保持你们适当的等级，因为天主为你们保留着荣耀。死人复活被比作天上的星辰。\BibleQuote{因为星辰与星辰的光荣不同}，如宗徒所说；\BibleQuote{死人复活也是这样}。因为那里童贞将用一种方式闪耀，婚配的贞洁用另一种方式闪耀，圣洁的寡居用另一种方式闪耀。它们闪耀得不同，但都在那里。光辉不等，天堂相同。\par

4. 因此，你们要想着自己的等级，持守自己的宣发，去接近主的体，接近主的血。凡知道自己并非如此的人，就不要接近。愿你们因我的话而感动悔改。因为那些知道自己为妻子保持贞洁、向妻子要求贞洁的人，那些知道自己完全保持节欲的人（如果他们已经向天主发了此愿），会因我的话感到喜乐；但那些听我说\Quote{你们中凡不保持贞洁的，就不要接近那饼}的人，则会悲伤。我不愿说这话，但我能做什么？我岂能怕人而抑制真理？若是那些仆人不敬畏主，难道我就因此而不敬畏吗？仿佛我不知道那话：\Quote{\BibleQuote{你这邪恶懒惰的仆人}，你应该分施，而我要求交账。}看，我已经分施了，我主我的天主；看，在你的眼前，在你圣天使的眼前，在这你子民的眼前，我拿出了你的银钱；因为我害怕你的审判。我已经分施，求你交账。即使我不说，你也会做。因此我宁愿说：我已经分施，求你使他们悔改，求你宽恕。使那些不贞洁的人成为贞洁，以便在你的眼前，当审判来临时，我们一同喜乐，无论是分施的人还是被分施的人。这话令你们喜悦吗？但愿如此！你们中凡不贞洁的人，趁你们还活着，要改正自己。因为我有权讲天主的话，但我没有权把那些在邪恶中坚持的不贞洁者从天主的审判和定罪中解救出来。\par

\SourceRecord{Translated by R.G. MacMullen.；Nicene and Post-Nicene Fathers, First Series；Vol. 6.；Edited by Philip Schaff.；Buffalo, NY: Christian Literature Publishing Co.,；1888.；<http://www.newadvent.org/fathers/160382.htm>.}

% structure: p0001 source=h1 target=section reason=page-title

\section{新约讲道集 第83篇}

\EditorialNote{【第一三三篇．本笃版】}\par

\Emph{论\BibleBook{若望}{福音}\BibleRef{若 7:6}等处的话，在那里耶稣说祂不上节期去，却仍照常去了。}\par

1. 我靠主的助佑，打算讲解刚刚诵读的这段福音；这里并非没有困难，以免真理受损，谎谬夸胜。并非真理可能消灭，或谎谬得胜。现在请暂时听听这段经文有什么困难；你们因着困难的提出而留心，请祈祷使我能充分解明。\BibleQuote{犹太人的帐棚节近了}；看来这些就是他们直到今日仍守的日子，那时他们搭盖帐棚。因为他们这个庆节因搭盖帐棚而得名；因为\OriginalTerm{帐棚}{\GreekText{σκηνὴ}}意思是\Quote{帐棚}，\OriginalTerm{搭帐棚}{\GreekText{σκηνοπηγία}}就是搭盖帐棚。这些日子在犹太人中被当作庆日；它被称为一个庆节，不是因为一天结束，而是因为持续地庆祝；正如逾越节庆日，和无酵饼庆日，然而，显然那庆节是连续多天举行的。这个周年节期临近了犹太，主耶稣在加里肋亚，祂在那里长大，那里也有祂的亲戚和亲属，圣经称他们为\BibleQuote{祂的弟兄}。\BibleQuote{因此，祂的弟兄}，正如我们听到的诵读，\BibleQuote{对祂说：你离开这里，往犹太去，好使你的门徒也看见你所行的事。因为没有人暗中行事，自己却想公开为人所知。你如果行这些事，就当将自己显示给世界。}然后福音作者接着说，\BibleQuote{原来连祂的弟兄也不信祂。}如果他们不信祂，他们说出的话是出于嫉妒。\BibleQuote{耶稣回答他们说：我的时候还没有到；你们的时候却常是准备好的。世界不能恨你们，却是恨我，因为我指证它所行的是恶的。你们上去过这个庆节吧。我不上去过这个庆节，因为我的时候还没有满。}接着福音作者说；\BibleQuote{祂说了这些话，自己仍住在加里肋亚。但祂的弟兄上去以后，祂也上去过节，不是明去，却似乎是暗中去的。}困难的界限就是这样，其余的都清楚。\par

2. 那么困难是什么？什么造成了困惑？什么处于危险中？恐怕主，是的，更直白地说，恐怕真理本身被认为说了谎。因为如果我们让人以为祂说了谎，软弱的人将得到说谎的权威。我们听说祂说了谎。因为那些认为祂说谎的人这样说：\Quote{祂说祂不上节期去，却上去了。}因此，首先，让我们尽量在时间的压力下看看，一个人说了话却不做，是否就说谎。例如，我告诉一个朋友：\Quote{我明天来看你}；一些更大的必要发生阻碍了我；我并非因此说了谎话。因为我承诺时，我是认真的。但后来发生的事阻止了我履行承诺，我并无意说谎，只是无法完成承诺。看，我想我没有费劲说服你们，我只是向你们的明智暗示：承诺某事而未履行者，并不说谎，如果他未履行是由于某事阻碍了承诺的实现，而不是虚假的证据。\par

3. 但有人听到我这么说，会问：\Quote{你岂能对基督说这样的话：祂不能完成祂所愿的，或祂不知道将来的事？}你说得好，提示得好，你的暗示正确；但是，人啊，请分担我的焦虑。我们岂敢说祂说谎，而我们不敢说祂在能力上软弱？我本人，凭我所想，按我所能判断的限度，我宁愿人在某件事上受骗，也不愿人在任何事上说谎。因为受骗是软弱的份，说谎是不义的份。他说：\BibleQuote{主，祢恨恶所有作恶的人。}紧接着：\BibleQuote{祢必消灭一切说谎的人。}要么\Quote{不义}和\Quote{谎言}是同一水平；要么\Quote{祢必消灭}比\Quote{祢恨恶}更重。因为被恨的人，不立即被毁灭。但那个问题：是否有必要说谎？我现在不讨论那个；那是一个晦暗的问题，有许多纠缠；我没有时间斩断它们直达核心。因此，它的处理推迟到另一个时间；也许不需要我的话，靠天主助佑就能治愈。但请注意并区分我今日想讨论的和推迟的。是否可以出于某种原因说谎，这个困难而最晦涩的问题我推迟。但基督是否说谎，真理是否说过虚假的话，由于福音经文提醒，我今天承担了。\par

4. 现在我要简要说明受骗与说谎之间的区别。受骗者认为他所说的是真的，因此他说出来，因为他认为是真的。如果受骗者所说的话是真的，他就不会受骗；如果它不仅是真的，而且他也知道它是真的，他就不会说谎。那么，他受骗，在于那是假的，而他认为是真的；但他之所以说出来，只是因为他认为是真的。错误在于人性的软弱，而不在于良心的正直。但凡认为那是假的，却断言为真的，他就是说谎。看，我的弟兄们，请分辨这个区别，你们这些在教会中长大、受教于主的圣经的人，不是无知、不是单纯、也不是愚昧的人。因为你们中间有博学多识的人，对各种文学并非平淡无奇地受过教导；至于你们中没有学过那种被称为自由文学的人，更是在天主的话语中滋养成长的。如果我在解释我的意思时费劲，请你们以聆听的注意和默想的深思来帮助我。你们若不先受帮助，也无法帮助。因此我们彼此代祷，并同样期待我们共同的助佑。受骗者，是当他说的话是假的，却以为是真；而说谎者，是认为一件事是假的，却把它当作真的说出来，不论那件事是真还是假。请注意我所添加的：\Quote{不论它是真还是假}；然而，谁认为是假的，却断言为真，他就说谎；他意在欺骗。因为如果那是真的，对他有什么好处呢？他始终认为那是假的，却说得好像真的一样。他所说的话本身是真的，它本身是真的；但就他而言，那是假的，他的良心并不认同他所说的话；他内心认为一件事是真的，却说出另一件事。他拥有\OriginalTerm{双重心}{double heart}，而不是单一的心；他没有把他内心所有的表达出来。双重心早已受到谴责。\BibleQuote{用欺谎的嘴唇，怀着双重心说了邪恶的话。} 如果说\BibleQuote{在一心中说了邪恶的话}就够了，那么\BibleQuote{欺谎的嘴唇}在哪里呢？什么是欺骗？就是一件事做了，另一件事假装。欺谎的嘴唇不是单一的心；因为不是单一的心，所以\BibleQuote{在一心中又在一心中}；因此两次\BibleQuote{在一心中}，因为心是双重的。\par

5. 那么我们如何看待主耶稣基督，说他曾说过谎？如果受骗比说谎是较小的恶，我们岂敢说他说谎，而我们却不敢说他受骗？但他既不受骗，也不说谎；而是诚然如经上所记（因为经上指的是他，也应当理解为指的是他）：\BibleQuote{没有一句虚假的话对君王说过，也没有一句虚假的话从他口中发出。} 如果这里的\InnerQuote{君王}是指任何一个人，让我们更愿将基督君王置于人君之上。但如果——这是更正确的理解——经上指的是基督，我说，如果这是更正确的理解，经上指的是基督（因为对他说来，确实没有虚假的话，因为他不受骗；从他口中没有虚假的话发出，因为他不说谎）；那么让我们看看该如何理解福音的这段经文，并且不要让我们，可以说，以属天的权威为谎言设下陷阱。但是，寻求解释真理，却又为谎言预备位置，这是极其荒谬的。我问你，这位向我解释这段经文的人，你要教我什么？我不知道你是否敢说\Quote{谎言}。因为如果你敢这样说，我就转过耳朵，用荆棘将其堵住，这样如果你想强行进入，我就可以借着荆棘的刺痛，在没有福音解释的情况下离开。告诉我你想要教我什么，你就解决了难题。请告诉我；我在这里，我的耳朵是敞开的，我的心准备好了，教我。但我问的是什么呢？我不愿绕许多圈子。你将要教我什么？无论你要提出什么学识，无论你在辩论中表现什么力量，只告诉我一件事，我问两件事之一：你要教我真理还是谎言？我们猜想他会如何回答，以免我离开；免得当他张着嘴努力说出他的话时，我立刻离开他：除了真理，他还会承诺什么？我正在聆听，站着，期待，极其热切地期待。看，这个承诺教我真理的人，却暗示关于基督的谎言。那么，一个说基督是假的的人，怎么能教我真理呢？如果基督是假的，我还能希望你对我说真话吗？\par

6. 再想想。祂说了什么？基督说了谎吗？我问你，在哪里？\Quote{祂说：\BibleQuote{我不上这个庆节去}；却上去了。}就我而言，我切愿仔细查考这处经文，或许我们可见基督并未说谎。诚然，我毫不怀疑基督没有说谎，我要么透彻查考并理解这段经文，要么不理解时就暂且搁下。但我绝不说基督说了谎。即使我不理解，我也将无知地离去。因为怀着虔敬无知，优于狂乱地下判断。不过，我们仍在努力查考，但愿借着祂（祂就是真理）的帮助，我们能发现些东西，我们自己也被发现是些东西，而这东西在真理中不会是谎言。因为如果在查考中我发现谎言，那我不是发现一个东西，而是一个虚无。那么，让我们来看看你说基督说谎的地方在哪里。他会说：\Quote{在于祂说：\BibleQuote{我不上这个庆节去}，却上去了。}你从哪里知道祂这样说了？假如我说——不，不是我，而是任何人（因为天主不许我说）——假如另一个人说：\Quote{基督没有说这话}，你如何反驳他，如何证明呢？你会翻开书，找到那处经文，指给那人看，甚至若他抗拒，你会极其自信地把书塞给他：\Quote{拿着，看，读，你手里拿的就是福音。}但为什么，我问你，你为什么如此粗暴地对待这个软弱的人？不要这样急切；更平静、更安宁地说。看，我手里拿的是福音；里面有什么？他回答：\Quote{福音宣告基督说了你所否认的话。}你会因为福音宣告了基督说了这话就相信吗？他说：\Quote{当然是因为这个理由。}我极其惊奇，你怎能说基督说谎，而福音不说谎。但恐怕当我说福音时，你想到的是那本书本身，以为羊皮纸和墨水就是福音；请看希腊文的意思：福音是\Quote{好信使}，或\Quote{好消息}。那么信使不说谎，而派遣他的那位会说谎吗？这位信使，即圣史，也提他的名号，这位写了这些的若望，他论基督是说了谎还是说了真话？你选择哪个，我都愿听。若他说谎，你便无法证明基督说了那些话。若他说实话，真理不能从谎言的泉源流出。谁是泉源？基督；让若望作溪流。溪流来到我这里，你对我说：\Quote{放心地喝罢}；然而你却使我警惕泉源本身，你告诉我在泉源中有谎话，你却对我说：\Quote{放心地喝罢。}我喝什么？若望说了什么，说基督说了谎？若望从哪里来？从基督来。那从祂来的，当从祂来的那位说了谎，他能对我说真话吗？我明明在福音中读到：\BibleQuote{若望靠在主的胸怀}；但我断定他饮入了真理。他靠在主胸怀时看见了什么？他饮入了什么？除了他倾流出来的，还有什么？\BibleQuote{在起初已有圣言，圣言与天主同在，圣言就是天主。这圣言在起初就与天主同在。万物是借着他造成的；凡受造的，没有一样不是借着他造成的。在他内有生命，这生命是人的光。光在黑暗中照耀，黑暗不能胜过他。}然而祂照耀，虽我偶有昏暗，不能完全理解祂，祂仍照耀。\BibleQuote{曾有一人，由天主派遣，名叫若翰；他来是为作证，为给光作证，为使众人借他而信。他不是那光：}谁？若翰？谁？若翰洗者。因为若望圣史论他说：\BibleQuote{他不是那光；}主论他说：\BibleQuote{他是燃烧而明亮的灯。}但灯可被点燃，也可熄灭。那么如何？你从哪里分辨？你在探问什么地方？灯所见证的那位，\BibleQuote{是真光。}若望加上了\Quote{真}，你却在寻找谎言。但请再听同一位圣史若望倾流他所饮入的：\BibleQuote{我们见了祂的光荣。}他见了什么？他见了什么光荣？\BibleQuote{正如父独生者的光荣，满溢恩宠和真理。}那么，看，看，我们是否应当约束软弱或轻率的辩论，不推断真理中有任何虚假，将主的应归于主；让我们光荣泉源，好能安全地充满自己。\BibleQuote{天主是真实的，而人人都是虚谎的。}这是什么？天主充满；人人空虚；若他要被充满，就让来到那充满者前吧。\BibleQuote{来就祂，并受光照。}此外，若人是空虚的，由于他是虚谎的，他寻求被充满，且急切地奔向泉源，他渴望被充满，他是空虚的。但你说：\Quote{提防那泉源，那里有谎话。}你还能说什么，除了\Quote{那里有毒。}？\par

7. 他說：\Quote{你已經說了，}\Quote{你已經說了這一切，已經責備了我，已經改正了我。但請告訴我，那位說\InnerQuote{我不上去}卻上了的人，如何沒有說謊？}我會告訴你們，如果我能的話；但請認為這不是小事：如果我沒有使你們在真理中站穩，至少我阻止了你們輕率行事。不過，我還是會告訴你們，我想你們已經知道的事，只要你們還記得我向你們陳述的話。這些話本身就解決了困難。那個慶節持續了好多天。祂說，在這個，就是現在的這個慶節日，今天，就是他們期望祂上去的日子，祂沒有上去；但等到祂自己決定要上去的時候。現在請注意下文：\BibleQuote{祂說了這些話，自己就留在加里肋亞。} 所以祂沒有在那個慶節日上去。因為祂的弟兄希望祂先上去；所以他們說了：\Quote{你從這裡上猶太去。} 他們沒有說：\Quote{我們上去，} 好像要作祂的同伴；也沒有說：\Quote{跟我們上猶太去，} 好像他們要先去；而是好像要打發祂在他們前面去。祂希望他們先去；祂避開了這個陷阱，表現出祂作為人的軟弱，隱藏了天主性；祂避開這事，如同祂逃往埃及一樣。因為這不是無能的後果，而是真理的後果，為要給我們一個謹慎的榜樣；使祂的僕人不會說：\Quote{我不逃跑，因為那是可恥的，} 而逃跑也許是有益的。正如祂將要向門徒說的：\BibleQuote{有人在這城迫害你們，你們就逃到另一城去；} 祂自己給了他們這個榜樣。因為祂是在祂願意的時候被捉拿的；祂是在祂願意的時候誕生的。那時，為了不讓他們預先知道並宣布祂要來，而且陰謀已備好，祂說了：\BibleQuote{我不上這個慶節日去。} 祂說\BibleQuote{我不上去，} 為了隱藏；祂加上\BibleQuote{這個，} 為了不說謊。祂表達了一些，隱藏了一些，抑制了一些；但祂沒有說任何假話，因為\Quote{從祂口中不出虛假。} 最後，祂說了這些話以後，\BibleQuote{祂的弟兄們上去以後；} 福音記載了這一點，請留意，請讀你們向我提出的異議；看這段本身是否不解決困難，看我是否從別處取來了要說的話。主正是等待這個，就是讓他們先上去，免得他們預先報告祂要來，\BibleQuote{祂的弟兄們上去以後，祂也上去了，不是公開地，而是好像秘密地。} 什麼是\Quote{好像秘密地}？祂在那裡好像秘密地行事。什麼是\Quote{好像秘密地}？因為這事實際上也不是秘密。因為祂並沒有真的努力隱藏，祂自己掌握著什麼時候被捉拿的權力。但在那種隱藏中，正如我所說，祂給了祂軟弱的門徒一個榜樣，他們還沒有能力在自己不願意的時候防止被捉拿，要提防敵人的陰謀。因為後來祂甚至公開地上去了，在聖殿裡教導他們；有些人說：\BibleQuote{\InnerQuote{看，這就是祂；看，祂在教導。} 我們的首長肯定說他們想要捉拿祂：\InnerQuote{看，祂公開講話，卻沒有人下手拿祂。}}\par

8. 但如今，如果我們轉向自己，如果我們思想祂的身體，我們就是祂。因為如果我們不是祂，那麼\BibleQuote{凡你們對我這些最小兄弟中的一個所做的，就是對我做的}就不會是真的。如果我們不是祂，那麼\BibleQuote{掃祿，掃祿，你為什麼迫害我？}就不會是真的。所以，我們是祂，因為我們是祂的肢體，我們是祂的身體，祂是我們的頭，整個基督既是頭又是身體。或許祂預見了我們不會遵守猶太人的慶節，這就是\BibleQuote{我不上這個慶節日去。}看，基督和聖史都沒有說謊；如果在這兩者中必須選擇一個，聖史會原諒我，我絕不會把那真實的人放在真理本身之前；我不會喜歡那被派遣的超過派遣他的那一位。但感謝天主，依我之見，那隱晦的事已經被揭示。你們的虔誠將在天主前幫助我。看，我已經盡我所能，解決了關於基督和聖史的疑問。你們這些愛真理的人，請與我一起持守真理，以愛德擁抱，不要爭執。\par

\SourceRecord{Translated by R.G. MacMullen.；Nicene and Post-Nicene Fathers, First Series；Vol. 6.；Edited by Philip Schaff.；Buffalo, NY: Christian Literature Publishing Co.,；1888.；<http://www.newadvent.org/fathers/160383.htm>.}

% structure: p0001 source=h1 target=section reason=page-title

\section{\WorkTitle{新约讲道集 第84篇}}

\Emph{[CXXXIV. Ben.]}\par

\Emph{《论福音的话》，\BibleRef{若 8:31}，\BibleQuote{如果你们存留在我话内，就真是我的门徒了}等等。}\par

1. 亲爱的诸位，你们深知，我们都有一个导师，在他之下都是同门。我们也不做你们的导师，因为我们在高处向你们说话；但他才是众人的导师，他住在我们众人内。他刚才在福音中向我们众人说话，也对我们说，我现在也对你们说；但他也为我们说，既为我们也为你们。\BibleQuote{如果你们存留在我话内}——当然不是现在向你们说话的我的话内，而是刚才在福音中说话的那一位的话内。\BibleQuote{如果你们存留在我话内}，他说，\BibleQuote{你们就真是我的门徒。}要成为门徒，仅仅来到还不够，还要存留。因此，他并没有说：\Quote{如果你们听我的话}；或\Quote{如果你们来到我的话前}；或\Quote{如果你们赞美我的话}；但请注意他所说的：\BibleQuote{如果你们存留在我话内，你们就真是我的门徒，并且你们将认识真理，而真理将使你们自由。}弟兄们，我们该说什么呢？存留在天主的话内是辛劳的还是不辛劳的呢？如果辛劳，看看那巨大的酬报；如果不辛劳，你就白白得到酬报。那么，让我们存留在那存留在我们内的那一位中。我们如果不在他内存留，就会跌倒；但他若不在我们内存留，并不因此失去住所。因为他懂得存留在自己内，他从不离开自己。但至于人，千万不要让他存留在自己内，因为他已迷失了自己。因此，我们因贫乏而存留在他内；他因仁慈而存留于我们内。\par

2. 既然已经说明了我们应当做什么，让我们看看我们要接受什么。因为他指定了一项工作，并应许了酬报。工作是什么？\Quote{如果你们存留在我内。}简短的工作：描述简短，实行起来却伟大。\Quote{如果你们建造在磐石上。}哦，弟兄们，建造在磐石上是多么伟大的事啊！\BibleQuote{洪水来了，风吹了，雨降了，冲击那座房屋，它却没有倒塌，因为基础建在磐石上。}那么，存留在天主的话内，岂不就是不向任何诱惑屈服吗？酬报是什么？\BibleQuote{你们将认识真理，而真理将使你们自由。}请容忍我，因为你们看出我的声音微弱；请用你们平静的注意助我。光荣的酬报！\BibleQuote{你们将认识真理。}或许有人会说：\Quote{认识真理对我有什么益处？}\BibleQuote{真理将使你们自由。}如果真理对你们没有吸引力，就让自由吸引你们吧。在拉丁语的用法中，\Quote{to be free}这个表达有两种含义；主要来说，我们习惯于听到这个词在这个意义上，即凡是自由的人可以被理解为摆脱了某种危险，解除了某种困境。但\Quote{to be free}的正确含义是\Quote{to be made free}；正如\Quote{to be saved}是\Quote{to be made safe}；\Quote{to be healed}是\Quote{to be made whole}；同样\Quote{to be freed}是\Quote{to be made free}。因此我说，\Quote{如果真理对你们没有吸引力，就让自由吸引你们吧。}这在希腊语中表达得更清楚，在那里不可能被理解为其他意义。为了使你们知道它不能被理解为其他意义：当主说话时，犹太人回答说：\Quote{我们从来没有受任何人的奴役；你怎么说真理要使你们自由呢？}那就是，\BibleQuote{真理要使你们自由}，你怎么对我们说，我们从未受过任何人的奴役？他们说：\Quote{你怎么把自由应许给他们，而如你所见，他们从未负过奴役的重轭？}\par

3. 他们听了应当听的，但没有做应当做的。他们听了什么？因为我说：\BibleQuote{真理要使你们自由}，你们就思想自己，没有受人的奴役，便说：\Quote{我们从未受过任何人的奴役。}每一个人，犹太人和希腊人，富人和穷人，有权位的和百姓，皇帝和乞丐，\BibleQuote{凡犯罪的，就是罪的奴仆。}他说：\BibleQuote{凡犯罪的，就是罪的奴仆。}如果人们承认自己的奴役，他们就会看到从哪里可以获得自由。某个生而自由的人被蛮族俘虏，从自由人变成了奴隶；另一个人听到，就怜悯他，考虑自己有钱，就成为他的赎身者，去到蛮族那里，付钱赎回了那人。他确实恢复了自由，如果他去除了不义的话。但曾有人去除过别人的不义吗？那个和蛮族一同为奴的人，被他的赎身者赎回了；赎身者和被赎者之间有很大区别；然而他们可能在罪恶的统治下同为奴隶。我问那被赎者：\Quote{你有罪吗？}他说：\Quote{我有。}我问赎身者：\Quote{你有罪吗？}他说：\Quote{我有。}所以，你既然被赎，不要自夸；你既然是赎身者，也不要抬高自己；但你们俩都逃向真正的拯救者。这还只是一小部分：那些在罪恶之下的，被称为奴仆；他们甚至被称为死人；人害怕被掳带来的事，早已被罪恶带来。为什么？因为他们看似活着，但说\BibleQuote{让死人埋葬他们的死人}的那一位，难道弄错了吗？因此，凡在罪下的都是死的：死的奴仆，在服役中死的，在死亡中为奴的。\par

4. 那么，除了那\BibleQuote{在死者中自由者}\BibleRef{咏 88:6}，谁能从死亡和束缚中解救出来？谁是\BibleQuote{在死者中自由者}\BibleRef{咏 88:6}，除了在罪人中无罪的祂？\BibleQuote{看，世界的首领来了，}\BibleRef{若 14:30}我们的救赎主、我们的解救者亲自说，\BibleQuote{看，世界的首领来了，在我身上他一无所获。}\BibleRef{若 14:30} 他紧抓着那些被他欺骗、诱惑、劝诱犯罪和死亡的人；\BibleQuote{在我身上他一无所获。}\BibleRef{若 14:30}来罢，主，救赎主，来罢，来；让被俘者认出祢，让俘获者逃避祢；作我的解救者。我本已失落，祂找到了我，在祂身上魔鬼找不到任何来自肉身的东西。世界的首领在祂身上找到了肉身，他找到了，但那是怎样的肉身？可朽的肉身，他能抓住、能钉十字架、能杀害的肉身。你错了，欺骗者，救赎主不受欺骗；你错了。你在主身上看到可朽的肉身，那不是罪恶的肉身，而是罪恶的肉身的模样。\BibleQuote{因为天主派遣了自己的儿子，带着罪恶肉身的模样。}\BibleRef{罗 8:3} 真实的肉身，可朽的肉身；但不是罪恶的肉身。\BibleQuote{因为天主派遣了自己的儿子，带着罪恶肉身的模样，以罪恶来惩罚罪恶在肉身中。}\BibleRef{罗 8:3} \BibleQuote{因为天主派遣了自己的儿子，带着罪恶肉身的模样；}在肉身中，但不是罪恶的肉身；而是\BibleQuote{带着罪恶肉身的模样。}为了什么目的？\BibleQuote{以罪恶，}祂身上当然毫无罪恶，\BibleQuote{来惩罚罪恶在肉身中，为使法律所要求的正义，成全在我们身上，我们不随从肉性，而随从圣神。}\BibleRef{罗 8:4}\par

5. 那么，如果那是\BibleQuote{罪恶肉身的模样}，而非罪恶的肉身，怎么是\BibleQuote{以罪恶来惩罚罪恶在肉身中}呢？因为一个模样通常得着它所模之物的名称。\Quote{人}这个词用于一个真实的人；但若你指着一幅画在墙上的画问那是什么，会得到回答：\Quote{一个人。}同样，具有罪恶肉身模样的肉身，为成为赎罪祭，被称为\Quote{罪}。同一宗徒在另一处说：\BibleQuote{祂使那不认识罪的，替我们成了罪。}\BibleRef{格后 5:21} \BibleQuote{那不认识罪的}：谁是不认识罪的，除了那说\BibleQuote{看，世界的首领来了，在我身上他一无所获}\BibleRef{若 14:30}的那位？\BibleQuote{祂使那不认识罪的，替我们成了罪；}\BibleRef{格后 5:21} 即基督自己，祂不认识罪，天主却使祂替我们成了罪。这意味着什么，弟兄们？若说\Quote{祂把罪放在祂身上}或\Quote{祂使祂有罪}，似乎难以忍受；我们怎能容忍所说\Quote{祂使祂成了罪}，基督本人竟成了罪？认识旧约圣经的人会明白我所说的。因为这不是一次使用的表达，而是多次、极其经常地，赎罪祭被称为\Quote{罪}。例如，一只山羊被献为赎罪祭，一只公绵羊，任何东西；为罪所献的牺牲本身被称为\Quote{罪}。因此，赎罪祭被称为\Quote{罪}，以至于法律在一处说：\BibleQuote{司祭们应该把手按在罪上。}\BibleRef{肋 4:15} 那么，\BibleQuote{祂使那不认识罪的，替我们成了罪}\BibleRef{格后 5:21}，即\Quote{祂被做成赎罪祭}。罪被献上，罪被取消。救赎主的血被倾流，债务人的债券被注销。这就是\BibleQuote{为许多人流出，使罪得赦的血。}\BibleRef{玛 26:28}\par

6. 那么，你这曾俘虏我的人，因我的解救者有可朽的肉身而愚蠢地欢腾，这有什么意义？看，祂若有罪；若你在我身上发现了属于你的任何东西，抓住祂罢。\BibleQuote{圣言成了血肉。}\BibleRef{若 1:14} 圣言是创造者，血肉是祂的受造物。这里有什么属于你的，仇敌？而且圣言是天主，祂的人性灵魂是祂的受造物，祂的人性血肉是祂的受造物，天主的可朽血肉是祂的受造物。在这里寻找罪。但你在找什么？真理说：\BibleQuote{这世界的首领要来，在我身上一无所获。}\BibleRef{若 14:30} 因此，祂并非没有血肉，而是没有属于他自己的东西，即没有罪。你欺骗了无辜者，使他们成为有罪。你杀害了无辜者；你消灭了那从你一无所欠的，交出你所紧抓的。那么，你为什么因在基督身上找到可朽的血肉而短暂欢腾？那是你的陷阱：你所欢喜的，正是你被擒之处。你曾因找到什么而得意，如今因失去你所拥有的而悲伤。因此，弟兄们，我们信基督的人，要存留祂的话里。因为我们若存留祂的话里，就真是祂的门徒。因为不仅是那十二位，凡存留祂话里的，都是祂的真门徒。而且\BibleQuote{我们将认识真理，真理必使我们自由}\BibleRef{若 8:32}；就是说，基督、天主子曾说\BibleQuote{我是真理}\BibleRef{若 14:6}，祂将使你们自由，即解救你们，不是从野蛮人手中，而是从魔鬼手中；不是从身体的囚禁，而是从灵魂的邪恶。唯有祂能这样解救。不要让人自称自由，免得仍做奴隶。我们的灵魂不会留在束缚中，因为我们的罪债天天被免除。\par

\SourceRecord{Translated by R.G. MacMullen.；Nicene and Post-Nicene Fathers, First Series；Vol. 6.；Edited by Philip Schaff.；Buffalo, NY: Christian Literature Publishing Co.,；1888.；<http://www.newadvent.org/fathers/160384.htm>.}

% structure: p0001 source=h1 target=section reason=page-title

\section{新约讲道集第85篇}

[第一三五篇。本。]\par

\Emph{关于福音的话，见\BibleRef{若 9:4}和\BibleRef{若 9:31}：\BibleQuote{我们必须做那派我来者的工作}等等。驳斥亚略派。以及关于那生来瞎眼并得见的人所说的：\BibleQuote{我们知道天主不听罪人}。}\par

1. 主耶稣，正如我们在恭读圣福音时所听到的，开了那生来瞎眼之人的眼。弟兄们，如果我们省察我们遗传的惩罚，整个世界都是瞎眼的。因此基督，那光照者，来了，因为魔鬼原是那弄瞎者。他使所有人生来瞎眼，因他诱惑了第一个人。让他们奔向那光照者，让他们奔跑，相信，接受用唾沫和成的泥。圣言好比唾液，肉体是泥土。让他们在史罗亚池中洗脸\BibleRef{若 9:7}。现在福音作者必须向我们解释史罗亚是什么意思，他说：\BibleQuote{这按解释即是被派遣者}。谁是那被派遣者，岂不就是祂在这同一篇训诲中所说的：\BibleQuote{我来做那派我来者的工作}\BibleRef{若 9:4}？看，史罗亚，洗脸，受洗，好使你受光照，使那先前看不见的，如今能看见。\par

2. 看，首先睁开你的眼睛去看那所说的；祂说：\BibleQuote{我来做那派我来者的工作}\BibleRef{若 9:4}。现在立刻有亚略派站出来，说：\Quote{你看，基督所做的不是祂自己的事，而是那派祂来的圣父的事。}他会这样说吗，如果他看见了，就是说，如果他在那被派遣者（如史罗亚）中洗了他的脸？那么你说什么？他说：\Quote{看，祂自己说的。}祂说了什么？\BibleQuote{我来做那派我来者的工作。}那么这些不都是祂自己的吗？不。那么那史罗亚自己所说的，那被派遣者自己所说的，那圣子自己所说的，那位独生子自己所说的，就是你们抱怨说是堕落的？祂说什么？\BibleQuote{凡圣父所有的，都是我的}\BibleRef{若 16:15}。你说祂做了别人的工作，因祂说：\BibleQuote{我必须做那派我来者的工作。}我说圣父有别人的东西：我是按你们的原则说话。为什么你要向我提出基督说\BibleQuote{我来做祂的工作}，如同\BibleQuote{不是我的，而是那派我来的}？\par

3. 我问你，主基督啊，解决困难，终止争论。祂说：\BibleQuote{凡圣父所有的，都是我的}\BibleRef{若 16:15}。那么，如果是你的，难道就不是圣父的吗？因为祂不是说：\Quote{凡圣父所有的，祂都赐给了我}；尽管如果祂这样说，祂也已表明了自己的平等。但困难在于祂说：\BibleQuote{凡圣父所有的，都是我的}。如果你正确理解，凡圣父所有的，都是圣子的；凡圣子所有的，都是圣父的。在另一处听祂说：\BibleQuote{凡我的都是你的，你的都是我的}\BibleRef{若 17:10}。关于圣父和圣子所有的，问题就结束了：他们同心合意地拥有，你不要引入纷争。祂称圣父的工作，就是祂自己的工作；因为\BibleQuote{你的也是我的}，因为祂论到那圣父的工作，祂曾对圣父说：\BibleQuote{凡我的都是你的，你的都是我的}。所以，我的工作是你的，你的工作是我的。\BibleQuote{因为凡圣父所做的，}祂自己说过，主说过，独生子说过，圣子说过，真理说过。祂说了什么？\BibleQuote{凡圣父所做的，圣子也同样做}\BibleRef{若 5:19}。显著的表达！显著的真理！显著的平等。\BibleQuote{凡圣父所做的，圣子也做}。说\Quote{凡圣父所做的，圣子也做}就够了吗？不够；我加上\Quote{同样}。为什么我加上\Quote{同样}？因为那些不理解、眼还未开的走路的人，往往说：\Quote{圣父以命令的方式做，圣子以服从的方式做，因此不是同样。}但如果同样，既是如何的一位，另一位也是如此；所以一位做什么，另一位也做什么。\par

4. \Quote{但是，}他说，\Quote{圣父命令，为的是让圣子执行。}你的想法确实是肉性的，但为了不损害真理，我容许你这样说。看，圣父命令，圣子服从；那么圣子是否因此就不是同一性体，因为一位命令，另一位服从？给我两个人，父亲和儿子；他们是两个人：发出命令的是一个人，服从的是一个人；发出命令的和服从的具有同一个性体。难道发出命令的不是生了与自己性体相同的儿子吗？难道服从的会因服从而失去他的性体吗？现在暂时这样来看，正如你设想两个人那样，圣父发出命令，圣子服从，然而两者都是天主。但是前两者合起来是两个人，后者合起来却是一位天主；这是一个神圣的奇迹。同时，如果你愿意我与你一同承认服从，那么你先与我一同承认性体。圣父生了与祂自身相同的。如果圣父生了与祂自身不同的任何东西，祂就没有生一个真正的儿子。圣父对圣子说：\BibleQuote{在晨星之前，我从母胎中生了祢。}什么是\BibleQuote{在晨星之前}？晨星象征时间。所以是在时间之前，在一切被称为\Quote{之前；}的事物之前，在一切不存在的事物之前，或是在一切存在的事物之前。因为福音并没有说：\Quote{在起初，天主创造了圣言}；如同经上所说：\BibleQuote{在起初，天主创造了天地}；或者说：\Quote{在起初，圣言被生下}；或者说：\Quote{在起初，天主生了圣言}。而是说什么？\Quote{祂是，祂是，祂是。}你听到\Quote{祂是；}就相信吧。\BibleQuote{在起初已有圣言，圣言与天主同在，圣言就是天主。}这么多遍你听到\Quote{是：}不要寻找时间，因为祂永远是\Quote{是。}那么，祂，既永远是，又一直与圣子同在，因为天主可以不藉着你而能生育；祂对圣子说：\BibleQuote{在晨星之前我从母胎中生了祢。}什么是\Quote{从母胎中}？天主有胎吗？难道我们以为天主有身体器官吗？决不可以！为何祂说\BibleQuote{从母胎中，}岂不正是为了让人明白，祂是从祂自己的实体生了祂？所以从母胎中出来的是与生育者自身相同的。如果生育者是另一种东西，而从母胎中出来的是另一种东西，那就是怪物，而不是儿子。\par

5. 因此，让圣子做派遣祂的那一位的工程，圣父也做圣子的工程。\Quote{无论如何，}你说，\Quote{圣父愿意，圣子执行。}看，我证明，圣子愿意，圣父执行。你说：\Quote{你在哪里证明这一点？}我立刻证明。\BibleQuote{父啊，我愿意。}现在如果我要挑剔，看，圣子命令，圣父执行。你愿意什么？\BibleQuote{我在哪里，他们也同我在一起。}我们逃脱了，我们将要在那里，就是祂所在的地方；我们将在那里，我们逃脱了。谁能取消全能的祂的\Quote{我愿意}？你听到祂能力的意愿，现在听祂意愿的能力。\BibleQuote{就如圣父}祂说\BibleQuote{唤起死者并使他们复生；同样，圣子也使祂所愿意的人复生。}\BibleQuote{祂所愿意的人。}不要说，圣子使他们复生，是圣父命令祂复生的。\BibleQuote{祂使自己愿意的人复生。}所以，无论圣父愿意的，还是祂自己愿意的，都是一样的：因为哪里有同一能力，哪里就有同一意愿。让我们现在在一颗不再盲目的人心中持守：圣父和圣子的性体是同一且相同的；因为圣父是真正的父，圣子是真正的子。祂是什么，就生了什么：因为被生者没有堕落。\par

6. 在那位盲人的话中，有一点可能会引起困惑，并可能使许多理解不正确的人绝望。因为那位眼睛被打开的人在他其余的话中说道：\BibleQuote{我们知道天主不听罪人。}我们该怎么办，如果天主不听罪人？我们敢向天主祈祷吗，如果祂不听罪人？给我一个可以祈祷的人；看，这里有一个人听。给我一个可以祈祷的人，仔细筛检整个人类，从不完全到完全。从春天上升到夏天；因为我们刚才咏唱了：\BibleQuote{祢创造了夏天和春天；}也就是说，\Quote{那些已经是属神的人和那些仍然是属血肉的人，祢创造了他们；}因为圣子自己说：\BibleQuote{祢的眼睛看见了我尚未形成的身体。}在我身体中尚未形成的，祢的眼睛看见了。那么怎样？尚未形成的人有希望吗？毫无疑问他们有希望。听接下来的话：\BibleQuote{而在祢的书上，所有人都被写上了。}但也许，弟兄们，属神的人祈祷并被垂听，因为他们不是罪人？那么属血肉的人该怎么办？他们该怎么办？他们要灭亡吗？他们不应该向天主祈祷吗？绝对不可！把福音中那个税吏给我。来，税吏，站出来，展示你的希望，使软弱的人不会失去希望。因为看，税吏同法利塞人上去祈祷，他眼目垂视地面，远远站着，捶着胸膛说：\BibleQuote{主啊，可怜我这个罪人吧！他下去时比那法利塞人更成义。}那位说\BibleQuote{可怜我这个罪人}的人，他说的是真话还是假话？如果他说的是真话，他就是一个罪人；然而他被垂听并被成义了。那么，主开你眼睛的人所说的\BibleQuote{我们知道天主不听罪人}是什么意思？看，天主确实听罪人。但你要洗你低下的脸，让你心中所成就的，如同你脸上所成就的；你就会看到天主听罪人。你心中的幻想欺骗了你。祂还要对你做一些事。我们看到这个人被赶出会堂；\Person{耶稣}听说了，就来找他，对他说：\BibleQuote{你信天主子吗？}他说：\BibleQuote{主，祂是谁，好让我信祂？}他看见了，却没有看见；他用眼睛看见，但心中还没有看见。主对他说：\BibleQuote{你已看见祂，}也就是说，用眼睛；\BibleQuote{且和祢谈话的就是祂。于是他俯伏朝拜了祂。}那时他洗了他心中的脸。\par

7. 你们要热心祈祷，罪人们：忏悔你们的罪过，祈求它们被涂抹，祈求它们减少，祈求当你们增长时，它们减少；但不要绝望，尽管你们是罪人，仍要祈祷。因为谁没有犯罪呢？从司铎开始。司铎们被告知：\Quote{首先为你们自己的罪过献祭，然后再为百姓。}祭献证明司铎们，若有人称自己为义人、没有罪，就可回答他：\Quote{我不看你说了什么，而看你献了什么；你的牺牲本身定了你的罪。你为何为自己的罪献祭，如果你没有罪？你在你的祭献中对天主撒谎吗？}但或许古时百姓的司铎是罪人；新百姓的司铎不是罪人。的确，弟兄们，因为天主这样愿意，我是祂的司铎；我是罪人；我与你们一起捶胸，与你们一起祈求宽恕，与你们一起希望天主会仁慈。但或许圣宗徒们，那些最初的、最高的羊群首领、牧者、牧者的肢体，或许这些人没有罪。是的，确实，他们也有，确实有；他们对此不生气，因为他们承认了。我不敢。首先听主亲自对宗徒们说：\BibleQuote{你们要这样祈祷。}正如那些其他司铎被祭献所定罪，这些也被祈祷所定罪。在祂命令他们祈求的其他事物中，祂也规定了这一点：\BibleQuote{宽免我们的罪债，如同我们也宽免我们的亏欠人。}宗徒们说什么？他们每天祈求罪债得到宽免。他们亏欠着来，被赦免后离开，又成为亏欠者回去祈祷。此生没有罪过，以至于祈祷多少次，罪过就应被宽免多少次。\par

8. 但我该说什么呢？或许当他们学习祈祷时，他们还是软弱的。有人，也许会说这个。当主耶稣教他们那个祈祷时，他们还是婴孩、软弱、属血肉的；他们还不是属神的，属神的人没有罪。那么，弟兄们？当他们变成属神时，他们停止祈祷了吗？那么基督应该说：\BibleQuote{现在你们要这样祈祷}；并给他们另一个属神的祈祷。它是同一个。赐予者是同一位；那么在教会中就使用它祈祷吧。但我们将除去一切争议，当你们说圣宗徒们是属神的，直到主受难时他们是属血肉的；你们必须这样说。而且，事实上，真理是，当主被悬挂时，他们惊慌失措，宗徒们那时绝望了，而强盗相信了。伯多禄胆敢跟随，当主被带去受苦时，他胆敢跟随，他来到院子，在宫廷中疲倦，站在火旁，感到寒冷；他站在火旁，被冰冷的恐惧冻住了。被使女询问，他第一次否认基督；第二次询问，他否认了；第三次询问，他否认了。感谢天主，询问停止了；如果询问没有停止，否认会长时间重复。所以，主复活后，祂坚定了他们，那时他们才成为属神的。那么那时他们没有罪吗？属神的宗徒们写了属神的书信，把它们送到各教会；\Quote{他们没有罪。}你这样说。我不相信你，我问他们自己。告诉我们，圣宗徒们，在主复活后，用从天上派遣的圣神坚定了你们，你们停止有罪了吗？请告诉我们。让我们听听，好让罪人不至于绝望，不至于停止向天主祈祷，因为他们不是无罪。告诉我们。他们中的一个说。是谁呢？就是主最爱的那个，曾靠在主胸怀上的那个，并汲取了天国奥秘、要再次倾泻出来的那个。我问祂：\Quote{你们有罪还是没有？}祂回答说：\BibleQuote{如果我们说我们没有罪，就是欺骗自己，真理不在我们内。}现在是同一位若望，他说：\BibleQuote{在起初已有圣言，圣言与天主同在，圣言就是天主。}你们看见他经过了多高的境界，以致能到达圣言！这样的一位，如此伟大，像鹰一样翱翔在云层之上，在心灵的晴朗宁静中看见了\BibleQuote{在起初已有圣言}；他说：\BibleQuote{如果我们说我们没有罪，就是欺骗自己，真理不在我们内。但如果我们承认我们的罪，祂是忠信正义的，必赦免我们的罪，并洗净我们的一切不义。}因此，你们要祈祷。\par

\SourceRecord{Translated by R.G. MacMullen.；Nicene and Post-Nicene Fathers, First Series；Vol. 6.；Edited by Philip Schaff.；Buffalo, NY: Christian Literature Publishing Co.,；1888.；<http://www.newadvent.org/fathers/160385.htm>.}

% structure: p0001 source=h1 target=section reason=page-title

\section{新约讲道集第86篇}

[CXXXVI. Ben.]\par

\Emph{论同一段福音}，\BibleRef{若 9}\Emph{，关于使生来瞎眼者复明。}\par

我们听到了我们惯常听到的圣经福音经文；但受提醒是好事：从遗忘的怠惰中更新记忆是好的。事实上，这极其古老的经文给我们带来的愉悦，如同新的一样。基督使一个生来瞎眼的人复明；我们为何惊奇？基督是救主；祂以一次慈悲的行为，补偿了在母腹中没有给予的东西。当祂不给那人眼睛时，这当然不是祂的失误；而是为了奇迹而延迟。你也许会说：\Quote{你如何知道这事？}我从祂本人听到的；祂刚才说了；我们一起听见了。因为当祂的门徒问祂说：\BibleQuote{主，是谁犯了罪，这人还是他的父母，竟使他生来就瞎呢？}祂怎样回答，你听见了，我也听见了。\BibleQuote{既不是这人犯了罪，也不是他的父母，而是为叫天主的作为在他身上显扬出来。}看，这就是为什么祂在他没有眼睛时不赐予眼睛的原因。祂没有赐予祂所能赐予的，祂没有赐予祂知道当需要时应赐予的。但你们不要以为，弟兄们，这人的父母没有罪，或者他本人出生时没有沾染原罪；为了赦免这罪，婴儿要领受圣洗以获得罪的赦免。然而，那瞎眼并非因为父母的罪，也非因为他自己的罪；\BibleQuote{而是为叫天主的作为在他身上显扬出来。}因为我们所有人出生时都沾染了原罪；然而我们并非生来瞎眼。但要仔细探究，我们生来就是瞎眼的。瞎眼，指的是心瞎。但主耶稣，因为祂创造了这两者，便治好了两者。\par

2. 你们用信德的眼睛已看见这人原为瞎子，也看见他从瞎眼变为看见；但你们听见他犯了错误。这瞎子错在哪里，我告诉你们：首先，他以为基督是先知，不知道祂是天主子。然后我们听见他一个完全错误的回答；因为他说：\BibleQuote{我们知道天主不听罪人。}如果天主不听罪人，我们还有什么希望？如果天主不听罪人，我们为什么祈祷，并捶胸公布我们的罪过？又那同法利塞人一起上圣殿去的税吏在哪里？法利塞人自夸炫耀自己的功德，而税吏远远站着，眼睛望着地上，捶着胸，承认自己的罪。这承认自己罪的人，从圣殿下去，成义了，而那个法利塞人却没有。那么，天主确实听罪人。但说这些话的人还没有在西罗亚池中洗过心的脸。圣事已运行在他的眼睛上；但心内尚未产生恩宠的祝福。这瞎子何时洗了他心的脸呢？当主被犹太人赶出去后，接纳他进入自己内时。因为主找到他，对他说，正如我们所听见的：\BibleQuote{你信天主子吗？}他回答：\BibleQuote{主，祂是谁，好使我能信祂？}的确，他用眼睛已经看见了；他心里已经看见了吗？不，还没有。等着；他即将看见。耶稣回答他说：\BibleQuote{同你讲话的就是我。}他犹豫了吗？不，他立即洗了脸。因为他正与那西罗亚说话，\BibleQuote{即被派遣者。}谁是被派遣者？不就是基督吗？祂常作证说：\BibleQuote{我承行派遣我来的父的旨意。}那么，祂自己就是西罗亚。那人来时心是盲的，他听了，信了，朝拜了；洗了脸，看见了。\par

3. 然而，那些赶出他的人仍然盲目，因为他们对主吹毛求疵，说祂在安息日用唾沫和泥，并抹瞎子的眼睛。因为当主用一句话治愈时，犹太人公开挑剔。因为祂在安息日没有说话行事，当祂说话时，事就成了。这是明显的挑剔；他们挑剔祂仅仅命令，他们挑剔祂说话；好像他们自己在整个安息日都不说话似的。我可以说，他们不仅在安息日不说话，而且没有一天说话，因为他们远离了对真天主的赞美。尽管如此，弟兄们，正如我所说，这是明显的挑剔。主对一个人说：\BibleQuote{伸出你的手}；那人便痊愈了，他们却因祂在安息日治病而挑剔。祂做了什么？祂做了什么工？祂担了什么重担？但在此例中，吐唾沫在地上、和泥、抹人的眼睛，是做了些工。无人怀疑，这是做了工。主确实违反了安息日，但并不因此有罪。我所说的\Quote{祂违反了安息日}是什么意思？祂，光明已经来到，祂正在除去阴影。因为安息日是由主天主所命定的，由基督亲自命定，祂与圣父同在，当那律法被赐下时；是祂命定的，但作为将来之事的阴影。\BibleQuote{所以，不拘在饮食上，或关于节期、月朔、安息日，都不可让人论断你们，这些原是将来之事的阴影}。现在已经来到，这些事所宣布的那位已经来到。为什么我们喜欢阴影？睁开你们的眼睛，犹太人；太阳就在眼前。\BibleQuote{我们知道}。你们知道什么，心里盲目的人？你们知道什么？\BibleQuote{这个人不是出于天主，因为祂如此违反安息日}。安息日，不幸的人啊，正是这个安息日，基督所设立的，你们却说祂不是出于天主。你们以属肉体的方式遵守安息日，你们没有基督的唾沫。在这个安息日的泥土中，也寻找基督的唾沫，你们就会明白，通过安息日，基督曾被预言。但你们，因为你们没有基督的唾沫在泥土中涂在你们的眼睛上，你们没有来到\Place{西罗亚}，没有洗脸，就继续盲目，对这位盲人的益处盲目，是的，现在不再身体或心灵盲目。他接受了掺有唾沫的泥，他的眼睛被抹了，他来到\Place{西罗亚}，洗了脸，相信了基督，看见了，没有留在那极其可怕的审判中；\BibleQuote{我为审判来到这个世界，叫那看不见的，可以看见；那看得见的，反成盲目}。\par

4. 极其可畏！\BibleQuote{叫那看不见的，可以看见}。好！这是救主的职分，医治能力的宣告，\BibleQuote{叫那看不见的，可以看见}。但是，主啊，祢所加的\BibleQuote{叫那看得见的，反成盲目}是什么意思？如果我们理解，这最真实、最公义。然而，\BibleQuote{那看得见的}是什么？他们就是犹太人。那么他们看见吗？按他们自己的话，他们看见；按真理，他们看不见。那么\BibleQuote{他们看见}是什么？他们以为自己看见，他们相信他们看见。因为他们相信他们确实看见了，当他们维护律法反对基督时。\BibleQuote{我们知道}；因此他们看见。什么是\BibleQuote{我们知道}，不就是我们看见吗？什么是\BibleQuote{这个人不是出于天主，因为祂如此违反安息日}？他们看见；他们读了律法所说的。因为律法规定，凡违反安息日的，要用石头打死。因此他们说祂不是出于天主；但虽然看见，他们对此却是盲目的：祂来到世界是为审判，祂将是生者死者的审判者；祂为什么来？\BibleQuote{叫那看不见的，可以看见}：叫那些承认自己看不见的，得以被光照。\BibleQuote{叫那看得见的，反成盲目}：也就是说，叫那些不承认自己盲目的，更加刚硬。事实上，\BibleQuote{叫那看得见的，反成盲目}已经应验了；律法的维护者，律法的博士，律法的教师，律法的理解者，把律法的创立者钉在十字架上。啊，盲目，这就是\BibleQuote{部分发生在以色列身上}。为使基督被钉十字架，并使外邦人的圆满得以进入，\BibleQuote{盲目部分发生在以色列身上}。什么是\BibleQuote{叫那看不见的，可以看见}？就是为使外邦人的圆满得以进入，\BibleQuote{盲目部分发生在以色列身上}。全世界都躺在盲目中；但祂来了，\BibleQuote{叫那看不见的，可以看见，叫那看得见的，反成盲目}。祂被犹太人弃绝，祂被犹太人钉在十字架上；祂用自己的血为瞎子做了眼药。那些夸口看见光的人，更加刚硬，成了盲目的，钉死了光。何等大的盲目？他们杀了光，但被钉的光照亮了瞎子。\par

5. 请听一位曾盲者如今看见之人。看！那些不愿向医生承认自己瞎眼的人，撞上了何等可怕的十字架！法律与他们同在。没有恩宠，法律有何用？不幸的人们！没有恩宠，法律能做什么？没有基督的唾沫，大地能做什么？没有恩宠的法律，除了使他们更有罪之外，还能做什么？为什么？因为他们是法律的听者而非行者，因此成为罪人、犯法者。天主的人的房东的儿子死了，他的仆人拿着他的棍杖，放在孩子的脸上，但没有复活。没有恩宠的法律有何用？宗徒如今看见了，从前是瞎眼，蒙受光照，说了什么？\BibleQuote{因为如果颁赐了能赋予生命的法律，那么正义就确确实实来自法律了。}注意！让我们回答并说：他说的这是什么？\BibleQuote{如果颁赐了能赋予生命的法律，那么正义就确确实实来自法律了。}如果它不能赋予生命，为什么被颁赐？他继续说并补充：\BibleQuote{但是圣经把所有的人都圈在罪中，为要使因耶稣基督的信德而来的应许赐给相信的人。}这光照的应许，爱的应许，因耶稣基督的信德，赐给相信的人，圣经，即法律，把所有的人都圈在罪中。\BibleQuote{把所有的人都圈在罪中}是什么意思？\BibleQuote{若不是法律说过：不可贪恋，我就不知道私欲偏情。}\BibleQuote{把所有的人都圈在罪中}是什么意思？它也使罪人成为违法者。因为它不能治愈罪人。\BibleQuote{它把所有的人都圈在罪中，}但怀着什么希望？恩宠的希望，仁慈的希望。你领受了法律，你想遵守它，却未能做到；你因骄傲而跌倒，看见了你的软弱。奔向医生，洗洗脸。渴慕基督，宣认基督，信仰基督；圣神加给文字，你便得救。因为如果你从文字中拿走圣神，\BibleQuote{文字叫人死。}如果它叫人死，希望在哪里？\BibleQuote{但是圣神赐予生命。}\par

6. 就让\Person{革哈齐}，\Person{厄里叟}的仆人，拿起棍杖，如同天主的仆人\Person{梅瑟}领受法律一样。让他拿起棍杖，拿起它，奔跑，前行，追上他，把棍杖放在死孩子的脸上。事情就这样发生了；他拿起了棍杖，奔跑，把棍杖放在死孩子的脸上。但有何用呢？棍杖有何用？\BibleQuote{如果颁赐了能赋予生命的法律，}孩子或许可以通过棍杖复活；但是看到\BibleQuote{圣经把所有的人都圈在罪中，}他仍然躺着死了。但为什么它把所有的人都圈在罪中？\BibleQuote{为要使因耶稣基督的信德而来的应许赐给相信的人。}让那通过仆人送出棍杖来证明孩子已死的\Person{厄里叟}来吧；让他自己来，亲自来，亲自进入那妇人的家，上楼到孩子那里，发现他死了，使自己与死孩子的肢体相符，他自己没有死，而是活的。因为他这样做了；他把自己的脸贴在他的脸上，眼睛贴着眼睛，手贴着手，脚贴着脚，他缩小自己，收敛自己，本是伟大的，却使自己成为渺小的。他收敛了自己，可以说，他缩小了自己。\BibleQuote{因为他虽具有天主的形体，却空虚了自己，取了奴仆的形体。}他使自己，活着的与死者相符，这是什么意思？你问这是什么意思吗？请听宗徒的话；\BibleQuote{天主派遣了自己的儿子。}他使自己与死者相符是什么意思？让他告诉我们，让他继续说并再次宣告；\BibleQuote{在罪恶肉身的模样中。}这就是使自己活着的与死者相符；在罪恶肉身的模样中来到我们这里，而不是在罪恶肉身中。人躺在罪恶的肉身中死了，罪恶肉身的模样使自己与他相符。因为他死了，他没有理由死。他死了，独自一人是\BibleQuote{在死者中自由的，}因为全人类的肉身确实是罪恶的肉身。如果那没有罪的那位，使自己与死者相符，在罪恶肉身的模样中来了，你们怎么能复活呢？\Person{主耶稣}啊，你为我们受苦，而不是为你自己，你本无罪，却忍受了罪的惩罚，为了能同时解除罪债和惩罚。\par

\SourceRecord{Translated by R.G. MacMullen.；Nicene and Post-Nicene Fathers, First Series；Vol. 6.；Edited by Philip Schaff.；Buffalo, NY: Christian Literature Publishing Co.,；1888.；<http://www.newadvent.org/fathers/160386.htm>.}

% structure: p0001 source=h1 target=section reason=page-title

\section{\WorkTitle{新约讲道集第87篇}}

\Emph{[CXXXVII. Ben.]}\par

\Emph{\BibleBook{若望}{福音}第十章。论牧人、佣工与盗贼。}\par

1. 诸位亲爱的，你们的信德并非无知；我知道你们借着那位来自天上的导师的教训学会了——你们已将希望寄托于祂——我们的主耶稣基督，祂如今已为我们受苦并复活了，是教会的头，教会是祂的身体；在祂的身体内，肢体的合一与爱德的联系，就是祂健全的健康。但凡在爱德上冷淡的人，便在基督的身体内变得衰弱。但那已经高举我们头颅的主，也能使软弱的肢体痊愈：只要他们不是因过分的邪恶而被切断，而是紧贴身体，直到康复。因为任何尚且附着于身体的，并非没有痊愈的希望；而被切断的，既不能在接受医治的过程中，也不能被治愈。既然祂是教会的头，教会是祂的身体，那么整个基督就是头与身体。祂已经复活了。因此，我们有头在天上。我们的头为我们代祷。我们的头无罪无死，如今为我们的罪向天主求情；好叫我们也在末日复活，改变为天上的荣耀，跟随我们的头。因为头在哪里，其余的肢体也在哪里。但我们尚在此世时，我们是肢体；不要绝望，因为我们必将跟随我们的头。\par

2. 弟兄们，请思考我们这头的爱。祂如今在天上，却仍在此受苦，只要祂的教会在此受苦。在此基督饥饿，在此基督口渴，赤身露体，作客旅，患病，坐监。因为凡祂的身体在此所受的，祂都说这是祂自己所受的；到了末了，祂将这身体置于右边，又将其余被祂践踏的置于左边，祂要对右边的人说：\BibleQuote{来，你们蒙我父降福的，承受自创世以来为你们预备的国度吧。}这是因为什么功德呢？\BibleQuote{因为我饿了，你们给了我吃的；}祂就这样逐一列举其余，仿佛祂亲自领受了；以致他们不明白，回答说：\BibleQuote{主啊，我们几时见你饿了、作客旅、或在监里呢？}祂对他们说：\BibleQuote{你们对我这些最小兄弟中的一个所做的，就是对我做的。}同样，在我们自己的身体里，头在上面，脚在地上；然而在人群拥挤时，若有人踩了你的脚，头岂不说：\Quote{你在踩我}？没有人踩你的头或舌头；头在上面，安全无恙，没有受到伤害；然而因着爱德的联系，从头到脚是合一的，舌头并不与它分离，反而说：\Quote{你在踩我}，虽然没有人碰它。正如那无人碰触的舌头说\Quote{你在踩我}，同样，基督这无人踩踏的头说：\BibleQuote{我饿了，你们给了我吃的。}对那些没有这样做的人，祂说：\BibleQuote{我饿了，你们没有给我吃的。}祂如何结束呢？\BibleQuote{这些人要进入永火，而义人进入永生。}\par

3. 当我们的主那时讲论时，祂说祂是\BibleQuote{牧人}，祂也说祂是\BibleQuote{门}。你们在那段经文中找到二者：\BibleQuote{我是门}和\BibleQuote{我是牧人}。在头内祂是门，在身体内祂是牧人。因为祂对伯多禄说——在伯多禄一人身上祂形成了教会——\BibleQuote{伯多禄，你爱我吗？}他回答说：\BibleQuote{主，我爱你。}\BibleQuote{牧放我的羊。}第三次：\BibleQuote{伯多禄，你爱我吗？}\BibleQuote{伯多禄因祂第三次问他而忧愁}；好像那看见否认者良心的一位，却看不见宣认者的信德。祂一直认识他，甚至在伯多禄不认识自己时也认识他。因为那时他并不认识自己，当他说\BibleQuote{我要与你同死}时，他不知道自己是多么软弱。正如实际上常发生在病人身上的那样：病人不知道身体内的情况，而医生知道；尽管病人正受着疾病的痛苦，医生却没有。医生能比病人自己更好地知道别人身体内的情况。伯多禄当时是病人，主是医生。前者声称他有力量，其实他没有；而主摸着他内心的脉搏，宣称他将三次否认祂。果然应验了，正如医生所预言的，而非如病人所假设的。因此，在祂复活后，主问他，并非不知道他怀着怎样的心来宣认对基督的爱，而是要以三次爱的宣认，抹去三次害怕的否认。\par

4. 因此，主向伯多禄问说：\BibleQuote{伯多禄，你爱我吗？}似乎是说：\Quote{你既然爱我，你要给我什么？你要为我做什么？}伯多禄能为已复活、升天、坐在圣父右边的天主做什么呢？主好像说：\BibleQuote{你若爱我，这就是你给我的，这就是你为我做的：喂养我的羊群；由门进入，不要从别处上去。}你们听福音诵读时听到：\BibleQuote{那由门进去的，才是牧人；但从别处爬上去的，便是贼，是强盗；他企图驱散、杀害、抢夺。}谁是那由门进入的？那藉着基督进入的。谁是这样的人？就是效法基督受苦、承认基督谦卑的人；因为天主为了我们而成为人，好让人承认自己不是神，而是人。因为谁若愿意在人前显得是神，而自己却是人，他就不效法那本来是天主却成为人的那一位。但对你说的并不是要你变得比你所是更小，而是要你承认你所是的。承认自己的软弱，承认自己是人，承认自己是罪人；承认那使人成义的是祂，承认你满身污点。让你的污点在你的告解中显露出来，这样你就属于基督的羊群。因为告明罪过邀请医生的医治；正如患病时，说\Quote{我健康}的人不寻找医生。法利塞人和税吏不是上到圣殿去了吗？一个夸耀自己的健康，另一个向医生显示自己的伤口。因为法利塞人说：\BibleQuote{天主，我感谢你，因为我不像这个税吏。}他因别人而自夸。如果那个税吏是健康的，法利塞人就会嫉妒他，因为他就没有人可以自夸了。那么怀着这种嫉妒之心而来的人，他是什么状态呢？他当然是不健康的；而当他称自己健康时，他就没有得着医治就下去了。但另一个人眼目低垂，不敢举目望天，捶着胸膛说：\BibleQuote{天主，可怜我这个罪人吧！}主怎么说呢？\BibleQuote{我实在告诉你们：那个税吏下去，成了义的，而不是那个法利塞人。因为凡高举自己的，必被贬抑；凡谦卑自己的，必被高举。}所以，那些高举自己的人，是想从别处爬上羊栈；但那些谦卑自己的人，是从门进入羊栈的。因此，主论到前者时说\Quote{他进入}，论到后者时说\Quote{他上去}。你们看，那上去的人，就是寻求高举的人，他并不进入，反而跌倒。但那谦卑自己、为要从门进入的人，并不跌倒，而是牧人。\par

5. 但主提到了三种人物，我们有责任在福音中查考它们：就是牧人、佣工和贼。我想你们在读经时注意到，祂指出来了牧人、佣工和贼。祂说：\BibleQuote{善牧为羊舍掉自己的性命}，并且由门进入。他说：\BibleQuote{贼和强盗由别处爬上去；佣工若看见狼或贼来，就逃跑了，因为他不关心羊；因为他是佣工，不是牧人。}前者由门进入，因为他是牧人；后者由别处上去，因为他是贼；第三者看见那些想要抢夺羊群的人，就害怕逃跑，因为他是佣工，不关心羊；因为他是佣工。如果我们找到了这三种人物，圣善的弟兄们，你们就找到了你们应当爱的、你们应当容忍的、以及你们应当提防的人。牧人应当被爱，佣工应当被容忍，强盗必须提防。在教会中有些人，如宗徒所说的，他们出于私心传福音，寻求人的好处，无论是金钱、荣誉还是人的称赞。他们传福音，希望以任何方式得到报酬，并不那么关心他们所宣讲之人的得救，而是关心自己的利益。但那个从没有得救的人那里听到得救之道的人，如果相信他所宣讲的那一位，而不寄望于向他宣讲得救的那人；那么，宣讲者将受亏损，而受宣讲者将获益。\par

6. 你们听主论及法利塞人说：\newline{}\newline{}\Quote{他们坐在梅瑟的座位上。}\newline{}\newline{}主并非只指他们；好像祂要那些相信基督的人到犹太人的学校去，在那里学习通往天国之路。难道主降来不是为此目的：为建立教会，将那些有良好信德、良好望德和良好爱德的犹太人如同麦子从糠秕中分别出来，使他们成为受割礼的一堵墙，另一堵墙则来自未受割礼的外邦人，这两堵来自不同方向的墙，祂自己成为角石？难道同一位主不是论及这两个将要合一的人民说：\newline{}\newline{}\Quote{我还有别的羊，不属于这一羊栈}？\newline{}\newline{}当时祂是对犹太人说的；\newline{}\newline{}\Quote{我该引领他们，}祂说，\newline{}\newline{}\Quote{使成为一牧一栈。}\newline{}\newline{}因此，祂从其中召叫门徒的两只船，就预表了这两个人民，当他们撒下网，捕获了如此之多、如此之大的鱼，以致网险些破裂。\newline{}\newline{}\Quote{他们便装满，}经上记着，\newline{}\newline{}\Quote{了两只船。}\newline{}\newline{}两只船预表了同一个教会，却由两个人民组成，在基督内结合，虽来自不同地方。这同样也是雅各伯的两位妻子肋阿和辣黑耳的预表。这两者，也有两位盲人的预表，他们坐在路旁，主使他们复明。如果你们留意圣经，会在许多地方发现这两个教会——并非两个而是一个——的预表。因为角石就是为了使二者合一。那位牧人就是为了使两群羊合一。那么，主既教导教会，并拥有超越犹太人的自己的学校，如我们目前所见，祂会打发信祂的人到犹太人那里去学习吗？祂却以经师和法利塞人的名字暗示在祂的教会中将有些只说不做的人；但以梅瑟的身份，祂指定了自己。因为梅瑟预表祂，正因如此，当他对百姓说话时，他脸上蒙了一块帕子；因为当他们遵行法律，沉溺于肉身的欢乐和享乐，盼望一个地上的国时，一块帕子就蒙在他们脸上，使他们看不到圣经中的基督。因为当主受难以后，帕子被除去，圣殿的奥秘就显明了。因此，当祂悬在十字架上时，圣殿的帐幔从上到下裂为两半；而保禄宗徒明确地说：\newline{}\newline{}\Quote{但当你们转向基督时，帕子就会除去。}\newline{}\newline{}对于那些不转向基督的人，尽管他们阅读梅瑟法律，帕子仍蒙在他们的心上，如宗徒所说。那么当主预先指明在祂的教会中将有这样的人时，祂说了什么？\newline{}\newline{}\Quote{经师和法利塞人坐在梅瑟的座位上。凡他们吩咐你们的，就照着做；但不要照他们的行为去做。}\par

7. 当邪恶的圣职人员听到这反对他们的话时，他们想要曲解它。因为我听说有些人的确想要曲解这句话。他们若能做到，岂不要从福音中把它抹去？但因他们不能抹去，便设法曲解它。但主的恩宠和仁慈临在，不容他们如此行；因为祂用祂的真理护卫了祂所有的宣言，并如此平衡了它们：若有人想从中割去什么，或用错误的读法或解释引入什么，任何正直的人都能将圣经中割去的那部分与圣经接合，读一读上文或下文，就会发现那人想要错误解释的意思。那么，你们以为，关于那被记载\newline{}\newline{}\Quote{照他们所说的去做}\newline{}\newline{}的人，他们说什么？那是（事实上正是如此）对平信徒说的。因为希望善度一生的平信徒，当他注意到一个邪恶的圣职人员时，他对自己说什么？主说：\newline{}\newline{}\Quote{照他们所说的去做，但不要照他们所做的去做。}\newline{}\newline{}让我行走在主道上，不追随此人的生活。让我从他那里听到的不是他的话，而是天主的话。我将跟随天主，让他跟随自己的情欲。因为若我想在天主面前这样为自己辩护说：\newline{}\newline{}\Quote{主，我见祢的圣职人员生活邪恶，因此我也生活邪恶，}\newline{}\newline{}祂岂不对我说：\newline{}\newline{}\Quote{你这恶仆，你没有听我说过：\InnerQuote{照他们所说的去做，但不要照他们所做的去做}？}\newline{}\newline{}但一个邪恶的平信徒，一个不信者，不属于基督的羊群，不属于基督的麦子，只是像糠秕一样在禾场上被容忍，当天主的话开始责备他时，他对自己说什么？\newline{}\newline{}\Quote{去！你为什么对我说？那些主教和圣职人员自己都不做，你强迫我做吗？}\newline{}\newline{}这样他为自己找的不是他坏案件的保护人，而是刑罚的同伴。因为那一位他选择模仿的恶人，无论他是谁，会在审判之日为他辩护吗？因为正如魔鬼诱惑所有的人，他不是引诱他们分享王国，而是分享他的永罚；同样，所有追随恶人的人，为自己寻找的不是通往天国的保护，而是下地狱的同伴。\par

8. 那么，当他们在邪恶的生活中听到这句话时，他们是如何曲解这宣告的呢？\Quote{主说得好：\BibleQuote{你们要照他们所说的去做，但不要照他们所行的去做}\BibleRef{玛 23:3}？}\Quote{说得好，}他们说。\Quote{因为这话是对你们说的：你们要照我们所说的去做，但不要照我们所行的去做。因为我们献祭，你们却不可。}看这些人的狡猾诡诈；我该称他们为什么？傭工。因为倘若他们是牧人，就不会说这样的话。因此，主为要封住他们的口，继续说：\BibleQuote{他们坐在梅瑟的座位上；凡他们所说的，你们要遵行；但不要照他们所行的去做；因为他们只说而不做。}\BibleRef{玛 23:2-3}那么，弟兄们，这是什么意思？如果祂是论及献祭，祂会说\Quote{因为他们只说而不做}吗？他们确实献祭，他们确实向天主献祭。那么，什么事他们只说而不做呢？请听下文：\BibleQuote{他们把沉重而难以负荷的重担捆起来，压在人的肩上，自己却连一个指头也不肯动。}\BibleRef{玛 23:4}祂如此公开地责备、描述并指出他们。但那些人想要曲解这段经文时，就明显表明他们在教会中只求自己的利益；并且他们没有读过福音；因为如果他们知道这一页，并读完全文，就绝不敢说这话。\par

9. 但请留意一个更清楚的证据，证明教会中有这样的人。免得有人对我们说：\Quote{祂完全是在说法利塞人，祂是说经师，祂是说犹太人；因为教会中并没有这样的人。}那么，主说：\BibleQuote{不是凡向我说\InnerQuote{主啊，主啊}的人，就能进入天国}\BibleRef{玛 7:21}，是指谁呢？祂又加上：\BibleQuote{当那日，许多人要对我说：\InnerQuote{主啊，主啊，我们不是因祢的名说过预言，因祢的名行过许多异能，因祢的名吃过喝过吗？}}\BibleRef{玛 7:22}什么！难道犹太人在基督的名里行这些事吗？显然，祂是指那些拥有基督之名的人。但下文呢？\BibleQuote{那时我要向他们宣告：我从来不认识你们，你们这些作恶的人，离开我去吧！}\BibleRef{玛 7:23}请听宗徒为这样的人叹息。他说有些人\BibleQuote{出于爱德}\BibleRef{斐 1:17}传福音，另一些人\BibleQuote{出于私心}\BibleRef{斐 1:17}；关于后者他说：\BibleQuote{他们传福音不正。}\BibleRef{斐 1:17}他们传的是正事，但他们自己不正。他们传的是对的，但传的人不对。为什么他不对？因为他在教会中寻求别的东西，不寻求天主。如果他寻求天主，他便是贞洁的；因为灵魂在天主内拥有她合法的丈夫。凡从天主寻求超过天主的东西，就不是贞洁地寻求天主。想一想，弟兄们：如果一个妻子爱她的丈夫是因为他富有，她就不贞洁。因为她爱的不是丈夫，而是丈夫的金子。但如果她爱丈夫，她就爱赤身露体和贫穷中的他。因为如果她爱他是因为他富有；那么，万一（人的际遇无常）他被剥夺法律保护，突然陷入困窘呢？也许她会放弃他；因为她所爱的不是丈夫，而是他的财产。但如果她确实爱丈夫，她在贫穷中会更爱他；因为她也以怜悯来爱。\par

10. 然而，弟兄们，我们的天主永不贫穷。祂是富有的，祂造了万物：天、地、海和天使。在天上，无论我们看见什么，无论我们看不见什么，都是祂造的。但尽管如此，我们不应爱这些财富，而应爱造它们的主。因为祂除了祂自己以外，没有应许你任何东西。若你找到更宝贵的事物，祂就会赐给你。大地是美丽的，天、天使也是美丽的；但造它们的主更美丽。那么，那些出于爱天主而宣讲天主的人，为天主而宣讲天主的人，就是牧养羊群，不是傭工。我们的主耶稣基督要求灵魂有这样的贞洁，当祂对伯多禄说：\BibleQuote{伯多禄，你爱我吗？}\BibleRef{若 21:15}\Quote{你爱我吗}是什么意思？你贞洁吗？你的心不是淫乱的吗？你在教会中寻求的不是你自己的东西，而是我的吗？如果你是这样的人，并且爱我，\BibleQuote{你牧放我的羊群。}\BibleRef{若 21:16}因为你不会是傭工，而是牧人。\par

11. 但那些宗徒为之叹息的人，并没有贞洁地宣讲。可是他说什么呢？\BibleQuote{那又怎样？无论如何，或是出于私心，或是出于真诚，基督总被传开了。}\BibleRef{斐 1:18}他容忍那里有傭工。牧人真诚地宣讲基督，傭工出于私心宣讲基督，寻求别的东西。然而，两者都宣讲基督。请听牧人保禄的声音：\BibleQuote{或是出于私心，或是出于真诚，基督总被传开了。}\BibleRef{斐 1:18}他本人是牧人，却甘愿有傭工。因为他们在他们所能之处行事，在他们所能之处有益。但当宗徒为了其他用途，寻找那些软弱者可效法其行为的人时，他说：\BibleQuote{我派了弟茂德到你们那里去，他必提醒你们我所行的道路。}\BibleRef{斐 2:19}他又说什么？\BibleQuote{我派了一位牧人到你们那里去，提醒你们我所行的道路；}\BibleRef{斐 2:19}也就是说，他自己也像我一样行事。在派这位牧人时，他说了什么？\BibleQuote{因为我没有别人像他那样与我同心，真诚地关心你们。}\BibleRef{斐 2:20}难道与他在一起的人不多吗？但下文是什么？\BibleQuote{因为众人都寻求自己的事，不寻求基督耶稣的事；}\BibleRef{斐 2:21}也就是说，\Quote{我曾想派一位牧人到你们那里去；因为有许多傭工；但派一个傭工是不合适的。}傭工被派去办理其他事务和生意；但对于保禄当时所渴望的，牧人是必要的。他在众多傭工中几乎找不到一个牧人；因为牧人少，傭工多。但论到傭工，有何话说？\BibleQuote{我实在告诉你们，他们已经获得了他们的赏报。}\BibleRef{玛 6:2}论到牧人，宗徒说了什么？\BibleQuote{所以，谁若洁净自己脱离这些事，就成为贵重的器皿，被圣化，为主所用，随时准备行各种善事。}\BibleRef{弟后 2:21}关于牧人，我已说了这么多。\par

12. 但现在我们要谈谈那些佣工。主说：\BibleQuote{佣工看见狼埋伏在羊群附近，就逃跑了。}为什么呢？\BibleQuote{因为他不顾念羊。}所以只要佣工还没看见狼来，还没看见贼和强盗，他是有用的；但当他看见他们，他就逃跑。那些佣工中，有谁在看见狼和强盗时，不从教会逃跑呢？狼和强盗很多。他们就是从别处上去的人。那些上去的人是谁？那些想藉多纳徒之道残害基督羊群的人，他们就是从别处上去的。他们不藉着基督进入，因为他们不谦卑。因为他们骄傲，所以他们上去。什么是\Quote{他们上去}？他们被高举。他们从何处上去？从别处：他们愿意按自己的道路被命名。那些不在合一中的人，是属另一条道路的，并藉此道路上去，即被高举，并想抢夺羊群。现在看他们如何上去。他们说：\Quote{是我们圣化，是我们使人成义，是我们使人成为义人。}看他们升到了何处。但\BibleQuote{凡高举自己的，必被贬抑。}我们的主天主能贬抑他们。狼就是魔鬼，他埋伏着欺骗，跟随他的人也是如此；因为经上说：\BibleQuote{他们外面披着羊皮，里面却是凶残的狼。}如果佣工观察到有人沉溺于邪恶言论、或有害灵魂的见解、或做任何可憎不洁的事，而且此人似乎在教会中拥有某种重要地位（佣工若指望从中得利，他就是佣工），却闭口不言，当他看见那人因罪而丧亡，看见狼跟随他，看见他的喉咙被狼的牙齿拖向惩罚，却不向他说：\Quote{你犯罪了}；不责备他，唯恐失去自己的利益。这就是：\Quote{他看见狼，就逃跑}；他不向那人说：\Quote{你在作恶。}这不是身体的逃跑，而是灵魂的逃跑。你看见身体站着不动的人，心里却在逃跑，当他看见罪人而不对他说\Quote{你犯罪了}时，甚至与他同流合污。\par

13. 我的弟兄们，难道有司铎或主教上到这高处，说些别的，而不是说不可掠夺他人财产、不可欺诈、不可行恶吗？凡坐在梅瑟座位上的人，不能说别的；正是那座位藉着他们说话，而非他们自己。那么，所谓\BibleQuote{人岂能从荆棘上摘葡萄，从蒺藜里摘无花果呢？}和\BibleQuote{每棵树凭它的果子就能被认出}是什么意思？一个法利塞人能说出好话吗？法利塞人是荆棘；我怎能从荆棘上摘葡萄呢？因为主，祢曾说：\BibleQuote{凡他们所吩咐你们的，你们要遵行；但不要效法他们的行为。}祢既说\BibleQuote{人岂能从荆棘上摘葡萄呢？}，怎能吩咐我从荆棘上摘葡萄呢？主回答你说：\Quote{我没有吩咐你从荆棘上摘葡萄；但你且看，要留意，正如常有的情形，当葡萄藤蔓匍匐在地时，会不会缠在荆棘中。}因为有时我们会发现，我的弟兄们，一株栽种在莎草旁的葡萄，那里有荆棘篱笆，它伸出枝条，缠绕在荆棘篱笆中，葡萄就悬在荆棘之间；看见的人摘下葡萄，却不是从荆棘上摘，而是从那缠在荆棘中的葡萄藤上摘。同样，法利塞人是多刺的；但因坐在梅瑟座位上，葡萄藤缠绕他们，葡萄——即好话、好教训——从他们身上垂下来。你若摘葡萄，荆棘不会刺伤你，当你读到\BibleQuote{凡他们所吩咐你们的，你们要遵行；但不要效法他们的行为}时。但如果你效法他们的行为，荆棘就会刺伤你。所以，为使你摘到葡萄而不被荆棘缠住，就要\BibleQuote{凡他们所吩咐你们的，你们要遵行；但不要效法他们的行为。}他们的行为是荆棘，他们的话是葡萄，但来自那葡萄藤，即来自梅瑟的座位。\par

14. 因此，这些人看见狼和强盗时就逃跑。我刚才已经开始说，从这高处的座位上，他们只能说：\Quote{要做好事}、\Quote{不可发虚誓}、\Quote{不可欺骗}、\Quote{不可诈骗任何人}。但有时人的行为如此败坏，以致有人向主教请教如何夺取他人产业，并且就是从主教那里寻求这样的意见。我们自己也遇到过这种情况，我们凭经验说话：因为我们本不该相信。许多人向我们要求邪恶的意见，要求说谎和欺诈的意见，以为这样能取悦我们。但靠着基督的名，如果我们所说的话蒙主喜悦，还没有这样的人来试探我们，并在我们身上找到他所希望的。因为由于那召叫我们的主的恩宠，我们是牧者，而非佣工。但正如宗徒所说：\BibleQuote{对我来说，受你们审断，或受人间法庭审断，都是极小的事；连我自己也不审断自己。因为我虽然自觉良心无愧，却不能因此就成义；那审断我的，是主。}因此，我的良心是好的，并非因为你称赞它。因为你怎能称赞你所看不见的呢？让那看见的来称赞吧；是的，让他来纠正，如果他看见那里有冒犯他目光的事。因为连我自己也不说我是完全健康的；但我捶着胸膛，向天主说：\Quote{求你宽恕，使我不犯罪。}然而，我确实认为——我是在他面前说话——我从你们那里寻求的，除了你们的救恩，别无他物；我不断为我弟兄们的罪而叹息，感到痛苦，内心受折磨，并且常常责备他们；是的，我从不停止责备他们。所有记得我所说的话的人都是证人，我多少次责备了犯罪的弟兄们，并且是严厉地责备。\par

15. 现在，圣善的弟兄们，我正与你们谈论我的劝告。因基督之名，你们是天主的子民，是公教的子民，是基督的肢体；你们没有分裂于合一。你们与宗徒的成员共融，与遍布全世界的圣殉道者的纪念共融，并且你们归属我的牧养，使我能为你们交出一份好账目。如今我全部的账目，你们知道是什么。\Quote{主，祢知道我讲了话，祢知道我没有缄默，祢知道我以何种精神讲话，祢知道当我在祢面前讲话时，我流泪了，却没有被听从。} 我想这就是我全部的账目。因为圣神借着先知\Person{厄则克耳}给了我确切的希望。你们知道关于守望者的这段经文：\BibleQuote{人子啊！} 祂说：\BibleQuote{我立你作以色列家的守望者；如果我对恶人说：恶人哪，你必要死，你却不讲话；} 也就是说（我向你们说话，为使你们讲话），\BibleQuote{如果你不宣告，刀剑}，也就是说，我所警告罪人的灾祸，\BibleQuote{来临并把他夺去；那个恶人必死在自己的罪恶中，但我要向守望者索讨他的血。} 为什么？因为他不讲话。\BibleQuote{但是如果守望者看见刀剑来临，就吹号角，} 使他可以逃跑，而他却没有谨慎自守，即没有改过自新，以免他处于天主所警告的惩罚之中，\BibleQuote{刀剑必来临并夺去任何一人；那个恶人必死在自己的罪恶中；但是你，} 祂说，\BibleQuote{已拯救了你的灵魂。} 在福音的那处，祂又对仆人说了什么？当他说：\BibleQuote{主，我知道祢是} 严厉或\BibleQuote{苛刻的人，因为祢收割了没有播种的地方，聚集了没有散播的地方；于是我害怕了，就去把祢的元宝埋在地里，看，祢收回了属于祢的。} 祂却说：\BibleQuote{\InnerQuote{你这个可恶而懒惰的仆人}，因为你已知道我是严厉而苛刻的人，没有播种却去收割，没有散播却去聚集；我的贪求本该更教导你，我期待从我的钱财获得盈利。\InnerQuote{你因此应当把我的钱交给兑换商，等我回来时我就能连本带利收回我的钱。}} 难道祂说：\Quote{你应当给，并索要}？所以，弟兄们，是我们给，祂要来索要。请祈祷，使祂发现我们准备好了。\par

\SourceRecord{Translated by R.G. MacMullen.；Nicene and Post-Nicene Fathers, First Series；Vol. 6.；Edited by Philip Schaff.；Buffalo, NY: Christian Literature Publishing Co.,；1888.；<http://www.newadvent.org/fathers/160387.htm>.}

% structure: p0001 source=h1 target=section reason=page-title

\section{论新约·讲道八十八}

\Emph{[CXXXVIII. Ben.]}\par

\Emph{论福音之言}，\BibleRef{若 10:14}\Emph{，\Quote{我是善牧}等。驳\OriginalTerm{多纳徒派}{Donatists}。}\par

1. 我们听到主耶稣向我们阐述善牧的职责。祂确已让我们知道，正如我们能理解的，有善牧存在。但为使众多牧人这一说法不被误解，祂说：\Quote{我是善牧。}祂如何是善牧，在接下来的话中表明：\Quote{善牧}，祂说，\Quote{为羊舍命。但雇工不是牧人，看见狼来，就撇下羊逃走，因为他是雇工，不关心羊。}基督，乃是善牧。那伯多禄是什么？他不也是一位善牧吗？他不也为羊舍命了吗？保禄是什么？其他宗徒呢？紧随其后的真福主教和殉道者呢？再如我们神圣的西彼廉？他们不都是善牧，而非雇工吗？关于雇工，经上说：\Quote{我实在告诉你们，他们已经获得了他们的赏报。}所有这些人都曾是善牧，不单因为他们流了血，而是因为他们为羊流了血。因为他们流血不是出于骄傲，而是出于爱德。\par

2. 因为即使在异端者中，那些为自己的不义和错误而遭受任何磨难的人，也以殉道的名义自夸，好借这美丽的伪装更加轻易地掠夺，因为他们是狼。若你欲知该将他们列入何等行列，请听善牧宗徒保禄的话：并非所有即使在苦难中将自己的身体交付火烧的人，都被算作为羊流血，反而是反对羊。\Quote{若，}他说，\Quote{我能说人间的语言，并说天使的语言，却没有爱德，我就成了鸣的锣、响的钹。若我能知一切奥秘，并有一切先知话和一切信德，甚至能移山，却没有爱德，我什么也不算。}移山的信德实在是一件大事。这些确实都是大事；但若我拥有它们却没有爱德，他说，不是它们，而是我一无所有。但至此，他尚未触及那些以虚假的殉道之名在苦难中自夸的人。且听他如何触及，甚至彻底刺透他们：\Quote{若我将我所有财产全施舍给穷人，并交出自己的身体被火烧。}他们就在这里。但请注意下文：\Quote{却没有爱德，于我毫无益处。}看，他们已到了受苦的地步，甚至到流血，甚至到火烧身体；然而这一切于他们毫无益处，因为缺少爱德。加上爱德，这一切都有益；除去爱德，其余一切毫无益处。\par

3. 这爱德是何等的美善，诸位弟兄！有什么更珍贵？有什么更能赐予光明？或力量？或益处？或安全？天主的恩赐很多，连恶人也拥有，他们将会说：\Quote{主，我们因祢的名字讲预言，因祢的名字驱魔，因祢的名字行了许多异能。}祂不会回答：\Quote{你们没有行这些事。}因为在如此伟大的审判者面前，他们不敢说谎，或夸耀没有行过的事。但因他们没有爱德，祂向他们全体回答：\Quote{我不认识你们。}一个即使被证实却仍不爱合一的人，怎能拥有哪怕最小的爱德呢？主正是为了给善牧们强调这个合一，才不愿提及众多的牧人。因为，正如我已经说过的，并非伯多禄不是善牧，保禄、其他宗徒、后来的圣主教们以及真福西彼廉不是善牧。所有这些人都曾是善牧；然而，主向善牧们所推荐的，不是善牧，而是一位善牧。\Quote{我}，祂说，\Quote{是善牧。}\par

4. 让我们以我们仅有的一点理解来质问主，并以最卑微的言谈与如此伟大的导师交谈。祢说什么，主，善牧？因为祢是善牧，也是善羔羊；既是牧人又是牧场，既是羔羊又是狮子。祢说什么？让我们倾听，并求祢帮助我们，使我们能明白。\Quote{我}，祂说，\Quote{是善牧。}那么伯多禄呢？他难道不是牧人，或是恶牧？让我们看看，他是否不是牧人。\Quote{你爱我吗？}祢对他说，主，\Quote{你爱我吗？}他回答：\Quote{我爱祢。}祢对他说：\Quote{喂养我的羊。}祢，主，藉祢自己的询问，藉祢自己话语的强力确证，使这爱者成为牧人。那么，祢将祢的羊托付给他牧放，他就是牧人。现在让我们看看他是否不是善牧。从这询问和他的回答，我们便可得知。祢问他是否爱祢；他回答：\Quote{我爱祢。}祢看见他的心，他回答的是真话。那么，爱如此伟大的善的他，难道不是善的吗？那回答是从他内心深处发出的。为何这伯多禄——他有祢的眼睛在他心里作证——因为祢不只一次，而是第二次和第三次问他而忧伤，好叫他以三次爱的宣认抹去三次否认的罪？为何，我说，因那知道自己在问什么并赐予他所听见的那位反复询问而忧伤的他，却回答：\Quote{主，祢知道一切，祢知道我爱祢。}什么？在做这样的宣认，更确切地说，这样的表白时，他会说谎吗？他真实地向祢表白了他的爱，从内心深处发出了爱的话语。如今祢说过：\Quote{善人从心里善的宝库发出善来。}那么，他既是牧人，又是善牧；当然比不上牧者之牧者的权能与良善，但他仍然是牧人，而且是善牧；所有其他这样的人也都是善牧。\par

5. 那么，你对善牧提出唯一牧人，这是何意？岂不是在唯一牧人中教导合一？主自己藉我的服务更清楚地解释了这事，藉这福音提醒你们，亲爱的，并说：\Quote{听我所提出的；我说过：\InnerQuote{我是善牧。}因为其余的一切善牧都是我的肢体。}一个头，一个身体，一个基督。所以，既是牧人之牧人，又是牧人的牧人，还有羊群与他们的牧人同在那牧人之下。这一切是什么，岂非宗徒所说的？\Quote{就如身体是一个，却有许多肢体，并且身体上所有的肢体虽多，仍是一个身体；基督也是这样。}因此，如果基督是这样，那么基督在自己内包含了一切善牧，提出唯一一位，就是有理由的，他说：\Quote{\InnerQuote{我是善牧。}\InnerQuote{我是，}我独自是，其余与我同在一体内。谁不与我牧养，就是反对我。\InnerQuote{凡不与我聚集的，就是分散的。}}那么听这合一更强烈地提出：\BibleQuote{别的羊，}他说，\BibleQuote{我还有，不属于这羊栈。}因为他是指着按血统属于肉身的以色列的第一个羊栈说的。但还有别的属于这以色列信德血统的人，他们还在外面，在外邦人中，已预定，却尚未聚集。他认识他们，因为他预定了他们；他知道他们，因为他曾来以自己的血救赎他们。他看见他们，他们那时尚未看见他；他认识他们，他们尚未信他。\BibleQuote{别的羊，}他说，\BibleQuote{我还有，不属于这羊栈；}因为按以色列的血统她们不属于这羊栈。但无论如何，他们不会在这羊栈之外，\BibleQuote{我也必须领他们来，好使只有一个羊栈，一个牧人。}\par

6. 那么，在这牧人之牧人身上，他的爱侣、他的新娘、他的美丽者——原是由他而美丽，先前因罪而丑陋，后来因宽恕和恩宠而漂亮——因爱慕他而热切地向他说话，并对他说：\Quote{你喂养在哪里？}请看，这神圣的爱情在这里被怎样的狂喜所激动。那些尝过这爱情一丝甘甜的人，更为这狂喜所愉悦。凡爱基督的人，正当地听见这话。因为教会是在他们内、并藉着他们，在\WorkTitle{雅歌}中歌唱：爱基督的人，视他如同没有俊美，却仍是唯一的美丽者。经上说：\Quote{我们看见了他，} \Quote{他没有俊美，也没有华丽。}他在十字架上就是如此，当他戴着茨冠时，他显示出自身，满脸创伤，毫无俊美，好像失去了能力，好像不是天主之子。在盲人眼中，他就是如此。因为依撒意亚是在犹太人的立场上说这话：\Quote{我们看见了他，他没有任何俊美和华丽。}当时有人说：\Quote{你如果是天主子，就从十字架上下来！他救了别人，却救不了自己。}他们用芦苇敲打他的头，说：\Quote{基督，向我们预言，打你的是谁？}因为\Quote{他没有俊美也没有华丽。}你犹太人就是这样看见他的。因为\Quote{部分的愚钝发生在以色列，直到外邦人的总数满了}，直到别的羊来到。因为当时愚钝发生了，所以你们看见那充满俊美的却没有俊美。\Quote{如果你们认识他，你们决不会把光荣的主钉在十字架上。}但你们做了，因为你们不认识他。然而他，虽说在你们面前好像没有俊美，但他自身充满美丽，却为你们祈祷：\Quote{父啊，}他说，\Quote{宽恕他们，因为他们不知道自己在做什么。}如果他毫无俊美，那么，那个说\Quote{告诉我，我的灵魂所爱的}的人，怎么会爱他？她怎么会爱他？她怎么会为他燃烧？她怎么会如此害怕离开他？她怎么会在他内有这么大的喜乐，以至于她的唯一惩罚就是没有他？如果他不是美丽的，还有什么可以使他被爱？但是，如果他显现给她的样子，与他向那些迫害他、不知道自己在做什么的盲人所显现的相同，她怎么能这样爱他？那么，她以什么作为爱他呢？以\Quote{你在世人中最为俊美，你的嘴唇满了恩宠}。那么，从这嘴唇：\Quote{告诉我，我的灵魂所爱的。告诉我，}她说，\Quote{我所爱的，}不是我的肉体，而是\Quote{我的灵魂所爱的。告诉我，你在哪里放牧，中午在哪里躺卧；免得我像蒙着脸的人，撞上你同伴的羊群。}\par

7. 这似乎隐晦，确实隐晦；因为这是神圣婚床的奥迹。因为她说：\BibleQuote{君王带我进入了祂的内室。}这样的内室是一个奥迹。但你们这些并非不敬而被隔绝于此内室之外的人，要听你们是什么，并与她同声说，如果你们与她一起爱（事实上你们与她一起爱，如果你们在她内）；大家同说，却让一人说，因为合一说话：\Quote{\BibleQuote{告诉我，我灵魂所爱的。} \BibleQuote{因为他们一心一意，归向天主，同心合意。} \BibleQuote{告诉我，祢在何处牧放，祢在何处午间躺卧？}}午间代表什么？\Quote{炎热，并极大的光明。}所以，\Quote{告诉我祢的智者是谁，}在心神上热切，在道理上明达。\Quote{告诉我祢的右手，和那些内心受教于智慧的人。}愿我能在祢的身体上依附他们，与他们结合，与他们共享祢。那么，告诉我，\BibleQuote{告诉我，祢在何处牧放，祢在何处午间躺卧；}免得我跌入那些对祢说不同之事、对祢怀有不同感受；相信关于祢的不同之事、宣讲关于祢的不同之事；他们有自己的羊群，却是祢的同伴；因为他们靠祢的餐桌生活，并处理祢餐桌的圣事。同伴之所以如此称呼，是因为他们一同吃饭，如同共食者。这样的人在圣咏中受责备：\BibleQuote{因为若是我的仇敌说我大话，我必向他隐藏；若是那恨我的向我夸大，我必向他隐藏；但你是与我同心的人，我的向导，我的知己，与我一同吃甜蜜食物，在天主殿中我们同心同行。}那么，为何现在他们与主殿不合而分裂呢？岂不就是因为\BibleQuote{他们从我们中间出去了，却不是属于我们的？}所以，\BibleQuote{我灵魂所爱的，}使我不至跌入这样的、祢的同伴，但如同三松的同伴，他们对朋友不守信，却想败坏他的妻子。因此，使我不至跌入像这样的人，\Quote{使我不至撞到他们，}即跌入他们，\Quote{如同蒙着面的人，}即被隐藏的人，晦暗不明，而非立于山上。\BibleQuote{告诉我，} \BibleQuote{我灵魂所爱的，祢在何处牧放，祢在何处午间躺卧；}谁是祢特别安息的智慧忠信者，免得我偶然如在黑暗中跌入羊群，不是祢的羊群，而是祢同伴的羊群。因为祢没有对伯多禄说：\BibleQuote{牧放你的羊，}而是：\BibleQuote{牧放我的羊。}\par

8. 那么，让\BibleQuote{善牧}和\BibleQuote{在世人中俊美出众者}回答这位蒙爱者；回答祂从世人中使其美丽的那一位。听祂如何回答，并要明白，警惕祂所警告的，爱祂所劝勉的。那么祂如何回答呢？何等缺乏温柔的爱抚，不，对祂的爱抚祂回报以严厉！祂尖锐是为要紧密地束缚她，好保守她。\BibleQuote{如果你不认识自己，}祂说，\BibleQuote{女子中最美丽者：}因为无论其他女子因你新郎的恩赐而多么美丽，她们是异端，外在装饰美丽，内心却不；她们外表美丽，外在发光，以正义之名伪装自己；\BibleQuote{但君王的女儿一切美丽都在内。}因此，\BibleQuote{如果你不认识自己，}即你是一体的，你遍及万民，你是贞洁的，你不应受邪恶同伴的混乱交往而败坏自己。\BibleQuote{如果你不认识自己，}即，在正直中，\BibleQuote{祂将你许配给我，要把你们当作贞洁的童女献给基督；}并且你应在正直中将你自己献给我，免得因邪恶的交往，\BibleQuote{正如蛇以诡诈诱惑了厄娃，同样你们的思想也会从我纯洁中被败坏。} \BibleQuote{如果，}我说，\BibleQuote{你不认识自己}是如此，\BibleQuote{你去吧；去吧。}因为对其他人我要说：\BibleQuote{进入你主人的福乐吧。}对你不说：\BibleQuote{进来；}而是：\BibleQuote{去吧；}使你可以成为那些\BibleQuote{从我们中间出去的人}中的一员。\BibleQuote{去吧。}就是说，\BibleQuote{如果你不认识自己，}那么，\BibleQuote{去吧。}但如果你认识自己，就进来。但是，\BibleQuote{如果你不认识自己，}那么，\BibleQuote{你去吧，跟着羊群的脚踪，在牧人们的帐幕旁牧放你的山羊。你去吧，跟着脚踪，}不是\BibleQuote{羊群的，}而是\BibleQuote{羊群的，并牧放，}不像伯多禄那样，\BibleQuote{我的羊，}而是\BibleQuote{你的山羊；在帐幕旁，}不是\BibleQuote{牧人的，}而是\BibleQuote{牧人们的；}不是合一，而是分裂；不是立在那里有\BibleQuote{一个羊群和一个牧人}的地方。这位蒙爱者被坚定、被建立、被加强，预备好为她新郎而死，与她新郎而活。\par

9. 我从神圣的 \BibleBook{}{雅歌}——新郎与新娘的一种婚歌（因为这是属灵的婚配，我们必须在其中以极大的纯洁生活，因为基督在精神上赐予了教会祂的母亲在肉体上所拥有的，即同时是母亲和童贞）——中所引用的这些话，我说，多纳徒派将这些话以极其不同的含义曲解为他们自己颠倒的意义。至于我如何不向你们隐瞒，以及你们可以如何回答他们，我将，靠主的帮助，尽可能简要地推荐给你们。当我们开始用遍布全世界的教会合一之光来压迫他们，并要求他们向我们出示圣经中的任何证据，证明天主曾预言教会应在非洲，仿佛所有其他民族都已失落；他们习惯于把这证据挂在嘴边，并说：\Quote{非洲在正午阳光之下；那么教会}他们说，\Quote{询问主祂在哪里牧放，祂在哪里躺卧；祂回答说：\InnerQuote{在正午阳光之下；}}仿佛那提问者的声音是：\BibleQuote{告诉我，我心所爱的祢在哪里牧放，祢在哪里躺卧；}而回答者的声音是：\BibleQuote{在正午阳光之下；}也就是说，在非洲。如果那么是教会提问，而主回答祂在哪里牧放，在非洲，因为教会是在非洲；那么提问者就不在非洲。\BibleQuote{告诉我，}她说，\BibleQuote{我心所爱的祢在哪里牧放，祢在哪里躺卧；}而祂向某个在非洲之外的教会回答说：\BibleQuote{在正午阳光之下，}在非洲我躺卧，在非洲我牧放，仿佛是说：\Quote{我不在你们中间牧放。}我重申，如果提问者是教会，这是无人争议，连他们自己也不否认的；而他们听到了关于非洲的事；那么提问者就在非洲之外；而因为它是教会，教会就在非洲之外。\par

10. 但是请看，我承认非洲在正午阳光之下；尽管埃及更在子午线下，比非洲更在正午阳光之下。现在这位牧人在埃及如何，知道的人会承认；不知道的人，让他们去询问祂在那里聚集了何等大的羊群，祂拥有何等多的完全轻视世界的圣男圣女。那羊群增长到如此地步，甚至从那里驱逐了迷信。且不提它如何在其增长中从那里驱逐了整个偶像迷信，那些偶像曾牢固地植根在那里；我承认你们所说，邪恶的同伴；我完全承认，我同意非洲在南方，并且非洲就是那所说的\BibleQuote{祢在哪里牧放，祢在哪里躺卧在正午阳光之下？}所预指的。但你们同样要注意，直到这一点，这些是新娘的话，而非新郎的话。到此为止是新娘在说：\BibleQuote{告诉我，我心所爱的祢在哪里牧放，祢在哪里躺卧在正午，免得我偶然像蒙着脸的一样遇见。}哦，你这又聋又瞎的人，如果在\Quote{正午}你看见了非洲，为什么在她\Quote{蒙着脸的}身上你看不见新娘？\BibleQuote{告诉我，}她说，\BibleQuote{我心所爱的祢。}无疑她是对她的配偶说话，当她说\Quote{whom} [阳性] \Quote{my soul loves。}正如如果说：告诉我，你（阴性）\Quote{my soul loves；}我们就当明白是新郎在对祂的新娘说这些话；所以当你听到\BibleQuote{告诉我，你（阳性）} \BibleQuote{我心所爱的祢在哪里牧放，祢在哪里躺卧；}再加上，以下也属于她的话：\BibleQuote{在正午。}我是在问：\BibleQuote{祢在哪里牧放在正午，免得我偶然像蒙着脸的一样落入祢同伴的羊群中。}我完全同意，我承认你们对非洲的理解；它由\Quote{正午}所预指。但正如你们所理解的，海外的基督教会正在对她的配偶说话，害怕落入非洲的错误中：\BibleQuote{我心所爱的祢，告诉我，}教导我。因为我听说\Quote{在正午，}即在非洲，有两派，实际上是许多分裂。\BibleQuote{告诉我，}那么，\BibleQuote{祢在哪里牧放，}哪些羊属于祢，祢命令我在那里爱哪个羊栈，我应该与什么联合。\BibleQuote{免得我偶然像蒙着脸的一样遇见。}因为他们嘲笑我好像我被隐藏，他们嘲笑我好像被毁灭，好像我别处不存在。\BibleQuote{免得，}那么，\BibleQuote{像蒙着脸的一样，}如同被隐藏，\BibleQuote{我落入羊群中，}即是，落在异端者的集会中，祢的同伴；多纳徒派、马克西米安派、罗加图斯派以及所有其他从外聚集、因而四散的害虫；\BibleQuote{告诉我，}我祈求祢，如果我必须在那里寻找我的牧人，使我不致落入重洗的深渊。我劝勉你们，我恳求你们，凭着如此婚配的神圣性，爱这教会，在这圣教会中，成为这教会；爱善牧，如此美丽的配偶，祂不欺骗任何人，不愿任何人丧亡。也要为四散的羊祈祷；使他们也能来，使他们也能认识祂，使他们也能爱祂；使成为一牧一栈。让我们转向主，等等。\par

\SourceRecord{Translated by R.G. MacMullen.；Nicene and Post-Nicene Fathers, First Series；Vol. 6.；Edited by Philip Schaff.；Buffalo, NY: Christian Literature Publishing Co.,；1888.；<http://www.newadvent.org/fathers/160388.htm>.}

% structure: p0001 source=h1 target=section reason=page-title

\section{新约讲道集89}

\Emph{[CXXXIX. Ben.]}\par

\Emph{论福音之言}，\BibleRef{若 10:30}，\Emph{，\BibleQuote{我与父原是一体。}}\par

1. 你们已经听见主天主耶稣基督——天主的独生子、无母而生、由童贞之母无父而生——所说的话：\BibleQuote{我与父原是一体。}领受这话，相信它，好使你们能达至理解。因为信德先于领悟，使领悟成为信德的赏报。先知说得最明确：\BibleQuote{除非你们相信，你们不能明白。}因此，凡单纯宣讲的当信；凡精确论究的当悟。首先，为以信德浸润你们的心智，我们向你们宣讲基督，天父的独生子。为何加上\Quote{独生子}？因为那位独生子的父，因恩宠有许多儿子。其余的一切圣人，因恩宠是天主的儿子，惟独祂因本性而是。因恩宠而成为天主子的人，并非父之所是。没有圣人曾敢说那独生子所说的话：\BibleQuote{我与父原是一体。}那么祂不是我们的父吗？若不是我们的父，我们祈祷时为何说\BibleQuote{我们的天父，你在天上}？但我们是祂按己意所立的儿子，并非祂本性所生的儿子。的确，祂也生了我们，但如经上所说，作为义子，由祂收养的恩宠所生，非由本性。我们也如此被称为，因为\BibleQuote{天主召叫我们成为义子}；我们虽是义子，却是人。祂被称为独生子、独生者，因为祂就是父之所是；但我们是人，父是天主。所以，祂是父之所是，祂说了并真说了：\BibleQuote{我与父原是一体。}什么是\Quote{原是一体}？同一性体。什么是\Quote{原是一体}？\OriginalTerm{同一实体}{Substance}。\par

2. 也许你们不完美地理解什么是\Quote{同一实体}。我们当努力使你们理解；愿天主助我说者及你们听者；助我说恰当而相宜的话，助你们首先且超越一切地相信，然后尽所能理解。那么什么是\Quote{同一实体}？让我用比喻向你们说明，使那不完全理解的得以借实例清晰。譬如，假设天主是金，祂的圣子也是金。若不可将天上之事与地上之事类比，怎会写着\BibleQuote{盘石就是基督}？所以，父是什么，圣子就是什么；如我所言，举例说：\Quote{父是金，圣子是金。}因为谁说\Quote{圣子并非父的同一实体}，就等于说\Quote{父是金，圣子是银}。若父是金，圣子是银，则独生子已从父堕落。人生人；生父是什么实体，所生的儿子也是同一实体。什么是\Quote{同一实体}？一人是男人，另一人也是男人；一人有灵魂，另一人也有灵魂；一人有身体，另一人也有身体；一人是什么，另一人也是什么。\par

3. 但\OriginalTerm{亚略异端}{Arian heresy}反驳说。它对说什么？\Quote{注意你所说的话}？我说了什么？\Quote{你将人之子比作天主子。}固然可以相比，但并非如你所想那样严格，而是作为比喻。但如今请告诉我，你对此作何解释。它说：\Quote{你难道不见，生父在年龄上更大，所生之子更小？那么你们如何说？告诉我，你们怎能说父与子、天主与基督平等？当你见人生子，子小父大？}你这聪明人，在永恒中寻求时间；在没有时间之处，寻求年龄的差别！当父年长、子幼时，二者都在时间中：一人成长，因为另一人衰老。因为就本性而言，父所生者并非更小，如我所说，不是因本性，而是因年龄。你愿知道如何就本性并非生得更小吗？且等他长大，他便与父同等。连一个小男孩也借成长达到父的体量。而你断言天主子生而更小，永不成长，甚至借成长达到父的体量。那么，人所生的人子，处境比天主子更好。如何？因为前者成长，达到父的体量。但基督，若如你所说，生而更小，以至于永为更小，且在此不可期望岁月的增长。如此，你便说本性上有差异。但为何这样说？只因你不信圣子与圣父是同一实体。最后，你先承认祂是同一实体，然后称祂更小。试看一人：他是什么实体？他是人。他所生者是什么？他更小，但他是人。年龄不等，本性相等。那么你也说：\Quote{父是什么，圣子就是什么，但圣子更小}？说吧，向前一步，说\Quote{同一实体，只是更小}；你便接近祂的平等。因为这不是一小步，不是一小步接近承认祂的平等，只要你承认祂虽然是更小，却是同一实体。\Quote{但祂并非同一实体}，这是你所说的。既然如此，这里就是金和银；你所说的，好比人生了一匹马。因为人是同一实体，马是另一实体。若圣子是另一实体，则父生了一个怪物。因为当一个受造物，即女人，生出人以外的任何东西，便称为怪物。但为使它不是怪物，所生者当与生他者同类，即人和人，马和马，鸽和鸽，雀和雀。\par

4. 祂赐给受造物生育与自己相同的后代。祂赐给受造物——终有一死的、地上的受造物——生育与自己相同的后代；难道你认为祂不能为自己保留这一能力吗？祂是谁？祂是从永远就存在的。那位没有时间起点的，难道要生一个儿子，与祂自身不同，一个堕落的儿子吗？听听说天主独生子属于另一本体是何等亵渎。如果真是这样，那祂就是堕落的。如果你对任何人的儿子说：\Quote{你是堕落的}，这是多大的冒犯！然而，人的儿子在何种意义上被称为堕落？例如，他的父亲勇敢，他却懦弱胆小。如果有人看见他，想要责骂他，当想到他勇敢的父亲时，会对他说什么？\Quote{滚开，你这堕落的东西！}什么是\Quote{堕落的东西}？\Quote{你的父亲是个勇敢的人，而你却因恐惧而发抖。}被这样说的人，是因某种过失而堕落，在本性上他是平等的。什么是\Quote{在本性上平等}？他是人，他的父亲也是人。但一个勇敢，一个懦弱；一个胆大，一个胆怯；然而都是人。所以他因某种过失而堕落，并非因本性。但当你说天父的独生子、唯一子堕落时，你无非是说祂不是父之所是；而且你并不是说祂已经出生后才变得堕落，而是祂被生下来就是这样。谁能忍受这种亵渎？如果他们能以某种方式看到这种亵渎，他们就会逃离它，并成为公教徒。\par

5. 但是，弟兄们，我该说什么呢？让我们不要对他们发怒；但要为他们祈祷，求天主赐给他们悟性；因为他们也许生来如此。什么是\Quote{生来如此}？他们从父母那里接受了他们所持守的。他们将自己的出身置于真理之上。让他们成为他们所不是的，这样他们才能保持他们所是的；也就是说，让他们成为公教徒，这样他们才能保持他们作为人的本性；好使天主在他们内的受造物不灭亡，让天主的恩宠加给他们。因为他们幻想通过侮辱圣子来尊敬圣父。当你对他说：\Quote{你在亵渎}；他回答说：\Quote{我为何亵渎？}\Quote{因为你说圣子不是父之所是。}他回答我说：\Quote{不，是你亵渎。}为什么？\Quote{因为你使圣子与圣父平等。}\Quote{我确实希望圣子与圣父平等，但这是使一个外人平等吗？当我使祂的独生子与祂平等时，圣父高兴；祂高兴，因为祂不嫉妒。既然天主不嫉妒祂的独生子，因此祂生了与祂自身相同的圣子。你对圣子和圣父本人都不义，因为为了尊敬祂而侮辱圣子。因为实际上，你正是为了不冒犯圣父才说圣子不是同一本体的。我很快就会向你证明，你冒犯了两者。}\Quote{怎么？}他说。\Quote{如果我对某人的儿子说：\InnerQuote{你是堕落的，你不像你的父亲；堕落的，你不是你父亲所是的。}儿子听见了，就生气，说：\InnerQuote{难道我生来就是堕落的吗？}父亲听见了，更加生气。他在愤怒中说什么？\InnerQuote{难道我生了一个堕落的儿子？如果我是一回事，而我生了另一回事，我就生了一个怪物。}那么，当你希望通过侮辱一位来尊敬另一位时，你侮辱了双方？你得罪了圣子，但你不会平息圣父。当你通过侮辱圣子来尊敬圣父时，你得罪了圣子和圣父。你将逃向谁？你将逃向谁？当圣父对你发怒时，你逃向圣子吗？祂对你说什么？\InnerQuote{你逃向谁？逃向我，你所认为堕落的那一位？}当圣子被冒犯时，你奔向圣父吗？祂也对你说：\InnerQuote{你逃向谁？逃向我，你曾说（我）生了一个堕落的儿子？}}这些话对你们足够了；持守它，记住它，铭刻在你们的信德中。但为使你们能领悟，请向天主——圣父和圣子，他们原是一体——倾心祈祷。\par

\SourceRecord{Translated by R.G. MacMullen.；Nicene and Post-Nicene Fathers, First Series；Vol. 6.；Edited by Philip Schaff.；Buffalo, NY: Christian Literature Publishing Co.,；1888.；<http://www.newadvent.org/fathers/160389.htm>.}

% structure: p0001 source=h1 target=section reason=page-title

\section{新约讲道集第九十篇}

\Emph{[CXL. Ben.]}\par

\Emph{论福音的话}，\BibleRef{若 12:44} \Emph{，\BibleQuote{凡信我的，不是信我，而是信那派遣我来的。}反对亚略派主教马克西米努斯的某种说法，他在非洲与伯爵塞吉斯武特同在时散布其亵渎之辞。}\par

1. 诸位弟兄，我们听主说：\BibleQuote{凡信我的，不是信我，而是信那派遣我来的}，这是什么意思？信基督是好的，何况祂还亲口说了你们刚才听到的话，即\BibleQuote{祂来到世界是光明，凡信祂的，不会在黑暗中行走，却要有生命的光}。因此，信基督是好的，不信基督是极大的恶。但因为基督是子，祂无论是什么，都是出自父；而父不是出自子，父是子的父；祂确实向我们推荐对祂的信德，却将尊荣归于祂的本源。\par

2. 如果你们愿意继续作天主教徒，就当持守这坚定而确切的真理：天主圣父在无时间之内生了天主圣子，并在时间内使祂由童贞女所生。第一个诞生超越时间；第二个诞生光照时间。两个诞生都是奇妙的：一个没有母亲，另一个没有父亲。当天主生圣子时，祂是由自己生的，并非由母亲；当母亲生她的儿子时，她是童贞女所生，并非由男子。祂由父所生，无始；祂由母亲所生，如同今日在指定的开始。祂由父所生，创造了我们；祂由母亲所生，重新创造了我们。祂由父所生，使我们存在；祂由母亲所生，使我们不至丧亡。父生了祂，与父同等；子所有的一切，都来自父。但天主父是什么，则并非来自子。因此，我们说父是天主，非由他者；子是天主，来自天主。所以，子所行的一切奇事、所说的一切真话，都归功于祂所自出者；但祂不能是别于祂所自出者。亚当被造为人，他有可能变得不同于他被造的样子，因为他是被造成义的，他也能成为不义的。但天主的独生子，祂是什么，是不能改变的；祂不能变成别的，不能缩小，祂不能不是祂所是的，祂不能不与父同等。但无疑，父藉着祂的生育，将一切赐给了子，赐给了一位一无所缺者；毫无疑问，这同等的本身，父也赐给了子。父是如何赐予的呢？祂是否生了较少的子，然后增加祂以完成祂的形态，使之相等？若祂这样做了，那就是赐予一位有需要者。但我已经告诉你们，你们最应持守的，即子所有的一切，都是父所赐予的，是藉着生育赐予的，并非因为祂有所需要。如果父藉着生育赐予祂，而且不是因祂有所需要，那么无疑，父既赐予了祂同等，在赐予同等时就生了祂为同等。尽管一位是一位位格，另一位是另一位位格，但二者并非一是物、另一是物；而是，一位是什么，另一位也是什么。那一位，并非另一位；但一位是什么，另一位也是什么。\par

3. \BibleQuote{那派遣我来的，}祂说，你们已听见了；\BibleQuote{那派遣我来的，}祂说，\BibleQuote{祂给了我一条诫命，要我说什么，讲什么；我知道祂的诫命是永生。}这是若望福音，要持守它。\BibleQuote{那派遣我来的，祂给了我一条诫命，要我说什么，讲什么；我知道祂的诫命是永生。}愿祂允许我说出我所愿意的！因为我的贫乏和祂的丰盛使我窘迫。祂说：\BibleQuote{祂给了我一条诫命，要我说什么，讲什么；我知道祂的诫命是永生。}你们在若望圣史的书信里查找他论基督所说的话。\BibleQuote{我们当信}，他说，\BibleQuote{祂的真子耶稣基督。这是真天主和永生。}\Quote{真天主和永生}是什么？天主的真子就是\Quote{真天主和永生}。祂为什么说\Quote{信祂的真子}？因为天主有许多儿子，所以要用加上\Quote{真子}来区分。不是简单地说祂是子，而是如我所说，加上\Quote{真子}；因为天主有许多儿子，所以要区分。我们是因恩宠而作子，祂是因本性而作子。我们由父藉着祂而造；祂本身就是父所是的；我们也同样是天主所是的吗？\par

4. 但有人遇到我们，却不知自己说什么，说：因此才说\Quote{\BibleQuote{我与父原是一体}；因为他们彼此有意志的一致，并非因为子的性体与父的性体完全相同。因为宗徒们（这是他说的话，不是我说）也与父和子是一体。}可怕的亵渎！\Quote{宗徒们，}他说，\Quote{与父和子是一体，因为他们服从了父和子的旨意。}他竟敢这样说！那么，让保禄说\Quote{我与天主是一体}，让伯多禄说，让每一位先知说\Quote{我与天主是一体}。他们并不这样说；但愿他们不这样说。他们知道自己是另一种本性，一种需要被拯救的本性；他们知道自己是另一种本性，一种需要被光照的本性。没有人说\Quote{我与天主是一体}。无论他取得多大进步，如何超越他人，在圣德上多么卓越，以何等崇高的德性出众，他永不会说\Quote{我与天主是一体}。因为如果他有了卓越，因而这样说，那么他因这样说就会丧失他所拥有的。\par

请相信，圣子与圣父同等；然而圣子出自圣父，但圣父并非出自圣子。本原在于圣父，同等在于圣子。因为若祂不相等，祂就不是真正的圣子。因为我们说什么呢，诸位弟兄？若祂不相等，祂就是较小的；若祂是较小的，我问那需要被拯救的本性，在其错误信仰中，\Quote{祂如何是生得较小？}回答，祂作为较小的，是成长呢，还是不成长？若祂成长，那么圣父就变老了。但若祂将永远是祂所生时的样子；若祂生来较小，祂将继续较小；带着这个缺失，祂将是完美的；生来完美，却带着圣父形态的缺失，祂永远不能达到圣父的形态。你等不敬之人就是这样攻击圣子；异端者就是这样亵渎圣子。那么公教信仰说什么？圣子是出自圣父天主的真天主；圣父天主，并非出自圣子天主；但是圣子天主与圣父同等，生而同等；非生而较小，非被造为同等，而是生而同等。圣父是什么，那所生的也是什么。圣父曾没有圣子吗？绝对不可！去掉你们的\Quote{永远}，那里没有时间。圣父永远，圣子永远；圣父无时间初始，圣子无时间初始；圣父从未先于圣子，圣父从未无圣子。但既然圣子是出自圣父天主的真天主，而圣父是天主，却非出自圣子天主，愿尊敬圣子在圣父内不令我们不悦；因为尊敬圣子就是尊敬圣父，不减损祂本身的天主性。\par

6. 因为那时我正在讲我所引用的，\BibleQuote{我认识}，祂说，\BibleQuote{祂的诫命就是永生。}注意，诸位弟兄，我所说的话；\BibleQuote{我知道祂的诫命就是永生。}我们在同一部若望福音中读到关于基督的话：\BibleQuote{祂是真天主和永生。}如果圣父的诫命是\Quote{永生}，而基督圣子本身是\Quote{永生}，那么圣子本身就是圣父的诫命。因为那圣父的圣言，怎么不是圣父的诫命呢？或者，如果你按肉身的方式理解圣父给圣子的诫命，好像圣父对圣子说：\Quote{我命令你这事，我愿你做那事；}那么祂对独一圣言说话时，用了什么话语？当祂给圣言诫命时，祂会寻找话语吗？那么，圣父的诫命既是\Quote{永生}，而圣子本身也是\Quote{永生}，请相信并接受，相信并领悟，因为先知说：\BibleQuote{你们若不相信，必不得领悟。}你们不理解吗？扩展心胸。听宗徒的话：\BibleQuote{扩展心胸，不要与不信者同负一轭。}那些在领悟之前不愿相信的人，是不信者。因为他们决意成为不信者，他们将停留在无知中。那么让他们相信，以便得以领悟。最确实的是，圣父的诫命就是\Quote{永生}。因此，圣父的诫命就是今天所生的圣子本身；一个不是按时赐予的诫命，而是生而存在的诫命。\BibleBook{若望}{福音}锻炼我们的心智，精炼它们，使它们脱去肉身化，好让我们对天主的思考不是按肉身的方式，而是属灵的方式。那么，诸位弟兄，愿这些已足够；免得在冗长的论述中，遗忘的睡意悄然侵袭你们。\par

\SourceRecord{Translated by R.G. MacMullen.；Nicene and Post-Nicene Fathers, First Series；Vol. 6.；Edited by Philip Schaff.；Buffalo, NY: Christian Literature Publishing Co.,；1888.；<http://www.newadvent.org/fathers/160390.htm>.}

% structure: p0001 source=h1 target=section reason=page-title

\section{新约讲道集 第九十一篇}

\Emph{[第一四一篇 本笃版]}\par

\Emph{论福音之言}，\BibleRef{若14:6}，\Emph{\BibleQuote{我是道路、真理、生命。}}\par

1. 此外，当圣福音被诵读时，你们听到主耶稣所说的话：\BibleQuote{我是道路、真理、生命。} 真理和生命是人人渴望的，但并非人人都找得到道路。那永恒的生命、不变、可理解、理智、智慧、使人智慧的天主，连这世界的某些哲学家也已看见。那固定、稳固、不变的真理，其中含有一切受造物的原理，他们确实看到了，但不过是远远地看到；他们看到了，但处在错误之中；因此，他们没有找到达到那如此伟大、不可言喻、福乐的财富的道路。因为他们甚至看见了（就人的能力所能见的）创造者借着受造物，工人借着工作，世界的塑造者借着世界，宗徒\Person{保禄}可以作证，基督徒当然应当相信他。因为他在谈到这种人时曾说：\BibleQuote{天主的义怒从天上显示出来，惩罚一切不虔敬的人。} 你们知道，这是宗徒\Person{保禄}的话：\BibleQuote{天主的义怒从天上显示出来，惩罚一切不虔敬和不义的人，他们以不义持有真理。} 难道他说他们不持有真理吗？不，而是：\BibleQuote{他们以不义持有真理。} 他们所持有的，是好的；但他们持有的方式，是坏的。\BibleQuote{他们以不义持有真理。}\par

2. 现在他想到，可能有人会对他说：\Quote{这些不虔敬的人从哪里持有真理？天主曾对他们任何人说过话吗？他们像\OriginalTerm{以色列子民}{Israelites}借着\Person{梅瑟}那样领受了法律吗？那么，他们从哪里持有真理，即使是在这不义之中呢？} 请听下文，他揭示出来。他说：\BibleQuote{因为天主所可认识的，在他们身上是明显的，因为天主给他们显明了。} 天主给他们显明了那些没有接受法律的人？请听他是如何显明的。\BibleQuote{因为他的不可见的事物，借着所造之物，就可以明明看见，被人理解。} 请问世界，天空的美，星辰的光辉与秩序，太阳，它足以供应白昼，月亮，夜晚的安慰；请问大地，出产蔬菜和树木，充满动物，以人类为装饰；请问海洋，充满多少和何种的鱼类；请问空气，充满多少飞鸟；请问万物，看看它们是否仿佛以自己的语言回答你：\Quote{天主造了我们。} 这些事物，卓越的哲学家们已经探求，并借着艺术认识了造物主。那么怎样？为什么天主的义怒显示出来惩罚这不虔敬？\BibleQuote{因为他们以不义持有真理。} 让他来，让他说明如何。因为他们如何认识他，他已经说过。\BibleQuote{他的不可见的事物，就是天主，由所造之物得以明白看见，他的永恒德能和神性，以致他们无可推诿。因为他们虽然认识天主，却不当作天主光荣他，也不感谢他，反倒在思想上成为虚妄，他们愚昧的心就昏暗了。} 这是宗徒的话，不是我的：\BibleQuote{他们愚昧的心就昏暗了；自称为聪明，反而成了愚拙。} 他们因好奇探索而找到的，因骄傲而丧失了。\BibleQuote{自称为聪明，}即把天主的恩赐归于自己，\BibleQuote{反而成了愚拙。} 我说，这是宗徒的话：\BibleQuote{自称为聪明，反而成了愚拙。}\par

3. 请指出，证明他们的愚拙。请指出，宗徒啊，正如你已向我们表明了他们如何能够达到对天主的认识，因为\BibleQuote{他的不可见的事物，借着所造之物，就可以明明看见，被人理解}；现在请指出如何\BibleQuote{自称为聪明，反而成了愚拙}。请听：因为\BibleQuote{他们}将\BibleQuote{不可朽坏的天主的荣耀，变成了可朽坏的人、飞禽、走兽和昆虫的形态。} 因为外教人用这些动物的形像为自己造了神明。你找到了天主，却敬拜\OriginalTerm{偶像}{idol}。你找到了真理，却以不义持有这真理。你借着天主的工程所认识的，借着人的工程丧失了。你曾思考宇宙，收集了天空、大地、海洋和一切元素的秩序；你却不注意这一点：世界是天主的工程，偶像是\OriginalTerm{木匠}{carpenter}的工程。如果木匠给了他形态之后，也能给他一颗心，那么木匠就会被他自己的偶像所崇拜。因为，人啊，如同天主是你的塑造者，同样偶像的塑造者是一个人。谁是你的天主？是那造你的。谁是木匠的天主？是那造他的。谁是偶像的天主？是那造它的。那么，如果偶像有心，它岂不是要崇拜造它的木匠吗？请看，他们是以何等的不义持有真理，却找不到通往他们所看见的那财富的道路。\par

4. 但是基督，因祂与圣父、真理、生命同在，祂是天主圣言，关于祂曾有话说：\Quote{生命是人的光}；因我说祂与圣父、真理、生命同在，而我们没有道路可以到达真理，圣子，祂永远在圣父内是真理和生命，因取了人性而成为道路。你们将祂作为人而行走，就来到天主面前。藉着祂你们去，向祂你们去。不要寻找任何别的道路来到祂那里，除了祂自己。因为如果祂未曾屈尊成为道路，我们就会永远迷失。于是祂成了你们应当前来的道路；我不对你们说，寻找道路。道路本身已来到你们面前，起来行走吧。行走，要用生命，而不是用双脚。因为许多人用脚行走得好，却用生命行走得坏。因为有时即使是行走得好的人，也跑道了路外。因此你们会找到生活得好却不是基督徒的人。他们跑得好，但他们没有跑在道路上。他们跑得越多，迷失得越多，因为他们离开了道路。但如果这样的人来到道路上，并坚持在道路上，哦，他们的安全是多么大！因为他们既行走得好，又不迷失！但如果他们不坚持在道路上，不论他们行走得多好，唉！他们真是可悲！因为在路上跛行，也比在路外健步行走更好。亲爱的诸位，这对你们已足够了。让我们转向主，等等。\par

\SourceRecord{Translated by R.G. MacMullen.；Nicene and Post-Nicene Fathers, First Series；Vol. 6.；Edited by Philip Schaff.；Buffalo, NY: Christian Literature Publishing Co.,；1888.；<http://www.newadvent.org/fathers/160391.htm>.}

% structure: p0001 source=h1 target=section reason=page-title

\section{\WorkTitle{新约讲道集第92篇}}

\Emph{[CXLII. Ben.]}\par

\Emph{论同一福音经文，\BibleRef{若 14:6}，\BibleQuote{我是道路}等。}\par

1. 神圣的读经提升我们，使我们不致因绝望而破碎；又再次使我们恐惧，以免因骄傲而摇摆不定。但要持守中道、真理之道、窄路，如同介于绝望的左手与傲慢的右手之间，对我们而言是极其困难的，若非基督说过：\BibleQuote{我是道路、真理、生命。}犹如祂说：\Quote{你们要走什么路？\InnerQuote{我是道路}。你们要去哪里？\InnerQuote{我是真理}。你们要住在哪里？\InnerQuote{我是生命}。}那么，让我们满怀信心地在道路上行走；但要提防路旁的陷阱。仇敌不敢在道路上设陷阱，因为基督就是道路；但毫无疑问，他在路旁却不停地这么做。为此，圣咏也说：\BibleQuote{他们在路旁为我设了绊脚石。}另一处经文说：\BibleQuote{要记住你行走在陷阱之中。}这些在我们行走中的陷阱并不在道路上，但它们却是在\Quote{路旁}。你怕什么？你为何惊慌？只要你行在道路上。那么，若你离弃道路，就害怕吧。因为仇敌甚至被允许在路旁设陷阱，正是为了免得你在狂喜的安逸中离弃道路，而陷入陷阱。\par

2. 谦卑的基督是道路；基督是真理和生命，是高高在上的基督和天主。你若行走在谦卑中，必能达到崇高。你若身为软弱者而不轻视谦卑，必能在崇高之中安然屹立。基督的谦卑有何原因？唯有你的软弱。因为你的软弱完全且无可救药地压迫着你，正是这情况使得如此伟大的医生来到你面前。因为即便你的疾病严重到你能亲自去找医生，这软弱或许还显得可以忍受。但因为你不能去他那里，他便到你这里来。他来教导谦卑，使我们能藉此回归；因为骄傲不容许我们回归生命，甚至使我们远离生命。人的心因着骄傲而反抗天主，在健康时忽视了祂救人的诫命，灵魂便堕入软弱；让这软弱的她学会听从那位在她强壮时所轻视的祂。让她听从祂好能兴起，她曾轻视祂以致跌倒。让她终于由于经验而学会听从，她曾因诫命的教导而不愿获得。因为她的痛苦教导了她，离开主而行淫是多么邪恶的事。因为离开那单纯而唯一的善，投入这众多的享乐、爱世界和世俗的腐朽，就是离开主而行淫。祂如同对一个妓女说话，以警告她回头：非常多次，祂通过先知责备她为妓女，但并未绝望，因为那责备妓女的手中也掌握着洁净妓女的能力。\par

3. 因为祂不是以侮辱的方式责备，而是为使她面容羞愧而治愈她。圣经的呐喊是强烈的，对那些祂要通过医治而恢复的人，祂并不用谄媚温和对待。\BibleQuote{你们这些犯奸淫的人，难道不知道与世俗为友就是与天主为敌吗？}爱世俗使灵魂犯奸淫，爱世界的创造者使灵魂贞洁；但除非她为自己的腐败而脸红，她便不会渴望回到那贞洁的怀抱。让她羞愧而回归吧，她曾自夸不会回归。因此，阻碍灵魂回归的是骄傲。但责备的人不造成罪，而是显示罪。灵魂不愿看见的，被放在她眼前；她希望放在背后的，被带到她面前。看看你自己吧。\BibleQuote{你为什么看见你兄弟眼中的木屑，却看不见自己眼中的大梁？}那离开自己的灵魂被召回到自身。正如她离开了自己，她也离开了她的主。因为她注视自己，喜欢自己，迷恋自己的能力。她离开了他，不再停留在自身内；她被迫离开自身，被关在自身之外，落入外在的事物。她爱世界，爱时间的事物，爱世俗的事物；如果她只爱自己而忽略那位创造她的主，她就会立刻变小，立刻因为爱那较小的事物而衰败。因为她小于天主；是的，远小于，且正如受造物小于造物主。因此，应该爱天主，是的，应该如此爱天主，以至于如果可能，我们应忘记自己。那么这转变是什么？灵魂忘记了自己，但通过爱世界；现在让她忘记自己，但通过爱世界的创造者。她甚至从自身被赶走，可以说迷失了自己，不善于看见自己的行为，她为自己的不义辩护；她自高自大，在傲慢、纵欲、荣誉、权位、财富、虚荣的权力中夸耀自己。她受责备、受斥责，被显给自己看，不喜欢自己，承认自己的丑陋，渴望自己最初的美，那位在放荡中离开的人，在羞愧中回归。\par

4. 看来他在祈祷反对她，还是为她祈祷？他说：\Quote{愿羞耻充满他们的脸。}这似乎是一个敌人，似乎是一个仇敌。请听下文，看一个朋友能否作这样的祈祷。他说：\Quote{愿羞耻充满他们的脸，好让他们寻求你的名，主。}他恨那些他希望其脸被羞耻充满的人吗？请看他是多么爱那些他希望寻求主名的人。他是只爱，还是只恨？还是既恨又爱？是的，他既恨又爱。他恨你所有的一切，他爱你。什么是\Quote{他恨你所有的一切，他爱你}？他恨你所造的，他爱天主所造的。因为你所有的一切，除了罪，还有什么？而你除了是天主按他自己的肖像和模样所造的人，又是什么？你忽略了你是怎样被造的，却爱你自己所造的。你爱你身外你自己的作品，却忽略了在你内的天主的工程。你理应离开，理应堕落，是的，理应甚至从你自己那里离开；理应听到这句话：\Quote{一个去了不再回来的神魂。}不如听那呼唤并说\Quote{转向我，我必转向你}的那一位。因为天主并不是真的转开又转回；他恒常不变，他责备时也是不变的。他转开了，是因为你转开了自己。你从他那里堕落了，他却没有从你那里堕落。那么就听他对你说：\Quote{转向我，我必转向你。}因为这就是\Quote{我转向你，因为你转向了我}。他追赶那逃离者，他光照那回归者的脸。因为逃离天主，你要逃到哪里去？逃离那不被任何地方所包含、又无处不在的那一位，你要逃到哪里去？那解救归向于他的人，惩罚背离他的人。你逃离时有一位审判官；你回归时有一位父亲。\par

5. 但他因骄傲而肿胀，因这肿胀无法从狭窄的道路回来。那成为道路的一位喊道：\Quote{由窄门进入。}他试图进入，肿胀阻碍了他；他的尝试越有害，肿胀就越成为更大的障碍。因为狭窄刺激了他的肿胀；受了刺激，他会更加肿胀；肿胀更甚，他何时才能进入？那么，就让他消除肿胀吧。如何消除？让他服下谦卑的药方；让他为对抗肿胀而饮下这苦口却有益的杯；饮下谦卑之杯。他为何要挤压自己？那体积，不是因其大小，而是因其肿胀，不容许他。因为大小有坚实，肿胀有膨胀。那肿胀的人不要自以为身材高大；为了使他变得高大、结实、坚固，让他消除肿胀。让他不要贪恋眼前的事物，不要以这些消逝腐败之物的浮华为荣；让他听从那说\Quote{由窄门进入}的一位，他也说：\Quote{我是道路。}因为好像有一个肿胀的人问：\Quote{我该如何进入？}他说：\Quote{\InnerQuote{我是道路。}你们要借着我进入；你们只能借着我行走，借着我进入门。}因为正如他说：\Quote{我是道路}，同样也说：\Quote{我是门。}你为什么寻求从哪里回归、回归到哪里、从哪里进入？以免你迷失任何方向，他为你成为了一切。因此他简短地说：\Quote{要谦卑，要温良。}让我们听他极其清楚地这样说，好让你看见哪里是道路、什么是道路、道路通向哪里。你要到哪里去？但也许出于贪婪，你想拥有一切。他说：\Quote{我父已将一切交付于我。}你可能会说：\Quote{它们被交付给了基督，但交付给我了吗？}听宗徒说话；听，如我先前所说，免得你因绝望而破碎；听你如何被爱——当你没有任何值得被爱之处时，听你如何被爱——当你丑陋、畸形、在你身上没有任何配得被爱的事物之前，你首先被爱了，好让你成为配得被爱的。因为基督，如宗徒所说，\Quote{为不虔敬的人死了}。什么！你会说不虔敬的人配得被爱吗？我问，不虔敬的人配得什么？配得被定罪。你在这里会回答：\Quote{然而基督为不虔敬的人死了。}看，当你仍不虔敬时，为你做了什么；现在你虔敬了，为你预备了什么？\Quote{基督为不虔敬的人死了。}你曾想拥有一切；不要出于贪婪去渴望，要透过虔敬寻求，透过谦卑寻求。因为如果你这样寻求，你必将得着。因为你将拥有那创造万物的他，并借着他拥有一切。\par

6. 我说这话，并非出于推理。请听宗徒自己说：\Quote{祂既然没有怜惜自己的儿子，反而为我们众人把祂交出来，怎会不把万有也同祂一起赐给我们呢？}\BibleRef{罗 8:32} 看，贪婪的人，你拥有一切。你所喜爱的一切，要轻视它们，免得你被阻隔于基督之外，并紧握祂，在祂内你才能拥有一切。这位医师本身并不需要这样的药剂，但为了鼓励病人，祂饮用了祂并不需要的；祂仿佛对拒绝饮用的病人说话，在病人恐惧中扶起他，祂先饮了。祂说：\Quote{我所要喝的杯}\BibleRef{玛 20:22}，\Quote{我身上没有任何需要这杯来医治的，但我仍要喝，使需要喝的人不致不屑于喝。} 现在，弟兄们，人类在接受了这样的药剂后，还能继续生病吗？天主已经谦卑自下，人却仍然骄傲吗？让他听，让他学。祂说：\Quote{一切都由我父交付给了我。}\BibleRef{玛 11:27} 如果你渴望一切，你将在我的内与我一并拥有它们；如果你渴望父，你藉着我并在我的内将拥有祂。\Quote{除了子，没有人认识父。}\BibleRef{玛 11:27} 不要绝望，到子这里来。听接下来所说：\Quote{和子愿意启示给谁。}\BibleRef{玛 11:27} 你说：\Quote{我不能。你召我走那狭窄的道路；我无法从窄路进入。} 祂说：\Quote{凡劳苦和负重担的，你们都到我这里来。}\BibleRef{玛 11:28} 你的负担就是你的膨胀。\Quote{凡劳苦和负重担的，你们都到我这里来，我要使你们安息。你们背起我的轭，跟我学。}\BibleRef{玛 11:28-29}\par

7. 天使的导师呼喊，天主的圣言，藉以喂养一切理性灵魂而不匮乏的食粮，那滋养且保持完整的食粮，呼喊说：\BibleQuote{你们跟我学。}\BibleRef{玛 11:29} 让百姓听祂说：\BibleQuote{你们跟我学。}\BibleRef{玛 11:29} 让他们回答：\Quote{我们跟你学什么？} 因为我们必定要从伟大的工匠那里听到我不知道是什么，当祂说：\BibleQuote{你们跟我学。}\BibleRef{玛 11:29} 是谁说：\BibleQuote{你们跟我学}\BibleRef{玛 11:29}？祂形成了大地，分开了海洋和陆地，创造了飞鸟，创造了地上的动物，创造了水中游动的一切，在天空安置星辰，分开了昼夜，建立了穹苍，从黑暗中分出了光明；祂说：\BibleQuote{你们跟我学。}\BibleRef{玛 11:29} 祂或许要告诉我们，要我们与祂一同做这些事？谁能这样做？唯有天主能做到。祂说：\Quote{不要害怕，我不是在给你们加重任何负担。\InnerQuote{你们跟我学}这个，是为了你们而成了的。} 祂说：\Quote{\InnerQuote{你们跟我学，}不是要学习我所造的受造物。我也不要你们学习那些我已经赐给某些人（不是所有人）的能力，比如复活死人、使瞎子复明、打开聋子的耳朵；也不要你们希求像某些大事一样跟我学这些。} 门徒们欢欣鼓舞地回来，说：\BibleQuote{主，因着你的名，连恶魔都屈服于我们。}\BibleRef{路 10:17} 主却对他们说：\BibleQuote{不要因为恶魔屈服于你们而欢喜，而要欢喜，因为你们的名字已登录在天上。}\BibleRef{路 10:20} 祂愿意给谁就给谁驱魔的能力，愿意给谁就给谁复活死人的能力。这样的奇迹甚至在主降生成人之前就已经发生；死人复活，麻风病人被洁净；我们读到过这些事。是谁当时做了这些事？是后来按肉身是达味后裔的基督，但按神性在亚巴郎之前就是天主的基督。祂赐予了这一切的能力，祂藉着人行了这些事；但祂没有把这能力赐给所有人。那些没有得到这能力的人，难道应该绝望，说自己与祂无分，因为他们没有蒙恩领受这些恩赐？身体上有不同的肢体：这个肢体能做一件事，那个能做另一件。天主将身体组合在一起，祂没有给耳朵看，也没有给眼睛听，没有给前额闻，没有给手尝；祂没有给它们这些功能；但所有肢体祂都给了健全，给了结合，给了合一，藉祂的圣神同样地赋予生命并联合一切。同样，在这里，祂没有给有些人复活死人的能力，没有给另一些人辩论的能力；但祂给了所有人什么？\BibleQuote{你们跟我学，因为我是良善心谦的。}\BibleRef{玛 11:29} 因为我们听到祂说：\BibleQuote{我是良善心谦的；}\BibleRef{玛 11:29} 看，我的弟兄们，这就是我们的全部药方：\BibleQuote{你们跟我学，因为我是良善心谦的。}\BibleRef{玛 11:29} 如果一个人行奇迹，却是骄傲的，不是良善心谦的，这对他有什么益处？他难道不会在最后一天被算在那些前来的人中间，说：\BibleQuote{我们不是因你的名说过预言，因你的名行过许多奇迹吗？}\BibleRef{玛 7:22} 但他们会听到什么？\BibleQuote{我从来不认识你们，你们这些作恶的人，离开我走吧！}\BibleRef{玛 7:23}\par

8. 那么，我们学这个有什么益处？祂说：\BibleQuote{因为我是良善心谦的。}\BibleRef{玛 11:29} 祂培植爱德，那最真诚的爱德，没有混乱，没有膨胀，没有高傲，没有欺骗；那位说\BibleQuote{你们跟我学，因为我是良善心谦的}\BibleRef{玛 11:29}的，培植了这爱德。一个骄傲自大的人怎能拥有真诚的爱德？他必定嫉妒。或许有人认为嫉妒的人也会有爱，而我们是错的？愿天主阻止任何人如此错误，说嫉妒的人有爱德。那么，宗徒怎么说？\BibleQuote{爱德不嫉妒。}\BibleRef{格前 13:4} 为什么不嫉妒？\BibleQuote{爱德不夸张，} 他立刻附加了原因，解释了为什么爱德不嫉妒。因为不夸张，所以不嫉妒。确实，他先说：\BibleQuote{爱德不嫉妒；} 但好像你问：\Quote{为什么不嫉妒？} 他加上：\BibleQuote{不夸张。} 如果爱德嫉妒是因为夸张；如果不夸张，就不嫉妒。如果爱德不夸张，因此不嫉妒；那么，那位说\BibleQuote{你们跟我学，因为我是良善心谦的}\BibleRef{玛 11:29}的，就培植了爱德。\par

9. 那么，无论人有什么，任凭他以什么自夸。\BibleQuote{我若能说人间的语言，和天使的语言，但我若没有爱德，我就成了个发声的锣，或发响的钹。}还有什么比各种方言的恩赐更崇高的呢？你若把爱德取走，它就成了\Quote{锣}、成了\Quote{发响的钹}。再听其他的恩赐：\BibleQuote{我若能明白一切奥秘。}有什么比这更卓越？有什么比这更伟大？再听另一项：\BibleQuote{我若有全先知之恩，和一切信德，甚至能移山，但我若没有爱德，我什么也不算。}祂提到更伟大的事，各位弟兄。祂还说了什么？\BibleQuote{我若把我所有的财产全施舍给穷人。}还有什么更完美的事可做呢？的确，主曾为成全之故命令那富贵少年说：\BibleQuote{你若愿意是成全的，去，变卖你所有的，施舍给穷人。}那么，他因为变卖了一切并施舍给穷人，就立刻成全了吗？不；因此祂又加上：\BibleQuote{并且来，跟随我。}\BibleQuote{变卖一切，}祂说，\BibleQuote{施舍给穷人，并且来追随我。}\Quote{我为什么该跟随祢？如今我变卖了一切，施舍给了穷人，难道我还不成全吗？我何必还需要跟随祢？}\BibleQuote{跟随我，}好使你学习\BibleQuote{我是良善心谦的。}为什么？一个人若还不是良善心谦的，能变卖一切施舍给穷人吗？他确实能。\BibleQuote{因为我若把我所有的财产全施舍给穷人。}再继续听。因为有些人，曾舍弃一切并已跟随了主，但尚未完全跟随祂（因为完全跟随祂就是效法祂），却不能忍受苦难的考验。伯多禄，各位弟兄，已是那些舍弃一切并跟随了主的人之一。因为正如那富贵少年忧愁地离去，门徒们感到困扰，问如何能成全，主安慰了他们，他们对主说：\BibleQuote{看，我们舍弃了一切，跟随了祢；那么，我们将得到什么呢？}主告诉他们祂今世将赐给他们什么，将来为他们保留什么。当时伯多禄已是这样做了的人之一。但到了苦难的关头，他竟在一个使女的问话下三次否认了祂，而他曾承诺情愿与祂同死。\par

10. 那么，亲爱的诸位，要留心：\BibleQuote{去，}祂说，\BibleQuote{变卖你所有的，施舍给穷人，你必有宝藏在天上，并且来跟随我。}伯多禄已成全了，如今主在天上坐在圣父的右边，那时他才达到了成全和成熟。因为当他跟随主受难时，他尚未成全；但当他在地上再也没有可跟随的人时，他才被成全。但你们却常有那可跟随的主；主在地上立了一个榜样，当祂把福音留给你们时，在福音中祂与你们同在。因为祂并非说谎，祂曾说：\BibleQuote{看，我同你们天天在一起，直到今世的终结。}因此，跟随主。什么是\Quote{跟随主}？效法主。什么是\Quote{效法主}？\BibleQuote{你们跟我学习，因为我是良善心谦的。}因为如果我把我所有的财产全施舍给穷人，并交出我的身体被火烧，但若没有爱德，于我毫无益处。因此，我劝勉你们的爱德；现在我劝勉爱德，但并非没有一些爱德。我劝勉，使那已开始的得以圆满；并祈祷，使那已开始的得以成全。我也恳求你们为我献上这祈祷，使我所劝告的也在内成全。因为我们如今都不成全，但将在那里成全，在那里一切都是成全的。圣保禄宗徒说：\BibleQuote{弟兄们，我并不认为我已经获得了。}他说：\BibleQuote{这不是说我已经达到，或已成了成全的。}那么，谁敢自夸成全呢？我们宁可承认自己的不完全，好能获得成全。\par

\SourceRecord{Translated by R.G. MacMullen.；Nicene and Post-Nicene Fathers, First Series；Vol. 6.；Edited by Philip Schaff.；Buffalo, NY: Christian Literature Publishing Co.,；1888.；<http://www.newadvent.org/fathers/160392.htm>.}

% structure: p0001 source=h1 target=section reason=page-title

\section{讲道九十三：论新约}

[CXLIII. Ben.]\par

\Emph{论福音之言}，\BibleRef{若 16:7}，\Emph{\BibleQuote{我告诉你们真理；我去对你们有益，}等等}\par

1. 灵魂所有创伤的药方，以及人类过犯的唯一赎罪祭，就是信基督；没有人能得洁净，无论是从原罪——他从亚当所承袭，在亚当内众人犯了罪，且按本性成了义怒之子——还是从他们自己所添加的罪，即没有抵抗肉身的贪欲，反而随从并侍奉它行了不洁和有害的事：除非藉着信德，他们被结合并连接到祂的身体内，祂是由童贞女所生，未曾受任何肉身的诱惑和致命的快乐，祂的母亲也在无罪的胎中孕育了祂，祂\BibleQuote{没有犯过罪，口中也找不到诡诈}。凡信祂的人，成为天主的子女，因为他们由天主而生，藉着义子之恩，即藉着我们主耶稣基督的信德。因此，亲爱的诸位，那位主和我们的救主有充分的理由只提及这一罪，就是圣神要指证世界的，即他们不信祂。祂说：\BibleQuote{我告诉你们真理：}祂说，\BibleQuote{我去为你们有益。因为我若不去，护慰者就不会到你们这里来；我若去了，我就派遣祂到你们这里来。当祂来到时，祂要指证世界关于罪、关于义德、关于审判。关于罪，因为他们不信我。关于义德，因为我往父那里去，你们将不再看见我。关于审判，因为这世界的首领已经受了审判。}\par

2. 因此，祂要世界被指证的，就是这唯一的罪，即他们不信祂；也就是说，因为藉着信祂，所有的罪都得释放，祂要这罪被归咎，而其余的罪就被捆绑。又因为藉着信，他们由天主而生，成为天主的子女；\BibleQuote{因为，}他说，\BibleQuote{祂赐给他们权能，成为天主的子女，就是那些信祂名的人。}所以，谁若信天主子，就他依附于祂，并因义子之恩也成为天主的儿子和继承人，与基督同为继承人而言，他就不会犯罪。因此若望说：\BibleQuote{凡由天主生的，就不犯罪。}所以，世界被指证的罪，就是他们不信祂。这也是祂所说的罪：\BibleQuote{我若没有来，他们就没有罪。}什么？难道他们没有无数其他的罪吗？但藉着祂的来临，这唯一的罪加给了那些不信的人，使得其余的罪被保留。而在那些信的人身上，因为这罪不存在，就使得一切罪都被赦免。保禄宗徒也是基于此说：\BibleQuote{众人都犯了罪，缺乏天主的荣耀}，为叫\BibleQuote{凡信祂的人，不致蒙羞}；正如圣咏也说：\BibleQuote{你们来就近祂，并受光照，你们的面容不会羞惭。}所以，凡在自己内夸耀的，必受羞惭；因为他必被发现有罪。因此，只有那在主内夸耀的，才不致蒙羞。\BibleQuote{因为众人都犯了罪，缺乏天主的荣耀。}因此，当他论及犹太人的不信时，他没有说：\BibleQuote{如果有些人犯了罪，难道他们的罪会使天主的信实失效吗？}他怎能说\BibleQuote{如果有些人犯了罪}呢？因为他自己说：\BibleQuote{众人都犯了罪。}但他却说：\BibleQuote{如果有些人不信，难道他们的不信会使天主的信实失效吗？}为更明确地指出这罪，唯独藉着这罪，其余的罪才被关闭，不能因天主的恩宠获得赦免。关于这唯一的罪，圣神降临——即祂恩宠的恩赐，赐给信者——世界被指证，如主的话：\BibleQuote{关于罪，因为他们不信我。}\par

3. 如果主总是以祂复活的身体显现于人的眼前，那么信祂就不会有伟大的功德和荣耀的福乐了。所以，圣神将这份大恩赐给了那些应该信的人，使他们虽不能用肉眼看见祂，却能用戒除肉欲、并被属灵的渴望所沉醉的心灵去渴慕祂。因此，当那位门徒——他曾说他除非亲手触摸祂的伤痕，否则决不相信——在触摸了主的身体后，仿佛从梦中醒来喊说：\BibleQuote{我的主，我的天主！}主就对他说：\BibleQuote{你因为看见了我才信；那些没有看见而信的人，是有福的。}这福乐是圣神护慰者带给我们的，即祂从童贞母胎中所取的仆人的形体从肉眼前被移开，使心灵被洁净的眼目得以转向天主性的形体，就是祂在与父同等时所具有的，甚至当祂屈尊降生成人时也是一样；以致被同一圣神充满的宗徒可以说：\BibleQuote{即使我们曾按肉体认识基督，但如今我们不再这样认识祂了。}因为他认识基督的肉体，不是按肉体，而是按圣神，他并不是出于好奇去触摸，而是出于确信的信德，承认祂复活的大能；他不在心里说：\BibleQuote{谁升到天上去？就是把基督带下来；或谁下到深渊？就是把基督从死者中带回来。}而是说：\Quote{但是，}他说\BibleQuote{话离你很近，就在你口里，就是耶稣是主；你若心里相信天主使祂从死者中复活，你必得救。因为心里相信，可使人成义；口里承认，可使人得救。}弟兄们，这些是宗徒的话，是被圣神本身的神圣沉醉所倾注的。\par

4. 既然如此，我们本不可能拥有这种不见而信的福分，除非我们从圣神领受了它；所以有充分的理由说：\BibleQuote{我去为你们是有益的，因为我若不去，护慰者就不会到你们这里来；但我若去了，就要派遣他到你们这里来。}的确，按其天主性，他常与我们同在；但除非他身体离开我们，我们总会按肉体看见他的身体，而从未按属灵的方式相信；藉此信德而成义和有福的我们，得以用洁净的心达到默观那真圣言，天主与天主同在，\BibleQuote{万物是藉着他造的}，并且\BibleQuote{他成了血肉，居住在我们中间。}若不是藉手的触摸，而是\BibleQuote{人心里相信，以致成义}；那么，世界——它只相信它所看见的——被说服我们的义德，就是有充分理由的。如今为使我们拥有那不信的世界应被说服的信德之义，主因此说：\BibleQuote{关于义德，因为我往圣父那里去，你们将不再见到我。}仿佛他说：这将是你们的义德：你们相信我——中保，你们将充分确信他复活了并到了圣父那里，虽然\newline{}你们不再按肉体看见他；藉着他和好，你们才能按神性看见天主。因此他对那代表教会的妇女，在复活后她俯伏在他脚前时说：\BibleQuote{不要摸我，因为我还没有升到圣父那里。}这句话奥秘地理解为：\BibleQuote{不要以肉体的方式藉着身体接触相信我；而要以属灵的方式相信；也就是说，当我升到圣父那里时，你们将以属灵的信德触摸我。}因为\BibleQuote{那些没有看见而相信的，是有福的。}这就是信德之义，世界没有它，却说服我们这些不缺乏它的人；因为\BibleQuote{义人因信德而生活。}那么，无论是作为在他内复活，在他内来到圣父那里，我们无形中并在成义上得以成全；还是作为不见而信，因信德生活，因为\BibleQuote{义人因信德而生活}，他就用这些含义说：\BibleQuote{关于义德，因为我往圣父那里去，你们将不再见到我。}\par

5. 也不要让世界以被魔鬼阻止相信基督为借口而自我开脱。因为对于信徒，今世的首领已被驱逐出去，使他不再在那些基督已藉信德开始拥有的人心中工作；正如他在悖逆之子心中工作一样；他不断煽动他们试探和骚扰义人。因为他已被驱逐出去，那曾内心掌权的，如今从外发动战争。虽然如此，藉着他的迫害，\BibleQuote{主领导温良的人秉公行义}；然而，在他被驱逐这一事实中，他\BibleQuote{已受了审判}。关于这一点，世界被\BibleQuote{审判}说服；因为那不愿相信基督的人抱怨魔鬼是徒劳的，这魔鬼已被审判，即被驱逐，并为了磨练我们而被允许从外面攻击我们，不仅男人，甚至女人、男孩和女孩——殉道者已战胜了他。然而是谁使他们战胜的呢？岂不是他们相信的那位吗？他们虽未看见却爱了他，并因他在他们心中的统治，摆脱了一个最压迫人的主。而这一切都是藉着恩宠，即圣神的恩赐。那么，恰好同一圣神说服世界：关于\BibleQuote{罪}，因为它不信基督；\BibleQuote{关于义德}，因为那些有意志的人已经相信，虽然他们未看见他们所相信的那位；并藉着他的复活，希望他们自己也应在复活中得以成全；\BibleQuote{关于审判}，因为如果他们愿意相信，没有人能阻止他们，\BibleQuote{因为今世的首领已受了审判}。\par

\SourceRecord{Translated by R.G. MacMullen.；Nicene and Post-Nicene Fathers, First Series；Vol. 6.；Edited by Philip Schaff.；Buffalo, NY: Christian Literature Publishing Co.,；1888.；<http://www.newadvent.org/fathers/160393.htm>.}

% structure: p0001 source=h1 target=section reason=page-title

\section{\WorkTitle{新约讲道集第九十四篇}}

\Emph{[CXLIV. Ben.]}\par

\Emph{关于同一福音之言}，\BibleRef{若 16:8}，\Emph{，\Quote{祂要指证世界关于罪、关于义、关于审判。}}\par

1. 当我们的主、救主耶稣基督详细谈论圣神的来临时，祂在其它话中说道：\Quote{祂要指证世界关于罪、关于义、关于审判。}祂说了这话之后，并没有转论别的题目，而是恩赐对这同一真理更为明确的说明。\Quote{关于罪，}祂说，\Quote{是因为他们不信我。关于义，是因为我往父那里去。关于审判，是因为这世界的元首已经受到了审判。}因此在我们心中生起了一种理解的渴望：为何仿佛人的唯一之罪是不信基督，祂却只说圣神要指证世界关于这罪；但若除了这不信之外，人还有众多其它的罪，为何圣神要指证世界唯独关于这罪？难道是因为所有的罪都因不信而被保留，因信德而被赦免，因此天主特将这罪置于其余之上，因着这罪，只要骄傲的人不信屈尊的天主，其余的罪便不得解脱？因为经上记载：\Quote{天主拒绝骄傲的人，却赐恩宠给谦卑的人。}如今，天主的这个恩宠是天主的恩赐。但最大的恩赐就是圣神本身；因此才称为恩宠。因为\Quote{所有的人都犯了罪，需要天主的荣耀；因为罪藉着一个人进入了世界，死亡藉着罪进入，众人都犯了罪；}因此这是恩宠，因为是白白赐予的。所以白白的赐予，因为不是按照严格的功过作为报酬偿还，而是在赦罪之后作为恩赐赐予。\par

2. 因此，不信的人，即世界的爱慕者，被指证为有罪；因为用\Quote{世界}之名来指代他们。因为经上说\Quote{祂要指证世界关于罪}，所指的不是别的罪，而是他们没有信基督。因为这罪若不存在，就没有罪会存留，因为义人靠信德生活时，一切都得赦免。如今，相信耶稣是基督，与相信基督，两者有很大区别。因为连魔鬼也相信耶稣是基督，然而魔鬼并不相信基督。因为相信基督的人，是既在基督内怀有望德又爱基督的人。如果他只有信德而没有望德和爱德，他相信基督存在，但并未相信基督。因此，凡相信基督的人，藉着相信基督，基督就来到他那里，并在某种程度上将自己与他结合，使他成为基督身体上的一个肢体。这只能藉着望德和爱德的加入才能实现。\par

3. 祂的话又是什么意思：\Quote{关于义，是因为我往父那里去}？首先必须探究：如果世界被指证关于罪，为何也被指证关于义？因为谁能被正确地指证为义呢？难道世界是被指证关于自己的罪，而关于基督的义吗？我看不出还能理解别的，因为祂说：\Quote{关于罪，是因为他们不信我。关于义，是因为我往父那里去。}他们不信，祂往父那里去。因此，这是他们的罪，并祂的义。但祂为何唯独在祂往父那里去这件事上称为义呢？难道祂从父那里来不是义吗？或者，那从父那里来到我们这里更是仁慈，而往父那里去才是义？\par

4. 所以，诸位弟兄，我认为在这如此深奥的圣经文句中，在那些或许隐藏着某种真理、可在适当时候揭示出来的话语中，我们应当一起忠实地探究，以便能健康地寻得答案。那么，为何祂称往父那里去为义，却不也称从父那里来为义呢？是否因为祂的来是仁慈，因此祂的去是义？为的是在我们自己身上也能学到：如果我们不肯先给仁慈留出地方，\Quote{不求自己的益处，也求别人的益处，}那么义在我们内就不能圆满。当宗徒给了这个劝告之后，他立刻附上了主自己的榜样：\Quote{不要做任何出于纷争或虚荣的事；却要存心谦卑，各人看别人比自己强。各人不要单顾自己的事，也要顾别人的事。}然后他立刻加上：\Quote{你们当以基督耶稣的心为心；祂虽具有天主的形体，却不以自己与天主同等为僭夺；反而空虚自己，取了奴仆的形体，成为人的模样；既显为人形，就卑抑自己，听命至死，且死在十字架上。}这就是祂从父那里来的仁慈。那么祂往父那里去的义是什么呢？祂继续说：\Quote{因此，天主极其举扬祂，赐给祂一个名号，超越一切名号；使天上、地上和地下的众膝，都因耶稣的名而屈膝，并且众口都要承认耶稣基督是主，归于天主父的光荣。}这就是祂往父那里去的义。\par

5. 但若唯祂独自往父那里去，这对我们有何益处？为何圣神能凭此\OriginalTerm{义德}{righteousness}说服世界？然而，若祂并非独自往父那里去，祂就不会在另一处说：\BibleQuote{沒有人上過天，除了那自天降下的人子。}但宗徒\Person{圣保禄}也说：\BibleQuote{我們的家乡是在天上。}这是为什么呢？因为他也说：\BibleQuote{你们既然与基督一同复活了，就该追求天上的事，在那里基督坐在天主的右边。你们要思念天上的事，不要思念地上的事。因为你们已经死了，你们的生命与基督一同藏在天主内。}那么祂怎么是独自的呢？难道祂独自是因为基督与祂所有的肢体是一体的，正如头与身体一样？那么祂的身体是什么，不就是教会吗？正如同一导师所说：\BibleQuote{你们便是基督的身体，各自做肢体。}因为我们既已堕落，祂为我们而降，那么所谓\BibleQuote{沒有人上過天，除了那下来的}是什么意思？不就是没有人上过天，除非与他结合为一，并作为肢体紧紧连于他那降下来的身体吗？因此祂对门徒说：\BibleQuote{你们离了我，什么也不能做。}因为祂与父为一是一种方式，与我们为一则是另一种方式。祂与父为一，是因为父与子的性体是同一的；祂与父为一，是因为\BibleQuote{祂虽具有天主的形体，却没有将与自己同等视为应当夺取的。}但祂与我们成为一体，是因为\BibleQuote{祂空虚了自己，取了奴仆的形体。}祂依\OriginalTerm{亚巴郎的后裔}{Seed of Abraham}与我们成为一体，\BibleQuote{万民都要因他蒙受祝福。}那地方当宗徒引证后，他说：\BibleQuote{不是说\InnerQuote{后裔们}，如同为多数；而是如同为单一：\InnerQuote{你的后裔}，就是基督。}而且我们也属于那\Person{基督}，通过我们与头的结合和联络，就成了一个基督。也因为他对我们说：\BibleQuote{所以你们是亚巴郎的后裔，按照恩许成为承继人。}因为如果亚巴郎的后裔只有一个，而且那唯一的亚巴郎后裔只能理解为基督，但我们也是亚巴郎的后裔，所以这个整体，即\OriginalTerm{头与身体}{Head and Body}，就是一个基督。\par

6. 因此我们不应自以为与那义德分离，就是主亲自提到的义德，说：\BibleQuote{关于义德，因为我往父那里去。}因为我们也与基督一同复活了，我们与我们的头基督在一起，现世暂时借着信德和望德；但我们的望德将在末日复活时圆满。当我们的望德圆满时，我们的\OriginalTerm{成义}{justification}也将圆满。那要使其圆满的主，在祂自己的肉身上（即我们的头上）指示了我们，祂在那里复活并升到父那里，我们应当希望什么。因为经上记载：\BibleQuote{祂为了我们的过犯被交付，又为了我们的成义而复活。}于是，世界在那些不信基督的人身上为\BibleQuote{关于罪}而被说服；在那些在基督肢体中复活的人身上则为\BibleQuote{关于义德}而被说服。因此经上说：\BibleQuote{使我们在祂内成为天主的义德。}因为若不在祂内，就绝不是义德。但若在祂内，祂就带着我们全体往父那里去，这完美的义德将在我们内实现。因此，世界也为审判而被说服，\BibleQuote{关于审判}，\BibleQuote{因为这今世的首领已经受了审判。}就是说，魔鬼，\OriginalTerm{今世的首领}{prince of this world}，不义者的首领，那些只在心中居住于他们所爱的这世界的人，因此被称为\Quote{世界}；正如我们的家乡在天上，如果我们已经与基督一同复活了。因此，正如基督与我们，即祂的身体，是一体；同样，魔鬼与所有不敬虔的人，他是他们的头，仿佛是他自己的身体，也是一体。因此，正如我们不与那义德分离，就是主所说的\BibleQuote{因为我往父那里去}的义德；同样，不敬虔的人也不与那审判分离，就是祂所说的\BibleQuote{因为这今世的首领已经受了审判}。\par

\SourceRecord{Translated by R.G. MacMullen.；Nicene and Post-Nicene Fathers, First Series；Vol. 6.；Edited by Philip Schaff.；Buffalo, NY: Christian Literature Publishing Co.,；1888.；<http://www.newadvent.org/fathers/160394.htm>.}

% structure: p0001 source=h1 target=section reason=page-title

\section{论新约讲道九十五}

\EditorialNote{[CXLV. Ben.]}\par

\Emph{论福音之言}，\BibleRef{若 16:24} \Emph{，\InnerQuote{到现在你们还没有因我的名求什么；}以及论} \BibleRef{路 10:17} \Emph{，\InnerQuote{主啊，因祢的名，连魔鬼都屈服于我们。}}\par

1. 当圣福音被宣读时，我们听到了那理应立刻激励每个热忱的灵魂去寻求，而不是懈怠的道理。因为凡不被激动者，便不会改变。但有一种危险的变动，关于它经上写道：\Quote{不要让我的脚动摇。}然而，另有一种变动的，是那寻求、叩门、祈求的人。刚才所读的，我们都听见了；但我猜想并非都理解了。它提到了那件事，你们应当与我一同寻求，与我一同祈求，为要获得它而与我一同叩门。因为我希望主的恩宠与我们同在，使我虽愿为你们服务，也配得领受。请问，我们刚刚听到主对祂的门徒说了什么？\Quote{到现在你们还没有因我的名求什么。}这不是对那些门徒说的吗？那些门徒，主派遣了他们，赐给他们权柄宣讲福音、行奇迹，他们欢欣归来，对祂说：\Quote{主啊，因祢的名，连魔鬼都屈服于我们。}你们认得、记得我引用的福音这话，它在每一段、每一句中都说真理，绝无虚假，绝不欺骗。那么，如何又是真的：\Quote{到现在你们还没有因我的名求什么}？以及\Quote{主啊，因祢的名，连魔鬼都屈服于我们}？这无疑激励人寻求这个难题的奥秘。因此，我们求、寻、叩门。愿我们内有忠实的虔敬，不是肉体的不安，而是心灵的顺服，使那看见我们叩门者为我们开门。\par

2. 那么，主可能赐予、用以侍奉你们的事，你们要怀着热切的关注，即如饥饿般领受；当我讲论时，你们无疑会以健全的品味，赞同那从主的仓库中摆在你们面前的食物。主耶稣知道人的灵魂，即按天主的肖像所造的理性心智，能因何满足：唯独因祂自己。祂知道这一点，也知道那灵魂尚未获得这圆满。祂知道祂是显明的，也知道祂是隐藏的。祂知道祂里面展示的是什么，隐藏的是什么。祂知道这一切。圣咏说：\Quote{上主，祢的慈爱何其伟大！祢为敬畏祢的人所珍藏的，为投靠祢的人所施行的！}\Quote{祢的慈爱}既伟大且多样，\Quote{祢为敬畏祢的人所珍藏。}若祢为敬畏祢的人珍藏，那么祢为谁开启呢？\Quote{祢为投靠祢的人所施行。}由此产生了一个双重的疑问，但每一个疑问都可以由另一个解决。若有人要问另一个：什么叫\Quote{祢为敬畏祢的人所珍藏，为投靠祢的人所施行}？敬畏的人与投靠的人不同吗？难道不是同一批人既敬畏天主又投靠天主吗？谁投靠那不敬畏祂的呢？谁以虔敬的方式敬畏祂，却没有投靠祂呢？那么，先解决这个。我想就那些投靠的人和那些敬畏的人说些什么。\par

3. 法律有恐惧，恩宠有望德。然而法律与恩宠有何区别，既然法律与恩宠同出一赐予者？法律警告那依靠自己的人，恩宠帮助那信靠天主的人。我说，法律警告；不要因其简短而轻视；仔细权衡，它至关重要。好好看一看我所说的，接受我们所提供的，证明我们从何处取得。法律警告那依靠自己的人，恩宠帮助那信靠天主的人。法律说什么？许多事；谁能枚举？我从其中提出一个微小而简短的诫命，就是宗徒所引用的，一个非常微小的；让我们看看谁能胜任它。\BibleQuote{不可贪恋。} 这是什么，弟兄们？我们听到了法律；若没有恩宠，你们听到了你们的惩罚。无论你是谁，听到这事而依靠自己的，你为何向我夸耀？为何向我夸耀无罪？为何因此自夸？你可以说，\Quote{我没有抢夺别人的财物；} 我听见，我相信，也许我还看见，你没有抢夺别人的财物。你听到了，\BibleQuote{不可贪恋。} \Quote{我不进别人的妻子；} 这也是我听见、相信、看见的。你听到了，\BibleQuote{不可贪恋。} 你为何全面检查自己外面，却不检查里面？往里看，你将看见你肢体中的另一条法律。往里看，你为何越过自己？进入你自己。你将\BibleQuote{看见你肢体中另一条法律，反抗你心神的法律，并将你掳去，使你隶属于那在你肢体内的罪的法律。} 那么，天主的甘饴对你隐藏是理所当然的。那在你肢体内的法律，反抗你心神的法律，将你掳去。那对你隐藏的甘饴，圣天使们饮用；你这被掳的人，不能饮用和品尝那甘饴。\BibleQuote{你原不认识贪欲，除非法律说：不可贪恋。} 你听见了，害怕了，试图战斗，却不能得胜。因为\BibleQuote{罪借着诫命得了机会，就产生了死亡。} 你们当然认得这些话，它们都是宗徒的话。\BibleQuote{罪借着诫命得了机会，就在我身上发起了各样的贪欲。} 你为何在骄傲中自夸？看，敌人用你自己的武器征服了你。你确实曾期望诫命作为防御；然而看，敌人却借着诫命找到了进入的机会。\BibleQuote{因为罪借着诫命得了机会，} 他说，\BibleQuote{欺骗了我，并且借着它杀了我。} 我所说的\Quote{敌人用你自己的武器征服了你} 是什么意思？请听同一宗徒接着说：\BibleQuote{所以法律是圣的，诫命也是圣的、公义的、善的。} 现在请答复那诋毁法律的人；凭着宗徒的权威回答：\BibleQuote{诫命是圣的，法律是圣的，诫命是公义而善的。那么，那善的竟成了我的死亡吗？绝对不可！但是罪，为显出它是罪，竟借着那善的在我身上产生了死亡。} 这是为何？岂不因为你接受诫命时，不是爱，而是恐惧？你害怕惩罚，你不爱正义。凡害怕惩罚的人，若可能，他愿意做自己所喜欢的事，而不愿遭受他所害怕的。天主禁止奸淫，你贪恋了别人的妻子，你不进到她那里，你没有这样做；机会给了你，你有时间，有利的地方开放，证人缺席，然而你没有做，为什么？因为你害怕惩罚。但没有人会知道。难道天主不知道吗？所以显然，因为天主知道你要做的事，你没有做；但在这里你害怕的是天主的威胁，而不是爱他的诫命。你为什么不做？因为如果你做了，你就会投入地狱之火。你所怕的是火。啊！如果你爱贞洁，即使全无惩罚，你也不会做。如果天主对你说：\Quote{看，去做吧，我不定你的罪，我不把你投入地狱之火，但我要向你掩面。} 如果你因这威胁而不做，那便是出于对天主的爱而不做，而非出于对审判的恐惧。但你会做的，也许我的意思是你会这样做；因为这不是我的位置来判断。如果你不这样做，是因为你憎恶奸淫的污秽，因为你爱他的诫命，为要获得他的许诺，而不是因为害怕他的定罪，那么，那使你成为圣人的恩宠在帮助你；一切都是恩宠，不要归功于你自己，不要归功于你的力量。你出于喜悦而行，好；你出于爱德而行，好；我赞同，我同意。当你按你的意愿行事时，爱德借着你在工作。你立刻就会尝到甘饴，只要你寄望于主。\par

4. 但是，你从何处得到这爱德，如果你还有的话？因为我恐怕，你至今仍因恐惧而不做，恐怕你在自己眼中显得伟大。现在，如果你因爱德而不做，你真是伟大的。你有爱德吗？你说：\Quote{我有。} 从何处？\Quote{从我自身。} 如果你从自身得到它，你就远离了甘饴。你会爱你自己，因为你将爱那来源。但我要证明你没有。因为你认为你从自己得到了如此伟大的东西，正因如此，我不相信你有。因为如果你有，你就会知道从何处得到。你从自己得到爱德，仿佛它是某种微小之物？\BibleQuote{如果你们能说人间的语言，并能说天使的语言，却没有爱德，你们就成了发声的锣、响的钹。如果你们能知悉一切奥秘，并拥有一切知识、一切预言、一切信心，甚至能移山，却没有爱德，这些对你们毫无益处。如果你们把全部财产分给穷人，并把自己的身体交给火焚，却没有爱德，你们仍然一无所是。} 这爱德何等伟大，若缺少它，一切皆无益处！不要把它与你的信德、你的知识、你的语言恩赐相比，与较小的东西，与你的眼、手、脚、肚腹，与任何最低的肢体相比，这些最小的东西能与爱德相比吗？眼和鼻你从天主领受，而爱德你从自己领受吗？如果你给了自己超越万物的爱德，你就把天主视为无足轻重。天主还能给你什么？无论祂给了什么，都更小。你给自己的爱德超越一切。但如果你有，你就不是自己给的。\BibleQuote{你有什么不是领受的呢？} 谁给了我，谁给了你？天主。在祂的恩赐中认识祂，以免你感受祂的惩罚。因信圣经，天主给了你爱德，一个伟大的恩赐，爱德，超越一切。天主给了你，\BibleQuote{因为天主的爱德已倾注在我们心中，} 也许是靠你自己？决不是；\BibleQuote{借着赐给我们的圣神。}\par

5. 与我一同回到那个俘虏那里，回到我的命题。\Quote{法律使依靠自己的人惊恐，恩宠却帮助信赖天主的人。} 请看那个俘虏。\BibleQuote{他看见在自己的肢体中，另有一种法律，反抗他理智的法律，并把他掳去，叫他隶属于那在他肢体内的罪恶的法律。} 看，他被捆绑，看，他被拖走，看，他被俘虏，看，他被制服。这对他有什么益处呢，\BibleQuote{不可贪恋}？他听到了，\BibleQuote{不可贪恋，} 好使他认识他的仇敌，而不是克服他。\BibleQuote{因为如果不是法律说：\BibleQuote{不可贪恋}，他就不知道私欲偏情。} 现在你看到了仇敌，战斗吧，释放你自己，维护你的自由，压制快感的诱惑，彻底消灭非法的欢乐。武装自己，你有法律，前进，如果你能，就征服。你已经借着天主恩宠的一小部分，\BibleQuote{按内心的人，喜爱天主的法律；但你也看见在你的肢体中，另有一种法律，反抗你理智的法律，} 不是 \Quote{反抗} 却毫无能力，而是 \BibleQuote{把你掳去，隶属于罪恶的法律。} 看，从何处，对于你害怕 \Quote{充满甘饴隐藏了} 的人，对于他害怕 \Quote{被隐藏}，对于他 \Quote{信赖} 的人，如何实现？在你的仇敌之下呼喊，因为你有攻击者，你也有助佑者，祂看着你战斗，在困难中帮助你；但只有祂发现你 \Quote{信赖}；因为祂憎恨骄傲者。那么，你会在这仇敌之下呼喊什么？\BibleQuote{我真苦啊！} 你已经看到了，因为你已经呼喊了。让你的呼喊这样，当你偶然在仇敌之下受困时，你说，在你的内心深处说，以纯正的信德说，\BibleQuote{我真苦啊！} 我真苦！\BibleQuote{因此我苦，} 因为 \Quote{我。} \BibleQuote{我真苦啊，} 既因为 \Quote{我，} 也因为 \Quote{人。} 因为 \Quote{他白白地受惊扰。} 虽然 \Quote{人按肖像行走，} 然而，\BibleQuote{我真苦啊！谁能救我脱离这该死的肉身呢？天主的恩宠，藉着我们的主耶稣基督。} 你自己吗？你的力量在哪里？你的自信在哪里？当然，你既呼喊又沉默；沉默，即是从自我夸耀中沉默，而不是从呼求天主中沉默。沉默，却要呼喊。因为天主自己也既沉默又大声呼喊；祂在审判中沉默，祂在诫命中不沉默；你们也当在得意中沉默，不在呼求中沉默；免得天主对你说：\Quote{我沉默，难道我会永远沉默吗？} 因此呼喊吧，\BibleQuote{我真苦啊！} 承认自己被征服，让你的力量羞愧，并说：\BibleQuote{我真苦啊！谁能救我脱离这该死的肉身呢？天主的恩宠，藉着我们的主耶稣基督。} 我上面说了什么？法律使依靠自己的人惊恐。看，人依靠自己，他试图战斗，不能取胜，被征服，被击倒，被制服，被俘虏。他学会了依靠天主，而法律曾使依靠自己的人惊恐，现在恩宠却帮助信赖天主的人。在这种信赖中，他说：\BibleQuote{谁能救我脱离这该死的肉身呢？天主的恩宠，藉着我们的主耶稣基督。} 现在看那甘饴，品尝它，赏味它；听\BibleBook{}{圣咏}：\BibleQuote{请你们体验并观看，上主是何等甘饴。} 祂已成了你的甘饴，因为祂拯救了你。当你依靠自己时，你对自己是苦涩的。畅饮甘饴，领受如此丰盛的预尝。\par

6. 那时主耶稣基督的门徒还在法律之下，仍需要被洁净，仍需要被滋养，仍需要被纠正，仍需要被引导。因为他们仍有私欲偏情；而法律却说：\Quote{你不可贪恋。} 请勿冒犯那些圣洁的公羊，羊群的领袖，请勿冒犯他们，我要说，因为我讲真理：福音记载，他们争论谁最大，而且当主还在世上时，他们因争权而骚动。这来自何处，岂非来自旧酵母？岂非来自肢体中的法律，反抗心灵的法律？他们寻求高位，是的，他们渴望它；他们想着谁该最大；因此他们的骄傲被一个小孩子羞辱。耶稣召叫谦卑的年龄来驯服膨胀的欲望。因此当他们回来时，也很有理由说：\Quote{主，请看，因祢的名号，连恶魔也屈服于我们。}（他们为虚无的事而欢喜；与天主所应许的相比，这算什么呢？）主，善师，平息恐惧，建立坚固的支撑，对他们说：\Quote{不要因恶魔屈服于你们而欢喜。} 为什么？因为\Quote{将来有许多人奉我的名来，说：看，我们因祢的名号赶走了魔鬼；但我要对他们说：我不认识你们。不要为此欢喜，而要因你们的名字已登记在天上而欢喜。} 你们还不能在那里，但你们已经登记在那里。因此\Quote{欢喜吧。} 所以那处经文又说：\Quote{直到现在，你们没有因我的名求什么。} 因为你们所求的，与我所愿意赐给的相比，算不得什么。你们因我的名求了什么？叫恶魔屈服于你们？\Quote{不要为此欢喜，} 就是说，你们所求的是虚无；因为如果它是任何东西，主就会叫他们欢喜。所以它并非绝对虚无，而是与天主恩赐的伟大相比是渺小的。因为宗徒保禄并非什么都不是；然而与天主相比，\Quote{栽种的不算什么，浇灌的也不算什么。} 因此我对你们说，也对我自己说，既对我自己又对你们说，当我们因基督的名求这些现世之物时，你们无疑是在求。因为谁不求呢？有人求健康，如果他病了；有人求释放，如果他在监里；有人求港口，如果他在海上颠簸；有人求胜利，如果他在与敌人交战中；而且他们因基督的名求一切，但所求的是虚无。那么该求什么呢？\Quote{因我的名求吧。} 他没有说什么，但通过这些话语我们明白该求什么。\Quote{求吧，你们必会得到，好使你们的喜乐得以圆满。求吧，你们必会得到，因我的名。} 但是什么？不是虚无；而是什么？\Quote{好使你们的喜乐得以圆满；} 也就是说，求那能令你们满足的。因为当你们求现世之物时，你们求的是虚无。\Quote{凡喝这水的，还要再渴。} 他将欲望的水桶放下井中，打水来喝，只是为再次口渴。\Quote{求，好使你们的喜乐得以圆满；} 就是说，叫你们得以满足，不只是暂时感到欢愉。求那能令你们满足的；说斐理伯的话：\Quote{主！把父显示给我们，我们就心满意足了。} 主对你们说：\Quote{斐理伯！这么长久的时间，我和你们在一起，而你还不认识我吗？谁看见了我，就是看见了父。} 因此，要感谢基督，祂为软弱的你们成为了软弱的，并预备你们的渴望去满足于基督的天主性。让我们转向主，等等。\par

\SourceRecord{Translated by R.G. MacMullen.；Nicene and Post-Nicene Fathers, First Series；Vol. 6.；Edited by Philip Schaff.；Buffalo, NY: Christian Literature Publishing Co.,；1888.；<http://www.newadvent.org/fathers/160395.htm>.}

% structure: p0001 source=h1 target=section reason=page-title

\section{论新约·讲道第九十六篇}

\Emph{[CXLVI. Ben.]}\par

\Emph{论\WorkTitle{福音}中的话}，\BibleRef{若 21:16} \Emph{，\BibleQuote{西满，若望的儿子，你爱我吗？}等等。}\par

1. 诸位亲爱的，你们已经注意到，在今天的读经中，主以提问的方式对伯多禄说：\BibleQuote{你爱我吗？}他回答说：\BibleQuote{主，你知道我爱你。}这样重复了第二次、第三次；每一次主都说：\BibleQuote{喂养我的羔羊。}基督将祂的羔羊托付给伯多禄喂养，而这位伯多禄也曾被基督喂养。伯多禄还能为主做什么呢？更何况如今主已拥有不朽的身体，即将升天？好像祂对他说：\BibleQuote{\InnerQuote{你爱我吗？}以此表示你爱我：\InnerQuote{喂养我的羊。}}所以，弟兄们，你们要顺从地聆听，你们是基督的羊；而我们却怀着恐惧聆听喂养我的羊这句话。如果我们怀着恐惧喂养，并为羊群恐惧；那么这些羊又该如何为自己恐惧呢？因此，愿谨慎是我们的份，顺从是你们的份；牧者的警醒是我们的份，羊群的谦卑是你们的份。尽管我们这些似乎从高处向你们讲话的人，也怀着恐惧匍匐在你们脚下；因为我们知道，这个看似崇高的席位将要作出何等危险的汇报。因此，亲爱的诸位，天主教的子民，基督的肢体，要想想你们拥有何等荣耀的头！天主的儿女，要想想你们找到了何等伟大的父亲！基督徒们，要想想应许给你们的是何等丰厚的产业！这产业不像世上那样，儿女只能在父母去世后才能继承。因为世上没有人能在父亲活着时继承父亲的产业，除非父亲死了。但我们却能在我们的父亲活着时继承祂所赐予的；因为我们的父亲不会死。我再加上一句，再说一句，并且说的是真理；我们的父亲本身就是我们的产业。\par

2. 你们要行事端正，尤其是基督的初信者，最近受洗、刚刚重生的人，正如我先前劝诫你们的，如今我也这样说，并表达我的忧虑；因为今天的福音读经使我更加惧怕：你们要谨慎，不要效法那些邪恶的基督徒。不要说\Quote{我要这样做，因为许多信友都这样做。}这并非为灵魂寻求庇护，而是为地狱寻找同伴。你们要在主的禾场上成长；在这里，如果你们自己善良，就会找到善良的人令你们喜悦。难道你们是我们的私产吗？异端者和分裂者从主那里偷窃，将其据为己有，他们要喂养的不是基督的羊群，而是他们自己的羊群，与基督对抗。的确，他们将祂的名号贴在这些掠夺物上，以便他们的抢劫似乎因祂权能的名号而得以维持。当这样的人悔改时，他们在教会之外领受了祂圣洗的印号，基督会做什么？祂赶走掠夺者，并不抹去那印号，而是占据那房屋；因为祂在那里找到了祂的印号。祂何必更改自己的名字？他们是否留意主对伯多禄所说的话：\BibleQuote{喂养我的羔羊，喂养我的羊？}祂是否对他说：\Quote{喂养你的羔羊；}或\Quote{喂养你的羊？}但对于那些被排除在外的人，祂在\WorkTitle{雅歌}中对教会说了什么？新郎对新娘说：\BibleQuote{如果你不认识自己，女子中最美丽的，出去吧。}好像祂说：\Quote{我并不赶你出去，\BibleQuote{如果你不认识自己，女子中最美丽的，出去吧，}如果你不认识自己，在圣经的镜子里，如果你不注意，美丽的女子，那镜子并无虚假的光彩欺骗你；如果你不知道经上论你所说的话：\BibleQuote{你的荣耀将超越全地；}论你所说的话：\BibleQuote{我要将万民赐给你作产业，地极作你的领土。}以及其他无数为天主教会作见证的经文。如果你不知道这些，你就于我无份，你不能使自己成为我的继承人。\BibleQuote{那么，出去跟随羊群的足迹吧}不要与羊群为伴；去喂养你的山羊，而不是像对伯多禄所说的：\BibleQuote{我的羊。}}对伯多禄所说的是：\BibleQuote{我的羊；}对分裂者所说的是：你的山羊。一处是\Quote{羊，}另一处是\Quote{山羊；}一处是\Quote{我的，}另一处是\Quote{你的。}请回想我们审判者的右手和左手；回想山羊将站在哪里，绵羊将站在哪里；那么你就会明白哪里是右边，哪里是左边，白色与黑色，光明与黑暗，美丽与丑陋，那将要承受国度的，和那将要遭受永罚的。\par

\SourceRecord{Translated by R.G. MacMullen.；Nicene and Post-Nicene Fathers, First Series；Vol. 6.；Edited by Philip Schaff.；Buffalo, NY: Christian Literature Publishing Co.,；1888.；<http://www.newadvent.org/fathers/160396.htm>.}

% structure: p0001 source=h1 target=section reason=page-title

\section{新约讲道集 第97篇}

[第147篇·本笃]\par

论同一个福音，即\BibleBook{若望}{福音}\BibleRef{若 21:15}的话：\BibleQuote{西满，若望的儿子，你爱我比这些更深吗？}等。\par

1. 你们记得，宗徒之长伯多禄曾在主的苦难中动摇。他因自己而动摇，却由基督得以更新。因为他先是鲁莽的自负者，后来成了胆怯的否认者。他曾许诺为主而死，而主却要先为他而死。那时他说：\BibleQuote{我愿与你同死}，又说：\BibleQuote{我要为你舍命}。主回答他说：\BibleQuote{你要为我舍命吗？我实实在在告诉你：鸡叫以前，你要三次否认我。}时辰来到，因为基督是天主，伯多禄是人，圣经便应验了：\BibleQuote{我曾在惊慌中说：人人都是说谎者。}宗徒也说：\BibleQuote{天主是真实的，人人都是说谎者。}基督真实，伯多禄说谎。\par

2. 但如今呢？正如你们在读福音时所听到的，主问他，对他说：\BibleQuote{西满，若望的儿子，你爱我比这些更深吗？}他回答说：\BibleQuote{主，是的，你知道我爱你。}主又问了这个问题，并且第三次问他。当他以爱回答时，主就把羊群托付给他。因为每一次主耶稣对伯多禄说：\BibleQuote{我爱你}，主便说：\BibleQuote{你喂养我的羔羊}，喂养我的\BibleQuote{小羊}。在这位伯多禄身上，预表了所有牧者的合一，即好牧者的合一，他们知道是为基督而不是为自己牧养基督的羊。此时伯多禄说谎了吗？还是他回答说自己爱主是不真实的？他回答得真实，因为他回答的是他内心所看见的。而当他曾说\BibleQuote{我要为你舍命}时，他是靠未来的力量自负。现在，每个人在说话时或许知道自己当时是怎样的人；但明天他要成为怎样的人，谁知道呢？因此，当主问他时，伯多禄回视自己的内心，并确信地回答他所看见的：\BibleQuote{主，是的，你知道我爱你。}我告诉你的，你知道；我在这心里所见的，你也看见。然而，他不敢回答主所问的。因为主不是简单地问：\BibleQuote{你爱我吗？}而是加了：\BibleQuote{你爱我比这些更深吗？}即\BibleQuote{你比这些人更爱我吗？}他是指其他门徒；伯多禄只能说\BibleQuote{我爱你}，却不敢说\BibleQuote{比这些更深}。他不愿第二次说谎。为自己内心作证就够了；判断别人的内心不是他的职责。\par

3. 因此伯多禄是真实的；或者更确切地说，基督在伯多禄内是真实的？当主耶稣基督愿意时，他离弃了伯多禄，伯多禄便显为凡人；但当主耶稣基督乐意时，他充满了伯多禄，伯多禄便显为真实。\OriginalTerm{盘石}{Petra}使伯多禄真实，因为盘石就是基督。当伯多禄第三次回答说他爱基督，主第三次将祂的小羊托付给伯多禄时，主向他预告了他的苦难：\BibleQuote{你年少时，自己束上带子，随意往来；但年老时，你要伸出手来，别人要给你束上带子，带你到你不愿意去的地方。}福音作者向我们解释了基督的意思：\BibleQuote{他说这话，是指明伯多禄要以怎样的死来光荣天主。}即他要为基督被钉在十字架上；因为\BibleQuote{你要伸出手来}就是指此。那否认者如今何在？此后，主基督说：\BibleQuote{跟随我。}这不同于他从前召唤门徒时的意义。因为那时他也说\BibleQuote{跟随我}，但那时是为训导，如今是为冠冕。当他否认基督时，难道他不怕死吗？他害怕遭受基督所受的苦。但现在他不再害怕了。因为他看见基督已复活在肉身中，就是他曾在木头上所看见的那位。基督的复活消除了对死亡的恐惧；既然消除了对死亡的恐惧，他就有理由询问伯多禄的爱。恐惧曾三次否认，爱德三次宣认。三次的否认是真理的离弃；三次的宣认是爱德的见证。\par

\SourceRecord{Translated by R.G. MacMullen.；Nicene and Post-Nicene Fathers, First Series；Vol. 6.；Edited by Philip Schaff.；Buffalo, NY: Christian Literature Publishing Co.,；1888.；<http://www.newadvent.org/fathers/160397.htm>.}
